% $Header: /Users/joseph/Documents/LaTeX/beamer/solutions/generic-talks/generic-ornate-15min-45min.en.tex,v 90e850259b8b 2007/01/28 20:48:30 tantau $

\documentclass{beamer}
\usepackage{pgfpages}

\mode<presentation>
{
 % \usetheme{cambridge}
  \usetheme{Warsaw}

  \setbeamertemplate{navigation symbols}{}
  \setbeamercovered{transparent}
  % or whatever (possibly just delete it)
}

\usepackage[english]{babel}
\usepackage[latin1]{inputenc}
\usepackage{multicol}

%\usepackage{times}
%\usepackage[T1]{fontenc}
% Or whatever. Note that the encoding and the font should match. If T1
% does not look nice, try deleting the line with the fontenc.


\title[HPC: An introduction] % (optional, use only with long paper titles)
{An Introduction to High Performance Computing}

%\subtitle
%{Presentation Subtitle} % (optional)

\author[P Sumption] % (optional, use only with lots of authors)
{Paul Sumption\\ \texttt{sjr20@cam.ac.uk}}
%{F.~Author\inst{1} \and S.~Another\inst{2}}
% - Use the \inst{?} command only if the authors have different
%   affiliation.

\institute[Research Computing Services, University of Cambridge] % (optional, but mostly needed)
{Research Computing Services (http://www.hpc.cam.ac.uk/)\\
University Information Services (http://www.uis.cam.ac.uk/)}
% - Use the \inst command only if there are several affiliations.
% - Keep it simple, no one is interested in your street address.

\date[22/10/2021] % (optional)
{22nd October 2021 / UIS Training}

\subject{Courses}


\begin{document}

\begin{frame}
  \titlepage
\end{frame}

\begin{frame}{Health and Safety}
\begin{columns}[c]
\begin{column}{0.33\textwidth}
\begin{center}
\includegraphics[width=0.8\textwidth,height=0.5\textheight,keepaspectratio]{imgs/health-safety-1.png}\\
\includegraphics[width=0.8\textwidth,height=0.5\textheight,keepaspectratio]{imgs/health-safety-4.png}
\end{center}
\end{column}
\begin{column}{0.33\textwidth}
\begin{center}
\includegraphics[width=0.8\textwidth,height=0.5\textheight,keepaspectratio]{imgs/health-safety-2.png}\\
\includegraphics[width=0.8\textwidth,height=0.5\textheight,keepaspectratio]{imgs/health-safety-5.png}
\end{center}
\end{column}
\begin{column}{0.33\textwidth}
\begin{center}
\includegraphics[width=0.8\textwidth,height=0.5\textheight,keepaspectratio]{imgs/health-safety-3.png}\\
\includegraphics[width=0.8\textwidth,height=0.5\textheight,keepaspectratio]{imgs/health-safety-6.png}
\end{center}
\end{column}
\end{columns}
\end{frame}

\begin{frame}<presentation>{Welcome}
\begin{itemize}
\item{Please sign in on the {\color{red}attendance sheet}.}
\item Please give your {\color{red}online feedback} at the end of the course:
      \url{http://feedback.training.cam.ac.uk/ucs/form.php}
\item{Keep your belongings with you.}
%\item Course files can be downloaded from:  \url{www.ucs.cam.ac.uk/training/files}
\item\alert{Please ask questions and let us know if you need assistance.}
\end{itemize}
\end{frame}

\section{Who are we?}
\begin{frame}{UIS: Research Computing Services}
Your trainers for today will be:\\
\begin{itemize}
\item{\alert{Stuart Rankin}\\{}\qquad\qquad Research Computing User Services}
\item{\alert{Mark Sharpley}\\{}\qquad\qquad Research Computing Platforms}
\item{We are generalists, but there is also the \alert{Research Software Engineering} team.}
\end{itemize}
\end{frame}

\section{Who are you?}
\begin{frame}{You may be $\ldots$}
  \begin{itemize}
  \item{Programmers (or not).}\pause
  \item{UNIX power users (or not).}\pause
  \item{Researchers wishing to run large, parallel code.}\pause
  \item{Researchers wishing to run many, non-parallel cases.}\pause
  \item{Researchers interested in big data, machine learning, AI.}\pause
  \item{Researchers requiring slightly more than an ordinary workstation.}\pause
  \item{\alert{Many different disciplines and requirements.}}
\end{itemize}
\end{frame}


\begin{frame}<presentation>{Plan of the Course}
\begin{description}
\item[Part 1:]{Basics}
\item[Part 2:]{Research Computing Services HPC}
\item[Part 3:]{Using HPC}
\medskip
\item<2>{\alert{09:30}\hbox to 0pt{\hphantom{{}-15:15}\quad WELCOME\hss}}
\item<2>{\alert{11:00-11:15}\quad Break}
\item<2>{\alert{12:30-13:30}\quad LUNCH}
\item<2>{\alert{15:30-15:45}\quad Break}
\item<2>{\alert{16:30}\hbox to 0pt{\hphantom{{}-15:15}\quad CLOSE\hss}}
\end{description}
\end{frame}

\documentclass{article}
\usepackage{amsmath}
\usepackage{color,pxfonts,fix-cm}
\usepackage{latexsym}
\usepackage[mathletters]{ucs}
\DeclareUnicodeCharacter{8211}{\textendash}
\DeclareUnicodeCharacter{46}{\textperiodcentered}
\DeclareUnicodeCharacter{58}{$\colon$}
\DeclareUnicodeCharacter{8226}{$\bullet$}
\DeclareUnicodeCharacter{8217}{\textquoteright}
\DeclareUnicodeCharacter{62}{\textgreater}
\DeclareUnicodeCharacter{32}{$\ $}
\usepackage[T1]{fontenc}
\usepackage[utf8x]{inputenc}
\usepackage{pict2e}
\usepackage{wasysym}
\usepackage[english]{babel}
\usepackage{tikz}
\pagestyle{empty}
\usepackage[margin=0in,paperwidth=720pt,paperheight=540pt]{geometry}
\begin{document}
\definecolor{color_283006}{rgb}{1,1,1}
\definecolor{color_29791}{rgb}{0,0,0}
\definecolor{color_167934}{rgb}{0.545098,0.545098,1}
\definecolor{color_167818}{rgb}{0.545098,0.545098,0.545098}
\definecolor{color_165832}{rgb}{0.537255,0.537255,0.537255}
\begin{tikzpicture}[overlay]
\path(0pt,0pt);
\filldraw[color_283006][even odd rule]
(-15pt, 10pt) -- (704.972pt, 10pt)
 -- (704.972pt, 10pt)
 -- (704.972pt, -529.972pt)
 -- (704.972pt, -529.972pt)
 -- (-15pt, -529.972pt)
 -- (-15pt, -529.972pt)
 -- (-15pt, 10pt) -- cycle
;
\end{tikzpicture}
\begin{picture}(-5,0)(2.5,0)
\put(-15,-529.972){\includegraphics[width=720pt,height=540pt]{latexImage_340e7acc1534b9ba351a92c1a1081d49.png}}
\put(48.411,-192.904){\fontsize{28.006}{1}\usefont{T1}{cmr}{m}{n}\selectfont\color{color_29791}B}
\put(67.00698,-192.904){\fontsize{28.006}{1}\usefont{T1}{cmr}{m}{n}\selectfont\color{color_29791}r}
\put(76.38899,-192.904){\fontsize{28.006}{1}\usefont{T1}{cmr}{m}{n}\selectfont\color{color_29791}i}
\put(82.66234,-192.904){\fontsize{28.006}{1}\usefont{T1}{cmr}{m}{n}\selectfont\color{color_29791}t}
\put(90.36399,-192.904){\fontsize{28.006}{1}\usefont{T1}{cmr}{m}{n}\selectfont\color{color_29791}i}
\put(96.63734,-192.904){\fontsize{28.006}{1}\usefont{T1}{cmr}{m}{n}\selectfont\color{color_29791}sh }
\put(132.401,-192.904){\fontsize{28.006}{1}\usefont{T1}{cmr}{m}{n}\selectfont\color{color_29791}A}
\put(151.081,-192.904){\fontsize{28.006}{1}\usefont{T1}{cmr}{m}{n}\selectfont\color{color_29791}nt}
\put(174.438,-192.904){\fontsize{28.006}{1}\usefont{T1}{cmr}{m}{n}\selectfont\color{color_29791}ar}
\put(199.3073,-192.904){\fontsize{28.006}{1}\usefont{T1}{cmr}{m}{n}\selectfont\color{color_29791}c}
\put(213.3944,-192.904){\fontsize{28.006}{1}\usefont{T1}{cmr}{m}{n}\selectfont\color{color_29791}t}
\put(221.18,-192.904){\fontsize{28.006}{1}\usefont{T1}{cmr}{m}{n}\selectfont\color{color_29791}i}
\put(227.3694,-192.904){\fontsize{28.006}{1}\usefont{T1}{cmr}{m}{n}\selectfont\color{color_29791}c }
\put(249.1581,-192.904){\fontsize{28.006}{1}\usefont{T1}{cmr}{m}{n}\selectfont\color{color_29791}S}
\put(267.754,-192.904){\fontsize{28.006}{1}\usefont{T1}{cmr}{m}{n}\selectfont\color{color_29791}ur}
\put(292.7074,-192.904){\fontsize{28.006}{1}\usefont{T1}{cmr}{m}{n}\selectfont\color{color_29791}vey}
\put(48.411,-260.113){\fontsize{28.006}{1}\usefont{T1}{cmr}{m}{n}\selectfont\color{color_29791}H}
\put(68.60332,-260.113){\fontsize{28.006}{1}\usefont{T1}{cmr}{m}{n}\selectfont\color{color_29791}P}
\put(87.19931,-260.113){\fontsize{28.006}{1}\usefont{T1}{cmr}{m}{n}\selectfont\color{color_29791}C}
\put(107.3916,-260.113){\fontsize{28.006}{1}\usefont{T1}{cmr}{m}{n}\selectfont\color{color_29791} }
\put(115.1773,-260.113){\fontsize{28.006}{1}\usefont{T1}{cmr}{m}{n}\selectfont\color{color_29791}I}
\put(122.963,-260.113){\fontsize{28.006}{1}\usefont{T1}{cmr}{m}{n}\selectfont\color{color_29791}nduct}
\put(191.4657,-260.113){\fontsize{28.006}{1}\usefont{T1}{cmr}{m}{n}\selectfont\color{color_29791}i}
\put(197.739,-260.113){\fontsize{28.006}{1}\usefont{T1}{cmr}{m}{n}\selectfont\color{color_29791}o}
\put(213.2543,-260.113){\fontsize{28.006}{1}\usefont{T1}{cmr}{m}{n}\selectfont\color{color_29791}n }
\put(236.6113,-260.113){\fontsize{28.006}{1}\usefont{T1}{cmr}{m}{n}\selectfont\color{color_29791}2021}
\end{picture}
\newpage
\begin{tikzpicture}[overlay]
\path(0pt,0pt);
\filldraw[color_283006][even odd rule]
(-15pt, 10pt) -- (704.972pt, 10pt)
 -- (704.972pt, 10pt)
 -- (704.972pt, -529.972pt)
 -- (704.972pt, -529.972pt)
 -- (-15pt, -529.972pt)
 -- (-15pt, -529.972pt)
 -- (-15pt, 10pt) -- cycle
;
\end{tikzpicture}
\begin{picture}(-5,0)(2.5,0)
\put(-15,-529.972){\includegraphics[width=720pt,height=540pt]{latexImage_15b9e6acda8be87c69673355bc94d374.png}}
\put(48.411,-60.69601){\fontsize{21.997}{1}\usefont{T1}{cmr}{b}{n}\selectfont\color{color_29791}T}
\put(60.31138,-60.69601){\fontsize{21.997}{1}\usefont{T1}{cmr}{b}{n}\selectfont\color{color_29791}o}
\put(73.61956,-60.69601){\fontsize{21.997}{1}\usefont{T1}{cmr}{b}{n}\selectfont\color{color_29791}p}
\put(87.10372,-60.69601){\fontsize{21.997}{1}\usefont{T1}{cmr}{b}{n}\selectfont\color{color_29791}ic}
\put(105.4052,-60.69601){\fontsize{21.997}{1}\usefont{T1}{cmr}{b}{n}\selectfont\color{color_29791}s}
\put(117.6796,-60.69601){\fontsize{21.997}{1}\usefont{T1}{cmr}{b}{n}\selectfont\color{color_29791} C}
\put(139.5666,-60.69601){\fontsize{21.997}{1}\usefont{T1}{cmr}{b}{n}\selectfont\color{color_29791}o}
\put(153.0507,-60.69601){\fontsize{21.997}{1}\usefont{T1}{cmr}{b}{n}\selectfont\color{color_29791}v}
\put(165.259,-60.69601){\fontsize{21.997}{1}\usefont{T1}{cmr}{b}{n}\selectfont\color{color_29791}e}
\put(177.5334,-60.69601){\fontsize{21.997}{1}\usefont{T1}{cmr}{b}{n}\selectfont\color{color_29791}r}
\put(186.0242,-60.69601){\fontsize{21.997}{1}\usefont{T1}{cmr}{b}{n}\selectfont\color{color_29791}e}
\put(198.2325,-60.69601){\fontsize{21.997}{1}\usefont{T1}{cmr}{b}{n}\selectfont\color{color_29791}d}
\put(48.213,-119.288){\fontsize{20.013}{1}\usefont{T1}{cmr}{m}{n}\selectfont\color{color_29791}•}
\put(62.187,-120.507){\fontsize{20.013}{1}\usefont{T1}{cmr}{m}{n}\selectfont\color{color_29791}H}
\put(76.67641,-120.507){\fontsize{20.013}{1}\usefont{T1}{cmr}{m}{n}\selectfont\color{color_29791}a}
\put(87.86368,-120.507){\fontsize{20.013}{1}\usefont{T1}{cmr}{m}{n}\selectfont\color{color_29791}r}
\put(94.54802,-120.507){\fontsize{20.013}{1}\usefont{T1}{cmr}{m}{n}\selectfont\color{color_29791}d}
\put(105.7353,-120.507){\fontsize{20.013}{1}\usefont{T1}{cmr}{m}{n}\selectfont\color{color_29791}w}
\put(120.2247,-120.507){\fontsize{20.013}{1}\usefont{T1}{cmr}{m}{n}\selectfont\color{color_29791}a}
\put(131.412,-120.507){\fontsize{20.013}{1}\usefont{T1}{cmr}{m}{n}\selectfont\color{color_29791}r}
\put(138.0963,-120.507){\fontsize{20.013}{1}\usefont{T1}{cmr}{m}{n}\selectfont\color{color_29791}e}
\put(48.213,-159.313){\fontsize{20.013}{1}\usefont{T1}{cmr}{m}{n}\selectfont\color{color_29791}•}
\put(62.187,-160.504){\fontsize{20.013}{1}\usefont{T1}{cmr}{m}{n}\selectfont\color{color_29791}S}
\put(75.49564,-160.504){\fontsize{20.013}{1}\usefont{T1}{cmr}{m}{n}\selectfont\color{color_29791}t}
\put(81.07927,-160.504){\fontsize{20.013}{1}\usefont{T1}{cmr}{m}{n}\selectfont\color{color_29791}o}
\put(92.18649,-160.504){\fontsize{20.013}{1}\usefont{T1}{cmr}{m}{n}\selectfont\color{color_29791}r}
\put(98.87083,-160.504){\fontsize{20.013}{1}\usefont{T1}{cmr}{m}{n}\selectfont\color{color_29791}a}
\put(110.0581,-160.504){\fontsize{20.013}{1}\usefont{T1}{cmr}{m}{n}\selectfont\color{color_29791}g}
\put(121.2454,-160.504){\fontsize{20.013}{1}\usefont{T1}{cmr}{m}{n}\selectfont\color{color_29791}e}
\put(48.213,-199.31){\fontsize{20.013}{1}\usefont{T1}{cmr}{m}{n}\selectfont\color{color_29791}•}
\put(62.187,-200.501){\fontsize{20.013}{1}\usefont{T1}{cmr}{m}{n}\selectfont\color{color_29791}A}
\put(75.49564,-200.501){\fontsize{20.013}{1}\usefont{T1}{cmr}{m}{n}\selectfont\color{color_29791}c}
\put(85.5822,-200.501){\fontsize{20.013}{1}\usefont{T1}{cmr}{m}{n}\selectfont\color{color_29791}c}
\put(95.66875,-200.501){\fontsize{20.013}{1}\usefont{T1}{cmr}{m}{n}\selectfont\color{color_29791}e}
\put(106.856,-200.501){\fontsize{20.013}{1}\usefont{T1}{cmr}{m}{n}\selectfont\color{color_29791}ss}
\put(48.213,-239.307){\fontsize{20.013}{1}\usefont{T1}{cmr}{m}{n}\selectfont\color{color_29791}•}
\put(62.187,-240.498){\fontsize{20.013}{1}\usefont{T1}{cmr}{m}{n}\selectfont\color{color_29791}D}
\put(76.67641,-240.498){\fontsize{20.013}{1}\usefont{T1}{cmr}{m}{n}\selectfont\color{color_29791}a}
\put(87.86368,-240.498){\fontsize{20.013}{1}\usefont{T1}{cmr}{m}{n}\selectfont\color{color_29791}t}
\put(93.36726,-240.498){\fontsize{20.013}{1}\usefont{T1}{cmr}{m}{n}\selectfont\color{color_29791}a}
\put(104.5545,-240.498){\fontsize{20.013}{1}\usefont{T1}{cmr}{m}{n}\selectfont\color{color_29791} }
\put(109.7579,-240.498){\fontsize{20.013}{1}\usefont{T1}{cmr}{m}{n}\selectfont\color{color_29791}T}
\put(121.1653,-240.498){\fontsize{20.013}{1}\usefont{T1}{cmr}{m}{n}\selectfont\color{color_29791}r}
\put(127.9497,-240.498){\fontsize{20.013}{1}\usefont{T1}{cmr}{m}{n}\selectfont\color{color_29791}a}
\put(139.0569,-240.498){\fontsize{20.013}{1}\usefont{T1}{cmr}{m}{n}\selectfont\color{color_29791}n}
\put(150.2442,-240.498){\fontsize{20.013}{1}\usefont{T1}{cmr}{m}{n}\selectfont\color{color_29791}s}
\put(160.3307,-240.498){\fontsize{20.013}{1}\usefont{T1}{cmr}{m}{n}\selectfont\color{color_29791}f}
\put(165.8343,-240.498){\fontsize{20.013}{1}\usefont{T1}{cmr}{m}{n}\selectfont\color{color_29791}e}
\put(176.9415,-240.498){\fontsize{20.013}{1}\usefont{T1}{cmr}{m}{n}\selectfont\color{color_29791}r}
\put(48.213,-279.304){\fontsize{20.013}{1}\usefont{T1}{cmr}{m}{n}\selectfont\color{color_29791}•}
\put(62.187,-280.494){\fontsize{20.013}{1}\usefont{T1}{cmr}{m}{n}\selectfont\color{color_29791}U}
\put(76.67641,-280.494){\fontsize{20.013}{1}\usefont{T1}{cmr}{m}{n}\selectfont\color{color_29791}s}
\put(86.76295,-280.494){\fontsize{20.013}{1}\usefont{T1}{cmr}{m}{n}\selectfont\color{color_29791}e}
\put(97.95023,-280.494){\fontsize{20.013}{1}\usefont{T1}{cmr}{m}{n}\selectfont\color{color_29791}r}
\put(104.6346,-280.494){\fontsize{20.013}{1}\usefont{T1}{cmr}{m}{n}\selectfont\color{color_29791} }
\put(110.1381,-280.494){\fontsize{20.013}{1}\usefont{T1}{cmr}{m}{n}\selectfont\color{color_29791}E}
\put(123.4468,-280.494){\fontsize{20.013}{1}\usefont{T1}{cmr}{m}{n}\selectfont\color{color_29791}n}
\put(134.6341,-280.494){\fontsize{20.013}{1}\usefont{T1}{cmr}{m}{n}\selectfont\color{color_29791}v}
\put(144.7206,-280.494){\fontsize{20.013}{1}\usefont{T1}{cmr}{m}{n}\selectfont\color{color_29791}i}
\put(149.1235,-280.494){\fontsize{20.013}{1}\usefont{T1}{cmr}{m}{n}\selectfont\color{color_29791}r}
\put(155.8078,-280.494){\fontsize{20.013}{1}\usefont{T1}{cmr}{m}{n}\selectfont\color{color_29791}o}
\put(166.9951,-280.494){\fontsize{20.013}{1}\usefont{T1}{cmr}{m}{n}\selectfont\color{color_29791}n}
\put(178.1824,-280.494){\fontsize{20.013}{1}\usefont{T1}{cmr}{m}{n}\selectfont\color{color_29791}m}
\put(194.7932,-280.494){\fontsize{20.013}{1}\usefont{T1}{cmr}{m}{n}\selectfont\color{color_29791}e}
\put(205.9804,-280.494){\fontsize{20.013}{1}\usefont{T1}{cmr}{m}{n}\selectfont\color{color_29791}n}
\put(217.0876,-280.494){\fontsize{20.013}{1}\usefont{T1}{cmr}{m}{n}\selectfont\color{color_29791}t}
\put(48.213,-319.301){\fontsize{20.013}{1}\usefont{T1}{cmr}{m}{n}\selectfont\color{color_29791}•}
\put(62.187,-320.491){\fontsize{20.013}{1}\usefont{T1}{cmr}{m}{n}\selectfont\color{color_29791}S}
\put(75.49564,-320.491){\fontsize{20.013}{1}\usefont{T1}{cmr}{m}{n}\selectfont\color{color_29791}o}
\put(86.68291,-320.491){\fontsize{20.013}{1}\usefont{T1}{cmr}{m}{n}\selectfont\color{color_29791}f}
\put(92.18649,-320.491){\fontsize{20.013}{1}\usefont{T1}{cmr}{m}{n}\selectfont\color{color_29791}t}
\put(97.69006,-320.491){\fontsize{20.013}{1}\usefont{T1}{cmr}{m}{n}\selectfont\color{color_29791}w}
\put(112.1795,-320.491){\fontsize{20.013}{1}\usefont{T1}{cmr}{m}{n}\selectfont\color{color_29791}a}
\put(123.3667,-320.491){\fontsize{20.013}{1}\usefont{T1}{cmr}{m}{n}\selectfont\color{color_29791}r}
\put(130.0511,-320.491){\fontsize{20.013}{1}\usefont{T1}{cmr}{m}{n}\selectfont\color{color_29791}e}
\put(349.592,-119.487){\fontsize{20.013}{1}\usefont{T1}{cmr}{m}{n}\selectfont\color{color_29791}•}
\put(363.595,-120.592){\fontsize{20.013}{1}\usefont{T1}{cmr}{m}{n}\selectfont\color{color_29791}C}
\put(378.0844,-120.592){\fontsize{20.013}{1}\usefont{T1}{cmr}{m}{n}\selectfont\color{color_29791}o}
\put(389.2717,-120.592){\fontsize{20.013}{1}\usefont{T1}{cmr}{m}{n}\selectfont\color{color_29791}n}
\put(400.4589,-120.592){\fontsize{20.013}{1}\usefont{T1}{cmr}{m}{n}\selectfont\color{color_29791}t}
\put(405.9625,-120.592){\fontsize{20.013}{1}\usefont{T1}{cmr}{m}{n}\selectfont\color{color_29791}a}
\put(417.0697,-120.592){\fontsize{20.013}{1}\usefont{T1}{cmr}{m}{n}\selectfont\color{color_29791}i}
\put(421.5526,-120.592){\fontsize{20.013}{1}\usefont{T1}{cmr}{m}{n}\selectfont\color{color_29791}n}
\put(432.6598,-120.592){\fontsize{20.013}{1}\usefont{T1}{cmr}{m}{n}\selectfont\color{color_29791}e}
\put(443.8471,-120.592){\fontsize{20.013}{1}\usefont{T1}{cmr}{m}{n}\selectfont\color{color_29791}r}
\put(450.5314,-120.592){\fontsize{20.013}{1}\usefont{T1}{cmr}{m}{n}\selectfont\color{color_29791}s}
\put(349.592,-159.512){\fontsize{20.013}{1}\usefont{T1}{cmr}{m}{n}\selectfont\color{color_29791}•}
\put(363.595,-160.589){\fontsize{20.013}{1}\usefont{T1}{cmr}{m}{n}\selectfont\color{color_29791}S}
\put(376.9037,-160.589){\fontsize{20.013}{1}\usefont{T1}{cmr}{m}{n}\selectfont\color{color_29791}L}
\put(388.0909,-160.589){\fontsize{20.013}{1}\usefont{T1}{cmr}{m}{n}\selectfont\color{color_29791}U}
\put(402.5803,-160.589){\fontsize{20.013}{1}\usefont{T1}{cmr}{m}{n}\selectfont\color{color_29791}R}
\put(417.0697,-160.589){\fontsize{20.013}{1}\usefont{T1}{cmr}{m}{n}\selectfont\color{color_29791}M}
\put(349.592,-199.509){\fontsize{20.013}{1}\usefont{T1}{cmr}{m}{n}\selectfont\color{color_29791}•}
\put(363.595,-200.614){\fontsize{20.013}{1}\usefont{T1}{cmr}{m}{n}\selectfont\color{color_29791}M}
\put(380.2858,-200.614){\fontsize{20.013}{1}\usefont{T1}{cmr}{m}{n}\selectfont\color{color_29791}o}
\put(391.4731,-200.614){\fontsize{20.013}{1}\usefont{T1}{cmr}{m}{n}\selectfont\color{color_29791}d}
\put(402.5803,-200.614){\fontsize{20.013}{1}\usefont{T1}{cmr}{m}{n}\selectfont\color{color_29791}e}
\put(413.7675,-200.614){\fontsize{20.013}{1}\usefont{T1}{cmr}{m}{n}\selectfont\color{color_29791}l}
\put(418.1704,-200.614){\fontsize{20.013}{1}\usefont{T1}{cmr}{m}{n}\selectfont\color{color_29791} }
\put(423.674,-200.614){\fontsize{20.013}{1}\usefont{T1}{cmr}{m}{n}\selectfont\color{color_29791}E}
\put(437.0627,-200.614){\fontsize{20.013}{1}\usefont{T1}{cmr}{m}{n}\selectfont\color{color_29791}n}
\put(448.1699,-200.614){\fontsize{20.013}{1}\usefont{T1}{cmr}{m}{n}\selectfont\color{color_29791}s}
\put(458.2564,-200.614){\fontsize{20.013}{1}\usefont{T1}{cmr}{m}{n}\selectfont\color{color_29791}e}
\put(469.4437,-200.614){\fontsize{20.013}{1}\usefont{T1}{cmr}{m}{n}\selectfont\color{color_29791}m}
\put(486.0545,-200.614){\fontsize{20.013}{1}\usefont{T1}{cmr}{m}{n}\selectfont\color{color_29791}b}
\put(497.2418,-200.614){\fontsize{20.013}{1}\usefont{T1}{cmr}{m}{n}\selectfont\color{color_29791}l}
\put(501.7247,-200.614){\fontsize{20.013}{1}\usefont{T1}{cmr}{m}{n}\selectfont\color{color_29791}e}
\put(512.8319,-200.614){\fontsize{20.013}{1}\usefont{T1}{cmr}{m}{n}\selectfont\color{color_29791}r}
\put(349.592,-239.506){\fontsize{20.013}{1}\usefont{T1}{cmr}{m}{n}\selectfont\color{color_29791}•}
\put(363.595,-240.611){\fontsize{20.013}{1}\usefont{T1}{cmr}{m}{n}\selectfont\color{color_29791}B}
\put(376.9037,-240.611){\fontsize{20.013}{1}\usefont{T1}{cmr}{m}{n}\selectfont\color{color_29791}e}
\put(388.0909,-240.611){\fontsize{20.013}{1}\usefont{T1}{cmr}{m}{n}\selectfont\color{color_29791}s}
\put(398.1775,-240.611){\fontsize{20.013}{1}\usefont{T1}{cmr}{m}{n}\selectfont\color{color_29791}t}
\put(403.681,-240.611){\fontsize{20.013}{1}\usefont{T1}{cmr}{m}{n}\selectfont\color{color_29791} }
\put(409.1846,-240.611){\fontsize{20.013}{1}\usefont{T1}{cmr}{m}{n}\selectfont\color{color_29791}P}
\put(422.5733,-240.611){\fontsize{20.013}{1}\usefont{T1}{cmr}{m}{n}\selectfont\color{color_29791}r}
\put(429.2577,-240.611){\fontsize{20.013}{1}\usefont{T1}{cmr}{m}{n}\selectfont\color{color_29791}a}
\put(440.3649,-240.611){\fontsize{20.013}{1}\usefont{T1}{cmr}{m}{n}\selectfont\color{color_29791}c}
\put(450.4514,-240.611){\fontsize{20.013}{1}\usefont{T1}{cmr}{m}{n}\selectfont\color{color_29791}t}
\put(455.955,-240.611){\fontsize{20.013}{1}\usefont{T1}{cmr}{m}{n}\selectfont\color{color_29791}i}
\put(460.4379,-240.611){\fontsize{20.013}{1}\usefont{T1}{cmr}{m}{n}\selectfont\color{color_29791}ce}
\put(349.592,-279.502){\fontsize{20.013}{1}\usefont{T1}{cmr}{m}{n}\selectfont\color{color_29791}•}
\put(363.595,-280.608){\fontsize{20.013}{1}\usefont{T1}{cmr}{m}{n}\selectfont\color{color_29791}H}
\put(378.0844,-280.608){\fontsize{20.013}{1}\usefont{T1}{cmr}{m}{n}\selectfont\color{color_29791}E}
\put(391.4731,-280.608){\fontsize{20.013}{1}\usefont{T1}{cmr}{m}{n}\selectfont\color{color_29791}L}
\put(402.5803,-280.608){\fontsize{20.013}{1}\usefont{T1}{cmr}{m}{n}\selectfont\color{color_29791}P}
\put(415.969,-280.608){\fontsize{20.013}{1}\usefont{T1}{cmr}{m}{n}\selectfont\color{color_29791}!}
\end{picture}
\newpage
\begin{tikzpicture}[overlay]
\path(0pt,0pt);
\filldraw[color_283006][even odd rule]
(-15pt, 10pt) -- (704.972pt, 10pt)
 -- (704.972pt, 10pt)
 -- (704.972pt, -529.972pt)
 -- (704.972pt, -529.972pt)
 -- (-15pt, -529.972pt)
 -- (-15pt, -529.972pt)
 -- (-15pt, 10pt) -- cycle
;
\end{tikzpicture}
\begin{picture}(-5,0)(2.5,0)
\put(-15,-529.972){\includegraphics[width=720pt,height=540pt]{latexImage_1aef5d9a345e7259c360833b5fc6008f.png}}
\put(48.411,-60.69601){\fontsize{21.997}{1}\usefont{T1}{cmr}{b}{n}\selectfont\color{color_29791}Ha}
\put(76.50117,-60.69601){\fontsize{21.997}{1}\usefont{T1}{cmr}{b}{n}\selectfont\color{color_29791}r}
\put(84.99201,-60.69601){\fontsize{21.997}{1}\usefont{T1}{cmr}{b}{n}\selectfont\color{color_29791}d}
\put(98.47617,-60.69601){\fontsize{21.997}{1}\usefont{T1}{cmr}{b}{n}\selectfont\color{color_29791}w}
\put(115.4799,-60.69601){\fontsize{21.997}{1}\usefont{T1}{cmr}{b}{n}\selectfont\color{color_29791}a}
\put(127.7542,-60.69601){\fontsize{21.997}{1}\usefont{T1}{cmr}{b}{n}\selectfont\color{color_29791}r}
\put(136.245,-60.69601){\fontsize{21.997}{1}\usefont{T1}{cmr}{b}{n}\selectfont\color{color_29791}e}
\put(48.213,-115.49){\fontsize{15.987}{1}\usefont{T1}{cmr}{m}{n}\selectfont\color{color_29791}•}
\put(62.187,-116.51){\fontsize{15.987}{1}\usefont{T1}{cmr}{m}{n}\selectfont\color{color_29791}G}
\put(74.57692,-116.51){\fontsize{15.987}{1}\usefont{T1}{cmr}{m}{n}\selectfont\color{color_29791}a}
\put(83.36977,-116.51){\fontsize{15.987}{1}\usefont{T1}{cmr}{m}{n}\selectfont\color{color_29791}t}
\put(87.84612,-116.51){\fontsize{15.987}{1}\usefont{T1}{cmr}{m}{n}\selectfont\color{color_29791}ew}
\put(108.2295,-116.51){\fontsize{15.987}{1}\usefont{T1}{cmr}{m}{n}\selectfont\color{color_29791}ay }
\put(129.5882,-116.51){\fontsize{15.987}{1}\usefont{T1}{cmr}{m}{n}\selectfont\color{color_29791}o}
\put(138.381,-116.51){\fontsize{15.987}{1}\usefont{T1}{cmr}{m}{n}\selectfont\color{color_29791}r}
\put(143.6727,-116.51){\fontsize{15.987}{1}\usefont{T1}{cmr}{m}{n}\selectfont\color{color_29791} }
\put(148.0692,-116.51){\fontsize{15.987}{1}\usefont{T1}{cmr}{m}{n}\selectfont\color{color_29791}B}
\put(158.7485,-116.51){\fontsize{15.987}{1}\usefont{T1}{cmr}{m}{n}\selectfont\color{color_29791}as}
\put(175.7107,-116.51){\fontsize{15.987}{1}\usefont{T1}{cmr}{m}{n}\selectfont\color{color_29791}t}
\put(180.1071,-116.51){\fontsize{15.987}{1}\usefont{T1}{cmr}{m}{n}\selectfont\color{color_29791}i}
\put(183.6882,-116.51){\fontsize{15.987}{1}\usefont{T1}{cmr}{m}{n}\selectfont\color{color_29791}on }
\put(205.8622,-116.51){\fontsize{15.987}{1}\usefont{T1}{cmr}{m}{n}\selectfont\color{color_29791}ho}
\put(223.5438,-116.51){\fontsize{15.987}{1}\usefont{T1}{cmr}{m}{n}\selectfont\color{color_29791}s}
\put(231.6172,-116.51){\fontsize{15.987}{1}\usefont{T1}{cmr}{m}{n}\selectfont\color{color_29791}t}
\put(236.0136,-116.51){\fontsize{15.987}{1}\usefont{T1}{cmr}{m}{n}\selectfont\color{color_29791}s}
\put(244.0871,-116.51){\fontsize{15.987}{1}\usefont{T1}{cmr}{m}{n}\selectfont\color{color_29791} }
\put(248.4835,-116.51){\fontsize{15.987}{1}\usefont{T1}{cmr}{m}{n}\selectfont\color{color_29791}(}
\put(254.093,-116.51){\fontsize{11.991}{1}\usefont{T1}{cmr}{m}{n}\selectfont\color{color_29791}b}
\put(260.772,-116.51){\fontsize{11.991}{1}\usefont{T1}{cmr}{m}{n}\selectfont\color{color_29791}sl}
\put(269.4535,-116.51){\fontsize{11.991}{1}\usefont{T1}{cmr}{m}{n}\selectfont\color{color_29791}ce}
\put(282.128,-116.51){\fontsize{11.991}{1}\usefont{T1}{cmr}{m}{n}\selectfont\color{color_29791}n}
\put(288.8069,-116.51){\fontsize{11.991}{1}\usefont{T1}{cmr}{m}{n}\selectfont\color{color_29791}b}
\put(295.4859,-116.51){\fontsize{11.991}{1}\usefont{T1}{cmr}{m}{n}\selectfont\color{color_29791} }
\put(298.8674,-116.51){\fontsize{11.991}{1}\usefont{T1}{cmr}{m}{n}\selectfont\color{color_29791}\& }
\put(310.2349,-116.51){\fontsize{11.991}{1}\usefont{T1}{cmr}{m}{n}\selectfont\color{color_29791}b}
\put(316.9138,-116.51){\fontsize{11.991}{1}\usefont{T1}{cmr}{m}{n}\selectfont\color{color_29791}sl}
\put(325.5114,-116.51){\fontsize{11.991}{1}\usefont{T1}{cmr}{m}{n}\selectfont\color{color_29791}ce}
\put(338.2938,-116.51){\fontsize{11.991}{1}\usefont{T1}{cmr}{m}{n}\selectfont\color{color_29791}n}
\put(344.9728,-116.51){\fontsize{11.991}{1}\usefont{T1}{cmr}{m}{n}\selectfont\color{color_29791}c)}
\put(75.709,-145.31){\fontsize{11.991}{1}\usefont{T1}{cmr}{m}{n}\selectfont\color{color_29791}•}
\put(98.187,-145.991){\fontsize{11.991}{1}\usefont{T1}{cmr}{m}{n}\selectfont\color{color_29791}O}
\put(107.564,-145.991){\fontsize{11.991}{1}\usefont{T1}{cmr}{m}{n}\selectfont\color{color_29791}n}
\put(114.243,-145.991){\fontsize{11.991}{1}\usefont{T1}{cmr}{m}{n}\selectfont\color{color_29791}l}
\put(116.845,-145.991){\fontsize{11.991}{1}\usefont{T1}{cmr}{m}{n}\selectfont\color{color_29791}y }
\put(126.222,-145.991){\fontsize{11.991}{1}\usefont{T1}{cmr}{m}{n}\selectfont\color{color_29791}u}
\put(132.901,-145.991){\fontsize{11.991}{1}\usefont{T1}{cmr}{m}{n}\selectfont\color{color_29791}se}
\put(145.5754,-145.991){\fontsize{11.991}{1}\usefont{T1}{cmr}{m}{n}\selectfont\color{color_29791} }
\put(148.9569,-145.991){\fontsize{11.991}{1}\usefont{T1}{cmr}{m}{n}\selectfont\color{color_29791}f}
\put(152.2544,-145.991){\fontsize{11.991}{1}\usefont{T1}{cmr}{m}{n}\selectfont\color{color_29791}o}
\put(159.0414,-145.991){\fontsize{11.991}{1}\usefont{T1}{cmr}{m}{n}\selectfont\color{color_29791}r}
\put(162.9384,-145.991){\fontsize{11.991}{1}\usefont{T1}{cmr}{m}{n}\selectfont\color{color_29791} }
\put(166.3199,-145.991){\fontsize{11.991}{1}\usefont{T1}{cmr}{m}{n}\selectfont\color{color_29791}a}
\put(172.9989,-145.991){\fontsize{11.991}{1}\usefont{T1}{cmr}{m}{n}\selectfont\color{color_29791}cce}
\put(191.6689,-145.991){\fontsize{11.991}{1}\usefont{T1}{cmr}{m}{n}\selectfont\color{color_29791}ss }
\put(207.0413,-145.991){\fontsize{11.991}{1}\usefont{T1}{cmr}{m}{n}\selectfont\color{color_29791}t}
\put(210.4228,-145.991){\fontsize{11.991}{1}\usefont{T1}{cmr}{m}{n}\selectfont\color{color_29791}o}
\put(217.1018,-145.991){\fontsize{11.991}{1}\usefont{T1}{cmr}{m}{n}\selectfont\color{color_29791} }
\put(220.3993,-145.991){\fontsize{11.991}{1}\usefont{T1}{cmr}{m}{n}\selectfont\color{color_29791}B}
\put(228.4812,-145.991){\fontsize{11.991}{1}\usefont{T1}{cmr}{m}{n}\selectfont\color{color_29791}AS}
\put(244.5492,-145.991){\fontsize{11.991}{1}\usefont{T1}{cmr}{m}{n}\selectfont\color{color_29791} }
\put(247.8467,-145.991){\fontsize{11.991}{1}\usefont{T1}{cmr}{m}{n}\selectfont\color{color_29791}o}
\put(254.5257,-145.991){\fontsize{11.991}{1}\usefont{T1}{cmr}{m}{n}\selectfont\color{color_29791}r }
\put(261.9001,-145.991){\fontsize{11.991}{1}\usefont{T1}{cmr}{m}{n}\selectfont\color{color_29791}t}
\put(265.1977,-145.991){\fontsize{11.991}{1}\usefont{T1}{cmr}{m}{n}\selectfont\color{color_29791}ra}
\put(275.8697,-145.991){\fontsize{11.991}{1}\usefont{T1}{cmr}{m}{n}\selectfont\color{color_29791}n}
\put(282.5486,-145.991){\fontsize{11.991}{1}\usefont{T1}{cmr}{m}{n}\selectfont\color{color_29791}sf}
\put(291.9256,-145.991){\fontsize{11.991}{1}\usefont{T1}{cmr}{m}{n}\selectfont\color{color_29791}e}
\put(298.6046,-145.991){\fontsize{11.991}{1}\usefont{T1}{cmr}{m}{n}\selectfont\color{color_29791}r}
\put(302.5016,-145.991){\fontsize{11.991}{1}\usefont{T1}{cmr}{m}{n}\selectfont\color{color_29791}ri}
\put(309.1806,-145.991){\fontsize{11.991}{1}\usefont{T1}{cmr}{m}{n}\selectfont\color{color_29791}n}
\put(315.8596,-145.991){\fontsize{11.991}{1}\usefont{T1}{cmr}{m}{n}\selectfont\color{color_29791}g}
\put(322.5386,-145.991){\fontsize{11.991}{1}\usefont{T1}{cmr}{m}{n}\selectfont\color{color_29791} }
\put(325.9201,-145.991){\fontsize{11.991}{1}\usefont{T1}{cmr}{m}{n}\selectfont\color{color_29791}f}
\put(329.2176,-145.991){\fontsize{11.991}{1}\usefont{T1}{cmr}{m}{n}\selectfont\color{color_29791}i}
\put(331.8196,-145.991){\fontsize{11.991}{1}\usefont{T1}{cmr}{m}{n}\selectfont\color{color_29791}l}
\put(334.5056,-145.991){\fontsize{11.991}{1}\usefont{T1}{cmr}{m}{n}\selectfont\color{color_29791}e}
\put(341.1846,-145.991){\fontsize{11.991}{1}\usefont{T1}{cmr}{m}{n}\selectfont\color{color_29791}s,}
\put(350.5616,-145.991){\fontsize{11.991}{1}\usefont{T1}{cmr}{m}{n}\selectfont\color{color_29791} }
\put(353.8591,-145.991){\fontsize{11.991}{1}\usefont{T1}{cmr}{m}{n}\selectfont\color{color_29791}d}
\put(360.5381,-145.991){\fontsize{11.991}{1}\usefont{T1}{cmr}{m}{n}\selectfont\color{color_29791}o}
\put(367.325,-145.991){\fontsize{11.991}{1}\usefont{T1}{cmr}{m}{n}\selectfont\color{color_29791}n}
\put(374.0039,-145.991){\fontsize{11.991}{1}\usefont{T1}{cmr}{m}{n}\selectfont\color{color_29791}’}
\put(376.606,-145.991){\fontsize{11.991}{1}\usefont{T1}{cmr}{m}{n}\selectfont\color{color_29791}t}
\put(379.9875,-145.991){\fontsize{11.991}{1}\usefont{T1}{cmr}{m}{n}\selectfont\color{color_29791} }
\put(383.3689,-145.991){\fontsize{11.991}{1}\usefont{T1}{cmr}{m}{n}\selectfont\color{color_29791}u}
\put(390.0479,-145.991){\fontsize{11.991}{1}\usefont{T1}{cmr}{m}{n}\selectfont\color{color_29791}se}
\put(402.7224,-145.991){\fontsize{11.991}{1}\usefont{T1}{cmr}{m}{n}\selectfont\color{color_29791} }
\put(406.1039,-145.991){\fontsize{11.991}{1}\usefont{T1}{cmr}{m}{n}\selectfont\color{color_29791}f}
\put(409.4014,-145.991){\fontsize{11.991}{1}\usefont{T1}{cmr}{m}{n}\selectfont\color{color_29791}o}
\put(416.0804,-145.991){\fontsize{11.991}{1}\usefont{T1}{cmr}{m}{n}\selectfont\color{color_29791}r }
\put(423.4548,-145.991){\fontsize{11.991}{1}\usefont{T1}{cmr}{m}{n}\selectfont\color{color_29791}r}
\put(427.3519,-145.991){\fontsize{11.991}{1}\usefont{T1}{cmr}{m}{n}\selectfont\color{color_29791}u}
\put(434.0309,-145.991){\fontsize{11.991}{1}\usefont{T1}{cmr}{m}{n}\selectfont\color{color_29791}n}
\put(440.8178,-145.991){\fontsize{11.991}{1}\usefont{T1}{cmr}{m}{n}\selectfont\color{color_29791}n}
\put(447.4968,-145.991){\fontsize{11.991}{1}\usefont{T1}{cmr}{m}{n}\selectfont\color{color_29791}i}
\put(450.0988,-145.991){\fontsize{11.991}{1}\usefont{T1}{cmr}{m}{n}\selectfont\color{color_29791}n}
\put(456.7778,-145.991){\fontsize{11.991}{1}\usefont{T1}{cmr}{m}{n}\selectfont\color{color_29791}g}
\put(463.5647,-145.991){\fontsize{11.991}{1}\usefont{T1}{cmr}{m}{n}\selectfont\color{color_29791} }
\put(466.8622,-145.991){\fontsize{11.991}{1}\usefont{T1}{cmr}{m}{n}\selectfont\color{color_29791}p}
\put(473.5412,-145.991){\fontsize{11.991}{1}\usefont{T1}{cmr}{m}{n}\selectfont\color{color_29791}ro}
\put(484.2132,-145.991){\fontsize{11.991}{1}\usefont{T1}{cmr}{m}{n}\selectfont\color{color_29791}g}
\put(490.8922,-145.991){\fontsize{11.991}{1}\usefont{T1}{cmr}{m}{n}\selectfont\color{color_29791}ra}
\put(501.5641,-145.991){\fontsize{11.991}{1}\usefont{T1}{cmr}{m}{n}\selectfont\color{color_29791}ms}
\put(518.594,-145.991){\fontsize{11.991}{1}\usefont{T1}{cmr}{m}{n}\selectfont\color{color_29791} }
\put(48.213,-178.589){\fontsize{15.987}{1}\usefont{T1}{cmr}{m}{n}\selectfont\color{color_29791}•}
\put(62.187,-179.609){\fontsize{15.987}{1}\usefont{T1}{cmr}{m}{n}\selectfont\color{color_29791}H}
\put(73.68165,-179.609){\fontsize{15.987}{1}\usefont{T1}{cmr}{m}{n}\selectfont\color{color_29791}eadn}
\put(109.1408,-179.609){\fontsize{15.987}{1}\usefont{T1}{cmr}{m}{n}\selectfont\color{color_29791}odes}
\put(75.709,-208.494){\fontsize{11.991}{1}\usefont{T1}{cmr}{m}{n}\selectfont\color{color_29791}•}
\put(98.187,-209.09){\fontsize{11.991}{1}\usefont{T1}{cmr}{m}{n}\selectfont\color{color_29791}N}
\put(106.8685,-209.09){\fontsize{11.991}{1}\usefont{T1}{cmr}{m}{n}\selectfont\color{color_29791}o}
\put(113.5475,-209.09){\fontsize{11.991}{1}\usefont{T1}{cmr}{m}{n}\selectfont\color{color_29791} }
\put(116.845,-209.09){\fontsize{11.991}{1}\usefont{T1}{cmr}{m}{n}\selectfont\color{color_29791}a}
\put(123.524,-209.09){\fontsize{11.991}{1}\usefont{T1}{cmr}{m}{n}\selectfont\color{color_29791}c}
\put(129.6034,-209.09){\fontsize{11.991}{1}\usefont{T1}{cmr}{m}{n}\selectfont\color{color_29791}ce}
\put(142.2779,-209.09){\fontsize{11.991}{1}\usefont{T1}{cmr}{m}{n}\selectfont\color{color_29791}ss,}
\put(157.5664,-209.09){\fontsize{11.991}{1}\usefont{T1}{cmr}{m}{n}\selectfont\color{color_29791} }
\put(160.9479,-209.09){\fontsize{11.991}{1}\usefont{T1}{cmr}{m}{n}\selectfont\color{color_29791}ma}
\put(177.6154,-209.09){\fontsize{11.991}{1}\usefont{T1}{cmr}{m}{n}\selectfont\color{color_29791}n}
\put(184.2944,-209.09){\fontsize{11.991}{1}\usefont{T1}{cmr}{m}{n}\selectfont\color{color_29791}a}
\put(190.9734,-209.09){\fontsize{11.991}{1}\usefont{T1}{cmr}{m}{n}\selectfont\color{color_29791}g}
\put(197.6524,-209.09){\fontsize{11.991}{1}\usefont{T1}{cmr}{m}{n}\selectfont\color{color_29791}e}
\put(204.4393,-209.09){\fontsize{11.991}{1}\usefont{T1}{cmr}{m}{n}\selectfont\color{color_29791}s }
\put(213.7323,-209.09){\fontsize{11.991}{1}\usefont{T1}{cmr}{m}{n}\selectfont\color{color_29791}j}
\put(216.3344,-209.09){\fontsize{11.991}{1}\usefont{T1}{cmr}{m}{n}\selectfont\color{color_29791}o}
\put(223.1213,-209.09){\fontsize{11.991}{1}\usefont{T1}{cmr}{m}{n}\selectfont\color{color_29791}b}
\put(229.8002,-209.09){\fontsize{11.991}{1}\usefont{T1}{cmr}{m}{n}\selectfont\color{color_29791} }
\put(233.1817,-209.09){\fontsize{11.991}{1}\usefont{T1}{cmr}{m}{n}\selectfont\color{color_29791}q}
\put(239.8607,-209.09){\fontsize{11.991}{1}\usefont{T1}{cmr}{m}{n}\selectfont\color{color_29791}u}
\put(246.5397,-209.09){\fontsize{11.991}{1}\usefont{T1}{cmr}{m}{n}\selectfont\color{color_29791}e}
\put(253.2187,-209.09){\fontsize{11.991}{1}\usefont{T1}{cmr}{m}{n}\selectfont\color{color_29791}u}
\put(259.8976,-209.09){\fontsize{11.991}{1}\usefont{T1}{cmr}{m}{n}\selectfont\color{color_29791}e}
\put(266.5766,-209.09){\fontsize{11.991}{1}\usefont{T1}{cmr}{m}{n}\selectfont\color{color_29791}s}
\put(272.6561,-209.09){\fontsize{11.991}{1}\usefont{T1}{cmr}{m}{n}\selectfont\color{color_29791} }
\put(275.9536,-209.09){\fontsize{11.991}{1}\usefont{T1}{cmr}{m}{n}\selectfont\color{color_29791}a}
\put(282.6326,-209.09){\fontsize{11.991}{1}\usefont{T1}{cmr}{m}{n}\selectfont\color{color_29791}n}
\put(289.4195,-209.09){\fontsize{11.991}{1}\usefont{T1}{cmr}{m}{n}\selectfont\color{color_29791}d}
\put(296.0984,-209.09){\fontsize{11.991}{1}\usefont{T1}{cmr}{m}{n}\selectfont\color{color_29791} }
\put(299.396,-209.09){\fontsize{11.991}{1}\usefont{T1}{cmr}{m}{n}\selectfont\color{color_29791}st}
\put(308.7729,-209.09){\fontsize{11.991}{1}\usefont{T1}{cmr}{m}{n}\selectfont\color{color_29791}o}
\put(315.4519,-209.09){\fontsize{11.991}{1}\usefont{T1}{cmr}{m}{n}\selectfont\color{color_29791}ra}
\put(326.1239,-209.09){\fontsize{11.991}{1}\usefont{T1}{cmr}{m}{n}\selectfont\color{color_29791}g}
\put(332.8029,-209.09){\fontsize{11.991}{1}\usefont{T1}{cmr}{m}{n}\selectfont\color{color_29791}e}
\put(339.4819,-209.09){\fontsize{11.991}{1}\usefont{T1}{cmr}{m}{n}\selectfont\color{color_29791} }
\put(342.8633,-209.09){\fontsize{11.991}{1}\usefont{T1}{cmr}{m}{n}\selectfont\color{color_29791}(/}
\put(350.1539,-209.09){\fontsize{11.991}{1}\usefont{T1}{cmr}{m}{n}\selectfont\color{color_29791}d}
\put(356.8329,-209.09){\fontsize{11.991}{1}\usefont{T1}{cmr}{m}{n}\selectfont\color{color_29791}a}
\put(363.5118,-209.09){\fontsize{11.991}{1}\usefont{T1}{cmr}{m}{n}\selectfont\color{color_29791}t}
\put(366.8933,-209.09){\fontsize{11.991}{1}\usefont{T1}{cmr}{m}{n}\selectfont\color{color_29791}a}
\put(373.5723,-209.09){\fontsize{11.991}{1}\usefont{T1}{cmr}{m}{n}\selectfont\color{color_29791}/}
\put(376.9538,-209.09){\fontsize{11.991}{1}\usefont{T1}{cmr}{m}{n}\selectfont\color{color_29791}h}
\put(383.6328,-209.09){\fontsize{11.991}{1}\usefont{T1}{cmr}{m}{n}\selectfont\color{color_29791}p}
\put(390.3117,-209.09){\fontsize{11.991}{1}\usefont{T1}{cmr}{m}{n}\selectfont\color{color_29791}cd}
\put(402.9862,-209.09){\fontsize{11.991}{1}\usefont{T1}{cmr}{m}{n}\selectfont\color{color_29791}a}
\put(409.7731,-209.09){\fontsize{11.991}{1}\usefont{T1}{cmr}{m}{n}\selectfont\color{color_29791}t}
\put(413.0706,-209.09){\fontsize{11.991}{1}\usefont{T1}{cmr}{m}{n}\selectfont\color{color_29791}a}
\put(419.8575,-209.09){\fontsize{11.991}{1}\usefont{T1}{cmr}{m}{n}\selectfont\color{color_29791})}
\put(48.213,-241.802){\fontsize{15.987}{1}\usefont{T1}{cmr}{m}{n}\selectfont\color{color_29791}•}
\put(62.187,-242.709){\fontsize{15.987}{1}\usefont{T1}{cmr}{m}{n}\selectfont\color{color_29791}G}
\put(74.57692,-242.709){\fontsize{15.987}{1}\usefont{T1}{cmr}{m}{n}\selectfont\color{color_29791}e}
\put(83.36977,-242.709){\fontsize{15.987}{1}\usefont{T1}{cmr}{m}{n}\selectfont\color{color_29791}ner}
\put(106.439,-242.709){\fontsize{15.987}{1}\usefont{T1}{cmr}{m}{n}\selectfont\color{color_29791}al}
\put(118.9089,-242.709){\fontsize{15.987}{1}\usefont{T1}{cmr}{m}{n}\selectfont\color{color_29791} }
\put(123.3053,-242.709){\fontsize{15.987}{1}\usefont{T1}{cmr}{m}{n}\selectfont\color{color_29791}U}
\put(134.8,-242.709){\fontsize{15.987}{1}\usefont{T1}{cmr}{m}{n}\selectfont\color{color_29791}s}
\put(142.8734,-242.709){\fontsize{15.987}{1}\usefont{T1}{cmr}{m}{n}\selectfont\color{color_29791}e}
\put(151.6662,-242.709){\fontsize{15.987}{1}\usefont{T1}{cmr}{m}{n}\selectfont\color{color_29791} }
\put(156.1426,-242.709){\fontsize{15.987}{1}\usefont{T1}{cmr}{m}{n}\selectfont\color{color_29791}W}
\put(170.9306,-242.709){\fontsize{15.987}{1}\usefont{T1}{cmr}{m}{n}\selectfont\color{color_29791}o}
\put(179.7234,-242.709){\fontsize{15.987}{1}\usefont{T1}{cmr}{m}{n}\selectfont\color{color_29791}r}
\put(185.0151,-242.709){\fontsize{15.987}{1}\usefont{T1}{cmr}{m}{n}\selectfont\color{color_29791}k}
\put(193.0886,-242.709){\fontsize{15.987}{1}\usefont{T1}{cmr}{m}{n}\selectfont\color{color_29791}st}
\put(205.4785,-242.709){\fontsize{15.987}{1}\usefont{T1}{cmr}{m}{n}\selectfont\color{color_29791}at}
\put(218.8436,-242.709){\fontsize{15.987}{1}\usefont{T1}{cmr}{m}{n}\selectfont\color{color_29791}i}
\put(222.4247,-242.709){\fontsize{15.987}{1}\usefont{T1}{cmr}{m}{n}\selectfont\color{color_29791}o}
\put(231.2176,-242.709){\fontsize{15.987}{1}\usefont{T1}{cmr}{m}{n}\selectfont\color{color_29791}ns}
\put(248.1798,-242.709){\fontsize{15.987}{1}\usefont{T1}{cmr}{m}{n}\selectfont\color{color_29791} }
\put(252.5762,-242.709){\fontsize{15.987}{1}\usefont{T1}{cmr}{m}{n}\selectfont\color{color_29791}\&}
\put(263.2555,-242.709){\fontsize{15.987}{1}\usefont{T1}{cmr}{m}{n}\selectfont\color{color_29791} }
\put(267.6519,-242.709){\fontsize{15.987}{1}\usefont{T1}{cmr}{m}{n}\selectfont\color{color_29791}P}
\put(278.3313,-242.709){\fontsize{15.987}{1}\usefont{T1}{cmr}{m}{n}\selectfont\color{color_29791}r}
\put(283.623,-242.709){\fontsize{15.987}{1}\usefont{T1}{cmr}{m}{n}\selectfont\color{color_29791}i}
\put(287.204,-242.709){\fontsize{15.987}{1}\usefont{T1}{cmr}{m}{n}\selectfont\color{color_29791}vat}
\put(308.4827,-242.709){\fontsize{15.987}{1}\usefont{T1}{cmr}{m}{n}\selectfont\color{color_29791}e }
\put(321.8479,-242.709){\fontsize{15.987}{1}\usefont{T1}{cmr}{m}{n}\selectfont\color{color_29791}W}
\put(336.5399,-242.709){\fontsize{15.987}{1}\usefont{T1}{cmr}{m}{n}\selectfont\color{color_29791}or}
\put(350.7204,-242.709){\fontsize{15.987}{1}\usefont{T1}{cmr}{m}{n}\selectfont\color{color_29791}ks}
\put(366.7873,-242.709){\fontsize{15.987}{1}\usefont{T1}{cmr}{m}{n}\selectfont\color{color_29791}t}
\put(371.1837,-242.709){\fontsize{15.987}{1}\usefont{T1}{cmr}{m}{n}\selectfont\color{color_29791}at}
\put(384.4689,-242.709){\fontsize{15.987}{1}\usefont{T1}{cmr}{m}{n}\selectfont\color{color_29791}i}
\put(388.05,-242.709){\fontsize{15.987}{1}\usefont{T1}{cmr}{m}{n}\selectfont\color{color_29791}ons}
\put(48.213,-276.101){\fontsize{15.987}{1}\usefont{T1}{cmr}{m}{n}\selectfont\color{color_29791}•}
\put(62.187,-277.093){\fontsize{15.987}{1}\usefont{T1}{cmr}{m}{n}\selectfont\color{color_29791}N}
\put(73.68165,-277.093){\fontsize{15.987}{1}\usefont{T1}{cmr}{m}{n}\selectfont\color{color_29791}odes}
\put(48.213,-310.513){\fontsize{15.987}{1}\usefont{T1}{cmr}{m}{n}\selectfont\color{color_29791}•}
\put(62.187,-311.392){\fontsize{15.987}{1}\usefont{T1}{cmr}{m}{n}\selectfont\color{color_29791}G}
\put(74.57692,-311.392){\fontsize{15.987}{1}\usefont{T1}{cmr}{m}{n}\selectfont\color{color_29791}P}
\put(85.25623,-311.392){\fontsize{15.987}{1}\usefont{T1}{cmr}{m}{n}\selectfont\color{color_29791}U}
\put(96.75089,-311.392){\fontsize{15.987}{1}\usefont{T1}{cmr}{m}{n}\selectfont\color{color_29791} }
\put(101.1473,-311.392){\fontsize{15.987}{1}\usefont{T1}{cmr}{m}{n}\selectfont\color{color_29791}N}
\put(112.642,-311.392){\fontsize{15.987}{1}\usefont{T1}{cmr}{m}{n}\selectfont\color{color_29791}odes}
\put(84.213,-340.306){\fontsize{11.991}{1}\usefont{T1}{cmr}{m}{n}\selectfont\color{color_29791}•}
\put(98.187,-341.014){\fontsize{11.991}{1}\usefont{T1}{cmr}{m}{n}\selectfont\color{color_29791}C}
\put(106.8685,-341.014){\fontsize{11.991}{1}\usefont{T1}{cmr}{m}{n}\selectfont\color{color_29791}u}
\put(113.5475,-341.014){\fontsize{11.991}{1}\usefont{T1}{cmr}{m}{n}\selectfont\color{color_29791}r}
\put(117.4445,-341.014){\fontsize{11.991}{1}\usefont{T1}{cmr}{m}{n}\selectfont\color{color_29791}re}
\put(128.1165,-341.014){\fontsize{11.991}{1}\usefont{T1}{cmr}{m}{n}\selectfont\color{color_29791}n}
\put(134.7955,-341.014){\fontsize{11.991}{1}\usefont{T1}{cmr}{m}{n}\selectfont\color{color_29791}t}
\put(138.177,-341.014){\fontsize{11.991}{1}\usefont{T1}{cmr}{m}{n}\selectfont\color{color_29791}l}
\put(140.7791,-341.014){\fontsize{11.991}{1}\usefont{T1}{cmr}{m}{n}\selectfont\color{color_29791}y }
\put(150.156,-341.014){\fontsize{11.991}{1}\usefont{T1}{cmr}{m}{n}\selectfont\color{color_29791}o}
\put(156.835,-341.014){\fontsize{11.991}{1}\usefont{T1}{cmr}{m}{n}\selectfont\color{color_29791}n}
\put(163.514,-341.014){\fontsize{11.991}{1}\usefont{T1}{cmr}{m}{n}\selectfont\color{color_29791}l}
\put(166.2,-341.014){\fontsize{11.991}{1}\usefont{T1}{cmr}{m}{n}\selectfont\color{color_29791}y }
\put(175.493,-341.014){\fontsize{11.991}{1}\usefont{T1}{cmr}{m}{n}\selectfont\color{color_29791}a}
\put(182.2799,-341.014){\fontsize{11.991}{1}\usefont{T1}{cmr}{m}{n}\selectfont\color{color_29791}v}
\put(188.1795,-341.014){\fontsize{11.991}{1}\usefont{T1}{cmr}{m}{n}\selectfont\color{color_29791}a}
\put(194.9664,-341.014){\fontsize{11.991}{1}\usefont{T1}{cmr}{m}{n}\selectfont\color{color_29791}i}
\put(197.5685,-341.014){\fontsize{11.991}{1}\usefont{T1}{cmr}{m}{n}\selectfont\color{color_29791}l}
\put(200.2544,-341.014){\fontsize{11.991}{1}\usefont{T1}{cmr}{m}{n}\selectfont\color{color_29791}a}
\put(206.9334,-341.014){\fontsize{11.991}{1}\usefont{T1}{cmr}{m}{n}\selectfont\color{color_29791}b}
\put(213.6124,-341.014){\fontsize{11.991}{1}\usefont{T1}{cmr}{m}{n}\selectfont\color{color_29791}l}
\put(216.2145,-341.014){\fontsize{11.991}{1}\usefont{T1}{cmr}{m}{n}\selectfont\color{color_29791}e}
\put(223.0014,-341.014){\fontsize{11.991}{1}\usefont{T1}{cmr}{m}{n}\selectfont\color{color_29791} }
\put(226.2989,-341.014){\fontsize{11.991}{1}\usefont{T1}{cmr}{m}{n}\selectfont\color{color_29791}f}
\put(229.6804,-341.014){\fontsize{11.991}{1}\usefont{T1}{cmr}{m}{n}\selectfont\color{color_29791}o}
\put(236.3593,-341.014){\fontsize{11.991}{1}\usefont{T1}{cmr}{m}{n}\selectfont\color{color_29791}r }
\put(243.6499,-341.014){\fontsize{11.991}{1}\usefont{T1}{cmr}{m}{n}\selectfont\color{color_29791}u}
\put(250.3289,-341.014){\fontsize{11.991}{1}\usefont{T1}{cmr}{m}{n}\selectfont\color{color_29791}se}
\put(263.0034,-341.014){\fontsize{11.991}{1}\usefont{T1}{cmr}{m}{n}\selectfont\color{color_29791} }
\put(266.3848,-341.014){\fontsize{11.991}{1}\usefont{T1}{cmr}{m}{n}\selectfont\color{color_29791}BA}
\put(282.4528,-341.014){\fontsize{11.991}{1}\usefont{T1}{cmr}{m}{n}\selectfont\color{color_29791}S }
\put(293.1367,-341.014){\fontsize{11.991}{1}\usefont{T1}{cmr}{m}{n}\selectfont\color{color_29791}AI}
\put(304.5042,-341.014){\fontsize{11.991}{1}\usefont{T1}{cmr}{m}{n}\selectfont\color{color_29791} }
\put(307.8017,-341.014){\fontsize{11.991}{1}\usefont{T1}{cmr}{m}{n}\selectfont\color{color_29791}L}
\put(314.5886,-341.014){\fontsize{11.991}{1}\usefont{T1}{cmr}{m}{n}\selectfont\color{color_29791}a}
\put(321.2676,-341.014){\fontsize{11.991}{1}\usefont{T1}{cmr}{m}{n}\selectfont\color{color_29791}b}
\put(327.9466,-341.014){\fontsize{11.991}{1}\usefont{T1}{cmr}{m}{n}\selectfont\color{color_29791} }
\put(331.2441,-341.014){\fontsize{11.991}{1}\usefont{T1}{cmr}{m}{n}\selectfont\color{color_29791}me}
\put(347.9116,-341.014){\fontsize{11.991}{1}\usefont{T1}{cmr}{m}{n}\selectfont\color{color_29791}mb}
\put(364.5791,-341.014){\fontsize{11.991}{1}\usefont{T1}{cmr}{m}{n}\selectfont\color{color_29791}e}
\put(371.2581,-341.014){\fontsize{11.991}{1}\usefont{T1}{cmr}{m}{n}\selectfont\color{color_29791}rs}
\put(48.213,-373.613){\fontsize{15.987}{1}\usefont{T1}{cmr}{m}{n}\selectfont\color{color_29791}•}
\put(62.187,-374.491){\fontsize{15.987}{1}\usefont{T1}{cmr}{m}{n}\selectfont\color{color_29791}D}
\put(73.68165,-374.491){\fontsize{15.987}{1}\usefont{T1}{cmr}{m}{n}\selectfont\color{color_29791}evel}
\put(103.0338,-374.491){\fontsize{15.987}{1}\usefont{T1}{cmr}{m}{n}\selectfont\color{color_29791}opm}
\put(134.0965,-374.491){\fontsize{15.987}{1}\usefont{T1}{cmr}{m}{n}\selectfont\color{color_29791}ent}
\put(156.2705,-374.491){\fontsize{15.987}{1}\usefont{T1}{cmr}{m}{n}\selectfont\color{color_29791} }
\put(160.7469,-374.491){\fontsize{15.987}{1}\usefont{T1}{cmr}{m}{n}\selectfont\color{color_29791}W}
\put(175.4389,-374.491){\fontsize{15.987}{1}\usefont{T1}{cmr}{m}{n}\selectfont\color{color_29791}or}
\put(189.6194,-374.491){\fontsize{15.987}{1}\usefont{T1}{cmr}{m}{n}\selectfont\color{color_29791}kst}
\put(210.0827,-374.491){\fontsize{15.987}{1}\usefont{T1}{cmr}{m}{n}\selectfont\color{color_29791}at}
\put(223.3679,-374.491){\fontsize{15.987}{1}\usefont{T1}{cmr}{m}{n}\selectfont\color{color_29791}i}
\put(226.949,-374.491){\fontsize{15.987}{1}\usefont{T1}{cmr}{m}{n}\selectfont\color{color_29791}on }
\put(249.123,-374.491){\fontsize{15.987}{1}\usefont{T1}{cmr}{m}{n}\selectfont\color{color_29791}and }
\put(280.1857,-374.491){\fontsize{15.987}{1}\usefont{T1}{cmr}{m}{n}\selectfont\color{color_29791}D}
\put(291.6804,-374.491){\fontsize{15.987}{1}\usefont{T1}{cmr}{m}{n}\selectfont\color{color_29791}evel}
\put(321.0325,-374.491){\fontsize{15.987}{1}\usefont{T1}{cmr}{m}{n}\selectfont\color{color_29791}opm}
\put(352.0953,-374.491){\fontsize{15.987}{1}\usefont{T1}{cmr}{m}{n}\selectfont\color{color_29791}ent}
\put(374.2693,-374.491){\fontsize{15.987}{1}\usefont{T1}{cmr}{m}{n}\selectfont\color{color_29791} }
\put(378.6657,-374.491){\fontsize{15.987}{1}\usefont{T1}{cmr}{m}{n}\selectfont\color{color_29791}N}
\put(390.2403,-374.491){\fontsize{15.987}{1}\usefont{T1}{cmr}{m}{n}\selectfont\color{color_29791}o}
\put(399.0331,-374.491){\fontsize{15.987}{1}\usefont{T1}{cmr}{m}{n}\selectfont\color{color_29791}de}
\put(84.213,-403.405){\fontsize{11.991}{1}\usefont{T1}{cmr}{m}{n}\selectfont\color{color_29791}•}
\put(98.187,-404.113){\fontsize{11.991}{1}\usefont{T1}{cmr}{m}{n}\selectfont\color{color_29791}N}
\put(106.8685,-404.113){\fontsize{11.991}{1}\usefont{T1}{cmr}{m}{n}\selectfont\color{color_29791}o}
\put(113.5475,-404.113){\fontsize{11.991}{1}\usefont{T1}{cmr}{m}{n}\selectfont\color{color_29791} }
\put(116.845,-404.113){\fontsize{11.991}{1}\usefont{T1}{cmr}{m}{n}\selectfont\color{color_29791}a}
\put(123.524,-404.113){\fontsize{11.991}{1}\usefont{T1}{cmr}{m}{n}\selectfont\color{color_29791}c}
\put(129.6034,-404.113){\fontsize{11.991}{1}\usefont{T1}{cmr}{m}{n}\selectfont\color{color_29791}ce}
\put(142.2779,-404.113){\fontsize{11.991}{1}\usefont{T1}{cmr}{m}{n}\selectfont\color{color_29791}ss,}
\put(157.5664,-404.113){\fontsize{11.991}{1}\usefont{T1}{cmr}{m}{n}\selectfont\color{color_29791} }
\put(160.9479,-404.113){\fontsize{11.991}{1}\usefont{T1}{cmr}{m}{n}\selectfont\color{color_29791}u}
\put(167.6269,-404.113){\fontsize{11.991}{1}\usefont{T1}{cmr}{m}{n}\selectfont\color{color_29791}se}
\put(180.3014,-404.113){\fontsize{11.991}{1}\usefont{T1}{cmr}{m}{n}\selectfont\color{color_29791}d}
\put(187.0883,-404.113){\fontsize{11.991}{1}\usefont{T1}{cmr}{m}{n}\selectfont\color{color_29791} }
\put(190.3858,-404.113){\fontsize{11.991}{1}\usefont{T1}{cmr}{m}{n}\selectfont\color{color_29791}f}
\put(193.7673,-404.113){\fontsize{11.991}{1}\usefont{T1}{cmr}{m}{n}\selectfont\color{color_29791}o}
\put(200.4463,-404.113){\fontsize{11.991}{1}\usefont{T1}{cmr}{m}{n}\selectfont\color{color_29791}r }
\put(207.7368,-404.113){\fontsize{11.991}{1}\usefont{T1}{cmr}{m}{n}\selectfont\color{color_29791}t}
\put(211.1183,-404.113){\fontsize{11.991}{1}\usefont{T1}{cmr}{m}{n}\selectfont\color{color_29791}e}
\put(217.7972,-404.113){\fontsize{11.991}{1}\usefont{T1}{cmr}{m}{n}\selectfont\color{color_29791}st}
\put(227.0903,-404.113){\fontsize{11.991}{1}\usefont{T1}{cmr}{m}{n}\selectfont\color{color_29791}i}
\put(229.7763,-404.113){\fontsize{11.991}{1}\usefont{T1}{cmr}{m}{n}\selectfont\color{color_29791}n}
\put(236.4552,-404.113){\fontsize{11.991}{1}\usefont{T1}{cmr}{m}{n}\selectfont\color{color_29791}g}
\put(243.1342,-404.113){\fontsize{11.991}{1}\usefont{T1}{cmr}{m}{n}\selectfont\color{color_29791} }
\put(246.5157,-404.113){\fontsize{11.991}{1}\usefont{T1}{cmr}{m}{n}\selectfont\color{color_29791}b}
\put(253.1947,-404.113){\fontsize{11.991}{1}\usefont{T1}{cmr}{m}{n}\selectfont\color{color_29791}y }
\put(262.5717,-404.113){\fontsize{11.991}{1}\usefont{T1}{cmr}{m}{n}\selectfont\color{color_29791}I}
\put(265.8692,-404.113){\fontsize{11.991}{1}\usefont{T1}{cmr}{m}{n}\selectfont\color{color_29791}T}
\end{picture}
\newpage
\begin{tikzpicture}[overlay]
\path(0pt,0pt);
\filldraw[color_283006][even odd rule]
(-15pt, 10pt) -- (704.972pt, 10pt)
 -- (704.972pt, 10pt)
 -- (704.972pt, -529.972pt)
 -- (704.972pt, -529.972pt)
 -- (-15pt, -529.972pt)
 -- (-15pt, -529.972pt)
 -- (-15pt, 10pt) -- cycle
;
\end{tikzpicture}
\begin{picture}(-5,0)(2.5,0)
\put(-15,-529.972){\includegraphics[width=720pt,height=540pt]{latexImage_15b9e6acda8be87c69673355bc94d374.png}}
\put(48.411,-60.69601){\fontsize{21.997}{1}\usefont{T1}{cmr}{b}{n}\selectfont\color{color_29791}Ac}
\put(76.50117,-60.69601){\fontsize{21.997}{1}\usefont{T1}{cmr}{b}{n}\selectfont\color{color_29791}c}
\put(88.7095,-60.69601){\fontsize{21.997}{1}\usefont{T1}{cmr}{b}{n}\selectfont\color{color_29791}e}
\put(100.9838,-60.69601){\fontsize{21.997}{1}\usefont{T1}{cmr}{b}{n}\selectfont\color{color_29791}s}
\put(113.1922,-60.69601){\fontsize{21.997}{1}\usefont{T1}{cmr}{b}{n}\selectfont\color{color_29791}s}
\put(48.213,-116.51){\fontsize{15.987}{1}\usefont{T1}{cmr}{m}{n}\selectfont\color{color_29791}A}
\put(58.89232,-116.51){\fontsize{15.987}{1}\usefont{T1}{cmr}{m}{n}\selectfont\color{color_29791}ut}
\put(72.17751,-116.51){\fontsize{15.987}{1}\usefont{T1}{cmr}{m}{n}\selectfont\color{color_29791}hen}
\put(98.74791,-116.51){\fontsize{15.987}{1}\usefont{T1}{cmr}{m}{n}\selectfont\color{color_29791}t}
\put(103.2243,-116.51){\fontsize{15.987}{1}\usefont{T1}{cmr}{m}{n}\selectfont\color{color_29791}i}
\put(106.8054,-116.51){\fontsize{15.987}{1}\usefont{T1}{cmr}{m}{n}\selectfont\color{color_29791}cat}
\put(128.164,-116.51){\fontsize{15.987}{1}\usefont{T1}{cmr}{m}{n}\selectfont\color{color_29791}i}
\put(131.7451,-116.51){\fontsize{15.987}{1}\usefont{T1}{cmr}{m}{n}\selectfont\color{color_29791}o}
\put(140.5379,-116.51){\fontsize{15.987}{1}\usefont{T1}{cmr}{m}{n}\selectfont\color{color_29791}n}
\put(84.213,-145.31){\fontsize{11.991}{1}\usefont{T1}{cmr}{m}{n}\selectfont\color{color_29791}•}
\put(98.187,-145.991){\fontsize{11.991}{1}\usefont{T1}{cmr}{m}{n}\selectfont\color{color_29791}T}
\put(105.4775,-145.991){\fontsize{11.991}{1}\usefont{T1}{cmr}{m}{n}\selectfont\color{color_29791}h}
\put(112.1565,-145.991){\fontsize{11.991}{1}\usefont{T1}{cmr}{m}{n}\selectfont\color{color_29791}re}
\put(122.8285,-145.991){\fontsize{11.991}{1}\usefont{T1}{cmr}{m}{n}\selectfont\color{color_29791}e}
\put(129.709,-145.991){\fontsize{11.991}{1}\usefont{T1}{cmr}{m}{n}\selectfont\color{color_29791} }
\put(132.997,-145.991){\fontsize{11.991}{1}\usefont{T1}{cmr}{m}{n}\selectfont\color{color_29791}p}
\put(139.676,-145.991){\fontsize{11.991}{1}\usefont{T1}{cmr}{m}{n}\selectfont\color{color_29791}a}
\put(146.4629,-145.991){\fontsize{11.991}{1}\usefont{T1}{cmr}{m}{n}\selectfont\color{color_29791}s}
\put(152.3625,-145.991){\fontsize{11.991}{1}\usefont{T1}{cmr}{m}{n}\selectfont\color{color_29791}s}
\put(158.4419,-145.991){\fontsize{11.991}{1}\usefont{T1}{cmr}{m}{n}\selectfont\color{color_29791}w}
\put(167.0394,-145.991){\fontsize{11.991}{1}\usefont{T1}{cmr}{m}{n}\selectfont\color{color_29791}o}
\put(173.7184,-145.991){\fontsize{11.991}{1}\usefont{T1}{cmr}{m}{n}\selectfont\color{color_29791}rd}
\put(184.3904,-145.991){\fontsize{11.991}{1}\usefont{T1}{cmr}{m}{n}\selectfont\color{color_29791}s }
\put(193.7674,-145.991){\fontsize{11.991}{1}\usefont{T1}{cmr}{m}{n}\selectfont\color{color_29791}–}
\put(200.4464,-145.991){\fontsize{11.991}{1}\usefont{T1}{cmr}{m}{n}\selectfont\color{color_29791} }
\put(203.7439,-145.991){\fontsize{11.991}{1}\usefont{T1}{cmr}{m}{n}\selectfont\color{color_29791}U}
\put(212.4254,-145.991){\fontsize{11.991}{1}\usefont{T1}{cmr}{m}{n}\selectfont\color{color_29791}N}
\put(221.0229,-145.991){\fontsize{11.991}{1}\usefont{T1}{cmr}{m}{n}\selectfont\color{color_29791}I}
\put(224.4044,-145.991){\fontsize{11.991}{1}\usefont{T1}{cmr}{m}{n}\selectfont\color{color_29791}X }
\put(235.7719,-145.991){\fontsize{11.991}{1}\usefont{T1}{cmr}{m}{n}\selectfont\color{color_29791}(}
\put(239.669,-145.991){\fontsize{11.991}{1}\usefont{T1}{cmr}{m}{n}\selectfont\color{color_29791}N}
\put(248.3504,-145.991){\fontsize{11.991}{1}\usefont{T1}{cmr}{m}{n}\selectfont\color{color_29791}I}
\put(251.6479,-145.991){\fontsize{11.991}{1}\usefont{T1}{cmr}{m}{n}\selectfont\color{color_29791}S}
\put(259.7299,-145.991){\fontsize{11.991}{1}\usefont{T1}{cmr}{m}{n}\selectfont\color{color_29791})}
\put(263.627,-145.991){\fontsize{11.991}{1}\usefont{T1}{cmr}{m}{n}\selectfont\color{color_29791},}
\put(267.0084,-145.991){\fontsize{11.991}{1}\usefont{T1}{cmr}{m}{n}\selectfont\color{color_29791} }
\put(270.3059,-145.991){\fontsize{11.991}{1}\usefont{T1}{cmr}{m}{n}\selectfont\color{color_29791}L}
\put(277.0928,-145.991){\fontsize{11.991}{1}\usefont{T1}{cmr}{m}{n}\selectfont\color{color_29791}D}
\put(285.6904,-145.991){\fontsize{11.991}{1}\usefont{T1}{cmr}{m}{n}\selectfont\color{color_29791}A}
\put(293.7723,-145.991){\fontsize{11.991}{1}\usefont{T1}{cmr}{m}{n}\selectfont\color{color_29791}P}
\put(301.5665,-145.991){\fontsize{11.991}{1}\usefont{T1}{cmr}{m}{n}\selectfont\color{color_29791} }
\put(304.864,-145.991){\fontsize{11.991}{1}\usefont{T1}{cmr}{m}{n}\selectfont\color{color_29791}a}
\put(311.543,-145.991){\fontsize{11.991}{1}\usefont{T1}{cmr}{m}{n}\selectfont\color{color_29791}n}
\put(318.3299,-145.991){\fontsize{11.991}{1}\usefont{T1}{cmr}{m}{n}\selectfont\color{color_29791}d}
\put(325.0089,-145.991){\fontsize{11.991}{1}\usefont{T1}{cmr}{m}{n}\selectfont\color{color_29791} }
\put(328.3064,-145.991){\fontsize{11.991}{1}\usefont{T1}{cmr}{m}{n}\selectfont\color{color_29791}S}
\put(336.3883,-145.991){\fontsize{11.991}{1}\usefont{T1}{cmr}{m}{n}\selectfont\color{color_29791}a}
\put(343.0673,-145.991){\fontsize{11.991}{1}\usefont{T1}{cmr}{m}{n}\selectfont\color{color_29791}m}
\put(352.9599,-145.991){\fontsize{11.991}{1}\usefont{T1}{cmr}{m}{n}\selectfont\color{color_29791}b}
\put(359.7468,-145.991){\fontsize{11.991}{1}\usefont{T1}{cmr}{m}{n}\selectfont\color{color_29791}a}
\put(366.4258,-145.991){\fontsize{11.991}{1}\usefont{T1}{cmr}{m}{n}\selectfont\color{color_29791},}
\put(370.313,-145.991){\fontsize{11.991}{1}\usefont{T1}{cmr}{m}{n}\selectfont\color{color_29791} }
\put(84.213,-174.11){\fontsize{11.991}{1}\usefont{T1}{cmr}{m}{n}\selectfont\color{color_29791}•}
\put(98.187,-174.791){\fontsize{11.991}{1}\usefont{T1}{cmr}{m}{n}\selectfont\color{color_29791}U}
\put(106.8685,-174.791){\fontsize{11.991}{1}\usefont{T1}{cmr}{m}{n}\selectfont\color{color_29791}N}
\put(115.466,-174.791){\fontsize{11.991}{1}\usefont{T1}{cmr}{m}{n}\selectfont\color{color_29791}I}
\put(118.8475,-174.791){\fontsize{11.991}{1}\usefont{T1}{cmr}{m}{n}\selectfont\color{color_29791}X }
\put(130.131,-174.791){\fontsize{11.991}{1}\usefont{T1}{cmr}{m}{n}\selectfont\color{color_29791}f}
\put(133.5125,-174.791){\fontsize{11.991}{1}\usefont{T1}{cmr}{m}{n}\selectfont\color{color_29791}o}
\put(140.1915,-174.791){\fontsize{11.991}{1}\usefont{T1}{cmr}{m}{n}\selectfont\color{color_29791}r }
\put(147.566,-174.791){\fontsize{11.991}{1}\usefont{T1}{cmr}{m}{n}\selectfont\color{color_29791}b}
\put(154.2449,-174.791){\fontsize{11.991}{1}\usefont{T1}{cmr}{m}{n}\selectfont\color{color_29791}sl}
\put(162.8425,-174.791){\fontsize{11.991}{1}\usefont{T1}{cmr}{m}{n}\selectfont\color{color_29791}ce}
\put(175.517,-174.791){\fontsize{11.991}{1}\usefont{T1}{cmr}{m}{n}\selectfont\color{color_29791}n}
\put(182.3039,-174.791){\fontsize{11.991}{1}\usefont{T1}{cmr}{m}{n}\selectfont\color{color_29791}b}
\put(188.9829,-174.791){\fontsize{11.991}{1}\usefont{T1}{cmr}{m}{n}\selectfont\color{color_29791} }
\put(192.2804,-174.791){\fontsize{11.991}{1}\usefont{T1}{cmr}{m}{n}\selectfont\color{color_29791}/}
\put(195.6619,-174.791){\fontsize{11.991}{1}\usefont{T1}{cmr}{m}{n}\selectfont\color{color_29791} }
\put(198.9594,-174.791){\fontsize{11.991}{1}\usefont{T1}{cmr}{m}{n}\selectfont\color{color_29791}b}
\put(205.7463,-174.791){\fontsize{11.991}{1}\usefont{T1}{cmr}{m}{n}\selectfont\color{color_29791}sl}
\put(214.3439,-174.791){\fontsize{11.991}{1}\usefont{T1}{cmr}{m}{n}\selectfont\color{color_29791}ce}
\put(227.0183,-174.791){\fontsize{11.991}{1}\usefont{T1}{cmr}{m}{n}\selectfont\color{color_29791}n}
\put(233.8053,-174.791){\fontsize{11.991}{1}\usefont{T1}{cmr}{m}{n}\selectfont\color{color_29791}c }
\put(243.0983,-174.791){\fontsize{11.991}{1}\usefont{T1}{cmr}{m}{n}\selectfont\color{color_29791}a}
\put(249.7773,-174.791){\fontsize{11.991}{1}\usefont{T1}{cmr}{m}{n}\selectfont\color{color_29791}n}
\put(256.4562,-174.791){\fontsize{11.991}{1}\usefont{T1}{cmr}{m}{n}\selectfont\color{color_29791}d}
\put(263.2431,-174.791){\fontsize{11.991}{1}\usefont{T1}{cmr}{m}{n}\selectfont\color{color_29791} }
\put(266.5406,-174.791){\fontsize{11.991}{1}\usefont{T1}{cmr}{m}{n}\selectfont\color{color_29791}L}
\put(273.3275,-174.791){\fontsize{11.991}{1}\usefont{T1}{cmr}{m}{n}\selectfont\color{color_29791}D}
\put(281.9251,-174.791){\fontsize{11.991}{1}\usefont{T1}{cmr}{m}{n}\selectfont\color{color_29791}AP}
\put(297.8012,-174.791){\fontsize{11.991}{1}\usefont{T1}{cmr}{m}{n}\selectfont\color{color_29791} }
\put(301.0987,-174.791){\fontsize{11.991}{1}\usefont{T1}{cmr}{m}{n}\selectfont\color{color_29791}f}
\put(304.4802,-174.791){\fontsize{11.991}{1}\usefont{T1}{cmr}{m}{n}\selectfont\color{color_29791}o}
\put(311.1591,-174.791){\fontsize{11.991}{1}\usefont{T1}{cmr}{m}{n}\selectfont\color{color_29791}r}
\put(315.0562,-174.791){\fontsize{11.991}{1}\usefont{T1}{cmr}{m}{n}\selectfont\color{color_29791} }
\put(318.4377,-174.791){\fontsize{11.991}{1}\usefont{T1}{cmr}{m}{n}\selectfont\color{color_29791}H}
\put(327.0352,-174.791){\fontsize{11.991}{1}\usefont{T1}{cmr}{m}{n}\selectfont\color{color_29791}P}
\put(335.1172,-174.791){\fontsize{11.991}{1}\usefont{T1}{cmr}{m}{n}\selectfont\color{color_29791}C}
\put(343.7147,-174.791){\fontsize{11.991}{1}\usefont{T1}{cmr}{m}{n}\selectfont\color{color_29791} }
\put(347.0962,-174.791){\fontsize{11.991}{1}\usefont{T1}{cmr}{m}{n}\selectfont\color{color_29791}w}
\put(355.6937,-174.791){\fontsize{11.991}{1}\usefont{T1}{cmr}{m}{n}\selectfont\color{color_29791}o}
\put(362.3727,-174.791){\fontsize{11.991}{1}\usefont{T1}{cmr}{m}{n}\selectfont\color{color_29791}rkst}
\put(381.6542,-174.791){\fontsize{11.991}{1}\usefont{T1}{cmr}{m}{n}\selectfont\color{color_29791}a}
\put(388.4411,-174.791){\fontsize{11.991}{1}\usefont{T1}{cmr}{m}{n}\selectfont\color{color_29791}t}
\put(391.7386,-174.791){\fontsize{11.991}{1}\usefont{T1}{cmr}{m}{n}\selectfont\color{color_29791}i}
\put(394.4246,-174.791){\fontsize{11.991}{1}\usefont{T1}{cmr}{m}{n}\selectfont\color{color_29791}o}
\put(401.1036,-174.791){\fontsize{11.991}{1}\usefont{T1}{cmr}{m}{n}\selectfont\color{color_29791}n}
\put(407.7826,-174.791){\fontsize{11.991}{1}\usefont{T1}{cmr}{m}{n}\selectfont\color{color_29791}s}
\put(84.213,-202.91){\fontsize{11.991}{1}\usefont{T1}{cmr}{m}{n}\selectfont\color{color_29791}•}
\put(98.187,-203.591){\fontsize{11.991}{1}\usefont{T1}{cmr}{m}{n}\selectfont\color{color_29791}T}
\put(105.0818,-203.591){\fontsize{11.991}{1}\usefont{T1}{cmr}{m}{n}\selectfont\color{color_29791}r}
\put(108.9789,-203.591){\fontsize{11.991}{1}\usefont{T1}{cmr}{m}{n}\selectfont\color{color_29791}y }
\put(118.3559,-203.591){\fontsize{11.991}{1}\usefont{T1}{cmr}{m}{n}\selectfont\color{color_29791}t}
\put(121.6534,-203.591){\fontsize{11.991}{1}\usefont{T1}{cmr}{m}{n}\selectfont\color{color_29791}o}
\put(128.4403,-203.591){\fontsize{11.991}{1}\usefont{T1}{cmr}{m}{n}\selectfont\color{color_29791} }
\put(131.7378,-203.591){\fontsize{11.991}{1}\usefont{T1}{cmr}{m}{n}\selectfont\color{color_29791}ke}
\put(144.4123,-203.591){\fontsize{11.991}{1}\usefont{T1}{cmr}{m}{n}\selectfont\color{color_29791}e}
\put(151.1992,-203.591){\fontsize{11.991}{1}\usefont{T1}{cmr}{m}{n}\selectfont\color{color_29791}p}
\put(157.8782,-203.591){\fontsize{11.991}{1}\usefont{T1}{cmr}{m}{n}\selectfont\color{color_29791} }
\put(161.1757,-203.591){\fontsize{11.991}{1}\usefont{T1}{cmr}{m}{n}\selectfont\color{color_29791}a}
\put(167.9626,-203.591){\fontsize{11.991}{1}\usefont{T1}{cmr}{m}{n}\selectfont\color{color_29791}l}
\put(170.5647,-203.591){\fontsize{11.991}{1}\usefont{T1}{cmr}{m}{n}\selectfont\color{color_29791}l}
\put(173.1667,-203.591){\fontsize{11.991}{1}\usefont{T1}{cmr}{m}{n}\selectfont\color{color_29791} }
\put(176.5482,-203.591){\fontsize{11.991}{1}\usefont{T1}{cmr}{m}{n}\selectfont\color{color_29791}t}
\put(179.9297,-203.591){\fontsize{11.991}{1}\usefont{T1}{cmr}{m}{n}\selectfont\color{color_29791}h}
\put(186.6087,-203.591){\fontsize{11.991}{1}\usefont{T1}{cmr}{m}{n}\selectfont\color{color_29791}e}
\put(193.2877,-203.591){\fontsize{11.991}{1}\usefont{T1}{cmr}{m}{n}\selectfont\color{color_29791}se}
\put(205.9621,-203.591){\fontsize{11.991}{1}\usefont{T1}{cmr}{m}{n}\selectfont\color{color_29791} }
\put(209.3436,-203.591){\fontsize{11.991}{1}\usefont{T1}{cmr}{m}{n}\selectfont\color{color_29791}p}
\put(216.0226,-203.591){\fontsize{11.991}{1}\usefont{T1}{cmr}{m}{n}\selectfont\color{color_29791}a}
\put(222.7016,-203.591){\fontsize{11.991}{1}\usefont{T1}{cmr}{m}{n}\selectfont\color{color_29791}ssw}
\put(243.3741,-203.591){\fontsize{11.991}{1}\usefont{T1}{cmr}{m}{n}\selectfont\color{color_29791}o}
\put(250.053,-203.591){\fontsize{11.991}{1}\usefont{T1}{cmr}{m}{n}\selectfont\color{color_29791}r}
\put(253.9501,-203.591){\fontsize{11.991}{1}\usefont{T1}{cmr}{m}{n}\selectfont\color{color_29791}d}
\put(260.737,-203.591){\fontsize{11.991}{1}\usefont{T1}{cmr}{m}{n}\selectfont\color{color_29791} }
\put(264.0345,-203.591){\fontsize{11.991}{1}\usefont{T1}{cmr}{m}{n}\selectfont\color{color_29791}syn}
\put(282.7045,-203.591){\fontsize{11.991}{1}\usefont{T1}{cmr}{m}{n}\selectfont\color{color_29791}ch}
\put(295.379,-203.591){\fontsize{11.991}{1}\usefont{T1}{cmr}{m}{n}\selectfont\color{color_29791}ro}
\put(306.051,-203.591){\fontsize{11.991}{1}\usefont{T1}{cmr}{m}{n}\selectfont\color{color_29791}n}
\put(312.8379,-203.591){\fontsize{11.991}{1}\usefont{T1}{cmr}{m}{n}\selectfont\color{color_29791}i}
\put(315.4399,-203.591){\fontsize{11.991}{1}\usefont{T1}{cmr}{m}{n}\selectfont\color{color_29791}se}
\put(328.1144,-203.591){\fontsize{11.991}{1}\usefont{T1}{cmr}{m}{n}\selectfont\color{color_29791}d}
\put(84.213,-231.71){\fontsize{11.991}{1}\usefont{T1}{cmr}{m}{n}\selectfont\color{color_29791}•}
\put(98.187,-232.306){\fontsize{11.991}{1}\usefont{T1}{cmr}{m}{n}\selectfont\color{color_29791}W}
\put(109.2787,-232.306){\fontsize{11.991}{1}\usefont{T1}{cmr}{m}{n}\selectfont\color{color_29791}e}
\put(115.9577,-232.306){\fontsize{11.991}{1}\usefont{T1}{cmr}{m}{n}\selectfont\color{color_29791} }
\put(119.2552,-232.306){\fontsize{11.991}{1}\usefont{T1}{cmr}{m}{n}\selectfont\color{color_29791}a}
\put(126.0421,-232.306){\fontsize{11.991}{1}\usefont{T1}{cmr}{m}{n}\selectfont\color{color_29791}r}
\put(129.9392,-232.306){\fontsize{11.991}{1}\usefont{T1}{cmr}{m}{n}\selectfont\color{color_29791}e}
\put(136.6182,-232.306){\fontsize{11.991}{1}\usefont{T1}{cmr}{m}{n}\selectfont\color{color_29791} }
\put(139.9996,-232.306){\fontsize{11.991}{1}\usefont{T1}{cmr}{m}{n}\selectfont\color{color_29791}w}
\put(148.5972,-232.306){\fontsize{11.991}{1}\usefont{T1}{cmr}{m}{n}\selectfont\color{color_29791}o}
\put(155.2762,-232.306){\fontsize{11.991}{1}\usefont{T1}{cmr}{m}{n}\selectfont\color{color_29791}rki}
\put(167.9507,-232.306){\fontsize{11.991}{1}\usefont{T1}{cmr}{m}{n}\selectfont\color{color_29791}n}
\put(174.6296,-232.306){\fontsize{11.991}{1}\usefont{T1}{cmr}{m}{n}\selectfont\color{color_29791}g}
\put(181.3086,-232.306){\fontsize{11.991}{1}\usefont{T1}{cmr}{m}{n}\selectfont\color{color_29791} }
\put(184.6062,-232.306){\fontsize{11.991}{1}\usefont{T1}{cmr}{m}{n}\selectfont\color{color_29791}t}
\put(187.9876,-232.306){\fontsize{11.991}{1}\usefont{T1}{cmr}{m}{n}\selectfont\color{color_29791}o}
\put(194.6666,-232.306){\fontsize{11.991}{1}\usefont{T1}{cmr}{m}{n}\selectfont\color{color_29791} }
\put(198.0481,-232.306){\fontsize{11.991}{1}\usefont{T1}{cmr}{m}{n}\selectfont\color{color_29791}si}
\put(206.6456,-232.306){\fontsize{11.991}{1}\usefont{T1}{cmr}{m}{n}\selectfont\color{color_29791}mp}
\put(223.3131,-232.306){\fontsize{11.991}{1}\usefont{T1}{cmr}{m}{n}\selectfont\color{color_29791}l}
\put(225.9991,-232.306){\fontsize{11.991}{1}\usefont{T1}{cmr}{m}{n}\selectfont\color{color_29791}i}
\put(228.6012,-232.306){\fontsize{11.991}{1}\usefont{T1}{cmr}{m}{n}\selectfont\color{color_29791}f}
\put(231.9826,-232.306){\fontsize{11.991}{1}\usefont{T1}{cmr}{m}{n}\selectfont\color{color_29791}y}
\put(237.8822,-232.306){\fontsize{11.991}{1}\usefont{T1}{cmr}{m}{n}\selectfont\color{color_29791} }
\put(241.2636,-232.306){\fontsize{11.991}{1}\usefont{T1}{cmr}{m}{n}\selectfont\color{color_29791}t}
\put(244.6451,-232.306){\fontsize{11.991}{1}\usefont{T1}{cmr}{m}{n}\selectfont\color{color_29791}h}
\put(251.3241,-232.306){\fontsize{11.991}{1}\usefont{T1}{cmr}{m}{n}\selectfont\color{color_29791}e}
\put(258.0031,-232.306){\fontsize{11.991}{1}\usefont{T1}{cmr}{m}{n}\selectfont\color{color_29791} }
\put(261.3846,-232.306){\fontsize{11.991}{1}\usefont{T1}{cmr}{m}{n}\selectfont\color{color_29791}si}
\put(269.9821,-232.306){\fontsize{11.991}{1}\usefont{T1}{cmr}{m}{n}\selectfont\color{color_29791}t}
\put(273.3636,-232.306){\fontsize{11.991}{1}\usefont{T1}{cmr}{m}{n}\selectfont\color{color_29791}u}
\put(280.0425,-232.306){\fontsize{11.991}{1}\usefont{T1}{cmr}{m}{n}\selectfont\color{color_29791}a}
\put(286.7215,-232.306){\fontsize{11.991}{1}\usefont{T1}{cmr}{m}{n}\selectfont\color{color_29791}t}
\put(290.103,-232.306){\fontsize{11.991}{1}\usefont{T1}{cmr}{m}{n}\selectfont\color{color_29791}i}
\put(292.705,-232.306){\fontsize{11.991}{1}\usefont{T1}{cmr}{m}{n}\selectfont\color{color_29791}o}
\put(299.384,-232.306){\fontsize{11.991}{1}\usefont{T1}{cmr}{m}{n}\selectfont\color{color_29791}n}
\put(48.213,-265.896){\fontsize{15.987}{1}\usefont{T1}{cmr}{m}{n}\selectfont\color{color_29791}S}
\put(58.89232,-265.896){\fontsize{15.987}{1}\usefont{T1}{cmr}{m}{n}\selectfont\color{color_29791}S}
\put(69.57163,-265.896){\fontsize{15.987}{1}\usefont{T1}{cmr}{m}{n}\selectfont\color{color_29791}H}
\put(84.213,-294.809){\fontsize{11.991}{1}\usefont{T1}{cmr}{m}{n}\selectfont\color{color_29791}•}
\put(98.187,-295.49){\fontsize{11.991}{1}\usefont{T1}{cmr}{m}{n}\selectfont\color{color_29791}F}
\put(105.4775,-295.49){\fontsize{11.991}{1}\usefont{T1}{cmr}{m}{n}\selectfont\color{color_29791}i}
\put(108.1635,-295.49){\fontsize{11.991}{1}\usefont{T1}{cmr}{m}{n}\selectfont\color{color_29791}r}
\put(112.0606,-295.49){\fontsize{11.991}{1}\usefont{T1}{cmr}{m}{n}\selectfont\color{color_29791}st}
\put(121.4375,-295.49){\fontsize{11.991}{1}\usefont{T1}{cmr}{m}{n}\selectfont\color{color_29791} }
\put(124.819,-295.49){\fontsize{11.991}{1}\usefont{T1}{cmr}{m}{n}\selectfont\color{color_29791}c}
\put(130.7186,-295.49){\fontsize{11.991}{1}\usefont{T1}{cmr}{m}{n}\selectfont\color{color_29791}o}
\put(137.5055,-295.49){\fontsize{11.991}{1}\usefont{T1}{cmr}{m}{n}\selectfont\color{color_29791}n}
\put(144.1845,-295.49){\fontsize{11.991}{1}\usefont{T1}{cmr}{m}{n}\selectfont\color{color_29791}n}
\put(150.8634,-295.49){\fontsize{11.991}{1}\usefont{T1}{cmr}{m}{n}\selectfont\color{color_29791}e}
\put(157.5424,-295.49){\fontsize{11.991}{1}\usefont{T1}{cmr}{m}{n}\selectfont\color{color_29791}ct}
\put(166.9194,-295.49){\fontsize{11.991}{1}\usefont{T1}{cmr}{m}{n}\selectfont\color{color_29791} }
\put(170.2169,-295.49){\fontsize{11.991}{1}\usefont{T1}{cmr}{m}{n}\selectfont\color{color_29791}t}
\put(173.5984,-295.49){\fontsize{11.991}{1}\usefont{T1}{cmr}{m}{n}\selectfont\color{color_29791}o}
\put(180.2774,-295.49){\fontsize{11.991}{1}\usefont{T1}{cmr}{m}{n}\selectfont\color{color_29791} }
\put(183.6588,-295.49){\fontsize{11.991}{1}\usefont{T1}{cmr}{m}{n}\selectfont\color{color_29791}g}
\put(190.3378,-295.49){\fontsize{11.991}{1}\usefont{T1}{cmr}{m}{n}\selectfont\color{color_29791}a}
\put(197.0168,-295.49){\fontsize{11.991}{1}\usefont{T1}{cmr}{m}{n}\selectfont\color{color_29791}t}
\put(200.3983,-295.49){\fontsize{11.991}{1}\usefont{T1}{cmr}{m}{n}\selectfont\color{color_29791}e}
\put(207.0773,-295.49){\fontsize{11.991}{1}\usefont{T1}{cmr}{m}{n}\selectfont\color{color_29791}w}
\put(215.7588,-295.49){\fontsize{11.991}{1}\usefont{T1}{cmr}{m}{n}\selectfont\color{color_29791}a}
\put(222.4377,-295.49){\fontsize{11.991}{1}\usefont{T1}{cmr}{m}{n}\selectfont\color{color_29791}y }
\put(231.7308,-295.49){\fontsize{11.991}{1}\usefont{T1}{cmr}{m}{n}\selectfont\color{color_29791}h}
\put(238.4098,-295.49){\fontsize{11.991}{1}\usefont{T1}{cmr}{m}{n}\selectfont\color{color_29791}o}
\put(245.1967,-295.49){\fontsize{11.991}{1}\usefont{T1}{cmr}{m}{n}\selectfont\color{color_29791}st}
\put(254.4897,-295.49){\fontsize{11.991}{1}\usefont{T1}{cmr}{m}{n}\selectfont\color{color_29791}s:}
\put(264.298,-295.49){\fontsize{11.991}{1}\usefont{T1}{cmr}{m}{n}\selectfont\color{color_29791} }
\put(267.614,-295.49){\fontsize{11.991}{1}\usefont{T1}{cmr}{m}{n}\selectfont\color{color_29791}b}
\put(274.4009,-295.49){\fontsize{11.991}{1}\usefont{T1}{cmr}{m}{n}\selectfont\color{color_29791}sl}
\put(282.9984,-295.49){\fontsize{11.991}{1}\usefont{T1}{cmr}{m}{n}\selectfont\color{color_29791}ce}
\put(295.6729,-295.49){\fontsize{11.991}{1}\usefont{T1}{cmr}{m}{n}\selectfont\color{color_29791}n}
\put(302.4598,-295.49){\fontsize{11.991}{1}\usefont{T1}{cmr}{m}{n}\selectfont\color{color_29791}b}
\put(309.1388,-295.49){\fontsize{11.991}{1}\usefont{T1}{cmr}{m}{n}\selectfont\color{color_29791}.}
\put(312.4363,-295.49){\fontsize{11.991}{1}\usefont{T1}{cmr}{m}{n}\selectfont\color{color_29791}n}
\put(319.2232,-295.49){\fontsize{11.991}{1}\usefont{T1}{cmr}{m}{n}\selectfont\color{color_29791}e}
\put(325.9022,-295.49){\fontsize{11.991}{1}\usefont{T1}{cmr}{m}{n}\selectfont\color{color_29791}r}
\put(329.7993,-295.49){\fontsize{11.991}{1}\usefont{T1}{cmr}{m}{n}\selectfont\color{color_29791}c-b}
\put(346.4668,-295.49){\fontsize{11.991}{1}\usefont{T1}{cmr}{m}{n}\selectfont\color{color_29791}a}
\put(353.1458,-295.49){\fontsize{11.991}{1}\usefont{T1}{cmr}{m}{n}\selectfont\color{color_29791}s.}
\put(362.5227,-295.49){\fontsize{11.991}{1}\usefont{T1}{cmr}{m}{n}\selectfont\color{color_29791}a}
\put(369.2017,-295.49){\fontsize{11.991}{1}\usefont{T1}{cmr}{m}{n}\selectfont\color{color_29791}c.}
\put(378.5786,-295.49){\fontsize{11.991}{1}\usefont{T1}{cmr}{m}{n}\selectfont\color{color_29791}u}
\put(385.2576,-295.49){\fontsize{11.991}{1}\usefont{T1}{cmr}{m}{n}\selectfont\color{color_29791}k }
\put(394.6346,-295.49){\fontsize{11.991}{1}\usefont{T1}{cmr}{m}{n}\selectfont\color{color_29791}/}
\put(398.206,-295.49){\fontsize{11.991}{1}\usefont{T1}{cmr}{m}{n}\selectfont\color{color_29791} }
\put(401.608,-295.49){\fontsize{11.991}{1}\usefont{T1}{cmr}{m}{n}\selectfont\color{color_29791}b}
\put(408.287,-295.49){\fontsize{11.991}{1}\usefont{T1}{cmr}{m}{n}\selectfont\color{color_29791}sl}
\put(416.8845,-295.49){\fontsize{11.991}{1}\usefont{T1}{cmr}{m}{n}\selectfont\color{color_29791}ce}
\put(429.559,-295.49){\fontsize{11.991}{1}\usefont{T1}{cmr}{m}{n}\selectfont\color{color_29791}n}
\put(436.3459,-295.49){\fontsize{11.991}{1}\usefont{T1}{cmr}{m}{n}\selectfont\color{color_29791}c.}
\put(445.6389,-295.49){\fontsize{11.991}{1}\usefont{T1}{cmr}{m}{n}\selectfont\color{color_29791}n}
\put(452.4258,-295.49){\fontsize{11.991}{1}\usefont{T1}{cmr}{m}{n}\selectfont\color{color_29791}e}
\put(459.1048,-295.49){\fontsize{11.991}{1}\usefont{T1}{cmr}{m}{n}\selectfont\color{color_29791}r}
\put(463.0019,-295.49){\fontsize{11.991}{1}\usefont{T1}{cmr}{m}{n}\selectfont\color{color_29791}c-b}
\put(479.6694,-295.49){\fontsize{11.991}{1}\usefont{T1}{cmr}{m}{n}\selectfont\color{color_29791}a}
\put(486.3484,-295.49){\fontsize{11.991}{1}\usefont{T1}{cmr}{m}{n}\selectfont\color{color_29791}s.}
\put(495.7253,-295.49){\fontsize{11.991}{1}\usefont{T1}{cmr}{m}{n}\selectfont\color{color_29791}a}
\put(502.4043,-295.49){\fontsize{11.991}{1}\usefont{T1}{cmr}{m}{n}\selectfont\color{color_29791}c.}
\put(511.7813,-295.49){\fontsize{11.991}{1}\usefont{T1}{cmr}{m}{n}\selectfont\color{color_29791}u}
\put(518.4603,-295.49){\fontsize{11.991}{1}\usefont{T1}{cmr}{m}{n}\selectfont\color{color_29791}k}
\put(524.688,-295.49){\fontsize{11.991}{1}\usefont{T1}{cmr}{m}{n}\selectfont\color{color_29791} }
\put(84.213,-323.609){\fontsize{11.991}{1}\usefont{T1}{cmr}{m}{n}\selectfont\color{color_29791}•}
\put(98.187,-324.205){\fontsize{11.991}{1}\usefont{T1}{cmr}{m}{n}\selectfont\color{color_29791}S}
\put(106.2689,-324.205){\fontsize{11.991}{1}\usefont{T1}{cmr}{m}{n}\selectfont\color{color_29791}e}
\put(112.9479,-324.205){\fontsize{11.991}{1}\usefont{T1}{cmr}{m}{n}\selectfont\color{color_29791}co}
\put(125.6224,-324.205){\fontsize{11.991}{1}\usefont{T1}{cmr}{m}{n}\selectfont\color{color_29791}n}
\put(132.3014,-324.205){\fontsize{11.991}{1}\usefont{T1}{cmr}{m}{n}\selectfont\color{color_29791}d}
\put(138.9804,-324.205){\fontsize{11.991}{1}\usefont{T1}{cmr}{m}{n}\selectfont\color{color_29791} }
\put(142.3618,-324.205){\fontsize{11.991}{1}\usefont{T1}{cmr}{m}{n}\selectfont\color{color_29791}co}
\put(155.0363,-324.205){\fontsize{11.991}{1}\usefont{T1}{cmr}{m}{n}\selectfont\color{color_29791}n}
\put(161.7153,-324.205){\fontsize{11.991}{1}\usefont{T1}{cmr}{m}{n}\selectfont\color{color_29791}n}
\put(168.5022,-324.205){\fontsize{11.991}{1}\usefont{T1}{cmr}{m}{n}\selectfont\color{color_29791}e}
\put(175.1812,-324.205){\fontsize{11.991}{1}\usefont{T1}{cmr}{m}{n}\selectfont\color{color_29791}ct}
\put(184.4742,-324.205){\fontsize{11.991}{1}\usefont{T1}{cmr}{m}{n}\selectfont\color{color_29791} }
\put(187.8557,-324.205){\fontsize{11.991}{1}\usefont{T1}{cmr}{m}{n}\selectfont\color{color_29791}t}
\put(191.1532,-324.205){\fontsize{11.991}{1}\usefont{T1}{cmr}{m}{n}\selectfont\color{color_29791}o}
\put(197.9402,-324.205){\fontsize{11.991}{1}\usefont{T1}{cmr}{m}{n}\selectfont\color{color_29791} }
\put(201.2377,-324.205){\fontsize{11.991}{1}\usefont{T1}{cmr}{m}{n}\selectfont\color{color_29791}H}
\put(209.9192,-324.205){\fontsize{11.991}{1}\usefont{T1}{cmr}{m}{n}\selectfont\color{color_29791}PC}
\put(226.5867,-324.205){\fontsize{11.991}{1}\usefont{T1}{cmr}{m}{n}\selectfont\color{color_29791} }
\put(229.8842,-324.205){\fontsize{11.991}{1}\usefont{T1}{cmr}{m}{n}\selectfont\color{color_29791}w}
\put(238.4818,-324.205){\fontsize{11.991}{1}\usefont{T1}{cmr}{m}{n}\selectfont\color{color_29791}o}
\put(245.2686,-324.205){\fontsize{11.991}{1}\usefont{T1}{cmr}{m}{n}\selectfont\color{color_29791}r}
\put(249.1657,-324.205){\fontsize{11.991}{1}\usefont{T1}{cmr}{m}{n}\selectfont\color{color_29791}kst}
\put(264.5382,-324.205){\fontsize{11.991}{1}\usefont{T1}{cmr}{m}{n}\selectfont\color{color_29791}a}
\put(271.2172,-324.205){\fontsize{11.991}{1}\usefont{T1}{cmr}{m}{n}\selectfont\color{color_29791}t}
\put(274.5986,-324.205){\fontsize{11.991}{1}\usefont{T1}{cmr}{m}{n}\selectfont\color{color_29791}i}
\put(277.2007,-324.205){\fontsize{11.991}{1}\usefont{T1}{cmr}{m}{n}\selectfont\color{color_29791}o}
\put(283.8797,-324.205){\fontsize{11.991}{1}\usefont{T1}{cmr}{m}{n}\selectfont\color{color_29791}n}
\put(290.6666,-324.205){\fontsize{11.991}{1}\usefont{T1}{cmr}{m}{n}\selectfont\color{color_29791}s:}
\put(299.9596,-324.205){\fontsize{11.991}{1}\usefont{T1}{cmr}{m}{n}\selectfont\color{color_29791} }
\put(303.3411,-324.205){\fontsize{11.991}{1}\usefont{T1}{cmr}{m}{n}\selectfont\color{color_29791}b}
\put(310.0201,-324.205){\fontsize{11.991}{1}\usefont{T1}{cmr}{m}{n}\selectfont\color{color_29791}sl}
\put(318.6176,-324.205){\fontsize{11.991}{1}\usefont{T1}{cmr}{m}{n}\selectfont\color{color_29791}w}
\put(327.2991,-324.205){\fontsize{11.991}{1}\usefont{T1}{cmr}{m}{n}\selectfont\color{color_29791}s0}
\put(339.9736,-324.205){\fontsize{11.991}{1}\usefont{T1}{cmr}{m}{n}\selectfont\color{color_29791}1}
\put(346.6526,-324.205){\fontsize{11.991}{1}\usefont{T1}{cmr}{m}{n}\selectfont\color{color_29791} }
\put(350.034,-324.205){\fontsize{11.991}{1}\usefont{T1}{cmr}{m}{n}\selectfont\color{color_29791}–}
\put(356.713,-324.205){\fontsize{11.991}{1}\usefont{T1}{cmr}{m}{n}\selectfont\color{color_29791} }
\put(360.0105,-324.205){\fontsize{11.991}{1}\usefont{T1}{cmr}{m}{n}\selectfont\color{color_29791}b}
\put(366.7974,-324.205){\fontsize{11.991}{1}\usefont{T1}{cmr}{m}{n}\selectfont\color{color_29791}sl}
\put(375.395,-324.205){\fontsize{11.991}{1}\usefont{T1}{cmr}{m}{n}\selectfont\color{color_29791}w}
\put(384.0764,-324.205){\fontsize{11.991}{1}\usefont{T1}{cmr}{m}{n}\selectfont\color{color_29791}s1}
\put(396.7509,-324.205){\fontsize{11.991}{1}\usefont{T1}{cmr}{m}{n}\selectfont\color{color_29791}2}
\put(84.213,-352.296){\fontsize{11.991}{1}\usefont{T1}{cmr}{m}{n}\selectfont\color{color_29791}•}
\put(98.187,-353.005){\fontsize{11.991}{1}\usefont{T1}{cmr}{m}{n}\selectfont\color{color_29791}O}
\put(107.564,-353.005){\fontsize{11.991}{1}\usefont{T1}{cmr}{m}{n}\selectfont\color{color_29791}p}
\put(114.243,-353.005){\fontsize{11.991}{1}\usefont{T1}{cmr}{m}{n}\selectfont\color{color_29791}e}
\put(120.9219,-353.005){\fontsize{11.991}{1}\usefont{T1}{cmr}{m}{n}\selectfont\color{color_29791}n}
\put(127.6009,-353.005){\fontsize{11.991}{1}\usefont{T1}{cmr}{m}{n}\selectfont\color{color_29791}S}
\put(135.6829,-353.005){\fontsize{11.991}{1}\usefont{T1}{cmr}{m}{n}\selectfont\color{color_29791}SH}
\put(152.2664,-353.005){\fontsize{11.991}{1}\usefont{T1}{cmr}{m}{n}\selectfont\color{color_29791} }
\put(155.6479,-353.005){\fontsize{11.991}{1}\usefont{T1}{cmr}{m}{n}\selectfont\color{color_29791}(a}
\put(166.3199,-353.005){\fontsize{11.991}{1}\usefont{T1}{cmr}{m}{n}\selectfont\color{color_29791}va}
\put(178.9944,-353.005){\fontsize{11.991}{1}\usefont{T1}{cmr}{m}{n}\selectfont\color{color_29791}i}
\put(181.6804,-353.005){\fontsize{11.991}{1}\usefont{T1}{cmr}{m}{n}\selectfont\color{color_29791}l}
\put(184.2824,-353.005){\fontsize{11.991}{1}\usefont{T1}{cmr}{m}{n}\selectfont\color{color_29791}a}
\put(190.9614,-353.005){\fontsize{11.991}{1}\usefont{T1}{cmr}{m}{n}\selectfont\color{color_29791}b}
\put(197.6404,-353.005){\fontsize{11.991}{1}\usefont{T1}{cmr}{m}{n}\selectfont\color{color_29791}l}
\put(200.3264,-353.005){\fontsize{11.991}{1}\usefont{T1}{cmr}{m}{n}\selectfont\color{color_29791}e}
\put(207.0054,-353.005){\fontsize{11.991}{1}\usefont{T1}{cmr}{m}{n}\selectfont\color{color_29791} }
\put(210.3868,-353.005){\fontsize{11.991}{1}\usefont{T1}{cmr}{m}{n}\selectfont\color{color_29791}f}
\put(213.6843,-353.005){\fontsize{11.991}{1}\usefont{T1}{cmr}{m}{n}\selectfont\color{color_29791}o}
\put(220.3633,-353.005){\fontsize{11.991}{1}\usefont{T1}{cmr}{m}{n}\selectfont\color{color_29791}r }
\put(227.6539,-353.005){\fontsize{11.991}{1}\usefont{T1}{cmr}{m}{n}\selectfont\color{color_29791}L}
\put(234.4408,-353.005){\fontsize{11.991}{1}\usefont{T1}{cmr}{m}{n}\selectfont\color{color_29791}i}
\put(237.0428,-353.005){\fontsize{11.991}{1}\usefont{T1}{cmr}{m}{n}\selectfont\color{color_29791}n}
\put(243.7218,-353.005){\fontsize{11.991}{1}\usefont{T1}{cmr}{m}{n}\selectfont\color{color_29791}u}
\put(250.5087,-353.005){\fontsize{11.991}{1}\usefont{T1}{cmr}{m}{n}\selectfont\color{color_29791}x}
\put(256.4083,-353.005){\fontsize{11.991}{1}\usefont{T1}{cmr}{m}{n}\selectfont\color{color_29791},}
\put(259.7897,-353.005){\fontsize{11.991}{1}\usefont{T1}{cmr}{m}{n}\selectfont\color{color_29791} }
\put(263.1712,-353.005){\fontsize{11.991}{1}\usefont{T1}{cmr}{m}{n}\selectfont\color{color_29791}M}
\put(273.0638,-353.005){\fontsize{11.991}{1}\usefont{T1}{cmr}{m}{n}\selectfont\color{color_29791}a}
\put(279.8507,-353.005){\fontsize{11.991}{1}\usefont{T1}{cmr}{m}{n}\selectfont\color{color_29791}c }
\put(289.1437,-353.005){\fontsize{11.991}{1}\usefont{T1}{cmr}{m}{n}\selectfont\color{color_29791}\&}
\put(297.2256,-353.005){\fontsize{11.991}{1}\usefont{T1}{cmr}{m}{n}\selectfont\color{color_29791} }
\put(300.5232,-353.005){\fontsize{11.991}{1}\usefont{T1}{cmr}{m}{n}\selectfont\color{color_29791}w}
\put(309.2047,-353.005){\fontsize{11.991}{1}\usefont{T1}{cmr}{m}{n}\selectfont\color{color_29791}i}
\put(311.8067,-353.005){\fontsize{11.991}{1}\usefont{T1}{cmr}{m}{n}\selectfont\color{color_29791}n}
\put(318.4857,-353.005){\fontsize{11.991}{1}\usefont{T1}{cmr}{m}{n}\selectfont\color{color_29791}d}
\put(325.1647,-353.005){\fontsize{11.991}{1}\usefont{T1}{cmr}{m}{n}\selectfont\color{color_29791}o}
\put(331.8437,-353.005){\fontsize{11.991}{1}\usefont{T1}{cmr}{m}{n}\selectfont\color{color_29791}w}
\put(340.5251,-353.005){\fontsize{11.991}{1}\usefont{T1}{cmr}{m}{n}\selectfont\color{color_29791}s),}
\put(354.411,-353.005){\fontsize{11.991}{1}\usefont{T1}{cmr}{m}{n}\selectfont\color{color_29791} }
\put(357.813,-353.005){\fontsize{11.991}{1}\usefont{T1}{cmr}{m}{n}\selectfont\color{color_29791}Pu}
\put(372.478,-353.005){\fontsize{11.991}{1}\usefont{T1}{cmr}{m}{n}\selectfont\color{color_29791}t}
\put(375.8594,-353.005){\fontsize{11.991}{1}\usefont{T1}{cmr}{m}{n}\selectfont\color{color_29791}t}
\put(379.157,-353.005){\fontsize{11.991}{1}\usefont{T1}{cmr}{m}{n}\selectfont\color{color_29791}y}
\put(384.3491,-353.005){\fontsize{11.991}{1}\usefont{T1}{cmr}{m}{n}\selectfont\color{color_29791},}
\put(387.6466,-353.005){\fontsize{11.991}{1}\usefont{T1}{cmr}{m}{n}\selectfont\color{color_29791} }
\put(390.9441,-353.005){\fontsize{11.991}{1}\usefont{T1}{cmr}{m}{n}\selectfont\color{color_29791}W}
\put(402.2396,-353.005){\fontsize{11.991}{1}\usefont{T1}{cmr}{m}{n}\selectfont\color{color_29791}SL}
\put(417.0125,-353.005){\fontsize{11.991}{1}\usefont{T1}{cmr}{m}{n}\selectfont\color{color_29791},}
\put(420.3101,-353.005){\fontsize{11.991}{1}\usefont{T1}{cmr}{m}{n}\selectfont\color{color_29791} }
\put(423.6915,-353.005){\fontsize{11.991}{1}\usefont{T1}{cmr}{m}{n}\selectfont\color{color_29791}M}
\put(433.5841,-353.005){\fontsize{11.991}{1}\usefont{T1}{cmr}{m}{n}\selectfont\color{color_29791}o}
\put(440.2631,-353.005){\fontsize{11.991}{1}\usefont{T1}{cmr}{m}{n}\selectfont\color{color_29791}b}
\put(447.05,-353.005){\fontsize{11.991}{1}\usefont{T1}{cmr}{m}{n}\selectfont\color{color_29791}a}
\put(453.729,-353.005){\fontsize{11.991}{1}\usefont{T1}{cmr}{m}{n}\selectfont\color{color_29791}Xt}
\put(465.0964,-353.005){\fontsize{11.991}{1}\usefont{T1}{cmr}{m}{n}\selectfont\color{color_29791}e}
\put(471.7754,-353.005){\fontsize{11.991}{1}\usefont{T1}{cmr}{m}{n}\selectfont\color{color_29791}rm}
\put(48.213,-386.595){\fontsize{15.987}{1}\usefont{T1}{cmr}{m}{n}\selectfont\color{color_29791}D}
\put(59.70766,-386.595){\fontsize{15.987}{1}\usefont{T1}{cmr}{m}{n}\selectfont\color{color_29791}em}
\put(81.88162,-386.595){\fontsize{15.987}{1}\usefont{T1}{cmr}{m}{n}\selectfont\color{color_29791}onst}
\put(112.129,-386.595){\fontsize{15.987}{1}\usefont{T1}{cmr}{m}{n}\selectfont\color{color_29791}r}
\put(117.4207,-386.595){\fontsize{15.987}{1}\usefont{T1}{cmr}{m}{n}\selectfont\color{color_29791}a}
\put(126.2136,-386.595){\fontsize{15.987}{1}\usefont{T1}{cmr}{m}{n}\selectfont\color{color_29791}t}
\put(130.6899,-386.595){\fontsize{15.987}{1}\usefont{T1}{cmr}{m}{n}\selectfont\color{color_29791}i}
\put(134.271,-386.595){\fontsize{15.987}{1}\usefont{T1}{cmr}{m}{n}\selectfont\color{color_29791}on}
\end{picture}
\newpage
\begin{tikzpicture}[overlay]
\path(0pt,0pt);
\filldraw[color_283006][even odd rule]
(-15pt, 10pt) -- (704.972pt, 10pt)
 -- (704.972pt, 10pt)
 -- (704.972pt, -529.972pt)
 -- (704.972pt, -529.972pt)
 -- (-15pt, -529.972pt)
 -- (-15pt, -529.972pt)
 -- (-15pt, 10pt) -- cycle
;
\end{tikzpicture}
\begin{picture}(-5,0)(2.5,0)
\put(-15,-529.972){\includegraphics[width=720pt,height=540pt]{latexImage_6f17c9ae575e6e3412bad9150a17b852.png}}
\put(48.411,-60.69601){\fontsize{21.997}{1}\usefont{T1}{cmr}{b}{n}\selectfont\color{color_29791}Ac}
\put(76.50117,-60.69601){\fontsize{21.997}{1}\usefont{T1}{cmr}{b}{n}\selectfont\color{color_29791}c}
\put(88.7095,-60.69601){\fontsize{21.997}{1}\usefont{T1}{cmr}{b}{n}\selectfont\color{color_29791}e}
\put(100.9838,-60.69601){\fontsize{21.997}{1}\usefont{T1}{cmr}{b}{n}\selectfont\color{color_29791}s}
\put(113.1922,-60.69601){\fontsize{21.997}{1}\usefont{T1}{cmr}{b}{n}\selectfont\color{color_29791}s}
\put(125.4665,-60.69601){\fontsize{21.997}{1}\usefont{T1}{cmr}{b}{n}\selectfont\color{color_29791} (}
\put(138.8627,-60.69601){\fontsize{21.997}{1}\usefont{T1}{cmr}{b}{n}\selectfont\color{color_29791}c}
\put(151.137,-60.69601){\fontsize{21.997}{1}\usefont{T1}{cmr}{b}{n}\selectfont\color{color_29791}o}
\put(164.5332,-60.69601){\fontsize{21.997}{1}\usefont{T1}{cmr}{b}{n}\selectfont\color{color_29791}n}
\put(177.9293,-60.69601){\fontsize{21.997}{1}\usefont{T1}{cmr}{b}{n}\selectfont\color{color_29791}t}
\put(185.2983,-60.69601){\fontsize{21.997}{1}\usefont{T1}{cmr}{b}{n}\selectfont\color{color_29791}in}
\put(204.7877,-60.69601){\fontsize{21.997}{1}\usefont{T1}{cmr}{b}{n}\selectfont\color{color_29791}u}
\put(218.1838,-60.69601){\fontsize{21.997}{1}\usefont{T1}{cmr}{b}{n}\selectfont\color{color_29791}e}
\put(230.4582,-60.69601){\fontsize{21.997}{1}\usefont{T1}{cmr}{b}{n}\selectfont\color{color_29791}d}
\put(243.8543,-60.69601){\fontsize{21.997}{1}\usefont{T1}{cmr}{b}{n}\selectfont\color{color_29791})}
\put(48.213,-116.397){\fontsize{15.987}{1}\usefont{T1}{cmr}{m}{n}\selectfont\color{color_29791}X}
\put(58.89232,-116.397){\fontsize{15.987}{1}\usefont{T1}{cmr}{m}{n}\selectfont\color{color_29791}2}
\put(67.68517,-116.397){\fontsize{15.987}{1}\usefont{T1}{cmr}{m}{n}\selectfont\color{color_29791}G}
\put(80.07509,-116.397){\fontsize{15.987}{1}\usefont{T1}{cmr}{m}{n}\selectfont\color{color_29791}o}
\put(84.213,-145.31){\fontsize{11.991}{1}\usefont{T1}{cmr}{m}{n}\selectfont\color{color_29791}•}
\put(106.691,-145.991){\fontsize{11.991}{1}\usefont{T1}{cmr}{m}{n}\selectfont\color{color_29791}Acce}
\put(133.4549,-145.991){\fontsize{11.991}{1}\usefont{T1}{cmr}{m}{n}\selectfont\color{color_29791}ss}
\put(145.35,-145.991){\fontsize{11.991}{1}\usefont{T1}{cmr}{m}{n}\selectfont\color{color_29791} }
\put(148.7314,-145.991){\fontsize{11.991}{1}\usefont{T1}{cmr}{m}{n}\selectfont\color{color_29791}H}
\put(157.4129,-145.991){\fontsize{11.991}{1}\usefont{T1}{cmr}{m}{n}\selectfont\color{color_29791}PC}
\put(173.9965,-145.991){\fontsize{11.991}{1}\usefont{T1}{cmr}{m}{n}\selectfont\color{color_29791} }
\put(177.378,-145.991){\fontsize{11.991}{1}\usefont{T1}{cmr}{m}{n}\selectfont\color{color_29791}d}
\put(184.0569,-145.991){\fontsize{11.991}{1}\usefont{T1}{cmr}{m}{n}\selectfont\color{color_29791}e}
\put(190.7359,-145.991){\fontsize{11.991}{1}\usefont{T1}{cmr}{m}{n}\selectfont\color{color_29791}skt}
\put(206.1084,-145.991){\fontsize{11.991}{1}\usefont{T1}{cmr}{m}{n}\selectfont\color{color_29791}o}
\put(212.7874,-145.991){\fontsize{11.991}{1}\usefont{T1}{cmr}{m}{n}\selectfont\color{color_29791}p}
\put(219.4664,-145.991){\fontsize{11.991}{1}\usefont{T1}{cmr}{m}{n}\selectfont\color{color_29791} }
\put(222.8478,-145.991){\fontsize{11.991}{1}\usefont{T1}{cmr}{m}{n}\selectfont\color{color_29791}i}
\put(225.4499,-145.991){\fontsize{11.991}{1}\usefont{T1}{cmr}{m}{n}\selectfont\color{color_29791}n}
\put(232.2368,-145.991){\fontsize{11.991}{1}\usefont{T1}{cmr}{m}{n}\selectfont\color{color_29791}t}
\put(235.5343,-145.991){\fontsize{11.991}{1}\usefont{T1}{cmr}{m}{n}\selectfont\color{color_29791}e}
\put(242.2133,-145.991){\fontsize{11.991}{1}\usefont{T1}{cmr}{m}{n}\selectfont\color{color_29791}rf}
\put(249.5038,-145.991){\fontsize{11.991}{1}\usefont{T1}{cmr}{m}{n}\selectfont\color{color_29791}a}
\put(256.2907,-145.991){\fontsize{11.991}{1}\usefont{T1}{cmr}{m}{n}\selectfont\color{color_29791}ce}
\put(268.9652,-145.991){\fontsize{11.991}{1}\usefont{T1}{cmr}{m}{n}\selectfont\color{color_29791} }
\put(272.3467,-145.991){\fontsize{11.991}{1}\usefont{T1}{cmr}{m}{n}\selectfont\color{color_29791}w}
\put(280.9442,-145.991){\fontsize{11.991}{1}\usefont{T1}{cmr}{m}{n}\selectfont\color{color_29791}i}
\put(283.5463,-145.991){\fontsize{11.991}{1}\usefont{T1}{cmr}{m}{n}\selectfont\color{color_29791}t}
\put(286.9278,-145.991){\fontsize{11.991}{1}\usefont{T1}{cmr}{m}{n}\selectfont\color{color_29791}h}
\put(293.6068,-145.991){\fontsize{11.991}{1}\usefont{T1}{cmr}{m}{n}\selectfont\color{color_29791} }
\put(296.9882,-145.991){\fontsize{11.991}{1}\usefont{T1}{cmr}{m}{n}\selectfont\color{color_29791}o}
\put(303.6672,-145.991){\fontsize{11.991}{1}\usefont{T1}{cmr}{m}{n}\selectfont\color{color_29791}r }
\put(310.9577,-145.991){\fontsize{11.991}{1}\usefont{T1}{cmr}{m}{n}\selectfont\color{color_29791}w}
\put(319.6392,-145.991){\fontsize{11.991}{1}\usefont{T1}{cmr}{m}{n}\selectfont\color{color_29791}i}
\put(322.2413,-145.991){\fontsize{11.991}{1}\usefont{T1}{cmr}{m}{n}\selectfont\color{color_29791}t}
\put(325.6227,-145.991){\fontsize{11.991}{1}\usefont{T1}{cmr}{m}{n}\selectfont\color{color_29791}h}
\put(332.3017,-145.991){\fontsize{11.991}{1}\usefont{T1}{cmr}{m}{n}\selectfont\color{color_29791}o}
\put(338.9807,-145.991){\fontsize{11.991}{1}\usefont{T1}{cmr}{m}{n}\selectfont\color{color_29791}u}
\put(345.6597,-145.991){\fontsize{11.991}{1}\usefont{T1}{cmr}{m}{n}\selectfont\color{color_29791}t}
\put(349.0412,-145.991){\fontsize{11.991}{1}\usefont{T1}{cmr}{m}{n}\selectfont\color{color_29791} }
\put(352.3387,-145.991){\fontsize{11.991}{1}\usefont{T1}{cmr}{m}{n}\selectfont\color{color_29791}V}
\put(360.4206,-145.991){\fontsize{11.991}{1}\usefont{T1}{cmr}{m}{n}\selectfont\color{color_29791}PN}
\put(377.0881,-145.991){\fontsize{11.991}{1}\usefont{T1}{cmr}{m}{n}\selectfont\color{color_29791} }
\put(380.3856,-145.991){\fontsize{11.991}{1}\usefont{T1}{cmr}{m}{n}\selectfont\color{color_29791}a}
\put(387.0646,-145.991){\fontsize{11.991}{1}\usefont{T1}{cmr}{m}{n}\selectfont\color{color_29791}cce}
\put(405.8425,-145.991){\fontsize{11.991}{1}\usefont{T1}{cmr}{m}{n}\selectfont\color{color_29791}ss}
\put(84.213,-174.11){\fontsize{11.991}{1}\usefont{T1}{cmr}{m}{n}\selectfont\color{color_29791}•}
\put(106.691,-174.791){\fontsize{11.991}{1}\usefont{T1}{cmr}{m}{n}\selectfont\color{color_29791}D}
\put(115.2885,-174.791){\fontsize{11.991}{1}\usefont{T1}{cmr}{m}{n}\selectfont\color{color_29791}i}
\put(117.9745,-174.791){\fontsize{11.991}{1}\usefont{T1}{cmr}{m}{n}\selectfont\color{color_29791}sco}
\put(136.6445,-174.791){\fontsize{11.991}{1}\usefont{T1}{cmr}{m}{n}\selectfont\color{color_29791}n}
\put(143.3235,-174.791){\fontsize{11.991}{1}\usefont{T1}{cmr}{m}{n}\selectfont\color{color_29791}n}
\put(150.0025,-174.791){\fontsize{11.991}{1}\usefont{T1}{cmr}{m}{n}\selectfont\color{color_29791}e}
\put(156.7894,-174.791){\fontsize{11.991}{1}\usefont{T1}{cmr}{m}{n}\selectfont\color{color_29791}c}
\put(162.689,-174.791){\fontsize{11.991}{1}\usefont{T1}{cmr}{m}{n}\selectfont\color{color_29791}t}
\put(166.0704,-174.791){\fontsize{11.991}{1}\usefont{T1}{cmr}{m}{n}\selectfont\color{color_29791}i}
\put(168.7564,-174.791){\fontsize{11.991}{1}\usefont{T1}{cmr}{m}{n}\selectfont\color{color_29791}n}
\put(175.4354,-174.791){\fontsize{11.991}{1}\usefont{T1}{cmr}{m}{n}\selectfont\color{color_29791}g}
\put(182.1144,-174.791){\fontsize{11.991}{1}\usefont{T1}{cmr}{m}{n}\selectfont\color{color_29791} }
\put(185.4958,-174.791){\fontsize{11.991}{1}\usefont{T1}{cmr}{m}{n}\selectfont\color{color_29791}a}
\put(192.1748,-174.791){\fontsize{11.991}{1}\usefont{T1}{cmr}{m}{n}\selectfont\color{color_29791}n}
\put(198.8538,-174.791){\fontsize{11.991}{1}\usefont{T1}{cmr}{m}{n}\selectfont\color{color_29791}d}
\put(205.5328,-174.791){\fontsize{11.991}{1}\usefont{T1}{cmr}{m}{n}\selectfont\color{color_29791} }
\put(208.9143,-174.791){\fontsize{11.991}{1}\usefont{T1}{cmr}{m}{n}\selectfont\color{color_29791}r}
\put(212.8114,-174.791){\fontsize{11.991}{1}\usefont{T1}{cmr}{m}{n}\selectfont\color{color_29791}e}
\put(219.5983,-174.791){\fontsize{11.991}{1}\usefont{T1}{cmr}{m}{n}\selectfont\color{color_29791}co}
\put(232.2728,-174.791){\fontsize{11.991}{1}\usefont{T1}{cmr}{m}{n}\selectfont\color{color_29791}n}
\put(238.9517,-174.791){\fontsize{11.991}{1}\usefont{T1}{cmr}{m}{n}\selectfont\color{color_29791}n}
\put(245.6307,-174.791){\fontsize{11.991}{1}\usefont{T1}{cmr}{m}{n}\selectfont\color{color_29791}e}
\put(252.3097,-174.791){\fontsize{11.991}{1}\usefont{T1}{cmr}{m}{n}\selectfont\color{color_29791}ct}
\put(261.6867,-174.791){\fontsize{11.991}{1}\usefont{T1}{cmr}{m}{n}\selectfont\color{color_29791}i}
\put(264.2887,-174.791){\fontsize{11.991}{1}\usefont{T1}{cmr}{m}{n}\selectfont\color{color_29791}n}
\put(270.9677,-174.791){\fontsize{11.991}{1}\usefont{T1}{cmr}{m}{n}\selectfont\color{color_29791}g}
\put(84.213,-202.91){\fontsize{11.991}{1}\usefont{T1}{cmr}{m}{n}\selectfont\color{color_29791}•}
\put(106.691,-203.506){\fontsize{11.991}{1}\usefont{T1}{cmr}{m}{n}\selectfont\color{color_29791}C}
\put(115.2885,-203.506){\fontsize{11.991}{1}\usefont{T1}{cmr}{m}{n}\selectfont\color{color_29791}o}
\put(122.0755,-203.506){\fontsize{11.991}{1}\usefont{T1}{cmr}{m}{n}\selectfont\color{color_29791}p}
\put(128.7544,-203.506){\fontsize{11.991}{1}\usefont{T1}{cmr}{m}{n}\selectfont\color{color_29791}y/}
\put(138.0475,-203.506){\fontsize{11.991}{1}\usefont{T1}{cmr}{m}{n}\selectfont\color{color_29791}p}
\put(144.7265,-203.506){\fontsize{11.991}{1}\usefont{T1}{cmr}{m}{n}\selectfont\color{color_29791}a}
\put(151.5134,-203.506){\fontsize{11.991}{1}\usefont{T1}{cmr}{m}{n}\selectfont\color{color_29791}st}
\put(160.8064,-203.506){\fontsize{11.991}{1}\usefont{T1}{cmr}{m}{n}\selectfont\color{color_29791}e}
\put(84.213,-231.597){\fontsize{11.991}{1}\usefont{T1}{cmr}{m}{n}\selectfont\color{color_29791}•}
\put(106.691,-232.306){\fontsize{11.991}{1}\usefont{T1}{cmr}{m}{n}\selectfont\color{color_29791}Sh}
\put(121.4639,-232.306){\fontsize{11.991}{1}\usefont{T1}{cmr}{m}{n}\selectfont\color{color_29791}a}
\put(128.1429,-232.306){\fontsize{11.991}{1}\usefont{T1}{cmr}{m}{n}\selectfont\color{color_29791}r}
\put(132.04,-232.306){\fontsize{11.991}{1}\usefont{T1}{cmr}{m}{n}\selectfont\color{color_29791}i}
\put(134.726,-232.306){\fontsize{11.991}{1}\usefont{T1}{cmr}{m}{n}\selectfont\color{color_29791}n}
\put(141.405,-232.306){\fontsize{11.991}{1}\usefont{T1}{cmr}{m}{n}\selectfont\color{color_29791}g}
\put(148.0839,-232.306){\fontsize{11.991}{1}\usefont{T1}{cmr}{m}{n}\selectfont\color{color_29791} }
\put(151.4654,-232.306){\fontsize{11.991}{1}\usefont{T1}{cmr}{m}{n}\selectfont\color{color_29791}f}
\put(154.7629,-232.306){\fontsize{11.991}{1}\usefont{T1}{cmr}{m}{n}\selectfont\color{color_29791}i}
\put(157.4489,-232.306){\fontsize{11.991}{1}\usefont{T1}{cmr}{m}{n}\selectfont\color{color_29791}l}
\put(160.051,-232.306){\fontsize{11.991}{1}\usefont{T1}{cmr}{m}{n}\selectfont\color{color_29791}e}
\put(166.7299,-232.306){\fontsize{11.991}{1}\usefont{T1}{cmr}{m}{n}\selectfont\color{color_29791}s }
\put(176.1069,-232.306){\fontsize{11.991}{1}\usefont{T1}{cmr}{m}{n}\selectfont\color{color_29791}f}
\put(179.4884,-232.306){\fontsize{11.991}{1}\usefont{T1}{cmr}{m}{n}\selectfont\color{color_29791}r}
\put(183.3855,-232.306){\fontsize{11.991}{1}\usefont{T1}{cmr}{m}{n}\selectfont\color{color_29791}o}
\put(190.0644,-232.306){\fontsize{11.991}{1}\usefont{T1}{cmr}{m}{n}\selectfont\color{color_29791}m }
\put(203.4344,-232.306){\fontsize{11.991}{1}\usefont{T1}{cmr}{m}{n}\selectfont\color{color_29791}y}
\put(209.334,-232.306){\fontsize{11.991}{1}\usefont{T1}{cmr}{m}{n}\selectfont\color{color_29791}o}
\put(216.1209,-232.306){\fontsize{11.991}{1}\usefont{T1}{cmr}{m}{n}\selectfont\color{color_29791}u}
\put(222.7999,-232.306){\fontsize{11.991}{1}\usefont{T1}{cmr}{m}{n}\selectfont\color{color_29791}r }
\put(230.0904,-232.306){\fontsize{11.991}{1}\usefont{T1}{cmr}{m}{n}\selectfont\color{color_29791}l}
\put(232.7764,-232.306){\fontsize{11.991}{1}\usefont{T1}{cmr}{m}{n}\selectfont\color{color_29791}a}
\put(239.4554,-232.306){\fontsize{11.991}{1}\usefont{T1}{cmr}{m}{n}\selectfont\color{color_29791}p}
\put(246.1344,-232.306){\fontsize{11.991}{1}\usefont{T1}{cmr}{m}{n}\selectfont\color{color_29791}t}
\put(249.4319,-232.306){\fontsize{11.991}{1}\usefont{T1}{cmr}{m}{n}\selectfont\color{color_29791}o}
\put(256.2188,-232.306){\fontsize{11.991}{1}\usefont{T1}{cmr}{m}{n}\selectfont\color{color_29791}p}
\put(262.8978,-232.306){\fontsize{11.991}{1}\usefont{T1}{cmr}{m}{n}\selectfont\color{color_29791} }
\put(266.2792,-232.306){\fontsize{11.991}{1}\usefont{T1}{cmr}{m}{n}\selectfont\color{color_29791}o}
\put(272.9582,-232.306){\fontsize{11.991}{1}\usefont{T1}{cmr}{m}{n}\selectfont\color{color_29791}r}
\put(276.8553,-232.306){\fontsize{11.991}{1}\usefont{T1}{cmr}{m}{n}\selectfont\color{color_29791} }
\put(280.2368,-232.306){\fontsize{11.991}{1}\usefont{T1}{cmr}{m}{n}\selectfont\color{color_29791}PC}
\put(84.213,-260.397){\fontsize{11.991}{1}\usefont{T1}{cmr}{m}{n}\selectfont\color{color_29791}•}
\put(106.691,-261.106){\fontsize{11.991}{1}\usefont{T1}{cmr}{m}{n}\selectfont\color{color_29791}M}
\put(116.5836,-261.106){\fontsize{11.991}{1}\usefont{T1}{cmr}{m}{n}\selectfont\color{color_29791}o}
\put(123.3705,-261.106){\fontsize{11.991}{1}\usefont{T1}{cmr}{m}{n}\selectfont\color{color_29791}r}
\put(127.2676,-261.106){\fontsize{11.991}{1}\usefont{T1}{cmr}{m}{n}\selectfont\color{color_29791}e}
\put(133.9465,-261.106){\fontsize{11.991}{1}\usefont{T1}{cmr}{m}{n}\selectfont\color{color_29791} }
\put(137.328,-261.106){\fontsize{11.991}{1}\usefont{T1}{cmr}{m}{n}\selectfont\color{color_29791}i}
\put(140.014,-261.106){\fontsize{11.991}{1}\usefont{T1}{cmr}{m}{n}\selectfont\color{color_29791}n}
\put(146.693,-261.106){\fontsize{11.991}{1}\usefont{T1}{cmr}{m}{n}\selectfont\color{color_29791}f}
\put(149.9905,-261.106){\fontsize{11.991}{1}\usefont{T1}{cmr}{m}{n}\selectfont\color{color_29791}o}
\put(156.7774,-261.106){\fontsize{11.991}{1}\usefont{T1}{cmr}{m}{n}\selectfont\color{color_29791}r}
\put(160.6745,-261.106){\fontsize{11.991}{1}\usefont{T1}{cmr}{m}{n}\selectfont\color{color_29791}ma}
\put(177.342,-261.106){\fontsize{11.991}{1}\usefont{T1}{cmr}{m}{n}\selectfont\color{color_29791}t}
\put(180.6395,-261.106){\fontsize{11.991}{1}\usefont{T1}{cmr}{m}{n}\selectfont\color{color_29791}i}
\put(183.3255,-261.106){\fontsize{11.991}{1}\usefont{T1}{cmr}{m}{n}\selectfont\color{color_29791}o}
\put(190.0045,-261.106){\fontsize{11.991}{1}\usefont{T1}{cmr}{m}{n}\selectfont\color{color_29791}n}
\put(196.6835,-261.106){\fontsize{11.991}{1}\usefont{T1}{cmr}{m}{n}\selectfont\color{color_29791}:}
\put(200.0649,-261.106){\fontsize{11.991}{1}\usefont{T1}{cmr}{m}{n}\selectfont\color{color_29791} }
\put(203.693,-261.106){\fontsize{11.991}{1}\usefont{T1}{cmr}{m}{n}\selectfont\color{color_167934}h}
\put(210.372,-261.106){\fontsize{11.991}{1}\usefont{T1}{cmr}{m}{n}\selectfont\color{color_167934}t}
\put(213.7534,-261.106){\fontsize{11.991}{1}\usefont{T1}{cmr}{m}{n}\selectfont\color{color_167934}t}
\put(217.051,-261.106){\fontsize{11.991}{1}\usefont{T1}{cmr}{m}{n}\selectfont\color{color_167934}p}
\put(223.73,-261.106){\fontsize{11.991}{1}\usefont{T1}{cmr}{m}{n}\selectfont\color{color_167934}:}
\put(227.1114,-261.106){\fontsize{11.991}{1}\usefont{T1}{cmr}{m}{n}\selectfont\color{color_167934}/}
\put(230.409,-261.106){\fontsize{11.991}{1}\usefont{T1}{cmr}{m}{n}\selectfont\color{color_167934}/}
\put(233.7904,-261.106){\fontsize{11.991}{1}\usefont{T1}{cmr}{m}{n}\selectfont\color{color_167934}i}
\put(236.3925,-261.106){\fontsize{11.991}{1}\usefont{T1}{cmr}{m}{n}\selectfont\color{color_167934}c}
\put(242.4719,-261.106){\fontsize{11.991}{1}\usefont{T1}{cmr}{m}{n}\selectfont\color{color_167934}t}
\put(245.7694,-261.106){\fontsize{11.991}{1}\usefont{T1}{cmr}{m}{n}\selectfont\color{color_167934}d}
\put(252.4484,-261.106){\fontsize{11.991}{1}\usefont{T1}{cmr}{m}{n}\selectfont\color{color_167934}o}
\put(259.1274,-261.106){\fontsize{11.991}{1}\usefont{T1}{cmr}{m}{n}\selectfont\color{color_167934}cs}
\put(271.2023,-261.106){\fontsize{11.991}{1}\usefont{T1}{cmr}{m}{n}\selectfont\color{color_167934}/}
\put(274.4998,-261.106){\fontsize{11.991}{1}\usefont{T1}{cmr}{m}{n}\selectfont\color{color_167934}w}
\put(283.0974,-261.106){\fontsize{11.991}{1}\usefont{T1}{cmr}{m}{n}\selectfont\color{color_167934}i}
\put(285.7834,-261.106){\fontsize{11.991}{1}\usefont{T1}{cmr}{m}{n}\selectfont\color{color_167934}ki}
\put(294.3809,-261.106){\fontsize{11.991}{1}\usefont{T1}{cmr}{m}{n}\selectfont\color{color_167934}/}
\put(297.7624,-261.106){\fontsize{11.991}{1}\usefont{T1}{cmr}{m}{n}\selectfont\color{color_167934}i}
\put(300.3644,-261.106){\fontsize{11.991}{1}\usefont{T1}{cmr}{m}{n}\selectfont\color{color_167934}n}
\put(307.1513,-261.106){\fontsize{11.991}{1}\usefont{T1}{cmr}{m}{n}\selectfont\color{color_167934}d}
\put(313.8303,-261.106){\fontsize{11.991}{1}\usefont{T1}{cmr}{m}{n}\selectfont\color{color_167934}e}
\put(320.5093,-261.106){\fontsize{11.991}{1}\usefont{T1}{cmr}{m}{n}\selectfont\color{color_167934}x.}
\put(329.8023,-261.106){\fontsize{11.991}{1}\usefont{T1}{cmr}{m}{n}\selectfont\color{color_167934}p}
\put(336.5892,-261.106){\fontsize{11.991}{1}\usefont{T1}{cmr}{m}{n}\selectfont\color{color_167934}h}
\put(343.2682,-261.106){\fontsize{11.991}{1}\usefont{T1}{cmr}{m}{n}\selectfont\color{color_167934}p}
\put(349.9472,-261.106){\fontsize{11.991}{1}\usefont{T1}{cmr}{m}{n}\selectfont\color{color_167934}/}
\put(353.3286,-261.106){\fontsize{11.991}{1}\usefont{T1}{cmr}{m}{n}\selectfont\color{color_167934}H}
\put(361.9262,-261.106){\fontsize{11.991}{1}\usefont{T1}{cmr}{m}{n}\selectfont\color{color_167934}PC}
\put(378.5937,-261.106){\fontsize{11.991}{1}\usefont{T1}{cmr}{m}{n}\selectfont\color{color_167934}:}
\put(381.8912,-261.106){\fontsize{11.991}{1}\usefont{T1}{cmr}{m}{n}\selectfont\color{color_167934}X}
\put(389.9731,-261.106){\fontsize{11.991}{1}\usefont{T1}{cmr}{m}{n}\selectfont\color{color_167934}2}
\put(396.6521,-261.106){\fontsize{11.991}{1}\usefont{T1}{cmr}{m}{n}\selectfont\color{color_167934}G}
\put(406.0291,-261.106){\fontsize{11.991}{1}\usefont{T1}{cmr}{m}{n}\selectfont\color{color_167934}O}
\end{picture}
\begin{tikzpicture}[overlay]
\path(0pt,0pt);
\draw[color_167934,line width=0.595pt]
(203.693pt, -262.098pt) -- (415.781pt, -262.098pt)
;
\end{tikzpicture}
\begin{picture}(-5,0)(2.5,0)
\put(84.213,-289.197){\fontsize{11.991}{1}\usefont{T1}{cmr}{m}{n}\selectfont\color{color_29791}•}
\put(106.691,-289.792){\fontsize{11.991}{1}\usefont{T1}{cmr}{m}{n}\selectfont\color{color_29791}D}
\put(115.2885,-289.792){\fontsize{11.991}{1}\usefont{T1}{cmr}{m}{n}\selectfont\color{color_29791}e}
\put(122.0755,-289.792){\fontsize{11.991}{1}\usefont{T1}{cmr}{m}{n}\selectfont\color{color_29791}m}
\put(131.968,-289.792){\fontsize{11.991}{1}\usefont{T1}{cmr}{m}{n}\selectfont\color{color_29791}o}
\put(138.647,-289.792){\fontsize{11.991}{1}\usefont{T1}{cmr}{m}{n}\selectfont\color{color_29791}n}
\put(145.326,-289.792){\fontsize{11.991}{1}\usefont{T1}{cmr}{m}{n}\selectfont\color{color_29791}s}
\put(151.4054,-289.792){\fontsize{11.991}{1}\usefont{T1}{cmr}{m}{n}\selectfont\color{color_29791}t}
\put(154.703,-289.792){\fontsize{11.991}{1}\usefont{T1}{cmr}{m}{n}\selectfont\color{color_29791}ra}
\put(165.375,-289.792){\fontsize{11.991}{1}\usefont{T1}{cmr}{m}{n}\selectfont\color{color_29791}t}
\put(168.7564,-289.792){\fontsize{11.991}{1}\usefont{T1}{cmr}{m}{n}\selectfont\color{color_29791}i}
\put(171.3585,-289.792){\fontsize{11.991}{1}\usefont{T1}{cmr}{m}{n}\selectfont\color{color_29791}o}
\put(178.0375,-289.792){\fontsize{11.991}{1}\usefont{T1}{cmr}{m}{n}\selectfont\color{color_29791}n}
\put(48.213,-323.411){\fontsize{15.987}{1}\usefont{T1}{cmr}{m}{n}\selectfont\color{color_29791}O}
\put(60.507,-323.411){\fontsize{15.987}{1}\usefont{T1}{cmr}{m}{n}\selectfont\color{color_29791}t}
\put(64.98336,-323.411){\fontsize{15.987}{1}\usefont{T1}{cmr}{m}{n}\selectfont\color{color_29791}he}
\put(82.66498,-323.411){\fontsize{15.987}{1}\usefont{T1}{cmr}{m}{n}\selectfont\color{color_29791}r}
\put(87.95667,-323.411){\fontsize{15.987}{1}\usefont{T1}{cmr}{m}{n}\selectfont\color{color_29791} }
\put(92.43303,-323.411){\fontsize{15.987}{1}\usefont{T1}{cmr}{m}{n}\selectfont\color{color_29791}O}
\put(104.727,-323.411){\fontsize{15.987}{1}\usefont{T1}{cmr}{m}{n}\selectfont\color{color_29791}pt}
\put(118.0122,-323.411){\fontsize{15.987}{1}\usefont{T1}{cmr}{m}{n}\selectfont\color{color_29791}i}
\put(121.5933,-323.411){\fontsize{15.987}{1}\usefont{T1}{cmr}{m}{n}\selectfont\color{color_29791}ons}
\put(84.213,-352.296){\fontsize{11.991}{1}\usefont{T1}{cmr}{m}{n}\selectfont\color{color_29791}•}
\put(106.691,-353.005){\fontsize{11.991}{1}\usefont{T1}{cmr}{m}{n}\selectfont\color{color_29791}Exce}
\put(133.4549,-353.005){\fontsize{11.991}{1}\usefont{T1}{cmr}{m}{n}\selectfont\color{color_29791}e}
\put(140.1339,-353.005){\fontsize{11.991}{1}\usefont{T1}{cmr}{m}{n}\selectfont\color{color_29791}d}
\put(146.8129,-353.005){\fontsize{11.991}{1}\usefont{T1}{cmr}{m}{n}\selectfont\color{color_29791} }
\put(150.1943,-353.005){\fontsize{11.991}{1}\usefont{T1}{cmr}{m}{n}\selectfont\color{color_29791}/}
\put(153.4919,-353.005){\fontsize{11.991}{1}\usefont{T1}{cmr}{m}{n}\selectfont\color{color_29791} }
\put(156.8733,-353.005){\fontsize{11.991}{1}\usefont{T1}{cmr}{m}{n}\selectfont\color{color_29791}XMi}
\put(177.4499,-353.005){\fontsize{11.991}{1}\usefont{T1}{cmr}{m}{n}\selectfont\color{color_29791}n}
\put(184.1289,-353.005){\fontsize{11.991}{1}\usefont{T1}{cmr}{m}{n}\selectfont\color{color_29791}g}
\put(84.213,-381.096){\fontsize{11.991}{1}\usefont{T1}{cmr}{m}{n}\selectfont\color{color_29791}•}
\put(106.691,-381.691){\fontsize{11.991}{1}\usefont{T1}{cmr}{m}{n}\selectfont\color{color_29791}M}
\put(116.5836,-381.691){\fontsize{11.991}{1}\usefont{T1}{cmr}{m}{n}\selectfont\color{color_29791}o}
\put(123.3705,-381.691){\fontsize{11.991}{1}\usefont{T1}{cmr}{m}{n}\selectfont\color{color_29791}b}
\put(130.0495,-381.691){\fontsize{11.991}{1}\usefont{T1}{cmr}{m}{n}\selectfont\color{color_29791}a}
\put(136.7285,-381.691){\fontsize{11.991}{1}\usefont{T1}{cmr}{m}{n}\selectfont\color{color_29791}Xt}
\put(148.0959,-381.691){\fontsize{11.991}{1}\usefont{T1}{cmr}{m}{n}\selectfont\color{color_29791}e}
\put(154.7749,-381.691){\fontsize{11.991}{1}\usefont{T1}{cmr}{m}{n}\selectfont\color{color_29791}rm}
\put(84.213,-409.811){\fontsize{11.991}{1}\usefont{T1}{cmr}{m}{n}\selectfont\color{color_29791}•}
\put(106.691,-410.491){\fontsize{11.991}{1}\usefont{T1}{cmr}{m}{n}\selectfont\color{color_29791}D}
\put(115.2885,-410.491){\fontsize{11.991}{1}\usefont{T1}{cmr}{m}{n}\selectfont\color{color_29791}e}
\put(122.0755,-410.491){\fontsize{11.991}{1}\usefont{T1}{cmr}{m}{n}\selectfont\color{color_29791}m}
\put(131.968,-410.491){\fontsize{11.991}{1}\usefont{T1}{cmr}{m}{n}\selectfont\color{color_29791}o}
\put(138.647,-410.491){\fontsize{11.991}{1}\usefont{T1}{cmr}{m}{n}\selectfont\color{color_29791}n}
\put(145.326,-410.491){\fontsize{11.991}{1}\usefont{T1}{cmr}{m}{n}\selectfont\color{color_29791}s}
\put(151.4054,-410.491){\fontsize{11.991}{1}\usefont{T1}{cmr}{m}{n}\selectfont\color{color_29791}t}
\put(154.703,-410.491){\fontsize{11.991}{1}\usefont{T1}{cmr}{m}{n}\selectfont\color{color_29791}ra}
\put(165.375,-410.491){\fontsize{11.991}{1}\usefont{T1}{cmr}{m}{n}\selectfont\color{color_29791}t}
\put(168.7564,-410.491){\fontsize{11.991}{1}\usefont{T1}{cmr}{m}{n}\selectfont\color{color_29791}i}
\put(171.3585,-410.491){\fontsize{11.991}{1}\usefont{T1}{cmr}{m}{n}\selectfont\color{color_29791}o}
\put(178.0375,-410.491){\fontsize{11.991}{1}\usefont{T1}{cmr}{m}{n}\selectfont\color{color_29791}n}
\end{picture}
\newpage
\begin{tikzpicture}[overlay]
\path(0pt,0pt);
\filldraw[color_283006][even odd rule]
(-15pt, 10pt) -- (704.972pt, 10pt)
 -- (704.972pt, 10pt)
 -- (704.972pt, -529.972pt)
 -- (704.972pt, -529.972pt)
 -- (-15pt, -529.972pt)
 -- (-15pt, -529.972pt)
 -- (-15pt, 10pt) -- cycle
;
\end{tikzpicture}
\begin{picture}(-5,0)(2.5,0)
\put(-15,-529.972){\includegraphics[width=720pt,height=540pt]{latexImage_2828171a63492229494bbe96f89381b1.png}}
\put(48.411,-60.69601){\fontsize{21.997}{1}\usefont{T1}{cmr}{b}{n}\selectfont\color{color_29791}S}
\put(63.083,-60.69601){\fontsize{21.997}{1}\usefont{T1}{cmr}{b}{n}\selectfont\color{color_29791}t}
\put(70.386,-60.69601){\fontsize{21.997}{1}\usefont{T1}{cmr}{b}{n}\selectfont\color{color_29791}o}
\put(83.78217,-60.69601){\fontsize{21.997}{1}\usefont{T1}{cmr}{b}{n}\selectfont\color{color_29791}r}
\put(92.27302,-60.69601){\fontsize{21.997}{1}\usefont{T1}{cmr}{b}{n}\selectfont\color{color_29791}a}
\put(104.5473,-60.69601){\fontsize{21.997}{1}\usefont{T1}{cmr}{b}{n}\selectfont\color{color_29791}g}
\put(117.9435,-60.69601){\fontsize{21.997}{1}\usefont{T1}{cmr}{b}{n}\selectfont\color{color_29791}e}
\put(48.213,-116.51){\fontsize{15.987}{1}\usefont{T1}{cmr}{m}{n}\selectfont\color{color_29791}U}
\put(59.70766,-116.51){\fontsize{15.987}{1}\usefont{T1}{cmr}{m}{n}\selectfont\color{color_29791}s}
\put(67.78109,-116.51){\fontsize{15.987}{1}\usefont{T1}{cmr}{m}{n}\selectfont\color{color_29791}e}
\put(76.57394,-116.51){\fontsize{15.987}{1}\usefont{T1}{cmr}{m}{n}\selectfont\color{color_29791}r}
\put(81.86563,-116.51){\fontsize{15.987}{1}\usefont{T1}{cmr}{m}{n}\selectfont\color{color_29791} }
\put(85.46271,-116.51){\fontsize{15.987}{1}\usefont{T1}{cmr}{m}{n}\selectfont\color{color_29791}A}
\put(96.0461,-116.51){\fontsize{15.987}{1}\usefont{T1}{cmr}{m}{n}\selectfont\color{color_29791}r}
\put(101.3378,-116.51){\fontsize{15.987}{1}\usefont{T1}{cmr}{m}{n}\selectfont\color{color_29791}ea}
\put(119.0194,-116.51){\fontsize{15.987}{1}\usefont{T1}{cmr}{m}{n}\selectfont\color{color_29791} }
\put(123.4958,-116.51){\fontsize{15.987}{1}\usefont{T1}{cmr}{m}{n}\selectfont\color{color_29791}-}
\put(128.7875,-116.51){\fontsize{15.987}{1}\usefont{T1}{cmr}{m}{n}\selectfont\color{color_29791} }
\put(133.1839,-116.51){\fontsize{15.987}{1}\usefont{T1}{cmr}{m}{n}\selectfont\color{color_29791}/}
\put(137.5803,-116.51){\fontsize{15.987}{1}\usefont{T1}{cmr}{m}{n}\selectfont\color{color_29791}us}
\put(154.5425,-116.51){\fontsize{15.987}{1}\usefont{T1}{cmr}{m}{n}\selectfont\color{color_29791}e}
\put(163.3354,-116.51){\fontsize{15.987}{1}\usefont{T1}{cmr}{m}{n}\selectfont\color{color_29791}r}
\put(168.6271,-116.51){\fontsize{15.987}{1}\usefont{T1}{cmr}{m}{n}\selectfont\color{color_29791}s/}
\put(181.097,-116.51){\fontsize{15.987}{1}\usefont{T1}{cmr}{m}{n}\selectfont\color{color_29791}user}
\put(212.1597,-116.51){\fontsize{15.987}{1}\usefont{T1}{cmr}{m}{n}\selectfont\color{color_29791}n}
\put(220.9525,-116.51){\fontsize{15.987}{1}\usefont{T1}{cmr}{m}{n}\selectfont\color{color_29791}am}
\put(243.2224,-116.51){\fontsize{15.987}{1}\usefont{T1}{cmr}{m}{n}\selectfont\color{color_29791}e}
\put(84.213,-145.31){\fontsize{11.991}{1}\usefont{T1}{cmr}{m}{n}\selectfont\color{color_29791}•}
\put(98.187,-145.991){\fontsize{11.991}{1}\usefont{T1}{cmr}{m}{n}\selectfont\color{color_29791}S}
\put(106.2689,-145.991){\fontsize{11.991}{1}\usefont{T1}{cmr}{m}{n}\selectfont\color{color_29791}m}
\put(116.1615,-145.991){\fontsize{11.991}{1}\usefont{T1}{cmr}{m}{n}\selectfont\color{color_29791}a}
\put(122.8405,-145.991){\fontsize{11.991}{1}\usefont{T1}{cmr}{m}{n}\selectfont\color{color_29791}l}
\put(125.5265,-145.991){\fontsize{11.991}{1}\usefont{T1}{cmr}{m}{n}\selectfont\color{color_29791}l}
\put(128.1285,-145.991){\fontsize{11.991}{1}\usefont{T1}{cmr}{m}{n}\selectfont\color{color_29791},}
\put(131.51,-145.991){\fontsize{11.991}{1}\usefont{T1}{cmr}{m}{n}\selectfont\color{color_29791} }
\put(134.8075,-145.991){\fontsize{11.991}{1}\usefont{T1}{cmr}{m}{n}\selectfont\color{color_29791}n}
\put(141.4865,-145.991){\fontsize{11.991}{1}\usefont{T1}{cmr}{m}{n}\selectfont\color{color_29791}o}
\put(148.2734,-145.991){\fontsize{11.991}{1}\usefont{T1}{cmr}{m}{n}\selectfont\color{color_29791}t}
\put(151.5709,-145.991){\fontsize{11.991}{1}\usefont{T1}{cmr}{m}{n}\selectfont\color{color_29791} }
\put(154.9524,-145.991){\fontsize{11.991}{1}\usefont{T1}{cmr}{m}{n}\selectfont\color{color_29791}i}
\put(157.5544,-145.991){\fontsize{11.991}{1}\usefont{T1}{cmr}{m}{n}\selectfont\color{color_29791}n}
\put(164.3414,-145.991){\fontsize{11.991}{1}\usefont{T1}{cmr}{m}{n}\selectfont\color{color_29791}t}
\put(167.6389,-145.991){\fontsize{11.991}{1}\usefont{T1}{cmr}{m}{n}\selectfont\color{color_29791}e}
\put(174.3179,-145.991){\fontsize{11.991}{1}\usefont{T1}{cmr}{m}{n}\selectfont\color{color_29791}n}
\put(180.9969,-145.991){\fontsize{11.991}{1}\usefont{T1}{cmr}{m}{n}\selectfont\color{color_29791}d}
\put(187.7838,-145.991){\fontsize{11.991}{1}\usefont{T1}{cmr}{m}{n}\selectfont\color{color_29791}e}
\put(194.4628,-145.991){\fontsize{11.991}{1}\usefont{T1}{cmr}{m}{n}\selectfont\color{color_29791}d}
\put(201.1417,-145.991){\fontsize{11.991}{1}\usefont{T1}{cmr}{m}{n}\selectfont\color{color_29791} }
\put(204.5232,-145.991){\fontsize{11.991}{1}\usefont{T1}{cmr}{m}{n}\selectfont\color{color_29791}f}
\put(207.8207,-145.991){\fontsize{11.991}{1}\usefont{T1}{cmr}{m}{n}\selectfont\color{color_29791}o}
\put(214.6076,-145.991){\fontsize{11.991}{1}\usefont{T1}{cmr}{m}{n}\selectfont\color{color_29791}r}
\put(218.5047,-145.991){\fontsize{11.991}{1}\usefont{T1}{cmr}{m}{n}\selectfont\color{color_29791} }
\put(221.8862,-145.991){\fontsize{11.991}{1}\usefont{T1}{cmr}{m}{n}\selectfont\color{color_29791}sh}
\put(234.5607,-145.991){\fontsize{11.991}{1}\usefont{T1}{cmr}{m}{n}\selectfont\color{color_29791}a}
\put(241.2397,-145.991){\fontsize{11.991}{1}\usefont{T1}{cmr}{m}{n}\selectfont\color{color_29791}ri}
\put(247.8347,-145.991){\fontsize{11.991}{1}\usefont{T1}{cmr}{m}{n}\selectfont\color{color_29791}n}
\put(254.5137,-145.991){\fontsize{11.991}{1}\usefont{T1}{cmr}{m}{n}\selectfont\color{color_29791}g}
\put(261.3006,-145.991){\fontsize{11.991}{1}\usefont{T1}{cmr}{m}{n}\selectfont\color{color_29791} }
\put(264.5981,-145.991){\fontsize{11.991}{1}\usefont{T1}{cmr}{m}{n}\selectfont\color{color_29791}d}
\put(271.2771,-145.991){\fontsize{11.991}{1}\usefont{T1}{cmr}{m}{n}\selectfont\color{color_29791}a}
\put(277.9561,-145.991){\fontsize{11.991}{1}\usefont{T1}{cmr}{m}{n}\selectfont\color{color_29791}t}
\put(281.3376,-145.991){\fontsize{11.991}{1}\usefont{T1}{cmr}{m}{n}\selectfont\color{color_29791}a}
\put(84.213,-174.11){\fontsize{11.991}{1}\usefont{T1}{cmr}{m}{n}\selectfont\color{color_29791}•}
\put(98.187,-174.791){\fontsize{11.991}{1}\usefont{T1}{cmr}{m}{n}\selectfont\color{color_29791}S}
\put(106.2689,-174.791){\fontsize{11.991}{1}\usefont{T1}{cmr}{m}{n}\selectfont\color{color_29791}p}
\put(112.9479,-174.791){\fontsize{11.991}{1}\usefont{T1}{cmr}{m}{n}\selectfont\color{color_29791}a}
\put(119.6269,-174.791){\fontsize{11.991}{1}\usefont{T1}{cmr}{m}{n}\selectfont\color{color_29791}ce}
\put(132.3014,-174.791){\fontsize{11.991}{1}\usefont{T1}{cmr}{m}{n}\selectfont\color{color_29791} }
\put(135.6829,-174.791){\fontsize{11.991}{1}\usefont{T1}{cmr}{m}{n}\selectfont\color{color_29791}r}
\put(139.5799,-174.791){\fontsize{11.991}{1}\usefont{T1}{cmr}{m}{n}\selectfont\color{color_29791}e}
\put(146.3669,-174.791){\fontsize{11.991}{1}\usefont{T1}{cmr}{m}{n}\selectfont\color{color_29791}s}
\put(152.2664,-174.791){\fontsize{11.991}{1}\usefont{T1}{cmr}{m}{n}\selectfont\color{color_29791}t}
\put(155.6479,-174.791){\fontsize{11.991}{1}\usefont{T1}{cmr}{m}{n}\selectfont\color{color_29791}ri}
\put(162.2429,-174.791){\fontsize{11.991}{1}\usefont{T1}{cmr}{m}{n}\selectfont\color{color_29791}ct}
\put(171.6199,-174.791){\fontsize{11.991}{1}\usefont{T1}{cmr}{m}{n}\selectfont\color{color_29791}e}
\put(178.2989,-174.791){\fontsize{11.991}{1}\usefont{T1}{cmr}{m}{n}\selectfont\color{color_29791}d}
\put(184.9779,-174.791){\fontsize{11.991}{1}\usefont{T1}{cmr}{m}{n}\selectfont\color{color_29791} }
\put(188.3593,-174.791){\fontsize{11.991}{1}\usefont{T1}{cmr}{m}{n}\selectfont\color{color_29791}vi}
\put(196.9569,-174.791){\fontsize{11.991}{1}\usefont{T1}{cmr}{m}{n}\selectfont\color{color_29791}a}
\put(203.6359,-174.791){\fontsize{11.991}{1}\usefont{T1}{cmr}{m}{n}\selectfont\color{color_29791} }
\put(207.0173,-174.791){\fontsize{11.991}{1}\usefont{T1}{cmr}{m}{n}\selectfont\color{color_29791}q}
\put(213.6963,-174.791){\fontsize{11.991}{1}\usefont{T1}{cmr}{m}{n}\selectfont\color{color_29791}u}
\put(220.3753,-174.791){\fontsize{11.991}{1}\usefont{T1}{cmr}{m}{n}\selectfont\color{color_29791}o}
\put(227.0543,-174.791){\fontsize{11.991}{1}\usefont{T1}{cmr}{m}{n}\selectfont\color{color_29791}t}
\put(230.4358,-174.791){\fontsize{11.991}{1}\usefont{T1}{cmr}{m}{n}\selectfont\color{color_29791}a}
\put(237.1147,-174.791){\fontsize{11.991}{1}\usefont{T1}{cmr}{m}{n}\selectfont\color{color_29791}s}
\put(48.213,-208.409){\fontsize{15.987}{1}\usefont{T1}{cmr}{m}{n}\selectfont\color{color_29791}S}
\put(58.89232,-208.409){\fontsize{15.987}{1}\usefont{T1}{cmr}{m}{n}\selectfont\color{color_29791}A}
\put(69.57163,-208.409){\fontsize{15.987}{1}\usefont{T1}{cmr}{m}{n}\selectfont\color{color_29791}N}
\put(81.06628,-208.409){\fontsize{15.987}{1}\usefont{T1}{cmr}{m}{n}\selectfont\color{color_29791} }
\put(85.46271,-208.409){\fontsize{15.987}{1}\usefont{T1}{cmr}{m}{n}\selectfont\color{color_29791}V}
\put(95.24675,-208.409){\fontsize{15.987}{1}\usefont{T1}{cmr}{m}{n}\selectfont\color{color_29791}o}
\put(104.0396,-208.409){\fontsize{15.987}{1}\usefont{T1}{cmr}{m}{n}\selectfont\color{color_29791}l}
\put(107.6207,-208.409){\fontsize{15.987}{1}\usefont{T1}{cmr}{m}{n}\selectfont\color{color_29791}um}
\put(129.7946,-208.409){\fontsize{15.987}{1}\usefont{T1}{cmr}{m}{n}\selectfont\color{color_29791}es}
\put(84.213,-237.209){\fontsize{11.991}{1}\usefont{T1}{cmr}{m}{n}\selectfont\color{color_29791}•}
\put(98.187,-237.89){\fontsize{11.991}{1}\usefont{T1}{cmr}{m}{n}\selectfont\color{color_29791}S}
\put(106.2689,-237.89){\fontsize{11.991}{1}\usefont{T1}{cmr}{m}{n}\selectfont\color{color_29791}e}
\put(112.9479,-237.89){\fontsize{11.991}{1}\usefont{T1}{cmr}{m}{n}\selectfont\color{color_29791}t}
\put(116.2454,-237.89){\fontsize{11.991}{1}\usefont{T1}{cmr}{m}{n}\selectfont\color{color_29791}u}
\put(122.9244,-237.89){\fontsize{11.991}{1}\usefont{T1}{cmr}{m}{n}\selectfont\color{color_29791}p}
\put(129.7113,-237.89){\fontsize{11.991}{1}\usefont{T1}{cmr}{m}{n}\selectfont\color{color_29791} }
\put(133.0089,-237.89){\fontsize{11.991}{1}\usefont{T1}{cmr}{m}{n}\selectfont\color{color_29791}f}
\put(136.3903,-237.89){\fontsize{11.991}{1}\usefont{T1}{cmr}{m}{n}\selectfont\color{color_29791}o}
\put(143.0693,-237.89){\fontsize{11.991}{1}\usefont{T1}{cmr}{m}{n}\selectfont\color{color_29791}r }
\put(150.3598,-237.89){\fontsize{11.991}{1}\usefont{T1}{cmr}{m}{n}\selectfont\color{color_29791}p}
\put(157.1467,-237.89){\fontsize{11.991}{1}\usefont{T1}{cmr}{m}{n}\selectfont\color{color_29791}r}
\put(161.0438,-237.89){\fontsize{11.991}{1}\usefont{T1}{cmr}{m}{n}\selectfont\color{color_29791}o}
\put(167.7228,-237.89){\fontsize{11.991}{1}\usefont{T1}{cmr}{m}{n}\selectfont\color{color_29791}j}
\put(170.4088,-237.89){\fontsize{11.991}{1}\usefont{T1}{cmr}{m}{n}\selectfont\color{color_29791}e}
\put(177.0878,-237.89){\fontsize{11.991}{1}\usefont{T1}{cmr}{m}{n}\selectfont\color{color_29791}ct}
\put(186.4648,-237.89){\fontsize{11.991}{1}\usefont{T1}{cmr}{m}{n}\selectfont\color{color_29791}s }
\put(195.7578,-237.89){\fontsize{11.991}{1}\usefont{T1}{cmr}{m}{n}\selectfont\color{color_29791}a}
\put(202.4368,-237.89){\fontsize{11.991}{1}\usefont{T1}{cmr}{m}{n}\selectfont\color{color_29791}n}
\put(209.1158,-237.89){\fontsize{11.991}{1}\usefont{T1}{cmr}{m}{n}\selectfont\color{color_29791}d}
\put(215.9027,-237.89){\fontsize{11.991}{1}\usefont{T1}{cmr}{m}{n}\selectfont\color{color_29791} }
\put(219.2002,-237.89){\fontsize{11.991}{1}\usefont{T1}{cmr}{m}{n}\selectfont\color{color_29791}d}
\put(225.9871,-237.89){\fontsize{11.991}{1}\usefont{T1}{cmr}{m}{n}\selectfont\color{color_29791}e}
\put(232.6661,-237.89){\fontsize{11.991}{1}\usefont{T1}{cmr}{m}{n}\selectfont\color{color_29791}p}
\put(239.3451,-237.89){\fontsize{11.991}{1}\usefont{T1}{cmr}{m}{n}\selectfont\color{color_29791}a}
\put(246.0241,-237.89){\fontsize{11.991}{1}\usefont{T1}{cmr}{m}{n}\selectfont\color{color_29791}rt}
\put(253.3146,-237.89){\fontsize{11.991}{1}\usefont{T1}{cmr}{m}{n}\selectfont\color{color_29791}me}
\put(269.9821,-237.89){\fontsize{11.991}{1}\usefont{T1}{cmr}{m}{n}\selectfont\color{color_29791}n}
\put(276.6611,-237.89){\fontsize{11.991}{1}\usefont{T1}{cmr}{m}{n}\selectfont\color{color_29791}t}
\put(280.0425,-237.89){\fontsize{11.991}{1}\usefont{T1}{cmr}{m}{n}\selectfont\color{color_29791}s,}
\put(289.4195,-237.89){\fontsize{11.991}{1}\usefont{T1}{cmr}{m}{n}\selectfont\color{color_29791} }
\put(292.717,-237.89){\fontsize{11.991}{1}\usefont{T1}{cmr}{m}{n}\selectfont\color{color_29791}e}
\put(299.396,-237.89){\fontsize{11.991}{1}\usefont{T1}{cmr}{m}{n}\selectfont\color{color_29791}g}
\put(306.075,-237.89){\fontsize{11.991}{1}\usefont{T1}{cmr}{m}{n}\selectfont\color{color_29791}:}
\put(309.4565,-237.89){\fontsize{11.991}{1}\usefont{T1}{cmr}{m}{n}\selectfont\color{color_29791} }
\put(312.8379,-237.89){\fontsize{11.991}{1}\usefont{T1}{cmr}{m}{n}\selectfont\color{color_29791}:}
\put(316.1354,-237.89){\fontsize{11.991}{1}\usefont{T1}{cmr}{m}{n}\selectfont\color{color_29791} }
\put(319.5169,-237.89){\fontsize{11.991}{1}\usefont{T1}{cmr}{m}{n}\selectfont\color{color_29791}/}
\put(322.8984,-237.89){\fontsize{11.991}{1}\usefont{T1}{cmr}{m}{n}\selectfont\color{color_29791}d}
\put(329.5774,-237.89){\fontsize{11.991}{1}\usefont{T1}{cmr}{m}{n}\selectfont\color{color_29791}a}
\put(336.2563,-237.89){\fontsize{11.991}{1}\usefont{T1}{cmr}{m}{n}\selectfont\color{color_29791}t}
\put(339.6378,-237.89){\fontsize{11.991}{1}\usefont{T1}{cmr}{m}{n}\selectfont\color{color_29791}a}
\put(346.3168,-237.89){\fontsize{11.991}{1}\usefont{T1}{cmr}{m}{n}\selectfont\color{color_29791}/}
\put(349.6143,-237.89){\fontsize{11.991}{1}\usefont{T1}{cmr}{m}{n}\selectfont\color{color_29791}cru}
\put(366.2818,-237.89){\fontsize{11.991}{1}\usefont{T1}{cmr}{m}{n}\selectfont\color{color_29791}i}
\put(368.9678,-237.89){\fontsize{11.991}{1}\usefont{T1}{cmr}{m}{n}\selectfont\color{color_29791}s}
\put(374.8673,-237.89){\fontsize{11.991}{1}\usefont{T1}{cmr}{m}{n}\selectfont\color{color_29791}e}
\put(381.6542,-237.89){\fontsize{11.991}{1}\usefont{T1}{cmr}{m}{n}\selectfont\color{color_29791},}
\put(384.9518,-237.89){\fontsize{11.991}{1}\usefont{T1}{cmr}{m}{n}\selectfont\color{color_29791} }
\put(388.3332,-237.89){\fontsize{11.991}{1}\usefont{T1}{cmr}{m}{n}\selectfont\color{color_29791}/}
\put(391.7147,-237.89){\fontsize{11.991}{1}\usefont{T1}{cmr}{m}{n}\selectfont\color{color_29791}d}
\put(398.3937,-237.89){\fontsize{11.991}{1}\usefont{T1}{cmr}{m}{n}\selectfont\color{color_29791}a}
\put(405.0727,-237.89){\fontsize{11.991}{1}\usefont{T1}{cmr}{m}{n}\selectfont\color{color_29791}t}
\put(408.4541,-237.89){\fontsize{11.991}{1}\usefont{T1}{cmr}{m}{n}\selectfont\color{color_29791}a}
\put(415.1331,-237.89){\fontsize{11.991}{1}\usefont{T1}{cmr}{m}{n}\selectfont\color{color_29791}/}
\put(418.4306,-237.89){\fontsize{11.991}{1}\usefont{T1}{cmr}{m}{n}\selectfont\color{color_29791}vl}
\put(427.1121,-237.89){\fontsize{11.991}{1}\usefont{T1}{cmr}{m}{n}\selectfont\color{color_29791}f}
\put(84.213,-266.009){\fontsize{11.991}{1}\usefont{T1}{cmr}{m}{n}\selectfont\color{color_29791}•}
\put(98.187,-266.69){\fontsize{11.991}{1}\usefont{T1}{cmr}{m}{n}\selectfont\color{color_29791}A}
\put(106.2689,-266.69){\fontsize{11.991}{1}\usefont{T1}{cmr}{m}{n}\selectfont\color{color_29791}c}
\put(112.1685,-266.69){\fontsize{11.991}{1}\usefont{T1}{cmr}{m}{n}\selectfont\color{color_29791}c}
\put(118.2479,-266.69){\fontsize{11.991}{1}\usefont{T1}{cmr}{m}{n}\selectfont\color{color_29791}e}
\put(124.9269,-266.69){\fontsize{11.991}{1}\usefont{T1}{cmr}{m}{n}\selectfont\color{color_29791}ssi}
\put(139.52,-266.69){\fontsize{11.991}{1}\usefont{T1}{cmr}{m}{n}\selectfont\color{color_29791}b}
\put(146.199,-266.69){\fontsize{11.991}{1}\usefont{T1}{cmr}{m}{n}\selectfont\color{color_29791}l}
\put(148.8849,-266.69){\fontsize{11.991}{1}\usefont{T1}{cmr}{m}{n}\selectfont\color{color_29791}e}
\put(155.5639,-266.69){\fontsize{11.991}{1}\usefont{T1}{cmr}{m}{n}\selectfont\color{color_29791} }
\put(158.8615,-266.69){\fontsize{11.991}{1}\usefont{T1}{cmr}{m}{n}\selectfont\color{color_29791}f}
\put(162.2429,-266.69){\fontsize{11.991}{1}\usefont{T1}{cmr}{m}{n}\selectfont\color{color_29791}ro}
\put(172.9149,-266.69){\fontsize{11.991}{1}\usefont{T1}{cmr}{m}{n}\selectfont\color{color_29791}m }
\put(186.2009,-266.69){\fontsize{11.991}{1}\usefont{T1}{cmr}{m}{n}\selectfont\color{color_29791}w}
\put(194.7985,-266.69){\fontsize{11.991}{1}\usefont{T1}{cmr}{m}{n}\selectfont\color{color_29791}o}
\put(201.5854,-266.69){\fontsize{11.991}{1}\usefont{T1}{cmr}{m}{n}\selectfont\color{color_29791}r}
\put(205.4825,-266.69){\fontsize{11.991}{1}\usefont{T1}{cmr}{m}{n}\selectfont\color{color_29791}kst}
\put(220.8549,-266.69){\fontsize{11.991}{1}\usefont{T1}{cmr}{m}{n}\selectfont\color{color_29791}a}
\put(227.5339,-266.69){\fontsize{11.991}{1}\usefont{T1}{cmr}{m}{n}\selectfont\color{color_29791}t}
\put(230.9154,-266.69){\fontsize{11.991}{1}\usefont{T1}{cmr}{m}{n}\selectfont\color{color_29791}i}
\put(233.5175,-266.69){\fontsize{11.991}{1}\usefont{T1}{cmr}{m}{n}\selectfont\color{color_29791}o}
\put(240.1964,-266.69){\fontsize{11.991}{1}\usefont{T1}{cmr}{m}{n}\selectfont\color{color_29791}n}
\put(246.9833,-266.69){\fontsize{11.991}{1}\usefont{T1}{cmr}{m}{n}\selectfont\color{color_29791}s}
\put(252.8829,-266.69){\fontsize{11.991}{1}\usefont{T1}{cmr}{m}{n}\selectfont\color{color_29791},}
\put(256.2644,-266.69){\fontsize{11.991}{1}\usefont{T1}{cmr}{m}{n}\selectfont\color{color_29791} }
\put(259.6458,-266.69){\fontsize{11.991}{1}\usefont{T1}{cmr}{m}{n}\selectfont\color{color_29791}b}
\put(266.3248,-266.69){\fontsize{11.991}{1}\usefont{T1}{cmr}{m}{n}\selectfont\color{color_29791}sl}
\put(274.9224,-266.69){\fontsize{11.991}{1}\usefont{T1}{cmr}{m}{n}\selectfont\color{color_29791}ce}
\put(287.7048,-266.69){\fontsize{11.991}{1}\usefont{T1}{cmr}{m}{n}\selectfont\color{color_29791}n}
\put(294.3838,-266.69){\fontsize{11.991}{1}\usefont{T1}{cmr}{m}{n}\selectfont\color{color_29791}b}
\put(301.0627,-266.69){\fontsize{11.991}{1}\usefont{T1}{cmr}{m}{n}\selectfont\color{color_29791},}
\put(304.4442,-266.69){\fontsize{11.991}{1}\usefont{T1}{cmr}{m}{n}\selectfont\color{color_29791} }
\put(307.7417,-266.69){\fontsize{11.991}{1}\usefont{T1}{cmr}{m}{n}\selectfont\color{color_29791}b}
\put(314.4207,-266.69){\fontsize{11.991}{1}\usefont{T1}{cmr}{m}{n}\selectfont\color{color_29791}sl}
\put(323.1022,-266.69){\fontsize{11.991}{1}\usefont{T1}{cmr}{m}{n}\selectfont\color{color_29791}ce}
\put(335.7767,-266.69){\fontsize{11.991}{1}\usefont{T1}{cmr}{m}{n}\selectfont\color{color_29791}n}
\put(342.4557,-266.69){\fontsize{11.991}{1}\usefont{T1}{cmr}{m}{n}\selectfont\color{color_29791}c}
\put(84.213,-294.809){\fontsize{11.991}{1}\usefont{T1}{cmr}{m}{n}\selectfont\color{color_29791}•}
\put(98.187,-295.49){\fontsize{11.991}{1}\usefont{T1}{cmr}{m}{n}\selectfont\color{color_29791}V}
\put(105.5854,-295.49){\fontsize{11.991}{1}\usefont{T1}{cmr}{m}{n}\selectfont\color{color_29791}o}
\put(112.1805,-295.49){\fontsize{11.991}{1}\usefont{T1}{cmr}{m}{n}\selectfont\color{color_29791}l}
\put(114.8665,-295.49){\fontsize{11.991}{1}\usefont{T1}{cmr}{m}{n}\selectfont\color{color_29791}u}
\put(121.5455,-295.49){\fontsize{11.991}{1}\usefont{T1}{cmr}{m}{n}\selectfont\color{color_29791}me}
\put(138.213,-295.49){\fontsize{11.991}{1}\usefont{T1}{cmr}{m}{n}\selectfont\color{color_29791} }
\put(141.5105,-295.49){\fontsize{11.991}{1}\usefont{T1}{cmr}{m}{n}\selectfont\color{color_29791}s}
\put(147.5899,-295.49){\fontsize{11.991}{1}\usefont{T1}{cmr}{m}{n}\selectfont\color{color_29791}h}
\put(154.2689,-295.49){\fontsize{11.991}{1}\usefont{T1}{cmr}{m}{n}\selectfont\color{color_29791}o}
\put(160.9479,-295.49){\fontsize{11.991}{1}\usefont{T1}{cmr}{m}{n}\selectfont\color{color_29791}u}
\put(167.6269,-295.49){\fontsize{11.991}{1}\usefont{T1}{cmr}{m}{n}\selectfont\color{color_29791}l}
\put(170.2289,-295.49){\fontsize{11.991}{1}\usefont{T1}{cmr}{m}{n}\selectfont\color{color_29791}d}
\put(177.0158,-295.49){\fontsize{11.991}{1}\usefont{T1}{cmr}{m}{n}\selectfont\color{color_29791} }
\put(180.3134,-295.49){\fontsize{11.991}{1}\usefont{T1}{cmr}{m}{n}\selectfont\color{color_29791}b}
\put(187.1003,-295.49){\fontsize{11.991}{1}\usefont{T1}{cmr}{m}{n}\selectfont\color{color_29791}e}
\put(193.7793,-295.49){\fontsize{11.991}{1}\usefont{T1}{cmr}{m}{n}\selectfont\color{color_29791} }
\put(197.0768,-295.49){\fontsize{11.991}{1}\usefont{T1}{cmr}{m}{n}\selectfont\color{color_29791}ma}
\put(213.7443,-295.49){\fontsize{11.991}{1}\usefont{T1}{cmr}{m}{n}\selectfont\color{color_29791}n}
\put(220.4233,-295.49){\fontsize{11.991}{1}\usefont{T1}{cmr}{m}{n}\selectfont\color{color_29791}a}
\put(227.1023,-295.49){\fontsize{11.991}{1}\usefont{T1}{cmr}{m}{n}\selectfont\color{color_29791}g}
\put(233.8892,-295.49){\fontsize{11.991}{1}\usefont{T1}{cmr}{m}{n}\selectfont\color{color_29791}e}
\put(240.5682,-295.49){\fontsize{11.991}{1}\usefont{T1}{cmr}{m}{n}\selectfont\color{color_29791}d}
\put(247.2472,-295.49){\fontsize{11.991}{1}\usefont{T1}{cmr}{m}{n}\selectfont\color{color_29791} }
\put(250.6286,-295.49){\fontsize{11.991}{1}\usefont{T1}{cmr}{m}{n}\selectfont\color{color_29791}a}
\put(257.3076,-295.49){\fontsize{11.991}{1}\usefont{T1}{cmr}{m}{n}\selectfont\color{color_29791}n}
\put(263.9866,-295.49){\fontsize{11.991}{1}\usefont{T1}{cmr}{m}{n}\selectfont\color{color_29791}d}
\put(270.6656,-295.49){\fontsize{11.991}{1}\usefont{T1}{cmr}{m}{n}\selectfont\color{color_29791} }
\put(274.0471,-295.49){\fontsize{11.991}{1}\usefont{T1}{cmr}{m}{n}\selectfont\color{color_29791}cu}
\put(286.7216,-295.49){\fontsize{11.991}{1}\usefont{T1}{cmr}{m}{n}\selectfont\color{color_29791}ra}
\put(297.3936,-295.49){\fontsize{11.991}{1}\usefont{T1}{cmr}{m}{n}\selectfont\color{color_29791}t}
\put(300.775,-295.49){\fontsize{11.991}{1}\usefont{T1}{cmr}{m}{n}\selectfont\color{color_29791}e}
\put(307.454,-295.49){\fontsize{11.991}{1}\usefont{T1}{cmr}{m}{n}\selectfont\color{color_29791}d}
\put(314.133,-295.49){\fontsize{11.991}{1}\usefont{T1}{cmr}{m}{n}\selectfont\color{color_29791} }
\put(317.4305,-295.49){\fontsize{11.991}{1}\usefont{T1}{cmr}{m}{n}\selectfont\color{color_29791}b}
\put(324.2174,-295.49){\fontsize{11.991}{1}\usefont{T1}{cmr}{m}{n}\selectfont\color{color_29791}y }
\put(333.5104,-295.49){\fontsize{11.991}{1}\usefont{T1}{cmr}{m}{n}\selectfont\color{color_29791}a}
\put(340.2973,-295.49){\fontsize{11.991}{1}\usefont{T1}{cmr}{m}{n}\selectfont\color{color_29791} }
\put(343.5948,-295.49){\fontsize{11.991}{1}\usefont{T1}{cmr}{m}{n}\selectfont\color{color_29791}d}
\put(350.2738,-295.49){\fontsize{11.991}{1}\usefont{T1}{cmr}{m}{n}\selectfont\color{color_29791}a}
\put(356.9528,-295.49){\fontsize{11.991}{1}\usefont{T1}{cmr}{m}{n}\selectfont\color{color_29791}t}
\put(360.3343,-295.49){\fontsize{11.991}{1}\usefont{T1}{cmr}{m}{n}\selectfont\color{color_29791}a}
\put(367.0133,-295.49){\fontsize{11.991}{1}\usefont{T1}{cmr}{m}{n}\selectfont\color{color_29791} }
\put(370.3947,-295.49){\fontsize{11.991}{1}\usefont{T1}{cmr}{m}{n}\selectfont\color{color_29791}ma}
\put(387.0622,-295.49){\fontsize{11.991}{1}\usefont{T1}{cmr}{m}{n}\selectfont\color{color_29791}n}
\put(393.7412,-295.49){\fontsize{11.991}{1}\usefont{T1}{cmr}{m}{n}\selectfont\color{color_29791}a}
\put(400.4202,-295.49){\fontsize{11.991}{1}\usefont{T1}{cmr}{m}{n}\selectfont\color{color_29791}g}
\put(407.0992,-295.49){\fontsize{11.991}{1}\usefont{T1}{cmr}{m}{n}\selectfont\color{color_29791}e}
\put(413.7782,-295.49){\fontsize{11.991}{1}\usefont{T1}{cmr}{m}{n}\selectfont\color{color_29791}r}
\put(84.213,-323.609){\fontsize{11.991}{1}\usefont{T1}{cmr}{m}{n}\selectfont\color{color_29791}•}
\put(98.187,-324.205){\fontsize{11.991}{1}\usefont{T1}{cmr}{m}{n}\selectfont\color{color_29791}S}
\put(106.2689,-324.205){\fontsize{11.991}{1}\usefont{T1}{cmr}{m}{n}\selectfont\color{color_29791}p}
\put(112.9479,-324.205){\fontsize{11.991}{1}\usefont{T1}{cmr}{m}{n}\selectfont\color{color_29791}a}
\put(119.6269,-324.205){\fontsize{11.991}{1}\usefont{T1}{cmr}{m}{n}\selectfont\color{color_29791}ce}
\put(132.3014,-324.205){\fontsize{11.991}{1}\usefont{T1}{cmr}{m}{n}\selectfont\color{color_29791} }
\put(135.6829,-324.205){\fontsize{11.991}{1}\usefont{T1}{cmr}{m}{n}\selectfont\color{color_29791}i}
\put(138.2849,-324.205){\fontsize{11.991}{1}\usefont{T1}{cmr}{m}{n}\selectfont\color{color_29791}s }
\put(147.6619,-324.205){\fontsize{11.991}{1}\usefont{T1}{cmr}{m}{n}\selectfont\color{color_29791}n}
\put(154.3409,-324.205){\fontsize{11.991}{1}\usefont{T1}{cmr}{m}{n}\selectfont\color{color_29791}o}
\put(161.0199,-324.205){\fontsize{11.991}{1}\usefont{T1}{cmr}{m}{n}\selectfont\color{color_29791}t}
\put(164.4013,-324.205){\fontsize{11.991}{1}\usefont{T1}{cmr}{m}{n}\selectfont\color{color_29791} }
\put(167.6988,-324.205){\fontsize{11.991}{1}\usefont{T1}{cmr}{m}{n}\selectfont\color{color_29791}co}
\put(180.3733,-324.205){\fontsize{11.991}{1}\usefont{T1}{cmr}{m}{n}\selectfont\color{color_29791}n}
\put(187.1602,-324.205){\fontsize{11.991}{1}\usefont{T1}{cmr}{m}{n}\selectfont\color{color_29791}t}
\put(190.4578,-324.205){\fontsize{11.991}{1}\usefont{T1}{cmr}{m}{n}\selectfont\color{color_29791}ro}
\put(201.1298,-324.205){\fontsize{11.991}{1}\usefont{T1}{cmr}{m}{n}\selectfont\color{color_29791}l}
\put(203.8158,-324.205){\fontsize{11.991}{1}\usefont{T1}{cmr}{m}{n}\selectfont\color{color_29791}l}
\put(206.4178,-324.205){\fontsize{11.991}{1}\usefont{T1}{cmr}{m}{n}\selectfont\color{color_29791}e}
\put(213.0968,-324.205){\fontsize{11.991}{1}\usefont{T1}{cmr}{m}{n}\selectfont\color{color_29791}d}
\put(219.7758,-324.205){\fontsize{11.991}{1}\usefont{T1}{cmr}{m}{n}\selectfont\color{color_29791} }
\put(223.1572,-324.205){\fontsize{11.991}{1}\usefont{T1}{cmr}{m}{n}\selectfont\color{color_29791}b}
\put(229.8362,-324.205){\fontsize{11.991}{1}\usefont{T1}{cmr}{m}{n}\selectfont\color{color_29791}y }
\put(239.2132,-324.205){\fontsize{11.991}{1}\usefont{T1}{cmr}{m}{n}\selectfont\color{color_29791}q}
\put(245.8922,-324.205){\fontsize{11.991}{1}\usefont{T1}{cmr}{m}{n}\selectfont\color{color_29791}u}
\put(252.5712,-324.205){\fontsize{11.991}{1}\usefont{T1}{cmr}{m}{n}\selectfont\color{color_29791}o}
\put(259.2502,-324.205){\fontsize{11.991}{1}\usefont{T1}{cmr}{m}{n}\selectfont\color{color_29791}t}
\put(262.6316,-324.205){\fontsize{11.991}{1}\usefont{T1}{cmr}{m}{n}\selectfont\color{color_29791}a}
\put(269.3106,-324.205){\fontsize{11.991}{1}\usefont{T1}{cmr}{m}{n}\selectfont\color{color_29791}'}
\put(271.6009,-324.205){\fontsize{11.991}{1}\usefont{T1}{cmr}{m}{n}\selectfont\color{color_29791}s}
\put(84.213,-352.296){\fontsize{11.991}{1}\usefont{T1}{cmr}{m}{n}\selectfont\color{color_29791}•}
\put(98.187,-353.005){\fontsize{11.991}{1}\usefont{T1}{cmr}{m}{n}\selectfont\color{color_29791}A}
\put(106.2689,-353.005){\fontsize{11.991}{1}\usefont{T1}{cmr}{m}{n}\selectfont\color{color_29791}d}
\put(112.9479,-353.005){\fontsize{11.991}{1}\usefont{T1}{cmr}{m}{n}\selectfont\color{color_29791}d}
\put(119.6269,-353.005){\fontsize{11.991}{1}\usefont{T1}{cmr}{m}{n}\selectfont\color{color_29791}i}
\put(122.229,-353.005){\fontsize{11.991}{1}\usefont{T1}{cmr}{m}{n}\selectfont\color{color_29791}n}
\put(129.0159,-353.005){\fontsize{11.991}{1}\usefont{T1}{cmr}{m}{n}\selectfont\color{color_29791}g}
\put(135.6949,-353.005){\fontsize{11.991}{1}\usefont{T1}{cmr}{m}{n}\selectfont\color{color_29791} }
\put(138.9924,-353.005){\fontsize{11.991}{1}\usefont{T1}{cmr}{m}{n}\selectfont\color{color_29791}a}
\put(145.7793,-353.005){\fontsize{11.991}{1}\usefont{T1}{cmr}{m}{n}\selectfont\color{color_29791}d}
\put(152.4583,-353.005){\fontsize{11.991}{1}\usefont{T1}{cmr}{m}{n}\selectfont\color{color_29791}d}
\put(159.1373,-353.005){\fontsize{11.991}{1}\usefont{T1}{cmr}{m}{n}\selectfont\color{color_29791}i}
\put(161.7393,-353.005){\fontsize{11.991}{1}\usefont{T1}{cmr}{m}{n}\selectfont\color{color_29791}t}
\put(165.1208,-353.005){\fontsize{11.991}{1}\usefont{T1}{cmr}{m}{n}\selectfont\color{color_29791}i}
\put(167.7228,-353.005){\fontsize{11.991}{1}\usefont{T1}{cmr}{m}{n}\selectfont\color{color_29791}o}
\put(174.5098,-353.005){\fontsize{11.991}{1}\usefont{T1}{cmr}{m}{n}\selectfont\color{color_29791}n}
\put(181.1887,-353.005){\fontsize{11.991}{1}\usefont{T1}{cmr}{m}{n}\selectfont\color{color_29791}a}
\put(187.8677,-353.005){\fontsize{11.991}{1}\usefont{T1}{cmr}{m}{n}\selectfont\color{color_29791}l}
\put(190.4698,-353.005){\fontsize{11.991}{1}\usefont{T1}{cmr}{m}{n}\selectfont\color{color_29791} }
\put(193.8512,-353.005){\fontsize{11.991}{1}\usefont{T1}{cmr}{m}{n}\selectfont\color{color_29791}sp}
\put(206.5257,-353.005){\fontsize{11.991}{1}\usefont{T1}{cmr}{m}{n}\selectfont\color{color_29791}a}
\put(213.2047,-353.005){\fontsize{11.991}{1}\usefont{T1}{cmr}{m}{n}\selectfont\color{color_29791}ce}
\put(225.9871,-353.005){\fontsize{11.991}{1}\usefont{T1}{cmr}{m}{n}\selectfont\color{color_29791} }
\put(229.2846,-353.005){\fontsize{11.991}{1}\usefont{T1}{cmr}{m}{n}\selectfont\color{color_29791}d}
\put(235.9636,-353.005){\fontsize{11.991}{1}\usefont{T1}{cmr}{m}{n}\selectfont\color{color_29791}e}
\put(242.7505,-353.005){\fontsize{11.991}{1}\usefont{T1}{cmr}{m}{n}\selectfont\color{color_29791}p}
\put(249.4295,-353.005){\fontsize{11.991}{1}\usefont{T1}{cmr}{m}{n}\selectfont\color{color_29791}e}
\put(256.1085,-353.005){\fontsize{11.991}{1}\usefont{T1}{cmr}{m}{n}\selectfont\color{color_29791}n}
\put(262.7875,-353.005){\fontsize{11.991}{1}\usefont{T1}{cmr}{m}{n}\selectfont\color{color_29791}d}
\put(269.4665,-353.005){\fontsize{11.991}{1}\usefont{T1}{cmr}{m}{n}\selectfont\color{color_29791}s }
\put(278.8434,-353.005){\fontsize{11.991}{1}\usefont{T1}{cmr}{m}{n}\selectfont\color{color_29791}a}
\put(285.5224,-353.005){\fontsize{11.991}{1}\usefont{T1}{cmr}{m}{n}\selectfont\color{color_29791}va}
\put(298.1969,-353.005){\fontsize{11.991}{1}\usefont{T1}{cmr}{m}{n}\selectfont\color{color_29791}i}
\put(300.8829,-353.005){\fontsize{11.991}{1}\usefont{T1}{cmr}{m}{n}\selectfont\color{color_29791}l}
\put(303.4849,-353.005){\fontsize{11.991}{1}\usefont{T1}{cmr}{m}{n}\selectfont\color{color_29791}a}
\put(310.1639,-353.005){\fontsize{11.991}{1}\usefont{T1}{cmr}{m}{n}\selectfont\color{color_29791}b}
\put(316.8429,-353.005){\fontsize{11.991}{1}\usefont{T1}{cmr}{m}{n}\selectfont\color{color_29791}i}
\put(319.5289,-353.005){\fontsize{11.991}{1}\usefont{T1}{cmr}{m}{n}\selectfont\color{color_29791}l}
\put(322.1309,-353.005){\fontsize{11.991}{1}\usefont{T1}{cmr}{m}{n}\selectfont\color{color_29791}i}
\put(324.8169,-353.005){\fontsize{11.991}{1}\usefont{T1}{cmr}{m}{n}\selectfont\color{color_29791}t}
\put(328.1144,-353.005){\fontsize{11.991}{1}\usefont{T1}{cmr}{m}{n}\selectfont\color{color_29791}y }
\put(337.4914,-353.005){\fontsize{11.991}{1}\usefont{T1}{cmr}{m}{n}\selectfont\color{color_29791}o}
\put(344.1703,-353.005){\fontsize{11.991}{1}\usefont{T1}{cmr}{m}{n}\selectfont\color{color_29791}f}
\put(347.5518,-353.005){\fontsize{11.991}{1}\usefont{T1}{cmr}{m}{n}\selectfont\color{color_29791} }
\put(350.8493,-353.005){\fontsize{11.991}{1}\usefont{T1}{cmr}{m}{n}\selectfont\color{color_29791}p}
\put(357.6362,-353.005){\fontsize{11.991}{1}\usefont{T1}{cmr}{m}{n}\selectfont\color{color_29791}h}
\put(364.3152,-353.005){\fontsize{11.991}{1}\usefont{T1}{cmr}{m}{n}\selectfont\color{color_29791}ysi}
\put(378.9083,-353.005){\fontsize{11.991}{1}\usefont{T1}{cmr}{m}{n}\selectfont\color{color_29791}ca}
\put(391.5828,-353.005){\fontsize{11.991}{1}\usefont{T1}{cmr}{m}{n}\selectfont\color{color_29791}l}
\put(394.2687,-353.005){\fontsize{11.991}{1}\usefont{T1}{cmr}{m}{n}\selectfont\color{color_29791} }
\put(397.5663,-353.005){\fontsize{11.991}{1}\usefont{T1}{cmr}{m}{n}\selectfont\color{color_29791}d}
\put(404.3531,-353.005){\fontsize{11.991}{1}\usefont{T1}{cmr}{m}{n}\selectfont\color{color_29791}i}
\put(406.9552,-353.005){\fontsize{11.991}{1}\usefont{T1}{cmr}{m}{n}\selectfont\color{color_29791}sk }
\put(422.3277,-353.005){\fontsize{11.991}{1}\usefont{T1}{cmr}{m}{n}\selectfont\color{color_29791}s}
\put(428.2272,-353.005){\fontsize{11.991}{1}\usefont{T1}{cmr}{m}{n}\selectfont\color{color_29791}p}
\put(435.0141,-353.005){\fontsize{11.991}{1}\usefont{T1}{cmr}{m}{n}\selectfont\color{color_29791}a}
\put(441.6931,-353.005){\fontsize{11.991}{1}\usefont{T1}{cmr}{m}{n}\selectfont\color{color_29791}ce}
\put(84.213,-381.096){\fontsize{11.991}{1}\usefont{T1}{cmr}{m}{n}\selectfont\color{color_29791}•}
\put(98.187,-381.805){\fontsize{11.991}{1}\usefont{T1}{cmr}{m}{n}\selectfont\color{color_29791}C}
\put(106.8685,-381.805){\fontsize{11.991}{1}\usefont{T1}{cmr}{m}{n}\selectfont\color{color_29791}o}
\put(113.5475,-381.805){\fontsize{11.991}{1}\usefont{T1}{cmr}{m}{n}\selectfont\color{color_29791}n}
\put(120.2265,-381.805){\fontsize{11.991}{1}\usefont{T1}{cmr}{m}{n}\selectfont\color{color_29791}t}
\put(123.524,-381.805){\fontsize{11.991}{1}\usefont{T1}{cmr}{m}{n}\selectfont\color{color_29791}a}
\put(130.3109,-381.805){\fontsize{11.991}{1}\usefont{T1}{cmr}{m}{n}\selectfont\color{color_29791}ct}
\put(139.6039,-381.805){\fontsize{11.991}{1}\usefont{T1}{cmr}{m}{n}\selectfont\color{color_29791} }
\put(142.9854,-381.805){\fontsize{11.991}{1}\usefont{T1}{cmr}{m}{n}\selectfont\color{color_29791}d}
\put(149.6644,-381.805){\fontsize{11.991}{1}\usefont{T1}{cmr}{m}{n}\selectfont\color{color_29791}a}
\put(156.3434,-381.805){\fontsize{11.991}{1}\usefont{T1}{cmr}{m}{n}\selectfont\color{color_29791}t}
\put(159.7248,-381.805){\fontsize{11.991}{1}\usefont{T1}{cmr}{m}{n}\selectfont\color{color_29791}a}
\put(166.4038,-381.805){\fontsize{11.991}{1}\usefont{T1}{cmr}{m}{n}\selectfont\color{color_29791} }
\put(169.7853,-381.805){\fontsize{11.991}{1}\usefont{T1}{cmr}{m}{n}\selectfont\color{color_29791}m}
\put(179.6779,-381.805){\fontsize{11.991}{1}\usefont{T1}{cmr}{m}{n}\selectfont\color{color_29791}a}
\put(186.4648,-381.805){\fontsize{11.991}{1}\usefont{T1}{cmr}{m}{n}\selectfont\color{color_29791}n}
\put(193.1438,-381.805){\fontsize{11.991}{1}\usefont{T1}{cmr}{m}{n}\selectfont\color{color_29791}a}
\put(199.8227,-381.805){\fontsize{11.991}{1}\usefont{T1}{cmr}{m}{n}\selectfont\color{color_29791}g}
\put(206.5017,-381.805){\fontsize{11.991}{1}\usefont{T1}{cmr}{m}{n}\selectfont\color{color_29791}e}
\put(213.1807,-381.805){\fontsize{11.991}{1}\usefont{T1}{cmr}{m}{n}\selectfont\color{color_29791}r }
\put(220.4712,-381.805){\fontsize{11.991}{1}\usefont{T1}{cmr}{m}{n}\selectfont\color{color_29791}f}
\put(223.8527,-381.805){\fontsize{11.991}{1}\usefont{T1}{cmr}{m}{n}\selectfont\color{color_29791}i}
\put(226.5387,-381.805){\fontsize{11.991}{1}\usefont{T1}{cmr}{m}{n}\selectfont\color{color_29791}r}
\put(230.4358,-381.805){\fontsize{11.991}{1}\usefont{T1}{cmr}{m}{n}\selectfont\color{color_29791}st}
\put(239.8127,-381.805){\fontsize{11.991}{1}\usefont{T1}{cmr}{m}{n}\selectfont\color{color_29791} }
\put(243.1103,-381.805){\fontsize{11.991}{1}\usefont{T1}{cmr}{m}{n}\selectfont\color{color_29791}i}
\put(245.7962,-381.805){\fontsize{11.991}{1}\usefont{T1}{cmr}{m}{n}\selectfont\color{color_29791}f}
\put(249.0938,-381.805){\fontsize{11.991}{1}\usefont{T1}{cmr}{m}{n}\selectfont\color{color_29791} }
\put(252.4752,-381.805){\fontsize{11.991}{1}\usefont{T1}{cmr}{m}{n}\selectfont\color{color_29791}yo}
\put(265.1497,-381.805){\fontsize{11.991}{1}\usefont{T1}{cmr}{m}{n}\selectfont\color{color_29791}u}
\put(271.8287,-381.805){\fontsize{11.991}{1}\usefont{T1}{cmr}{m}{n}\selectfont\color{color_29791} }
\put(275.2102,-381.805){\fontsize{11.991}{1}\usefont{T1}{cmr}{m}{n}\selectfont\color{color_29791}t}
\put(278.5916,-381.805){\fontsize{11.991}{1}\usefont{T1}{cmr}{m}{n}\selectfont\color{color_29791}h}
\put(285.2706,-381.805){\fontsize{11.991}{1}\usefont{T1}{cmr}{m}{n}\selectfont\color{color_29791}i}
\put(287.8727,-381.805){\fontsize{11.991}{1}\usefont{T1}{cmr}{m}{n}\selectfont\color{color_29791}n}
\put(294.5517,-381.805){\fontsize{11.991}{1}\usefont{T1}{cmr}{m}{n}\selectfont\color{color_29791}k }
\put(303.9286,-381.805){\fontsize{11.991}{1}\usefont{T1}{cmr}{m}{n}\selectfont\color{color_29791}yo}
\put(316.6031,-381.805){\fontsize{11.991}{1}\usefont{T1}{cmr}{m}{n}\selectfont\color{color_29791}u}
\put(323.2821,-381.805){\fontsize{11.991}{1}\usefont{T1}{cmr}{m}{n}\selectfont\color{color_29791} }
\put(326.6636,-381.805){\fontsize{11.991}{1}\usefont{T1}{cmr}{m}{n}\selectfont\color{color_29791}re}
\put(337.3356,-381.805){\fontsize{11.991}{1}\usefont{T1}{cmr}{m}{n}\selectfont\color{color_29791}q}
\put(344.0146,-381.805){\fontsize{11.991}{1}\usefont{T1}{cmr}{m}{n}\selectfont\color{color_29791}u}
\put(350.6935,-381.805){\fontsize{11.991}{1}\usefont{T1}{cmr}{m}{n}\selectfont\color{color_29791}i}
\put(353.2956,-381.805){\fontsize{11.991}{1}\usefont{T1}{cmr}{m}{n}\selectfont\color{color_29791}re}
\put(363.9676,-381.805){\fontsize{11.991}{1}\usefont{T1}{cmr}{m}{n}\selectfont\color{color_29791} }
\put(367.3491,-381.805){\fontsize{11.991}{1}\usefont{T1}{cmr}{m}{n}\selectfont\color{color_29791}a}
\put(374.028,-381.805){\fontsize{11.991}{1}\usefont{T1}{cmr}{m}{n}\selectfont\color{color_29791}d}
\put(380.707,-381.805){\fontsize{11.991}{1}\usefont{T1}{cmr}{m}{n}\selectfont\color{color_29791}d}
\put(387.4939,-381.805){\fontsize{11.991}{1}\usefont{T1}{cmr}{m}{n}\selectfont\color{color_29791}i}
\put(390.096,-381.805){\fontsize{11.991}{1}\usefont{T1}{cmr}{m}{n}\selectfont\color{color_29791}t}
\put(393.4774,-381.805){\fontsize{11.991}{1}\usefont{T1}{cmr}{m}{n}\selectfont\color{color_29791}i}
\put(396.0795,-381.805){\fontsize{11.991}{1}\usefont{T1}{cmr}{m}{n}\selectfont\color{color_29791}o}
\put(402.7585,-381.805){\fontsize{11.991}{1}\usefont{T1}{cmr}{m}{n}\selectfont\color{color_29791}n}
\put(409.4375,-381.805){\fontsize{11.991}{1}\usefont{T1}{cmr}{m}{n}\selectfont\color{color_29791}a}
\put(416.2244,-381.805){\fontsize{11.991}{1}\usefont{T1}{cmr}{m}{n}\selectfont\color{color_29791}l}
\put(418.8264,-381.805){\fontsize{11.991}{1}\usefont{T1}{cmr}{m}{n}\selectfont\color{color_29791} }
\put(422.2079,-381.805){\fontsize{11.991}{1}\usefont{T1}{cmr}{m}{n}\selectfont\color{color_29791}s}
\put(428.1075,-381.805){\fontsize{11.991}{1}\usefont{T1}{cmr}{m}{n}\selectfont\color{color_29791}t}
\put(431.4889,-381.805){\fontsize{11.991}{1}\usefont{T1}{cmr}{m}{n}\selectfont\color{color_29791}o}
\put(438.1679,-381.805){\fontsize{11.991}{1}\usefont{T1}{cmr}{m}{n}\selectfont\color{color_29791}ra}
\put(448.8399,-381.805){\fontsize{11.991}{1}\usefont{T1}{cmr}{m}{n}\selectfont\color{color_29791}g}
\put(455.5189,-381.805){\fontsize{11.991}{1}\usefont{T1}{cmr}{m}{n}\selectfont\color{color_29791}e}
\end{picture}
\newpage
\begin{tikzpicture}[overlay]
\path(0pt,0pt);
\filldraw[color_283006][even odd rule]
(-15pt, 10pt) -- (704.972pt, 10pt)
 -- (704.972pt, 10pt)
 -- (704.972pt, -529.972pt)
 -- (704.972pt, -529.972pt)
 -- (-15pt, -529.972pt)
 -- (-15pt, -529.972pt)
 -- (-15pt, 10pt) -- cycle
;
\end{tikzpicture}
\begin{picture}(-5,0)(2.5,0)
\put(-15,-529.972){\includegraphics[width=720pt,height=540pt]{latexImage_43fb9acc39826b86b6493904abaeacaa.png}}
\put(48.411,-60.69601){\fontsize{21.997}{1}\usefont{T1}{cmr}{b}{n}\selectfont\color{color_29791}S}
\put(63.083,-60.69601){\fontsize{21.997}{1}\usefont{T1}{cmr}{b}{n}\selectfont\color{color_29791}t}
\put(70.386,-60.69601){\fontsize{21.997}{1}\usefont{T1}{cmr}{b}{n}\selectfont\color{color_29791}o}
\put(83.78217,-60.69601){\fontsize{21.997}{1}\usefont{T1}{cmr}{b}{n}\selectfont\color{color_29791}r}
\put(92.27302,-60.69601){\fontsize{21.997}{1}\usefont{T1}{cmr}{b}{n}\selectfont\color{color_29791}a}
\put(104.5473,-60.69601){\fontsize{21.997}{1}\usefont{T1}{cmr}{b}{n}\selectfont\color{color_29791}g}
\put(117.9435,-60.69601){\fontsize{21.997}{1}\usefont{T1}{cmr}{b}{n}\selectfont\color{color_29791}e}
\put(130.2178,-60.69601){\fontsize{21.997}{1}\usefont{T1}{cmr}{b}{n}\selectfont\color{color_29791} (}
\put(143.614,-60.69601){\fontsize{21.997}{1}\usefont{T1}{cmr}{b}{n}\selectfont\color{color_29791}c}
\put(155.8884,-60.69601){\fontsize{21.997}{1}\usefont{T1}{cmr}{b}{n}\selectfont\color{color_29791}o}
\put(169.2845,-60.69601){\fontsize{21.997}{1}\usefont{T1}{cmr}{b}{n}\selectfont\color{color_29791}n}
\put(182.6807,-60.69601){\fontsize{21.997}{1}\usefont{T1}{cmr}{b}{n}\selectfont\color{color_29791}t}
\put(190.0497,-60.69601){\fontsize{21.997}{1}\usefont{T1}{cmr}{b}{n}\selectfont\color{color_29791}in}
\put(209.539,-60.69601){\fontsize{21.997}{1}\usefont{T1}{cmr}{b}{n}\selectfont\color{color_29791}u}
\put(223.0232,-60.69601){\fontsize{21.997}{1}\usefont{T1}{cmr}{b}{n}\selectfont\color{color_29791}e}
\put(235.2315,-60.69601){\fontsize{21.997}{1}\usefont{T1}{cmr}{b}{n}\selectfont\color{color_29791}d}
\put(248.6277,-60.69601){\fontsize{21.997}{1}\usefont{T1}{cmr}{b}{n}\selectfont\color{color_29791})}
\put(48.213,-116.51){\fontsize{15.987}{1}\usefont{T1}{cmr}{m}{n}\selectfont\color{color_29791}H}
\put(59.70766,-116.51){\fontsize{15.987}{1}\usefont{T1}{cmr}{m}{n}\selectfont\color{color_29791}P}
\put(70.38697,-116.51){\fontsize{15.987}{1}\usefont{T1}{cmr}{m}{n}\selectfont\color{color_29791}C}
\put(81.88162,-116.51){\fontsize{15.987}{1}\usefont{T1}{cmr}{m}{n}\selectfont\color{color_29791} }
\put(86.35798,-116.51){\fontsize{15.987}{1}\usefont{T1}{cmr}{m}{n}\selectfont\color{color_29791}S}
\put(97.03729,-116.51){\fontsize{15.987}{1}\usefont{T1}{cmr}{m}{n}\selectfont\color{color_29791}t}
\put(101.4337,-116.51){\fontsize{15.987}{1}\usefont{T1}{cmr}{m}{n}\selectfont\color{color_29791}or}
\put(115.5183,-116.51){\fontsize{15.987}{1}\usefont{T1}{cmr}{m}{n}\selectfont\color{color_29791}age }
\put(146.581,-116.51){\fontsize{15.987}{1}\usefont{T1}{cmr}{m}{n}\selectfont\color{color_29791}-}
\put(151.8727,-116.51){\fontsize{15.987}{1}\usefont{T1}{cmr}{m}{n}\selectfont\color{color_29791} }
\put(156.2691,-116.51){\fontsize{15.987}{1}\usefont{T1}{cmr}{m}{n}\selectfont\color{color_29791}/}
\put(160.7455,-116.51){\fontsize{15.987}{1}\usefont{T1}{cmr}{m}{n}\selectfont\color{color_29791}dat}
\put(182.9195,-116.51){\fontsize{15.987}{1}\usefont{T1}{cmr}{m}{n}\selectfont\color{color_29791}a/}
\put(196.2047,-116.51){\fontsize{15.987}{1}\usefont{T1}{cmr}{m}{n}\selectfont\color{color_29791}hpcdat}
\put(244.1497,-116.51){\fontsize{15.987}{1}\usefont{T1}{cmr}{m}{n}\selectfont\color{color_29791}a/}
\put(257.4349,-116.51){\fontsize{15.987}{1}\usefont{T1}{cmr}{m}{n}\selectfont\color{color_29791} }
\put(261.9112,-116.51){\fontsize{15.987}{1}\usefont{T1}{cmr}{m}{n}\selectfont\color{color_29791}(}
\put(267.107,-116.51){\fontsize{15.987}{1}\usefont{T1}{cmr}{m}{n}\selectfont\color{color_29791}us}
\put(284.0692,-116.51){\fontsize{15.987}{1}\usefont{T1}{cmr}{m}{n}\selectfont\color{color_29791}e}
\put(292.8621,-116.51){\fontsize{15.987}{1}\usefont{T1}{cmr}{m}{n}\selectfont\color{color_29791}r}
\put(298.1537,-116.51){\fontsize{15.987}{1}\usefont{T1}{cmr}{m}{n}\selectfont\color{color_29791}s}
\put(306.2272,-116.51){\fontsize{15.987}{1}\usefont{T1}{cmr}{m}{n}\selectfont\color{color_29791},}
\put(311.013,-116.51){\fontsize{15.987}{1}\usefont{T1}{cmr}{m}{n}\selectfont\color{color_29791} }
\put(315.491,-116.51){\fontsize{15.987}{1}\usefont{T1}{cmr}{m}{n}\selectfont\color{color_29791}d}
\put(324.2838,-116.51){\fontsize{15.987}{1}\usefont{T1}{cmr}{m}{n}\selectfont\color{color_29791}at}
\put(337.649,-116.51){\fontsize{15.987}{1}\usefont{T1}{cmr}{m}{n}\selectfont\color{color_29791}a}
\put(346.4418,-116.51){\fontsize{15.987}{1}\usefont{T1}{cmr}{m}{n}\selectfont\color{color_29791})}
\put(351.7335,-116.51){\fontsize{15.987}{1}\usefont{T1}{cmr}{m}{n}\selectfont\color{color_29791} }
\put(356.2099,-116.51){\fontsize{15.987}{1}\usefont{T1}{cmr}{m}{n}\selectfont\color{color_29791}\&}
\put(366.7933,-116.51){\fontsize{15.987}{1}\usefont{T1}{cmr}{m}{n}\selectfont\color{color_29791} }
\put(371.2696,-116.51){\fontsize{15.987}{1}\usefont{T1}{cmr}{m}{n}\selectfont\color{color_29791}/}
\put(375.666,-116.51){\fontsize{15.987}{1}\usefont{T1}{cmr}{m}{n}\selectfont\color{color_29791}dat}
\put(397.84,-116.51){\fontsize{15.987}{1}\usefont{T1}{cmr}{m}{n}\selectfont\color{color_29791}a/}
\put(411.1252,-116.51){\fontsize{15.987}{1}\usefont{T1}{cmr}{m}{n}\selectfont\color{color_29791}hpcf}
\put(441.3726,-116.51){\fontsize{15.987}{1}\usefont{T1}{cmr}{m}{n}\selectfont\color{color_29791}l}
\put(444.9537,-116.51){\fontsize{15.987}{1}\usefont{T1}{cmr}{m}{n}\selectfont\color{color_29791}a}
\put(453.7466,-116.51){\fontsize{15.987}{1}\usefont{T1}{cmr}{m}{n}\selectfont\color{color_29791}s}
\put(461.82,-116.51){\fontsize{15.987}{1}\usefont{T1}{cmr}{m}{n}\selectfont\color{color_29791}h}
\put(84.213,-145.31){\fontsize{11.991}{1}\usefont{T1}{cmr}{m}{n}\selectfont\color{color_29791}•}
\put(98.187,-145.991){\fontsize{11.991}{1}\usefont{T1}{cmr}{m}{n}\selectfont\color{color_29791}A}
\put(106.2689,-145.991){\fontsize{11.991}{1}\usefont{T1}{cmr}{m}{n}\selectfont\color{color_29791}c}
\put(112.1685,-145.991){\fontsize{11.991}{1}\usefont{T1}{cmr}{m}{n}\selectfont\color{color_29791}c}
\put(118.2479,-145.991){\fontsize{11.991}{1}\usefont{T1}{cmr}{m}{n}\selectfont\color{color_29791}e}
\put(124.9269,-145.991){\fontsize{11.991}{1}\usefont{T1}{cmr}{m}{n}\selectfont\color{color_29791}ssi}
\put(139.52,-145.991){\fontsize{11.991}{1}\usefont{T1}{cmr}{m}{n}\selectfont\color{color_29791}b}
\put(146.199,-145.991){\fontsize{11.991}{1}\usefont{T1}{cmr}{m}{n}\selectfont\color{color_29791}l}
\put(148.8849,-145.991){\fontsize{11.991}{1}\usefont{T1}{cmr}{m}{n}\selectfont\color{color_29791}e}
\put(155.5639,-145.991){\fontsize{11.991}{1}\usefont{T1}{cmr}{m}{n}\selectfont\color{color_29791} }
\put(158.8615,-145.991){\fontsize{11.991}{1}\usefont{T1}{cmr}{m}{n}\selectfont\color{color_29791}f}
\put(162.2429,-145.991){\fontsize{11.991}{1}\usefont{T1}{cmr}{m}{n}\selectfont\color{color_29791}ro}
\put(172.9149,-145.991){\fontsize{11.991}{1}\usefont{T1}{cmr}{m}{n}\selectfont\color{color_29791}m }
\put(186.2009,-145.991){\fontsize{11.991}{1}\usefont{T1}{cmr}{m}{n}\selectfont\color{color_29791}n}
\put(192.8799,-145.991){\fontsize{11.991}{1}\usefont{T1}{cmr}{m}{n}\selectfont\color{color_29791}o}
\put(199.6668,-145.991){\fontsize{11.991}{1}\usefont{T1}{cmr}{m}{n}\selectfont\color{color_29791}d}
\put(206.3458,-145.991){\fontsize{11.991}{1}\usefont{T1}{cmr}{m}{n}\selectfont\color{color_29791}e}
\put(213.0248,-145.991){\fontsize{11.991}{1}\usefont{T1}{cmr}{m}{n}\selectfont\color{color_29791}s }
\put(222.4018,-145.991){\fontsize{11.991}{1}\usefont{T1}{cmr}{m}{n}\selectfont\color{color_29791}a}
\put(229.0808,-145.991){\fontsize{11.991}{1}\usefont{T1}{cmr}{m}{n}\selectfont\color{color_29791}n}
\put(235.7598,-145.991){\fontsize{11.991}{1}\usefont{T1}{cmr}{m}{n}\selectfont\color{color_29791}d}
\put(242.4388,-145.991){\fontsize{11.991}{1}\usefont{T1}{cmr}{m}{n}\selectfont\color{color_29791} }
\put(245.8202,-145.991){\fontsize{11.991}{1}\usefont{T1}{cmr}{m}{n}\selectfont\color{color_29791}w}
\put(254.4178,-145.991){\fontsize{11.991}{1}\usefont{T1}{cmr}{m}{n}\selectfont\color{color_29791}o}
\put(261.0967,-145.991){\fontsize{11.991}{1}\usefont{T1}{cmr}{m}{n}\selectfont\color{color_29791}rkst}
\put(280.4622,-145.991){\fontsize{11.991}{1}\usefont{T1}{cmr}{m}{n}\selectfont\color{color_29791}a}
\put(287.1412,-145.991){\fontsize{11.991}{1}\usefont{T1}{cmr}{m}{n}\selectfont\color{color_29791}t}
\put(290.4387,-145.991){\fontsize{11.991}{1}\usefont{T1}{cmr}{m}{n}\selectfont\color{color_29791}i}
\put(293.1247,-145.991){\fontsize{11.991}{1}\usefont{T1}{cmr}{m}{n}\selectfont\color{color_29791}o}
\put(299.8037,-145.991){\fontsize{11.991}{1}\usefont{T1}{cmr}{m}{n}\selectfont\color{color_29791}n}
\put(306.4827,-145.991){\fontsize{11.991}{1}\usefont{T1}{cmr}{m}{n}\selectfont\color{color_29791}s,}
\put(315.8596,-145.991){\fontsize{11.991}{1}\usefont{T1}{cmr}{m}{n}\selectfont\color{color_29791} }
\put(319.1571,-145.991){\fontsize{11.991}{1}\usefont{T1}{cmr}{m}{n}\selectfont\color{color_29791}b}
\put(325.944,-145.991){\fontsize{11.991}{1}\usefont{T1}{cmr}{m}{n}\selectfont\color{color_29791}s}
\put(331.8436,-145.991){\fontsize{11.991}{1}\usefont{T1}{cmr}{m}{n}\selectfont\color{color_29791}l}
\put(334.5296,-145.991){\fontsize{11.991}{1}\usefont{T1}{cmr}{m}{n}\selectfont\color{color_29791}ce}
\put(347.2041,-145.991){\fontsize{11.991}{1}\usefont{T1}{cmr}{m}{n}\selectfont\color{color_29791}n}
\put(353.8831,-145.991){\fontsize{11.991}{1}\usefont{T1}{cmr}{m}{n}\selectfont\color{color_29791}b}
\put(360.562,-145.991){\fontsize{11.991}{1}\usefont{T1}{cmr}{m}{n}\selectfont\color{color_29791},}
\put(363.9435,-145.991){\fontsize{11.991}{1}\usefont{T1}{cmr}{m}{n}\selectfont\color{color_29791} }
\put(367.325,-145.991){\fontsize{11.991}{1}\usefont{T1}{cmr}{m}{n}\selectfont\color{color_29791}b}
\put(374.004,-145.991){\fontsize{11.991}{1}\usefont{T1}{cmr}{m}{n}\selectfont\color{color_29791}sl}
\put(382.6015,-145.991){\fontsize{11.991}{1}\usefont{T1}{cmr}{m}{n}\selectfont\color{color_29791}ce}
\put(395.3839,-145.991){\fontsize{11.991}{1}\usefont{T1}{cmr}{m}{n}\selectfont\color{color_29791}n}
\put(402.0629,-145.991){\fontsize{11.991}{1}\usefont{T1}{cmr}{m}{n}\selectfont\color{color_29791}c.}
\put(84.213,-174.11){\fontsize{11.991}{1}\usefont{T1}{cmr}{m}{n}\selectfont\color{color_29791}•}
\put(98.187,-174.791){\fontsize{11.991}{1}\usefont{T1}{cmr}{m}{n}\selectfont\color{color_29791}U}
\put(106.8685,-174.791){\fontsize{11.991}{1}\usefont{T1}{cmr}{m}{n}\selectfont\color{color_29791}s}
\put(112.7681,-174.791){\fontsize{11.991}{1}\usefont{T1}{cmr}{m}{n}\selectfont\color{color_29791}a}
\put(119.555,-174.791){\fontsize{11.991}{1}\usefont{T1}{cmr}{m}{n}\selectfont\color{color_29791}g}
\put(126.2339,-174.791){\fontsize{11.991}{1}\usefont{T1}{cmr}{m}{n}\selectfont\color{color_29791}e}
\put(132.9129,-174.791){\fontsize{11.991}{1}\usefont{T1}{cmr}{m}{n}\selectfont\color{color_29791} }
\put(136.2944,-174.791){\fontsize{11.991}{1}\usefont{T1}{cmr}{m}{n}\selectfont\color{color_29791}l}
\put(138.8965,-174.791){\fontsize{11.991}{1}\usefont{T1}{cmr}{m}{n}\selectfont\color{color_29791}i}
\put(141.4985,-174.791){\fontsize{11.991}{1}\usefont{T1}{cmr}{m}{n}\selectfont\color{color_29791}mi}
\put(154.173,-174.791){\fontsize{11.991}{1}\usefont{T1}{cmr}{m}{n}\selectfont\color{color_29791}t}
\put(157.4705,-174.791){\fontsize{11.991}{1}\usefont{T1}{cmr}{m}{n}\selectfont\color{color_29791}e}
\put(164.1495,-174.791){\fontsize{11.991}{1}\usefont{T1}{cmr}{m}{n}\selectfont\color{color_29791}d}
\put(170.9364,-174.791){\fontsize{11.991}{1}\usefont{T1}{cmr}{m}{n}\selectfont\color{color_29791} }
\put(174.2339,-174.791){\fontsize{11.991}{1}\usefont{T1}{cmr}{m}{n}\selectfont\color{color_29791}vi}
\put(182.9154,-174.791){\fontsize{11.991}{1}\usefont{T1}{cmr}{m}{n}\selectfont\color{color_29791}a}
\put(189.5944,-174.791){\fontsize{11.991}{1}\usefont{T1}{cmr}{m}{n}\selectfont\color{color_29791} }
\put(192.892,-174.791){\fontsize{11.991}{1}\usefont{T1}{cmr}{m}{n}\selectfont\color{color_29791}q}
\put(199.6789,-174.791){\fontsize{11.991}{1}\usefont{T1}{cmr}{m}{n}\selectfont\color{color_29791}u}
\put(206.3578,-174.791){\fontsize{11.991}{1}\usefont{T1}{cmr}{m}{n}\selectfont\color{color_29791}o}
\put(213.0368,-174.791){\fontsize{11.991}{1}\usefont{T1}{cmr}{m}{n}\selectfont\color{color_29791}t}
\put(216.3344,-174.791){\fontsize{11.991}{1}\usefont{T1}{cmr}{m}{n}\selectfont\color{color_29791}a}
\put(223.1213,-174.791){\fontsize{11.991}{1}\usefont{T1}{cmr}{m}{n}\selectfont\color{color_29791}s}
\put(48.213,-208.409){\fontsize{15.987}{1}\usefont{T1}{cmr}{m}{n}\selectfont\color{color_29791}Q}
\put(60.60293,-208.409){\fontsize{15.987}{1}\usefont{T1}{cmr}{m}{n}\selectfont\color{color_29791}u}
\put(69.39577,-208.409){\fontsize{15.987}{1}\usefont{T1}{cmr}{m}{n}\selectfont\color{color_29791}ot}
\put(82.76091,-208.409){\fontsize{15.987}{1}\usefont{T1}{cmr}{m}{n}\selectfont\color{color_29791}a}
\put(91.55376,-208.409){\fontsize{15.987}{1}\usefont{T1}{cmr}{m}{n}\selectfont\color{color_29791}s}
\put(84.213,-237.209){\fontsize{11.991}{1}\usefont{T1}{cmr}{m}{n}\selectfont\color{color_29791}•}
\put(97.706,-237.89){\fontsize{11.991}{1}\usefont{T1}{cmr}{m}{n}\selectfont\color{color_29791}O}
\put(107.083,-237.89){\fontsize{11.991}{1}\usefont{T1}{cmr}{m}{n}\selectfont\color{color_29791}n}
\put(113.7619,-237.89){\fontsize{11.991}{1}\usefont{T1}{cmr}{m}{n}\selectfont\color{color_29791} }
\put(117.1434,-237.89){\fontsize{11.991}{1}\usefont{T1}{cmr}{m}{n}\selectfont\color{color_29791}H}
\put(125.741,-237.89){\fontsize{11.991}{1}\usefont{T1}{cmr}{m}{n}\selectfont\color{color_29791}PC}
\put(142.4084,-237.89){\fontsize{11.991}{1}\usefont{T1}{cmr}{m}{n}\selectfont\color{color_29791} }
\put(145.706,-237.89){\fontsize{11.991}{1}\usefont{T1}{cmr}{m}{n}\selectfont\color{color_29791}yo}
\put(158.3805,-237.89){\fontsize{11.991}{1}\usefont{T1}{cmr}{m}{n}\selectfont\color{color_29791}u}
\put(165.1674,-237.89){\fontsize{11.991}{1}\usefont{T1}{cmr}{m}{n}\selectfont\color{color_29791} }
\put(168.4649,-237.89){\fontsize{11.991}{1}\usefont{T1}{cmr}{m}{n}\selectfont\color{color_29791}ca}
\put(181.2473,-237.89){\fontsize{11.991}{1}\usefont{T1}{cmr}{m}{n}\selectfont\color{color_29791}n}
\put(187.9263,-237.89){\fontsize{11.991}{1}\usefont{T1}{cmr}{m}{n}\selectfont\color{color_29791} }
\put(191.2238,-237.89){\fontsize{11.991}{1}\usefont{T1}{cmr}{m}{n}\selectfont\color{color_29791}ch}
\put(203.8983,-237.89){\fontsize{11.991}{1}\usefont{T1}{cmr}{m}{n}\selectfont\color{color_29791}e}
\put(210.6852,-237.89){\fontsize{11.991}{1}\usefont{T1}{cmr}{m}{n}\selectfont\color{color_29791}ck }
\put(225.9738,-237.89){\fontsize{11.991}{1}\usefont{T1}{cmr}{m}{n}\selectfont\color{color_29791}yo}
\put(238.7562,-237.89){\fontsize{11.991}{1}\usefont{T1}{cmr}{m}{n}\selectfont\color{color_29791}u}
\put(245.4352,-237.89){\fontsize{11.991}{1}\usefont{T1}{cmr}{m}{n}\selectfont\color{color_29791}r}
\put(249.3322,-237.89){\fontsize{11.991}{1}\usefont{T1}{cmr}{m}{n}\selectfont\color{color_29791} }
\put(252.7137,-237.89){\fontsize{11.991}{1}\usefont{T1}{cmr}{m}{n}\selectfont\color{color_29791}q}
\put(259.3927,-237.89){\fontsize{11.991}{1}\usefont{T1}{cmr}{m}{n}\selectfont\color{color_29791}u}
\put(266.0717,-237.89){\fontsize{11.991}{1}\usefont{T1}{cmr}{m}{n}\selectfont\color{color_29791}o}
\put(272.7506,-237.89){\fontsize{11.991}{1}\usefont{T1}{cmr}{m}{n}\selectfont\color{color_29791}t}
\put(276.1321,-237.89){\fontsize{11.991}{1}\usefont{T1}{cmr}{m}{n}\selectfont\color{color_29791}a}
\put(282.8111,-237.89){\fontsize{11.991}{1}\usefont{T1}{cmr}{m}{n}\selectfont\color{color_29791}s }
\put(292.188,-237.89){\fontsize{11.991}{1}\usefont{T1}{cmr}{m}{n}\selectfont\color{color_29791}u}
\put(298.867,-237.89){\fontsize{11.991}{1}\usefont{T1}{cmr}{m}{n}\selectfont\color{color_29791}si}
\put(307.5485,-237.89){\fontsize{11.991}{1}\usefont{T1}{cmr}{m}{n}\selectfont\color{color_29791}n}
\put(314.2275,-237.89){\fontsize{11.991}{1}\usefont{T1}{cmr}{m}{n}\selectfont\color{color_29791}g}
\put(320.9065,-237.89){\fontsize{11.991}{1}\usefont{T1}{cmr}{m}{n}\selectfont\color{color_29791}:}
\put(324.204,-237.89){\fontsize{11.991}{1}\usefont{T1}{cmr}{m}{n}\selectfont\color{color_29791} }
\put(327.5855,-237.89){\fontsize{11.991}{1}\usefont{T1}{cmr}{m}{n}\selectfont\color{color_29791}myq}
\put(350.2485,-237.89){\fontsize{11.991}{1}\usefont{T1}{cmr}{m}{n}\selectfont\color{color_29791}u}
\put(356.9275,-237.89){\fontsize{11.991}{1}\usefont{T1}{cmr}{m}{n}\selectfont\color{color_29791}o}
\put(363.6064,-237.89){\fontsize{11.991}{1}\usefont{T1}{cmr}{m}{n}\selectfont\color{color_29791}t}
\put(366.9879,-237.89){\fontsize{11.991}{1}\usefont{T1}{cmr}{m}{n}\selectfont\color{color_29791}a}
\put(84.213,-266.009){\fontsize{11.991}{1}\usefont{T1}{cmr}{m}{n}\selectfont\color{color_29791}•}
\put(97.706,-266.69){\fontsize{11.991}{1}\usefont{T1}{cmr}{m}{n}\selectfont\color{color_29791}N}
\put(106.3875,-266.69){\fontsize{11.991}{1}\usefont{T1}{cmr}{m}{n}\selectfont\color{color_29791}e}
\put(113.0665,-266.69){\fontsize{11.991}{1}\usefont{T1}{cmr}{m}{n}\selectfont\color{color_29791}e}
\put(119.7455,-266.69){\fontsize{11.991}{1}\usefont{T1}{cmr}{m}{n}\selectfont\color{color_29791}d}
\put(126.4244,-266.69){\fontsize{11.991}{1}\usefont{T1}{cmr}{m}{n}\selectfont\color{color_29791} }
\put(129.8059,-266.69){\fontsize{11.991}{1}\usefont{T1}{cmr}{m}{n}\selectfont\color{color_29791}m}
\put(139.6985,-266.69){\fontsize{11.991}{1}\usefont{T1}{cmr}{m}{n}\selectfont\color{color_29791}o}
\put(146.3775,-266.69){\fontsize{11.991}{1}\usefont{T1}{cmr}{m}{n}\selectfont\color{color_29791}re}
\put(157.0495,-266.69){\fontsize{11.991}{1}\usefont{T1}{cmr}{m}{n}\selectfont\color{color_29791} }
\put(160.4309,-266.69){\fontsize{11.991}{1}\usefont{T1}{cmr}{m}{n}\selectfont\color{color_29791}sp}
\put(173.1054,-266.69){\fontsize{11.991}{1}\usefont{T1}{cmr}{m}{n}\selectfont\color{color_29791}a}
\put(179.7844,-266.69){\fontsize{11.991}{1}\usefont{T1}{cmr}{m}{n}\selectfont\color{color_29791}ce}
\put(192.5668,-266.69){\fontsize{11.991}{1}\usefont{T1}{cmr}{m}{n}\selectfont\color{color_29791} }
\put(195.8643,-266.69){\fontsize{11.991}{1}\usefont{T1}{cmr}{m}{n}\selectfont\color{color_29791}-}
\put(200.093,-266.69){\fontsize{11.991}{1}\usefont{T1}{cmr}{m}{n}\selectfont\color{color_29791} }
\put(203.409,-266.69){\fontsize{11.991}{1}\usefont{T1}{cmr}{m}{n}\selectfont\color{color_29791}co}
\put(216.0835,-266.69){\fontsize{11.991}{1}\usefont{T1}{cmr}{m}{n}\selectfont\color{color_29791}n}
\put(222.8704,-266.69){\fontsize{11.991}{1}\usefont{T1}{cmr}{m}{n}\selectfont\color{color_29791}t}
\put(226.1679,-266.69){\fontsize{11.991}{1}\usefont{T1}{cmr}{m}{n}\selectfont\color{color_29791}a}
\put(232.8469,-266.69){\fontsize{11.991}{1}\usefont{T1}{cmr}{m}{n}\selectfont\color{color_29791}c}
\put(238.9263,-266.69){\fontsize{11.991}{1}\usefont{T1}{cmr}{m}{n}\selectfont\color{color_29791}t}
\put(242.2239,-266.69){\fontsize{11.991}{1}\usefont{T1}{cmr}{m}{n}\selectfont\color{color_29791} }
\put(245.6053,-266.69){\fontsize{11.991}{1}\usefont{T1}{cmr}{m}{n}\selectfont\color{color_29791}t}
\put(248.9029,-266.69){\fontsize{11.991}{1}\usefont{T1}{cmr}{m}{n}\selectfont\color{color_29791}h}
\put(255.5818,-266.69){\fontsize{11.991}{1}\usefont{T1}{cmr}{m}{n}\selectfont\color{color_29791}e}
\put(262.3687,-266.69){\fontsize{11.991}{1}\usefont{T1}{cmr}{m}{n}\selectfont\color{color_29791} }
\put(265.6663,-266.69){\fontsize{11.991}{1}\usefont{T1}{cmr}{m}{n}\selectfont\color{color_29791}se}
\put(278.4487,-266.69){\fontsize{11.991}{1}\usefont{T1}{cmr}{m}{n}\selectfont\color{color_29791}r}
\put(282.3457,-266.69){\fontsize{11.991}{1}\usefont{T1}{cmr}{m}{n}\selectfont\color{color_29791}vi}
\put(290.9433,-266.69){\fontsize{11.991}{1}\usefont{T1}{cmr}{m}{n}\selectfont\color{color_29791}c}
\put(297.0227,-266.69){\fontsize{11.991}{1}\usefont{T1}{cmr}{m}{n}\selectfont\color{color_29791}e}
\put(303.7017,-266.69){\fontsize{11.991}{1}\usefont{T1}{cmr}{m}{n}\selectfont\color{color_29791} }
\put(306.9992,-266.69){\fontsize{11.991}{1}\usefont{T1}{cmr}{m}{n}\selectfont\color{color_29791}d}
\put(313.6782,-266.69){\fontsize{11.991}{1}\usefont{T1}{cmr}{m}{n}\selectfont\color{color_29791}e}
\put(320.4651,-266.69){\fontsize{11.991}{1}\usefont{T1}{cmr}{m}{n}\selectfont\color{color_29791}sk}
\put(48.213,-300.309){\fontsize{15.987}{1}\usefont{T1}{cmr}{m}{n}\selectfont\color{color_29791}B}
\put(58.89232,-300.309){\fontsize{15.987}{1}\usefont{T1}{cmr}{m}{n}\selectfont\color{color_29791}ackups}
\put(109.611,-300.309){\fontsize{15.987}{1}\usefont{T1}{cmr}{m}{n}\selectfont\color{color_29791} }
\put(84.213,-329.109){\fontsize{11.991}{1}\usefont{T1}{cmr}{m}{n}\selectfont\color{color_29791}•}
\put(98.187,-329.789){\fontsize{11.991}{1}\usefont{T1}{cmr}{m}{n}\selectfont\color{color_29791}D}
\put(106.8685,-329.789){\fontsize{11.991}{1}\usefont{T1}{cmr}{m}{n}\selectfont\color{color_29791}a}
\put(113.5475,-329.789){\fontsize{11.991}{1}\usefont{T1}{cmr}{m}{n}\selectfont\color{color_29791}i}
\put(116.1495,-329.789){\fontsize{11.991}{1}\usefont{T1}{cmr}{m}{n}\selectfont\color{color_29791}l}
\put(118.8355,-329.789){\fontsize{11.991}{1}\usefont{T1}{cmr}{m}{n}\selectfont\color{color_29791}y }
\put(128.1285,-329.789){\fontsize{11.991}{1}\usefont{T1}{cmr}{m}{n}\selectfont\color{color_29791}a}
\put(134.8075,-329.789){\fontsize{11.991}{1}\usefont{T1}{cmr}{m}{n}\selectfont\color{color_29791}t}
\put(138.189,-329.789){\fontsize{11.991}{1}\usefont{T1}{cmr}{m}{n}\selectfont\color{color_29791} }
\put(141.4865,-329.789){\fontsize{11.991}{1}\usefont{T1}{cmr}{m}{n}\selectfont\color{color_29791}6}
\put(148.2734,-329.789){\fontsize{11.991}{1}\usefont{T1}{cmr}{m}{n}\selectfont\color{color_29791}p}
\put(154.9524,-329.789){\fontsize{11.991}{1}\usefont{T1}{cmr}{m}{n}\selectfont\color{color_29791}m}
\put(84.213,-357.909){\fontsize{11.991}{1}\usefont{T1}{cmr}{m}{n}\selectfont\color{color_29791}•}
\put(98.187,-358.589){\fontsize{11.991}{1}\usefont{T1}{cmr}{m}{n}\selectfont\color{color_29791}A}
\put(106.2689,-358.589){\fontsize{11.991}{1}\usefont{T1}{cmr}{m}{n}\selectfont\color{color_29791}l}
\put(108.871,-358.589){\fontsize{11.991}{1}\usefont{T1}{cmr}{m}{n}\selectfont\color{color_29791}l}
\put(111.473,-358.589){\fontsize{11.991}{1}\usefont{T1}{cmr}{m}{n}\selectfont\color{color_29791} }
\put(114.8545,-358.589){\fontsize{11.991}{1}\usefont{T1}{cmr}{m}{n}\selectfont\color{color_29791}SA}
\put(130.9224,-358.589){\fontsize{11.991}{1}\usefont{T1}{cmr}{m}{n}\selectfont\color{color_29791}N}
\put(139.52,-358.589){\fontsize{11.991}{1}\usefont{T1}{cmr}{m}{n}\selectfont\color{color_29791} }
\put(142.9014,-358.589){\fontsize{11.991}{1}\usefont{T1}{cmr}{m}{n}\selectfont\color{color_29791}a}
\put(149.5804,-358.589){\fontsize{11.991}{1}\usefont{T1}{cmr}{m}{n}\selectfont\color{color_29791}n}
\put(156.2594,-358.589){\fontsize{11.991}{1}\usefont{T1}{cmr}{m}{n}\selectfont\color{color_29791}d}
\put(162.9384,-358.589){\fontsize{11.991}{1}\usefont{T1}{cmr}{m}{n}\selectfont\color{color_29791} }
\put(166.3199,-358.589){\fontsize{11.991}{1}\usefont{T1}{cmr}{m}{n}\selectfont\color{color_29791}H}
\put(174.9174,-358.589){\fontsize{11.991}{1}\usefont{T1}{cmr}{m}{n}\selectfont\color{color_29791}P}
\put(182.9993,-358.589){\fontsize{11.991}{1}\usefont{T1}{cmr}{m}{n}\selectfont\color{color_29791}C}
\put(191.5969,-358.589){\fontsize{11.991}{1}\usefont{T1}{cmr}{m}{n}\selectfont\color{color_29791} }
\put(194.9784,-358.589){\fontsize{11.991}{1}\usefont{T1}{cmr}{m}{n}\selectfont\color{color_29791}vo}
\put(207.6528,-358.589){\fontsize{11.991}{1}\usefont{T1}{cmr}{m}{n}\selectfont\color{color_29791}l}
\put(210.2549,-358.589){\fontsize{11.991}{1}\usefont{T1}{cmr}{m}{n}\selectfont\color{color_29791}u}
\put(216.9339,-358.589){\fontsize{11.991}{1}\usefont{T1}{cmr}{m}{n}\selectfont\color{color_29791}me}
\put(233.6014,-358.589){\fontsize{11.991}{1}\usefont{T1}{cmr}{m}{n}\selectfont\color{color_29791}s }
\put(242.9783,-358.589){\fontsize{11.991}{1}\usefont{T1}{cmr}{m}{n}\selectfont\color{color_29791}b}
\put(249.6573,-358.589){\fontsize{11.991}{1}\usefont{T1}{cmr}{m}{n}\selectfont\color{color_29791}a}
\put(256.3363,-358.589){\fontsize{11.991}{1}\usefont{T1}{cmr}{m}{n}\selectfont\color{color_29791}cke}
\put(275.1142,-358.589){\fontsize{11.991}{1}\usefont{T1}{cmr}{m}{n}\selectfont\color{color_29791}d}
\put(281.7932,-358.589){\fontsize{11.991}{1}\usefont{T1}{cmr}{m}{n}\selectfont\color{color_29791} }
\put(285.0907,-358.589){\fontsize{11.991}{1}\usefont{T1}{cmr}{m}{n}\selectfont\color{color_29791}u}
\put(291.7697,-358.589){\fontsize{11.991}{1}\usefont{T1}{cmr}{m}{n}\selectfont\color{color_29791}p}
\put(84.213,-386.709){\fontsize{11.991}{1}\usefont{T1}{cmr}{m}{n}\selectfont\color{color_29791}•}
\put(98.187,-387.389){\fontsize{11.991}{1}\usefont{T1}{cmr}{m}{n}\selectfont\color{color_29791}B}
\put(106.2689,-387.389){\fontsize{11.991}{1}\usefont{T1}{cmr}{m}{n}\selectfont\color{color_29791}a}
\put(112.9479,-387.389){\fontsize{11.991}{1}\usefont{T1}{cmr}{m}{n}\selectfont\color{color_29791}cku}
\put(131.6179,-387.389){\fontsize{11.991}{1}\usefont{T1}{cmr}{m}{n}\selectfont\color{color_29791}p}
\put(138.2969,-387.389){\fontsize{11.991}{1}\usefont{T1}{cmr}{m}{n}\selectfont\color{color_29791}s }
\put(147.6739,-387.389){\fontsize{11.991}{1}\usefont{T1}{cmr}{m}{n}\selectfont\color{color_29791}a}
\put(154.3528,-387.389){\fontsize{11.991}{1}\usefont{T1}{cmr}{m}{n}\selectfont\color{color_29791}re}
\put(165.0248,-387.389){\fontsize{11.991}{1}\usefont{T1}{cmr}{m}{n}\selectfont\color{color_29791} }
\put(168.3224,-387.389){\fontsize{11.991}{1}\usefont{T1}{cmr}{m}{n}\selectfont\color{color_29791}b}
\put(175.1093,-387.389){\fontsize{11.991}{1}\usefont{T1}{cmr}{m}{n}\selectfont\color{color_29791}o}
\put(181.7883,-387.389){\fontsize{11.991}{1}\usefont{T1}{cmr}{m}{n}\selectfont\color{color_29791}t}
\put(185.0858,-387.389){\fontsize{11.991}{1}\usefont{T1}{cmr}{m}{n}\selectfont\color{color_29791}h}
\put(191.7648,-387.389){\fontsize{11.991}{1}\usefont{T1}{cmr}{m}{n}\selectfont\color{color_29791} }
\put(195.1462,-387.389){\fontsize{11.991}{1}\usefont{T1}{cmr}{m}{n}\selectfont\color{color_29791}o}
\put(201.8252,-387.389){\fontsize{11.991}{1}\usefont{T1}{cmr}{m}{n}\selectfont\color{color_29791}n}
\put(208.6121,-387.389){\fontsize{11.991}{1}\usefont{T1}{cmr}{m}{n}\selectfont\color{color_29791}si}
\put(217.2097,-387.389){\fontsize{11.991}{1}\usefont{T1}{cmr}{m}{n}\selectfont\color{color_29791}t}
\put(220.5072,-387.389){\fontsize{11.991}{1}\usefont{T1}{cmr}{m}{n}\selectfont\color{color_29791}e}
\put(227.2941,-387.389){\fontsize{11.991}{1}\usefont{T1}{cmr}{m}{n}\selectfont\color{color_29791} }
\put(230.5917,-387.389){\fontsize{11.991}{1}\usefont{T1}{cmr}{m}{n}\selectfont\color{color_29791}a}
\put(237.3786,-387.389){\fontsize{11.991}{1}\usefont{T1}{cmr}{m}{n}\selectfont\color{color_29791}n}
\put(244.0576,-387.389){\fontsize{11.991}{1}\usefont{T1}{cmr}{m}{n}\selectfont\color{color_29791}d}
\put(250.7365,-387.389){\fontsize{11.991}{1}\usefont{T1}{cmr}{m}{n}\selectfont\color{color_29791} }
\put(254.118,-387.389){\fontsize{11.991}{1}\usefont{T1}{cmr}{m}{n}\selectfont\color{color_29791}o}
\put(260.797,-387.389){\fontsize{11.991}{1}\usefont{T1}{cmr}{m}{n}\selectfont\color{color_29791}f}
\put(263.8907,-387.389){\fontsize{11.991}{1}\usefont{T1}{cmr}{m}{n}\selectfont\color{color_29791}f}
\put(267.1882,-387.389){\fontsize{11.991}{1}\usefont{T1}{cmr}{m}{n}\selectfont\color{color_29791}s}
\put(273.2676,-387.389){\fontsize{11.991}{1}\usefont{T1}{cmr}{m}{n}\selectfont\color{color_29791}i}
\put(275.8697,-387.389){\fontsize{11.991}{1}\usefont{T1}{cmr}{m}{n}\selectfont\color{color_29791}t}
\put(279.2512,-387.389){\fontsize{11.991}{1}\usefont{T1}{cmr}{m}{n}\selectfont\color{color_29791}e}
\put(285.9301,-387.389){\fontsize{11.991}{1}\usefont{T1}{cmr}{m}{n}\selectfont\color{color_29791},}
\put(289.696,-387.389){\fontsize{11.991}{1}\usefont{T1}{cmr}{m}{n}\selectfont\color{color_29791} }
\put(293.098,-387.389){\fontsize{11.991}{1}\usefont{T1}{cmr}{m}{n}\selectfont\color{color_29791}vi}
\put(301.6955,-387.389){\fontsize{11.991}{1}\usefont{T1}{cmr}{m}{n}\selectfont\color{color_29791}a}
\put(308.4824,-387.389){\fontsize{11.991}{1}\usefont{T1}{cmr}{m}{n}\selectfont\color{color_29791} }
\put(311.7799,-387.389){\fontsize{11.991}{1}\usefont{T1}{cmr}{m}{n}\selectfont\color{color_29791}t}
\put(315.1614,-387.389){\fontsize{11.991}{1}\usefont{T1}{cmr}{m}{n}\selectfont\color{color_29791}a}
\put(321.8404,-387.389){\fontsize{11.991}{1}\usefont{T1}{cmr}{m}{n}\selectfont\color{color_29791}p}
\put(328.5194,-387.389){\fontsize{11.991}{1}\usefont{T1}{cmr}{m}{n}\selectfont\color{color_29791}e}
\put(335.1984,-387.389){\fontsize{11.991}{1}\usefont{T1}{cmr}{m}{n}\selectfont\color{color_29791}s }
\put(344.5753,-387.389){\fontsize{11.991}{1}\usefont{T1}{cmr}{m}{n}\selectfont\color{color_29791}\& }
\put(355.9428,-387.389){\fontsize{11.991}{1}\usefont{T1}{cmr}{m}{n}\selectfont\color{color_29791}d}
\put(362.6218,-387.389){\fontsize{11.991}{1}\usefont{T1}{cmr}{m}{n}\selectfont\color{color_29791}i}
\put(365.3077,-387.389){\fontsize{11.991}{1}\usefont{T1}{cmr}{m}{n}\selectfont\color{color_29791}s}
\put(371.2073,-387.389){\fontsize{11.991}{1}\usefont{T1}{cmr}{m}{n}\selectfont\color{color_29791}k}
\put(84.213,-415.509){\fontsize{11.991}{1}\usefont{T1}{cmr}{m}{n}\selectfont\color{color_29791}•}
\put(98.187,-416.104){\fontsize{11.991}{1}\usefont{T1}{cmr}{m}{n}\selectfont\color{color_29791}I}
\put(101.5685,-416.104){\fontsize{11.991}{1}\usefont{T1}{cmr}{m}{n}\selectfont\color{color_29791}f}
\put(104.866,-416.104){\fontsize{11.991}{1}\usefont{T1}{cmr}{m}{n}\selectfont\color{color_29791} }
\put(108.2474,-416.104){\fontsize{11.991}{1}\usefont{T1}{cmr}{m}{n}\selectfont\color{color_29791}yo}
\put(120.9219,-416.104){\fontsize{11.991}{1}\usefont{T1}{cmr}{m}{n}\selectfont\color{color_29791}u}
\put(127.6009,-416.104){\fontsize{11.991}{1}\usefont{T1}{cmr}{m}{n}\selectfont\color{color_29791} }
\put(130.9824,-416.104){\fontsize{11.991}{1}\usefont{T1}{cmr}{m}{n}\selectfont\color{color_29791}n}
\put(137.6614,-416.104){\fontsize{11.991}{1}\usefont{T1}{cmr}{m}{n}\selectfont\color{color_29791}e}
\put(144.3404,-416.104){\fontsize{11.991}{1}\usefont{T1}{cmr}{m}{n}\selectfont\color{color_29791}e}
\put(151.1273,-416.104){\fontsize{11.991}{1}\usefont{T1}{cmr}{m}{n}\selectfont\color{color_29791}d}
\put(157.8063,-416.104){\fontsize{11.991}{1}\usefont{T1}{cmr}{m}{n}\selectfont\color{color_29791} }
\put(161.1038,-416.104){\fontsize{11.991}{1}\usefont{T1}{cmr}{m}{n}\selectfont\color{color_29791}a}
\put(167.8907,-416.104){\fontsize{11.991}{1}\usefont{T1}{cmr}{m}{n}\selectfont\color{color_29791} }
\put(171.1882,-416.104){\fontsize{11.991}{1}\usefont{T1}{cmr}{m}{n}\selectfont\color{color_29791}f}
\put(174.5697,-416.104){\fontsize{11.991}{1}\usefont{T1}{cmr}{m}{n}\selectfont\color{color_29791}i}
\put(177.1718,-416.104){\fontsize{11.991}{1}\usefont{T1}{cmr}{m}{n}\selectfont\color{color_29791}l}
\put(179.8577,-416.104){\fontsize{11.991}{1}\usefont{T1}{cmr}{m}{n}\selectfont\color{color_29791}e}
\put(186.5367,-416.104){\fontsize{11.991}{1}\usefont{T1}{cmr}{m}{n}\selectfont\color{color_29791} }
\put(189.8343,-416.104){\fontsize{11.991}{1}\usefont{T1}{cmr}{m}{n}\selectfont\color{color_29791}re}
\put(200.5063,-416.104){\fontsize{11.991}{1}\usefont{T1}{cmr}{m}{n}\selectfont\color{color_29791}st}
\put(209.8832,-416.104){\fontsize{11.991}{1}\usefont{T1}{cmr}{m}{n}\selectfont\color{color_29791}o}
\put(216.5622,-416.104){\fontsize{11.991}{1}\usefont{T1}{cmr}{m}{n}\selectfont\color{color_29791}r}
\put(220.4593,-416.104){\fontsize{11.991}{1}\usefont{T1}{cmr}{m}{n}\selectfont\color{color_29791}e}
\put(227.2462,-416.104){\fontsize{11.991}{1}\usefont{T1}{cmr}{m}{n}\selectfont\color{color_29791}d}
\put(233.9252,-416.104){\fontsize{11.991}{1}\usefont{T1}{cmr}{m}{n}\selectfont\color{color_29791},}
\put(237.3066,-416.104){\fontsize{11.991}{1}\usefont{T1}{cmr}{m}{n}\selectfont\color{color_29791} }
\put(240.6042,-416.104){\fontsize{11.991}{1}\usefont{T1}{cmr}{m}{n}\selectfont\color{color_29791}co}
\put(253.2787,-416.104){\fontsize{11.991}{1}\usefont{T1}{cmr}{m}{n}\selectfont\color{color_29791}n}
\put(259.9576,-416.104){\fontsize{11.991}{1}\usefont{T1}{cmr}{m}{n}\selectfont\color{color_29791}t}
\put(263.3391,-416.104){\fontsize{11.991}{1}\usefont{T1}{cmr}{m}{n}\selectfont\color{color_29791}a}
\put(270.0181,-416.104){\fontsize{11.991}{1}\usefont{T1}{cmr}{m}{n}\selectfont\color{color_29791}ct}
\put(279.3951,-416.104){\fontsize{11.991}{1}\usefont{T1}{cmr}{m}{n}\selectfont\color{color_29791} }
\put(282.7765,-416.104){\fontsize{11.991}{1}\usefont{T1}{cmr}{m}{n}\selectfont\color{color_29791}t}
\put(286.074,-416.104){\fontsize{11.991}{1}\usefont{T1}{cmr}{m}{n}\selectfont\color{color_29791}h}
\put(292.753,-416.104){\fontsize{11.991}{1}\usefont{T1}{cmr}{m}{n}\selectfont\color{color_29791}e}
\put(299.432,-416.104){\fontsize{11.991}{1}\usefont{T1}{cmr}{m}{n}\selectfont\color{color_29791} }
\put(302.8135,-416.104){\fontsize{11.991}{1}\usefont{T1}{cmr}{m}{n}\selectfont\color{color_29791}se}
\put(315.488,-416.104){\fontsize{11.991}{1}\usefont{T1}{cmr}{m}{n}\selectfont\color{color_29791}rvi}
\put(328.0785,-416.104){\fontsize{11.991}{1}\usefont{T1}{cmr}{m}{n}\selectfont\color{color_29791}ce}
\put(340.8609,-416.104){\fontsize{11.991}{1}\usefont{T1}{cmr}{m}{n}\selectfont\color{color_29791} }
\put(344.1584,-416.104){\fontsize{11.991}{1}\usefont{T1}{cmr}{m}{n}\selectfont\color{color_29791}d}
\put(350.8374,-416.104){\fontsize{11.991}{1}\usefont{T1}{cmr}{m}{n}\selectfont\color{color_29791}e}
\put(357.6243,-416.104){\fontsize{11.991}{1}\usefont{T1}{cmr}{m}{n}\selectfont\color{color_29791}s}
\put(363.5239,-416.104){\fontsize{11.991}{1}\usefont{T1}{cmr}{m}{n}\selectfont\color{color_29791}k}
\end{picture}
\newpage
\begin{tikzpicture}[overlay]
\path(0pt,0pt);
\filldraw[color_283006][even odd rule]
(-15pt, 10pt) -- (704.972pt, 10pt)
 -- (704.972pt, 10pt)
 -- (704.972pt, -529.972pt)
 -- (704.972pt, -529.972pt)
 -- (-15pt, -529.972pt)
 -- (-15pt, -529.972pt)
 -- (-15pt, 10pt) -- cycle
;
\end{tikzpicture}
\begin{picture}(-5,0)(2.5,0)
\put(-15,-529.972){\includegraphics[width=720pt,height=540pt]{latexImage_d5964f0fe7de530ea43707adcb11ad51.png}}
\put(48.411,-60.69601){\fontsize{21.997}{1}\usefont{T1}{cmr}{b}{n}\selectfont\color{color_29791}Da}
\put(76.50117,-60.69601){\fontsize{21.997}{1}\usefont{T1}{cmr}{b}{n}\selectfont\color{color_29791}t}
\put(83.80417,-60.69601){\fontsize{21.997}{1}\usefont{T1}{cmr}{b}{n}\selectfont\color{color_29791}a}
\put(96.07849,-60.69601){\fontsize{21.997}{1}\usefont{T1}{cmr}{b}{n}\selectfont\color{color_29791} }
\put(101.2698,-60.69601){\fontsize{21.997}{1}\usefont{T1}{cmr}{b}{n}\selectfont\color{color_29791}Ac}
\put(129.4259,-60.69601){\fontsize{21.997}{1}\usefont{T1}{cmr}{b}{n}\selectfont\color{color_29791}c}
\put(141.6343,-60.69601){\fontsize{21.997}{1}\usefont{T1}{cmr}{b}{n}\selectfont\color{color_29791}e}
\put(153.8426,-60.69601){\fontsize{21.997}{1}\usefont{T1}{cmr}{b}{n}\selectfont\color{color_29791}s}
\put(166.1169,-60.69601){\fontsize{21.997}{1}\usefont{T1}{cmr}{b}{n}\selectfont\color{color_29791}s}
\put(178.3253,-60.69601){\fontsize{21.997}{1}\usefont{T1}{cmr}{b}{n}\selectfont\color{color_29791} a}
\put(196.6268,-60.69601){\fontsize{21.997}{1}\usefont{T1}{cmr}{b}{n}\selectfont\color{color_29791}n}
\put(210.1109,-60.69601){\fontsize{21.997}{1}\usefont{T1}{cmr}{b}{n}\selectfont\color{color_29791}d}
\put(223.5071,-60.69601){\fontsize{21.997}{1}\usefont{T1}{cmr}{b}{n}\selectfont\color{color_29791} }
\put(229.6882,-60.69601){\fontsize{21.997}{1}\usefont{T1}{cmr}{b}{n}\selectfont\color{color_29791}T}
\put(241.8966,-60.69601){\fontsize{21.997}{1}\usefont{T1}{cmr}{b}{n}\selectfont\color{color_29791}r}
\put(250.3874,-60.69601){\fontsize{21.997}{1}\usefont{T1}{cmr}{b}{n}\selectfont\color{color_29791}a}
\put(262.6617,-60.69601){\fontsize{21.997}{1}\usefont{T1}{cmr}{b}{n}\selectfont\color{color_29791}n}
\put(276.0579,-60.69601){\fontsize{21.997}{1}\usefont{T1}{cmr}{b}{n}\selectfont\color{color_29791}s}
\put(288.2663,-60.69601){\fontsize{21.997}{1}\usefont{T1}{cmr}{b}{n}\selectfont\color{color_29791}f}
\put(295.6353,-60.69601){\fontsize{21.997}{1}\usefont{T1}{cmr}{b}{n}\selectfont\color{color_29791}e}
\put(307.8436,-60.69601){\fontsize{21.997}{1}\usefont{T1}{cmr}{b}{n}\selectfont\color{color_29791}r}
\put(48.213,-116.51){\fontsize{15.987}{1}\usefont{T1}{cmr}{m}{n}\selectfont\color{color_29791}S}
\put(58.89232,-116.51){\fontsize{15.987}{1}\usefont{T1}{cmr}{m}{n}\selectfont\color{color_29791}am}
\put(81.06628,-116.51){\fontsize{15.987}{1}\usefont{T1}{cmr}{m}{n}\selectfont\color{color_29791}ba}
\put(98.811,-116.51){\fontsize{15.987}{1}\usefont{T1}{cmr}{m}{n}\selectfont\color{color_29791} }
\put(84.213,-145.31){\fontsize{11.991}{1}\usefont{T1}{cmr}{m}{n}\selectfont\color{color_29791}•}
\put(98.187,-145.991){\fontsize{11.991}{1}\usefont{T1}{cmr}{m}{n}\selectfont\color{color_29791}A}
\put(106.2689,-145.991){\fontsize{11.991}{1}\usefont{T1}{cmr}{m}{n}\selectfont\color{color_29791}l}
\put(108.871,-145.991){\fontsize{11.991}{1}\usefont{T1}{cmr}{m}{n}\selectfont\color{color_29791}l}
\put(111.473,-145.991){\fontsize{11.991}{1}\usefont{T1}{cmr}{m}{n}\selectfont\color{color_29791}o}
\put(118.2599,-145.991){\fontsize{11.991}{1}\usefont{T1}{cmr}{m}{n}\selectfont\color{color_29791}w}
\put(126.8575,-145.991){\fontsize{11.991}{1}\usefont{T1}{cmr}{m}{n}\selectfont\color{color_29791}s }
\put(136.2345,-145.991){\fontsize{11.991}{1}\usefont{T1}{cmr}{m}{n}\selectfont\color{color_29791}cl}
\put(144.832,-145.991){\fontsize{11.991}{1}\usefont{T1}{cmr}{m}{n}\selectfont\color{color_29791}i}
\put(147.4341,-145.991){\fontsize{11.991}{1}\usefont{T1}{cmr}{m}{n}\selectfont\color{color_29791}e}
\put(154.221,-145.991){\fontsize{11.991}{1}\usefont{T1}{cmr}{m}{n}\selectfont\color{color_29791}n}
\put(160.8999,-145.991){\fontsize{11.991}{1}\usefont{T1}{cmr}{m}{n}\selectfont\color{color_29791}t}
\put(164.1975,-145.991){\fontsize{11.991}{1}\usefont{T1}{cmr}{m}{n}\selectfont\color{color_29791}s }
\put(173.5744,-145.991){\fontsize{11.991}{1}\usefont{T1}{cmr}{m}{n}\selectfont\color{color_29791}t}
\put(176.872,-145.991){\fontsize{11.991}{1}\usefont{T1}{cmr}{m}{n}\selectfont\color{color_29791}o}
\put(183.6589,-145.991){\fontsize{11.991}{1}\usefont{T1}{cmr}{m}{n}\selectfont\color{color_29791} }
\put(186.9564,-145.991){\fontsize{11.991}{1}\usefont{T1}{cmr}{m}{n}\selectfont\color{color_29791}co}
\put(199.7388,-145.991){\fontsize{11.991}{1}\usefont{T1}{cmr}{m}{n}\selectfont\color{color_29791}n}
\put(206.4178,-145.991){\fontsize{11.991}{1}\usefont{T1}{cmr}{m}{n}\selectfont\color{color_29791}n}
\put(213.0968,-145.991){\fontsize{11.991}{1}\usefont{T1}{cmr}{m}{n}\selectfont\color{color_29791}e}
\put(219.7758,-145.991){\fontsize{11.991}{1}\usefont{T1}{cmr}{m}{n}\selectfont\color{color_29791}ct}
\put(229.1528,-145.991){\fontsize{11.991}{1}\usefont{T1}{cmr}{m}{n}\selectfont\color{color_29791} }
\put(232.4503,-145.991){\fontsize{11.991}{1}\usefont{T1}{cmr}{m}{n}\selectfont\color{color_29791}t}
\put(235.8318,-145.991){\fontsize{11.991}{1}\usefont{T1}{cmr}{m}{n}\selectfont\color{color_29791}o}
\put(242.5107,-145.991){\fontsize{11.991}{1}\usefont{T1}{cmr}{m}{n}\selectfont\color{color_29791} }
\put(245.8922,-145.991){\fontsize{11.991}{1}\usefont{T1}{cmr}{m}{n}\selectfont\color{color_29791}U}
\put(254.4897,-145.991){\fontsize{11.991}{1}\usefont{T1}{cmr}{m}{n}\selectfont\color{color_29791}N}
\put(263.0873,-145.991){\fontsize{11.991}{1}\usefont{T1}{cmr}{m}{n}\selectfont\color{color_29791}I}
\put(266.4688,-145.991){\fontsize{11.991}{1}\usefont{T1}{cmr}{m}{n}\selectfont\color{color_29791}X }
\put(277.8362,-145.991){\fontsize{11.991}{1}\usefont{T1}{cmr}{m}{n}\selectfont\color{color_29791}st}
\put(287.2132,-145.991){\fontsize{11.991}{1}\usefont{T1}{cmr}{m}{n}\selectfont\color{color_29791}o}
\put(293.8922,-145.991){\fontsize{11.991}{1}\usefont{T1}{cmr}{m}{n}\selectfont\color{color_29791}ra}
\put(304.5641,-145.991){\fontsize{11.991}{1}\usefont{T1}{cmr}{m}{n}\selectfont\color{color_29791}g}
\put(311.2431,-145.991){\fontsize{11.991}{1}\usefont{T1}{cmr}{m}{n}\selectfont\color{color_29791}e}
\put(317.9221,-145.991){\fontsize{11.991}{1}\usefont{T1}{cmr}{m}{n}\selectfont\color{color_29791} }
\put(321.3036,-145.991){\fontsize{11.991}{1}\usefont{T1}{cmr}{m}{n}\selectfont\color{color_29791}a}
\put(327.9826,-145.991){\fontsize{11.991}{1}\usefont{T1}{cmr}{m}{n}\selectfont\color{color_29791}s }
\put(337.3595,-145.991){\fontsize{11.991}{1}\usefont{T1}{cmr}{m}{n}\selectfont\color{color_29791}i}
\put(339.9616,-145.991){\fontsize{11.991}{1}\usefont{T1}{cmr}{m}{n}\selectfont\color{color_29791}f}
\put(343.343,-145.991){\fontsize{11.991}{1}\usefont{T1}{cmr}{m}{n}\selectfont\color{color_29791} }
\put(346.6406,-145.991){\fontsize{11.991}{1}\usefont{T1}{cmr}{m}{n}\selectfont\color{color_29791}i}
\put(349.2426,-145.991){\fontsize{11.991}{1}\usefont{T1}{cmr}{m}{n}\selectfont\color{color_29791}t}
\put(352.6241,-145.991){\fontsize{11.991}{1}\usefont{T1}{cmr}{m}{n}\selectfont\color{color_29791} }
\put(356.0056,-145.991){\fontsize{11.991}{1}\usefont{T1}{cmr}{m}{n}\selectfont\color{color_29791}w}
\put(364.6031,-145.991){\fontsize{11.991}{1}\usefont{T1}{cmr}{m}{n}\selectfont\color{color_29791}e}
\put(371.2821,-145.991){\fontsize{11.991}{1}\usefont{T1}{cmr}{m}{n}\selectfont\color{color_29791}re}
\put(381.9541,-145.991){\fontsize{11.991}{1}\usefont{T1}{cmr}{m}{n}\selectfont\color{color_29791} }
\put(385.3355,-145.991){\fontsize{11.991}{1}\usefont{T1}{cmr}{m}{n}\selectfont\color{color_29791}a}
\put(392.0145,-145.991){\fontsize{11.991}{1}\usefont{T1}{cmr}{m}{n}\selectfont\color{color_29791} }
\put(395.396,-145.991){\fontsize{11.991}{1}\usefont{T1}{cmr}{m}{n}\selectfont\color{color_29791}w}
\put(403.9935,-145.991){\fontsize{11.991}{1}\usefont{T1}{cmr}{m}{n}\selectfont\color{color_29791}i}
\put(406.5956,-145.991){\fontsize{11.991}{1}\usefont{T1}{cmr}{m}{n}\selectfont\color{color_29791}n}
\put(413.3825,-145.991){\fontsize{11.991}{1}\usefont{T1}{cmr}{m}{n}\selectfont\color{color_29791}d}
\put(420.0615,-145.991){\fontsize{11.991}{1}\usefont{T1}{cmr}{m}{n}\selectfont\color{color_29791}o}
\put(426.7404,-145.991){\fontsize{11.991}{1}\usefont{T1}{cmr}{m}{n}\selectfont\color{color_29791}w}
\put(435.338,-145.991){\fontsize{11.991}{1}\usefont{T1}{cmr}{m}{n}\selectfont\color{color_29791}s }
\put(444.7149,-145.991){\fontsize{11.991}{1}\usefont{T1}{cmr}{m}{n}\selectfont\color{color_29791}n}
\put(451.3939,-145.991){\fontsize{11.991}{1}\usefont{T1}{cmr}{m}{n}\selectfont\color{color_29791}e}
\put(458.0729,-145.991){\fontsize{11.991}{1}\usefont{T1}{cmr}{m}{n}\selectfont\color{color_29791}t}
\put(461.4544,-145.991){\fontsize{11.991}{1}\usefont{T1}{cmr}{m}{n}\selectfont\color{color_29791}w}
\put(470.0519,-145.991){\fontsize{11.991}{1}\usefont{T1}{cmr}{m}{n}\selectfont\color{color_29791}o}
\put(476.8388,-145.991){\fontsize{11.991}{1}\usefont{T1}{cmr}{m}{n}\selectfont\color{color_29791}r}
\put(480.7359,-145.991){\fontsize{11.991}{1}\usefont{T1}{cmr}{m}{n}\selectfont\color{color_29791}k }
\put(490.1128,-145.991){\fontsize{11.991}{1}\usefont{T1}{cmr}{m}{n}\selectfont\color{color_29791}sh}
\put(502.7873,-145.991){\fontsize{11.991}{1}\usefont{T1}{cmr}{m}{n}\selectfont\color{color_29791}a}
\put(509.4663,-145.991){\fontsize{11.991}{1}\usefont{T1}{cmr}{m}{n}\selectfont\color{color_29791}re}
\put(520.1383,-145.991){\fontsize{11.991}{1}\usefont{T1}{cmr}{m}{n}\selectfont\color{color_29791}.}
\put(84.213,-174.11){\fontsize{11.991}{1}\usefont{T1}{cmr}{m}{n}\selectfont\color{color_29791}•}
\put(98.187,-174.791){\fontsize{11.991}{1}\usefont{T1}{cmr}{m}{n}\selectfont\color{color_29791}A}
\put(106.2689,-174.791){\fontsize{11.991}{1}\usefont{T1}{cmr}{m}{n}\selectfont\color{color_29791}l}
\put(108.871,-174.791){\fontsize{11.991}{1}\usefont{T1}{cmr}{m}{n}\selectfont\color{color_29791}l}
\put(111.473,-174.791){\fontsize{11.991}{1}\usefont{T1}{cmr}{m}{n}\selectfont\color{color_29791}o}
\put(118.2599,-174.791){\fontsize{11.991}{1}\usefont{T1}{cmr}{m}{n}\selectfont\color{color_29791}w}
\put(126.8575,-174.791){\fontsize{11.991}{1}\usefont{T1}{cmr}{m}{n}\selectfont\color{color_29791}s }
\put(136.2345,-174.791){\fontsize{11.991}{1}\usefont{T1}{cmr}{m}{n}\selectfont\color{color_29791}a}
\put(142.9134,-174.791){\fontsize{11.991}{1}\usefont{T1}{cmr}{m}{n}\selectfont\color{color_29791}cce}
\put(161.5834,-174.791){\fontsize{11.991}{1}\usefont{T1}{cmr}{m}{n}\selectfont\color{color_29791}ss }
\put(176.8719,-174.791){\fontsize{11.991}{1}\usefont{T1}{cmr}{m}{n}\selectfont\color{color_29791}t}
\put(180.2534,-174.791){\fontsize{11.991}{1}\usefont{T1}{cmr}{m}{n}\selectfont\color{color_29791}o}
\put(186.9324,-174.791){\fontsize{11.991}{1}\usefont{T1}{cmr}{m}{n}\selectfont\color{color_29791} }
\put(190.3139,-174.791){\fontsize{11.991}{1}\usefont{T1}{cmr}{m}{n}\selectfont\color{color_29791}SA}
\put(206.3818,-174.791){\fontsize{11.991}{1}\usefont{T1}{cmr}{m}{n}\selectfont\color{color_29791}N}
\put(214.9794,-174.791){\fontsize{11.991}{1}\usefont{T1}{cmr}{m}{n}\selectfont\color{color_29791} }
\put(218.3608,-174.791){\fontsize{11.991}{1}\usefont{T1}{cmr}{m}{n}\selectfont\color{color_29791}vo}
\put(231.0353,-174.791){\fontsize{11.991}{1}\usefont{T1}{cmr}{m}{n}\selectfont\color{color_29791}l}
\put(233.6373,-174.791){\fontsize{11.991}{1}\usefont{T1}{cmr}{m}{n}\selectfont\color{color_29791}u}
\put(240.4243,-174.791){\fontsize{11.991}{1}\usefont{T1}{cmr}{m}{n}\selectfont\color{color_29791}m}
\put(250.3168,-174.791){\fontsize{11.991}{1}\usefont{T1}{cmr}{m}{n}\selectfont\color{color_29791}e}
\put(256.9958,-174.791){\fontsize{11.991}{1}\usefont{T1}{cmr}{m}{n}\selectfont\color{color_29791}s,}
\put(266.3728,-174.791){\fontsize{11.991}{1}\usefont{T1}{cmr}{m}{n}\selectfont\color{color_29791} }
\put(269.7542,-174.791){\fontsize{11.991}{1}\usefont{T1}{cmr}{m}{n}\selectfont\color{color_29791}/}
\put(273.0518,-174.791){\fontsize{11.991}{1}\usefont{T1}{cmr}{m}{n}\selectfont\color{color_29791}u}
\put(279.7307,-174.791){\fontsize{11.991}{1}\usefont{T1}{cmr}{m}{n}\selectfont\color{color_29791}se}
\put(292.4052,-174.791){\fontsize{11.991}{1}\usefont{T1}{cmr}{m}{n}\selectfont\color{color_29791}rs }
\put(305.7752,-174.791){\fontsize{11.991}{1}\usefont{T1}{cmr}{m}{n}\selectfont\color{color_29791}a}
\put(312.4542,-174.791){\fontsize{11.991}{1}\usefont{T1}{cmr}{m}{n}\selectfont\color{color_29791}n}
\put(319.1332,-174.791){\fontsize{11.991}{1}\usefont{T1}{cmr}{m}{n}\selectfont\color{color_29791}d}
\put(325.8122,-174.791){\fontsize{11.991}{1}\usefont{T1}{cmr}{m}{n}\selectfont\color{color_29791} }
\put(329.1936,-174.791){\fontsize{11.991}{1}\usefont{T1}{cmr}{m}{n}\selectfont\color{color_29791}/}
\put(332.5751,-174.791){\fontsize{11.991}{1}\usefont{T1}{cmr}{m}{n}\selectfont\color{color_29791}d}
\put(339.2541,-174.791){\fontsize{11.991}{1}\usefont{T1}{cmr}{m}{n}\selectfont\color{color_29791}a}
\put(345.9331,-174.791){\fontsize{11.991}{1}\usefont{T1}{cmr}{m}{n}\selectfont\color{color_29791}t}
\put(349.2306,-174.791){\fontsize{11.991}{1}\usefont{T1}{cmr}{m}{n}\selectfont\color{color_29791}a}
\put(356.0175,-174.791){\fontsize{11.991}{1}\usefont{T1}{cmr}{m}{n}\selectfont\color{color_29791}/}
\put(359.315,-174.791){\fontsize{11.991}{1}\usefont{T1}{cmr}{m}{n}\selectfont\color{color_29791}h}
\put(366.1019,-174.791){\fontsize{11.991}{1}\usefont{T1}{cmr}{m}{n}\selectfont\color{color_29791}p}
\put(372.7809,-174.791){\fontsize{11.991}{1}\usefont{T1}{cmr}{m}{n}\selectfont\color{color_29791}cd}
\put(385.4554,-174.791){\fontsize{11.991}{1}\usefont{T1}{cmr}{m}{n}\selectfont\color{color_29791}a}
\put(392.1344,-174.791){\fontsize{11.991}{1}\usefont{T1}{cmr}{m}{n}\selectfont\color{color_29791}t}
\put(395.5158,-174.791){\fontsize{11.991}{1}\usefont{T1}{cmr}{m}{n}\selectfont\color{color_29791}a}
\put(84.213,-202.91){\fontsize{11.991}{1}\usefont{T1}{cmr}{m}{n}\selectfont\color{color_29791}•}
\put(98.187,-203.591){\fontsize{11.991}{1}\usefont{T1}{cmr}{m}{n}\selectfont\color{color_29791}N}
\put(106.8685,-203.591){\fontsize{11.991}{1}\usefont{T1}{cmr}{m}{n}\selectfont\color{color_29791}o}
\put(113.5475,-203.591){\fontsize{11.991}{1}\usefont{T1}{cmr}{m}{n}\selectfont\color{color_29791} }
\put(116.845,-203.591){\fontsize{11.991}{1}\usefont{T1}{cmr}{m}{n}\selectfont\color{color_29791}a}
\put(123.524,-203.591){\fontsize{11.991}{1}\usefont{T1}{cmr}{m}{n}\selectfont\color{color_29791}c}
\put(129.6034,-203.591){\fontsize{11.991}{1}\usefont{T1}{cmr}{m}{n}\selectfont\color{color_29791}ce}
\put(142.2779,-203.591){\fontsize{11.991}{1}\usefont{T1}{cmr}{m}{n}\selectfont\color{color_29791}ss }
\put(157.5664,-203.591){\fontsize{11.991}{1}\usefont{T1}{cmr}{m}{n}\selectfont\color{color_29791}t}
\put(160.9479,-203.591){\fontsize{11.991}{1}\usefont{T1}{cmr}{m}{n}\selectfont\color{color_29791}o}
\put(167.6269,-203.591){\fontsize{11.991}{1}\usefont{T1}{cmr}{m}{n}\selectfont\color{color_29791} }
\put(171.0084,-203.591){\fontsize{11.991}{1}\usefont{T1}{cmr}{m}{n}\selectfont\color{color_29791}/}
\put(174.3059,-203.591){\fontsize{11.991}{1}\usefont{T1}{cmr}{m}{n}\selectfont\color{color_29791}d}
\put(180.9849,-203.591){\fontsize{11.991}{1}\usefont{T1}{cmr}{m}{n}\selectfont\color{color_29791}a}
\put(187.7718,-203.591){\fontsize{11.991}{1}\usefont{T1}{cmr}{m}{n}\selectfont\color{color_29791}t}
\put(191.0693,-203.591){\fontsize{11.991}{1}\usefont{T1}{cmr}{m}{n}\selectfont\color{color_29791}a}
\put(197.8562,-203.591){\fontsize{11.991}{1}\usefont{T1}{cmr}{m}{n}\selectfont\color{color_29791}/}
\put(201.1538,-203.591){\fontsize{11.991}{1}\usefont{T1}{cmr}{m}{n}\selectfont\color{color_29791}h}
\put(207.8327,-203.591){\fontsize{11.991}{1}\usefont{T1}{cmr}{m}{n}\selectfont\color{color_29791}p}
\put(214.6197,-203.591){\fontsize{11.991}{1}\usefont{T1}{cmr}{m}{n}\selectfont\color{color_29791}c}
\put(220.5192,-203.591){\fontsize{11.991}{1}\usefont{T1}{cmr}{m}{n}\selectfont\color{color_29791}f}
\put(223.9007,-203.591){\fontsize{11.991}{1}\usefont{T1}{cmr}{m}{n}\selectfont\color{color_29791}l}
\put(226.5867,-203.591){\fontsize{11.991}{1}\usefont{T1}{cmr}{m}{n}\selectfont\color{color_29791}a}
\put(233.2657,-203.591){\fontsize{11.991}{1}\usefont{T1}{cmr}{m}{n}\selectfont\color{color_29791}sh}
\put(48.213,-237.096){\fontsize{15.987}{1}\usefont{T1}{cmr}{m}{n}\selectfont\color{color_29791}F}
\put(57.90112,-237.096){\fontsize{15.987}{1}\usefont{T1}{cmr}{m}{n}\selectfont\color{color_29791}T}
\put(67.68517,-237.096){\fontsize{15.987}{1}\usefont{T1}{cmr}{m}{n}\selectfont\color{color_29791}P}
\put(84.213,-266.009){\fontsize{11.991}{1}\usefont{T1}{cmr}{m}{n}\selectfont\color{color_29791}•}
\put(98.187,-266.69){\fontsize{11.991}{1}\usefont{T1}{cmr}{m}{n}\selectfont\color{color_29791}A}
\put(106.2689,-266.69){\fontsize{11.991}{1}\usefont{T1}{cmr}{m}{n}\selectfont\color{color_29791}l}
\put(108.871,-266.69){\fontsize{11.991}{1}\usefont{T1}{cmr}{m}{n}\selectfont\color{color_29791}l}
\put(111.473,-266.69){\fontsize{11.991}{1}\usefont{T1}{cmr}{m}{n}\selectfont\color{color_29791}o}
\put(118.2599,-266.69){\fontsize{11.991}{1}\usefont{T1}{cmr}{m}{n}\selectfont\color{color_29791}w}
\put(126.8575,-266.69){\fontsize{11.991}{1}\usefont{T1}{cmr}{m}{n}\selectfont\color{color_29791}s }
\put(136.2345,-266.69){\fontsize{11.991}{1}\usefont{T1}{cmr}{m}{n}\selectfont\color{color_29791}n}
\put(142.9134,-266.69){\fontsize{11.991}{1}\usefont{T1}{cmr}{m}{n}\selectfont\color{color_29791}o}
\put(149.5924,-266.69){\fontsize{11.991}{1}\usefont{T1}{cmr}{m}{n}\selectfont\color{color_29791}n}
\put(156.2714,-266.69){\fontsize{11.991}{1}\usefont{T1}{cmr}{m}{n}\selectfont\color{color_29791}-BA}
\put(176.3324,-266.69){\fontsize{11.991}{1}\usefont{T1}{cmr}{m}{n}\selectfont\color{color_29791}S }
\put(187.6159,-266.69){\fontsize{11.991}{1}\usefont{T1}{cmr}{m}{n}\selectfont\color{color_29791}u}
\put(194.4028,-266.69){\fontsize{11.991}{1}\usefont{T1}{cmr}{m}{n}\selectfont\color{color_29791}se}
\put(207.0773,-266.69){\fontsize{11.991}{1}\usefont{T1}{cmr}{m}{n}\selectfont\color{color_29791}rs}
\put(216.9698,-266.69){\fontsize{11.991}{1}\usefont{T1}{cmr}{m}{n}\selectfont\color{color_29791} }
\put(220.3513,-266.69){\fontsize{11.991}{1}\usefont{T1}{cmr}{m}{n}\selectfont\color{color_29791}t}
\put(223.7328,-266.69){\fontsize{11.991}{1}\usefont{T1}{cmr}{m}{n}\selectfont\color{color_29791}o}
\put(230.4118,-266.69){\fontsize{11.991}{1}\usefont{T1}{cmr}{m}{n}\selectfont\color{color_29791} }
\put(233.7932,-266.69){\fontsize{11.991}{1}\usefont{T1}{cmr}{m}{n}\selectfont\color{color_29791}r}
\put(237.6903,-266.69){\fontsize{11.991}{1}\usefont{T1}{cmr}{m}{n}\selectfont\color{color_29791}e}
\put(244.3693,-266.69){\fontsize{11.991}{1}\usefont{T1}{cmr}{m}{n}\selectfont\color{color_29791}t}
\put(247.7508,-266.69){\fontsize{11.991}{1}\usefont{T1}{cmr}{m}{n}\selectfont\color{color_29791}ri}
\put(254.3459,-266.69){\fontsize{11.991}{1}\usefont{T1}{cmr}{m}{n}\selectfont\color{color_29791}e}
\put(261.0248,-266.69){\fontsize{11.991}{1}\usefont{T1}{cmr}{m}{n}\selectfont\color{color_29791}ve}
\put(273.6993,-266.69){\fontsize{11.991}{1}\usefont{T1}{cmr}{m}{n}\selectfont\color{color_29791} }
\put(277.0808,-266.69){\fontsize{11.991}{1}\usefont{T1}{cmr}{m}{n}\selectfont\color{color_29791}f}
\put(280.4623,-266.69){\fontsize{11.991}{1}\usefont{T1}{cmr}{m}{n}\selectfont\color{color_29791}i}
\put(283.0643,-266.69){\fontsize{11.991}{1}\usefont{T1}{cmr}{m}{n}\selectfont\color{color_29791}l}
\put(285.6664,-266.69){\fontsize{11.991}{1}\usefont{T1}{cmr}{m}{n}\selectfont\color{color_29791}e}
\put(292.3454,-266.69){\fontsize{11.991}{1}\usefont{T1}{cmr}{m}{n}\selectfont\color{color_29791}s}
\put(298.4248,-266.69){\fontsize{11.991}{1}\usefont{T1}{cmr}{m}{n}\selectfont\color{color_29791} }
\put(301.7223,-266.69){\fontsize{11.991}{1}\usefont{T1}{cmr}{m}{n}\selectfont\color{color_29791}f}
\put(305.1038,-266.69){\fontsize{11.991}{1}\usefont{T1}{cmr}{m}{n}\selectfont\color{color_29791}r}
\put(309.0009,-266.69){\fontsize{11.991}{1}\usefont{T1}{cmr}{m}{n}\selectfont\color{color_29791}o}
\put(315.7878,-266.69){\fontsize{11.991}{1}\usefont{T1}{cmr}{m}{n}\selectfont\color{color_29791}m}
\put(325.6803,-266.69){\fontsize{11.991}{1}\usefont{T1}{cmr}{m}{n}\selectfont\color{color_29791} }
\put(329.0618,-266.69){\fontsize{11.991}{1}\usefont{T1}{cmr}{m}{n}\selectfont\color{color_29791}t}
\put(332.3593,-266.69){\fontsize{11.991}{1}\usefont{T1}{cmr}{m}{n}\selectfont\color{color_29791}h}
\put(339.0383,-266.69){\fontsize{11.991}{1}\usefont{T1}{cmr}{m}{n}\selectfont\color{color_29791}e}
\put(345.8252,-266.69){\fontsize{11.991}{1}\usefont{T1}{cmr}{m}{n}\selectfont\color{color_29791} }
\put(349.1227,-266.69){\fontsize{11.991}{1}\usefont{T1}{cmr}{m}{n}\selectfont\color{color_29791}F}
\put(356.4132,-266.69){\fontsize{11.991}{1}\usefont{T1}{cmr}{m}{n}\selectfont\color{color_29791}T}
\put(363.7997,-266.69){\fontsize{11.991}{1}\usefont{T1}{cmr}{m}{n}\selectfont\color{color_29791}P}
\put(371.5938,-266.69){\fontsize{11.991}{1}\usefont{T1}{cmr}{m}{n}\selectfont\color{color_29791} }
\put(374.8914,-266.69){\fontsize{11.991}{1}\usefont{T1}{cmr}{m}{n}\selectfont\color{color_29791}a}
\put(381.5703,-266.69){\fontsize{11.991}{1}\usefont{T1}{cmr}{m}{n}\selectfont\color{color_29791}re}
\put(392.2423,-266.69){\fontsize{11.991}{1}\usefont{T1}{cmr}{m}{n}\selectfont\color{color_29791}a}
\put(399.0292,-266.69){\fontsize{11.991}{1}\usefont{T1}{cmr}{m}{n}\selectfont\color{color_29791} }
\put(402.997,-266.69){\fontsize{11.991}{1}\usefont{T1}{cmr}{m}{n}\selectfont\color{color_29791}f}
\put(406.3785,-266.69){\fontsize{11.991}{1}\usefont{T1}{cmr}{m}{n}\selectfont\color{color_29791}t}
\put(409.676,-266.69){\fontsize{11.991}{1}\usefont{T1}{cmr}{m}{n}\selectfont\color{color_29791}p}
\put(416.4629,-266.69){\fontsize{11.991}{1}\usefont{T1}{cmr}{m}{n}\selectfont\color{color_29791}:}
\put(419.7604,-266.69){\fontsize{11.991}{1}\usefont{T1}{cmr}{m}{n}\selectfont\color{color_29791}/}
\put(423.1419,-266.69){\fontsize{11.991}{1}\usefont{T1}{cmr}{m}{n}\selectfont\color{color_29791}/}
\put(426.5233,-266.69){\fontsize{11.991}{1}\usefont{T1}{cmr}{m}{n}\selectfont\color{color_29791}f}
\put(429.8209,-266.69){\fontsize{11.991}{1}\usefont{T1}{cmr}{m}{n}\selectfont\color{color_29791}t}
\put(433.2023,-266.69){\fontsize{11.991}{1}\usefont{T1}{cmr}{m}{n}\selectfont\color{color_29791}p}
\put(439.8813,-266.69){\fontsize{11.991}{1}\usefont{T1}{cmr}{m}{n}\selectfont\color{color_29791}.}
\put(443.1788,-266.69){\fontsize{11.991}{1}\usefont{T1}{cmr}{m}{n}\selectfont\color{color_29791}b}
\put(449.9657,-266.69){\fontsize{11.991}{1}\usefont{T1}{cmr}{m}{n}\selectfont\color{color_29791}a}
\put(456.6447,-266.69){\fontsize{11.991}{1}\usefont{T1}{cmr}{m}{n}\selectfont\color{color_29791}s.}
\put(466.0217,-266.69){\fontsize{11.991}{1}\usefont{T1}{cmr}{m}{n}\selectfont\color{color_29791}a}
\put(472.7007,-266.69){\fontsize{11.991}{1}\usefont{T1}{cmr}{m}{n}\selectfont\color{color_29791}c.}
\put(481.9937,-266.69){\fontsize{11.991}{1}\usefont{T1}{cmr}{m}{n}\selectfont\color{color_29791}u}
\put(488.6727,-266.69){\fontsize{11.991}{1}\usefont{T1}{cmr}{m}{n}\selectfont\color{color_29791}k}
\put(494.7521,-266.69){\fontsize{11.991}{1}\usefont{T1}{cmr}{m}{n}\selectfont\color{color_29791}/}
\end{picture}
\begin{tikzpicture}[overlay]
\path(0pt,0pt);
\draw[color_29791,line width=0.595pt]
(402.997pt, -267.682pt) -- (498.298pt, -267.682pt)
;
\end{tikzpicture}
\begin{picture}(-5,0)(2.5,0)
\put(498.298,-266.69){\fontsize{11.991}{1}\usefont{T1}{cmr}{m}{n}\selectfont\color{color_29791} }
\put(84.213,-294.809){\fontsize{11.991}{1}\usefont{T1}{cmr}{m}{n}\selectfont\color{color_29791}•}
\put(98.187,-295.49){\fontsize{11.991}{1}\usefont{T1}{cmr}{m}{n}\selectfont\color{color_29791}U}
\put(106.8685,-295.49){\fontsize{11.991}{1}\usefont{T1}{cmr}{m}{n}\selectfont\color{color_29791}s}
\put(112.7681,-295.49){\fontsize{11.991}{1}\usefont{T1}{cmr}{m}{n}\selectfont\color{color_29791}e}
\put(119.555,-295.49){\fontsize{11.991}{1}\usefont{T1}{cmr}{m}{n}\selectfont\color{color_29791}r}
\put(123.452,-295.49){\fontsize{11.991}{1}\usefont{T1}{cmr}{m}{n}\selectfont\color{color_29791}s }
\put(132.829,-295.49){\fontsize{11.991}{1}\usefont{T1}{cmr}{m}{n}\selectfont\color{color_29791}w}
\put(141.4266,-295.49){\fontsize{11.991}{1}\usefont{T1}{cmr}{m}{n}\selectfont\color{color_29791}i}
\put(144.1125,-295.49){\fontsize{11.991}{1}\usefont{T1}{cmr}{m}{n}\selectfont\color{color_29791}t}
\put(147.4101,-295.49){\fontsize{11.991}{1}\usefont{T1}{cmr}{m}{n}\selectfont\color{color_29791}h}
\put(154.197,-295.49){\fontsize{11.991}{1}\usefont{T1}{cmr}{m}{n}\selectfont\color{color_29791}i}
\put(156.799,-295.49){\fontsize{11.991}{1}\usefont{T1}{cmr}{m}{n}\selectfont\color{color_29791}n}
\put(163.478,-295.49){\fontsize{11.991}{1}\usefont{T1}{cmr}{m}{n}\selectfont\color{color_29791} }
\put(166.8595,-295.49){\fontsize{11.991}{1}\usefont{T1}{cmr}{m}{n}\selectfont\color{color_29791}BA}
\put(182.9274,-295.49){\fontsize{11.991}{1}\usefont{T1}{cmr}{m}{n}\selectfont\color{color_29791}S }
\put(194.2949,-295.49){\fontsize{11.991}{1}\usefont{T1}{cmr}{m}{n}\selectfont\color{color_29791}ca}
\put(206.9694,-295.49){\fontsize{11.991}{1}\usefont{T1}{cmr}{m}{n}\selectfont\color{color_29791}n}
\put(213.6484,-295.49){\fontsize{11.991}{1}\usefont{T1}{cmr}{m}{n}\selectfont\color{color_29791} }
\put(216.9459,-295.49){\fontsize{11.991}{1}\usefont{T1}{cmr}{m}{n}\selectfont\color{color_29791}g}
\put(223.7328,-295.49){\fontsize{11.991}{1}\usefont{T1}{cmr}{m}{n}\selectfont\color{color_29791}a}
\put(230.4118,-295.49){\fontsize{11.991}{1}\usefont{T1}{cmr}{m}{n}\selectfont\color{color_29791}i}
\put(233.0138,-295.49){\fontsize{11.991}{1}\usefont{T1}{cmr}{m}{n}\selectfont\color{color_29791}n}
\put(239.8008,-295.49){\fontsize{11.991}{1}\usefont{T1}{cmr}{m}{n}\selectfont\color{color_29791} }
\put(243.0983,-295.49){\fontsize{11.991}{1}\usefont{T1}{cmr}{m}{n}\selectfont\color{color_29791}a}
\put(249.7773,-295.49){\fontsize{11.991}{1}\usefont{T1}{cmr}{m}{n}\selectfont\color{color_29791}cce}
\put(268.5551,-295.49){\fontsize{11.991}{1}\usefont{T1}{cmr}{m}{n}\selectfont\color{color_29791}s}
\put(274.4547,-295.49){\fontsize{11.991}{1}\usefont{T1}{cmr}{m}{n}\selectfont\color{color_29791}s}
\put(280.5341,-295.49){\fontsize{11.991}{1}\usefont{T1}{cmr}{m}{n}\selectfont\color{color_29791} }
\put(283.8317,-295.49){\fontsize{11.991}{1}\usefont{T1}{cmr}{m}{n}\selectfont\color{color_29791}t}
\put(287.2131,-295.49){\fontsize{11.991}{1}\usefont{T1}{cmr}{m}{n}\selectfont\color{color_29791}o}
\put(293.8921,-295.49){\fontsize{11.991}{1}\usefont{T1}{cmr}{m}{n}\selectfont\color{color_29791} }
\put(297.2736,-295.49){\fontsize{11.991}{1}\usefont{T1}{cmr}{m}{n}\selectfont\color{color_29791}t}
\put(300.5711,-295.49){\fontsize{11.991}{1}\usefont{T1}{cmr}{m}{n}\selectfont\color{color_29791}h}
\put(307.2501,-295.49){\fontsize{11.991}{1}\usefont{T1}{cmr}{m}{n}\selectfont\color{color_29791}i}
\put(309.9361,-295.49){\fontsize{11.991}{1}\usefont{T1}{cmr}{m}{n}\selectfont\color{color_29791}s }
\put(319.2291,-295.49){\fontsize{11.991}{1}\usefont{T1}{cmr}{m}{n}\selectfont\color{color_29791}a}
\put(326.016,-295.49){\fontsize{11.991}{1}\usefont{T1}{cmr}{m}{n}\selectfont\color{color_29791}r}
\put(329.9131,-295.49){\fontsize{11.991}{1}\usefont{T1}{cmr}{m}{n}\selectfont\color{color_29791}e}
\put(336.592,-295.49){\fontsize{11.991}{1}\usefont{T1}{cmr}{m}{n}\selectfont\color{color_29791}a}
\put(343.894,-295.49){\fontsize{11.991}{1}\usefont{T1}{cmr}{m}{n}\selectfont\color{color_29791} }
\put(347.211,-295.49){\fontsize{11.991}{1}\usefont{T1}{cmr}{m}{n}\selectfont\color{color_29791}a}
\put(353.89,-295.49){\fontsize{11.991}{1}\usefont{T1}{cmr}{m}{n}\selectfont\color{color_29791}n}
\put(360.569,-295.49){\fontsize{11.991}{1}\usefont{T1}{cmr}{m}{n}\selectfont\color{color_29791}d}
\put(367.3559,-295.49){\fontsize{11.991}{1}\usefont{T1}{cmr}{m}{n}\selectfont\color{color_29791} }
\put(370.6534,-295.49){\fontsize{11.991}{1}\usefont{T1}{cmr}{m}{n}\selectfont\color{color_29791}d}
\put(377.4403,-295.49){\fontsize{11.991}{1}\usefont{T1}{cmr}{m}{n}\selectfont\color{color_29791}e}
\put(384.1193,-295.49){\fontsize{11.991}{1}\usefont{T1}{cmr}{m}{n}\selectfont\color{color_29791}p}
\put(390.7982,-295.49){\fontsize{11.991}{1}\usefont{T1}{cmr}{m}{n}\selectfont\color{color_29791}o}
\put(397.4772,-295.49){\fontsize{11.991}{1}\usefont{T1}{cmr}{m}{n}\selectfont\color{color_29791}si}
\put(406.1587,-295.49){\fontsize{11.991}{1}\usefont{T1}{cmr}{m}{n}\selectfont\color{color_29791}t}
\put(409.4562,-295.49){\fontsize{11.991}{1}\usefont{T1}{cmr}{m}{n}\selectfont\color{color_29791} }
\put(412.8377,-295.49){\fontsize{11.991}{1}\usefont{T1}{cmr}{m}{n}\selectfont\color{color_29791}f}
\put(416.1352,-295.49){\fontsize{11.991}{1}\usefont{T1}{cmr}{m}{n}\selectfont\color{color_29791}i}
\put(418.8212,-295.49){\fontsize{11.991}{1}\usefont{T1}{cmr}{m}{n}\selectfont\color{color_29791}l}
\put(421.4232,-295.49){\fontsize{11.991}{1}\usefont{T1}{cmr}{m}{n}\selectfont\color{color_29791}e}
\put(428.1022,-295.49){\fontsize{11.991}{1}\usefont{T1}{cmr}{m}{n}\selectfont\color{color_29791}s}
\put(84.213,-323.609){\fontsize{11.991}{1}\usefont{T1}{cmr}{m}{n}\selectfont\color{color_29791}•}
\put(98.187,-324.205){\fontsize{11.991}{1}\usefont{T1}{cmr}{m}{n}\selectfont\color{color_29791}P}
\put(106.2689,-324.205){\fontsize{11.991}{1}\usefont{T1}{cmr}{m}{n}\selectfont\color{color_29791}l}
\put(108.871,-324.205){\fontsize{11.991}{1}\usefont{T1}{cmr}{m}{n}\selectfont\color{color_29791}e}
\put(115.55,-324.205){\fontsize{11.991}{1}\usefont{T1}{cmr}{m}{n}\selectfont\color{color_29791}a}
\put(122.2289,-324.205){\fontsize{11.991}{1}\usefont{T1}{cmr}{m}{n}\selectfont\color{color_29791}se}
\put(135.0114,-324.205){\fontsize{11.991}{1}\usefont{T1}{cmr}{m}{n}\selectfont\color{color_29791} }
\put(138.3089,-324.205){\fontsize{11.991}{1}\usefont{T1}{cmr}{m}{n}\selectfont\color{color_29791}co}
\put(150.9834,-324.205){\fontsize{11.991}{1}\usefont{T1}{cmr}{m}{n}\selectfont\color{color_29791}n}
\put(157.7703,-324.205){\fontsize{11.991}{1}\usefont{T1}{cmr}{m}{n}\selectfont\color{color_29791}t}
\put(161.0678,-324.205){\fontsize{11.991}{1}\usefont{T1}{cmr}{m}{n}\selectfont\color{color_29791}a}
\put(167.7468,-324.205){\fontsize{11.991}{1}\usefont{T1}{cmr}{m}{n}\selectfont\color{color_29791}ct}
\put(177.1238,-324.205){\fontsize{11.991}{1}\usefont{T1}{cmr}{m}{n}\selectfont\color{color_29791} }
\put(180.4213,-324.205){\fontsize{11.991}{1}\usefont{T1}{cmr}{m}{n}\selectfont\color{color_29791}t}
\put(183.8028,-324.205){\fontsize{11.991}{1}\usefont{T1}{cmr}{m}{n}\selectfont\color{color_29791}h}
\put(190.4818,-324.205){\fontsize{11.991}{1}\usefont{T1}{cmr}{m}{n}\selectfont\color{color_29791}e}
\put(197.2687,-324.205){\fontsize{11.991}{1}\usefont{T1}{cmr}{m}{n}\selectfont\color{color_29791} }
\put(200.5662,-324.205){\fontsize{11.991}{1}\usefont{T1}{cmr}{m}{n}\selectfont\color{color_29791}I}
\put(203.9477,-324.205){\fontsize{11.991}{1}\usefont{T1}{cmr}{m}{n}\selectfont\color{color_29791}T}
\put(211.0463,-324.205){\fontsize{11.991}{1}\usefont{T1}{cmr}{m}{n}\selectfont\color{color_29791} }
\put(214.3439,-324.205){\fontsize{11.991}{1}\usefont{T1}{cmr}{m}{n}\selectfont\color{color_29791}S}
\put(222.4258,-324.205){\fontsize{11.991}{1}\usefont{T1}{cmr}{m}{n}\selectfont\color{color_29791}e}
\put(229.1048,-324.205){\fontsize{11.991}{1}\usefont{T1}{cmr}{m}{n}\selectfont\color{color_29791}r}
\put(233.0019,-324.205){\fontsize{11.991}{1}\usefont{T1}{cmr}{m}{n}\selectfont\color{color_29791}vi}
\put(241.6833,-324.205){\fontsize{11.991}{1}\usefont{T1}{cmr}{m}{n}\selectfont\color{color_29791}ce}
\put(254.3578,-324.205){\fontsize{11.991}{1}\usefont{T1}{cmr}{m}{n}\selectfont\color{color_29791}D}
\put(262.9554,-324.205){\fontsize{11.991}{1}\usefont{T1}{cmr}{m}{n}\selectfont\color{color_29791}e}
\put(269.7423,-324.205){\fontsize{11.991}{1}\usefont{T1}{cmr}{m}{n}\selectfont\color{color_29791}sk}
\put(281.6374,-324.205){\fontsize{11.991}{1}\usefont{T1}{cmr}{m}{n}\selectfont\color{color_29791} }
\put(285.0188,-324.205){\fontsize{11.991}{1}\usefont{T1}{cmr}{m}{n}\selectfont\color{color_29791}t}
\put(288.4003,-324.205){\fontsize{11.991}{1}\usefont{T1}{cmr}{m}{n}\selectfont\color{color_29791}o}
\put(295.0793,-324.205){\fontsize{11.991}{1}\usefont{T1}{cmr}{m}{n}\selectfont\color{color_29791} }
\put(298.4608,-324.205){\fontsize{11.991}{1}\usefont{T1}{cmr}{m}{n}\selectfont\color{color_29791}h}
\put(305.1397,-324.205){\fontsize{11.991}{1}\usefont{T1}{cmr}{m}{n}\selectfont\color{color_29791}a}
\put(311.8187,-324.205){\fontsize{11.991}{1}\usefont{T1}{cmr}{m}{n}\selectfont\color{color_29791}ve}
\put(324.4932,-324.205){\fontsize{11.991}{1}\usefont{T1}{cmr}{m}{n}\selectfont\color{color_29791} }
\put(327.8747,-324.205){\fontsize{11.991}{1}\usefont{T1}{cmr}{m}{n}\selectfont\color{color_29791}a}
\put(334.5537,-324.205){\fontsize{11.991}{1}\usefont{T1}{cmr}{m}{n}\selectfont\color{color_29791} }
\put(337.9352,-324.205){\fontsize{11.991}{1}\usefont{T1}{cmr}{m}{n}\selectfont\color{color_29791}d}
\put(344.6141,-324.205){\fontsize{11.991}{1}\usefont{T1}{cmr}{m}{n}\selectfont\color{color_29791}i}
\put(347.2162,-324.205){\fontsize{11.991}{1}\usefont{T1}{cmr}{m}{n}\selectfont\color{color_29791}re}
\put(357.8882,-324.205){\fontsize{11.991}{1}\usefont{T1}{cmr}{m}{n}\selectfont\color{color_29791}ct}
\put(367.1812,-324.205){\fontsize{11.991}{1}\usefont{T1}{cmr}{m}{n}\selectfont\color{color_29791}o}
\put(373.9681,-324.205){\fontsize{11.991}{1}\usefont{T1}{cmr}{m}{n}\selectfont\color{color_29791}r}
\put(377.8652,-324.205){\fontsize{11.991}{1}\usefont{T1}{cmr}{m}{n}\selectfont\color{color_29791}y }
\put(387.2421,-324.205){\fontsize{11.991}{1}\usefont{T1}{cmr}{m}{n}\selectfont\color{color_29791}se}
\put(399.9166,-324.205){\fontsize{11.991}{1}\usefont{T1}{cmr}{m}{n}\selectfont\color{color_29791}t}
\put(403.2981,-324.205){\fontsize{11.991}{1}\usefont{T1}{cmr}{m}{n}\selectfont\color{color_29791}u}
\put(409.9771,-324.205){\fontsize{11.991}{1}\usefont{T1}{cmr}{m}{n}\selectfont\color{color_29791}p}
\put(416.6561,-324.205){\fontsize{11.991}{1}\usefont{T1}{cmr}{m}{n}\selectfont\color{color_29791} }
\put(420.0375,-324.205){\fontsize{11.991}{1}\usefont{T1}{cmr}{m}{n}\selectfont\color{color_29791}i}
\put(422.6396,-324.205){\fontsize{11.991}{1}\usefont{T1}{cmr}{m}{n}\selectfont\color{color_29791}e}
\put(429.3186,-324.205){\fontsize{11.991}{1}\usefont{T1}{cmr}{m}{n}\selectfont\color{color_29791}.}
\put(433.498,-324.205){\fontsize{11.991}{1}\usefont{T1}{cmr}{m}{n}\selectfont\color{color_29791} }
\put(436.814,-324.205){\fontsize{11.991}{1}\usefont{T1}{cmr}{m}{n}\selectfont\color{color_29791}/}
\put(440.1955,-324.205){\fontsize{11.991}{1}\usefont{T1}{cmr}{m}{n}\selectfont\color{color_29791}d}
\put(446.8745,-324.205){\fontsize{11.991}{1}\usefont{T1}{cmr}{m}{n}\selectfont\color{color_29791}a}
\put(453.6613,-324.205){\fontsize{11.991}{1}\usefont{T1}{cmr}{m}{n}\selectfont\color{color_29791}t}
\put(456.9589,-324.205){\fontsize{11.991}{1}\usefont{T1}{cmr}{m}{n}\selectfont\color{color_29791}a}
\put(463.6378,-324.205){\fontsize{11.991}{1}\usefont{T1}{cmr}{m}{n}\selectfont\color{color_29791}/}
\put(467.0193,-324.205){\fontsize{11.991}{1}\usefont{T1}{cmr}{m}{n}\selectfont\color{color_29791}f}
\put(470.3168,-324.205){\fontsize{11.991}{1}\usefont{T1}{cmr}{m}{n}\selectfont\color{color_29791}t}
\put(473.6983,-324.205){\fontsize{11.991}{1}\usefont{T1}{cmr}{m}{n}\selectfont\color{color_29791}p}
\put(480.3773,-324.205){\fontsize{11.991}{1}\usefont{T1}{cmr}{m}{n}\selectfont\color{color_29791}/}
\put(483.7588,-324.205){\fontsize{11.991}{1}\usefont{T1}{cmr}{m}{n}\selectfont\color{color_29791}u}
\put(490.4377,-324.205){\fontsize{11.991}{1}\usefont{T1}{cmr}{m}{n}\selectfont\color{color_29791}se}
\put(503.1122,-324.205){\fontsize{11.991}{1}\usefont{T1}{cmr}{m}{n}\selectfont\color{color_29791}rn}
\put(513.7842,-324.205){\fontsize{11.991}{1}\usefont{T1}{cmr}{m}{n}\selectfont\color{color_29791}a}
\put(520.4633,-324.205){\fontsize{11.991}{1}\usefont{T1}{cmr}{m}{n}\selectfont\color{color_29791}me}
\put(537.387,-324.205){\fontsize{11.991}{1}\usefont{T1}{cmr}{m}{n}\selectfont\color{color_29791} }
\put(48.213,-357.795){\fontsize{15.987}{1}\usefont{T1}{cmr}{m}{n}\selectfont\color{color_29791}W}
\put(63.00098,-357.795){\fontsize{15.987}{1}\usefont{T1}{cmr}{m}{n}\selectfont\color{color_29791}r}
\put(68.29267,-357.795){\fontsize{15.987}{1}\usefont{T1}{cmr}{m}{n}\selectfont\color{color_29791}i}
\put(71.87376,-357.795){\fontsize{15.987}{1}\usefont{T1}{cmr}{m}{n}\selectfont\color{color_29791}t}
\put(76.27018,-357.795){\fontsize{15.987}{1}\usefont{T1}{cmr}{m}{n}\selectfont\color{color_29791}ea}
\put(93.9518,-357.795){\fontsize{15.987}{1}\usefont{T1}{cmr}{m}{n}\selectfont\color{color_29791}bl}
\put(106.4217,-357.795){\fontsize{15.987}{1}\usefont{T1}{cmr}{m}{n}\selectfont\color{color_29791}e }
\put(119.7069,-357.795){\fontsize{15.987}{1}\usefont{T1}{cmr}{m}{n}\selectfont\color{color_29791}F}
\put(129.4909,-357.795){\fontsize{15.987}{1}\usefont{T1}{cmr}{m}{n}\selectfont\color{color_29791}T}
\put(139.179,-357.795){\fontsize{15.987}{1}\usefont{T1}{cmr}{m}{n}\selectfont\color{color_29791}P}
\put(149.4747,-357.795){\fontsize{15.987}{1}\usefont{T1}{cmr}{m}{n}\selectfont\color{color_29791} }
\put(153.0717,-357.795){\fontsize{15.987}{1}\usefont{T1}{cmr}{m}{n}\selectfont\color{color_29791}A}
\put(163.6551,-357.795){\fontsize{15.987}{1}\usefont{T1}{cmr}{m}{n}\selectfont\color{color_29791}r}
\put(168.9468,-357.795){\fontsize{15.987}{1}\usefont{T1}{cmr}{m}{n}\selectfont\color{color_29791}e}
\put(177.7397,-357.795){\fontsize{15.987}{1}\usefont{T1}{cmr}{m}{n}\selectfont\color{color_29791}a}
\put(84.213,-386.709){\fontsize{11.991}{1}\usefont{T1}{cmr}{m}{n}\selectfont\color{color_29791}•}
\put(98.187,-387.389){\fontsize{11.991}{1}\usefont{T1}{cmr}{m}{n}\selectfont\color{color_29791}P}
\put(106.2689,-387.389){\fontsize{11.991}{1}\usefont{T1}{cmr}{m}{n}\selectfont\color{color_29791}o}
\put(112.9479,-387.389){\fontsize{11.991}{1}\usefont{T1}{cmr}{m}{n}\selectfont\color{color_29791}ssi}
\put(127.541,-387.389){\fontsize{11.991}{1}\usefont{T1}{cmr}{m}{n}\selectfont\color{color_29791}b}
\put(134.2199,-387.389){\fontsize{11.991}{1}\usefont{T1}{cmr}{m}{n}\selectfont\color{color_29791}l}
\put(136.9059,-387.389){\fontsize{11.991}{1}\usefont{T1}{cmr}{m}{n}\selectfont\color{color_29791}e}
\put(143.5849,-387.389){\fontsize{11.991}{1}\usefont{T1}{cmr}{m}{n}\selectfont\color{color_29791} }
\put(146.9664,-387.389){\fontsize{11.991}{1}\usefont{T1}{cmr}{m}{n}\selectfont\color{color_29791}f}
\put(150.2639,-387.389){\fontsize{11.991}{1}\usefont{T1}{cmr}{m}{n}\selectfont\color{color_29791}o}
\put(156.9429,-387.389){\fontsize{11.991}{1}\usefont{T1}{cmr}{m}{n}\selectfont\color{color_29791}r }
\put(164.3174,-387.389){\fontsize{11.991}{1}\usefont{T1}{cmr}{m}{n}\selectfont\color{color_29791}n}
\put(170.9963,-387.389){\fontsize{11.991}{1}\usefont{T1}{cmr}{m}{n}\selectfont\color{color_29791}o}
\put(177.6753,-387.389){\fontsize{11.991}{1}\usefont{T1}{cmr}{m}{n}\selectfont\color{color_29791}n}
\put(184.3543,-387.389){\fontsize{11.991}{1}\usefont{T1}{cmr}{m}{n}\selectfont\color{color_29791}-BA}
\put(204.4153,-387.389){\fontsize{11.991}{1}\usefont{T1}{cmr}{m}{n}\selectfont\color{color_29791}S }
\put(215.7827,-387.389){\fontsize{11.991}{1}\usefont{T1}{cmr}{m}{n}\selectfont\color{color_29791}u}
\put(222.4617,-387.389){\fontsize{11.991}{1}\usefont{T1}{cmr}{m}{n}\selectfont\color{color_29791}se}
\put(235.1362,-387.389){\fontsize{11.991}{1}\usefont{T1}{cmr}{m}{n}\selectfont\color{color_29791}rs}
\put(245.0287,-387.389){\fontsize{11.991}{1}\usefont{T1}{cmr}{m}{n}\selectfont\color{color_29791} }
\put(248.4102,-387.389){\fontsize{11.991}{1}\usefont{T1}{cmr}{m}{n}\selectfont\color{color_29791}t}
\put(251.7917,-387.389){\fontsize{11.991}{1}\usefont{T1}{cmr}{m}{n}\selectfont\color{color_29791}o}
\put(258.912,-387.389){\fontsize{11.991}{1}\usefont{T1}{cmr}{m}{n}\selectfont\color{color_29791} }
\put(262.313,-387.389){\fontsize{11.991}{1}\usefont{T1}{cmr}{m}{n}\selectfont\color{color_29791}u}
\put(268.992,-387.389){\fontsize{11.991}{1}\usefont{T1}{cmr}{m}{n}\selectfont\color{color_29791}p}
\put(275.671,-387.389){\fontsize{11.991}{1}\usefont{T1}{cmr}{m}{n}\selectfont\color{color_29791}l}
\put(278.273,-387.389){\fontsize{11.991}{1}\usefont{T1}{cmr}{m}{n}\selectfont\color{color_29791}o}
\put(284.952,-387.389){\fontsize{11.991}{1}\usefont{T1}{cmr}{m}{n}\selectfont\color{color_29791}a}
\put(291.7389,-387.389){\fontsize{11.991}{1}\usefont{T1}{cmr}{m}{n}\selectfont\color{color_29791}d}
\put(298.4179,-387.389){\fontsize{11.991}{1}\usefont{T1}{cmr}{m}{n}\selectfont\color{color_29791} }
\put(301.7993,-387.389){\fontsize{11.991}{1}\usefont{T1}{cmr}{m}{n}\selectfont\color{color_29791}f}
\put(305.0969,-387.389){\fontsize{11.991}{1}\usefont{T1}{cmr}{m}{n}\selectfont\color{color_29791}i}
\put(307.7828,-387.389){\fontsize{11.991}{1}\usefont{T1}{cmr}{m}{n}\selectfont\color{color_29791}l}
\put(310.3849,-387.389){\fontsize{11.991}{1}\usefont{T1}{cmr}{m}{n}\selectfont\color{color_29791}e}
\put(317.0639,-387.389){\fontsize{11.991}{1}\usefont{T1}{cmr}{m}{n}\selectfont\color{color_29791}s }
\put(326.4408,-387.389){\fontsize{11.991}{1}\usefont{T1}{cmr}{m}{n}\selectfont\color{color_29791}a}
\put(333.1198,-387.389){\fontsize{11.991}{1}\usefont{T1}{cmr}{m}{n}\selectfont\color{color_29791}s }
\put(342.4128,-387.389){\fontsize{11.991}{1}\usefont{T1}{cmr}{m}{n}\selectfont\color{color_29791}w}
\put(351.0943,-387.389){\fontsize{11.991}{1}\usefont{T1}{cmr}{m}{n}\selectfont\color{color_29791}e}
\put(357.7733,-387.389){\fontsize{11.991}{1}\usefont{T1}{cmr}{m}{n}\selectfont\color{color_29791}l}
\put(360.3754,-387.389){\fontsize{11.991}{1}\usefont{T1}{cmr}{m}{n}\selectfont\color{color_29791}l}
\put(363.0613,-387.389){\fontsize{11.991}{1}\usefont{T1}{cmr}{m}{n}\selectfont\color{color_29791},}
\put(366.4428,-387.389){\fontsize{11.991}{1}\usefont{T1}{cmr}{m}{n}\selectfont\color{color_29791} }
\put(369.7403,-387.389){\fontsize{11.991}{1}\usefont{T1}{cmr}{m}{n}\selectfont\color{color_29791}p}
\put(376.691,-387.389){\fontsize{11.991}{1}\usefont{T1}{cmr}{m}{n}\selectfont\color{color_29791}l}
\put(379.377,-387.389){\fontsize{11.991}{1}\usefont{T1}{cmr}{m}{n}\selectfont\color{color_29791}e}
\put(386.056,-387.389){\fontsize{11.991}{1}\usefont{T1}{cmr}{m}{n}\selectfont\color{color_29791}a}
\put(392.735,-387.389){\fontsize{11.991}{1}\usefont{T1}{cmr}{m}{n}\selectfont\color{color_29791}se}
\put(405.4095,-387.389){\fontsize{11.991}{1}\usefont{T1}{cmr}{m}{n}\selectfont\color{color_29791} }
\put(408.7909,-387.389){\fontsize{11.991}{1}\usefont{T1}{cmr}{m}{n}\selectfont\color{color_29791}co}
\put(421.4654,-387.389){\fontsize{11.991}{1}\usefont{T1}{cmr}{m}{n}\selectfont\color{color_29791}n}
\put(428.1444,-387.389){\fontsize{11.991}{1}\usefont{T1}{cmr}{m}{n}\selectfont\color{color_29791}t}
\put(431.5259,-387.389){\fontsize{11.991}{1}\usefont{T1}{cmr}{m}{n}\selectfont\color{color_29791}a}
\put(438.2049,-387.389){\fontsize{11.991}{1}\usefont{T1}{cmr}{m}{n}\selectfont\color{color_29791}ct}
\put(447.5818,-387.389){\fontsize{11.991}{1}\usefont{T1}{cmr}{m}{n}\selectfont\color{color_29791} }
\put(450.9633,-387.389){\fontsize{11.991}{1}\usefont{T1}{cmr}{m}{n}\selectfont\color{color_29791}t}
\put(454.2608,-387.389){\fontsize{11.991}{1}\usefont{T1}{cmr}{m}{n}\selectfont\color{color_29791}h}
\put(460.9398,-387.389){\fontsize{11.991}{1}\usefont{T1}{cmr}{m}{n}\selectfont\color{color_29791}e}
\put(467.6188,-387.389){\fontsize{11.991}{1}\usefont{T1}{cmr}{m}{n}\selectfont\color{color_29791} }
\put(471.0002,-387.389){\fontsize{11.991}{1}\usefont{T1}{cmr}{m}{n}\selectfont\color{color_29791}I}
\put(474.3817,-387.389){\fontsize{11.991}{1}\usefont{T1}{cmr}{m}{n}\selectfont\color{color_29791}T}
\put(481.4804,-387.389){\fontsize{11.991}{1}\usefont{T1}{cmr}{m}{n}\selectfont\color{color_29791} }
\put(484.7779,-387.389){\fontsize{11.991}{1}\usefont{T1}{cmr}{m}{n}\selectfont\color{color_29791}Se}
\put(499.5508,-387.389){\fontsize{11.991}{1}\usefont{T1}{cmr}{m}{n}\selectfont\color{color_29791}r}
\put(503.4479,-387.389){\fontsize{11.991}{1}\usefont{T1}{cmr}{m}{n}\selectfont\color{color_29791}vi}
\put(512.1293,-387.389){\fontsize{11.991}{1}\usefont{T1}{cmr}{m}{n}\selectfont\color{color_29791}ce}
\put(524.8038,-387.389){\fontsize{11.991}{1}\usefont{T1}{cmr}{m}{n}\selectfont\color{color_29791}D}
\put(533.4014,-387.389){\fontsize{11.991}{1}\usefont{T1}{cmr}{m}{n}\selectfont\color{color_29791}e}
\put(540.1883,-387.389){\fontsize{11.991}{1}\usefont{T1}{cmr}{m}{n}\selectfont\color{color_29791}s}
\put(546.0879,-387.389){\fontsize{11.991}{1}\usefont{T1}{cmr}{m}{n}\selectfont\color{color_29791}k}
\put(552.1673,-387.389){\fontsize{11.991}{1}\usefont{T1}{cmr}{m}{n}\selectfont\color{color_29791} }
\put(555.4648,-387.389){\fontsize{11.991}{1}\usefont{T1}{cmr}{m}{n}\selectfont\color{color_29791}f}
\put(558.8463,-387.389){\fontsize{11.991}{1}\usefont{T1}{cmr}{m}{n}\selectfont\color{color_29791}o}
\put(565.5253,-387.389){\fontsize{11.991}{1}\usefont{T1}{cmr}{m}{n}\selectfont\color{color_29791}r }
\put(572.8159,-387.389){\fontsize{11.991}{1}\usefont{T1}{cmr}{m}{n}\selectfont\color{color_29791}d}
\put(579.4949,-387.389){\fontsize{11.991}{1}\usefont{T1}{cmr}{m}{n}\selectfont\color{color_29791}e}
\put(586.1739,-387.389){\fontsize{11.991}{1}\usefont{T1}{cmr}{m}{n}\selectfont\color{color_29791}t}
\put(589.5554,-387.389){\fontsize{11.991}{1}\usefont{T1}{cmr}{m}{n}\selectfont\color{color_29791}a}
\put(596.2344,-387.389){\fontsize{11.991}{1}\usefont{T1}{cmr}{m}{n}\selectfont\color{color_29791}i}
\put(598.9203,-387.389){\fontsize{11.991}{1}\usefont{T1}{cmr}{m}{n}\selectfont\color{color_29791}l}
\put(601.5224,-387.389){\fontsize{11.991}{1}\usefont{T1}{cmr}{m}{n}\selectfont\color{color_29791}s}
\end{picture}
\newpage
\begin{tikzpicture}[overlay]
\path(0pt,0pt);
\filldraw[color_283006][even odd rule]
(-15pt, 10pt) -- (704.972pt, 10pt)
 -- (704.972pt, 10pt)
 -- (704.972pt, -529.972pt)
 -- (704.972pt, -529.972pt)
 -- (-15pt, -529.972pt)
 -- (-15pt, -529.972pt)
 -- (-15pt, 10pt) -- cycle
;
\end{tikzpicture}
\begin{picture}(-5,0)(2.5,0)
\put(-15,-529.972){\includegraphics[width=720pt,height=540pt]{latexImage_70f318a8f2de3f58eebb26eed4f2123e.png}}
\put(48.411,-60.69601){\fontsize{21.997}{1}\usefont{T1}{cmr}{b}{n}\selectfont\color{color_29791}Da}
\put(76.50117,-60.69601){\fontsize{21.997}{1}\usefont{T1}{cmr}{b}{n}\selectfont\color{color_29791}t}
\put(83.80417,-60.69601){\fontsize{21.997}{1}\usefont{T1}{cmr}{b}{n}\selectfont\color{color_29791}a}
\put(96.07849,-60.69601){\fontsize{21.997}{1}\usefont{T1}{cmr}{b}{n}\selectfont\color{color_29791} }
\put(101.2698,-60.69601){\fontsize{21.997}{1}\usefont{T1}{cmr}{b}{n}\selectfont\color{color_29791}Ac}
\put(129.4259,-60.69601){\fontsize{21.997}{1}\usefont{T1}{cmr}{b}{n}\selectfont\color{color_29791}c}
\put(141.6343,-60.69601){\fontsize{21.997}{1}\usefont{T1}{cmr}{b}{n}\selectfont\color{color_29791}e}
\put(153.8426,-60.69601){\fontsize{21.997}{1}\usefont{T1}{cmr}{b}{n}\selectfont\color{color_29791}s}
\put(166.1169,-60.69601){\fontsize{21.997}{1}\usefont{T1}{cmr}{b}{n}\selectfont\color{color_29791}s}
\put(178.3253,-60.69601){\fontsize{21.997}{1}\usefont{T1}{cmr}{b}{n}\selectfont\color{color_29791} a}
\put(196.6268,-60.69601){\fontsize{21.997}{1}\usefont{T1}{cmr}{b}{n}\selectfont\color{color_29791}n}
\put(210.1109,-60.69601){\fontsize{21.997}{1}\usefont{T1}{cmr}{b}{n}\selectfont\color{color_29791}d}
\put(223.5071,-60.69601){\fontsize{21.997}{1}\usefont{T1}{cmr}{b}{n}\selectfont\color{color_29791} }
\put(229.6882,-60.69601){\fontsize{21.997}{1}\usefont{T1}{cmr}{b}{n}\selectfont\color{color_29791}T}
\put(241.8966,-60.69601){\fontsize{21.997}{1}\usefont{T1}{cmr}{b}{n}\selectfont\color{color_29791}r}
\put(250.3874,-60.69601){\fontsize{21.997}{1}\usefont{T1}{cmr}{b}{n}\selectfont\color{color_29791}a}
\put(262.6617,-60.69601){\fontsize{21.997}{1}\usefont{T1}{cmr}{b}{n}\selectfont\color{color_29791}n}
\put(276.0579,-60.69601){\fontsize{21.997}{1}\usefont{T1}{cmr}{b}{n}\selectfont\color{color_29791}s}
\put(288.2663,-60.69601){\fontsize{21.997}{1}\usefont{T1}{cmr}{b}{n}\selectfont\color{color_29791}f}
\put(295.6353,-60.69601){\fontsize{21.997}{1}\usefont{T1}{cmr}{b}{n}\selectfont\color{color_29791}e}
\put(307.8436,-60.69601){\fontsize{21.997}{1}\usefont{T1}{cmr}{b}{n}\selectfont\color{color_29791}r}
\put(316.3344,-60.69601){\fontsize{21.997}{1}\usefont{T1}{cmr}{b}{n}\selectfont\color{color_29791} (}
\put(329.7306,-60.69601){\fontsize{21.997}{1}\usefont{T1}{cmr}{b}{n}\selectfont\color{color_29791}c}
\put(342.0049,-60.69601){\fontsize{21.997}{1}\usefont{T1}{cmr}{b}{n}\selectfont\color{color_29791}o}
\put(355.4011,-60.69601){\fontsize{21.997}{1}\usefont{T1}{cmr}{b}{n}\selectfont\color{color_29791}n}
\put(368.8853,-60.69601){\fontsize{21.997}{1}\usefont{T1}{cmr}{b}{n}\selectfont\color{color_29791}t}
\put(376.1883,-60.69601){\fontsize{21.997}{1}\usefont{T1}{cmr}{b}{n}\selectfont\color{color_29791}in}
\put(395.7656,-60.69601){\fontsize{21.997}{1}\usefont{T1}{cmr}{b}{n}\selectfont\color{color_29791}u}
\put(409.1618,-60.69601){\fontsize{21.997}{1}\usefont{T1}{cmr}{b}{n}\selectfont\color{color_29791}e}
\put(421.3701,-60.69601){\fontsize{21.997}{1}\usefont{T1}{cmr}{b}{n}\selectfont\color{color_29791}d}
\put(434.8543,-60.69601){\fontsize{21.997}{1}\usefont{T1}{cmr}{b}{n}\selectfont\color{color_29791})}
\put(48.213,-115.49){\fontsize{15.987}{1}\usefont{T1}{cmr}{m}{n}\selectfont\color{color_29791}•}
\put(62.187,-116.51){\fontsize{15.987}{1}\usefont{T1}{cmr}{m}{n}\selectfont\color{color_29791}r}
\put(67.47869,-116.51){\fontsize{15.987}{1}\usefont{T1}{cmr}{m}{n}\selectfont\color{color_29791}sy}
\put(83.54562,-116.51){\fontsize{15.987}{1}\usefont{T1}{cmr}{m}{n}\selectfont\color{color_29791}nc}
\put(84.213,-145.31){\fontsize{11.991}{1}\usefont{T1}{cmr}{m}{n}\selectfont\color{color_29791}•}
\put(106.691,-145.991){\fontsize{11.991}{1}\usefont{T1}{cmr}{m}{n}\selectfont\color{color_29791}P}
\put(114.7729,-145.991){\fontsize{11.991}{1}\usefont{T1}{cmr}{m}{n}\selectfont\color{color_29791}e}
\put(121.4519,-145.991){\fontsize{11.991}{1}\usefont{T1}{cmr}{m}{n}\selectfont\color{color_29791}r}
\put(125.349,-145.991){\fontsize{11.991}{1}\usefont{T1}{cmr}{m}{n}\selectfont\color{color_29791}f}
\put(128.7305,-145.991){\fontsize{11.991}{1}\usefont{T1}{cmr}{m}{n}\selectfont\color{color_29791}e}
\put(135.4095,-145.991){\fontsize{11.991}{1}\usefont{T1}{cmr}{m}{n}\selectfont\color{color_29791}ct}
\put(144.7864,-145.991){\fontsize{11.991}{1}\usefont{T1}{cmr}{m}{n}\selectfont\color{color_29791} }
\put(148.084,-145.991){\fontsize{11.991}{1}\usefont{T1}{cmr}{m}{n}\selectfont\color{color_29791}t}
\put(151.4654,-145.991){\fontsize{11.991}{1}\usefont{T1}{cmr}{m}{n}\selectfont\color{color_29791}o}
\put(158.1444,-145.991){\fontsize{11.991}{1}\usefont{T1}{cmr}{m}{n}\selectfont\color{color_29791}o}
\put(164.8234,-145.991){\fontsize{11.991}{1}\usefont{T1}{cmr}{m}{n}\selectfont\color{color_29791}l}
\put(167.5094,-145.991){\fontsize{11.991}{1}\usefont{T1}{cmr}{m}{n}\selectfont\color{color_29791} }
\put(170.8069,-145.991){\fontsize{11.991}{1}\usefont{T1}{cmr}{m}{n}\selectfont\color{color_29791}f}
\put(174.1884,-145.991){\fontsize{11.991}{1}\usefont{T1}{cmr}{m}{n}\selectfont\color{color_29791}o}
\put(180.8674,-145.991){\fontsize{11.991}{1}\usefont{T1}{cmr}{m}{n}\selectfont\color{color_29791}r }
\put(188.1579,-145.991){\fontsize{11.991}{1}\usefont{T1}{cmr}{m}{n}\selectfont\color{color_29791}t}
\put(191.5394,-145.991){\fontsize{11.991}{1}\usefont{T1}{cmr}{m}{n}\selectfont\color{color_29791}ra}
\put(202.2114,-145.991){\fontsize{11.991}{1}\usefont{T1}{cmr}{m}{n}\selectfont\color{color_29791}n}
\put(208.8904,-145.991){\fontsize{11.991}{1}\usefont{T1}{cmr}{m}{n}\selectfont\color{color_29791}sf}
\put(218.1834,-145.991){\fontsize{11.991}{1}\usefont{T1}{cmr}{m}{n}\selectfont\color{color_29791}e}
\put(224.9703,-145.991){\fontsize{11.991}{1}\usefont{T1}{cmr}{m}{n}\selectfont\color{color_29791}r}
\put(228.8674,-145.991){\fontsize{11.991}{1}\usefont{T1}{cmr}{m}{n}\selectfont\color{color_29791}ri}
\put(235.4624,-145.991){\fontsize{11.991}{1}\usefont{T1}{cmr}{m}{n}\selectfont\color{color_29791}n}
\put(242.2493,-145.991){\fontsize{11.991}{1}\usefont{T1}{cmr}{m}{n}\selectfont\color{color_29791}g}
\put(248.9283,-145.991){\fontsize{11.991}{1}\usefont{T1}{cmr}{m}{n}\selectfont\color{color_29791} }
\put(252.2258,-145.991){\fontsize{11.991}{1}\usefont{T1}{cmr}{m}{n}\selectfont\color{color_29791}f}
\put(255.6073,-145.991){\fontsize{11.991}{1}\usefont{T1}{cmr}{m}{n}\selectfont\color{color_29791}i}
\put(258.2094,-145.991){\fontsize{11.991}{1}\usefont{T1}{cmr}{m}{n}\selectfont\color{color_29791}l}
\put(260.8953,-145.991){\fontsize{11.991}{1}\usefont{T1}{cmr}{m}{n}\selectfont\color{color_29791}e}
\put(267.5743,-145.991){\fontsize{11.991}{1}\usefont{T1}{cmr}{m}{n}\selectfont\color{color_29791} }
\put(270.9558,-145.991){\fontsize{11.991}{1}\usefont{T1}{cmr}{m}{n}\selectfont\color{color_29791}l}
\put(273.5578,-145.991){\fontsize{11.991}{1}\usefont{T1}{cmr}{m}{n}\selectfont\color{color_29791}o}
\put(280.2368,-145.991){\fontsize{11.991}{1}\usefont{T1}{cmr}{m}{n}\selectfont\color{color_29791}ca}
\put(292.9113,-145.991){\fontsize{11.991}{1}\usefont{T1}{cmr}{m}{n}\selectfont\color{color_29791}l}
\put(295.5973,-145.991){\fontsize{11.991}{1}\usefont{T1}{cmr}{m}{n}\selectfont\color{color_29791}l}
\put(298.1993,-145.991){\fontsize{11.991}{1}\usefont{T1}{cmr}{m}{n}\selectfont\color{color_29791}y }
\put(307.5763,-145.991){\fontsize{11.991}{1}\usefont{T1}{cmr}{m}{n}\selectfont\color{color_29791}a}
\put(314.2553,-145.991){\fontsize{11.991}{1}\usefont{T1}{cmr}{m}{n}\selectfont\color{color_29791}n}
\put(320.9343,-145.991){\fontsize{11.991}{1}\usefont{T1}{cmr}{m}{n}\selectfont\color{color_29791}d}
\put(327.6133,-145.991){\fontsize{11.991}{1}\usefont{T1}{cmr}{m}{n}\selectfont\color{color_29791} }
\put(330.9947,-145.991){\fontsize{11.991}{1}\usefont{T1}{cmr}{m}{n}\selectfont\color{color_29791}se}
\put(343.6692,-145.991){\fontsize{11.991}{1}\usefont{T1}{cmr}{m}{n}\selectfont\color{color_29791}cu}
\put(356.3437,-145.991){\fontsize{11.991}{1}\usefont{T1}{cmr}{m}{n}\selectfont\color{color_29791}re}
\put(367.0157,-145.991){\fontsize{11.991}{1}\usefont{T1}{cmr}{m}{n}\selectfont\color{color_29791}l}
\put(369.7017,-145.991){\fontsize{11.991}{1}\usefont{T1}{cmr}{m}{n}\selectfont\color{color_29791}y }
\put(378.9947,-145.991){\fontsize{11.991}{1}\usefont{T1}{cmr}{m}{n}\selectfont\color{color_29791}o}
\put(385.6737,-145.991){\fontsize{11.991}{1}\usefont{T1}{cmr}{m}{n}\selectfont\color{color_29791}v}
\put(391.7531,-145.991){\fontsize{11.991}{1}\usefont{T1}{cmr}{m}{n}\selectfont\color{color_29791}e}
\put(398.4321,-145.991){\fontsize{11.991}{1}\usefont{T1}{cmr}{m}{n}\selectfont\color{color_29791}r}
\put(402.3292,-145.991){\fontsize{11.991}{1}\usefont{T1}{cmr}{m}{n}\selectfont\color{color_29791} }
\put(405.7107,-145.991){\fontsize{11.991}{1}\usefont{T1}{cmr}{m}{n}\selectfont\color{color_29791}t}
\put(409.0082,-145.991){\fontsize{11.991}{1}\usefont{T1}{cmr}{m}{n}\selectfont\color{color_29791}h}
\put(415.6872,-145.991){\fontsize{11.991}{1}\usefont{T1}{cmr}{m}{n}\selectfont\color{color_29791}e}
\put(422.4741,-145.991){\fontsize{11.991}{1}\usefont{T1}{cmr}{m}{n}\selectfont\color{color_29791} }
\put(425.7716,-145.991){\fontsize{11.991}{1}\usefont{T1}{cmr}{m}{n}\selectfont\color{color_29791}i}
\put(428.4576,-145.991){\fontsize{11.991}{1}\usefont{T1}{cmr}{m}{n}\selectfont\color{color_29791}n}
\put(435.1365,-145.991){\fontsize{11.991}{1}\usefont{T1}{cmr}{m}{n}\selectfont\color{color_29791}t}
\put(438.518,-145.991){\fontsize{11.991}{1}\usefont{T1}{cmr}{m}{n}\selectfont\color{color_29791}e}
\put(445.197,-145.991){\fontsize{11.991}{1}\usefont{T1}{cmr}{m}{n}\selectfont\color{color_29791}rn}
\put(455.869,-145.991){\fontsize{11.991}{1}\usefont{T1}{cmr}{m}{n}\selectfont\color{color_29791}e}
\put(462.548,-145.991){\fontsize{11.991}{1}\usefont{T1}{cmr}{m}{n}\selectfont\color{color_29791}t}
\put(84.213,-174.11){\fontsize{11.991}{1}\usefont{T1}{cmr}{m}{n}\selectfont\color{color_29791}•}
\put(106.691,-174.791){\fontsize{11.991}{1}\usefont{T1}{cmr}{m}{n}\selectfont\color{color_29791}O}
\put(116.068,-174.791){\fontsize{11.991}{1}\usefont{T1}{cmr}{m}{n}\selectfont\color{color_29791}p}
\put(122.7469,-174.791){\fontsize{11.991}{1}\usefont{T1}{cmr}{m}{n}\selectfont\color{color_29791}t}
\put(126.1284,-174.791){\fontsize{11.991}{1}\usefont{T1}{cmr}{m}{n}\selectfont\color{color_29791}i}
\put(128.7305,-174.791){\fontsize{11.991}{1}\usefont{T1}{cmr}{m}{n}\selectfont\color{color_29791}o}
\put(135.4095,-174.791){\fontsize{11.991}{1}\usefont{T1}{cmr}{m}{n}\selectfont\color{color_29791}n}
\put(142.0884,-174.791){\fontsize{11.991}{1}\usefont{T1}{cmr}{m}{n}\selectfont\color{color_29791}s }
\put(151.4654,-174.791){\fontsize{11.991}{1}\usefont{T1}{cmr}{m}{n}\selectfont\color{color_29791}t}
\put(154.8469,-174.791){\fontsize{11.991}{1}\usefont{T1}{cmr}{m}{n}\selectfont\color{color_29791}o}
\put(161.5258,-174.791){\fontsize{11.991}{1}\usefont{T1}{cmr}{m}{n}\selectfont\color{color_29791} }
\put(164.8234,-174.791){\fontsize{11.991}{1}\usefont{T1}{cmr}{m}{n}\selectfont\color{color_29791}re}
\put(175.4954,-174.791){\fontsize{11.991}{1}\usefont{T1}{cmr}{m}{n}\selectfont\color{color_29791}su}
\put(188.1699,-174.791){\fontsize{11.991}{1}\usefont{T1}{cmr}{m}{n}\selectfont\color{color_29791}me}
\put(204.8374,-174.791){\fontsize{11.991}{1}\usefont{T1}{cmr}{m}{n}\selectfont\color{color_29791},}
\put(208.2188,-174.791){\fontsize{11.991}{1}\usefont{T1}{cmr}{m}{n}\selectfont\color{color_29791} }
\put(211.5163,-174.791){\fontsize{11.991}{1}\usefont{T1}{cmr}{m}{n}\selectfont\color{color_29791}re}
\put(222.1883,-174.791){\fontsize{11.991}{1}\usefont{T1}{cmr}{m}{n}\selectfont\color{color_29791}co}
\put(234.8628,-174.791){\fontsize{11.991}{1}\usefont{T1}{cmr}{m}{n}\selectfont\color{color_29791}n}
\put(241.6497,-174.791){\fontsize{11.991}{1}\usefont{T1}{cmr}{m}{n}\selectfont\color{color_29791}n}
\put(248.3287,-174.791){\fontsize{11.991}{1}\usefont{T1}{cmr}{m}{n}\selectfont\color{color_29791}e}
\put(255.0077,-174.791){\fontsize{11.991}{1}\usefont{T1}{cmr}{m}{n}\selectfont\color{color_29791}ct}
\put(264.3846,-174.791){\fontsize{11.991}{1}\usefont{T1}{cmr}{m}{n}\selectfont\color{color_29791},}
\put(267.6822,-174.791){\fontsize{11.991}{1}\usefont{T1}{cmr}{m}{n}\selectfont\color{color_29791} }
\put(271.0636,-174.791){\fontsize{11.991}{1}\usefont{T1}{cmr}{m}{n}\selectfont\color{color_29791}co}
\put(283.7381,-174.791){\fontsize{11.991}{1}\usefont{T1}{cmr}{m}{n}\selectfont\color{color_29791}m}
\put(293.6307,-174.791){\fontsize{11.991}{1}\usefont{T1}{cmr}{m}{n}\selectfont\color{color_29791}p}
\put(300.4176,-174.791){\fontsize{11.991}{1}\usefont{T1}{cmr}{m}{n}\selectfont\color{color_29791}r}
\put(304.3147,-174.791){\fontsize{11.991}{1}\usefont{T1}{cmr}{m}{n}\selectfont\color{color_29791}e}
\put(311.1016,-174.791){\fontsize{11.991}{1}\usefont{T1}{cmr}{m}{n}\selectfont\color{color_29791}ssi}
\put(325.6946,-174.791){\fontsize{11.991}{1}\usefont{T1}{cmr}{m}{n}\selectfont\color{color_29791}o}
\put(332.3736,-174.791){\fontsize{11.991}{1}\usefont{T1}{cmr}{m}{n}\selectfont\color{color_29791}n}
\put(339.0526,-174.791){\fontsize{11.991}{1}\usefont{T1}{cmr}{m}{n}\selectfont\color{color_29791},}
\put(342.4341,-174.791){\fontsize{11.991}{1}\usefont{T1}{cmr}{m}{n}\selectfont\color{color_29791} }
\put(345.8155,-174.791){\fontsize{11.991}{1}\usefont{T1}{cmr}{m}{n}\selectfont\color{color_29791}l}
\put(348.4176,-174.791){\fontsize{11.991}{1}\usefont{T1}{cmr}{m}{n}\selectfont\color{color_29791}i}
\put(351.0196,-174.791){\fontsize{11.991}{1}\usefont{T1}{cmr}{m}{n}\selectfont\color{color_29791}mi}
\put(363.6102,-174.791){\fontsize{11.991}{1}\usefont{T1}{cmr}{m}{n}\selectfont\color{color_29791}t}
\put(366.9916,-174.791){\fontsize{11.991}{1}\usefont{T1}{cmr}{m}{n}\selectfont\color{color_29791} }
\put(370.3731,-174.791){\fontsize{11.991}{1}\usefont{T1}{cmr}{m}{n}\selectfont\color{color_29791}t}
\put(373.6706,-174.791){\fontsize{11.991}{1}\usefont{T1}{cmr}{m}{n}\selectfont\color{color_29791}ra}
\put(384.3426,-174.791){\fontsize{11.991}{1}\usefont{T1}{cmr}{m}{n}\selectfont\color{color_29791}n}
\put(391.0216,-174.791){\fontsize{11.991}{1}\usefont{T1}{cmr}{m}{n}\selectfont\color{color_29791}sf}
\put(400.3986,-174.791){\fontsize{11.991}{1}\usefont{T1}{cmr}{m}{n}\selectfont\color{color_29791}e}
\put(407.0775,-174.791){\fontsize{11.991}{1}\usefont{T1}{cmr}{m}{n}\selectfont\color{color_29791}rr}
\put(414.9676,-174.791){\fontsize{11.991}{1}\usefont{T1}{cmr}{m}{n}\selectfont\color{color_29791}e}
\put(421.6466,-174.791){\fontsize{11.991}{1}\usefont{T1}{cmr}{m}{n}\selectfont\color{color_29791}d}
\put(428.4335,-174.791){\fontsize{11.991}{1}\usefont{T1}{cmr}{m}{n}\selectfont\color{color_29791} }
\put(431.731,-174.791){\fontsize{11.991}{1}\usefont{T1}{cmr}{m}{n}\selectfont\color{color_29791}ra}
\put(442.403,-174.791){\fontsize{11.991}{1}\usefont{T1}{cmr}{m}{n}\selectfont\color{color_29791}t}
\put(445.7845,-174.791){\fontsize{11.991}{1}\usefont{T1}{cmr}{m}{n}\selectfont\color{color_29791}e}
\put(452.4635,-174.791){\fontsize{11.991}{1}\usefont{T1}{cmr}{m}{n}\selectfont\color{color_29791}s.}
\put(48.213,-207.389){\fontsize{15.987}{1}\usefont{T1}{cmr}{m}{n}\selectfont\color{color_29791}•}
\put(62.187,-208.409){\fontsize{15.987}{1}\usefont{T1}{cmr}{m}{n}\selectfont\color{color_29791}s}
\put(70.26043,-208.409){\fontsize{15.987}{1}\usefont{T1}{cmr}{m}{n}\selectfont\color{color_29791}cp}
\put(48.213,-241.802){\fontsize{15.987}{1}\usefont{T1}{cmr}{m}{n}\selectfont\color{color_29791}•}
\put(62.187,-242.709){\fontsize{15.987}{1}\usefont{T1}{cmr}{m}{n}\selectfont\color{color_29791}s}
\put(70.26043,-242.709){\fontsize{15.987}{1}\usefont{T1}{cmr}{m}{n}\selectfont\color{color_29791}shf}
\put(91.53912,-242.709){\fontsize{15.987}{1}\usefont{T1}{cmr}{m}{n}\selectfont\color{color_29791}s}
\end{picture}
\newpage
\begin{tikzpicture}[overlay]
\path(0pt,0pt);
\filldraw[color_283006][even odd rule]
(-15pt, 10pt) -- (704.972pt, 10pt)
 -- (704.972pt, 10pt)
 -- (704.972pt, -529.972pt)
 -- (704.972pt, -529.972pt)
 -- (-15pt, -529.972pt)
 -- (-15pt, -529.972pt)
 -- (-15pt, 10pt) -- cycle
;
\end{tikzpicture}
\begin{picture}(-5,0)(2.5,0)
\put(-15,-529.972){\includegraphics[width=720pt,height=540pt]{latexImage_2828171a63492229494bbe96f89381b1.png}}
\put(48.411,-60.69601){\fontsize{21.997}{1}\usefont{T1}{cmr}{b}{n}\selectfont\color{color_29791}Us}
\put(76.50117,-60.69601){\fontsize{21.997}{1}\usefont{T1}{cmr}{b}{n}\selectfont\color{color_29791}e}
\put(88.7095,-60.69601){\fontsize{21.997}{1}\usefont{T1}{cmr}{b}{n}\selectfont\color{color_29791}r}
\put(97.28833,-60.69601){\fontsize{21.997}{1}\usefont{T1}{cmr}{b}{n}\selectfont\color{color_29791} E}
\put(117.9875,-60.69601){\fontsize{21.997}{1}\usefont{T1}{cmr}{b}{n}\selectfont\color{color_29791}n}
\put(131.4717,-60.69601){\fontsize{21.997}{1}\usefont{T1}{cmr}{b}{n}\selectfont\color{color_29791}v}
\put(143.68,-60.69601){\fontsize{21.997}{1}\usefont{T1}{cmr}{b}{n}\selectfont\color{color_29791}ir}
\put(158.352,-60.69601){\fontsize{21.997}{1}\usefont{T1}{cmr}{b}{n}\selectfont\color{color_29791}o}
\put(171.7482,-60.69601){\fontsize{21.997}{1}\usefont{T1}{cmr}{b}{n}\selectfont\color{color_29791}n}
\put(185.1443,-60.69601){\fontsize{21.997}{1}\usefont{T1}{cmr}{b}{n}\selectfont\color{color_29791}m}
\put(204.7217,-60.69601){\fontsize{21.997}{1}\usefont{T1}{cmr}{b}{n}\selectfont\color{color_29791}e}
\put(216.93,-60.69601){\fontsize{21.997}{1}\usefont{T1}{cmr}{b}{n}\selectfont\color{color_29791}n}
\put(230.4142,-60.69601){\fontsize{21.997}{1}\usefont{T1}{cmr}{b}{n}\selectfont\color{color_29791}t}
\put(48.213,-115.49){\fontsize{15.987}{1}\usefont{T1}{cmr}{m}{n}\selectfont\color{color_29791}•}
\put(62.187,-116.51){\fontsize{15.987}{1}\usefont{T1}{cmr}{m}{n}\selectfont\color{color_29791}S}
\put(72.86631,-116.51){\fontsize{15.987}{1}\usefont{T1}{cmr}{m}{n}\selectfont\color{color_29791}he}
\put(90.54793,-116.51){\fontsize{15.987}{1}\usefont{T1}{cmr}{m}{n}\selectfont\color{color_29791}l}
\put(94.12901,-116.51){\fontsize{15.987}{1}\usefont{T1}{cmr}{m}{n}\selectfont\color{color_29791}l}
\put(97.7101,-116.51){\fontsize{15.987}{1}\usefont{T1}{cmr}{m}{n}\selectfont\color{color_29791} }
\put(84.213,-145.31){\fontsize{11.991}{1}\usefont{T1}{cmr}{m}{n}\selectfont\color{color_29791}•}
\put(98.187,-145.991){\fontsize{11.991}{1}\usefont{T1}{cmr}{m}{n}\selectfont\color{color_29791}O}
\put(107.564,-145.991){\fontsize{11.991}{1}\usefont{T1}{cmr}{m}{n}\selectfont\color{color_29791}u}
\put(114.243,-145.991){\fontsize{11.991}{1}\usefont{T1}{cmr}{m}{n}\selectfont\color{color_29791}r}
\put(118.313,-145.991){\fontsize{11.991}{1}\usefont{T1}{cmr}{m}{n}\selectfont\color{color_29791} }
\put(121.602,-145.991){\fontsize{11.991}{1}\usefont{T1}{cmr}{m}{n}\selectfont\color{color_29791}d}
\put(128.281,-145.991){\fontsize{11.991}{1}\usefont{T1}{cmr}{m}{n}\selectfont\color{color_29791}e}
\put(135.0679,-145.991){\fontsize{11.991}{1}\usefont{T1}{cmr}{m}{n}\selectfont\color{color_29791}f}
\put(138.3654,-145.991){\fontsize{11.991}{1}\usefont{T1}{cmr}{m}{n}\selectfont\color{color_29791}a}
\put(145.0444,-145.991){\fontsize{11.991}{1}\usefont{T1}{cmr}{m}{n}\selectfont\color{color_29791}u}
\put(151.7234,-145.991){\fontsize{11.991}{1}\usefont{T1}{cmr}{m}{n}\selectfont\color{color_29791}l}
\put(154.4094,-145.991){\fontsize{11.991}{1}\usefont{T1}{cmr}{m}{n}\selectfont\color{color_29791}t}
\put(157.7909,-145.991){\fontsize{11.991}{1}\usefont{T1}{cmr}{m}{n}\selectfont\color{color_29791} }
\put(161.0884,-145.991){\fontsize{11.991}{1}\usefont{T1}{cmr}{m}{n}\selectfont\color{color_29791}sh}
\put(173.7629,-145.991){\fontsize{11.991}{1}\usefont{T1}{cmr}{m}{n}\selectfont\color{color_29791}e}
\put(180.4419,-145.991){\fontsize{11.991}{1}\usefont{T1}{cmr}{m}{n}\selectfont\color{color_29791}l}
\put(183.1279,-145.991){\fontsize{11.991}{1}\usefont{T1}{cmr}{m}{n}\selectfont\color{color_29791}l}
\put(185.7299,-145.991){\fontsize{11.991}{1}\usefont{T1}{cmr}{m}{n}\selectfont\color{color_29791} }
\put(189.1114,-145.991){\fontsize{11.991}{1}\usefont{T1}{cmr}{m}{n}\selectfont\color{color_29791}i}
\put(191.7134,-145.991){\fontsize{11.991}{1}\usefont{T1}{cmr}{m}{n}\selectfont\color{color_29791}s }
\put(201.0904,-145.991){\fontsize{11.991}{1}\usefont{T1}{cmr}{m}{n}\selectfont\color{color_29791}t}
\put(204.4718,-145.991){\fontsize{11.991}{1}\usefont{T1}{cmr}{m}{n}\selectfont\color{color_29791}cs}
\put(216.3669,-145.991){\fontsize{11.991}{1}\usefont{T1}{cmr}{m}{n}\selectfont\color{color_29791}h}
\put(84.213,-174.11){\fontsize{11.991}{1}\usefont{T1}{cmr}{m}{n}\selectfont\color{color_29791}•}
\put(98.187,-174.791){\fontsize{11.991}{1}\usefont{T1}{cmr}{m}{n}\selectfont\color{color_29791}I}
\put(101.5685,-174.791){\fontsize{11.991}{1}\usefont{T1}{cmr}{m}{n}\selectfont\color{color_29791}f}
\put(104.866,-174.791){\fontsize{11.991}{1}\usefont{T1}{cmr}{m}{n}\selectfont\color{color_29791} }
\put(108.2474,-174.791){\fontsize{11.991}{1}\usefont{T1}{cmr}{m}{n}\selectfont\color{color_29791}yo}
\put(120.9219,-174.791){\fontsize{11.991}{1}\usefont{T1}{cmr}{m}{n}\selectfont\color{color_29791}u}
\put(127.6009,-174.791){\fontsize{11.991}{1}\usefont{T1}{cmr}{m}{n}\selectfont\color{color_29791} }
\put(130.9824,-174.791){\fontsize{11.991}{1}\usefont{T1}{cmr}{m}{n}\selectfont\color{color_29791}p}
\put(137.6614,-174.791){\fontsize{11.991}{1}\usefont{T1}{cmr}{m}{n}\selectfont\color{color_29791}re}
\put(148.3334,-174.791){\fontsize{11.991}{1}\usefont{T1}{cmr}{m}{n}\selectfont\color{color_29791}f}
\put(151.6309,-174.791){\fontsize{11.991}{1}\usefont{T1}{cmr}{m}{n}\selectfont\color{color_29791}e}
\put(158.4178,-174.791){\fontsize{11.991}{1}\usefont{T1}{cmr}{m}{n}\selectfont\color{color_29791}r}
\put(162.3149,-174.791){\fontsize{11.991}{1}\usefont{T1}{cmr}{m}{n}\selectfont\color{color_29791} }
\put(165.6964,-174.791){\fontsize{11.991}{1}\usefont{T1}{cmr}{m}{n}\selectfont\color{color_29791}so}
\put(178.3708,-174.791){\fontsize{11.991}{1}\usefont{T1}{cmr}{m}{n}\selectfont\color{color_29791}me}
\put(195.0383,-174.791){\fontsize{11.991}{1}\usefont{T1}{cmr}{m}{n}\selectfont\color{color_29791}t}
\put(198.3359,-174.791){\fontsize{11.991}{1}\usefont{T1}{cmr}{m}{n}\selectfont\color{color_29791}h}
\put(205.1228,-174.791){\fontsize{11.991}{1}\usefont{T1}{cmr}{m}{n}\selectfont\color{color_29791}i}
\put(207.7248,-174.791){\fontsize{11.991}{1}\usefont{T1}{cmr}{m}{n}\selectfont\color{color_29791}n}
\put(214.4038,-174.791){\fontsize{11.991}{1}\usefont{T1}{cmr}{m}{n}\selectfont\color{color_29791}g}
\put(221.1907,-174.791){\fontsize{11.991}{1}\usefont{T1}{cmr}{m}{n}\selectfont\color{color_29791} }
\put(224.4883,-174.791){\fontsize{11.991}{1}\usefont{T1}{cmr}{m}{n}\selectfont\color{color_29791}d}
\put(231.1672,-174.791){\fontsize{11.991}{1}\usefont{T1}{cmr}{m}{n}\selectfont\color{color_29791}i}
\put(233.8532,-174.791){\fontsize{11.991}{1}\usefont{T1}{cmr}{m}{n}\selectfont\color{color_29791}f}
\put(236.9469,-174.791){\fontsize{11.991}{1}\usefont{T1}{cmr}{m}{n}\selectfont\color{color_29791}f}
\put(240.3284,-174.791){\fontsize{11.991}{1}\usefont{T1}{cmr}{m}{n}\selectfont\color{color_29791}e}
\put(247.0074,-174.791){\fontsize{11.991}{1}\usefont{T1}{cmr}{m}{n}\selectfont\color{color_29791}r}
\put(250.9044,-174.791){\fontsize{11.991}{1}\usefont{T1}{cmr}{m}{n}\selectfont\color{color_29791}e}
\put(257.6913,-174.791){\fontsize{11.991}{1}\usefont{T1}{cmr}{m}{n}\selectfont\color{color_29791}n}
\put(264.3703,-174.791){\fontsize{11.991}{1}\usefont{T1}{cmr}{m}{n}\selectfont\color{color_29791}t}
\put(267.6678,-174.791){\fontsize{11.991}{1}\usefont{T1}{cmr}{m}{n}\selectfont\color{color_29791} }
\put(271.0493,-174.791){\fontsize{11.991}{1}\usefont{T1}{cmr}{m}{n}\selectfont\color{color_29791}su}
\put(283.7238,-174.791){\fontsize{11.991}{1}\usefont{T1}{cmr}{m}{n}\selectfont\color{color_29791}ch}
\put(296.3983,-174.791){\fontsize{11.991}{1}\usefont{T1}{cmr}{m}{n}\selectfont\color{color_29791} }
\put(299.7798,-174.791){\fontsize{11.991}{1}\usefont{T1}{cmr}{m}{n}\selectfont\color{color_29791}a}
\put(306.4587,-174.791){\fontsize{11.991}{1}\usefont{T1}{cmr}{m}{n}\selectfont\color{color_29791}s }
\put(315.8357,-174.791){\fontsize{11.991}{1}\usefont{T1}{cmr}{m}{n}\selectfont\color{color_29791}b}
\put(322.5147,-174.791){\fontsize{11.991}{1}\usefont{T1}{cmr}{m}{n}\selectfont\color{color_29791}a}
\put(329.1937,-174.791){\fontsize{11.991}{1}\usefont{T1}{cmr}{m}{n}\selectfont\color{color_29791}sh}
\put(341.8682,-174.791){\fontsize{11.991}{1}\usefont{T1}{cmr}{m}{n}\selectfont\color{color_29791},}
\put(345.2496,-174.791){\fontsize{11.991}{1}\usefont{T1}{cmr}{m}{n}\selectfont\color{color_29791} }
\put(348.6311,-174.791){\fontsize{11.991}{1}\usefont{T1}{cmr}{m}{n}\selectfont\color{color_29791}co}
\put(361.3056,-174.791){\fontsize{11.991}{1}\usefont{T1}{cmr}{m}{n}\selectfont\color{color_29791}n}
\put(367.9846,-174.791){\fontsize{11.991}{1}\usefont{T1}{cmr}{m}{n}\selectfont\color{color_29791}t}
\put(371.2821,-174.791){\fontsize{11.991}{1}\usefont{T1}{cmr}{m}{n}\selectfont\color{color_29791}a}
\put(378.069,-174.791){\fontsize{11.991}{1}\usefont{T1}{cmr}{m}{n}\selectfont\color{color_29791}ct}
\put(387.362,-174.791){\fontsize{11.991}{1}\usefont{T1}{cmr}{m}{n}\selectfont\color{color_29791} }
\put(390.7435,-174.791){\fontsize{11.991}{1}\usefont{T1}{cmr}{m}{n}\selectfont\color{color_29791}t}
\put(394.125,-174.791){\fontsize{11.991}{1}\usefont{T1}{cmr}{m}{n}\selectfont\color{color_29791}h}
\put(400.804,-174.791){\fontsize{11.991}{1}\usefont{T1}{cmr}{m}{n}\selectfont\color{color_29791}e}
\put(407.4829,-174.791){\fontsize{11.991}{1}\usefont{T1}{cmr}{m}{n}\selectfont\color{color_29791} }
\put(410.7805,-174.791){\fontsize{11.991}{1}\usefont{T1}{cmr}{m}{n}\selectfont\color{color_29791}s}
\put(416.8599,-174.791){\fontsize{11.991}{1}\usefont{T1}{cmr}{m}{n}\selectfont\color{color_29791}e}
\put(423.5389,-174.791){\fontsize{11.991}{1}\usefont{T1}{cmr}{m}{n}\selectfont\color{color_29791}r}
\put(427.4359,-174.791){\fontsize{11.991}{1}\usefont{T1}{cmr}{m}{n}\selectfont\color{color_29791}vi}
\put(436.1174,-174.791){\fontsize{11.991}{1}\usefont{T1}{cmr}{m}{n}\selectfont\color{color_29791}ce}
\put(448.7919,-174.791){\fontsize{11.991}{1}\usefont{T1}{cmr}{m}{n}\selectfont\color{color_29791} }
\put(452.1734,-174.791){\fontsize{11.991}{1}\usefont{T1}{cmr}{m}{n}\selectfont\color{color_29791}d}
\put(458.8524,-174.791){\fontsize{11.991}{1}\usefont{T1}{cmr}{m}{n}\selectfont\color{color_29791}e}
\put(465.5314,-174.791){\fontsize{11.991}{1}\usefont{T1}{cmr}{m}{n}\selectfont\color{color_29791}sk}
\put(48.213,-207.389){\fontsize{15.987}{1}\usefont{T1}{cmr}{m}{n}\selectfont\color{color_29791}•}
\put(62.187,-208.409){\fontsize{15.987}{1}\usefont{T1}{cmr}{m}{n}\selectfont\color{color_29791}S}
\put(72.86631,-208.409){\fontsize{15.987}{1}\usefont{T1}{cmr}{m}{n}\selectfont\color{color_29791}S}
\put(83.54562,-208.409){\fontsize{15.987}{1}\usefont{T1}{cmr}{m}{n}\selectfont\color{color_29791}H}
\put(95.04028,-208.409){\fontsize{15.987}{1}\usefont{T1}{cmr}{m}{n}\selectfont\color{color_29791} }
\put(99.4367,-208.409){\fontsize{15.987}{1}\usefont{T1}{cmr}{m}{n}\selectfont\color{color_29791}k}
\put(107.5101,-208.409){\fontsize{15.987}{1}\usefont{T1}{cmr}{m}{n}\selectfont\color{color_29791}e}
\put(116.303,-208.409){\fontsize{15.987}{1}\usefont{T1}{cmr}{m}{n}\selectfont\color{color_29791}y}
\put(124.3764,-208.409){\fontsize{15.987}{1}\usefont{T1}{cmr}{m}{n}\selectfont\color{color_29791}s}
\put(84.213,-237.209){\fontsize{11.991}{1}\usefont{T1}{cmr}{m}{n}\selectfont\color{color_29791}•}
\put(98.187,-237.89){\fontsize{11.991}{1}\usefont{T1}{cmr}{m}{n}\selectfont\color{color_29791}C}
\put(106.8685,-237.89){\fontsize{11.991}{1}\usefont{T1}{cmr}{m}{n}\selectfont\color{color_29791}o}
\put(113.5475,-237.89){\fontsize{11.991}{1}\usefont{T1}{cmr}{m}{n}\selectfont\color{color_29791}n}
\put(120.2265,-237.89){\fontsize{11.991}{1}\usefont{T1}{cmr}{m}{n}\selectfont\color{color_29791}n}
\put(126.9054,-237.89){\fontsize{11.991}{1}\usefont{T1}{cmr}{m}{n}\selectfont\color{color_29791}e}
\put(133.5844,-237.89){\fontsize{11.991}{1}\usefont{T1}{cmr}{m}{n}\selectfont\color{color_29791}ct}
\put(142.9614,-237.89){\fontsize{11.991}{1}\usefont{T1}{cmr}{m}{n}\selectfont\color{color_29791} }
\put(146.3429,-237.89){\fontsize{11.991}{1}\usefont{T1}{cmr}{m}{n}\selectfont\color{color_29791}t}
\put(149.6404,-237.89){\fontsize{11.991}{1}\usefont{T1}{cmr}{m}{n}\selectfont\color{color_29791}o}
\put(156.3194,-237.89){\fontsize{11.991}{1}\usefont{T1}{cmr}{m}{n}\selectfont\color{color_29791} }
\put(159.7008,-237.89){\fontsize{11.991}{1}\usefont{T1}{cmr}{m}{n}\selectfont\color{color_29791}BA}
\put(175.7688,-237.89){\fontsize{11.991}{1}\usefont{T1}{cmr}{m}{n}\selectfont\color{color_29791}S }
\put(187.1362,-237.89){\fontsize{11.991}{1}\usefont{T1}{cmr}{m}{n}\selectfont\color{color_29791}syst}
\put(208.4203,-237.89){\fontsize{11.991}{1}\usefont{T1}{cmr}{m}{n}\selectfont\color{color_29791}e}
\put(215.2072,-237.89){\fontsize{11.991}{1}\usefont{T1}{cmr}{m}{n}\selectfont\color{color_29791}m}
\put(225.0997,-237.89){\fontsize{11.991}{1}\usefont{T1}{cmr}{m}{n}\selectfont\color{color_29791}s }
\put(234.4767,-237.89){\fontsize{11.991}{1}\usefont{T1}{cmr}{m}{n}\selectfont\color{color_29791}w}
\put(243.0742,-237.89){\fontsize{11.991}{1}\usefont{T1}{cmr}{m}{n}\selectfont\color{color_29791}i}
\put(245.6763,-237.89){\fontsize{11.991}{1}\usefont{T1}{cmr}{m}{n}\selectfont\color{color_29791}t}
\put(249.0578,-237.89){\fontsize{11.991}{1}\usefont{T1}{cmr}{m}{n}\selectfont\color{color_29791}h}
\put(255.7368,-237.89){\fontsize{11.991}{1}\usefont{T1}{cmr}{m}{n}\selectfont\color{color_29791}o}
\put(262.5237,-237.89){\fontsize{11.991}{1}\usefont{T1}{cmr}{m}{n}\selectfont\color{color_29791}u}
\put(269.2026,-237.89){\fontsize{11.991}{1}\usefont{T1}{cmr}{m}{n}\selectfont\color{color_29791}t}
\put(272.5002,-237.89){\fontsize{11.991}{1}\usefont{T1}{cmr}{m}{n}\selectfont\color{color_29791} }
\put(275.8816,-237.89){\fontsize{11.991}{1}\usefont{T1}{cmr}{m}{n}\selectfont\color{color_29791}t}
\put(279.2631,-237.89){\fontsize{11.991}{1}\usefont{T1}{cmr}{m}{n}\selectfont\color{color_29791}y}
\put(285.1627,-237.89){\fontsize{11.991}{1}\usefont{T1}{cmr}{m}{n}\selectfont\color{color_29791}p}
\put(291.9496,-237.89){\fontsize{11.991}{1}\usefont{T1}{cmr}{m}{n}\selectfont\color{color_29791}i}
\put(294.5516,-237.89){\fontsize{11.991}{1}\usefont{T1}{cmr}{m}{n}\selectfont\color{color_29791}n}
\put(301.2306,-237.89){\fontsize{11.991}{1}\usefont{T1}{cmr}{m}{n}\selectfont\color{color_29791}g}
\put(308.0175,-237.89){\fontsize{11.991}{1}\usefont{T1}{cmr}{m}{n}\selectfont\color{color_29791} }
\put(311.315,-237.89){\fontsize{11.991}{1}\usefont{T1}{cmr}{m}{n}\selectfont\color{color_29791}p}
\put(317.994,-237.89){\fontsize{11.991}{1}\usefont{T1}{cmr}{m}{n}\selectfont\color{color_29791}a}
\put(324.673,-237.89){\fontsize{11.991}{1}\usefont{T1}{cmr}{m}{n}\selectfont\color{color_29791}s}
\put(330.7524,-237.89){\fontsize{11.991}{1}\usefont{T1}{cmr}{m}{n}\selectfont\color{color_29791}sw}
\put(345.3455,-237.89){\fontsize{11.991}{1}\usefont{T1}{cmr}{m}{n}\selectfont\color{color_29791}o}
\put(352.0244,-237.89){\fontsize{11.991}{1}\usefont{T1}{cmr}{m}{n}\selectfont\color{color_29791}rd}
\put(362.6964,-237.89){\fontsize{11.991}{1}\usefont{T1}{cmr}{m}{n}\selectfont\color{color_29791}s}
\put(84.213,-266.009){\fontsize{11.991}{1}\usefont{T1}{cmr}{m}{n}\selectfont\color{color_29791}•}
\put(98.187,-266.69){\fontsize{11.991}{1}\usefont{T1}{cmr}{m}{n}\selectfont\color{color_29791}ssh}
\put(116.857,-266.69){\fontsize{11.991}{1}\usefont{T1}{cmr}{m}{n}\selectfont\color{color_29791}-ke}
\put(133.5245,-266.69){\fontsize{11.991}{1}\usefont{T1}{cmr}{m}{n}\selectfont\color{color_29791}yg}
\put(146.199,-266.69){\fontsize{11.991}{1}\usefont{T1}{cmr}{m}{n}\selectfont\color{color_29791}e}
\put(152.9859,-266.69){\fontsize{11.991}{1}\usefont{T1}{cmr}{m}{n}\selectfont\color{color_29791}n}
\put(159.6649,-266.69){\fontsize{11.991}{1}\usefont{T1}{cmr}{m}{n}\selectfont\color{color_29791} }
\put(162.9624,-266.69){\fontsize{11.991}{1}\usefont{T1}{cmr}{m}{n}\selectfont\color{color_29791}–}
\put(169.6414,-266.69){\fontsize{11.991}{1}\usefont{T1}{cmr}{m}{n}\selectfont\color{color_29791} }
\put(172.4353,-266.69){\fontsize{11.991}{1}\usefont{T1}{cmr}{m}{n}\selectfont\color{color_29791}A}
\put(180.3253,-266.69){\fontsize{11.991}{1}\usefont{T1}{cmr}{m}{n}\selectfont\color{color_29791}l}
\put(183.0113,-266.69){\fontsize{11.991}{1}\usefont{T1}{cmr}{m}{n}\selectfont\color{color_29791}w}
\put(191.6089,-266.69){\fontsize{11.991}{1}\usefont{T1}{cmr}{m}{n}\selectfont\color{color_29791}a}
\put(198.2879,-266.69){\fontsize{11.991}{1}\usefont{T1}{cmr}{m}{n}\selectfont\color{color_29791}ys }
\put(213.6603,-266.69){\fontsize{11.991}{1}\usefont{T1}{cmr}{m}{n}\selectfont\color{color_29791}cr}
\put(223.5529,-266.69){\fontsize{11.991}{1}\usefont{T1}{cmr}{m}{n}\selectfont\color{color_29791}e}
\put(230.3398,-266.69){\fontsize{11.991}{1}\usefont{T1}{cmr}{m}{n}\selectfont\color{color_29791}a}
\put(237.0188,-266.69){\fontsize{11.991}{1}\usefont{T1}{cmr}{m}{n}\selectfont\color{color_29791}t}
\put(240.4003,-266.69){\fontsize{11.991}{1}\usefont{T1}{cmr}{m}{n}\selectfont\color{color_29791}e}
\put(247.0793,-266.69){\fontsize{11.991}{1}\usefont{T1}{cmr}{m}{n}\selectfont\color{color_29791} }
\put(250.3768,-266.69){\fontsize{11.991}{1}\usefont{T1}{cmr}{m}{n}\selectfont\color{color_29791}w}
\put(259.0583,-266.69){\fontsize{11.991}{1}\usefont{T1}{cmr}{m}{n}\selectfont\color{color_29791}i}
\put(261.6603,-266.69){\fontsize{11.991}{1}\usefont{T1}{cmr}{m}{n}\selectfont\color{color_29791}t}
\put(265.0418,-266.69){\fontsize{11.991}{1}\usefont{T1}{cmr}{m}{n}\selectfont\color{color_29791}h}
\put(271.7208,-266.69){\fontsize{11.991}{1}\usefont{T1}{cmr}{m}{n}\selectfont\color{color_29791} }
\put(275.1022,-266.69){\fontsize{11.991}{1}\usefont{T1}{cmr}{m}{n}\selectfont\color{color_29791}a}
\put(281.7812,-266.69){\fontsize{11.991}{1}\usefont{T1}{cmr}{m}{n}\selectfont\color{color_29791} }
\put(285.0787,-266.69){\fontsize{11.991}{1}\usefont{T1}{cmr}{m}{n}\selectfont\color{color_29791}p}
\put(291.7577,-266.69){\fontsize{11.991}{1}\usefont{T1}{cmr}{m}{n}\selectfont\color{color_29791}a}
\put(298.5446,-266.69){\fontsize{11.991}{1}\usefont{T1}{cmr}{m}{n}\selectfont\color{color_29791}ssp}
\put(317.2146,-266.69){\fontsize{11.991}{1}\usefont{T1}{cmr}{m}{n}\selectfont\color{color_29791}h}
\put(323.8936,-266.69){\fontsize{11.991}{1}\usefont{T1}{cmr}{m}{n}\selectfont\color{color_29791}ra}
\put(334.5656,-266.69){\fontsize{11.991}{1}\usefont{T1}{cmr}{m}{n}\selectfont\color{color_29791}se}
\put(84.213,-294.809){\fontsize{11.991}{1}\usefont{T1}{cmr}{m}{n}\selectfont\color{color_29791}•}
\put(98.187,-295.49){\fontsize{11.991}{1}\usefont{T1}{cmr}{m}{n}\selectfont\color{color_29791}ssh}
\put(116.857,-295.49){\fontsize{11.991}{1}\usefont{T1}{cmr}{m}{n}\selectfont\color{color_29791}-a}
\put(127.529,-295.49){\fontsize{11.991}{1}\usefont{T1}{cmr}{m}{n}\selectfont\color{color_29791}g}
\put(134.208,-295.49){\fontsize{11.991}{1}\usefont{T1}{cmr}{m}{n}\selectfont\color{color_29791}e}
\put(140.8869,-295.49){\fontsize{11.991}{1}\usefont{T1}{cmr}{m}{n}\selectfont\color{color_29791}n}
\put(147.6739,-295.49){\fontsize{11.991}{1}\usefont{T1}{cmr}{m}{n}\selectfont\color{color_29791}t}
\put(84.213,-323.609){\fontsize{11.991}{1}\usefont{T1}{cmr}{m}{n}\selectfont\color{color_29791}•}
\put(98.187,-324.205){\fontsize{11.991}{1}\usefont{T1}{cmr}{m}{n}\selectfont\color{color_29791}.}
\put(101.5685,-324.205){\fontsize{11.991}{1}\usefont{T1}{cmr}{m}{n}\selectfont\color{color_29791}/}
\put(104.866,-324.205){\fontsize{11.991}{1}\usefont{T1}{cmr}{m}{n}\selectfont\color{color_29791}ssh}
\put(123.536,-324.205){\fontsize{11.991}{1}\usefont{T1}{cmr}{m}{n}\selectfont\color{color_29791}/}
\put(126.9174,-324.205){\fontsize{11.991}{1}\usefont{T1}{cmr}{m}{n}\selectfont\color{color_29791}co}
\put(139.5919,-324.205){\fontsize{11.991}{1}\usefont{T1}{cmr}{m}{n}\selectfont\color{color_29791}n}
\put(146.3788,-324.205){\fontsize{11.991}{1}\usefont{T1}{cmr}{m}{n}\selectfont\color{color_29791}f}
\put(149.6763,-324.205){\fontsize{11.991}{1}\usefont{T1}{cmr}{m}{n}\selectfont\color{color_29791}i}
\put(152.2784,-324.205){\fontsize{11.991}{1}\usefont{T1}{cmr}{m}{n}\selectfont\color{color_29791}g}
\put(84.213,-352.296){\fontsize{11.991}{1}\usefont{T1}{cmr}{m}{n}\selectfont\color{color_29791}•}
\put(98.187,-353.005){\fontsize{11.991}{1}\usefont{T1}{cmr}{m}{n}\selectfont\color{color_29791}Mo}
\put(114.8545,-353.005){\fontsize{11.991}{1}\usefont{T1}{cmr}{m}{n}\selectfont\color{color_29791}re}
\put(125.5265,-353.005){\fontsize{11.991}{1}\usefont{T1}{cmr}{m}{n}\selectfont\color{color_29791} }
\put(128.824,-353.005){\fontsize{11.991}{1}\usefont{T1}{cmr}{m}{n}\selectfont\color{color_29791}i}
\put(131.51,-353.005){\fontsize{11.991}{1}\usefont{T1}{cmr}{m}{n}\selectfont\color{color_29791}n}
\put(138.189,-353.005){\fontsize{11.991}{1}\usefont{T1}{cmr}{m}{n}\selectfont\color{color_29791}f}
\put(141.4865,-353.005){\fontsize{11.991}{1}\usefont{T1}{cmr}{m}{n}\selectfont\color{color_29791}o}
\put(148.2734,-353.005){\fontsize{11.991}{1}\usefont{T1}{cmr}{m}{n}\selectfont\color{color_29791}r}
\put(152.1705,-353.005){\fontsize{11.991}{1}\usefont{T1}{cmr}{m}{n}\selectfont\color{color_29791}ma}
\put(168.838,-353.005){\fontsize{11.991}{1}\usefont{T1}{cmr}{m}{n}\selectfont\color{color_29791}t}
\put(172.2195,-353.005){\fontsize{11.991}{1}\usefont{T1}{cmr}{m}{n}\selectfont\color{color_29791}i}
\put(174.8215,-353.005){\fontsize{11.991}{1}\usefont{T1}{cmr}{m}{n}\selectfont\color{color_29791}o}
\put(181.5005,-353.005){\fontsize{11.991}{1}\usefont{T1}{cmr}{m}{n}\selectfont\color{color_29791}n}
\put(188.1795,-353.005){\fontsize{11.991}{1}\usefont{T1}{cmr}{m}{n}\selectfont\color{color_29791}:}
\put(191.561,-353.005){\fontsize{11.991}{1}\usefont{T1}{cmr}{m}{n}\selectfont\color{color_29791} }
\put(195.189,-353.005){\fontsize{11.991}{1}\usefont{T1}{cmr}{m}{n}\selectfont\color{color_167934}h}
\put(201.3883,-353.005){\fontsize{11.991}{1}\usefont{T1}{cmr}{m}{n}\selectfont\color{color_167934}tt}
\put(208.9786,-353.005){\fontsize{11.991}{1}\usefont{T1}{cmr}{m}{n}\selectfont\color{color_167934}p:}
\put(218.5594,-353.005){\fontsize{11.991}{1}\usefont{T1}{cmr}{m}{n}\selectfont\color{color_167934}/}
\put(223.2479,-353.005){\fontsize{11.991}{1}\usefont{T1}{cmr}{m}{n}\selectfont\color{color_167934}/}
\put(227.9364,-353.005){\fontsize{11.991}{1}\usefont{T1}{cmr}{m}{n}\selectfont\color{color_167934}i}
\put(230.6344,-353.005){\fontsize{11.991}{1}\usefont{T1}{cmr}{m}{n}\selectfont\color{color_167934}c}
\put(235.7186,-353.005){\fontsize{11.991}{1}\usefont{T1}{cmr}{m}{n}\selectfont\color{color_167934}t}
\put(239.5197,-353.005){\fontsize{11.991}{1}\usefont{T1}{cmr}{m}{n}\selectfont\color{color_167934}d}
\put(245.8989,-353.005){\fontsize{11.991}{1}\usefont{T1}{cmr}{m}{n}\selectfont\color{color_167934}o}
\put(252.1942,-353.005){\fontsize{11.991}{1}\usefont{T1}{cmr}{m}{n}\selectfont\color{color_167934}c}
\put(257.2784,-353.005){\fontsize{11.991}{1}\usefont{T1}{cmr}{m}{n}\selectfont\color{color_167934}s}
\put(261.871,-353.005){\fontsize{11.991}{1}\usefont{T1}{cmr}{m}{n}\selectfont\color{color_167934}/}
\put(266.5594,-353.005){\fontsize{11.991}{1}\usefont{T1}{cmr}{m}{n}\selectfont\color{color_167934}w}
\put(275.0491,-353.005){\fontsize{11.991}{1}\usefont{T1}{cmr}{m}{n}\selectfont\color{color_167934}i}
\put(277.831,-353.005){\fontsize{11.991}{1}\usefont{T1}{cmr}{m}{n}\selectfont\color{color_167934}k}
\put(283.2269,-353.005){\fontsize{11.991}{1}\usefont{T1}{cmr}{m}{n}\selectfont\color{color_167934}i}
\put(286.0089,-353.005){\fontsize{11.991}{1}\usefont{T1}{cmr}{m}{n}\selectfont\color{color_167934}/}
\put(290.6014,-353.005){\fontsize{11.991}{1}\usefont{T1}{cmr}{m}{n}\selectfont\color{color_167934}i}
\put(293.3833,-353.005){\fontsize{11.991}{1}\usefont{T1}{cmr}{m}{n}\selectfont\color{color_167934}n}
\put(299.7625,-353.005){\fontsize{11.991}{1}\usefont{T1}{cmr}{m}{n}\selectfont\color{color_167934}de}
\put(311.8495,-353.005){\fontsize{11.991}{1}\usefont{T1}{cmr}{m}{n}\selectfont\color{color_167934}x}
\put(316.9456,-353.005){\fontsize{11.991}{1}\usefont{T1}{cmr}{m}{n}\selectfont\color{color_167934}.}
\put(319.9434,-353.005){\fontsize{11.991}{1}\usefont{T1}{cmr}{m}{n}\selectfont\color{color_167934}p}
\put(326.3226,-353.005){\fontsize{11.991}{1}\usefont{T1}{cmr}{m}{n}\selectfont\color{color_167934}h}
\put(332.7018,-353.005){\fontsize{11.991}{1}\usefont{T1}{cmr}{m}{n}\selectfont\color{color_167934}p/}
\put(343.6855,-353.005){\fontsize{11.991}{1}\usefont{T1}{cmr}{m}{n}\selectfont\color{color_167934}S}
\put(349.1774,-353.005){\fontsize{11.991}{1}\usefont{T1}{cmr}{m}{n}\selectfont\color{color_167934}S}
\put(354.6693,-353.005){\fontsize{11.991}{1}\usefont{T1}{cmr}{m}{n}\selectfont\color{color_167934}H}
\put(362.1517,-353.005){\fontsize{11.991}{1}\usefont{T1}{cmr}{m}{n}\selectfont\color{color_167934}\_}
\put(368.1352,-353.005){\fontsize{11.991}{1}\usefont{T1}{cmr}{m}{n}\selectfont\color{color_167934}K}
\put(374.1307,-353.005){\fontsize{11.991}{1}\usefont{T1}{cmr}{m}{n}\selectfont\color{color_167934}e}
\put(379.9224,-353.005){\fontsize{11.991}{1}\usefont{T1}{cmr}{m}{n}\selectfont\color{color_167934}y}
\put(385.2224,-353.005){\fontsize{11.991}{1}\usefont{T1}{cmr}{m}{n}\selectfont\color{color_167934}s}
\end{picture}
\begin{tikzpicture}[overlay]
\path(0pt,0pt);
\draw[color_167934,line width=0.794pt]
(195.189pt, -354.196pt) -- (390.298pt, -354.196pt)
;
\end{tikzpicture}
\begin{picture}(-5,0)(2.5,0)
\put(84.213,-381.096){\fontsize{11.991}{1}\usefont{T1}{cmr}{m}{n}\selectfont\color{color_29791}•}
\put(98.187,-381.805){\fontsize{11.991}{1}\usefont{T1}{cmr}{m}{n}\selectfont\color{color_29791}D}
\put(106.8685,-381.805){\fontsize{11.991}{1}\usefont{T1}{cmr}{m}{n}\selectfont\color{color_29791}e}
\put(113.5475,-381.805){\fontsize{11.991}{1}\usefont{T1}{cmr}{m}{n}\selectfont\color{color_29791}m}
\put(123.44,-381.805){\fontsize{11.991}{1}\usefont{T1}{cmr}{m}{n}\selectfont\color{color_29791}o}
\put(130.119,-381.805){\fontsize{11.991}{1}\usefont{T1}{cmr}{m}{n}\selectfont\color{color_29791}n}
\put(136.9059,-381.805){\fontsize{11.991}{1}\usefont{T1}{cmr}{m}{n}\selectfont\color{color_29791}st}
\put(146.199,-381.805){\fontsize{11.991}{1}\usefont{T1}{cmr}{m}{n}\selectfont\color{color_29791}ra}
\put(156.871,-381.805){\fontsize{11.991}{1}\usefont{T1}{cmr}{m}{n}\selectfont\color{color_29791}t}
\put(160.2524,-381.805){\fontsize{11.991}{1}\usefont{T1}{cmr}{m}{n}\selectfont\color{color_29791}i}
\put(162.8545,-381.805){\fontsize{11.991}{1}\usefont{T1}{cmr}{m}{n}\selectfont\color{color_29791}o}
\put(169.5335,-381.805){\fontsize{11.991}{1}\usefont{T1}{cmr}{m}{n}\selectfont\color{color_29791}n}
\end{picture}
\newpage
\begin{tikzpicture}[overlay]
\path(0pt,0pt);
\filldraw[color_283006][even odd rule]
(-15pt, 10pt) -- (704.972pt, 10pt)
 -- (704.972pt, 10pt)
 -- (704.972pt, -529.972pt)
 -- (704.972pt, -529.972pt)
 -- (-15pt, -529.972pt)
 -- (-15pt, -529.972pt)
 -- (-15pt, 10pt) -- cycle
;
\end{tikzpicture}
\begin{picture}(-5,0)(2.5,0)
\put(-15,-529.972){\includegraphics[width=720pt,height=540pt]{latexImage_1aef5d9a345e7259c360833b5fc6008f.png}}
\put(48.411,-60.69601){\fontsize{21.997}{1}\usefont{T1}{cmr}{b}{n}\selectfont\color{color_29791}Us}
\put(76.50117,-60.69601){\fontsize{21.997}{1}\usefont{T1}{cmr}{b}{n}\selectfont\color{color_29791}e}
\put(88.7095,-60.69601){\fontsize{21.997}{1}\usefont{T1}{cmr}{b}{n}\selectfont\color{color_29791}r}
\put(97.28833,-60.69601){\fontsize{21.997}{1}\usefont{T1}{cmr}{b}{n}\selectfont\color{color_29791} E}
\put(117.9875,-60.69601){\fontsize{21.997}{1}\usefont{T1}{cmr}{b}{n}\selectfont\color{color_29791}n}
\put(131.4717,-60.69601){\fontsize{21.997}{1}\usefont{T1}{cmr}{b}{n}\selectfont\color{color_29791}v}
\put(143.68,-60.69601){\fontsize{21.997}{1}\usefont{T1}{cmr}{b}{n}\selectfont\color{color_29791}ir}
\put(158.352,-60.69601){\fontsize{21.997}{1}\usefont{T1}{cmr}{b}{n}\selectfont\color{color_29791}o}
\put(171.7482,-60.69601){\fontsize{21.997}{1}\usefont{T1}{cmr}{b}{n}\selectfont\color{color_29791}n}
\put(185.1443,-60.69601){\fontsize{21.997}{1}\usefont{T1}{cmr}{b}{n}\selectfont\color{color_29791}m}
\put(204.7217,-60.69601){\fontsize{21.997}{1}\usefont{T1}{cmr}{b}{n}\selectfont\color{color_29791}e}
\put(216.93,-60.69601){\fontsize{21.997}{1}\usefont{T1}{cmr}{b}{n}\selectfont\color{color_29791}n}
\put(230.4142,-60.69601){\fontsize{21.997}{1}\usefont{T1}{cmr}{b}{n}\selectfont\color{color_29791}t}
\put(48.213,-115.49){\fontsize{15.987}{1}\usefont{T1}{cmr}{m}{n}\selectfont\color{color_29791}•}
\put(62.187,-116.51){\fontsize{15.987}{1}\usefont{T1}{cmr}{m}{n}\selectfont\color{color_29791}t}
\put(66.66335,-116.51){\fontsize{15.987}{1}\usefont{T1}{cmr}{m}{n}\selectfont\color{color_29791}m}
\put(79.94855,-116.51){\fontsize{15.987}{1}\usefont{T1}{cmr}{m}{n}\selectfont\color{color_29791}ux}
\put(84.213,-145.31){\fontsize{11.991}{1}\usefont{T1}{cmr}{m}{n}\selectfont\color{color_29791}•}
\put(98.187,-145.991){\fontsize{11.991}{1}\usefont{T1}{cmr}{m}{n}\selectfont\color{color_29791}K}
\put(106.2689,-145.991){\fontsize{11.991}{1}\usefont{T1}{cmr}{m}{n}\selectfont\color{color_29791}e}
\put(112.9479,-145.991){\fontsize{11.991}{1}\usefont{T1}{cmr}{m}{n}\selectfont\color{color_29791}e}
\put(119.6269,-145.991){\fontsize{11.991}{1}\usefont{T1}{cmr}{m}{n}\selectfont\color{color_29791}p}
\put(126.3059,-145.991){\fontsize{11.991}{1}\usefont{T1}{cmr}{m}{n}\selectfont\color{color_29791}s }
\put(135.6829,-145.991){\fontsize{11.991}{1}\usefont{T1}{cmr}{m}{n}\selectfont\color{color_29791}l}
\put(138.2849,-145.991){\fontsize{11.991}{1}\usefont{T1}{cmr}{m}{n}\selectfont\color{color_29791}o}
\put(144.9639,-145.991){\fontsize{11.991}{1}\usefont{T1}{cmr}{m}{n}\selectfont\color{color_29791}n}
\put(151.6429,-145.991){\fontsize{11.991}{1}\usefont{T1}{cmr}{m}{n}\selectfont\color{color_29791}g}
\put(158.4298,-145.991){\fontsize{11.991}{1}\usefont{T1}{cmr}{m}{n}\selectfont\color{color_29791} }
\put(161.7273,-145.991){\fontsize{11.991}{1}\usefont{T1}{cmr}{m}{n}\selectfont\color{color_29791}ru}
\put(172.3993,-145.991){\fontsize{11.991}{1}\usefont{T1}{cmr}{m}{n}\selectfont\color{color_29791}n}
\put(179.0783,-145.991){\fontsize{11.991}{1}\usefont{T1}{cmr}{m}{n}\selectfont\color{color_29791}n}
\put(185.8652,-145.991){\fontsize{11.991}{1}\usefont{T1}{cmr}{m}{n}\selectfont\color{color_29791}i}
\put(188.4673,-145.991){\fontsize{11.991}{1}\usefont{T1}{cmr}{m}{n}\selectfont\color{color_29791}n}
\put(195.1463,-145.991){\fontsize{11.991}{1}\usefont{T1}{cmr}{m}{n}\selectfont\color{color_29791}g}
\put(201.8252,-145.991){\fontsize{11.991}{1}\usefont{T1}{cmr}{m}{n}\selectfont\color{color_29791} }
\put(205.2067,-145.991){\fontsize{11.991}{1}\usefont{T1}{cmr}{m}{n}\selectfont\color{color_29791}co}
\put(217.8812,-145.991){\fontsize{11.991}{1}\usefont{T1}{cmr}{m}{n}\selectfont\color{color_29791}mm}
\put(237.7623,-145.991){\fontsize{11.991}{1}\usefont{T1}{cmr}{m}{n}\selectfont\color{color_29791}a}
\put(244.5492,-145.991){\fontsize{11.991}{1}\usefont{T1}{cmr}{m}{n}\selectfont\color{color_29791}n}
\put(251.2281,-145.991){\fontsize{11.991}{1}\usefont{T1}{cmr}{m}{n}\selectfont\color{color_29791}d}
\put(257.9071,-145.991){\fontsize{11.991}{1}\usefont{T1}{cmr}{m}{n}\selectfont\color{color_29791} }
\put(261.2886,-145.991){\fontsize{11.991}{1}\usefont{T1}{cmr}{m}{n}\selectfont\color{color_29791}l}
\put(263.8907,-145.991){\fontsize{11.991}{1}\usefont{T1}{cmr}{m}{n}\selectfont\color{color_29791}i}
\put(266.4927,-145.991){\fontsize{11.991}{1}\usefont{T1}{cmr}{m}{n}\selectfont\color{color_29791}n}
\put(273.2796,-145.991){\fontsize{11.991}{1}\usefont{T1}{cmr}{m}{n}\selectfont\color{color_29791}e}
\put(280.398,-145.991){\fontsize{11.991}{1}\usefont{T1}{cmr}{m}{n}\selectfont\color{color_29791} }
\put(283.687,-145.991){\fontsize{11.991}{1}\usefont{T1}{cmr}{m}{n}\selectfont\color{color_29791}se}
\put(296.3615,-145.991){\fontsize{11.991}{1}\usefont{T1}{cmr}{m}{n}\selectfont\color{color_29791}s}
\put(302.4409,-145.991){\fontsize{11.991}{1}\usefont{T1}{cmr}{m}{n}\selectfont\color{color_29791}si}
\put(311.0385,-145.991){\fontsize{11.991}{1}\usefont{T1}{cmr}{m}{n}\selectfont\color{color_29791}o}
\put(317.7175,-145.991){\fontsize{11.991}{1}\usefont{T1}{cmr}{m}{n}\selectfont\color{color_29791}n}
\put(324.3965,-145.991){\fontsize{11.991}{1}\usefont{T1}{cmr}{m}{n}\selectfont\color{color_29791}s }
\put(333.7734,-145.991){\fontsize{11.991}{1}\usefont{T1}{cmr}{m}{n}\selectfont\color{color_29791}ru}
\put(344.4454,-145.991){\fontsize{11.991}{1}\usefont{T1}{cmr}{m}{n}\selectfont\color{color_29791}n}
\put(351.1244,-145.991){\fontsize{11.991}{1}\usefont{T1}{cmr}{m}{n}\selectfont\color{color_29791}n}
\put(357.8034,-145.991){\fontsize{11.991}{1}\usefont{T1}{cmr}{m}{n}\selectfont\color{color_29791}i}
\put(360.4054,-145.991){\fontsize{11.991}{1}\usefont{T1}{cmr}{m}{n}\selectfont\color{color_29791}n}
\put(367.1923,-145.991){\fontsize{11.991}{1}\usefont{T1}{cmr}{m}{n}\selectfont\color{color_29791}g}
\put(374.112,-145.991){\fontsize{11.991}{1}\usefont{T1}{cmr}{m}{n}\selectfont\color{color_29791} }
\put(84.213,-174.11){\fontsize{11.991}{1}\usefont{T1}{cmr}{m}{n}\selectfont\color{color_29791}•}
\put(98.187,-174.791){\fontsize{11.991}{1}\usefont{T1}{cmr}{m}{n}\selectfont\color{color_29791}A}
\put(106.2689,-174.791){\fontsize{11.991}{1}\usefont{T1}{cmr}{m}{n}\selectfont\color{color_29791}l}
\put(108.871,-174.791){\fontsize{11.991}{1}\usefont{T1}{cmr}{m}{n}\selectfont\color{color_29791}l}
\put(111.473,-174.791){\fontsize{11.991}{1}\usefont{T1}{cmr}{m}{n}\selectfont\color{color_29791}o}
\put(118.2599,-174.791){\fontsize{11.991}{1}\usefont{T1}{cmr}{m}{n}\selectfont\color{color_29791}w}
\put(126.8575,-174.791){\fontsize{11.991}{1}\usefont{T1}{cmr}{m}{n}\selectfont\color{color_29791}s }
\put(136.2345,-174.791){\fontsize{11.991}{1}\usefont{T1}{cmr}{m}{n}\selectfont\color{color_29791}d}
\put(142.9134,-174.791){\fontsize{11.991}{1}\usefont{T1}{cmr}{m}{n}\selectfont\color{color_29791}i}
\put(145.5155,-174.791){\fontsize{11.991}{1}\usefont{T1}{cmr}{m}{n}\selectfont\color{color_29791}sco}
\put(164.1855,-174.791){\fontsize{11.991}{1}\usefont{T1}{cmr}{m}{n}\selectfont\color{color_29791}n}
\put(170.9724,-174.791){\fontsize{11.991}{1}\usefont{T1}{cmr}{m}{n}\selectfont\color{color_29791}n}
\put(177.6514,-174.791){\fontsize{11.991}{1}\usefont{T1}{cmr}{m}{n}\selectfont\color{color_29791}e}
\put(184.3304,-174.791){\fontsize{11.991}{1}\usefont{T1}{cmr}{m}{n}\selectfont\color{color_29791}ct}
\put(193.7073,-174.791){\fontsize{11.991}{1}\usefont{T1}{cmr}{m}{n}\selectfont\color{color_29791}i}
\put(196.3094,-174.791){\fontsize{11.991}{1}\usefont{T1}{cmr}{m}{n}\selectfont\color{color_29791}n}
\put(202.9884,-174.791){\fontsize{11.991}{1}\usefont{T1}{cmr}{m}{n}\selectfont\color{color_29791}g}
\put(209.6673,-174.791){\fontsize{11.991}{1}\usefont{T1}{cmr}{m}{n}\selectfont\color{color_29791} }
\put(213.0488,-174.791){\fontsize{11.991}{1}\usefont{T1}{cmr}{m}{n}\selectfont\color{color_29791}a}
\put(219.7278,-174.791){\fontsize{11.991}{1}\usefont{T1}{cmr}{m}{n}\selectfont\color{color_29791}n}
\put(226.4068,-174.791){\fontsize{11.991}{1}\usefont{T1}{cmr}{m}{n}\selectfont\color{color_29791}d}
\put(233.1937,-174.791){\fontsize{11.991}{1}\usefont{T1}{cmr}{m}{n}\selectfont\color{color_29791} }
\put(236.4912,-174.791){\fontsize{11.991}{1}\usefont{T1}{cmr}{m}{n}\selectfont\color{color_29791}re}
\put(247.1632,-174.791){\fontsize{11.991}{1}\usefont{T1}{cmr}{m}{n}\selectfont\color{color_29791}co}
\put(259.8377,-174.791){\fontsize{11.991}{1}\usefont{T1}{cmr}{m}{n}\selectfont\color{color_29791}n}
\put(266.5167,-174.791){\fontsize{11.991}{1}\usefont{T1}{cmr}{m}{n}\selectfont\color{color_29791}n}
\put(273.3036,-174.791){\fontsize{11.991}{1}\usefont{T1}{cmr}{m}{n}\selectfont\color{color_29791}e}
\put(279.9826,-174.791){\fontsize{11.991}{1}\usefont{T1}{cmr}{m}{n}\selectfont\color{color_29791}ct}
\put(289.2756,-174.791){\fontsize{11.991}{1}\usefont{T1}{cmr}{m}{n}\selectfont\color{color_29791}i}
\put(291.9616,-174.791){\fontsize{11.991}{1}\usefont{T1}{cmr}{m}{n}\selectfont\color{color_29791}n}
\put(298.6406,-174.791){\fontsize{11.991}{1}\usefont{T1}{cmr}{m}{n}\selectfont\color{color_29791}g}
\put(84.213,-202.91){\fontsize{11.991}{1}\usefont{T1}{cmr}{m}{n}\selectfont\color{color_29791}•}
\put(98.187,-203.591){\fontsize{11.991}{1}\usefont{T1}{cmr}{m}{n}\selectfont\color{color_29791}Mu}
\put(114.8545,-203.591){\fontsize{11.991}{1}\usefont{T1}{cmr}{m}{n}\selectfont\color{color_29791}l}
\put(117.4565,-203.591){\fontsize{11.991}{1}\usefont{T1}{cmr}{m}{n}\selectfont\color{color_29791}t}
\put(120.838,-203.591){\fontsize{11.991}{1}\usefont{T1}{cmr}{m}{n}\selectfont\color{color_29791}i}
\put(123.4401,-203.591){\fontsize{11.991}{1}\usefont{T1}{cmr}{m}{n}\selectfont\color{color_29791}p}
\put(130.119,-203.591){\fontsize{11.991}{1}\usefont{T1}{cmr}{m}{n}\selectfont\color{color_29791}l}
\put(132.805,-203.591){\fontsize{11.991}{1}\usefont{T1}{cmr}{m}{n}\selectfont\color{color_29791}e}
\put(139.484,-203.591){\fontsize{11.991}{1}\usefont{T1}{cmr}{m}{n}\selectfont\color{color_29791} }
\put(142.8655,-203.591){\fontsize{11.991}{1}\usefont{T1}{cmr}{m}{n}\selectfont\color{color_29791}co}
\put(155.54,-203.591){\fontsize{11.991}{1}\usefont{T1}{cmr}{m}{n}\selectfont\color{color_29791}mm}
\put(175.421,-203.591){\fontsize{11.991}{1}\usefont{T1}{cmr}{m}{n}\selectfont\color{color_29791}a}
\put(182.1,-203.591){\fontsize{11.991}{1}\usefont{T1}{cmr}{m}{n}\selectfont\color{color_29791}n}
\put(188.8869,-203.591){\fontsize{11.991}{1}\usefont{T1}{cmr}{m}{n}\selectfont\color{color_29791}d}
\put(195.5659,-203.591){\fontsize{11.991}{1}\usefont{T1}{cmr}{m}{n}\selectfont\color{color_29791} }
\put(198.8634,-203.591){\fontsize{11.991}{1}\usefont{T1}{cmr}{m}{n}\selectfont\color{color_29791}l}
\put(201.5494,-203.591){\fontsize{11.991}{1}\usefont{T1}{cmr}{m}{n}\selectfont\color{color_29791}i}
\put(204.1515,-203.591){\fontsize{11.991}{1}\usefont{T1}{cmr}{m}{n}\selectfont\color{color_29791}n}
\put(210.8305,-203.591){\fontsize{11.991}{1}\usefont{T1}{cmr}{m}{n}\selectfont\color{color_29791}e}
\put(217.894,-203.591){\fontsize{11.991}{1}\usefont{T1}{cmr}{m}{n}\selectfont\color{color_29791} }
\put(221.211,-203.591){\fontsize{11.991}{1}\usefont{T1}{cmr}{m}{n}\selectfont\color{color_29791}se}
\put(233.8855,-203.591){\fontsize{11.991}{1}\usefont{T1}{cmr}{m}{n}\selectfont\color{color_29791}ssi}
\put(248.5625,-203.591){\fontsize{11.991}{1}\usefont{T1}{cmr}{m}{n}\selectfont\color{color_29791}o}
\put(255.2415,-203.591){\fontsize{11.991}{1}\usefont{T1}{cmr}{m}{n}\selectfont\color{color_29791}n}
\put(261.9204,-203.591){\fontsize{11.991}{1}\usefont{T1}{cmr}{m}{n}\selectfont\color{color_29791}s }
\put(271.2974,-203.591){\fontsize{11.991}{1}\usefont{T1}{cmr}{m}{n}\selectfont\color{color_29791}a}
\put(277.9764,-203.591){\fontsize{11.991}{1}\usefont{T1}{cmr}{m}{n}\selectfont\color{color_29791}n}
\put(284.6554,-203.591){\fontsize{11.991}{1}\usefont{T1}{cmr}{m}{n}\selectfont\color{color_29791}d}
\put(291.3344,-203.591){\fontsize{11.991}{1}\usefont{T1}{cmr}{m}{n}\selectfont\color{color_29791} }
\put(294.7158,-203.591){\fontsize{11.991}{1}\usefont{T1}{cmr}{m}{n}\selectfont\color{color_29791}co}
\put(307.3903,-203.591){\fontsize{11.991}{1}\usefont{T1}{cmr}{m}{n}\selectfont\color{color_29791}n}
\put(314.0693,-203.591){\fontsize{11.991}{1}\usefont{T1}{cmr}{m}{n}\selectfont\color{color_29791}so}
\put(326.8517,-203.591){\fontsize{11.991}{1}\usefont{T1}{cmr}{m}{n}\selectfont\color{color_29791}l}
\put(329.4538,-203.591){\fontsize{11.991}{1}\usefont{T1}{cmr}{m}{n}\selectfont\color{color_29791}e}
\put(336.1328,-203.591){\fontsize{11.991}{1}\usefont{T1}{cmr}{m}{n}\selectfont\color{color_29791} }
\put(339.5142,-203.591){\fontsize{11.991}{1}\usefont{T1}{cmr}{m}{n}\selectfont\color{color_29791}sp}
\put(352.1887,-203.591){\fontsize{11.991}{1}\usefont{T1}{cmr}{m}{n}\selectfont\color{color_29791}l}
\put(354.7908,-203.591){\fontsize{11.991}{1}\usefont{T1}{cmr}{m}{n}\selectfont\color{color_29791}i}
\put(357.4767,-203.591){\fontsize{11.991}{1}\usefont{T1}{cmr}{m}{n}\selectfont\color{color_29791}t}
\put(360.7743,-203.591){\fontsize{11.991}{1}\usefont{T1}{cmr}{m}{n}\selectfont\color{color_29791}t}
\put(364.1557,-203.591){\fontsize{11.991}{1}\usefont{T1}{cmr}{m}{n}\selectfont\color{color_29791}i}
\put(366.7578,-203.591){\fontsize{11.991}{1}\usefont{T1}{cmr}{m}{n}\selectfont\color{color_29791}n}
\put(373.5447,-203.591){\fontsize{11.991}{1}\usefont{T1}{cmr}{m}{n}\selectfont\color{color_29791}g}
\put(84.213,-231.71){\fontsize{11.991}{1}\usefont{T1}{cmr}{m}{n}\selectfont\color{color_29791}•}
\put(98.187,-232.306){\fontsize{11.991}{1}\usefont{T1}{cmr}{m}{n}\selectfont\color{color_29791}Mo}
\put(114.8545,-232.306){\fontsize{11.991}{1}\usefont{T1}{cmr}{m}{n}\selectfont\color{color_29791}re}
\put(125.5265,-232.306){\fontsize{11.991}{1}\usefont{T1}{cmr}{m}{n}\selectfont\color{color_29791} }
\put(128.824,-232.306){\fontsize{11.991}{1}\usefont{T1}{cmr}{m}{n}\selectfont\color{color_29791}i}
\put(131.51,-232.306){\fontsize{11.991}{1}\usefont{T1}{cmr}{m}{n}\selectfont\color{color_29791}n}
\put(138.189,-232.306){\fontsize{11.991}{1}\usefont{T1}{cmr}{m}{n}\selectfont\color{color_29791}f}
\put(141.4865,-232.306){\fontsize{11.991}{1}\usefont{T1}{cmr}{m}{n}\selectfont\color{color_29791}o}
\put(148.2734,-232.306){\fontsize{11.991}{1}\usefont{T1}{cmr}{m}{n}\selectfont\color{color_29791}r}
\put(152.1705,-232.306){\fontsize{11.991}{1}\usefont{T1}{cmr}{m}{n}\selectfont\color{color_29791}ma}
\put(168.838,-232.306){\fontsize{11.991}{1}\usefont{T1}{cmr}{m}{n}\selectfont\color{color_29791}t}
\put(172.2195,-232.306){\fontsize{11.991}{1}\usefont{T1}{cmr}{m}{n}\selectfont\color{color_29791}i}
\put(174.8215,-232.306){\fontsize{11.991}{1}\usefont{T1}{cmr}{m}{n}\selectfont\color{color_29791}o}
\put(181.5005,-232.306){\fontsize{11.991}{1}\usefont{T1}{cmr}{m}{n}\selectfont\color{color_29791}n}
\put(188.1795,-232.306){\fontsize{11.991}{1}\usefont{T1}{cmr}{m}{n}\selectfont\color{color_29791}:}
\put(191.561,-232.306){\fontsize{11.991}{1}\usefont{T1}{cmr}{m}{n}\selectfont\color{color_29791} }
\put(195.189,-232.306){\fontsize{11.991}{1}\usefont{T1}{cmr}{m}{n}\selectfont\color{color_167934}h}
\put(201.3883,-232.306){\fontsize{11.991}{1}\usefont{T1}{cmr}{m}{n}\selectfont\color{color_167934}tt}
\put(208.9786,-232.306){\fontsize{11.991}{1}\usefont{T1}{cmr}{m}{n}\selectfont\color{color_167934}p:}
\put(218.5594,-232.306){\fontsize{11.991}{1}\usefont{T1}{cmr}{m}{n}\selectfont\color{color_167934}/}
\put(223.2479,-232.306){\fontsize{11.991}{1}\usefont{T1}{cmr}{m}{n}\selectfont\color{color_167934}/}
\put(227.9364,-232.306){\fontsize{11.991}{1}\usefont{T1}{cmr}{m}{n}\selectfont\color{color_167934}i}
\put(230.6344,-232.306){\fontsize{11.991}{1}\usefont{T1}{cmr}{m}{n}\selectfont\color{color_167934}c}
\put(235.7186,-232.306){\fontsize{11.991}{1}\usefont{T1}{cmr}{m}{n}\selectfont\color{color_167934}t}
\put(239.5197,-232.306){\fontsize{11.991}{1}\usefont{T1}{cmr}{m}{n}\selectfont\color{color_167934}d}
\put(245.8989,-232.306){\fontsize{11.991}{1}\usefont{T1}{cmr}{m}{n}\selectfont\color{color_167934}o}
\put(252.1942,-232.306){\fontsize{11.991}{1}\usefont{T1}{cmr}{m}{n}\selectfont\color{color_167934}c}
\put(257.2784,-232.306){\fontsize{11.991}{1}\usefont{T1}{cmr}{m}{n}\selectfont\color{color_167934}s}
\put(261.871,-232.306){\fontsize{11.991}{1}\usefont{T1}{cmr}{m}{n}\selectfont\color{color_167934}/}
\put(266.5594,-232.306){\fontsize{11.991}{1}\usefont{T1}{cmr}{m}{n}\selectfont\color{color_167934}w}
\put(275.0491,-232.306){\fontsize{11.991}{1}\usefont{T1}{cmr}{m}{n}\selectfont\color{color_167934}i}
\put(277.831,-232.306){\fontsize{11.991}{1}\usefont{T1}{cmr}{m}{n}\selectfont\color{color_167934}k}
\put(283.2269,-232.306){\fontsize{11.991}{1}\usefont{T1}{cmr}{m}{n}\selectfont\color{color_167934}i}
\put(286.0089,-232.306){\fontsize{11.991}{1}\usefont{T1}{cmr}{m}{n}\selectfont\color{color_167934}/}
\put(290.6014,-232.306){\fontsize{11.991}{1}\usefont{T1}{cmr}{m}{n}\selectfont\color{color_167934}i}
\put(293.3833,-232.306){\fontsize{11.991}{1}\usefont{T1}{cmr}{m}{n}\selectfont\color{color_167934}n}
\put(299.7625,-232.306){\fontsize{11.991}{1}\usefont{T1}{cmr}{m}{n}\selectfont\color{color_167934}de}
\put(311.8495,-232.306){\fontsize{11.991}{1}\usefont{T1}{cmr}{m}{n}\selectfont\color{color_167934}x}
\put(316.9456,-232.306){\fontsize{11.991}{1}\usefont{T1}{cmr}{m}{n}\selectfont\color{color_167934}.}
\put(319.9434,-232.306){\fontsize{11.991}{1}\usefont{T1}{cmr}{m}{n}\selectfont\color{color_167934}p}
\put(326.3226,-232.306){\fontsize{11.991}{1}\usefont{T1}{cmr}{m}{n}\selectfont\color{color_167934}h}
\put(332.7018,-232.306){\fontsize{11.991}{1}\usefont{T1}{cmr}{m}{n}\selectfont\color{color_167934}p/}
\put(343.6855,-232.306){\fontsize{11.991}{1}\usefont{T1}{cmr}{m}{n}\selectfont\color{color_167934}t}
\put(347.5826,-232.306){\fontsize{11.991}{1}\usefont{T1}{cmr}{m}{n}\selectfont\color{color_167934}m}
\put(357.2593,-232.306){\fontsize{11.991}{1}\usefont{T1}{cmr}{m}{n}\selectfont\color{color_167934}ux}
\end{picture}
\begin{tikzpicture}[overlay]
\path(0pt,0pt);
\draw[color_167934,line width=0.794pt]
(195.189pt, -233.497pt) -- (369.095pt, -233.497pt)
;
\end{tikzpicture}
\begin{picture}(-5,0)(2.5,0)
\put(84.213,-260.397){\fontsize{11.991}{1}\usefont{T1}{cmr}{m}{n}\selectfont\color{color_29791}•}
\put(98.187,-261.106){\fontsize{11.991}{1}\usefont{T1}{cmr}{m}{n}\selectfont\color{color_29791}D}
\put(106.8685,-261.106){\fontsize{11.991}{1}\usefont{T1}{cmr}{m}{n}\selectfont\color{color_29791}e}
\put(113.5475,-261.106){\fontsize{11.991}{1}\usefont{T1}{cmr}{m}{n}\selectfont\color{color_29791}m}
\put(123.44,-261.106){\fontsize{11.991}{1}\usefont{T1}{cmr}{m}{n}\selectfont\color{color_29791}o}
\put(130.119,-261.106){\fontsize{11.991}{1}\usefont{T1}{cmr}{m}{n}\selectfont\color{color_29791}n}
\put(136.9059,-261.106){\fontsize{11.991}{1}\usefont{T1}{cmr}{m}{n}\selectfont\color{color_29791}st}
\put(146.199,-261.106){\fontsize{11.991}{1}\usefont{T1}{cmr}{m}{n}\selectfont\color{color_29791}ra}
\put(156.871,-261.106){\fontsize{11.991}{1}\usefont{T1}{cmr}{m}{n}\selectfont\color{color_29791}t}
\put(160.2524,-261.106){\fontsize{11.991}{1}\usefont{T1}{cmr}{m}{n}\selectfont\color{color_29791}i}
\put(162.8545,-261.106){\fontsize{11.991}{1}\usefont{T1}{cmr}{m}{n}\selectfont\color{color_29791}o}
\put(169.5335,-261.106){\fontsize{11.991}{1}\usefont{T1}{cmr}{m}{n}\selectfont\color{color_29791}n}
\end{picture}
\newpage
\begin{tikzpicture}[overlay]
\path(0pt,0pt);
\filldraw[color_283006][even odd rule]
(-15pt, 10pt) -- (704.972pt, 10pt)
 -- (704.972pt, 10pt)
 -- (704.972pt, -529.972pt)
 -- (704.972pt, -529.972pt)
 -- (-15pt, -529.972pt)
 -- (-15pt, -529.972pt)
 -- (-15pt, 10pt) -- cycle
;
\end{tikzpicture}
\begin{picture}(-5,0)(2.5,0)
\put(-15,-529.972){\includegraphics[width=720pt,height=540pt]{latexImage_15b9e6acda8be87c69673355bc94d374.png}}
\put(48.411,-60.69601){\fontsize{21.997}{1}\usefont{T1}{cmr}{b}{n}\selectfont\color{color_29791}S}
\put(63.083,-60.69601){\fontsize{21.997}{1}\usefont{T1}{cmr}{b}{n}\selectfont\color{color_29791}o}
\put(76.47917,-60.69601){\fontsize{21.997}{1}\usefont{T1}{cmr}{b}{n}\selectfont\color{color_29791}f}
\put(83.78217,-60.69601){\fontsize{21.997}{1}\usefont{T1}{cmr}{b}{n}\selectfont\color{color_29791}t}
\put(91.15117,-60.69601){\fontsize{21.997}{1}\usefont{T1}{cmr}{b}{n}\selectfont\color{color_29791}w}
\put(108.1548,-60.69601){\fontsize{21.997}{1}\usefont{T1}{cmr}{b}{n}\selectfont\color{color_29791}a}
\put(120.3632,-60.69601){\fontsize{21.997}{1}\usefont{T1}{cmr}{b}{n}\selectfont\color{color_29791}r}
\put(128.942,-60.69601){\fontsize{21.997}{1}\usefont{T1}{cmr}{b}{n}\selectfont\color{color_29791}e}
\put(48.213,-116.397){\fontsize{15.987}{1}\usefont{T1}{cmr}{m}{n}\selectfont\color{color_29791}O}
\put(60.507,-116.397){\fontsize{15.987}{1}\usefont{T1}{cmr}{m}{n}\selectfont\color{color_29791}per}
\put(83.57625,-116.397){\fontsize{15.987}{1}\usefont{T1}{cmr}{m}{n}\selectfont\color{color_29791}at}
\put(96.86144,-116.397){\fontsize{15.987}{1}\usefont{T1}{cmr}{m}{n}\selectfont\color{color_29791}i}
\put(100.4425,-116.397){\fontsize{15.987}{1}\usefont{T1}{cmr}{m}{n}\selectfont\color{color_29791}ng }
\put(122.6165,-116.397){\fontsize{15.987}{1}\usefont{T1}{cmr}{m}{n}\selectfont\color{color_29791}S}
\put(133.2958,-116.397){\fontsize{15.987}{1}\usefont{T1}{cmr}{m}{n}\selectfont\color{color_29791}ys}
\put(149.3628,-116.397){\fontsize{15.987}{1}\usefont{T1}{cmr}{m}{n}\selectfont\color{color_29791}t}
\put(153.7592,-116.397){\fontsize{15.987}{1}\usefont{T1}{cmr}{m}{n}\selectfont\color{color_29791}em}
\put(175.9332,-116.397){\fontsize{15.987}{1}\usefont{T1}{cmr}{m}{n}\selectfont\color{color_29791} }
\put(180.3296,-116.397){\fontsize{15.987}{1}\usefont{T1}{cmr}{m}{n}\selectfont\color{color_29791}s}
\put(188.403,-116.397){\fontsize{15.987}{1}\usefont{T1}{cmr}{m}{n}\selectfont\color{color_29791}of}
\put(201.6882,-116.397){\fontsize{15.987}{1}\usefont{T1}{cmr}{m}{n}\selectfont\color{color_29791}t}
\put(206.1646,-116.397){\fontsize{15.987}{1}\usefont{T1}{cmr}{m}{n}\selectfont\color{color_29791}w}
\put(217.6592,-116.397){\fontsize{15.987}{1}\usefont{T1}{cmr}{m}{n}\selectfont\color{color_29791}a}
\put(226.4521,-116.397){\fontsize{15.987}{1}\usefont{T1}{cmr}{m}{n}\selectfont\color{color_29791}r}
\put(231.7438,-116.397){\fontsize{15.987}{1}\usefont{T1}{cmr}{m}{n}\selectfont\color{color_29791}e}
\put(84.213,-145.31){\fontsize{11.991}{1}\usefont{T1}{cmr}{m}{n}\selectfont\color{color_29791}•}
\put(106.691,-145.991){\fontsize{11.991}{1}\usefont{T1}{cmr}{m}{n}\selectfont\color{color_29791}T}
\put(113.382,-145.991){\fontsize{11.991}{1}\usefont{T1}{cmr}{m}{n}\selectfont\color{color_29791}y}
\put(119.2815,-145.991){\fontsize{11.991}{1}\usefont{T1}{cmr}{m}{n}\selectfont\color{color_29791}p}
\put(125.9605,-145.991){\fontsize{11.991}{1}\usefont{T1}{cmr}{m}{n}\selectfont\color{color_29791}i}
\put(128.6465,-145.991){\fontsize{11.991}{1}\usefont{T1}{cmr}{m}{n}\selectfont\color{color_29791}c}
\put(134.5461,-145.991){\fontsize{11.991}{1}\usefont{T1}{cmr}{m}{n}\selectfont\color{color_29791}a}
\put(141.333,-145.991){\fontsize{11.991}{1}\usefont{T1}{cmr}{m}{n}\selectfont\color{color_29791}l}
\put(143.935,-145.991){\fontsize{11.991}{1}\usefont{T1}{cmr}{m}{n}\selectfont\color{color_29791} }
\put(147.3165,-145.991){\fontsize{11.991}{1}\usefont{T1}{cmr}{m}{n}\selectfont\color{color_29791}l}
\put(149.9185,-145.991){\fontsize{11.991}{1}\usefont{T1}{cmr}{m}{n}\selectfont\color{color_29791}i}
\put(152.5206,-145.991){\fontsize{11.991}{1}\usefont{T1}{cmr}{m}{n}\selectfont\color{color_29791}n}
\put(159.3075,-145.991){\fontsize{11.991}{1}\usefont{T1}{cmr}{m}{n}\selectfont\color{color_29791}u}
\put(165.9865,-145.991){\fontsize{11.991}{1}\usefont{T1}{cmr}{m}{n}\selectfont\color{color_29791}x }
\put(175.3635,-145.991){\fontsize{11.991}{1}\usefont{T1}{cmr}{m}{n}\selectfont\color{color_29791}c}
\put(181.263,-145.991){\fontsize{11.991}{1}\usefont{T1}{cmr}{m}{n}\selectfont\color{color_29791}o}
\put(188.0499,-145.991){\fontsize{11.991}{1}\usefont{T1}{cmr}{m}{n}\selectfont\color{color_29791}m}
\put(197.9425,-145.991){\fontsize{11.991}{1}\usefont{T1}{cmr}{m}{n}\selectfont\color{color_29791}ma}
\put(214.61,-145.991){\fontsize{11.991}{1}\usefont{T1}{cmr}{m}{n}\selectfont\color{color_29791}n}
\put(221.289,-145.991){\fontsize{11.991}{1}\usefont{T1}{cmr}{m}{n}\selectfont\color{color_29791}d}
\put(227.968,-145.991){\fontsize{11.991}{1}\usefont{T1}{cmr}{m}{n}\selectfont\color{color_29791}s}
\put(234.0474,-145.991){\fontsize{11.991}{1}\usefont{T1}{cmr}{m}{n}\selectfont\color{color_29791} }
\put(237.345,-145.991){\fontsize{11.991}{1}\usefont{T1}{cmr}{m}{n}\selectfont\color{color_29791}a}
\put(244.024,-145.991){\fontsize{11.991}{1}\usefont{T1}{cmr}{m}{n}\selectfont\color{color_29791}n}
\put(250.8109,-145.991){\fontsize{11.991}{1}\usefont{T1}{cmr}{m}{n}\selectfont\color{color_29791}d}
\put(257.4898,-145.991){\fontsize{11.991}{1}\usefont{T1}{cmr}{m}{n}\selectfont\color{color_29791} }
\put(260.7874,-145.991){\fontsize{11.991}{1}\usefont{T1}{cmr}{m}{n}\selectfont\color{color_29791}so}
\put(273.5698,-145.991){\fontsize{11.991}{1}\usefont{T1}{cmr}{m}{n}\selectfont\color{color_29791}m}
\put(283.4623,-145.991){\fontsize{11.991}{1}\usefont{T1}{cmr}{m}{n}\selectfont\color{color_29791}e}
\put(290.1413,-145.991){\fontsize{11.991}{1}\usefont{T1}{cmr}{m}{n}\selectfont\color{color_29791} }
\put(293.5228,-145.991){\fontsize{11.991}{1}\usefont{T1}{cmr}{m}{n}\selectfont\color{color_29791}g}
\put(300.2018,-145.991){\fontsize{11.991}{1}\usefont{T1}{cmr}{m}{n}\selectfont\color{color_29791}ra}
\put(310.8738,-145.991){\fontsize{11.991}{1}\usefont{T1}{cmr}{m}{n}\selectfont\color{color_29791}p}
\put(317.5528,-145.991){\fontsize{11.991}{1}\usefont{T1}{cmr}{m}{n}\selectfont\color{color_29791}h}
\put(324.2318,-145.991){\fontsize{11.991}{1}\usefont{T1}{cmr}{m}{n}\selectfont\color{color_29791}i}
\put(326.9177,-145.991){\fontsize{11.991}{1}\usefont{T1}{cmr}{m}{n}\selectfont\color{color_29791}ca}
\put(339.5922,-145.991){\fontsize{11.991}{1}\usefont{T1}{cmr}{m}{n}\selectfont\color{color_29791}l}
\put(342.1943,-145.991){\fontsize{11.991}{1}\usefont{T1}{cmr}{m}{n}\selectfont\color{color_29791} }
\put(345.5757,-145.991){\fontsize{11.991}{1}\usefont{T1}{cmr}{m}{n}\selectfont\color{color_29791}p}
\put(352.2547,-145.991){\fontsize{11.991}{1}\usefont{T1}{cmr}{m}{n}\selectfont\color{color_29791}a}
\put(358.9337,-145.991){\fontsize{11.991}{1}\usefont{T1}{cmr}{m}{n}\selectfont\color{color_29791}cka}
\put(377.7116,-145.991){\fontsize{11.991}{1}\usefont{T1}{cmr}{m}{n}\selectfont\color{color_29791}g}
\put(384.3906,-145.991){\fontsize{11.991}{1}\usefont{T1}{cmr}{m}{n}\selectfont\color{color_29791}e}
\put(391.0696,-145.991){\fontsize{11.991}{1}\usefont{T1}{cmr}{m}{n}\selectfont\color{color_29791}s }
\put(400.3626,-145.991){\fontsize{11.991}{1}\usefont{T1}{cmr}{m}{n}\selectfont\color{color_29791}a}
\put(407.1495,-145.991){\fontsize{11.991}{1}\usefont{T1}{cmr}{m}{n}\selectfont\color{color_29791}r}
\put(411.0466,-145.991){\fontsize{11.991}{1}\usefont{T1}{cmr}{m}{n}\selectfont\color{color_29791}e}
\put(417.7256,-145.991){\fontsize{11.991}{1}\usefont{T1}{cmr}{m}{n}\selectfont\color{color_29791} }
\put(421.107,-145.991){\fontsize{11.991}{1}\usefont{T1}{cmr}{m}{n}\selectfont\color{color_29791}i}
\put(423.793,-145.991){\fontsize{11.991}{1}\usefont{T1}{cmr}{m}{n}\selectfont\color{color_29791}n}
\put(430.472,-145.991){\fontsize{11.991}{1}\usefont{T1}{cmr}{m}{n}\selectfont\color{color_29791}st}
\put(439.765,-145.991){\fontsize{11.991}{1}\usefont{T1}{cmr}{m}{n}\selectfont\color{color_29791}a}
\put(446.444,-145.991){\fontsize{11.991}{1}\usefont{T1}{cmr}{m}{n}\selectfont\color{color_29791}l}
\put(449.13,-145.991){\fontsize{11.991}{1}\usefont{T1}{cmr}{m}{n}\selectfont\color{color_29791}l}
\put(451.732,-145.991){\fontsize{11.991}{1}\usefont{T1}{cmr}{m}{n}\selectfont\color{color_29791}e}
\put(458.5189,-145.991){\fontsize{11.991}{1}\usefont{T1}{cmr}{m}{n}\selectfont\color{color_29791}d}
\put(465.1979,-145.991){\fontsize{11.991}{1}\usefont{T1}{cmr}{m}{n}\selectfont\color{color_29791} }
\put(468.4954,-145.991){\fontsize{11.991}{1}\usefont{T1}{cmr}{m}{n}\selectfont\color{color_29791}a}
\put(475.1744,-145.991){\fontsize{11.991}{1}\usefont{T1}{cmr}{m}{n}\selectfont\color{color_29791}s}
\put(481.2538,-145.991){\fontsize{11.991}{1}\usefont{T1}{cmr}{m}{n}\selectfont\color{color_29791} }
\put(484.5514,-145.991){\fontsize{11.991}{1}\usefont{T1}{cmr}{m}{n}\selectfont\color{color_29791}p}
\put(491.2303,-145.991){\fontsize{11.991}{1}\usefont{T1}{cmr}{m}{n}\selectfont\color{color_29791}a}
\put(498.0173,-145.991){\fontsize{11.991}{1}\usefont{T1}{cmr}{m}{n}\selectfont\color{color_29791}r}
\put(501.9144,-145.991){\fontsize{11.991}{1}\usefont{T1}{cmr}{m}{n}\selectfont\color{color_29791}t}
\put(505.2958,-145.991){\fontsize{11.991}{1}\usefont{T1}{cmr}{m}{n}\selectfont\color{color_29791} }
\put(508.5934,-145.991){\fontsize{11.991}{1}\usefont{T1}{cmr}{m}{n}\selectfont\color{color_29791}o}
\put(515.2724,-145.991){\fontsize{11.991}{1}\usefont{T1}{cmr}{m}{n}\selectfont\color{color_29791}f}
\put(518.6539,-145.991){\fontsize{11.991}{1}\usefont{T1}{cmr}{m}{n}\selectfont\color{color_29791} }
\put(522.0353,-145.991){\fontsize{11.991}{1}\usefont{T1}{cmr}{m}{n}\selectfont\color{color_29791}O}
\put(531.3284,-145.991){\fontsize{11.991}{1}\usefont{T1}{cmr}{m}{n}\selectfont\color{color_29791}S}
\put(539.4103,-145.991){\fontsize{11.991}{1}\usefont{T1}{cmr}{m}{n}\selectfont\color{color_29791}.}
\put(84.213,-174.11){\fontsize{11.991}{1}\usefont{T1}{cmr}{m}{n}\selectfont\color{color_29791}•}
\put(106.691,-174.791){\fontsize{11.991}{1}\usefont{T1}{cmr}{m}{n}\selectfont\color{color_29791}T}
\put(113.9815,-174.791){\fontsize{11.991}{1}\usefont{T1}{cmr}{m}{n}\selectfont\color{color_29791}h}
\put(120.6605,-174.791){\fontsize{11.991}{1}\usefont{T1}{cmr}{m}{n}\selectfont\color{color_29791}e}
\put(127.3395,-174.791){\fontsize{11.991}{1}\usefont{T1}{cmr}{m}{n}\selectfont\color{color_29791}s}
\put(133.4189,-174.791){\fontsize{11.991}{1}\usefont{T1}{cmr}{m}{n}\selectfont\color{color_29791}e}
\put(140.0979,-174.791){\fontsize{11.991}{1}\usefont{T1}{cmr}{m}{n}\selectfont\color{color_29791} }
\put(143.3955,-174.791){\fontsize{11.991}{1}\usefont{T1}{cmr}{m}{n}\selectfont\color{color_29791}ca}
\put(156.0699,-174.791){\fontsize{11.991}{1}\usefont{T1}{cmr}{m}{n}\selectfont\color{color_29791}n}
\put(162.8569,-174.791){\fontsize{11.991}{1}\usefont{T1}{cmr}{m}{n}\selectfont\color{color_29791} }
\put(166.1544,-174.791){\fontsize{11.991}{1}\usefont{T1}{cmr}{m}{n}\selectfont\color{color_29791}b}
\put(172.9413,-174.791){\fontsize{11.991}{1}\usefont{T1}{cmr}{m}{n}\selectfont\color{color_29791}e}
\put(179.6203,-174.791){\fontsize{11.991}{1}\usefont{T1}{cmr}{m}{n}\selectfont\color{color_29791} }
\put(182.9178,-174.791){\fontsize{11.991}{1}\usefont{T1}{cmr}{m}{n}\selectfont\color{color_29791}ru}
\put(193.5898,-174.791){\fontsize{11.991}{1}\usefont{T1}{cmr}{m}{n}\selectfont\color{color_29791}n}
\put(200.2688,-174.791){\fontsize{11.991}{1}\usefont{T1}{cmr}{m}{n}\selectfont\color{color_29791} }
\put(203.6503,-174.791){\fontsize{11.991}{1}\usefont{T1}{cmr}{m}{n}\selectfont\color{color_29791}f}
\put(207.0317,-174.791){\fontsize{11.991}{1}\usefont{T1}{cmr}{m}{n}\selectfont\color{color_29791}r}
\put(210.9288,-174.791){\fontsize{11.991}{1}\usefont{T1}{cmr}{m}{n}\selectfont\color{color_29791}o}
\put(217.6078,-174.791){\fontsize{11.991}{1}\usefont{T1}{cmr}{m}{n}\selectfont\color{color_29791}m }
\put(230.8938,-174.791){\fontsize{11.991}{1}\usefont{T1}{cmr}{m}{n}\selectfont\color{color_29791}t}
\put(234.2753,-174.791){\fontsize{11.991}{1}\usefont{T1}{cmr}{m}{n}\selectfont\color{color_29791}h}
\put(240.9543,-174.791){\fontsize{11.991}{1}\usefont{T1}{cmr}{m}{n}\selectfont\color{color_29791}e}
\put(247.7412,-174.791){\fontsize{11.991}{1}\usefont{T1}{cmr}{m}{n}\selectfont\color{color_29791} }
\put(251.0387,-174.791){\fontsize{11.991}{1}\usefont{T1}{cmr}{m}{n}\selectfont\color{color_29791}co}
\put(263.7132,-174.791){\fontsize{11.991}{1}\usefont{T1}{cmr}{m}{n}\selectfont\color{color_29791}mma}
\put(290.3692,-174.791){\fontsize{11.991}{1}\usefont{T1}{cmr}{m}{n}\selectfont\color{color_29791}n}
\put(297.0482,-174.791){\fontsize{11.991}{1}\usefont{T1}{cmr}{m}{n}\selectfont\color{color_29791}d}
\put(303.7272,-174.791){\fontsize{11.991}{1}\usefont{T1}{cmr}{m}{n}\selectfont\color{color_29791} }
\put(307.1086,-174.791){\fontsize{11.991}{1}\usefont{T1}{cmr}{m}{n}\selectfont\color{color_29791}l}
\put(309.7107,-174.791){\fontsize{11.991}{1}\usefont{T1}{cmr}{m}{n}\selectfont\color{color_29791}i}
\put(312.3967,-174.791){\fontsize{11.991}{1}\usefont{T1}{cmr}{m}{n}\selectfont\color{color_29791}n}
\put(319.0757,-174.791){\fontsize{11.991}{1}\usefont{T1}{cmr}{m}{n}\selectfont\color{color_29791}e}
\put(325.7546,-174.791){\fontsize{11.991}{1}\usefont{T1}{cmr}{m}{n}\selectfont\color{color_29791} }
\put(329.1361,-174.791){\fontsize{11.991}{1}\usefont{T1}{cmr}{m}{n}\selectfont\color{color_29791}a}
\put(335.8151,-174.791){\fontsize{11.991}{1}\usefont{T1}{cmr}{m}{n}\selectfont\color{color_29791}n}
\put(342.4941,-174.791){\fontsize{11.991}{1}\usefont{T1}{cmr}{m}{n}\selectfont\color{color_29791}d}
\put(349.1731,-174.791){\fontsize{11.991}{1}\usefont{T1}{cmr}{m}{n}\selectfont\color{color_29791} }
\put(352.5545,-174.791){\fontsize{11.991}{1}\usefont{T1}{cmr}{m}{n}\selectfont\color{color_29791}d}
\put(359.2335,-174.791){\fontsize{11.991}{1}\usefont{T1}{cmr}{m}{n}\selectfont\color{color_29791}e}
\put(365.9125,-174.791){\fontsize{11.991}{1}\usefont{T1}{cmr}{m}{n}\selectfont\color{color_29791}skt}
\put(381.285,-174.791){\fontsize{11.991}{1}\usefont{T1}{cmr}{m}{n}\selectfont\color{color_29791}o}
\put(387.964,-174.791){\fontsize{11.991}{1}\usefont{T1}{cmr}{m}{n}\selectfont\color{color_29791}p}
\put(394.6429,-174.791){\fontsize{11.991}{1}\usefont{T1}{cmr}{m}{n}\selectfont\color{color_29791} }
\put(398.0244,-174.791){\fontsize{11.991}{1}\usefont{T1}{cmr}{m}{n}\selectfont\color{color_29791}i}
\put(400.6265,-174.791){\fontsize{11.991}{1}\usefont{T1}{cmr}{m}{n}\selectfont\color{color_29791}n}
\put(407.3055,-174.791){\fontsize{11.991}{1}\usefont{T1}{cmr}{m}{n}\selectfont\color{color_29791}t}
\put(410.6869,-174.791){\fontsize{11.991}{1}\usefont{T1}{cmr}{m}{n}\selectfont\color{color_29791}e}
\put(417.3659,-174.791){\fontsize{11.991}{1}\usefont{T1}{cmr}{m}{n}\selectfont\color{color_29791}rf}
\put(424.6564,-174.791){\fontsize{11.991}{1}\usefont{T1}{cmr}{m}{n}\selectfont\color{color_29791}a}
\put(431.4433,-174.791){\fontsize{11.991}{1}\usefont{T1}{cmr}{m}{n}\selectfont\color{color_29791}ce}
\put(48.213,-205.093){\fontsize{15.987}{1}\usefont{T1}{cmr}{m}{n}\selectfont\color{color_29791}M}
\put(61.4982,-205.093){\fontsize{15.987}{1}\usefont{T1}{cmr}{m}{n}\selectfont\color{color_29791}odul}
\put(91.74561,-205.093){\fontsize{15.987}{1}\usefont{T1}{cmr}{m}{n}\selectfont\color{color_29791}e}
\put(100.5385,-205.093){\fontsize{15.987}{1}\usefont{T1}{cmr}{m}{n}\selectfont\color{color_29791}s}
\put(84.213,-221.987){\fontsize{11.991}{1}\usefont{T1}{cmr}{m}{n}\selectfont\color{color_29791}•}
\put(106.691,-222.696){\fontsize{11.991}{1}\usefont{T1}{cmr}{m}{n}\selectfont\color{color_29791}D}
\put(115.2885,-222.696){\fontsize{11.991}{1}\usefont{T1}{cmr}{m}{n}\selectfont\color{color_29791}o}
\put(122.0755,-222.696){\fontsize{11.991}{1}\usefont{T1}{cmr}{m}{n}\selectfont\color{color_29791} }
\put(125.373,-222.696){\fontsize{11.991}{1}\usefont{T1}{cmr}{m}{n}\selectfont\color{color_29791}n}
\put(132.052,-222.696){\fontsize{11.991}{1}\usefont{T1}{cmr}{m}{n}\selectfont\color{color_29791}o}
\put(138.8389,-222.696){\fontsize{11.991}{1}\usefont{T1}{cmr}{m}{n}\selectfont\color{color_29791}t}
\put(142.1364,-222.696){\fontsize{11.991}{1}\usefont{T1}{cmr}{m}{n}\selectfont\color{color_29791} }
\put(145.5179,-222.696){\fontsize{11.991}{1}\usefont{T1}{cmr}{m}{n}\selectfont\color{color_29791}w}
\put(154.1154,-222.696){\fontsize{11.991}{1}\usefont{T1}{cmr}{m}{n}\selectfont\color{color_29791}o}
\put(160.7944,-222.696){\fontsize{11.991}{1}\usefont{T1}{cmr}{m}{n}\selectfont\color{color_29791}rk }
\put(174.0804,-222.696){\fontsize{11.991}{1}\usefont{T1}{cmr}{m}{n}\selectfont\color{color_29791}o}
\put(180.8674,-222.696){\fontsize{11.991}{1}\usefont{T1}{cmr}{m}{n}\selectfont\color{color_29791}n}
\put(187.5463,-222.696){\fontsize{11.991}{1}\usefont{T1}{cmr}{m}{n}\selectfont\color{color_29791} }
\put(190.9278,-222.696){\fontsize{11.991}{1}\usefont{T1}{cmr}{m}{n}\selectfont\color{color_29791}b}
\put(197.6068,-222.696){\fontsize{11.991}{1}\usefont{T1}{cmr}{m}{n}\selectfont\color{color_29791}sl}
\put(206.2043,-222.696){\fontsize{11.991}{1}\usefont{T1}{cmr}{m}{n}\selectfont\color{color_29791}ce}
\put(218.8788,-222.696){\fontsize{11.991}{1}\usefont{T1}{cmr}{m}{n}\selectfont\color{color_29791}n}
\put(225.6657,-222.696){\fontsize{11.991}{1}\usefont{T1}{cmr}{m}{n}\selectfont\color{color_29791}b}
\put(232.3447,-222.696){\fontsize{11.991}{1}\usefont{T1}{cmr}{m}{n}\selectfont\color{color_29791} }
\put(235.6422,-222.696){\fontsize{11.991}{1}\usefont{T1}{cmr}{m}{n}\selectfont\color{color_29791}o}
\put(242.3212,-222.696){\fontsize{11.991}{1}\usefont{T1}{cmr}{m}{n}\selectfont\color{color_29791}r }
\put(249.6957,-222.696){\fontsize{11.991}{1}\usefont{T1}{cmr}{m}{n}\selectfont\color{color_29791}b}
\put(256.3747,-222.696){\fontsize{11.991}{1}\usefont{T1}{cmr}{m}{n}\selectfont\color{color_29791}sl}
\put(264.9722,-222.696){\fontsize{11.991}{1}\usefont{T1}{cmr}{m}{n}\selectfont\color{color_29791}ce}
\put(277.6467,-222.696){\fontsize{11.991}{1}\usefont{T1}{cmr}{m}{n}\selectfont\color{color_29791}n}
\put(284.4336,-222.696){\fontsize{11.991}{1}\usefont{T1}{cmr}{m}{n}\selectfont\color{color_29791}c}
\put(84.213,-250.787){\fontsize{11.991}{1}\usefont{T1}{cmr}{m}{n}\selectfont\color{color_29791}•}
\put(106.691,-251.496){\fontsize{11.991}{1}\usefont{T1}{cmr}{m}{n}\selectfont\color{color_29791}T}
\put(113.9815,-251.496){\fontsize{11.991}{1}\usefont{T1}{cmr}{m}{n}\selectfont\color{color_29791}h}
\put(120.6605,-251.496){\fontsize{11.991}{1}\usefont{T1}{cmr}{m}{n}\selectfont\color{color_29791}e}
\put(127.3395,-251.496){\fontsize{11.991}{1}\usefont{T1}{cmr}{m}{n}\selectfont\color{color_29791}re}
\put(138.0115,-251.496){\fontsize{11.991}{1}\usefont{T1}{cmr}{m}{n}\selectfont\color{color_29791} }
\put(141.393,-251.496){\fontsize{11.991}{1}\usefont{T1}{cmr}{m}{n}\selectfont\color{color_29791}a}
\put(148.072,-251.496){\fontsize{11.991}{1}\usefont{T1}{cmr}{m}{n}\selectfont\color{color_29791}re}
\put(158.744,-251.496){\fontsize{11.991}{1}\usefont{T1}{cmr}{m}{n}\selectfont\color{color_29791} }
\put(162.1254,-251.496){\fontsize{11.991}{1}\usefont{T1}{cmr}{m}{n}\selectfont\color{color_29791}t}
\put(165.423,-251.496){\fontsize{11.991}{1}\usefont{T1}{cmr}{m}{n}\selectfont\color{color_29791}w}
\put(174.0205,-251.496){\fontsize{11.991}{1}\usefont{T1}{cmr}{m}{n}\selectfont\color{color_29791}o}
\put(180.8074,-251.496){\fontsize{11.991}{1}\usefont{T1}{cmr}{m}{n}\selectfont\color{color_29791} }
\put(184.1049,-251.496){\fontsize{11.991}{1}\usefont{T1}{cmr}{m}{n}\selectfont\color{color_29791}mo}
\put(200.7724,-251.496){\fontsize{11.991}{1}\usefont{T1}{cmr}{m}{n}\selectfont\color{color_29791}d}
\put(207.4514,-251.496){\fontsize{11.991}{1}\usefont{T1}{cmr}{m}{n}\selectfont\color{color_29791}u}
\put(214.2383,-251.496){\fontsize{11.991}{1}\usefont{T1}{cmr}{m}{n}\selectfont\color{color_29791}l}
\put(216.8404,-251.496){\fontsize{11.991}{1}\usefont{T1}{cmr}{m}{n}\selectfont\color{color_29791}e}
\put(223.5194,-251.496){\fontsize{11.991}{1}\usefont{T1}{cmr}{m}{n}\selectfont\color{color_29791} }
\put(226.9008,-251.496){\fontsize{11.991}{1}\usefont{T1}{cmr}{m}{n}\selectfont\color{color_29791}r}
\put(230.7979,-251.496){\fontsize{11.991}{1}\usefont{T1}{cmr}{m}{n}\selectfont\color{color_29791}e}
\put(237.5848,-251.496){\fontsize{11.991}{1}\usefont{T1}{cmr}{m}{n}\selectfont\color{color_29791}p}
\put(244.2638,-251.496){\fontsize{11.991}{1}\usefont{T1}{cmr}{m}{n}\selectfont\color{color_29791}o}
\put(250.9428,-251.496){\fontsize{11.991}{1}\usefont{T1}{cmr}{m}{n}\selectfont\color{color_29791}si}
\put(259.5403,-251.496){\fontsize{11.991}{1}\usefont{T1}{cmr}{m}{n}\selectfont\color{color_29791}t}
\put(262.9218,-251.496){\fontsize{11.991}{1}\usefont{T1}{cmr}{m}{n}\selectfont\color{color_29791}o}
\put(269.6008,-251.496){\fontsize{11.991}{1}\usefont{T1}{cmr}{m}{n}\selectfont\color{color_29791}ri}
\put(276.1959,-251.496){\fontsize{11.991}{1}\usefont{T1}{cmr}{m}{n}\selectfont\color{color_29791}e}
\put(282.8748,-251.496){\fontsize{11.991}{1}\usefont{T1}{cmr}{m}{n}\selectfont\color{color_29791}s:}
\put(292.2518,-251.496){\fontsize{11.991}{1}\usefont{T1}{cmr}{m}{n}\selectfont\color{color_29791} }
\put(295.5493,-251.496){\fontsize{11.991}{1}\usefont{T1}{cmr}{m}{n}\selectfont\color{color_29791}/}
\put(298.9308,-251.496){\fontsize{11.991}{1}\usefont{T1}{cmr}{m}{n}\selectfont\color{color_29791}p}
\put(305.6098,-251.496){\fontsize{11.991}{1}\usefont{T1}{cmr}{m}{n}\selectfont\color{color_29791}a}
\put(312.3967,-251.496){\fontsize{11.991}{1}\usefont{T1}{cmr}{m}{n}\selectfont\color{color_29791}ck}
\put(324.2917,-251.496){\fontsize{11.991}{1}\usefont{T1}{cmr}{m}{n}\selectfont\color{color_29791}a}
\put(331.0786,-251.496){\fontsize{11.991}{1}\usefont{T1}{cmr}{m}{n}\selectfont\color{color_29791}g}
\put(337.7576,-251.496){\fontsize{11.991}{1}\usefont{T1}{cmr}{m}{n}\selectfont\color{color_29791}e}
\put(344.4366,-251.496){\fontsize{11.991}{1}\usefont{T1}{cmr}{m}{n}\selectfont\color{color_29791}s/}
\put(353.8136,-251.496){\fontsize{11.991}{1}\usefont{T1}{cmr}{m}{n}\selectfont\color{color_29791}m}
\put(363.7061,-251.496){\fontsize{11.991}{1}\usefont{T1}{cmr}{m}{n}\selectfont\color{color_29791}o}
\put(370.493,-251.496){\fontsize{11.991}{1}\usefont{T1}{cmr}{m}{n}\selectfont\color{color_29791}d}
\put(377.172,-251.496){\fontsize{11.991}{1}\usefont{T1}{cmr}{m}{n}\selectfont\color{color_29791}u}
\put(383.851,-251.496){\fontsize{11.991}{1}\usefont{T1}{cmr}{m}{n}\selectfont\color{color_29791}l}
\put(386.4531,-251.496){\fontsize{11.991}{1}\usefont{T1}{cmr}{m}{n}\selectfont\color{color_29791}e}
\put(393.132,-251.496){\fontsize{11.991}{1}\usefont{T1}{cmr}{m}{n}\selectfont\color{color_29791}s}
\put(399.2115,-251.496){\fontsize{11.991}{1}\usefont{T1}{cmr}{m}{n}\selectfont\color{color_29791} }
\put(402.509,-251.496){\fontsize{11.991}{1}\usefont{T1}{cmr}{m}{n}\selectfont\color{color_29791}\& }
\put(413.8765,-251.496){\fontsize{11.991}{1}\usefont{T1}{cmr}{m}{n}\selectfont\color{color_29791}/}
\put(417.2579,-251.496){\fontsize{11.991}{1}\usefont{T1}{cmr}{m}{n}\selectfont\color{color_29791}h}
\put(423.9369,-251.496){\fontsize{11.991}{1}\usefont{T1}{cmr}{m}{n}\selectfont\color{color_29791}p}
\put(430.6159,-251.496){\fontsize{11.991}{1}\usefont{T1}{cmr}{m}{n}\selectfont\color{color_29791}cp}
\put(443.2904,-251.496){\fontsize{11.991}{1}\usefont{T1}{cmr}{m}{n}\selectfont\color{color_29791}a}
\put(449.9694,-251.496){\fontsize{11.991}{1}\usefont{T1}{cmr}{m}{n}\selectfont\color{color_29791}c}
\put(456.0488,-251.496){\fontsize{11.991}{1}\usefont{T1}{cmr}{m}{n}\selectfont\color{color_29791}ka}
\put(468.7233,-251.496){\fontsize{11.991}{1}\usefont{T1}{cmr}{m}{n}\selectfont\color{color_29791}g}
\put(475.4023,-251.496){\fontsize{11.991}{1}\usefont{T1}{cmr}{m}{n}\selectfont\color{color_29791}e}
\put(482.0813,-251.496){\fontsize{11.991}{1}\usefont{T1}{cmr}{m}{n}\selectfont\color{color_29791}s/}
\put(491.4583,-251.496){\fontsize{11.991}{1}\usefont{T1}{cmr}{m}{n}\selectfont\color{color_29791}m}
\put(501.3508,-251.496){\fontsize{11.991}{1}\usefont{T1}{cmr}{m}{n}\selectfont\color{color_29791}o}
\put(508.0298,-251.496){\fontsize{11.991}{1}\usefont{T1}{cmr}{m}{n}\selectfont\color{color_29791}d}
\put(514.8168,-251.496){\fontsize{11.991}{1}\usefont{T1}{cmr}{m}{n}\selectfont\color{color_29791}u}
\put(521.4958,-251.496){\fontsize{11.991}{1}\usefont{T1}{cmr}{m}{n}\selectfont\color{color_29791}l}
\put(524.0978,-251.496){\fontsize{11.991}{1}\usefont{T1}{cmr}{m}{n}\selectfont\color{color_29791}e}
\put(530.8848,-251.496){\fontsize{11.991}{1}\usefont{T1}{cmr}{m}{n}\selectfont\color{color_29791}s}
\put(537.813,-251.496){\fontsize{11.991}{1}\usefont{T1}{cmr}{m}{n}\selectfont\color{color_29791} }
\put(84.213,-279.587){\fontsize{11.991}{1}\usefont{T1}{cmr}{m}{n}\selectfont\color{color_29791}•}
\put(106.691,-280.211){\fontsize{11.991}{1}\usefont{T1}{cmr}{m}{n}\selectfont\color{color_29791}Pre}
\put(125.349,-280.211){\fontsize{11.991}{1}\usefont{T1}{cmr}{m}{n}\selectfont\color{color_29791}f}
\put(128.7305,-280.211){\fontsize{11.991}{1}\usefont{T1}{cmr}{m}{n}\selectfont\color{color_29791}e}
\put(135.4095,-280.211){\fontsize{11.991}{1}\usefont{T1}{cmr}{m}{n}\selectfont\color{color_29791}r }
\put(142.7,-280.211){\fontsize{11.991}{1}\usefont{T1}{cmr}{m}{n}\selectfont\color{color_29791}/}
\put(146.0815,-280.211){\fontsize{11.991}{1}\usefont{T1}{cmr}{m}{n}\selectfont\color{color_29791}h}
\put(152.7604,-280.211){\fontsize{11.991}{1}\usefont{T1}{cmr}{m}{n}\selectfont\color{color_29791}p}
\put(159.4394,-280.211){\fontsize{11.991}{1}\usefont{T1}{cmr}{m}{n}\selectfont\color{color_29791}cp}
\put(172.1139,-280.211){\fontsize{11.991}{1}\usefont{T1}{cmr}{m}{n}\selectfont\color{color_29791}a}
\put(178.9008,-280.211){\fontsize{11.991}{1}\usefont{T1}{cmr}{m}{n}\selectfont\color{color_29791}c}
\put(184.8004,-280.211){\fontsize{11.991}{1}\usefont{T1}{cmr}{m}{n}\selectfont\color{color_29791}k}
\put(190.8798,-280.211){\fontsize{11.991}{1}\usefont{T1}{cmr}{m}{n}\selectfont\color{color_29791}a}
\put(197.5588,-280.211){\fontsize{11.991}{1}\usefont{T1}{cmr}{m}{n}\selectfont\color{color_29791}g}
\put(204.2378,-280.211){\fontsize{11.991}{1}\usefont{T1}{cmr}{m}{n}\selectfont\color{color_29791}e}
\put(210.9168,-280.211){\fontsize{11.991}{1}\usefont{T1}{cmr}{m}{n}\selectfont\color{color_29791}s/}
\put(220.2937,-280.211){\fontsize{11.991}{1}\usefont{T1}{cmr}{m}{n}\selectfont\color{color_29791}m}
\put(230.1863,-280.211){\fontsize{11.991}{1}\usefont{T1}{cmr}{m}{n}\selectfont\color{color_29791}o}
\put(236.9732,-280.211){\fontsize{11.991}{1}\usefont{T1}{cmr}{m}{n}\selectfont\color{color_29791}d}
\put(243.6522,-280.211){\fontsize{11.991}{1}\usefont{T1}{cmr}{m}{n}\selectfont\color{color_29791}u}
\put(250.3312,-280.211){\fontsize{11.991}{1}\usefont{T1}{cmr}{m}{n}\selectfont\color{color_29791}l}
\put(252.9333,-280.211){\fontsize{11.991}{1}\usefont{T1}{cmr}{m}{n}\selectfont\color{color_29791}e}
\put(259.6122,-280.211){\fontsize{11.991}{1}\usefont{T1}{cmr}{m}{n}\selectfont\color{color_29791}s}
\put(265.6917,-280.211){\fontsize{11.991}{1}\usefont{T1}{cmr}{m}{n}\selectfont\color{color_29791} }
\put(268.9892,-280.211){\fontsize{11.991}{1}\usefont{T1}{cmr}{m}{n}\selectfont\color{color_29791}-}
\put(273.397,-280.211){\fontsize{11.991}{1}\usefont{T1}{cmr}{m}{n}\selectfont\color{color_29791} }
\put(276.713,-280.211){\fontsize{11.991}{1}\usefont{T1}{cmr}{m}{n}\selectfont\color{color_29791}w}
\put(285.3945,-280.211){\fontsize{11.991}{1}\usefont{T1}{cmr}{m}{n}\selectfont\color{color_29791}o}
\put(292.0735,-280.211){\fontsize{11.991}{1}\usefont{T1}{cmr}{m}{n}\selectfont\color{color_29791}r}
\put(295.9706,-280.211){\fontsize{11.991}{1}\usefont{T1}{cmr}{m}{n}\selectfont\color{color_29791}k}
\put(302.05,-280.211){\fontsize{11.991}{1}\usefont{T1}{cmr}{m}{n}\selectfont\color{color_29791}s }
\put(311.343,-280.211){\fontsize{11.991}{1}\usefont{T1}{cmr}{m}{n}\selectfont\color{color_29791}w}
\put(320.0245,-280.211){\fontsize{11.991}{1}\usefont{T1}{cmr}{m}{n}\selectfont\color{color_29791}i}
\put(322.6266,-280.211){\fontsize{11.991}{1}\usefont{T1}{cmr}{m}{n}\selectfont\color{color_29791}t}
\put(326.008,-280.211){\fontsize{11.991}{1}\usefont{T1}{cmr}{m}{n}\selectfont\color{color_29791}h}
\put(332.687,-280.211){\fontsize{11.991}{1}\usefont{T1}{cmr}{m}{n}\selectfont\color{color_29791} }
\put(335.9845,-280.211){\fontsize{11.991}{1}\usefont{T1}{cmr}{m}{n}\selectfont\color{color_29791}n}
\put(342.6635,-280.211){\fontsize{11.991}{1}\usefont{T1}{cmr}{m}{n}\selectfont\color{color_29791}o}
\put(349.4504,-280.211){\fontsize{11.991}{1}\usefont{T1}{cmr}{m}{n}\selectfont\color{color_29791}d}
\put(356.1294,-280.211){\fontsize{11.991}{1}\usefont{T1}{cmr}{m}{n}\selectfont\color{color_29791}e}
\put(362.8084,-280.211){\fontsize{11.991}{1}\usefont{T1}{cmr}{m}{n}\selectfont\color{color_29791}s }
\put(372.1853,-280.211){\fontsize{11.991}{1}\usefont{T1}{cmr}{m}{n}\selectfont\color{color_29791}a}
\put(378.8643,-280.211){\fontsize{11.991}{1}\usefont{T1}{cmr}{m}{n}\selectfont\color{color_29791}n}
\put(385.5433,-280.211){\fontsize{11.991}{1}\usefont{T1}{cmr}{m}{n}\selectfont\color{color_29791}d}
\put(392.2223,-280.211){\fontsize{11.991}{1}\usefont{T1}{cmr}{m}{n}\selectfont\color{color_29791} }
\put(395.6038,-280.211){\fontsize{11.991}{1}\usefont{T1}{cmr}{m}{n}\selectfont\color{color_29791}w}
\put(404.2013,-280.211){\fontsize{11.991}{1}\usefont{T1}{cmr}{m}{n}\selectfont\color{color_29791}o}
\put(410.8803,-280.211){\fontsize{11.991}{1}\usefont{T1}{cmr}{m}{n}\selectfont\color{color_29791}rkst}
\put(430.2458,-280.211){\fontsize{11.991}{1}\usefont{T1}{cmr}{m}{n}\selectfont\color{color_29791}a}
\put(436.9247,-280.211){\fontsize{11.991}{1}\usefont{T1}{cmr}{m}{n}\selectfont\color{color_29791}t}
\put(440.2223,-280.211){\fontsize{11.991}{1}\usefont{T1}{cmr}{m}{n}\selectfont\color{color_29791}i}
\put(442.9082,-280.211){\fontsize{11.991}{1}\usefont{T1}{cmr}{m}{n}\selectfont\color{color_29791}o}
\put(449.5872,-280.211){\fontsize{11.991}{1}\usefont{T1}{cmr}{m}{n}\selectfont\color{color_29791}n}
\put(456.2662,-280.211){\fontsize{11.991}{1}\usefont{T1}{cmr}{m}{n}\selectfont\color{color_29791}s}
\put(84.213,-308.302){\fontsize{11.991}{1}\usefont{T1}{cmr}{m}{n}\selectfont\color{color_29791}•}
\put(106.691,-309.011){\fontsize{11.991}{1}\usefont{T1}{cmr}{m}{n}\selectfont\color{color_29791}M}
\put(116.5836,-309.011){\fontsize{11.991}{1}\usefont{T1}{cmr}{m}{n}\selectfont\color{color_29791}o}
\put(123.3705,-309.011){\fontsize{11.991}{1}\usefont{T1}{cmr}{m}{n}\selectfont\color{color_29791}d}
\put(130.0495,-309.011){\fontsize{11.991}{1}\usefont{T1}{cmr}{m}{n}\selectfont\color{color_29791}u}
\put(136.7285,-309.011){\fontsize{11.991}{1}\usefont{T1}{cmr}{m}{n}\selectfont\color{color_29791}l}
\put(139.4144,-309.011){\fontsize{11.991}{1}\usefont{T1}{cmr}{m}{n}\selectfont\color{color_29791}e}
\put(146.0934,-309.011){\fontsize{11.991}{1}\usefont{T1}{cmr}{m}{n}\selectfont\color{color_29791}s }
\put(155.3865,-309.011){\fontsize{11.991}{1}\usefont{T1}{cmr}{m}{n}\selectfont\color{color_29791}so}
\put(168.1689,-309.011){\fontsize{11.991}{1}\usefont{T1}{cmr}{m}{n}\selectfont\color{color_29791}m}
\put(178.0614,-309.011){\fontsize{11.991}{1}\usefont{T1}{cmr}{m}{n}\selectfont\color{color_29791}e}
\put(184.7404,-309.011){\fontsize{11.991}{1}\usefont{T1}{cmr}{m}{n}\selectfont\color{color_29791}t}
\put(188.1219,-309.011){\fontsize{11.991}{1}\usefont{T1}{cmr}{m}{n}\selectfont\color{color_29791}i}
\put(190.7239,-309.011){\fontsize{11.991}{1}\usefont{T1}{cmr}{m}{n}\selectfont\color{color_29791}me}
\put(207.3914,-309.011){\fontsize{11.991}{1}\usefont{T1}{cmr}{m}{n}\selectfont\color{color_29791}s }
\put(216.7684,-309.011){\fontsize{11.991}{1}\usefont{T1}{cmr}{m}{n}\selectfont\color{color_29791}i}
\put(219.3705,-309.011){\fontsize{11.991}{1}\usefont{T1}{cmr}{m}{n}\selectfont\color{color_29791}n}
\put(226.0494,-309.011){\fontsize{11.991}{1}\usefont{T1}{cmr}{m}{n}\selectfont\color{color_29791}cl}
\put(234.7309,-309.011){\fontsize{11.991}{1}\usefont{T1}{cmr}{m}{n}\selectfont\color{color_29791}u}
\put(241.4099,-309.011){\fontsize{11.991}{1}\usefont{T1}{cmr}{m}{n}\selectfont\color{color_29791}d}
\put(248.0889,-309.011){\fontsize{11.991}{1}\usefont{T1}{cmr}{m}{n}\selectfont\color{color_29791}e}
\put(254.7679,-309.011){\fontsize{11.991}{1}\usefont{T1}{cmr}{m}{n}\selectfont\color{color_29791} }
\put(258.1494,-309.011){\fontsize{11.991}{1}\usefont{T1}{cmr}{m}{n}\selectfont\color{color_29791}t}
\put(261.5308,-309.011){\fontsize{11.991}{1}\usefont{T1}{cmr}{m}{n}\selectfont\color{color_29791}h}
\put(268.2098,-309.011){\fontsize{11.991}{1}\usefont{T1}{cmr}{m}{n}\selectfont\color{color_29791}e}
\put(274.8888,-309.011){\fontsize{11.991}{1}\usefont{T1}{cmr}{m}{n}\selectfont\color{color_29791} }
\put(278.1863,-309.011){\fontsize{11.991}{1}\usefont{T1}{cmr}{m}{n}\selectfont\color{color_29791}c}
\put(284.2657,-309.011){\fontsize{11.991}{1}\usefont{T1}{cmr}{m}{n}\selectfont\color{color_29791}o}
\put(290.9447,-309.011){\fontsize{11.991}{1}\usefont{T1}{cmr}{m}{n}\selectfont\color{color_29791}m}
\put(300.8373,-309.011){\fontsize{11.991}{1}\usefont{T1}{cmr}{m}{n}\selectfont\color{color_29791}p}
\put(307.6242,-309.011){\fontsize{11.991}{1}\usefont{T1}{cmr}{m}{n}\selectfont\color{color_29791}i}
\put(310.2263,-309.011){\fontsize{11.991}{1}\usefont{T1}{cmr}{m}{n}\selectfont\color{color_29791}l}
\put(312.8283,-309.011){\fontsize{11.991}{1}\usefont{T1}{cmr}{m}{n}\selectfont\color{color_29791}e}
\put(319.6152,-309.011){\fontsize{11.991}{1}\usefont{T1}{cmr}{m}{n}\selectfont\color{color_29791}r}
\put(323.5123,-309.011){\fontsize{11.991}{1}\usefont{T1}{cmr}{m}{n}\selectfont\color{color_29791} }
\put(326.8937,-309.011){\fontsize{11.991}{1}\usefont{T1}{cmr}{m}{n}\selectfont\color{color_29791}u}
\put(333.5727,-309.011){\fontsize{11.991}{1}\usefont{T1}{cmr}{m}{n}\selectfont\color{color_29791}se}
\put(346.2472,-309.011){\fontsize{11.991}{1}\usefont{T1}{cmr}{m}{n}\selectfont\color{color_29791}d}
\put(352.9262,-309.011){\fontsize{11.991}{1}\usefont{T1}{cmr}{m}{n}\selectfont\color{color_29791} }
\put(356.3077,-309.011){\fontsize{11.991}{1}\usefont{T1}{cmr}{m}{n}\selectfont\color{color_29791}i}
\put(358.9097,-309.011){\fontsize{11.991}{1}\usefont{T1}{cmr}{m}{n}\selectfont\color{color_29791}n}
\put(365.6966,-309.011){\fontsize{11.991}{1}\usefont{T1}{cmr}{m}{n}\selectfont\color{color_29791} }
\put(368.9941,-309.011){\fontsize{11.991}{1}\usefont{T1}{cmr}{m}{n}\selectfont\color{color_29791}t}
\put(372.3756,-309.011){\fontsize{11.991}{1}\usefont{T1}{cmr}{m}{n}\selectfont\color{color_29791}h}
\put(379.0546,-309.011){\fontsize{11.991}{1}\usefont{T1}{cmr}{m}{n}\selectfont\color{color_29791}e}
\put(385.7336,-309.011){\fontsize{11.991}{1}\usefont{T1}{cmr}{m}{n}\selectfont\color{color_29791}i}
\put(388.4196,-309.011){\fontsize{11.991}{1}\usefont{T1}{cmr}{m}{n}\selectfont\color{color_29791}r}
\put(392.3166,-309.011){\fontsize{11.991}{1}\usefont{T1}{cmr}{m}{n}\selectfont\color{color_29791} }
\put(395.6981,-309.011){\fontsize{11.991}{1}\usefont{T1}{cmr}{m}{n}\selectfont\color{color_29791}n}
\put(402.3771,-309.011){\fontsize{11.991}{1}\usefont{T1}{cmr}{m}{n}\selectfont\color{color_29791}a}
\put(409.0561,-309.011){\fontsize{11.991}{1}\usefont{T1}{cmr}{m}{n}\selectfont\color{color_29791}me}
\put(425.7235,-309.011){\fontsize{11.991}{1}\usefont{T1}{cmr}{m}{n}\selectfont\color{color_29791} }
\put(429.0211,-309.011){\fontsize{11.991}{1}\usefont{T1}{cmr}{m}{n}\selectfont\color{color_29791}e}
\put(435.808,-309.011){\fontsize{11.991}{1}\usefont{T1}{cmr}{m}{n}\selectfont\color{color_29791}g}
\put(442.4869,-309.011){\fontsize{11.991}{1}\usefont{T1}{cmr}{m}{n}\selectfont\color{color_29791}.}
\put(445.8684,-309.011){\fontsize{11.991}{1}\usefont{T1}{cmr}{m}{n}\selectfont\color{color_29791} }
\put(449.1659,-309.011){\fontsize{11.991}{1}\usefont{T1}{cmr}{m}{n}\selectfont\color{color_29791}h}
\put(455.8449,-309.011){\fontsize{11.991}{1}\usefont{T1}{cmr}{m}{n}\selectfont\color{color_29791}p}
\put(462.6318,-309.011){\fontsize{11.991}{1}\usefont{T1}{cmr}{m}{n}\selectfont\color{color_29791}c}
\put(468.5314,-309.011){\fontsize{11.991}{1}\usefont{T1}{cmr}{m}{n}\selectfont\color{color_29791}/}
\put(471.9128,-309.011){\fontsize{11.991}{1}\usefont{T1}{cmr}{m}{n}\selectfont\color{color_29791}n}
\put(478.5918,-309.011){\fontsize{11.991}{1}\usefont{T1}{cmr}{m}{n}\selectfont\color{color_29791}e}
\put(485.3787,-309.011){\fontsize{11.991}{1}\usefont{T1}{cmr}{m}{n}\selectfont\color{color_29791}t}
\put(488.6762,-309.011){\fontsize{11.991}{1}\usefont{T1}{cmr}{m}{n}\selectfont\color{color_29791}cd}
\put(501.3507,-309.011){\fontsize{11.991}{1}\usefont{T1}{cmr}{m}{n}\selectfont\color{color_29791}f}
\put(504.7322,-309.011){\fontsize{11.991}{1}\usefont{T1}{cmr}{m}{n}\selectfont\color{color_29791}/}
\put(508.0297,-309.011){\fontsize{11.991}{1}\usefont{T1}{cmr}{m}{n}\selectfont\color{color_29791}i}
\put(510.7157,-309.011){\fontsize{11.991}{1}\usefont{T1}{cmr}{m}{n}\selectfont\color{color_29791}n}
\put(517.3947,-309.011){\fontsize{11.991}{1}\usefont{T1}{cmr}{m}{n}\selectfont\color{color_29791}t}
\put(520.7762,-309.011){\fontsize{11.991}{1}\usefont{T1}{cmr}{m}{n}\selectfont\color{color_29791}e}
\put(527.4552,-309.011){\fontsize{11.991}{1}\usefont{T1}{cmr}{m}{n}\selectfont\color{color_29791}l}
\put(530.0573,-309.011){\fontsize{11.991}{1}\usefont{T1}{cmr}{m}{n}\selectfont\color{color_29791}/}
\put(533.4387,-309.011){\fontsize{11.991}{1}\usefont{T1}{cmr}{m}{n}\selectfont\color{color_29791}4}
\put(540.1177,-309.011){\fontsize{11.991}{1}\usefont{T1}{cmr}{m}{n}\selectfont\color{color_29791}.}
\put(543.4992,-309.011){\fontsize{11.991}{1}\usefont{T1}{cmr}{m}{n}\selectfont\color{color_29791}4}
\put(550.1782,-309.011){\fontsize{11.991}{1}\usefont{T1}{cmr}{m}{n}\selectfont\color{color_29791}.}
\put(553.5597,-309.011){\fontsize{11.991}{1}\usefont{T1}{cmr}{m}{n}\selectfont\color{color_29791}1}
\put(560.2387,-309.011){\fontsize{11.991}{1}\usefont{T1}{cmr}{m}{n}\selectfont\color{color_29791}.}
\put(563.5363,-309.011){\fontsize{11.991}{1}\usefont{T1}{cmr}{m}{n}\selectfont\color{color_29791}1}
\put(84.213,-337.102){\fontsize{11.991}{1}\usefont{T1}{cmr}{m}{n}\selectfont\color{color_29791}•}
\put(106.691,-337.811){\fontsize{11.991}{1}\usefont{T1}{cmr}{m}{n}\selectfont\color{color_29791}W}
\put(117.7827,-337.811){\fontsize{11.991}{1}\usefont{T1}{cmr}{m}{n}\selectfont\color{color_29791}o}
\put(124.4617,-337.811){\fontsize{11.991}{1}\usefont{T1}{cmr}{m}{n}\selectfont\color{color_29791}r}
\put(128.3587,-337.811){\fontsize{11.991}{1}\usefont{T1}{cmr}{m}{n}\selectfont\color{color_29791}ks }
\put(143.7312,-337.811){\fontsize{11.991}{1}\usefont{T1}{cmr}{m}{n}\selectfont\color{color_29791}b}
\put(150.4102,-337.811){\fontsize{11.991}{1}\usefont{T1}{cmr}{m}{n}\selectfont\color{color_29791}y }
\put(159.7032,-337.811){\fontsize{11.991}{1}\usefont{T1}{cmr}{m}{n}\selectfont\color{color_29791}a}
\put(166.4901,-337.811){\fontsize{11.991}{1}\usefont{T1}{cmr}{m}{n}\selectfont\color{color_29791}d}
\put(173.1691,-337.811){\fontsize{11.991}{1}\usefont{T1}{cmr}{m}{n}\selectfont\color{color_29791}j}
\put(175.7712,-337.811){\fontsize{11.991}{1}\usefont{T1}{cmr}{m}{n}\selectfont\color{color_29791}u}
\put(182.5581,-337.811){\fontsize{11.991}{1}\usefont{T1}{cmr}{m}{n}\selectfont\color{color_29791}s}
\put(188.4577,-337.811){\fontsize{11.991}{1}\usefont{T1}{cmr}{m}{n}\selectfont\color{color_29791}t}
\put(191.8391,-337.811){\fontsize{11.991}{1}\usefont{T1}{cmr}{m}{n}\selectfont\color{color_29791}i}
\put(194.5251,-337.811){\fontsize{11.991}{1}\usefont{T1}{cmr}{m}{n}\selectfont\color{color_29791}n}
\put(201.2041,-337.811){\fontsize{11.991}{1}\usefont{T1}{cmr}{m}{n}\selectfont\color{color_29791}g}
\put(207.8831,-337.811){\fontsize{11.991}{1}\usefont{T1}{cmr}{m}{n}\selectfont\color{color_29791} }
\put(211.2645,-337.811){\fontsize{11.991}{1}\usefont{T1}{cmr}{m}{n}\selectfont\color{color_29791}s}
\put(217.1641,-337.811){\fontsize{11.991}{1}\usefont{T1}{cmr}{m}{n}\selectfont\color{color_29791}h}
\put(223.951,-337.811){\fontsize{11.991}{1}\usefont{T1}{cmr}{m}{n}\selectfont\color{color_29791}e}
\put(230.63,-337.811){\fontsize{11.991}{1}\usefont{T1}{cmr}{m}{n}\selectfont\color{color_29791}l}
\put(233.232,-337.811){\fontsize{11.991}{1}\usefont{T1}{cmr}{m}{n}\selectfont\color{color_29791}l}
\put(235.918,-337.811){\fontsize{11.991}{1}\usefont{T1}{cmr}{m}{n}\selectfont\color{color_29791} }
\put(239.2156,-337.811){\fontsize{11.991}{1}\usefont{T1}{cmr}{m}{n}\selectfont\color{color_29791}va}
\put(251.998,-337.811){\fontsize{11.991}{1}\usefont{T1}{cmr}{m}{n}\selectfont\color{color_29791}r}
\put(255.895,-337.811){\fontsize{11.991}{1}\usefont{T1}{cmr}{m}{n}\selectfont\color{color_29791}i}
\put(258.581,-337.811){\fontsize{11.991}{1}\usefont{T1}{cmr}{m}{n}\selectfont\color{color_29791}a}
\put(265.26,-337.811){\fontsize{11.991}{1}\usefont{T1}{cmr}{m}{n}\selectfont\color{color_29791}b}
\put(271.939,-337.811){\fontsize{11.991}{1}\usefont{T1}{cmr}{m}{n}\selectfont\color{color_29791}l}
\put(274.541,-337.811){\fontsize{11.991}{1}\usefont{T1}{cmr}{m}{n}\selectfont\color{color_29791}e}
\put(281.22,-337.811){\fontsize{11.991}{1}\usefont{T1}{cmr}{m}{n}\selectfont\color{color_29791}s}
\put(287.2994,-337.811){\fontsize{11.991}{1}\usefont{T1}{cmr}{m}{n}\selectfont\color{color_29791} }
\put(290.597,-337.811){\fontsize{11.991}{1}\usefont{T1}{cmr}{m}{n}\selectfont\color{color_29791}e}
\put(297.2759,-337.811){\fontsize{11.991}{1}\usefont{T1}{cmr}{m}{n}\selectfont\color{color_29791}g}
\put(304.0628,-337.811){\fontsize{11.991}{1}\usefont{T1}{cmr}{m}{n}\selectfont\color{color_29791}.}
\put(307.809,-337.811){\fontsize{11.991}{1}\usefont{T1}{cmr}{m}{n}\selectfont\color{color_29791} }
\put(311.211,-337.811){\fontsize{11.991}{1}\usefont{T1}{cmr}{m}{n}\selectfont\color{color_29791}P}
\put(318.3097,-337.811){\fontsize{11.991}{1}\usefont{T1}{cmr}{m}{n}\selectfont\color{color_29791}A}
\put(325.4083,-337.811){\fontsize{11.991}{1}\usefont{T1}{cmr}{m}{n}\selectfont\color{color_29791}T}
\put(332.6029,-337.811){\fontsize{11.991}{1}\usefont{T1}{cmr}{m}{n}\selectfont\color{color_29791}H}
\put(341.2844,-337.811){\fontsize{11.991}{1}\usefont{T1}{cmr}{m}{n}\selectfont\color{color_29791},}
\put(344.688,-337.811){\fontsize{11.991}{1}\usefont{T1}{cmr}{m}{n}\selectfont\color{color_29791} }
\put(348.005,-337.811){\fontsize{11.991}{1}\usefont{T1}{cmr}{m}{n}\selectfont\color{color_29791}L}
\put(354.684,-337.811){\fontsize{11.991}{1}\usefont{T1}{cmr}{m}{n}\selectfont\color{color_29791}D}
\put(363.3655,-337.811){\fontsize{11.991}{1}\usefont{T1}{cmr}{m}{n}\selectfont\color{color_29791}\_}
\put(370.0445,-337.811){\fontsize{11.991}{1}\usefont{T1}{cmr}{m}{n}\selectfont\color{color_29791}L}
\put(376.7234,-337.811){\fontsize{11.991}{1}\usefont{T1}{cmr}{m}{n}\selectfont\color{color_29791}I}
\put(380.1049,-337.811){\fontsize{11.991}{1}\usefont{T1}{cmr}{m}{n}\selectfont\color{color_29791}BR}
\put(396.6885,-337.811){\fontsize{11.991}{1}\usefont{T1}{cmr}{m}{n}\selectfont\color{color_29791}A}
\put(404.7704,-337.811){\fontsize{11.991}{1}\usefont{T1}{cmr}{m}{n}\selectfont\color{color_29791}R}
\put(413.1641,-337.811){\fontsize{11.991}{1}\usefont{T1}{cmr}{m}{n}\selectfont\color{color_29791}Y\_}
\put(427.937,-337.811){\fontsize{11.991}{1}\usefont{T1}{cmr}{m}{n}\selectfont\color{color_29791}P}
\put(435.0357,-337.811){\fontsize{11.991}{1}\usefont{T1}{cmr}{m}{n}\selectfont\color{color_29791}A}
\put(442.1344,-337.811){\fontsize{11.991}{1}\usefont{T1}{cmr}{m}{n}\selectfont\color{color_29791}T}
\put(449.329,-337.811){\fontsize{11.991}{1}\usefont{T1}{cmr}{m}{n}\selectfont\color{color_29791}H}
\put(84.213,-365.902){\fontsize{11.991}{1}\usefont{T1}{cmr}{m}{n}\selectfont\color{color_29791}•}
\put(106.691,-366.611){\fontsize{11.991}{1}\usefont{T1}{cmr}{m}{n}\selectfont\color{color_29791}L}
\put(113.37,-366.611){\fontsize{11.991}{1}\usefont{T1}{cmr}{m}{n}\selectfont\color{color_29791}o}
\put(120.049,-366.611){\fontsize{11.991}{1}\usefont{T1}{cmr}{m}{n}\selectfont\color{color_29791}a}
\put(126.728,-366.611){\fontsize{11.991}{1}\usefont{T1}{cmr}{m}{n}\selectfont\color{color_29791}d}
\put(133.5149,-366.611){\fontsize{11.991}{1}\usefont{T1}{cmr}{m}{n}\selectfont\color{color_29791}e}
\put(140.1939,-366.611){\fontsize{11.991}{1}\usefont{T1}{cmr}{m}{n}\selectfont\color{color_29791}d}
\put(146.8728,-366.611){\fontsize{11.991}{1}\usefont{T1}{cmr}{m}{n}\selectfont\color{color_29791} }
\put(150.2543,-366.611){\fontsize{11.991}{1}\usefont{T1}{cmr}{m}{n}\selectfont\color{color_29791}m}
\put(160.1469,-366.611){\fontsize{11.991}{1}\usefont{T1}{cmr}{m}{n}\selectfont\color{color_29791}o}
\put(166.8259,-366.611){\fontsize{11.991}{1}\usefont{T1}{cmr}{m}{n}\selectfont\color{color_29791}d}
\put(173.6128,-366.611){\fontsize{11.991}{1}\usefont{T1}{cmr}{m}{n}\selectfont\color{color_29791}u}
\put(180.2918,-366.611){\fontsize{11.991}{1}\usefont{T1}{cmr}{m}{n}\selectfont\color{color_29791}l}
\put(182.8938,-366.611){\fontsize{11.991}{1}\usefont{T1}{cmr}{m}{n}\selectfont\color{color_29791}e}
\put(189.6807,-366.611){\fontsize{11.991}{1}\usefont{T1}{cmr}{m}{n}\selectfont\color{color_29791}s}
\put(195.5803,-366.611){\fontsize{11.991}{1}\usefont{T1}{cmr}{m}{n}\selectfont\color{color_29791} }
\put(198.9618,-366.611){\fontsize{11.991}{1}\usefont{T1}{cmr}{m}{n}\selectfont\color{color_29791}o}
\put(205.6407,-366.611){\fontsize{11.991}{1}\usefont{T1}{cmr}{m}{n}\selectfont\color{color_29791}n}
\put(212.4277,-366.611){\fontsize{11.991}{1}\usefont{T1}{cmr}{m}{n}\selectfont\color{color_29791}l}
\put(215.0297,-366.611){\fontsize{11.991}{1}\usefont{T1}{cmr}{m}{n}\selectfont\color{color_29791}y }
\put(224.3227,-366.611){\fontsize{11.991}{1}\usefont{T1}{cmr}{m}{n}\selectfont\color{color_29791}a}
\put(231.1096,-366.611){\fontsize{11.991}{1}\usefont{T1}{cmr}{m}{n}\selectfont\color{color_29791}f}
\put(234.2033,-366.611){\fontsize{11.991}{1}\usefont{T1}{cmr}{m}{n}\selectfont\color{color_29791}f}
\put(237.5848,-366.611){\fontsize{11.991}{1}\usefont{T1}{cmr}{m}{n}\selectfont\color{color_29791}e}
\put(244.2638,-366.611){\fontsize{11.991}{1}\usefont{T1}{cmr}{m}{n}\selectfont\color{color_29791}ct}
\put(253.5568,-366.611){\fontsize{11.991}{1}\usefont{T1}{cmr}{m}{n}\selectfont\color{color_29791} }
\put(256.9383,-366.611){\fontsize{11.991}{1}\usefont{T1}{cmr}{m}{n}\selectfont\color{color_29791}t}
\put(260.2358,-366.611){\fontsize{11.991}{1}\usefont{T1}{cmr}{m}{n}\selectfont\color{color_29791}h}
\put(267.0227,-366.611){\fontsize{11.991}{1}\usefont{T1}{cmr}{m}{n}\selectfont\color{color_29791}e}
\put(273.7017,-366.611){\fontsize{11.991}{1}\usefont{T1}{cmr}{m}{n}\selectfont\color{color_29791} }
\put(277.0831,-366.611){\fontsize{11.991}{1}\usefont{T1}{cmr}{m}{n}\selectfont\color{color_29791}t}
\put(280.3806,-366.611){\fontsize{11.991}{1}\usefont{T1}{cmr}{m}{n}\selectfont\color{color_29791}e}
\put(287.0596,-366.611){\fontsize{11.991}{1}\usefont{T1}{cmr}{m}{n}\selectfont\color{color_29791}rmi}
\put(303.6432,-366.611){\fontsize{11.991}{1}\usefont{T1}{cmr}{m}{n}\selectfont\color{color_29791}n}
\put(310.3222,-366.611){\fontsize{11.991}{1}\usefont{T1}{cmr}{m}{n}\selectfont\color{color_29791}a}
\put(317.0012,-366.611){\fontsize{11.991}{1}\usefont{T1}{cmr}{m}{n}\selectfont\color{color_29791}l}
\put(319.6871,-366.611){\fontsize{11.991}{1}\usefont{T1}{cmr}{m}{n}\selectfont\color{color_29791} }
\put(322.9846,-366.611){\fontsize{11.991}{1}\usefont{T1}{cmr}{m}{n}\selectfont\color{color_29791}yo}
\put(335.6591,-366.611){\fontsize{11.991}{1}\usefont{T1}{cmr}{m}{n}\selectfont\color{color_29791}u}
\put(342.446,-366.611){\fontsize{11.991}{1}\usefont{T1}{cmr}{m}{n}\selectfont\color{color_29791}r}
\put(346.3431,-366.611){\fontsize{11.991}{1}\usefont{T1}{cmr}{m}{n}\selectfont\color{color_29791} }
\put(349.7246,-366.611){\fontsize{11.991}{1}\usefont{T1}{cmr}{m}{n}\selectfont\color{color_29791}l}
\put(352.3266,-366.611){\fontsize{11.991}{1}\usefont{T1}{cmr}{m}{n}\selectfont\color{color_29791}o}
\put(359.1135,-366.611){\fontsize{11.991}{1}\usefont{T1}{cmr}{m}{n}\selectfont\color{color_29791}a}
\put(365.7925,-366.611){\fontsize{11.991}{1}\usefont{T1}{cmr}{m}{n}\selectfont\color{color_29791}d}
\put(372.4715,-366.611){\fontsize{11.991}{1}\usefont{T1}{cmr}{m}{n}\selectfont\color{color_29791}e}
\put(379.1505,-366.611){\fontsize{11.991}{1}\usefont{T1}{cmr}{m}{n}\selectfont\color{color_29791}d}
\put(385.8295,-366.611){\fontsize{11.991}{1}\usefont{T1}{cmr}{m}{n}\selectfont\color{color_29791} }
\put(389.2109,-366.611){\fontsize{11.991}{1}\usefont{T1}{cmr}{m}{n}\selectfont\color{color_29791}t}
\put(392.5085,-366.611){\fontsize{11.991}{1}\usefont{T1}{cmr}{m}{n}\selectfont\color{color_29791}h}
\put(399.2953,-366.611){\fontsize{11.991}{1}\usefont{T1}{cmr}{m}{n}\selectfont\color{color_29791}e}
\put(405.9743,-366.611){\fontsize{11.991}{1}\usefont{T1}{cmr}{m}{n}\selectfont\color{color_29791}m }
\put(419.2604,-366.611){\fontsize{11.991}{1}\usefont{T1}{cmr}{m}{n}\selectfont\color{color_29791}i}
\put(421.8624,-366.611){\fontsize{11.991}{1}\usefont{T1}{cmr}{m}{n}\selectfont\color{color_29791}n}
\end{picture}
\newpage
\begin{tikzpicture}[overlay]
\path(0pt,0pt);
\filldraw[color_283006][even odd rule]
(-15pt, 10pt) -- (704.972pt, 10pt)
 -- (704.972pt, 10pt)
 -- (704.972pt, -529.972pt)
 -- (704.972pt, -529.972pt)
 -- (-15pt, -529.972pt)
 -- (-15pt, -529.972pt)
 -- (-15pt, 10pt) -- cycle
;
\end{tikzpicture}
\begin{picture}(-5,0)(2.5,0)
\put(-15,-529.972){\includegraphics[width=720pt,height=540pt]{latexImage_43fb9acc39826b86b6493904abaeacaa.png}}
\put(48.411,-60.69601){\fontsize{21.997}{1}\usefont{T1}{cmr}{b}{n}\selectfont\color{color_29791}S}
\put(63.083,-60.69601){\fontsize{21.997}{1}\usefont{T1}{cmr}{b}{n}\selectfont\color{color_29791}o}
\put(76.47917,-60.69601){\fontsize{21.997}{1}\usefont{T1}{cmr}{b}{n}\selectfont\color{color_29791}f}
\put(83.78217,-60.69601){\fontsize{21.997}{1}\usefont{T1}{cmr}{b}{n}\selectfont\color{color_29791}t}
\put(91.15117,-60.69601){\fontsize{21.997}{1}\usefont{T1}{cmr}{b}{n}\selectfont\color{color_29791}w}
\put(108.1548,-60.69601){\fontsize{21.997}{1}\usefont{T1}{cmr}{b}{n}\selectfont\color{color_29791}a}
\put(120.3632,-60.69601){\fontsize{21.997}{1}\usefont{T1}{cmr}{b}{n}\selectfont\color{color_29791}r}
\put(128.942,-60.69601){\fontsize{21.997}{1}\usefont{T1}{cmr}{b}{n}\selectfont\color{color_29791}e}
\put(141.1503,-60.69601){\fontsize{21.997}{1}\usefont{T1}{cmr}{b}{n}\selectfont\color{color_29791} (}
\put(154.6125,-60.69601){\fontsize{21.997}{1}\usefont{T1}{cmr}{b}{n}\selectfont\color{color_29791}c}
\put(166.8208,-60.69601){\fontsize{21.997}{1}\usefont{T1}{cmr}{b}{n}\selectfont\color{color_29791}o}
\put(180.217,-60.69601){\fontsize{21.997}{1}\usefont{T1}{cmr}{b}{n}\selectfont\color{color_29791}n}
\put(193.7012,-60.69601){\fontsize{21.997}{1}\usefont{T1}{cmr}{b}{n}\selectfont\color{color_29791}t}
\put(201.0042,-60.69601){\fontsize{21.997}{1}\usefont{T1}{cmr}{b}{n}\selectfont\color{color_29791}in}
\put(220.4935,-60.69601){\fontsize{21.997}{1}\usefont{T1}{cmr}{b}{n}\selectfont\color{color_29791}u}
\put(233.9777,-60.69601){\fontsize{21.997}{1}\usefont{T1}{cmr}{b}{n}\selectfont\color{color_29791}e}
\put(246.186,-60.69601){\fontsize{21.997}{1}\usefont{T1}{cmr}{b}{n}\selectfont\color{color_29791}d}
\put(259.6702,-60.69601){\fontsize{21.997}{1}\usefont{T1}{cmr}{b}{n}\selectfont\color{color_29791})}
\put(48.213,-116.397){\fontsize{15.987}{1}\usefont{T1}{cmr}{m}{n}\selectfont\color{color_29791}M}
\put(61.4982,-116.397){\fontsize{15.987}{1}\usefont{T1}{cmr}{m}{n}\selectfont\color{color_29791}odul}
\put(91.74561,-116.397){\fontsize{15.987}{1}\usefont{T1}{cmr}{m}{n}\selectfont\color{color_29791}e}
\put(100.5385,-116.397){\fontsize{15.987}{1}\usefont{T1}{cmr}{m}{n}\selectfont\color{color_29791}s}
\put(84.213,-133.291){\fontsize{11.991}{1}\usefont{T1}{cmr}{m}{n}\selectfont\color{color_29791}•}
\put(106.691,-134){\fontsize{11.991}{1}\usefont{T1}{cmr}{m}{n}\selectfont\color{color_29791}U}
\put(115.2885,-134){\fontsize{11.991}{1}\usefont{T1}{cmr}{m}{n}\selectfont\color{color_29791}se}
\put(128.071,-134){\fontsize{11.991}{1}\usefont{T1}{cmr}{m}{n}\selectfont\color{color_29791}f}
\put(131.3685,-134){\fontsize{11.991}{1}\usefont{T1}{cmr}{m}{n}\selectfont\color{color_29791}u}
\put(138.0475,-134){\fontsize{11.991}{1}\usefont{T1}{cmr}{m}{n}\selectfont\color{color_29791}l}
\put(140.7335,-134){\fontsize{11.991}{1}\usefont{T1}{cmr}{m}{n}\selectfont\color{color_29791} }
\put(144.031,-134){\fontsize{11.991}{1}\usefont{T1}{cmr}{m}{n}\selectfont\color{color_29791}mo}
\put(160.6985,-134){\fontsize{11.991}{1}\usefont{T1}{cmr}{m}{n}\selectfont\color{color_29791}d}
\put(167.3775,-134){\fontsize{11.991}{1}\usefont{T1}{cmr}{m}{n}\selectfont\color{color_29791}u}
\put(174.0565,-134){\fontsize{11.991}{1}\usefont{T1}{cmr}{m}{n}\selectfont\color{color_29791}l}
\put(176.7424,-134){\fontsize{11.991}{1}\usefont{T1}{cmr}{m}{n}\selectfont\color{color_29791}e}
\put(183.4214,-134){\fontsize{11.991}{1}\usefont{T1}{cmr}{m}{n}\selectfont\color{color_29791} }
\put(186.8029,-134){\fontsize{11.991}{1}\usefont{T1}{cmr}{m}{n}\selectfont\color{color_29791}co}
\put(199.4774,-134){\fontsize{11.991}{1}\usefont{T1}{cmr}{m}{n}\selectfont\color{color_29791}m}
\put(209.3699,-134){\fontsize{11.991}{1}\usefont{T1}{cmr}{m}{n}\selectfont\color{color_29791}ma}
\put(226.0374,-134){\fontsize{11.991}{1}\usefont{T1}{cmr}{m}{n}\selectfont\color{color_29791}n}
\put(232.8244,-134){\fontsize{11.991}{1}\usefont{T1}{cmr}{m}{n}\selectfont\color{color_29791}d}
\put(239.5033,-134){\fontsize{11.991}{1}\usefont{T1}{cmr}{m}{n}\selectfont\color{color_29791}s:}
\put(249.104,-134){\fontsize{11.991}{1}\usefont{T1}{cmr}{m}{n}\selectfont\color{color_29791} }
\put(120.213,-160.391){\fontsize{10.006}{1}\usefont{T1}{cmr}{m}{n}\selectfont\color{color_29791}module avail}
\put(192.213,-160.391){\fontsize{10.006}{1}\usefont{T1}{cmr}{m}{n}\selectfont\color{color_29791} }
\put(198.194,-160.391){\fontsize{10.006}{1}\usefont{T1}{cmr}{m}{n}\selectfont\color{color_29791} }
\put(204.203,-160.391){\fontsize{10.006}{1}\usefont{T1}{cmr}{m}{n}\selectfont\color{color_29791} }
\put(210.213,-160.391){\fontsize{10.006}{1}\usefont{T1}{cmr}{m}{n}\selectfont\color{color_29791} }
\put(216.194,-160.391){\fontsize{10.006}{1}\usefont{T1}{cmr}{m}{n}\selectfont\color{color_29791} }
\put(222.203,-160.391){\fontsize{10.006}{1}\usefont{T1}{cmr}{m}{n}\selectfont\color{color_29791} }
\put(228.213,-160.391){\fontsize{10.006}{1}\usefont{T1}{cmr}{m}{n}\selectfont\color{color_29791} }
\put(234.194,-160.391){\fontsize{10.006}{1}\usefont{T1}{cmr}{m}{n}\selectfont\color{color_29791} }
\put(240.203,-160.391){\fontsize{10.006}{1}\usefont{T1}{cmr}{m}{n}\selectfont\color{color_29791} }
\put(246.213,-160.391){\fontsize{10.006}{1}\usefont{T1}{cmr}{m}{n}\selectfont\color{color_29791} }
\put(252.194,-160.391){\fontsize{10.006}{1}\usefont{T1}{cmr}{m}{n}\selectfont\color{color_29791} }
\put(258.203,-160.391){\fontsize{10.006}{1}\usefont{T1}{cmr}{m}{n}\selectfont\color{color_29791} }
\put(264.213,-160.391){\fontsize{10.006}{1}\usefont{T1}{cmr}{m}{n}\selectfont\color{color_29791} }
\put(270.194,-160.391){\fontsize{10.006}{1}\usefont{T1}{cmr}{m}{n}\selectfont\color{color_29791} }
\put(276.203,-160.391){\fontsize{10.006}{1}\usefont{T1}{cmr}{m}{n}\selectfont\color{color_29791} }
\put(282.213,-160.391){\fontsize{10.006}{1}\usefont{T1}{cmr}{m}{n}\selectfont\color{color_29791} }
\put(288.194,-160.391){\fontsize{10.006}{1}\usefont{T1}{cmr}{m}{n}\selectfont\color{color_29791} }
\put(294.203,-160.391){\fontsize{10.006}{1}\usefont{T1}{cmr}{m}{n}\selectfont\color{color_29791} }
\put(300.213,-160.391){\fontsize{10.006}{1}\usefont{T1}{cmr}{m}{n}\selectfont\color{color_29791} }
\put(306.194,-160.391){\fontsize{10.006}{1}\usefont{T1}{cmr}{m}{n}\selectfont\color{color_29791} }
\put(312.203,-160.391){\fontsize{10.006}{1}\usefont{T1}{cmr}{m}{n}\selectfont\color{color_29791} }
\put(318.213,-160.391){\fontsize{10.006}{1}\usefont{T1}{cmr}{m}{n}\selectfont\color{color_29791} }
\put(324.194,-160.391){\fontsize{10.006}{1}\usefont{T1}{cmr}{m}{n}\selectfont\color{color_29791} }
\put(330.203,-160.391){\fontsize{10.006}{1}\usefont{T1}{cmr}{m}{n}\selectfont\color{color_29791} }
\put(336.213,-160.391){\fontsize{10.006}{1}\usefont{T1}{cmr}{m}{n}\selectfont\color{color_29791} }
\put(342.194,-160.391){\fontsize{10.006}{1}\usefont{T1}{cmr}{m}{n}\selectfont\color{color_29791} }
\put(348.203,-160.391){\fontsize{10.006}{1}\usefont{T1}{cmr}{m}{n}\selectfont\color{color_29791} }
\put(354.213,-160.391){\fontsize{10.006}{1}\usefont{T1}{cmr}{m}{n}\selectfont\color{color_29791} module display name/version}
\put(120.213,-186.413){\fontsize{10.006}{1}\usefont{T1}{cmr}{m}{n}\selectfont\color{color_29791}module load name/version}
\put(264.213,-186.413){\fontsize{10.006}{1}\usefont{T1}{cmr}{m}{n}\selectfont\color{color_29791} }
\put(270.194,-186.413){\fontsize{10.006}{1}\usefont{T1}{cmr}{m}{n}\selectfont\color{color_29791} }
\put(276.203,-186.413){\fontsize{10.006}{1}\usefont{T1}{cmr}{m}{n}\selectfont\color{color_29791} }
\put(282.213,-186.413){\fontsize{10.006}{1}\usefont{T1}{cmr}{m}{n}\selectfont\color{color_29791} }
\put(288.194,-186.413){\fontsize{10.006}{1}\usefont{T1}{cmr}{m}{n}\selectfont\color{color_29791} }
\put(294.203,-186.413){\fontsize{10.006}{1}\usefont{T1}{cmr}{m}{n}\selectfont\color{color_29791} }
\put(300.213,-186.413){\fontsize{10.006}{1}\usefont{T1}{cmr}{m}{n}\selectfont\color{color_29791} }
\put(306.194,-186.413){\fontsize{10.006}{1}\usefont{T1}{cmr}{m}{n}\selectfont\color{color_29791} }
\put(312.203,-186.413){\fontsize{10.006}{1}\usefont{T1}{cmr}{m}{n}\selectfont\color{color_29791} }
\put(318.213,-186.413){\fontsize{10.006}{1}\usefont{T1}{cmr}{m}{n}\selectfont\color{color_29791} }
\put(324.194,-186.413){\fontsize{10.006}{1}\usefont{T1}{cmr}{m}{n}\selectfont\color{color_29791} }
\put(330.203,-186.413){\fontsize{10.006}{1}\usefont{T1}{cmr}{m}{n}\selectfont\color{color_29791} }
\put(336.213,-186.413){\fontsize{10.006}{1}\usefont{T1}{cmr}{m}{n}\selectfont\color{color_29791} }
\put(342.194,-186.413){\fontsize{10.006}{1}\usefont{T1}{cmr}{m}{n}\selectfont\color{color_29791} }
\put(348.203,-186.413){\fontsize{10.006}{1}\usefont{T1}{cmr}{m}{n}\selectfont\color{color_29791} }
\put(354.213,-186.413){\fontsize{10.006}{1}\usefont{T1}{cmr}{m}{n}\selectfont\color{color_29791} module list}
\put(426.213,-186.413){\fontsize{10.006}{1}\usefont{T1}{cmr}{m}{n}\selectfont\color{color_29791} }
\put(120.213,-212.406){\fontsize{10.006}{1}\usefont{T1}{cmr}{m}{n}\selectfont\color{color_29791}module unload name/version}
\put(276.203,-212.406){\fontsize{10.006}{1}\usefont{T1}{cmr}{m}{n}\selectfont\color{color_29791} }
\put(282.213,-212.406){\fontsize{10.006}{1}\usefont{T1}{cmr}{m}{n}\selectfont\color{color_29791} }
\put(288.194,-212.406){\fontsize{10.006}{1}\usefont{T1}{cmr}{m}{n}\selectfont\color{color_29791} }
\put(294.203,-212.406){\fontsize{10.006}{1}\usefont{T1}{cmr}{m}{n}\selectfont\color{color_29791} }
\put(300.213,-212.406){\fontsize{10.006}{1}\usefont{T1}{cmr}{m}{n}\selectfont\color{color_29791} }
\put(306.194,-212.406){\fontsize{10.006}{1}\usefont{T1}{cmr}{m}{n}\selectfont\color{color_29791} }
\put(312.203,-212.406){\fontsize{10.006}{1}\usefont{T1}{cmr}{m}{n}\selectfont\color{color_29791} }
\put(318.213,-212.406){\fontsize{10.006}{1}\usefont{T1}{cmr}{m}{n}\selectfont\color{color_29791} }
\put(324.194,-212.406){\fontsize{10.006}{1}\usefont{T1}{cmr}{m}{n}\selectfont\color{color_29791} }
\put(330.203,-212.406){\fontsize{10.006}{1}\usefont{T1}{cmr}{m}{n}\selectfont\color{color_29791} }
\put(336.213,-212.406){\fontsize{10.006}{1}\usefont{T1}{cmr}{m}{n}\selectfont\color{color_29791} }
\put(342.194,-212.406){\fontsize{10.006}{1}\usefont{T1}{cmr}{m}{n}\selectfont\color{color_29791} }
\put(348.203,-212.406){\fontsize{10.006}{1}\usefont{T1}{cmr}{m}{n}\selectfont\color{color_29791} }
\put(354.213,-212.406){\fontsize{10.006}{1}\usefont{T1}{cmr}{m}{n}\selectfont\color{color_29791} module purge}
\put(84.213,-240.101){\fontsize{11.991}{1}\usefont{T1}{cmr}{m}{n}\selectfont\color{color_29791}•}
\put(106.691,-240.809){\fontsize{11.991}{1}\usefont{T1}{cmr}{m}{n}\selectfont\color{color_29791}C}
\put(115.2885,-240.809){\fontsize{11.991}{1}\usefont{T1}{cmr}{m}{n}\selectfont\color{color_29791}o}
\put(122.0755,-240.809){\fontsize{11.991}{1}\usefont{T1}{cmr}{m}{n}\selectfont\color{color_29791}m}
\put(131.968,-240.809){\fontsize{11.991}{1}\usefont{T1}{cmr}{m}{n}\selectfont\color{color_29791}mo}
\put(148.6355,-240.809){\fontsize{11.991}{1}\usefont{T1}{cmr}{m}{n}\selectfont\color{color_29791}n}
\put(155.3145,-240.809){\fontsize{11.991}{1}\usefont{T1}{cmr}{m}{n}\selectfont\color{color_29791} }
\put(158.696,-240.809){\fontsize{11.991}{1}\usefont{T1}{cmr}{m}{n}\selectfont\color{color_29791}m}
\put(168.5886,-240.809){\fontsize{11.991}{1}\usefont{T1}{cmr}{m}{n}\selectfont\color{color_29791}i}
\put(171.2746,-240.809){\fontsize{11.991}{1}\usefont{T1}{cmr}{m}{n}\selectfont\color{color_29791}st}
\put(180.5676,-240.809){\fontsize{11.991}{1}\usefont{T1}{cmr}{m}{n}\selectfont\color{color_29791}a}
\put(187.3545,-240.809){\fontsize{11.991}{1}\usefont{T1}{cmr}{m}{n}\selectfont\color{color_29791}ke}
\put(200.029,-240.809){\fontsize{11.991}{1}\usefont{T1}{cmr}{m}{n}\selectfont\color{color_29791}s}
\put(129.198,-268.787){\fontsize{11.991}{1}\usefont{T1}{cmr}{m}{n}\selectfont\color{color_29791}•}
\put(142.691,-269.496){\fontsize{11.991}{1}\usefont{T1}{cmr}{m}{n}\selectfont\color{color_29791}F}
\put(149.9815,-269.496){\fontsize{11.991}{1}\usefont{T1}{cmr}{m}{n}\selectfont\color{color_29791}o}
\put(156.6605,-269.496){\fontsize{11.991}{1}\usefont{T1}{cmr}{m}{n}\selectfont\color{color_29791}rg}
\put(167.3325,-269.496){\fontsize{11.991}{1}\usefont{T1}{cmr}{m}{n}\selectfont\color{color_29791}e}
\put(174.0115,-269.496){\fontsize{11.991}{1}\usefont{T1}{cmr}{m}{n}\selectfont\color{color_29791}t}
\put(177.3929,-269.496){\fontsize{11.991}{1}\usefont{T1}{cmr}{m}{n}\selectfont\color{color_29791}t}
\put(180.6905,-269.496){\fontsize{11.991}{1}\usefont{T1}{cmr}{m}{n}\selectfont\color{color_29791}i}
\put(183.3765,-269.496){\fontsize{11.991}{1}\usefont{T1}{cmr}{m}{n}\selectfont\color{color_29791}n}
\put(190.0555,-269.496){\fontsize{11.991}{1}\usefont{T1}{cmr}{m}{n}\selectfont\color{color_29791}g}
\put(196.7344,-269.496){\fontsize{11.991}{1}\usefont{T1}{cmr}{m}{n}\selectfont\color{color_29791} }
\put(200.1159,-269.496){\fontsize{11.991}{1}\usefont{T1}{cmr}{m}{n}\selectfont\color{color_29791}t}
\put(203.4974,-269.496){\fontsize{11.991}{1}\usefont{T1}{cmr}{m}{n}\selectfont\color{color_29791}o}
\put(210.411,-269.496){\fontsize{11.991}{1}\usefont{T1}{cmr}{m}{n}\selectfont\color{color_29791} }
\put(213.699,-269.496){\fontsize{11.991}{1}\usefont{T1}{cmr}{m}{n}\selectfont\color{color_29791}u}
\put(220.378,-269.496){\fontsize{11.991}{1}\usefont{T1}{cmr}{m}{n}\selectfont\color{color_29791}s}
\put(226.4574,-269.496){\fontsize{11.991}{1}\usefont{T1}{cmr}{m}{n}\selectfont\color{color_29791}e}
\put(233.1364,-269.496){\fontsize{11.991}{1}\usefont{T1}{cmr}{m}{n}\selectfont\color{color_29791} }
\put(236.4339,-269.496){\fontsize{11.991}{1}\usefont{T1}{cmr}{m}{n}\selectfont\color{color_29791}h}
\put(243.2209,-269.496){\fontsize{11.991}{1}\usefont{T1}{cmr}{m}{n}\selectfont\color{color_29791}p}
\put(249.8998,-269.496){\fontsize{11.991}{1}\usefont{T1}{cmr}{m}{n}\selectfont\color{color_29791}c }
\put(259.1929,-269.496){\fontsize{11.991}{1}\usefont{T1}{cmr}{m}{n}\selectfont\color{color_29791}mo}
\put(275.8604,-269.496){\fontsize{11.991}{1}\usefont{T1}{cmr}{m}{n}\selectfont\color{color_29791}d}
\put(282.5393,-269.496){\fontsize{11.991}{1}\usefont{T1}{cmr}{m}{n}\selectfont\color{color_29791}u}
\put(289.2183,-269.496){\fontsize{11.991}{1}\usefont{T1}{cmr}{m}{n}\selectfont\color{color_29791}l}
\put(291.9043,-269.496){\fontsize{11.991}{1}\usefont{T1}{cmr}{m}{n}\selectfont\color{color_29791}e}
\put(298.5833,-269.496){\fontsize{11.991}{1}\usefont{T1}{cmr}{m}{n}\selectfont\color{color_29791}s }
\put(307.9602,-269.496){\fontsize{11.991}{1}\usefont{T1}{cmr}{m}{n}\selectfont\color{color_29791}o}
\put(314.6392,-269.496){\fontsize{11.991}{1}\usefont{T1}{cmr}{m}{n}\selectfont\color{color_29791}n}
\put(321.3182,-269.496){\fontsize{11.991}{1}\usefont{T1}{cmr}{m}{n}\selectfont\color{color_29791} }
\put(324.6997,-269.496){\fontsize{11.991}{1}\usefont{T1}{cmr}{m}{n}\selectfont\color{color_29791}n}
\put(331.3787,-269.496){\fontsize{11.991}{1}\usefont{T1}{cmr}{m}{n}\selectfont\color{color_29791}o}
\put(338.0576,-269.496){\fontsize{11.991}{1}\usefont{T1}{cmr}{m}{n}\selectfont\color{color_29791}d}
\put(344.7366,-269.496){\fontsize{11.991}{1}\usefont{T1}{cmr}{m}{n}\selectfont\color{color_29791}e}
\put(351.5235,-269.496){\fontsize{11.991}{1}\usefont{T1}{cmr}{m}{n}\selectfont\color{color_29791}s}
\put(129.198,-297.587){\fontsize{11.991}{1}\usefont{T1}{cmr}{m}{n}\selectfont\color{color_29791}•}
\put(142.691,-298.296){\fontsize{11.991}{1}\usefont{T1}{cmr}{m}{n}\selectfont\color{color_29791}M}
\put(152.5836,-298.296){\fontsize{11.991}{1}\usefont{T1}{cmr}{m}{n}\selectfont\color{color_29791}i}
\put(155.2696,-298.296){\fontsize{11.991}{1}\usefont{T1}{cmr}{m}{n}\selectfont\color{color_29791}xi}
\put(163.8671,-298.296){\fontsize{11.991}{1}\usefont{T1}{cmr}{m}{n}\selectfont\color{color_29791}n}
\put(170.5461,-298.296){\fontsize{11.991}{1}\usefont{T1}{cmr}{m}{n}\selectfont\color{color_29791}g}
\put(177.333,-298.296){\fontsize{11.991}{1}\usefont{T1}{cmr}{m}{n}\selectfont\color{color_29791} }
\put(180.6305,-298.296){\fontsize{11.991}{1}\usefont{T1}{cmr}{m}{n}\selectfont\color{color_29791}mo}
\put(197.298,-298.296){\fontsize{11.991}{1}\usefont{T1}{cmr}{m}{n}\selectfont\color{color_29791}d}
\put(203.977,-298.296){\fontsize{11.991}{1}\usefont{T1}{cmr}{m}{n}\selectfont\color{color_29791}u}
\put(210.7639,-298.296){\fontsize{11.991}{1}\usefont{T1}{cmr}{m}{n}\selectfont\color{color_29791}l}
\put(213.366,-298.296){\fontsize{11.991}{1}\usefont{T1}{cmr}{m}{n}\selectfont\color{color_29791}e}
\put(220.045,-298.296){\fontsize{11.991}{1}\usefont{T1}{cmr}{m}{n}\selectfont\color{color_29791}s }
\put(229.4219,-298.296){\fontsize{11.991}{1}\usefont{T1}{cmr}{m}{n}\selectfont\color{color_29791}cr}
\put(239.3145,-298.296){\fontsize{11.991}{1}\usefont{T1}{cmr}{m}{n}\selectfont\color{color_29791}e}
\put(245.9935,-298.296){\fontsize{11.991}{1}\usefont{T1}{cmr}{m}{n}\selectfont\color{color_29791}a}
\put(252.7804,-298.296){\fontsize{11.991}{1}\usefont{T1}{cmr}{m}{n}\selectfont\color{color_29791}t}
\put(256.0779,-298.296){\fontsize{11.991}{1}\usefont{T1}{cmr}{m}{n}\selectfont\color{color_29791}e}
\put(262.8648,-298.296){\fontsize{11.991}{1}\usefont{T1}{cmr}{m}{n}\selectfont\color{color_29791}d}
\put(269.5438,-298.296){\fontsize{11.991}{1}\usefont{T1}{cmr}{m}{n}\selectfont\color{color_29791} }
\put(272.8413,-298.296){\fontsize{11.991}{1}\usefont{T1}{cmr}{m}{n}\selectfont\color{color_29791}u}
\put(279.6282,-298.296){\fontsize{11.991}{1}\usefont{T1}{cmr}{m}{n}\selectfont\color{color_29791}s}
\put(285.5278,-298.296){\fontsize{11.991}{1}\usefont{T1}{cmr}{m}{n}\selectfont\color{color_29791}i}
\put(288.2137,-298.296){\fontsize{11.991}{1}\usefont{T1}{cmr}{m}{n}\selectfont\color{color_29791}n}
\put(294.8927,-298.296){\fontsize{11.991}{1}\usefont{T1}{cmr}{m}{n}\selectfont\color{color_29791}g}
\put(301.5717,-298.296){\fontsize{11.991}{1}\usefont{T1}{cmr}{m}{n}\selectfont\color{color_29791} }
\put(304.9532,-298.296){\fontsize{11.991}{1}\usefont{T1}{cmr}{m}{n}\selectfont\color{color_29791}d}
\put(311.6322,-298.296){\fontsize{11.991}{1}\usefont{T1}{cmr}{m}{n}\selectfont\color{color_29791}i}
\put(314.2342,-298.296){\fontsize{11.991}{1}\usefont{T1}{cmr}{m}{n}\selectfont\color{color_29791}f}
\put(317.4358,-298.296){\fontsize{11.991}{1}\usefont{T1}{cmr}{m}{n}\selectfont\color{color_29791}f}
\put(320.7333,-298.296){\fontsize{11.991}{1}\usefont{T1}{cmr}{m}{n}\selectfont\color{color_29791}e}
\put(327.5202,-298.296){\fontsize{11.991}{1}\usefont{T1}{cmr}{m}{n}\selectfont\color{color_29791}r}
\put(331.4173,-298.296){\fontsize{11.991}{1}\usefont{T1}{cmr}{m}{n}\selectfont\color{color_29791}e}
\put(338.0963,-298.296){\fontsize{11.991}{1}\usefont{T1}{cmr}{m}{n}\selectfont\color{color_29791}n}
\put(344.8832,-298.296){\fontsize{11.991}{1}\usefont{T1}{cmr}{m}{n}\selectfont\color{color_29791}t}
\put(348.1807,-298.296){\fontsize{11.991}{1}\usefont{T1}{cmr}{m}{n}\selectfont\color{color_29791} }
\put(351.5622,-298.296){\fontsize{11.991}{1}\usefont{T1}{cmr}{m}{n}\selectfont\color{color_29791}co}
\put(364.2367,-298.296){\fontsize{11.991}{1}\usefont{T1}{cmr}{m}{n}\selectfont\color{color_29791}mp}
\put(380.9041,-298.296){\fontsize{11.991}{1}\usefont{T1}{cmr}{m}{n}\selectfont\color{color_29791}i}
\put(383.5062,-298.296){\fontsize{11.991}{1}\usefont{T1}{cmr}{m}{n}\selectfont\color{color_29791}l}
\put(386.1082,-298.296){\fontsize{11.991}{1}\usefont{T1}{cmr}{m}{n}\selectfont\color{color_29791}e}
\put(392.8951,-298.296){\fontsize{11.991}{1}\usefont{T1}{cmr}{m}{n}\selectfont\color{color_29791}r}
\put(396.7922,-298.296){\fontsize{11.991}{1}\usefont{T1}{cmr}{m}{n}\selectfont\color{color_29791}s}
\put(129.198,-326.387){\fontsize{11.991}{1}\usefont{T1}{cmr}{m}{n}\selectfont\color{color_29791}•}
\put(142.691,-327.096){\fontsize{11.991}{1}\usefont{T1}{cmr}{m}{n}\selectfont\color{color_29791}L}
\put(149.37,-327.096){\fontsize{11.991}{1}\usefont{T1}{cmr}{m}{n}\selectfont\color{color_29791}o}
\put(156.049,-327.096){\fontsize{11.991}{1}\usefont{T1}{cmr}{m}{n}\selectfont\color{color_29791}a}
\put(162.728,-327.096){\fontsize{11.991}{1}\usefont{T1}{cmr}{m}{n}\selectfont\color{color_29791}d}
\put(169.5149,-327.096){\fontsize{11.991}{1}\usefont{T1}{cmr}{m}{n}\selectfont\color{color_29791}i}
\put(172.1169,-327.096){\fontsize{11.991}{1}\usefont{T1}{cmr}{m}{n}\selectfont\color{color_29791}n}
\put(178.7959,-327.096){\fontsize{11.991}{1}\usefont{T1}{cmr}{m}{n}\selectfont\color{color_29791}g}
\put(185.5828,-327.096){\fontsize{11.991}{1}\usefont{T1}{cmr}{m}{n}\selectfont\color{color_29791} }
\put(188.8803,-327.096){\fontsize{11.991}{1}\usefont{T1}{cmr}{m}{n}\selectfont\color{color_29791}cl}
\put(197.5618,-327.096){\fontsize{11.991}{1}\usefont{T1}{cmr}{m}{n}\selectfont\color{color_29791}a}
\put(204.2408,-327.096){\fontsize{11.991}{1}\usefont{T1}{cmr}{m}{n}\selectfont\color{color_29791}sh}
\put(216.9153,-327.096){\fontsize{11.991}{1}\usefont{T1}{cmr}{m}{n}\selectfont\color{color_29791}i}
\put(219.5173,-327.096){\fontsize{11.991}{1}\usefont{T1}{cmr}{m}{n}\selectfont\color{color_29791}n}
\put(226.3043,-327.096){\fontsize{11.991}{1}\usefont{T1}{cmr}{m}{n}\selectfont\color{color_29791}g}
\put(232.9832,-327.096){\fontsize{11.991}{1}\usefont{T1}{cmr}{m}{n}\selectfont\color{color_29791} }
\put(236.2808,-327.096){\fontsize{11.991}{1}\usefont{T1}{cmr}{m}{n}\selectfont\color{color_29791}mo}
\put(252.9483,-327.096){\fontsize{11.991}{1}\usefont{T1}{cmr}{m}{n}\selectfont\color{color_29791}d}
\put(259.6273,-327.096){\fontsize{11.991}{1}\usefont{T1}{cmr}{m}{n}\selectfont\color{color_29791}u}
\put(266.4142,-327.096){\fontsize{11.991}{1}\usefont{T1}{cmr}{m}{n}\selectfont\color{color_29791}l}
\put(269.0162,-327.096){\fontsize{11.991}{1}\usefont{T1}{cmr}{m}{n}\selectfont\color{color_29791}e}
\put(275.6952,-327.096){\fontsize{11.991}{1}\usefont{T1}{cmr}{m}{n}\selectfont\color{color_29791}s}
\put(84.213,-355.187){\fontsize{11.991}{1}\usefont{T1}{cmr}{m}{n}\selectfont\color{color_29791}•}
\put(106.691,-355.811){\fontsize{11.991}{1}\usefont{T1}{cmr}{m}{n}\selectfont\color{color_29791}M}
\put(116.5836,-355.811){\fontsize{11.991}{1}\usefont{T1}{cmr}{m}{n}\selectfont\color{color_29791}o}
\put(123.3705,-355.811){\fontsize{11.991}{1}\usefont{T1}{cmr}{m}{n}\selectfont\color{color_29791}r}
\put(127.2676,-355.811){\fontsize{11.991}{1}\usefont{T1}{cmr}{m}{n}\selectfont\color{color_29791}e}
\put(133.9465,-355.811){\fontsize{11.991}{1}\usefont{T1}{cmr}{m}{n}\selectfont\color{color_29791} }
\put(137.328,-355.811){\fontsize{11.991}{1}\usefont{T1}{cmr}{m}{n}\selectfont\color{color_29791}i}
\put(140.014,-355.811){\fontsize{11.991}{1}\usefont{T1}{cmr}{m}{n}\selectfont\color{color_29791}n}
\put(146.693,-355.811){\fontsize{11.991}{1}\usefont{T1}{cmr}{m}{n}\selectfont\color{color_29791}f}
\put(149.9905,-355.811){\fontsize{11.991}{1}\usefont{T1}{cmr}{m}{n}\selectfont\color{color_29791}o}
\put(156.7774,-355.811){\fontsize{11.991}{1}\usefont{T1}{cmr}{m}{n}\selectfont\color{color_29791}r}
\put(160.6745,-355.811){\fontsize{11.991}{1}\usefont{T1}{cmr}{m}{n}\selectfont\color{color_29791}ma}
\put(177.342,-355.811){\fontsize{11.991}{1}\usefont{T1}{cmr}{m}{n}\selectfont\color{color_29791}t}
\put(180.6395,-355.811){\fontsize{11.991}{1}\usefont{T1}{cmr}{m}{n}\selectfont\color{color_29791}i}
\put(183.3255,-355.811){\fontsize{11.991}{1}\usefont{T1}{cmr}{m}{n}\selectfont\color{color_29791}o}
\put(190.0045,-355.811){\fontsize{11.991}{1}\usefont{T1}{cmr}{m}{n}\selectfont\color{color_29791}n}
\put(196.6835,-355.811){\fontsize{11.991}{1}\usefont{T1}{cmr}{m}{n}\selectfont\color{color_29791}:}
\put(200.0649,-355.811){\fontsize{11.991}{1}\usefont{T1}{cmr}{m}{n}\selectfont\color{color_29791} }
\put(203.693,-355.811){\fontsize{11.991}{1}\usefont{T1}{cmr}{m}{n}\selectfont\color{color_167934}h}
\put(209.8923,-355.811){\fontsize{11.991}{1}\usefont{T1}{cmr}{m}{n}\selectfont\color{color_167934}tt}
\put(217.4826,-355.811){\fontsize{11.991}{1}\usefont{T1}{cmr}{m}{n}\selectfont\color{color_167934}p:}
\put(227.0634,-355.811){\fontsize{11.991}{1}\usefont{T1}{cmr}{m}{n}\selectfont\color{color_167934}/}
\put(231.7519,-355.811){\fontsize{11.991}{1}\usefont{T1}{cmr}{m}{n}\selectfont\color{color_167934}/}
\put(236.3445,-355.811){\fontsize{11.991}{1}\usefont{T1}{cmr}{m}{n}\selectfont\color{color_167934}i}
\put(239.1264,-355.811){\fontsize{11.991}{1}\usefont{T1}{cmr}{m}{n}\selectfont\color{color_167934}c}
\put(244.1266,-355.811){\fontsize{11.991}{1}\usefont{T1}{cmr}{m}{n}\selectfont\color{color_167934}t}
\put(248.0237,-355.811){\fontsize{11.991}{1}\usefont{T1}{cmr}{m}{n}\selectfont\color{color_167934}do}
\put(260.6982,-355.811){\fontsize{11.991}{1}\usefont{T1}{cmr}{m}{n}\selectfont\color{color_167934}c}
\put(265.6985,-355.811){\fontsize{11.991}{1}\usefont{T1}{cmr}{m}{n}\selectfont\color{color_167934}s/}
\put(275.0754,-355.811){\fontsize{11.991}{1}\usefont{T1}{cmr}{m}{n}\selectfont\color{color_167934}w}
\put(283.565,-355.811){\fontsize{11.991}{1}\usefont{T1}{cmr}{m}{n}\selectfont\color{color_167934}i}
\put(286.347,-355.811){\fontsize{11.991}{1}\usefont{T1}{cmr}{m}{n}\selectfont\color{color_167934}k}
\put(291.7429,-355.811){\fontsize{11.991}{1}\usefont{T1}{cmr}{m}{n}\selectfont\color{color_167934}i}
\put(294.4409,-355.811){\fontsize{11.991}{1}\usefont{T1}{cmr}{m}{n}\selectfont\color{color_167934}/}
\put(299.1293,-355.811){\fontsize{11.991}{1}\usefont{T1}{cmr}{m}{n}\selectfont\color{color_167934}i}
\put(301.9113,-355.811){\fontsize{11.991}{1}\usefont{T1}{cmr}{m}{n}\selectfont\color{color_167934}nd}
\put(314.5858,-355.811){\fontsize{11.991}{1}\usefont{T1}{cmr}{m}{n}\selectfont\color{color_167934}e}
\put(320.2815,-355.811){\fontsize{11.991}{1}\usefont{T1}{cmr}{m}{n}\selectfont\color{color_167934}x.}
\put(328.4713,-355.811){\fontsize{11.991}{1}\usefont{T1}{cmr}{m}{n}\selectfont\color{color_167934}p}
\put(334.8506,-355.811){\fontsize{11.991}{1}\usefont{T1}{cmr}{m}{n}\selectfont\color{color_167934}hp}
\put(347.5251,-355.811){\fontsize{11.991}{1}\usefont{T1}{cmr}{m}{n}\selectfont\color{color_167934}?}
\put(353.0169,-355.811){\fontsize{11.991}{1}\usefont{T1}{cmr}{m}{n}\selectfont\color{color_167934}tit}
\put(363.6889,-355.811){\fontsize{11.991}{1}\usefont{T1}{cmr}{m}{n}\selectfont\color{color_167934}l}
\put(366.3869,-355.811){\fontsize{11.991}{1}\usefont{T1}{cmr}{m}{n}\selectfont\color{color_167934}e}
\put(372.4663,-355.811){\fontsize{11.991}{1}\usefont{T1}{cmr}{m}{n}\selectfont\color{color_167934}=}
\put(378.3659,-355.811){\fontsize{11.991}{1}\usefont{T1}{cmr}{m}{n}\selectfont\color{color_167934}H}
\put(385.8483,-355.811){\fontsize{11.991}{1}\usefont{T1}{cmr}{m}{n}\selectfont\color{color_167934}PC:Use}
\put(419.9867,-355.811){\fontsize{11.991}{1}\usefont{T1}{cmr}{m}{n}\selectfont\color{color_167934}r}
\put(424.1715,-355.811){\fontsize{11.991}{1}\usefont{T1}{cmr}{m}{n}\selectfont\color{color_167934}\_}
\put(430.155,-355.811){\fontsize{11.991}{1}\usefont{T1}{cmr}{m}{n}\selectfont\color{color_167934}G}
\put(437.6494,-355.811){\fontsize{11.991}{1}\usefont{T1}{cmr}{m}{n}\selectfont\color{color_167934}u}
\put(444.0286,-355.811){\fontsize{11.991}{1}\usefont{T1}{cmr}{m}{n}\selectfont\color{color_167934}i}
\put(446.8105,-355.811){\fontsize{11.991}{1}\usefont{T1}{cmr}{m}{n}\selectfont\color{color_167934}de}
\end{picture}
\begin{tikzpicture}[overlay]
\path(0pt,0pt);
\draw[color_167934,line width=0.794pt]
(203.693pt, -357.002pt) -- (459.605pt, -357.002pt)
;
\end{tikzpicture}
\begin{picture}(-5,0)(2.5,0)
\put(84.213,-383.902){\fontsize{11.991}{1}\usefont{T1}{cmr}{m}{n}\selectfont\color{color_29791}•}
\put(106.691,-384.611){\fontsize{11.991}{1}\usefont{T1}{cmr}{m}{n}\selectfont\color{color_29791}D}
\put(115.2885,-384.611){\fontsize{11.991}{1}\usefont{T1}{cmr}{m}{n}\selectfont\color{color_29791}e}
\put(122.0755,-384.611){\fontsize{11.991}{1}\usefont{T1}{cmr}{m}{n}\selectfont\color{color_29791}m}
\put(131.968,-384.611){\fontsize{11.991}{1}\usefont{T1}{cmr}{m}{n}\selectfont\color{color_29791}o}
\put(138.647,-384.611){\fontsize{11.991}{1}\usefont{T1}{cmr}{m}{n}\selectfont\color{color_29791}n}
\put(145.326,-384.611){\fontsize{11.991}{1}\usefont{T1}{cmr}{m}{n}\selectfont\color{color_29791}s}
\put(151.4054,-384.611){\fontsize{11.991}{1}\usefont{T1}{cmr}{m}{n}\selectfont\color{color_29791}t}
\put(154.703,-384.611){\fontsize{11.991}{1}\usefont{T1}{cmr}{m}{n}\selectfont\color{color_29791}ra}
\put(165.375,-384.611){\fontsize{11.991}{1}\usefont{T1}{cmr}{m}{n}\selectfont\color{color_29791}t}
\put(168.7564,-384.611){\fontsize{11.991}{1}\usefont{T1}{cmr}{m}{n}\selectfont\color{color_29791}i}
\put(171.3585,-384.611){\fontsize{11.991}{1}\usefont{T1}{cmr}{m}{n}\selectfont\color{color_29791}o}
\put(178.0375,-384.611){\fontsize{11.991}{1}\usefont{T1}{cmr}{m}{n}\selectfont\color{color_29791}n}
\end{picture}
\newpage
\begin{tikzpicture}[overlay]
\path(0pt,0pt);
\filldraw[color_283006][even odd rule]
(-15pt, 10pt) -- (704.972pt, 10pt)
 -- (704.972pt, 10pt)
 -- (704.972pt, -529.972pt)
 -- (704.972pt, -529.972pt)
 -- (-15pt, -529.972pt)
 -- (-15pt, -529.972pt)
 -- (-15pt, 10pt) -- cycle
;
\end{tikzpicture}
\begin{picture}(-5,0)(2.5,0)
\put(-15,-529.972){\includegraphics[width=720pt,height=540pt]{latexImage_d5964f0fe7de530ea43707adcb11ad51.png}}
\put(48.411,-60.69601){\fontsize{21.997}{1}\usefont{T1}{cmr}{b}{n}\selectfont\color{color_29791}J}
\put(60.68533,-60.69601){\fontsize{21.997}{1}\usefont{T1}{cmr}{b}{n}\selectfont\color{color_29791}u}
\put(74.0815,-60.69601){\fontsize{21.997}{1}\usefont{T1}{cmr}{b}{n}\selectfont\color{color_29791}p}
\put(87.56566,-60.69601){\fontsize{21.997}{1}\usefont{T1}{cmr}{b}{n}\selectfont\color{color_29791}y}
\put(99.77399,-60.69601){\fontsize{21.997}{1}\usefont{T1}{cmr}{b}{n}\selectfont\color{color_29791}t}
\put(107.077,-60.69601){\fontsize{21.997}{1}\usefont{T1}{cmr}{b}{n}\selectfont\color{color_29791}e}
\put(119.3513,-60.69601){\fontsize{21.997}{1}\usefont{T1}{cmr}{b}{n}\selectfont\color{color_29791}r}
\put(127.8422,-60.69601){\fontsize{21.997}{1}\usefont{T1}{cmr}{b}{n}\selectfont\color{color_29791} N}
\put(149.7292,-60.69601){\fontsize{21.997}{1}\usefont{T1}{cmr}{b}{n}\selectfont\color{color_29791}o}
\put(163.2133,-60.69601){\fontsize{21.997}{1}\usefont{T1}{cmr}{b}{n}\selectfont\color{color_29791}t}
\put(170.5163,-60.69601){\fontsize{21.997}{1}\usefont{T1}{cmr}{b}{n}\selectfont\color{color_29791}e}
\put(182.7247,-60.69601){\fontsize{21.997}{1}\usefont{T1}{cmr}{b}{n}\selectfont\color{color_29791}b}
\put(196.2088,-60.69601){\fontsize{21.997}{1}\usefont{T1}{cmr}{b}{n}\selectfont\color{color_29791}o}
\put(209.605,-60.69601){\fontsize{21.997}{1}\usefont{T1}{cmr}{b}{n}\selectfont\color{color_29791}o}
\put(223.0892,-60.69601){\fontsize{21.997}{1}\usefont{T1}{cmr}{b}{n}\selectfont\color{color_29791}k}
\put(235.2975,-60.69601){\fontsize{21.997}{1}\usefont{T1}{cmr}{b}{n}\selectfont\color{color_29791}s}
\put(48.213,-115.49){\fontsize{15.987}{1}\usefont{T1}{cmr}{m}{n}\selectfont\color{color_29791}•}
\put(62.187,-116.51){\fontsize{15.987}{1}\usefont{T1}{cmr}{m}{n}\selectfont\color{color_29791}J}
\put(70.26043,-116.51){\fontsize{15.987}{1}\usefont{T1}{cmr}{m}{n}\selectfont\color{color_29791}u}
\put(79.05328,-116.51){\fontsize{15.987}{1}\usefont{T1}{cmr}{m}{n}\selectfont\color{color_29791}pt}
\put(92.41841,-116.51){\fontsize{15.987}{1}\usefont{T1}{cmr}{m}{n}\selectfont\color{color_29791}e}
\put(101.2113,-116.51){\fontsize{15.987}{1}\usefont{T1}{cmr}{m}{n}\selectfont\color{color_29791}r}
\put(106.503,-116.51){\fontsize{15.987}{1}\usefont{T1}{cmr}{m}{n}\selectfont\color{color_29791} }
\put(110.9793,-116.51){\fontsize{15.987}{1}\usefont{T1}{cmr}{m}{n}\selectfont\color{color_29791}n}
\put(119.7722,-116.51){\fontsize{15.987}{1}\usefont{T1}{cmr}{m}{n}\selectfont\color{color_29791}ot}
\put(133.1373,-116.51){\fontsize{15.987}{1}\usefont{T1}{cmr}{m}{n}\selectfont\color{color_29791}e}
\put(141.9301,-116.51){\fontsize{15.987}{1}\usefont{T1}{cmr}{m}{n}\selectfont\color{color_29791}books }
\put(189.0598,-116.51){\fontsize{15.987}{1}\usefont{T1}{cmr}{m}{n}\selectfont\color{color_29791}r}
\put(194.3515,-116.51){\fontsize{15.987}{1}\usefont{T1}{cmr}{m}{n}\selectfont\color{color_29791}un}
\put(212.0331,-116.51){\fontsize{15.987}{1}\usefont{T1}{cmr}{m}{n}\selectfont\color{color_29791}ni}
\put(224.503,-116.51){\fontsize{15.987}{1}\usefont{T1}{cmr}{m}{n}\selectfont\color{color_29791}ng }
\put(246.677,-116.51){\fontsize{15.987}{1}\usefont{T1}{cmr}{m}{n}\selectfont\color{color_29791}on }
\put(268.851,-116.51){\fontsize{15.987}{1}\usefont{T1}{cmr}{m}{n}\selectfont\color{color_29791}w}
\put(280.3456,-116.51){\fontsize{15.987}{1}\usefont{T1}{cmr}{m}{n}\selectfont\color{color_29791}or}
\put(294.5261,-116.51){\fontsize{15.987}{1}\usefont{T1}{cmr}{m}{n}\selectfont\color{color_29791}ks}
\put(310.593,-116.51){\fontsize{15.987}{1}\usefont{T1}{cmr}{m}{n}\selectfont\color{color_29791}t}
\put(314.9894,-116.51){\fontsize{15.987}{1}\usefont{T1}{cmr}{m}{n}\selectfont\color{color_29791}at}
\put(328.2746,-116.51){\fontsize{15.987}{1}\usefont{T1}{cmr}{m}{n}\selectfont\color{color_29791}i}
\put(331.8557,-116.51){\fontsize{15.987}{1}\usefont{T1}{cmr}{m}{n}\selectfont\color{color_29791}ons:}
\put(362.603,-116.51){\fontsize{15.987}{1}\usefont{T1}{cmr}{m}{n}\selectfont\color{color_167818} }
\put(366.203,-116.51){\fontsize{15.987}{1}\usefont{T1}{cmr}{m}{n}\selectfont\color{color_167934}h}
\put(374.5002,-116.51){\fontsize{15.987}{1}\usefont{T1}{cmr}{m}{n}\selectfont\color{color_167934}tt}
\put(384.4921,-116.51){\fontsize{15.987}{1}\usefont{T1}{cmr}{m}{n}\selectfont\color{color_167934}p:}
\put(397.1698,-116.51){\fontsize{15.987}{1}\usefont{T1}{cmr}{m}{n}\selectfont\color{color_167934}/}
\put(403.3568,-116.51){\fontsize{15.987}{1}\usefont{T1}{cmr}{m}{n}\selectfont\color{color_167934}/}
\put(409.4478,-116.51){\fontsize{15.987}{1}\usefont{T1}{cmr}{m}{n}\selectfont\color{color_167934}j}
\put(413.2368,-116.51){\fontsize{15.987}{1}\usefont{T1}{cmr}{m}{n}\selectfont\color{color_167934}up}
\put(429.9272,-116.51){\fontsize{15.987}{1}\usefont{T1}{cmr}{m}{n}\selectfont\color{color_167934}y}
\put(437.2013,-116.51){\fontsize{15.987}{1}\usefont{T1}{cmr}{m}{n}\selectfont\color{color_167934}t}
\put(442.397,-116.51){\fontsize{15.987}{1}\usefont{T1}{cmr}{m}{n}\selectfont\color{color_167934}e}
\put(450.2946,-116.51){\fontsize{15.987}{1}\usefont{T1}{cmr}{m}{n}\selectfont\color{color_167934}r}
\put(455.7942,-116.51){\fontsize{15.987}{1}\usefont{T1}{cmr}{m}{n}\selectfont\color{color_167934}hub.}
\put(485.0504,-116.51){\fontsize{15.987}{1}\usefont{T1}{cmr}{m}{n}\selectfont\color{color_167934}ne}
\put(501.3411,-116.51){\fontsize{15.987}{1}\usefont{T1}{cmr}{m}{n}\selectfont\color{color_167934}r}
\put(506.6328,-116.51){\fontsize{15.987}{1}\usefont{T1}{cmr}{m}{n}\selectfont\color{color_167934}c}
\put(513.3314,-116.51){\fontsize{15.987}{1}\usefont{T1}{cmr}{m}{n}\selectfont\color{color_167934}-}
\put(518.3193,-116.51){\fontsize{15.987}{1}\usefont{T1}{cmr}{m}{n}\selectfont\color{color_167934}b}
\put(526.6165,-116.51){\fontsize{15.987}{1}\usefont{T1}{cmr}{m}{n}\selectfont\color{color_167934}a}
\put(534.2903,-116.51){\fontsize{15.987}{1}\usefont{T1}{cmr}{m}{n}\selectfont\color{color_167934}s}
\put(540.5732,-116.51){\fontsize{15.987}{1}\usefont{T1}{cmr}{m}{n}\selectfont\color{color_167934}.}
\put(544.5699,-116.51){\fontsize{15.987}{1}\usefont{T1}{cmr}{m}{n}\selectfont\color{color_167934}a}
\put(552.2437,-116.51){\fontsize{15.987}{1}\usefont{T1}{cmr}{m}{n}\selectfont\color{color_167934}c}
\put(558.9423,-116.51){\fontsize{15.987}{1}\usefont{T1}{cmr}{m}{n}\selectfont\color{color_167934}.}
\put(563.0189,-116.51){\fontsize{15.987}{1}\usefont{T1}{cmr}{m}{n}\selectfont\color{color_167934}uk}
\end{picture}
\begin{tikzpicture}[overlay]
\path(0pt,0pt);
\draw[color_167934,line width=0.992pt]
(366.203pt, -118.097pt) -- (578.915pt, -118.097pt)
;
\end{tikzpicture}
\begin{picture}(-5,0)(2.5,0)
\put(48.213,-149.902){\fontsize{15.987}{1}\usefont{T1}{cmr}{m}{n}\selectfont\color{color_29791}•}
\put(62.187,-150.809){\fontsize{15.987}{1}\usefont{T1}{cmr}{m}{n}\selectfont\color{color_29791}M}
\put(75.47219,-150.809){\fontsize{15.987}{1}\usefont{T1}{cmr}{m}{n}\selectfont\color{color_29791}or}
\put(89.65266,-150.809){\fontsize{15.987}{1}\usefont{T1}{cmr}{m}{n}\selectfont\color{color_29791}e }
\put(102.9379,-150.809){\fontsize{15.987}{1}\usefont{T1}{cmr}{m}{n}\selectfont\color{color_29791}i}
\put(106.5189,-150.809){\fontsize{15.987}{1}\usefont{T1}{cmr}{m}{n}\selectfont\color{color_29791}nf}
\put(119.8041,-150.809){\fontsize{15.987}{1}\usefont{T1}{cmr}{m}{n}\selectfont\color{color_29791}or}
\put(133.9846,-150.809){\fontsize{15.987}{1}\usefont{T1}{cmr}{m}{n}\selectfont\color{color_29791}m}
\put(147.2698,-150.809){\fontsize{15.987}{1}\usefont{T1}{cmr}{m}{n}\selectfont\color{color_29791}at}
\put(160.555,-150.809){\fontsize{15.987}{1}\usefont{T1}{cmr}{m}{n}\selectfont\color{color_29791}i}
\put(164.1361,-150.809){\fontsize{15.987}{1}\usefont{T1}{cmr}{m}{n}\selectfont\color{color_29791}on:}
\put(186.487,-150.809){\fontsize{15.987}{1}\usefont{T1}{cmr}{m}{n}\selectfont\color{color_165832} }
\put(190.087,-150.809){\fontsize{15.987}{1}\usefont{T1}{cmr}{m}{n}\selectfont\color{color_167934}h}
\put(198.3843,-150.809){\fontsize{15.987}{1}\usefont{T1}{cmr}{m}{n}\selectfont\color{color_167934}tt}
\put(208.4721,-150.809){\fontsize{15.987}{1}\usefont{T1}{cmr}{m}{n}\selectfont\color{color_167934}p:}
\put(221.1497,-150.809){\fontsize{15.987}{1}\usefont{T1}{cmr}{m}{n}\selectfont\color{color_167934}/}
\put(227.2408,-150.809){\fontsize{15.987}{1}\usefont{T1}{cmr}{m}{n}\selectfont\color{color_167934}/}
\put(233.4278,-150.809){\fontsize{15.987}{1}\usefont{T1}{cmr}{m}{n}\selectfont\color{color_167934}i}
\put(237.1048,-150.809){\fontsize{15.987}{1}\usefont{T1}{cmr}{m}{n}\selectfont\color{color_167934}c}
\put(243.8033,-150.809){\fontsize{15.987}{1}\usefont{T1}{cmr}{m}{n}\selectfont\color{color_167934}t}
\put(248.9991,-150.809){\fontsize{15.987}{1}\usefont{T1}{cmr}{m}{n}\selectfont\color{color_167934}do}
\put(265.7855,-150.809){\fontsize{15.987}{1}\usefont{T1}{cmr}{m}{n}\selectfont\color{color_167934}c}
\put(272.484,-150.809){\fontsize{15.987}{1}\usefont{T1}{cmr}{m}{n}\selectfont\color{color_167934}s}
\put(278.671,-150.809){\fontsize{15.987}{1}\usefont{T1}{cmr}{m}{n}\selectfont\color{color_167934}/}
\put(284.762,-150.809){\fontsize{15.987}{1}\usefont{T1}{cmr}{m}{n}\selectfont\color{color_167934}w}
\put(296.2407,-150.809){\fontsize{15.987}{1}\usefont{T1}{cmr}{m}{n}\selectfont\color{color_167934}i}
\put(299.9177,-150.809){\fontsize{15.987}{1}\usefont{T1}{cmr}{m}{n}\selectfont\color{color_167934}k}
\put(307.1118,-150.809){\fontsize{15.987}{1}\usefont{T1}{cmr}{m}{n}\selectfont\color{color_167934}i}
\put(310.7888,-150.809){\fontsize{15.987}{1}\usefont{T1}{cmr}{m}{n}\selectfont\color{color_167934}/}
\put(316.8799,-150.809){\fontsize{15.987}{1}\usefont{T1}{cmr}{m}{n}\selectfont\color{color_167934}i}
\put(320.5569,-150.809){\fontsize{15.987}{1}\usefont{T1}{cmr}{m}{n}\selectfont\color{color_167934}nde}
\put(345.033,-150.809){\fontsize{15.987}{1}\usefont{T1}{cmr}{m}{n}\selectfont\color{color_167934}x}
\put(352.0193,-150.809){\fontsize{15.987}{1}\usefont{T1}{cmr}{m}{n}\selectfont\color{color_167934}.}
\put(356.0161,-150.809){\fontsize{15.987}{1}\usefont{T1}{cmr}{m}{n}\selectfont\color{color_167934}p}
\put(364.4892,-150.809){\fontsize{15.987}{1}\usefont{T1}{cmr}{m}{n}\selectfont\color{color_167934}h}
\put(372.7864,-150.809){\fontsize{15.987}{1}\usefont{T1}{cmr}{m}{n}\selectfont\color{color_167934}p/}
\put(387.3666,-150.809){\fontsize{15.987}{1}\usefont{T1}{cmr}{m}{n}\selectfont\color{color_167934}H}
\put(397.2625,-150.809){\fontsize{15.987}{1}\usefont{T1}{cmr}{m}{n}\selectfont\color{color_167934}P}
\put(405.5438,-150.809){\fontsize{15.987}{1}\usefont{T1}{cmr}{m}{n}\selectfont\color{color_167934}C}
\put(414.0329,-150.809){\fontsize{15.987}{1}\usefont{T1}{cmr}{m}{n}\selectfont\color{color_167934}:}
\put(418.4134,-150.809){\fontsize{15.987}{1}\usefont{T1}{cmr}{m}{n}\selectfont\color{color_167934}J}
\put(423.4013,-150.809){\fontsize{15.987}{1}\usefont{T1}{cmr}{m}{n}\selectfont\color{color_167934}up}
\put(440.0917,-150.809){\fontsize{15.987}{1}\usefont{T1}{cmr}{m}{n}\selectfont\color{color_167934}y}
\put(447.3658,-150.809){\fontsize{15.987}{1}\usefont{T1}{cmr}{m}{n}\selectfont\color{color_167934}t}
\put(452.4657,-150.809){\fontsize{15.987}{1}\usefont{T1}{cmr}{m}{n}\selectfont\color{color_167934}e}
\put(460.4432,-150.809){\fontsize{15.987}{1}\usefont{T1}{cmr}{m}{n}\selectfont\color{color_167934}r}
\put(465.9427,-150.809){\fontsize{15.987}{1}\usefont{T1}{cmr}{m}{n}\selectfont\color{color_167934}H}
\put(475.9186,-150.809){\fontsize{15.987}{1}\usefont{T1}{cmr}{m}{n}\selectfont\color{color_167934}u}
\put(484.2159,-150.809){\fontsize{15.987}{1}\usefont{T1}{cmr}{m}{n}\selectfont\color{color_167934}b}
\end{picture}
\begin{tikzpicture}[overlay]
\path(0pt,0pt);
\end{tikzpicture}
\newpage
\begin{tikzpicture}[overlay]
\path(0pt,0pt);
\draw[color_167934,line width=0.992pt]
(190.087pt, -152.396pt) -- (493.082pt, -152.396pt)
;
\filldraw[color_283006][even odd rule]
(-15pt, 10pt) -- (704.972pt, 10pt)
 -- (704.972pt, 10pt)
 -- (704.972pt, -529.972pt)
 -- (704.972pt, -529.972pt)
 -- (-15pt, -529.972pt)
 -- (-15pt, -529.972pt)
 -- (-15pt, 10pt) -- cycle
;
\end{tikzpicture}
\begin{picture}(-5,0)(2.5,0)
\put(-15,-529.972){\includegraphics[width=720pt,height=540pt]{latexImage_3f3a920d1f055ee96cf3917d47957291.png}}
\put(48.411,-60.69601){\fontsize{21.997}{1}\usefont{T1}{cmr}{b}{n}\selectfont\color{color_29791}Co}
\put(77.68901,-60.69601){\fontsize{21.997}{1}\usefont{T1}{cmr}{b}{n}\selectfont\color{color_29791}n}
\put(91.17317,-60.69601){\fontsize{21.997}{1}\usefont{T1}{cmr}{b}{n}\selectfont\color{color_29791}t}
\put(98.47617,-60.69601){\fontsize{21.997}{1}\usefont{T1}{cmr}{b}{n}\selectfont\color{color_29791}a}
\put(110.6845,-60.69601){\fontsize{21.997}{1}\usefont{T1}{cmr}{b}{n}\selectfont\color{color_29791}in}
\put(130.2618,-60.69601){\fontsize{21.997}{1}\usefont{T1}{cmr}{b}{n}\selectfont\color{color_29791}e}
\put(142.4702,-60.69601){\fontsize{21.997}{1}\usefont{T1}{cmr}{b}{n}\selectfont\color{color_29791}r}
\put(151.049,-60.69601){\fontsize{21.997}{1}\usefont{T1}{cmr}{b}{n}\selectfont\color{color_29791}s}
\put(48.213,-115.49){\fontsize{15.987}{1}\usefont{T1}{cmr}{m}{n}\selectfont\color{color_29791}•}
\put(62.187,-116.51){\fontsize{15.987}{1}\usefont{T1}{cmr}{m}{n}\selectfont\color{color_29791}C}
\put(73.68165,-116.51){\fontsize{15.987}{1}\usefont{T1}{cmr}{m}{n}\selectfont\color{color_29791}ont}
\put(95.85562,-116.51){\fontsize{15.987}{1}\usefont{T1}{cmr}{m}{n}\selectfont\color{color_29791}ai}
\put(108.3255,-116.51){\fontsize{15.987}{1}\usefont{T1}{cmr}{m}{n}\selectfont\color{color_29791}ner}
\put(131.3947,-116.51){\fontsize{15.987}{1}\usefont{T1}{cmr}{m}{n}\selectfont\color{color_29791}s }
\put(143.7847,-116.51){\fontsize{15.987}{1}\usefont{T1}{cmr}{m}{n}\selectfont\color{color_29791}at}
\put(157.1498,-116.51){\fontsize{15.987}{1}\usefont{T1}{cmr}{m}{n}\selectfont\color{color_29791} }
\put(161.5462,-116.51){\fontsize{15.987}{1}\usefont{T1}{cmr}{m}{n}\selectfont\color{color_29791}B}
\put(172.2255,-116.51){\fontsize{15.987}{1}\usefont{T1}{cmr}{m}{n}\selectfont\color{color_29791}A}
\put(182.9048,-116.51){\fontsize{15.987}{1}\usefont{T1}{cmr}{m}{n}\selectfont\color{color_29791}S}
\put(193.4882,-116.51){\fontsize{15.987}{1}\usefont{T1}{cmr}{m}{n}\selectfont\color{color_29791} }
\put(197.9646,-116.51){\fontsize{15.987}{1}\usefont{T1}{cmr}{m}{n}\selectfont\color{color_29791}ar}
\put(212.0491,-116.51){\fontsize{15.987}{1}\usefont{T1}{cmr}{m}{n}\selectfont\color{color_29791}e }
\put(225.4143,-116.51){\fontsize{15.987}{1}\usefont{T1}{cmr}{m}{n}\selectfont\color{color_29791}st}
\put(237.8042,-116.51){\fontsize{15.987}{1}\usefont{T1}{cmr}{m}{n}\selectfont\color{color_29791}i}
\put(241.3853,-116.51){\fontsize{15.987}{1}\usefont{T1}{cmr}{m}{n}\selectfont\color{color_29791}l}
\put(244.9664,-116.51){\fontsize{15.987}{1}\usefont{T1}{cmr}{m}{n}\selectfont\color{color_29791}l}
\put(248.5475,-116.51){\fontsize{15.987}{1}\usefont{T1}{cmr}{m}{n}\selectfont\color{color_29791} }
\put(253.0238,-116.51){\fontsize{15.987}{1}\usefont{T1}{cmr}{m}{n}\selectfont\color{color_29791}a }
\put(266.309,-116.51){\fontsize{15.987}{1}\usefont{T1}{cmr}{m}{n}\selectfont\color{color_29791}w}
\put(277.8036,-116.51){\fontsize{15.987}{1}\usefont{T1}{cmr}{m}{n}\selectfont\color{color_29791}or}
\put(291.9841,-116.51){\fontsize{15.987}{1}\usefont{T1}{cmr}{m}{n}\selectfont\color{color_29791}k }
\put(304.454,-116.51){\fontsize{15.987}{1}\usefont{T1}{cmr}{m}{n}\selectfont\color{color_29791}i}
\put(308.0351,-116.51){\fontsize{15.987}{1}\usefont{T1}{cmr}{m}{n}\selectfont\color{color_29791}n}
\put(316.8279,-116.51){\fontsize{15.987}{1}\usefont{T1}{cmr}{m}{n}\selectfont\color{color_29791} }
\put(321.3043,-116.51){\fontsize{15.987}{1}\usefont{T1}{cmr}{m}{n}\selectfont\color{color_29791}pr}
\put(335.3888,-116.51){\fontsize{15.987}{1}\usefont{T1}{cmr}{m}{n}\selectfont\color{color_29791}ogr}
\put(358.458,-116.51){\fontsize{15.987}{1}\usefont{T1}{cmr}{m}{n}\selectfont\color{color_29791}e}
\put(367.2509,-116.51){\fontsize{15.987}{1}\usefont{T1}{cmr}{m}{n}\selectfont\color{color_29791}s}
\put(375.3243,-116.51){\fontsize{15.987}{1}\usefont{T1}{cmr}{m}{n}\selectfont\color{color_29791}s}
\put(48.213,-149.902){\fontsize{15.987}{1}\usefont{T1}{cmr}{m}{n}\selectfont\color{color_29791}•}
\put(62.187,-150.809){\fontsize{15.987}{1}\usefont{T1}{cmr}{m}{n}\selectfont\color{color_29791}What}
\put(99.43671,-150.809){\fontsize{15.987}{1}\usefont{T1}{cmr}{m}{n}\selectfont\color{color_29791} }
\put(103.9131,-150.809){\fontsize{15.987}{1}\usefont{T1}{cmr}{m}{n}\selectfont\color{color_29791}a}
\put(112.7059,-150.809){\fontsize{15.987}{1}\usefont{T1}{cmr}{m}{n}\selectfont\color{color_29791}r}
\put(117.9976,-150.809){\fontsize{15.987}{1}\usefont{T1}{cmr}{m}{n}\selectfont\color{color_29791}e }
\put(131.3627,-150.809){\fontsize{15.987}{1}\usefont{T1}{cmr}{m}{n}\selectfont\color{color_29791}cont}
\put(161.5302,-150.809){\fontsize{15.987}{1}\usefont{T1}{cmr}{m}{n}\selectfont\color{color_29791}ai}
\put(174.0001,-150.809){\fontsize{15.987}{1}\usefont{T1}{cmr}{m}{n}\selectfont\color{color_29791}ne}
\put(191.6817,-150.809){\fontsize{15.987}{1}\usefont{T1}{cmr}{m}{n}\selectfont\color{color_29791}r}
\put(196.9734,-150.809){\fontsize{15.987}{1}\usefont{T1}{cmr}{m}{n}\selectfont\color{color_29791}s}
\put(205.0468,-150.809){\fontsize{15.987}{1}\usefont{T1}{cmr}{m}{n}\selectfont\color{color_29791}?}
\put(48.213,-184.202){\fontsize{15.987}{1}\usefont{T1}{cmr}{m}{n}\selectfont\color{color_29791}•}
\put(62.187,-185.194){\fontsize{15.987}{1}\usefont{T1}{cmr}{m}{n}\selectfont\color{color_29791}P}
\put(72.86631,-185.194){\fontsize{15.987}{1}\usefont{T1}{cmr}{m}{n}\selectfont\color{color_29791}od}
\put(90.54793,-185.194){\fontsize{15.987}{1}\usefont{T1}{cmr}{m}{n}\selectfont\color{color_29791}m}
\put(103.929,-185.194){\fontsize{15.987}{1}\usefont{T1}{cmr}{m}{n}\selectfont\color{color_29791}a}
\put(112.7219,-185.194){\fontsize{15.987}{1}\usefont{T1}{cmr}{m}{n}\selectfont\color{color_29791}n}
\put(75.709,-214.107){\fontsize{11.991}{1}\usefont{T1}{cmr}{m}{n}\selectfont\color{color_29791}•}
\put(98.187,-214.702){\fontsize{11.991}{1}\usefont{T1}{cmr}{m}{n}\selectfont\color{color_29791}T}
\put(104.1825,-214.702){\fontsize{11.991}{1}\usefont{T1}{cmr}{m}{n}\selectfont\color{color_29791}o}
\put(110.8615,-214.702){\fontsize{11.991}{1}\usefont{T1}{cmr}{m}{n}\selectfont\color{color_29791} }
\put(114.243,-214.702){\fontsize{11.991}{1}\usefont{T1}{cmr}{m}{n}\selectfont\color{color_29791}b}
\put(120.9219,-214.702){\fontsize{11.991}{1}\usefont{T1}{cmr}{m}{n}\selectfont\color{color_29791}e}
\put(127.6009,-214.702){\fontsize{11.991}{1}\usefont{T1}{cmr}{m}{n}\selectfont\color{color_29791} }
\put(130.9824,-214.702){\fontsize{11.991}{1}\usefont{T1}{cmr}{m}{n}\selectfont\color{color_29791}a}
\put(137.6614,-214.702){\fontsize{11.991}{1}\usefont{T1}{cmr}{m}{n}\selectfont\color{color_29791}b}
\put(144.3404,-214.702){\fontsize{11.991}{1}\usefont{T1}{cmr}{m}{n}\selectfont\color{color_29791}l}
\put(147.0264,-214.702){\fontsize{11.991}{1}\usefont{T1}{cmr}{m}{n}\selectfont\color{color_29791}e}
\put(153.7053,-214.702){\fontsize{11.991}{1}\usefont{T1}{cmr}{m}{n}\selectfont\color{color_29791} }
\put(157.0868,-214.702){\fontsize{11.991}{1}\usefont{T1}{cmr}{m}{n}\selectfont\color{color_29791}t}
\put(160.3843,-214.702){\fontsize{11.991}{1}\usefont{T1}{cmr}{m}{n}\selectfont\color{color_29791}o}
\put(167.0633,-214.702){\fontsize{11.991}{1}\usefont{T1}{cmr}{m}{n}\selectfont\color{color_29791} }
\put(170.4448,-214.702){\fontsize{11.991}{1}\usefont{T1}{cmr}{m}{n}\selectfont\color{color_29791}u}
\put(177.1237,-214.702){\fontsize{11.991}{1}\usefont{T1}{cmr}{m}{n}\selectfont\color{color_29791}se}
\put(189.7982,-214.702){\fontsize{11.991}{1}\usefont{T1}{cmr}{m}{n}\selectfont\color{color_29791},}
\put(193.1797,-214.702){\fontsize{11.991}{1}\usefont{T1}{cmr}{m}{n}\selectfont\color{color_29791} }
\put(196.5612,-214.702){\fontsize{11.991}{1}\usefont{T1}{cmr}{m}{n}\selectfont\color{color_29791}y}
\put(202.4607,-214.702){\fontsize{11.991}{1}\usefont{T1}{cmr}{m}{n}\selectfont\color{color_29791}o}
\put(209.2476,-214.702){\fontsize{11.991}{1}\usefont{T1}{cmr}{m}{n}\selectfont\color{color_29791}u}
\put(215.9266,-214.702){\fontsize{11.991}{1}\usefont{T1}{cmr}{m}{n}\selectfont\color{color_29791} }
\put(219.3081,-214.702){\fontsize{11.991}{1}\usefont{T1}{cmr}{m}{n}\selectfont\color{color_29791}n}
\put(225.9871,-214.702){\fontsize{11.991}{1}\usefont{T1}{cmr}{m}{n}\selectfont\color{color_29791}e}
\put(232.666,-214.702){\fontsize{11.991}{1}\usefont{T1}{cmr}{m}{n}\selectfont\color{color_29791}e}
\put(239.345,-214.702){\fontsize{11.991}{1}\usefont{T1}{cmr}{m}{n}\selectfont\color{color_29791}d}
\put(246.024,-214.702){\fontsize{11.991}{1}\usefont{T1}{cmr}{m}{n}\selectfont\color{color_29791} }
\put(249.4055,-214.702){\fontsize{11.991}{1}\usefont{T1}{cmr}{m}{n}\selectfont\color{color_29791}t}
\put(252.703,-214.702){\fontsize{11.991}{1}\usefont{T1}{cmr}{m}{n}\selectfont\color{color_29791}o}
\put(259.4899,-214.702){\fontsize{11.991}{1}\usefont{T1}{cmr}{m}{n}\selectfont\color{color_29791} }
\put(262.7874,-214.702){\fontsize{11.991}{1}\usefont{T1}{cmr}{m}{n}\selectfont\color{color_29791}co}
\put(275.5698,-214.702){\fontsize{11.991}{1}\usefont{T1}{cmr}{m}{n}\selectfont\color{color_29791}n}
\put(282.2488,-214.702){\fontsize{11.991}{1}\usefont{T1}{cmr}{m}{n}\selectfont\color{color_29791}t}
\put(285.5463,-214.702){\fontsize{11.991}{1}\usefont{T1}{cmr}{m}{n}\selectfont\color{color_29791}a}
\put(292.2253,-214.702){\fontsize{11.991}{1}\usefont{T1}{cmr}{m}{n}\selectfont\color{color_29791}c}
\put(298.3047,-214.702){\fontsize{11.991}{1}\usefont{T1}{cmr}{m}{n}\selectfont\color{color_29791}t}
\put(301.6023,-214.702){\fontsize{11.991}{1}\usefont{T1}{cmr}{m}{n}\selectfont\color{color_29791} }
\put(304.9837,-214.702){\fontsize{11.991}{1}\usefont{T1}{cmr}{m}{n}\selectfont\color{color_29791}t}
\put(308.2813,-214.702){\fontsize{11.991}{1}\usefont{T1}{cmr}{m}{n}\selectfont\color{color_29791}h}
\put(315.0681,-214.702){\fontsize{11.991}{1}\usefont{T1}{cmr}{m}{n}\selectfont\color{color_29791}e}
\put(321.7471,-214.702){\fontsize{11.991}{1}\usefont{T1}{cmr}{m}{n}\selectfont\color{color_29791} }
\put(325.0446,-214.702){\fontsize{11.991}{1}\usefont{T1}{cmr}{m}{n}\selectfont\color{color_29791}se}
\put(337.8271,-214.702){\fontsize{11.991}{1}\usefont{T1}{cmr}{m}{n}\selectfont\color{color_29791}r}
\put(341.7241,-214.702){\fontsize{11.991}{1}\usefont{T1}{cmr}{m}{n}\selectfont\color{color_29791}vi}
\put(350.4056,-214.702){\fontsize{11.991}{1}\usefont{T1}{cmr}{m}{n}\selectfont\color{color_29791}c}
\put(356.3052,-214.702){\fontsize{11.991}{1}\usefont{T1}{cmr}{m}{n}\selectfont\color{color_29791}e}
\put(363.0921,-214.702){\fontsize{11.991}{1}\usefont{T1}{cmr}{m}{n}\selectfont\color{color_29791} }
\put(366.3896,-214.702){\fontsize{11.991}{1}\usefont{T1}{cmr}{m}{n}\selectfont\color{color_29791}d}
\put(373.1765,-214.702){\fontsize{11.991}{1}\usefont{T1}{cmr}{m}{n}\selectfont\color{color_29791}e}
\put(379.8555,-214.702){\fontsize{11.991}{1}\usefont{T1}{cmr}{m}{n}\selectfont\color{color_29791}sk}
\put(75.709,-242.794){\fontsize{11.991}{1}\usefont{T1}{cmr}{m}{n}\selectfont\color{color_29791}•}
\put(98.187,-243.502){\fontsize{11.991}{1}\usefont{T1}{cmr}{m}{n}\selectfont\color{color_29791}C}
\put(106.8685,-243.502){\fontsize{11.991}{1}\usefont{T1}{cmr}{m}{n}\selectfont\color{color_29791}o}
\put(113.5475,-243.502){\fontsize{11.991}{1}\usefont{T1}{cmr}{m}{n}\selectfont\color{color_29791}n}
\put(120.2265,-243.502){\fontsize{11.991}{1}\usefont{T1}{cmr}{m}{n}\selectfont\color{color_29791}t}
\put(123.524,-243.502){\fontsize{11.991}{1}\usefont{T1}{cmr}{m}{n}\selectfont\color{color_29791}a}
\put(130.3109,-243.502){\fontsize{11.991}{1}\usefont{T1}{cmr}{m}{n}\selectfont\color{color_29791}i}
\put(132.9129,-243.502){\fontsize{11.991}{1}\usefont{T1}{cmr}{m}{n}\selectfont\color{color_29791}n}
\put(139.5919,-243.502){\fontsize{11.991}{1}\usefont{T1}{cmr}{m}{n}\selectfont\color{color_29791}e}
\put(146.3788,-243.502){\fontsize{11.991}{1}\usefont{T1}{cmr}{m}{n}\selectfont\color{color_29791}r}
\put(150.2759,-243.502){\fontsize{11.991}{1}\usefont{T1}{cmr}{m}{n}\selectfont\color{color_29791} }
\put(153.6574,-243.502){\fontsize{11.991}{1}\usefont{T1}{cmr}{m}{n}\selectfont\color{color_29791}i}
\put(156.2594,-243.502){\fontsize{11.991}{1}\usefont{T1}{cmr}{m}{n}\selectfont\color{color_29791}ma}
\put(172.9269,-243.502){\fontsize{11.991}{1}\usefont{T1}{cmr}{m}{n}\selectfont\color{color_29791}g}
\put(179.6059,-243.502){\fontsize{11.991}{1}\usefont{T1}{cmr}{m}{n}\selectfont\color{color_29791}e}
\put(186.2849,-243.502){\fontsize{11.991}{1}\usefont{T1}{cmr}{m}{n}\selectfont\color{color_29791}s }
\put(195.6619,-243.502){\fontsize{11.991}{1}\usefont{T1}{cmr}{m}{n}\selectfont\color{color_29791}m}
\put(205.5544,-243.502){\fontsize{11.991}{1}\usefont{T1}{cmr}{m}{n}\selectfont\color{color_29791}u}
\put(212.3414,-243.502){\fontsize{11.991}{1}\usefont{T1}{cmr}{m}{n}\selectfont\color{color_29791}st}
\put(221.6344,-243.502){\fontsize{11.991}{1}\usefont{T1}{cmr}{m}{n}\selectfont\color{color_29791} }
\put(225.0158,-243.502){\fontsize{11.991}{1}\usefont{T1}{cmr}{m}{n}\selectfont\color{color_29791}b}
\put(231.6948,-243.502){\fontsize{11.991}{1}\usefont{T1}{cmr}{m}{n}\selectfont\color{color_29791}e}
\put(238.3738,-243.502){\fontsize{11.991}{1}\usefont{T1}{cmr}{m}{n}\selectfont\color{color_29791} }
\put(241.7553,-243.502){\fontsize{11.991}{1}\usefont{T1}{cmr}{m}{n}\selectfont\color{color_29791}d}
\put(248.4343,-243.502){\fontsize{11.991}{1}\usefont{T1}{cmr}{m}{n}\selectfont\color{color_29791}o}
\put(255.1133,-243.502){\fontsize{11.991}{1}\usefont{T1}{cmr}{m}{n}\selectfont\color{color_29791}w}
\put(263.7947,-243.502){\fontsize{11.991}{1}\usefont{T1}{cmr}{m}{n}\selectfont\color{color_29791}n}
\put(270.4737,-243.502){\fontsize{11.991}{1}\usefont{T1}{cmr}{m}{n}\selectfont\color{color_29791}l}
\put(273.0758,-243.502){\fontsize{11.991}{1}\usefont{T1}{cmr}{m}{n}\selectfont\color{color_29791}o}
\put(279.8627,-243.502){\fontsize{11.991}{1}\usefont{T1}{cmr}{m}{n}\selectfont\color{color_29791}a}
\put(286.5417,-243.502){\fontsize{11.991}{1}\usefont{T1}{cmr}{m}{n}\selectfont\color{color_29791}d}
\put(293.2206,-243.502){\fontsize{11.991}{1}\usefont{T1}{cmr}{m}{n}\selectfont\color{color_29791}e}
\put(299.8996,-243.502){\fontsize{11.991}{1}\usefont{T1}{cmr}{m}{n}\selectfont\color{color_29791}d}
\put(306.5786,-243.502){\fontsize{11.991}{1}\usefont{T1}{cmr}{m}{n}\selectfont\color{color_29791} }
\put(309.9601,-243.502){\fontsize{11.991}{1}\usefont{T1}{cmr}{m}{n}\selectfont\color{color_29791}t}
\put(313.2576,-243.502){\fontsize{11.991}{1}\usefont{T1}{cmr}{m}{n}\selectfont\color{color_29791}o}
\put(320.0445,-243.502){\fontsize{11.991}{1}\usefont{T1}{cmr}{m}{n}\selectfont\color{color_29791} }
\put(323.342,-243.502){\fontsize{11.991}{1}\usefont{T1}{cmr}{m}{n}\selectfont\color{color_29791}e}
\put(330.1289,-243.502){\fontsize{11.991}{1}\usefont{T1}{cmr}{m}{n}\selectfont\color{color_29791}a}
\put(336.8079,-243.502){\fontsize{11.991}{1}\usefont{T1}{cmr}{m}{n}\selectfont\color{color_29791}ch}
\put(349.4824,-243.502){\fontsize{11.991}{1}\usefont{T1}{cmr}{m}{n}\selectfont\color{color_29791} }
\put(352.7799,-243.502){\fontsize{11.991}{1}\usefont{T1}{cmr}{m}{n}\selectfont\color{color_29791}n}
\put(359.5668,-243.502){\fontsize{11.991}{1}\usefont{T1}{cmr}{m}{n}\selectfont\color{color_29791}o}
\put(366.2458,-243.502){\fontsize{11.991}{1}\usefont{T1}{cmr}{m}{n}\selectfont\color{color_29791}d}
\put(372.9248,-243.502){\fontsize{11.991}{1}\usefont{T1}{cmr}{m}{n}\selectfont\color{color_29791}e}
\put(379.6038,-243.502){\fontsize{11.991}{1}\usefont{T1}{cmr}{m}{n}\selectfont\color{color_29791} }
\put(382.9852,-243.502){\fontsize{11.991}{1}\usefont{T1}{cmr}{m}{n}\selectfont\color{color_29791}o}
\put(389.6642,-243.502){\fontsize{11.991}{1}\usefont{T1}{cmr}{m}{n}\selectfont\color{color_29791}r }
\put(396.9547,-243.502){\fontsize{11.991}{1}\usefont{T1}{cmr}{m}{n}\selectfont\color{color_29791}w}
\put(405.6362,-243.502){\fontsize{11.991}{1}\usefont{T1}{cmr}{m}{n}\selectfont\color{color_29791}o}
\put(412.3152,-243.502){\fontsize{11.991}{1}\usefont{T1}{cmr}{m}{n}\selectfont\color{color_29791}rks}
\put(428.2033,-243.502){\fontsize{11.991}{1}\usefont{T1}{cmr}{m}{n}\selectfont\color{color_29791}t}
\put(431.5847,-243.502){\fontsize{11.991}{1}\usefont{T1}{cmr}{m}{n}\selectfont\color{color_29791}a}
\put(438.2637,-243.502){\fontsize{11.991}{1}\usefont{T1}{cmr}{m}{n}\selectfont\color{color_29791}t}
\put(441.6452,-243.502){\fontsize{11.991}{1}\usefont{T1}{cmr}{m}{n}\selectfont\color{color_29791}i}
\put(444.2473,-243.502){\fontsize{11.991}{1}\usefont{T1}{cmr}{m}{n}\selectfont\color{color_29791}o}
\put(451.0341,-243.502){\fontsize{11.991}{1}\usefont{T1}{cmr}{m}{n}\selectfont\color{color_29791}n}
\put(48.213,-276.101){\fontsize{15.987}{1}\usefont{T1}{cmr}{m}{n}\selectfont\color{color_29791}•}
\put(62.187,-277.093){\fontsize{15.987}{1}\usefont{T1}{cmr}{m}{n}\selectfont\color{color_29791}S}
\put(72.86631,-277.093){\fontsize{15.987}{1}\usefont{T1}{cmr}{m}{n}\selectfont\color{color_29791}i}
\put(76.4474,-277.093){\fontsize{15.987}{1}\usefont{T1}{cmr}{m}{n}\selectfont\color{color_29791}ng}
\put(94.12901,-277.093){\fontsize{15.987}{1}\usefont{T1}{cmr}{m}{n}\selectfont\color{color_29791}ul}
\put(106.5989,-277.093){\fontsize{15.987}{1}\usefont{T1}{cmr}{m}{n}\selectfont\color{color_29791}ar}
\put(120.7793,-277.093){\fontsize{15.987}{1}\usefont{T1}{cmr}{m}{n}\selectfont\color{color_29791}i}
\put(124.3604,-277.093){\fontsize{15.987}{1}\usefont{T1}{cmr}{m}{n}\selectfont\color{color_29791}t}
\put(128.7569,-277.093){\fontsize{15.987}{1}\usefont{T1}{cmr}{m}{n}\selectfont\color{color_29791}y}
\put(75.709,-305.893){\fontsize{11.991}{1}\usefont{T1}{cmr}{m}{n}\selectfont\color{color_29791}•}
\put(98.187,-306.602){\fontsize{11.991}{1}\usefont{T1}{cmr}{m}{n}\selectfont\color{color_29791}D}
\put(106.8685,-306.602){\fontsize{11.991}{1}\usefont{T1}{cmr}{m}{n}\selectfont\color{color_29791}e}
\put(113.5475,-306.602){\fontsize{11.991}{1}\usefont{T1}{cmr}{m}{n}\selectfont\color{color_29791}si}
\put(122.145,-306.602){\fontsize{11.991}{1}\usefont{T1}{cmr}{m}{n}\selectfont\color{color_29791}g}
\put(128.824,-306.602){\fontsize{11.991}{1}\usefont{T1}{cmr}{m}{n}\selectfont\color{color_29791}n}
\put(135.6109,-306.602){\fontsize{11.991}{1}\usefont{T1}{cmr}{m}{n}\selectfont\color{color_29791}e}
\put(142.2899,-306.602){\fontsize{11.991}{1}\usefont{T1}{cmr}{m}{n}\selectfont\color{color_29791}d}
\put(148.9689,-306.602){\fontsize{11.991}{1}\usefont{T1}{cmr}{m}{n}\selectfont\color{color_29791} }
\put(152.2664,-306.602){\fontsize{11.991}{1}\usefont{T1}{cmr}{m}{n}\selectfont\color{color_29791}w}
\put(160.9479,-306.602){\fontsize{11.991}{1}\usefont{T1}{cmr}{m}{n}\selectfont\color{color_29791}i}
\put(163.55,-306.602){\fontsize{11.991}{1}\usefont{T1}{cmr}{m}{n}\selectfont\color{color_29791}t}
\put(166.9314,-306.602){\fontsize{11.991}{1}\usefont{T1}{cmr}{m}{n}\selectfont\color{color_29791}h}
\put(173.787,-306.602){\fontsize{11.991}{1}\usefont{T1}{cmr}{m}{n}\selectfont\color{color_29791} }
\put(177.189,-306.602){\fontsize{11.991}{1}\usefont{T1}{cmr}{m}{n}\selectfont\color{color_29791}H}
\put(185.7865,-306.602){\fontsize{11.991}{1}\usefont{T1}{cmr}{m}{n}\selectfont\color{color_29791}P}
\put(193.8685,-306.602){\fontsize{11.991}{1}\usefont{T1}{cmr}{m}{n}\selectfont\color{color_29791}C}
\put(202.466,-306.602){\fontsize{11.991}{1}\usefont{T1}{cmr}{m}{n}\selectfont\color{color_29791} }
\put(205.7636,-306.602){\fontsize{11.991}{1}\usefont{T1}{cmr}{m}{n}\selectfont\color{color_29791}u}
\put(212.5505,-306.602){\fontsize{11.991}{1}\usefont{T1}{cmr}{m}{n}\selectfont\color{color_29791}sa}
\put(225.225,-306.602){\fontsize{11.991}{1}\usefont{T1}{cmr}{m}{n}\selectfont\color{color_29791}g}
\put(231.9039,-306.602){\fontsize{11.991}{1}\usefont{T1}{cmr}{m}{n}\selectfont\color{color_29791}e}
\put(238.5829,-306.602){\fontsize{11.991}{1}\usefont{T1}{cmr}{m}{n}\selectfont\color{color_29791} }
\put(241.9644,-306.602){\fontsize{11.991}{1}\usefont{T1}{cmr}{m}{n}\selectfont\color{color_29791}i}
\put(244.5664,-306.602){\fontsize{11.991}{1}\usefont{T1}{cmr}{m}{n}\selectfont\color{color_29791}n}
\put(251.3533,-306.602){\fontsize{11.991}{1}\usefont{T1}{cmr}{m}{n}\selectfont\color{color_29791} }
\put(254.6508,-306.602){\fontsize{11.991}{1}\usefont{T1}{cmr}{m}{n}\selectfont\color{color_29791}mi}
\put(267.2414,-306.602){\fontsize{11.991}{1}\usefont{T1}{cmr}{m}{n}\selectfont\color{color_29791}n}
\put(273.9204,-306.602){\fontsize{11.991}{1}\usefont{T1}{cmr}{m}{n}\selectfont\color{color_29791}d}
\put(75.709,-334.693){\fontsize{11.991}{1}\usefont{T1}{cmr}{m}{n}\selectfont\color{color_29791}•}
\put(98.187,-335.402){\fontsize{11.991}{1}\usefont{T1}{cmr}{m}{n}\selectfont\color{color_29791}R}
\put(106.8685,-335.402){\fontsize{11.991}{1}\usefont{T1}{cmr}{m}{n}\selectfont\color{color_29791}e}
\put(113.5475,-335.402){\fontsize{11.991}{1}\usefont{T1}{cmr}{m}{n}\selectfont\color{color_29791}a}
\put(120.2265,-335.402){\fontsize{11.991}{1}\usefont{T1}{cmr}{m}{n}\selectfont\color{color_29791}d}
\put(126.9054,-335.402){\fontsize{11.991}{1}\usefont{T1}{cmr}{m}{n}\selectfont\color{color_29791}y }
\put(136.2824,-335.402){\fontsize{11.991}{1}\usefont{T1}{cmr}{m}{n}\selectfont\color{color_29791}t}
\put(139.5799,-335.402){\fontsize{11.991}{1}\usefont{T1}{cmr}{m}{n}\selectfont\color{color_29791}o}
\put(146.3669,-335.402){\fontsize{11.991}{1}\usefont{T1}{cmr}{m}{n}\selectfont\color{color_29791} }
\put(149.6644,-335.402){\fontsize{11.991}{1}\usefont{T1}{cmr}{m}{n}\selectfont\color{color_29791}u}
\put(156.3434,-335.402){\fontsize{11.991}{1}\usefont{T1}{cmr}{m}{n}\selectfont\color{color_29791}se}
\put(169.1258,-335.402){\fontsize{11.991}{1}\usefont{T1}{cmr}{m}{n}\selectfont\color{color_29791} }
\put(172.4233,-335.402){\fontsize{11.991}{1}\usefont{T1}{cmr}{m}{n}\selectfont\color{color_29791}o}
\put(179.1023,-335.402){\fontsize{11.991}{1}\usefont{T1}{cmr}{m}{n}\selectfont\color{color_29791}n}
\put(185.8892,-335.402){\fontsize{11.991}{1}\usefont{T1}{cmr}{m}{n}\selectfont\color{color_29791} }
\put(189.1867,-335.402){\fontsize{11.991}{1}\usefont{T1}{cmr}{m}{n}\selectfont\color{color_29791}w}
\put(197.8682,-335.402){\fontsize{11.991}{1}\usefont{T1}{cmr}{m}{n}\selectfont\color{color_29791}o}
\put(204.5472,-335.402){\fontsize{11.991}{1}\usefont{T1}{cmr}{m}{n}\selectfont\color{color_29791}r}
\put(208.4443,-335.402){\fontsize{11.991}{1}\usefont{T1}{cmr}{m}{n}\selectfont\color{color_29791}kst}
\put(223.8168,-335.402){\fontsize{11.991}{1}\usefont{T1}{cmr}{m}{n}\selectfont\color{color_29791}a}
\put(230.4957,-335.402){\fontsize{11.991}{1}\usefont{T1}{cmr}{m}{n}\selectfont\color{color_29791}t}
\put(233.8772,-335.402){\fontsize{11.991}{1}\usefont{T1}{cmr}{m}{n}\selectfont\color{color_29791}i}
\put(236.4792,-335.402){\fontsize{11.991}{1}\usefont{T1}{cmr}{m}{n}\selectfont\color{color_29791}o}
\put(243.1582,-335.402){\fontsize{11.991}{1}\usefont{T1}{cmr}{m}{n}\selectfont\color{color_29791}n}
\put(249.9451,-335.402){\fontsize{11.991}{1}\usefont{T1}{cmr}{m}{n}\selectfont\color{color_29791}s}
\put(255.8447,-335.402){\fontsize{11.991}{1}\usefont{T1}{cmr}{m}{n}\selectfont\color{color_29791} }
\put(259.2262,-335.402){\fontsize{11.991}{1}\usefont{T1}{cmr}{m}{n}\selectfont\color{color_29791}a}
\put(265.9052,-335.402){\fontsize{11.991}{1}\usefont{T1}{cmr}{m}{n}\selectfont\color{color_29791}n}
\put(272.692,-335.402){\fontsize{11.991}{1}\usefont{T1}{cmr}{m}{n}\selectfont\color{color_29791}d}
\put(279.371,-335.402){\fontsize{11.991}{1}\usefont{T1}{cmr}{m}{n}\selectfont\color{color_29791} }
\put(282.6685,-335.402){\fontsize{11.991}{1}\usefont{T1}{cmr}{m}{n}\selectfont\color{color_29791}n}
\put(289.4554,-335.402){\fontsize{11.991}{1}\usefont{T1}{cmr}{m}{n}\selectfont\color{color_29791}o}
\put(296.1344,-335.402){\fontsize{11.991}{1}\usefont{T1}{cmr}{m}{n}\selectfont\color{color_29791}d}
\put(302.8134,-335.402){\fontsize{11.991}{1}\usefont{T1}{cmr}{m}{n}\selectfont\color{color_29791}e}
\put(309.4924,-335.402){\fontsize{11.991}{1}\usefont{T1}{cmr}{m}{n}\selectfont\color{color_29791}s}
\put(48.213,-368){\fontsize{15.987}{1}\usefont{T1}{cmr}{m}{n}\selectfont\color{color_29791}•}
\put(62.187,-368.992){\fontsize{15.987}{1}\usefont{T1}{cmr}{m}{n}\selectfont\color{color_29791}F}
\put(71.87512,-368.992){\fontsize{15.987}{1}\usefont{T1}{cmr}{m}{n}\selectfont\color{color_29791}or}
\put(86.05559,-368.992){\fontsize{15.987}{1}\usefont{T1}{cmr}{m}{n}\selectfont\color{color_29791} }
\put(90.45201,-368.992){\fontsize{15.987}{1}\usefont{T1}{cmr}{m}{n}\selectfont\color{color_29791}m}
\put(103.8331,-368.992){\fontsize{15.987}{1}\usefont{T1}{cmr}{m}{n}\selectfont\color{color_29791}o}
\put(112.626,-368.992){\fontsize{15.987}{1}\usefont{T1}{cmr}{m}{n}\selectfont\color{color_29791}r}
\put(117.9177,-368.992){\fontsize{15.987}{1}\usefont{T1}{cmr}{m}{n}\selectfont\color{color_29791}e }
\put(131.2029,-368.992){\fontsize{15.987}{1}\usefont{T1}{cmr}{m}{n}\selectfont\color{color_29791}i}
\put(134.784,-368.992){\fontsize{15.987}{1}\usefont{T1}{cmr}{m}{n}\selectfont\color{color_29791}nf}
\put(148.0692,-368.992){\fontsize{15.987}{1}\usefont{T1}{cmr}{m}{n}\selectfont\color{color_29791}or}
\put(162.2496,-368.992){\fontsize{15.987}{1}\usefont{T1}{cmr}{m}{n}\selectfont\color{color_29791}m}
\put(175.5348,-368.992){\fontsize{15.987}{1}\usefont{T1}{cmr}{m}{n}\selectfont\color{color_29791}at}
\put(188.9,-368.992){\fontsize{15.987}{1}\usefont{T1}{cmr}{m}{n}\selectfont\color{color_29791}i}
\put(192.481,-368.992){\fontsize{15.987}{1}\usefont{T1}{cmr}{m}{n}\selectfont\color{color_29791}o}
\put(201.2739,-368.992){\fontsize{15.987}{1}\usefont{T1}{cmr}{m}{n}\selectfont\color{color_29791}n:}
\put(214.639,-368.992){\fontsize{15.987}{1}\usefont{T1}{cmr}{m}{n}\selectfont\color{color_29791} }
\put(219.312,-368.992){\fontsize{15.987}{1}\usefont{T1}{cmr}{m}{n}\selectfont\color{color_167934}h}
\put(227.6093,-368.992){\fontsize{15.987}{1}\usefont{T1}{cmr}{m}{n}\selectfont\color{color_167934}tt}
\put(237.6971,-368.992){\fontsize{15.987}{1}\usefont{T1}{cmr}{m}{n}\selectfont\color{color_167934}p}
\put(245.9943,-368.992){\fontsize{15.987}{1}\usefont{T1}{cmr}{m}{n}\selectfont\color{color_167934}:}
\put(250.3747,-368.992){\fontsize{15.987}{1}\usefont{T1}{cmr}{m}{n}\selectfont\color{color_167934}/}
\put(256.4658,-368.992){\fontsize{15.987}{1}\usefont{T1}{cmr}{m}{n}\selectfont\color{color_167934}/}
\put(262.5568,-368.992){\fontsize{15.987}{1}\usefont{T1}{cmr}{m}{n}\selectfont\color{color_167934}i}
\put(266.2338,-368.992){\fontsize{15.987}{1}\usefont{T1}{cmr}{m}{n}\selectfont\color{color_167934}c}
\put(273.0123,-368.992){\fontsize{15.987}{1}\usefont{T1}{cmr}{m}{n}\selectfont\color{color_167934}t}
\put(278.1121,-368.992){\fontsize{15.987}{1}\usefont{T1}{cmr}{m}{n}\selectfont\color{color_167934}do}
\put(294.8985,-368.992){\fontsize{15.987}{1}\usefont{T1}{cmr}{m}{n}\selectfont\color{color_167934}c}
\put(301.597,-368.992){\fontsize{15.987}{1}\usefont{T1}{cmr}{m}{n}\selectfont\color{color_167934}s}
\put(307.8799,-368.992){\fontsize{15.987}{1}\usefont{T1}{cmr}{m}{n}\selectfont\color{color_167934}/}
\put(313.9709,-368.992){\fontsize{15.987}{1}\usefont{T1}{cmr}{m}{n}\selectfont\color{color_167934}w}
\put(325.3537,-368.992){\fontsize{15.987}{1}\usefont{T1}{cmr}{m}{n}\selectfont\color{color_167934}i}
\put(329.0307,-368.992){\fontsize{15.987}{1}\usefont{T1}{cmr}{m}{n}\selectfont\color{color_167934}k}
\put(336.2249,-368.992){\fontsize{15.987}{1}\usefont{T1}{cmr}{m}{n}\selectfont\color{color_167934}i}
\put(339.9019,-368.992){\fontsize{15.987}{1}\usefont{T1}{cmr}{m}{n}\selectfont\color{color_167934}/}
\put(345.9929,-368.992){\fontsize{15.987}{1}\usefont{T1}{cmr}{m}{n}\selectfont\color{color_167934}i}
\put(349.7818,-368.992){\fontsize{15.987}{1}\usefont{T1}{cmr}{m}{n}\selectfont\color{color_167934}nde}
\put(374.2579,-368.992){\fontsize{15.987}{1}\usefont{T1}{cmr}{m}{n}\selectfont\color{color_167934}x}
\put(381.1483,-368.992){\fontsize{15.987}{1}\usefont{T1}{cmr}{m}{n}\selectfont\color{color_167934}.}
\put(385.225,-368.992){\fontsize{15.987}{1}\usefont{T1}{cmr}{m}{n}\selectfont\color{color_167934}php/}
\put(416.4956,-368.992){\fontsize{15.987}{1}\usefont{T1}{cmr}{m}{n}\selectfont\color{color_167934}H}
\put(426.4715,-368.992){\fontsize{15.987}{1}\usefont{T1}{cmr}{m}{n}\selectfont\color{color_167934}P}
\put(434.7527,-368.992){\fontsize{15.987}{1}\usefont{T1}{cmr}{m}{n}\selectfont\color{color_167934}C}
\put(443.2419,-368.992){\fontsize{15.987}{1}\usefont{T1}{cmr}{m}{n}\selectfont\color{color_167934}:}
\put(447.5264,-368.992){\fontsize{15.987}{1}\usefont{T1}{cmr}{m}{n}\selectfont\color{color_167934}C}
\put(456.0155,-368.992){\fontsize{15.987}{1}\usefont{T1}{cmr}{m}{n}\selectfont\color{color_167934}o}
\put(464.4087,-368.992){\fontsize{15.987}{1}\usefont{T1}{cmr}{m}{n}\selectfont\color{color_167934}n}
\put(472.7059,-368.992){\fontsize{15.987}{1}\usefont{T1}{cmr}{m}{n}\selectfont\color{color_167934}t}
\put(477.9017,-368.992){\fontsize{15.987}{1}\usefont{T1}{cmr}{m}{n}\selectfont\color{color_167934}a}
\put(485.4955,-368.992){\fontsize{15.987}{1}\usefont{T1}{cmr}{m}{n}\selectfont\color{color_167934}i}
\put(489.2844,-368.992){\fontsize{15.987}{1}\usefont{T1}{cmr}{m}{n}\selectfont\color{color_167934}n}
\put(497.5817,-368.992){\fontsize{15.987}{1}\usefont{T1}{cmr}{m}{n}\selectfont\color{color_167934}e}
\put(505.5592,-368.992){\fontsize{15.987}{1}\usefont{T1}{cmr}{m}{n}\selectfont\color{color_167934}r}
\put(510.8509,-368.992){\fontsize{15.987}{1}\usefont{T1}{cmr}{m}{n}\selectfont\color{color_167934}s}
\end{picture}
\begin{tikzpicture}[overlay]
\path(0pt,0pt);
\draw[color_167934,line width=0.992pt]
(219.312pt, -370.579pt) -- (517.517pt, -370.579pt)
;
\end{tikzpicture}
\begin{picture}(-5,0)(2.5,0)
\put(517.488,-368.992){\fontsize{15.987}{1}\usefont{T1}{cmr}{m}{n}\selectfont\color{color_167818} }
\put(48.213,-402.413){\fontsize{15.987}{1}\usefont{T1}{cmr}{m}{n}\selectfont\color{color_29791}•}
\put(62.187,-403.291){\fontsize{15.987}{1}\usefont{T1}{cmr}{m}{n}\selectfont\color{color_29791}D}
\put(73.68165,-403.291){\fontsize{15.987}{1}\usefont{T1}{cmr}{m}{n}\selectfont\color{color_29791}em}
\put(95.85562,-403.291){\fontsize{15.987}{1}\usefont{T1}{cmr}{m}{n}\selectfont\color{color_29791}onst}
\put(126.103,-403.291){\fontsize{15.987}{1}\usefont{T1}{cmr}{m}{n}\selectfont\color{color_29791}r}
\put(131.3947,-403.291){\fontsize{15.987}{1}\usefont{T1}{cmr}{m}{n}\selectfont\color{color_29791}a}
\put(140.1876,-403.291){\fontsize{15.987}{1}\usefont{T1}{cmr}{m}{n}\selectfont\color{color_29791}t}
\put(144.6639,-403.291){\fontsize{15.987}{1}\usefont{T1}{cmr}{m}{n}\selectfont\color{color_29791}i}
\put(148.245,-403.291){\fontsize{15.987}{1}\usefont{T1}{cmr}{m}{n}\selectfont\color{color_29791}on}
\end{picture}
\newpage
\begin{tikzpicture}[overlay]
\path(0pt,0pt);
\filldraw[color_283006][even odd rule]
(-15pt, 10pt) -- (704.972pt, 10pt)
 -- (704.972pt, 10pt)
 -- (704.972pt, -529.972pt)
 -- (704.972pt, -529.972pt)
 -- (-15pt, -529.972pt)
 -- (-15pt, -529.972pt)
 -- (-15pt, 10pt) -- cycle
;
\end{tikzpicture}
\begin{picture}(-5,0)(2.5,0)
\put(-15,-529.972){\includegraphics[width=720pt,height=540pt]{latexImage_bdabda61317014bb20dc1cc743a20acb.png}}
\put(48.411,-60.69601){\fontsize{21.997}{1}\usefont{T1}{cmr}{b}{n}\selectfont\color{color_29791}S}
\put(63.083,-60.69601){\fontsize{21.997}{1}\usefont{T1}{cmr}{b}{n}\selectfont\color{color_29791}L}
\put(76.47917,-60.69601){\fontsize{21.997}{1}\usefont{T1}{cmr}{b}{n}\selectfont\color{color_29791}U}
\put(92.27302,-60.69601){\fontsize{21.997}{1}\usefont{T1}{cmr}{b}{n}\selectfont\color{color_29791}RM}
\put(48.213,-115.49){\fontsize{15.987}{1}\usefont{T1}{cmr}{m}{n}\selectfont\color{color_29791}•}
\put(62.187,-116.51){\fontsize{15.987}{1}\usefont{T1}{cmr}{m}{n}\selectfont\color{color_29791}What}
\put(99.43671,-116.51){\fontsize{15.987}{1}\usefont{T1}{cmr}{m}{n}\selectfont\color{color_29791} }
\put(103.9131,-116.51){\fontsize{15.987}{1}\usefont{T1}{cmr}{m}{n}\selectfont\color{color_29791}i}
\put(107.4941,-116.51){\fontsize{15.987}{1}\usefont{T1}{cmr}{m}{n}\selectfont\color{color_29791}s }
\put(119.8841,-116.51){\fontsize{15.987}{1}\usefont{T1}{cmr}{m}{n}\selectfont\color{color_29791}i}
\put(123.4652,-116.51){\fontsize{15.987}{1}\usefont{T1}{cmr}{m}{n}\selectfont\color{color_29791}t}
\put(127.9415,-116.51){\fontsize{15.987}{1}\usefont{T1}{cmr}{m}{n}\selectfont\color{color_29791}?}
\put(84.213,-145.31){\fontsize{11.991}{1}\usefont{T1}{cmr}{m}{n}\selectfont\color{color_29791}•}
\put(98.187,-145.991){\fontsize{11.991}{1}\usefont{T1}{cmr}{m}{n}\selectfont\color{color_29791}S}
\put(106.2689,-145.991){\fontsize{11.991}{1}\usefont{T1}{cmr}{m}{n}\selectfont\color{color_29791}i}
\put(108.871,-145.991){\fontsize{11.991}{1}\usefont{T1}{cmr}{m}{n}\selectfont\color{color_29791}mp}
\put(125.5385,-145.991){\fontsize{11.991}{1}\usefont{T1}{cmr}{m}{n}\selectfont\color{color_29791}l}
\put(128.1405,-145.991){\fontsize{11.991}{1}\usefont{T1}{cmr}{m}{n}\selectfont\color{color_29791}e}
\put(134.8195,-145.991){\fontsize{11.991}{1}\usefont{T1}{cmr}{m}{n}\selectfont\color{color_29791} }
\put(138.201,-145.991){\fontsize{11.991}{1}\usefont{T1}{cmr}{m}{n}\selectfont\color{color_29791}L}
\put(144.8799,-145.991){\fontsize{11.991}{1}\usefont{T1}{cmr}{m}{n}\selectfont\color{color_29791}i}
\put(147.5659,-145.991){\fontsize{11.991}{1}\usefont{T1}{cmr}{m}{n}\selectfont\color{color_29791}n}
\put(154.2449,-145.991){\fontsize{11.991}{1}\usefont{T1}{cmr}{m}{n}\selectfont\color{color_29791}u}
\put(160.9239,-145.991){\fontsize{11.991}{1}\usefont{T1}{cmr}{m}{n}\selectfont\color{color_29791}x }
\put(170.2169,-145.991){\fontsize{11.991}{1}\usefont{T1}{cmr}{m}{n}\selectfont\color{color_29791}U}
\put(178.8984,-145.991){\fontsize{11.991}{1}\usefont{T1}{cmr}{m}{n}\selectfont\color{color_29791}t}
\put(182.2799,-145.991){\fontsize{11.991}{1}\usefont{T1}{cmr}{m}{n}\selectfont\color{color_29791}i}
\put(184.8819,-145.991){\fontsize{11.991}{1}\usefont{T1}{cmr}{m}{n}\selectfont\color{color_29791}l}
\put(187.484,-145.991){\fontsize{11.991}{1}\usefont{T1}{cmr}{m}{n}\selectfont\color{color_29791}i}
\put(190.17,-145.991){\fontsize{11.991}{1}\usefont{T1}{cmr}{m}{n}\selectfont\color{color_29791}t}
\put(193.4675,-145.991){\fontsize{11.991}{1}\usefont{T1}{cmr}{m}{n}\selectfont\color{color_29791}y }
\put(202.8445,-145.991){\fontsize{11.991}{1}\usefont{T1}{cmr}{m}{n}\selectfont\color{color_29791}f}
\put(206.142,-145.991){\fontsize{11.991}{1}\usefont{T1}{cmr}{m}{n}\selectfont\color{color_29791}o}
\put(212.9289,-145.991){\fontsize{11.991}{1}\usefont{T1}{cmr}{m}{n}\selectfont\color{color_29791}r}
\put(216.826,-145.991){\fontsize{11.991}{1}\usefont{T1}{cmr}{m}{n}\selectfont\color{color_29791} }
\put(220.2075,-145.991){\fontsize{11.991}{1}\usefont{T1}{cmr}{m}{n}\selectfont\color{color_29791}R}
\put(228.805,-145.991){\fontsize{11.991}{1}\usefont{T1}{cmr}{m}{n}\selectfont\color{color_29791}e}
\put(235.5919,-145.991){\fontsize{11.991}{1}\usefont{T1}{cmr}{m}{n}\selectfont\color{color_29791}s}
\put(241.4915,-145.991){\fontsize{11.991}{1}\usefont{T1}{cmr}{m}{n}\selectfont\color{color_29791}o}
\put(248.2784,-145.991){\fontsize{11.991}{1}\usefont{T1}{cmr}{m}{n}\selectfont\color{color_29791}u}
\put(254.9574,-145.991){\fontsize{11.991}{1}\usefont{T1}{cmr}{m}{n}\selectfont\color{color_29791}rce}
\put(271.6248,-145.991){\fontsize{11.991}{1}\usefont{T1}{cmr}{m}{n}\selectfont\color{color_29791} }
\put(274.9224,-145.991){\fontsize{11.991}{1}\usefont{T1}{cmr}{m}{n}\selectfont\color{color_29791}Ma}
\put(291.5898,-145.991){\fontsize{11.991}{1}\usefont{T1}{cmr}{m}{n}\selectfont\color{color_29791}n}
\put(298.2688,-145.991){\fontsize{11.991}{1}\usefont{T1}{cmr}{m}{n}\selectfont\color{color_29791}a}
\put(305.0557,-145.991){\fontsize{11.991}{1}\usefont{T1}{cmr}{m}{n}\selectfont\color{color_29791}g}
\put(311.7347,-145.991){\fontsize{11.991}{1}\usefont{T1}{cmr}{m}{n}\selectfont\color{color_29791}e}
\put(318.4137,-145.991){\fontsize{11.991}{1}\usefont{T1}{cmr}{m}{n}\selectfont\color{color_29791}m}
\put(328.3063,-145.991){\fontsize{11.991}{1}\usefont{T1}{cmr}{m}{n}\selectfont\color{color_29791}e}
\put(335.0932,-145.991){\fontsize{11.991}{1}\usefont{T1}{cmr}{m}{n}\selectfont\color{color_29791}n}
\put(341.7722,-145.991){\fontsize{11.991}{1}\usefont{T1}{cmr}{m}{n}\selectfont\color{color_29791}t}
\put(345.1536,-145.991){\fontsize{11.991}{1}\usefont{T1}{cmr}{m}{n}\selectfont\color{color_29791} }
\put(348.4511,-145.991){\fontsize{11.991}{1}\usefont{T1}{cmr}{m}{n}\selectfont\color{color_29791}–}
\put(355.1301,-145.991){\fontsize{11.991}{1}\usefont{T1}{cmr}{m}{n}\selectfont\color{color_29791} }
\put(358.5116,-145.991){\fontsize{11.991}{1}\usefont{T1}{cmr}{m}{n}\selectfont\color{color_29791}o}
\put(365.1906,-145.991){\fontsize{11.991}{1}\usefont{T1}{cmr}{m}{n}\selectfont\color{color_29791}u}
\put(371.8696,-145.991){\fontsize{11.991}{1}\usefont{T1}{cmr}{m}{n}\selectfont\color{color_29791}r }
\put(379.244,-145.991){\fontsize{11.991}{1}\usefont{T1}{cmr}{m}{n}\selectfont\color{color_29791}H}
\put(387.8416,-145.991){\fontsize{11.991}{1}\usefont{T1}{cmr}{m}{n}\selectfont\color{color_29791}PC}
\put(404.509,-145.991){\fontsize{11.991}{1}\usefont{T1}{cmr}{m}{n}\selectfont\color{color_29791} }
\put(407.8065,-145.991){\fontsize{11.991}{1}\usefont{T1}{cmr}{m}{n}\selectfont\color{color_29791}re}
\put(418.4785,-145.991){\fontsize{11.991}{1}\usefont{T1}{cmr}{m}{n}\selectfont\color{color_29791}so}
\put(431.153,-145.991){\fontsize{11.991}{1}\usefont{T1}{cmr}{m}{n}\selectfont\color{color_29791}u}
\put(437.9399,-145.991){\fontsize{11.991}{1}\usefont{T1}{cmr}{m}{n}\selectfont\color{color_29791}r}
\put(441.837,-145.991){\fontsize{11.991}{1}\usefont{T1}{cmr}{m}{n}\selectfont\color{color_29791}ce}
\put(454.5115,-145.991){\fontsize{11.991}{1}\usefont{T1}{cmr}{m}{n}\selectfont\color{color_29791} }
\put(457.893,-145.991){\fontsize{11.991}{1}\usefont{T1}{cmr}{m}{n}\selectfont\color{color_29791}m}
\put(467.7856,-145.991){\fontsize{11.991}{1}\usefont{T1}{cmr}{m}{n}\selectfont\color{color_29791}a}
\put(474.5724,-145.991){\fontsize{11.991}{1}\usefont{T1}{cmr}{m}{n}\selectfont\color{color_29791}n}
\put(481.2514,-145.991){\fontsize{11.991}{1}\usefont{T1}{cmr}{m}{n}\selectfont\color{color_29791}a}
\put(487.9304,-145.991){\fontsize{11.991}{1}\usefont{T1}{cmr}{m}{n}\selectfont\color{color_29791}g}
\put(494.6094,-145.991){\fontsize{11.991}{1}\usefont{T1}{cmr}{m}{n}\selectfont\color{color_29791}e}
\put(501.3963,-145.991){\fontsize{11.991}{1}\usefont{T1}{cmr}{m}{n}\selectfont\color{color_29791}r}
\put(84.213,-174.11){\fontsize{11.991}{1}\usefont{T1}{cmr}{m}{n}\selectfont\color{color_29791}•}
\put(98.187,-174.791){\fontsize{11.991}{1}\usefont{T1}{cmr}{m}{n}\selectfont\color{color_29791}S}
\put(106.2689,-174.791){\fontsize{11.991}{1}\usefont{T1}{cmr}{m}{n}\selectfont\color{color_29791}c}
\put(112.1685,-174.791){\fontsize{11.991}{1}\usefont{T1}{cmr}{m}{n}\selectfont\color{color_29791}h}
\put(118.9554,-174.791){\fontsize{11.991}{1}\usefont{T1}{cmr}{m}{n}\selectfont\color{color_29791}e}
\put(125.6344,-174.791){\fontsize{11.991}{1}\usefont{T1}{cmr}{m}{n}\selectfont\color{color_29791}d}
\put(132.3134,-174.791){\fontsize{11.991}{1}\usefont{T1}{cmr}{m}{n}\selectfont\color{color_29791}u}
\put(138.9924,-174.791){\fontsize{11.991}{1}\usefont{T1}{cmr}{m}{n}\selectfont\color{color_29791}l}
\put(141.6784,-174.791){\fontsize{11.991}{1}\usefont{T1}{cmr}{m}{n}\selectfont\color{color_29791}e}
\put(148.3573,-174.791){\fontsize{11.991}{1}\usefont{T1}{cmr}{m}{n}\selectfont\color{color_29791}s }
\put(157.7343,-174.791){\fontsize{11.991}{1}\usefont{T1}{cmr}{m}{n}\selectfont\color{color_29791}j}
\put(160.3364,-174.791){\fontsize{11.991}{1}\usefont{T1}{cmr}{m}{n}\selectfont\color{color_29791}o}
\put(167.0154,-174.791){\fontsize{11.991}{1}\usefont{T1}{cmr}{m}{n}\selectfont\color{color_29791}b}
\put(173.6943,-174.791){\fontsize{11.991}{1}\usefont{T1}{cmr}{m}{n}\selectfont\color{color_29791}s }
\put(183.0713,-174.791){\fontsize{11.991}{1}\usefont{T1}{cmr}{m}{n}\selectfont\color{color_29791}b}
\put(189.7503,-174.791){\fontsize{11.991}{1}\usefont{T1}{cmr}{m}{n}\selectfont\color{color_29791}a}
\put(196.4293,-174.791){\fontsize{11.991}{1}\usefont{T1}{cmr}{m}{n}\selectfont\color{color_29791}se}
\put(209.1038,-174.791){\fontsize{11.991}{1}\usefont{T1}{cmr}{m}{n}\selectfont\color{color_29791}d}
\put(215.8907,-174.791){\fontsize{11.991}{1}\usefont{T1}{cmr}{m}{n}\selectfont\color{color_29791} }
\put(219.1882,-174.791){\fontsize{11.991}{1}\usefont{T1}{cmr}{m}{n}\selectfont\color{color_29791}o}
\put(225.9751,-174.791){\fontsize{11.991}{1}\usefont{T1}{cmr}{m}{n}\selectfont\color{color_29791}n}
\put(232.6541,-174.791){\fontsize{11.991}{1}\usefont{T1}{cmr}{m}{n}\selectfont\color{color_29791} }
\put(235.9516,-174.791){\fontsize{11.991}{1}\usefont{T1}{cmr}{m}{n}\selectfont\color{color_29791}t}
\put(239.3331,-174.791){\fontsize{11.991}{1}\usefont{T1}{cmr}{m}{n}\selectfont\color{color_29791}h}
\put(246.0121,-174.791){\fontsize{11.991}{1}\usefont{T1}{cmr}{m}{n}\selectfont\color{color_29791}e}
\put(252.6911,-174.791){\fontsize{11.991}{1}\usefont{T1}{cmr}{m}{n}\selectfont\color{color_29791} }
\put(256.0725,-174.791){\fontsize{11.991}{1}\usefont{T1}{cmr}{m}{n}\selectfont\color{color_29791}r}
\put(259.9696,-174.791){\fontsize{11.991}{1}\usefont{T1}{cmr}{m}{n}\selectfont\color{color_29791}e}
\put(266.7565,-174.791){\fontsize{11.991}{1}\usefont{T1}{cmr}{m}{n}\selectfont\color{color_29791}so}
\put(279.431,-174.791){\fontsize{11.991}{1}\usefont{T1}{cmr}{m}{n}\selectfont\color{color_29791}u}
\put(286.11,-174.791){\fontsize{11.991}{1}\usefont{T1}{cmr}{m}{n}\selectfont\color{color_29791}rce}
\put(302.7775,-174.791){\fontsize{11.991}{1}\usefont{T1}{cmr}{m}{n}\selectfont\color{color_29791}s }
\put(312.1544,-174.791){\fontsize{11.991}{1}\usefont{T1}{cmr}{m}{n}\selectfont\color{color_29791}t}
\put(315.4519,-174.791){\fontsize{11.991}{1}\usefont{T1}{cmr}{m}{n}\selectfont\color{color_29791}h}
\put(322.1309,-174.791){\fontsize{11.991}{1}\usefont{T1}{cmr}{m}{n}\selectfont\color{color_29791}e}
\put(328.9178,-174.791){\fontsize{11.991}{1}\usefont{T1}{cmr}{m}{n}\selectfont\color{color_29791}y}
\put(334.8174,-174.791){\fontsize{11.991}{1}\usefont{T1}{cmr}{m}{n}\selectfont\color{color_29791} }
\put(338.1989,-174.791){\fontsize{11.991}{1}\usefont{T1}{cmr}{m}{n}\selectfont\color{color_29791}n}
\put(344.8778,-174.791){\fontsize{11.991}{1}\usefont{T1}{cmr}{m}{n}\selectfont\color{color_29791}e}
\put(351.6647,-174.791){\fontsize{11.991}{1}\usefont{T1}{cmr}{m}{n}\selectfont\color{color_29791}e}
\put(358.3437,-174.791){\fontsize{11.991}{1}\usefont{T1}{cmr}{m}{n}\selectfont\color{color_29791}d}
\put(48.213,-207.389){\fontsize{15.987}{1}\usefont{T1}{cmr}{m}{n}\selectfont\color{color_29791}•}
\put(62.187,-208.409){\fontsize{15.987}{1}\usefont{T1}{cmr}{m}{n}\selectfont\color{color_29791}D}
\put(73.68165,-208.409){\fontsize{15.987}{1}\usefont{T1}{cmr}{m}{n}\selectfont\color{color_29791}i}
\put(77.26273,-208.409){\fontsize{15.987}{1}\usefont{T1}{cmr}{m}{n}\selectfont\color{color_29791}f}
\put(81.56323,-208.409){\fontsize{15.987}{1}\usefont{T1}{cmr}{m}{n}\selectfont\color{color_29791}f}
\put(85.95966,-208.409){\fontsize{15.987}{1}\usefont{T1}{cmr}{m}{n}\selectfont\color{color_29791}e}
\put(94.7525,-208.409){\fontsize{15.987}{1}\usefont{T1}{cmr}{m}{n}\selectfont\color{color_29791}r}
\put(100.0442,-208.409){\fontsize{15.987}{1}\usefont{T1}{cmr}{m}{n}\selectfont\color{color_29791}ent}
\put(122.2182,-208.409){\fontsize{15.987}{1}\usefont{T1}{cmr}{m}{n}\selectfont\color{color_29791} }
\put(126.6945,-208.409){\fontsize{15.987}{1}\usefont{T1}{cmr}{m}{n}\selectfont\color{color_29791}qu}
\put(144.3762,-208.409){\fontsize{15.987}{1}\usefont{T1}{cmr}{m}{n}\selectfont\color{color_29791}eues}
\put(84.213,-237.209){\fontsize{11.991}{1}\usefont{T1}{cmr}{m}{n}\selectfont\color{color_29791}•}
\put(98.187,-237.89){\fontsize{11.991}{1}\usefont{T1}{cmr}{m}{n}\selectfont\color{color_29791}S}
\put(106.2689,-237.89){\fontsize{11.991}{1}\usefont{T1}{cmr}{m}{n}\selectfont\color{color_29791}h}
\put(112.9479,-237.89){\fontsize{11.991}{1}\usefont{T1}{cmr}{m}{n}\selectfont\color{color_29791}o}
\put(119.6269,-237.89){\fontsize{11.991}{1}\usefont{T1}{cmr}{m}{n}\selectfont\color{color_29791}r}
\put(123.524,-237.89){\fontsize{11.991}{1}\usefont{T1}{cmr}{m}{n}\selectfont\color{color_29791}t}
\put(126.9055,-237.89){\fontsize{11.991}{1}\usefont{T1}{cmr}{m}{n}\selectfont\color{color_29791} }
\put(130.2869,-237.89){\fontsize{11.991}{1}\usefont{T1}{cmr}{m}{n}\selectfont\color{color_29791}-}
\put(134.184,-237.89){\fontsize{11.991}{1}\usefont{T1}{cmr}{m}{n}\selectfont\color{color_29791} }
\put(84.213,-266.009){\fontsize{11.991}{1}\usefont{T1}{cmr}{m}{n}\selectfont\color{color_29791}•}
\put(98.187,-266.69){\fontsize{11.991}{1}\usefont{T1}{cmr}{m}{n}\selectfont\color{color_29791}Me}
\put(114.8545,-266.69){\fontsize{11.991}{1}\usefont{T1}{cmr}{m}{n}\selectfont\color{color_29791}d}
\put(121.5335,-266.69){\fontsize{11.991}{1}\usefont{T1}{cmr}{m}{n}\selectfont\color{color_29791}i}
\put(124.2195,-266.69){\fontsize{11.991}{1}\usefont{T1}{cmr}{m}{n}\selectfont\color{color_29791}u}
\put(130.8985,-266.69){\fontsize{11.991}{1}\usefont{T1}{cmr}{m}{n}\selectfont\color{color_29791}m}
\put(140.791,-266.69){\fontsize{11.991}{1}\usefont{T1}{cmr}{m}{n}\selectfont\color{color_29791} }
\put(144.1725,-266.69){\fontsize{11.991}{1}\usefont{T1}{cmr}{m}{n}\selectfont\color{color_29791}-}
\put(148.0696,-266.69){\fontsize{11.991}{1}\usefont{T1}{cmr}{m}{n}\selectfont\color{color_29791} }
\put(84.213,-294.809){\fontsize{11.991}{1}\usefont{T1}{cmr}{m}{n}\selectfont\color{color_29791}•}
\put(98.187,-295.49){\fontsize{11.991}{1}\usefont{T1}{cmr}{m}{n}\selectfont\color{color_29791}L}
\put(104.866,-295.49){\fontsize{11.991}{1}\usefont{T1}{cmr}{m}{n}\selectfont\color{color_29791}o}
\put(111.6529,-295.49){\fontsize{11.991}{1}\usefont{T1}{cmr}{m}{n}\selectfont\color{color_29791}n}
\put(118.3319,-295.49){\fontsize{11.991}{1}\usefont{T1}{cmr}{m}{n}\selectfont\color{color_29791}g}
\put(125.0109,-295.49){\fontsize{11.991}{1}\usefont{T1}{cmr}{m}{n}\selectfont\color{color_29791} }
\put(128.3923,-295.49){\fontsize{11.991}{1}\usefont{T1}{cmr}{m}{n}\selectfont\color{color_29791}-}
\put(132.2894,-295.49){\fontsize{11.991}{1}\usefont{T1}{cmr}{m}{n}\selectfont\color{color_29791} }
\put(84.213,-323.609){\fontsize{11.991}{1}\usefont{T1}{cmr}{m}{n}\selectfont\color{color_29791}•}
\put(98.187,-324.205){\fontsize{11.991}{1}\usefont{T1}{cmr}{m}{n}\selectfont\color{color_29791}G}
\put(107.564,-324.205){\fontsize{11.991}{1}\usefont{T1}{cmr}{m}{n}\selectfont\color{color_29791}PU}
\put(124.2314,-324.205){\fontsize{11.991}{1}\usefont{T1}{cmr}{m}{n}\selectfont\color{color_29791} }
\put(127.529,-324.205){\fontsize{11.991}{1}\usefont{T1}{cmr}{m}{n}\selectfont\color{color_29791}- }
\put(48.213,-356.888){\fontsize{15.987}{1}\usefont{T1}{cmr}{m}{n}\selectfont\color{color_29791}•}
\put(62.187,-357.795){\fontsize{15.987}{1}\usefont{T1}{cmr}{m}{n}\selectfont\color{color_29791}F}
\put(71.87512,-357.795){\fontsize{15.987}{1}\usefont{T1}{cmr}{m}{n}\selectfont\color{color_29791}ai}
\put(84.34499,-357.795){\fontsize{15.987}{1}\usefont{T1}{cmr}{m}{n}\selectfont\color{color_29791}r}
\put(89.63668,-357.795){\fontsize{15.987}{1}\usefont{T1}{cmr}{m}{n}\selectfont\color{color_29791} }
\put(94.0331,-357.795){\fontsize{15.987}{1}\usefont{T1}{cmr}{m}{n}\selectfont\color{color_29791}U}
\put(105.5278,-357.795){\fontsize{15.987}{1}\usefont{T1}{cmr}{m}{n}\selectfont\color{color_29791}s}
\put(113.6012,-357.795){\fontsize{15.987}{1}\usefont{T1}{cmr}{m}{n}\selectfont\color{color_29791}age}
\put(84.213,-386.709){\fontsize{11.991}{1}\usefont{T1}{cmr}{m}{n}\selectfont\color{color_29791}•}
\put(98.187,-387.389){\fontsize{11.991}{1}\usefont{T1}{cmr}{m}{n}\selectfont\color{color_29791}E}
\put(106.2689,-387.389){\fontsize{11.991}{1}\usefont{T1}{cmr}{m}{n}\selectfont\color{color_29791}n}
\put(112.9479,-387.389){\fontsize{11.991}{1}\usefont{T1}{cmr}{m}{n}\selectfont\color{color_29791}su}
\put(125.6224,-387.389){\fontsize{11.991}{1}\usefont{T1}{cmr}{m}{n}\selectfont\color{color_29791}re}
\put(136.2944,-387.389){\fontsize{11.991}{1}\usefont{T1}{cmr}{m}{n}\selectfont\color{color_29791}s }
\put(145.5874,-387.389){\fontsize{11.991}{1}\usefont{T1}{cmr}{m}{n}\selectfont\color{color_29791}e}
\put(152.2664,-387.389){\fontsize{11.991}{1}\usefont{T1}{cmr}{m}{n}\selectfont\color{color_29791}a}
\put(159.0533,-387.389){\fontsize{11.991}{1}\usefont{T1}{cmr}{m}{n}\selectfont\color{color_29791}ch}
\put(171.7278,-387.389){\fontsize{11.991}{1}\usefont{T1}{cmr}{m}{n}\selectfont\color{color_29791} }
\put(175.1093,-387.389){\fontsize{11.991}{1}\usefont{T1}{cmr}{m}{n}\selectfont\color{color_29791}u}
\put(181.7883,-387.389){\fontsize{11.991}{1}\usefont{T1}{cmr}{m}{n}\selectfont\color{color_29791}se}
\put(194.4628,-387.389){\fontsize{11.991}{1}\usefont{T1}{cmr}{m}{n}\selectfont\color{color_29791}r}
\put(198.3598,-387.389){\fontsize{11.991}{1}\usefont{T1}{cmr}{m}{n}\selectfont\color{color_29791} }
\put(201.7413,-387.389){\fontsize{11.991}{1}\usefont{T1}{cmr}{m}{n}\selectfont\color{color_29791}g}
\put(208.4203,-387.389){\fontsize{11.991}{1}\usefont{T1}{cmr}{m}{n}\selectfont\color{color_29791}e}
\put(215.2072,-387.389){\fontsize{11.991}{1}\usefont{T1}{cmr}{m}{n}\selectfont\color{color_29791}t}
\put(218.5047,-387.389){\fontsize{11.991}{1}\usefont{T1}{cmr}{m}{n}\selectfont\color{color_29791}s }
\put(227.8817,-387.389){\fontsize{11.991}{1}\usefont{T1}{cmr}{m}{n}\selectfont\color{color_29791}f}
\put(231.1792,-387.389){\fontsize{11.991}{1}\usefont{T1}{cmr}{m}{n}\selectfont\color{color_29791}a}
\put(237.8582,-387.389){\fontsize{11.991}{1}\usefont{T1}{cmr}{m}{n}\selectfont\color{color_29791}i}
\put(240.5442,-387.389){\fontsize{11.991}{1}\usefont{T1}{cmr}{m}{n}\selectfont\color{color_29791}r }
\put(247.8347,-387.389){\fontsize{11.991}{1}\usefont{T1}{cmr}{m}{n}\selectfont\color{color_29791}u}
\put(254.5137,-387.389){\fontsize{11.991}{1}\usefont{T1}{cmr}{m}{n}\selectfont\color{color_29791}sa}
\put(267.1882,-387.389){\fontsize{11.991}{1}\usefont{T1}{cmr}{m}{n}\selectfont\color{color_29791}g}
\put(273.9751,-387.389){\fontsize{11.991}{1}\usefont{T1}{cmr}{m}{n}\selectfont\color{color_29791}e}
\put(280.6541,-387.389){\fontsize{11.991}{1}\usefont{T1}{cmr}{m}{n}\selectfont\color{color_29791} }
\put(284.0356,-387.389){\fontsize{11.991}{1}\usefont{T1}{cmr}{m}{n}\selectfont\color{color_29791}o}
\put(290.7146,-387.389){\fontsize{11.991}{1}\usefont{T1}{cmr}{m}{n}\selectfont\color{color_29791}f}
\put(294.0121,-387.389){\fontsize{11.991}{1}\usefont{T1}{cmr}{m}{n}\selectfont\color{color_29791} }
\put(297.3936,-387.389){\fontsize{11.991}{1}\usefont{T1}{cmr}{m}{n}\selectfont\color{color_29791}e}
\put(304.0725,-387.389){\fontsize{11.991}{1}\usefont{T1}{cmr}{m}{n}\selectfont\color{color_29791}a}
\put(310.7515,-387.389){\fontsize{11.991}{1}\usefont{T1}{cmr}{m}{n}\selectfont\color{color_29791}ch}
\put(324.109,-387.389){\fontsize{11.991}{1}\usefont{T1}{cmr}{m}{n}\selectfont\color{color_29791} }
\put(327.397,-387.389){\fontsize{11.991}{1}\usefont{T1}{cmr}{m}{n}\selectfont\color{color_29791}H}
\put(335.9945,-387.389){\fontsize{11.991}{1}\usefont{T1}{cmr}{m}{n}\selectfont\color{color_29791}P}
\put(344.0765,-387.389){\fontsize{11.991}{1}\usefont{T1}{cmr}{m}{n}\selectfont\color{color_29791}C}
\put(352.674,-387.389){\fontsize{11.991}{1}\usefont{T1}{cmr}{m}{n}\selectfont\color{color_29791} }
\put(356.0555,-387.389){\fontsize{11.991}{1}\usefont{T1}{cmr}{m}{n}\selectfont\color{color_29791}q}
\put(362.7345,-387.389){\fontsize{11.991}{1}\usefont{T1}{cmr}{m}{n}\selectfont\color{color_29791}u}
\put(369.4135,-387.389){\fontsize{11.991}{1}\usefont{T1}{cmr}{m}{n}\selectfont\color{color_29791}e}
\put(376.2003,-387.389){\fontsize{11.991}{1}\usefont{T1}{cmr}{m}{n}\selectfont\color{color_29791}u}
\put(382.8793,-387.389){\fontsize{11.991}{1}\usefont{T1}{cmr}{m}{n}\selectfont\color{color_29791}e}
\put(84.213,-415.509){\fontsize{11.991}{1}\usefont{T1}{cmr}{m}{n}\selectfont\color{color_29791}•}
\put(98.187,-416.104){\fontsize{11.991}{1}\usefont{T1}{cmr}{m}{n}\selectfont\color{color_29791}A}
\put(106.2689,-416.104){\fontsize{11.991}{1}\usefont{T1}{cmr}{m}{n}\selectfont\color{color_29791}d}
\put(112.9479,-416.104){\fontsize{11.991}{1}\usefont{T1}{cmr}{m}{n}\selectfont\color{color_29791}j}
\put(115.55,-416.104){\fontsize{11.991}{1}\usefont{T1}{cmr}{m}{n}\selectfont\color{color_29791}u}
\put(122.2289,-416.104){\fontsize{11.991}{1}\usefont{T1}{cmr}{m}{n}\selectfont\color{color_29791}st}
\put(131.6059,-416.104){\fontsize{11.991}{1}\usefont{T1}{cmr}{m}{n}\selectfont\color{color_29791}s }
\put(140.8989,-416.104){\fontsize{11.991}{1}\usefont{T1}{cmr}{m}{n}\selectfont\color{color_29791}p}
\put(147.6859,-416.104){\fontsize{11.991}{1}\usefont{T1}{cmr}{m}{n}\selectfont\color{color_29791}r}
\put(151.5829,-416.104){\fontsize{11.991}{1}\usefont{T1}{cmr}{m}{n}\selectfont\color{color_29791}i}
\put(154.2689,-416.104){\fontsize{11.991}{1}\usefont{T1}{cmr}{m}{n}\selectfont\color{color_29791}o}
\put(160.9479,-416.104){\fontsize{11.991}{1}\usefont{T1}{cmr}{m}{n}\selectfont\color{color_29791}ri}
\put(167.543,-416.104){\fontsize{11.991}{1}\usefont{T1}{cmr}{m}{n}\selectfont\color{color_29791}t}
\put(170.9244,-416.104){\fontsize{11.991}{1}\usefont{T1}{cmr}{m}{n}\selectfont\color{color_29791}i}
\put(173.5265,-416.104){\fontsize{11.991}{1}\usefont{T1}{cmr}{m}{n}\selectfont\color{color_29791}e}
\put(180.2054,-416.104){\fontsize{11.991}{1}\usefont{T1}{cmr}{m}{n}\selectfont\color{color_29791}s }
\put(189.5824,-416.104){\fontsize{11.991}{1}\usefont{T1}{cmr}{m}{n}\selectfont\color{color_29791}o}
\put(196.2614,-416.104){\fontsize{11.991}{1}\usefont{T1}{cmr}{m}{n}\selectfont\color{color_29791}f}
\put(199.6429,-416.104){\fontsize{11.991}{1}\usefont{T1}{cmr}{m}{n}\selectfont\color{color_29791} }
\put(202.9404,-416.104){\fontsize{11.991}{1}\usefont{T1}{cmr}{m}{n}\selectfont\color{color_29791}su}
\put(215.6149,-416.104){\fontsize{11.991}{1}\usefont{T1}{cmr}{m}{n}\selectfont\color{color_29791}b}
\put(222.4018,-416.104){\fontsize{11.991}{1}\usefont{T1}{cmr}{m}{n}\selectfont\color{color_29791}m}
\put(232.2944,-416.104){\fontsize{11.991}{1}\usefont{T1}{cmr}{m}{n}\selectfont\color{color_29791}i}
\put(234.8964,-416.104){\fontsize{11.991}{1}\usefont{T1}{cmr}{m}{n}\selectfont\color{color_29791}t}
\put(238.2779,-416.104){\fontsize{11.991}{1}\usefont{T1}{cmr}{m}{n}\selectfont\color{color_29791}t}
\put(241.6593,-416.104){\fontsize{11.991}{1}\usefont{T1}{cmr}{m}{n}\selectfont\color{color_29791}e}
\put(248.3383,-416.104){\fontsize{11.991}{1}\usefont{T1}{cmr}{m}{n}\selectfont\color{color_29791}d}
\put(255.0173,-416.104){\fontsize{11.991}{1}\usefont{T1}{cmr}{m}{n}\selectfont\color{color_29791} }
\put(258.3988,-416.104){\fontsize{11.991}{1}\usefont{T1}{cmr}{m}{n}\selectfont\color{color_29791}j}
\put(261.0008,-416.104){\fontsize{11.991}{1}\usefont{T1}{cmr}{m}{n}\selectfont\color{color_29791}o}
\put(267.6798,-416.104){\fontsize{11.991}{1}\usefont{T1}{cmr}{m}{n}\selectfont\color{color_29791}b}
\put(274.3588,-416.104){\fontsize{11.991}{1}\usefont{T1}{cmr}{m}{n}\selectfont\color{color_29791}s}
\put(280.4382,-416.104){\fontsize{11.991}{1}\usefont{T1}{cmr}{m}{n}\selectfont\color{color_29791} }
\put(283.7357,-416.104){\fontsize{11.991}{1}\usefont{T1}{cmr}{m}{n}\selectfont\color{color_29791}b}
\put(290.4147,-416.104){\fontsize{11.991}{1}\usefont{T1}{cmr}{m}{n}\selectfont\color{color_29791}a}
\put(297.2016,-416.104){\fontsize{11.991}{1}\usefont{T1}{cmr}{m}{n}\selectfont\color{color_29791}s}
\put(303.1012,-416.104){\fontsize{11.991}{1}\usefont{T1}{cmr}{m}{n}\selectfont\color{color_29791}e}
\put(309.8881,-416.104){\fontsize{11.991}{1}\usefont{T1}{cmr}{m}{n}\selectfont\color{color_29791}d}
\put(316.5671,-416.104){\fontsize{11.991}{1}\usefont{T1}{cmr}{m}{n}\selectfont\color{color_29791} }
\put(319.9485,-416.104){\fontsize{11.991}{1}\usefont{T1}{cmr}{m}{n}\selectfont\color{color_29791}o}
\put(326.6275,-416.104){\fontsize{11.991}{1}\usefont{T1}{cmr}{m}{n}\selectfont\color{color_29791}n}
\put(333.3065,-416.104){\fontsize{11.991}{1}\usefont{T1}{cmr}{m}{n}\selectfont\color{color_29791} }
\put(336.688,-416.104){\fontsize{11.991}{1}\usefont{T1}{cmr}{m}{n}\selectfont\color{color_29791}p}
\put(343.367,-416.104){\fontsize{11.991}{1}\usefont{T1}{cmr}{m}{n}\selectfont\color{color_29791}r}
\put(347.264,-416.104){\fontsize{11.991}{1}\usefont{T1}{cmr}{m}{n}\selectfont\color{color_29791}e}
\put(354.0509,-416.104){\fontsize{11.991}{1}\usefont{T1}{cmr}{m}{n}\selectfont\color{color_29791}v}
\put(359.9505,-416.104){\fontsize{11.991}{1}\usefont{T1}{cmr}{m}{n}\selectfont\color{color_29791}i}
\put(362.6365,-416.104){\fontsize{11.991}{1}\usefont{T1}{cmr}{m}{n}\selectfont\color{color_29791}o}
\put(369.3155,-416.104){\fontsize{11.991}{1}\usefont{T1}{cmr}{m}{n}\selectfont\color{color_29791}u}
\put(375.9944,-416.104){\fontsize{11.991}{1}\usefont{T1}{cmr}{m}{n}\selectfont\color{color_29791}s }
\put(385.3714,-416.104){\fontsize{11.991}{1}\usefont{T1}{cmr}{m}{n}\selectfont\color{color_29791}u}
\put(392.0504,-416.104){\fontsize{11.991}{1}\usefont{T1}{cmr}{m}{n}\selectfont\color{color_29791}sa}
\put(404.7249,-416.104){\fontsize{11.991}{1}\usefont{T1}{cmr}{m}{n}\selectfont\color{color_29791}g}
\put(411.5118,-416.104){\fontsize{11.991}{1}\usefont{T1}{cmr}{m}{n}\selectfont\color{color_29791}e}
\end{picture}
\newpage
\begin{tikzpicture}[overlay]
\path(0pt,0pt);
\filldraw[color_283006][even odd rule]
(-15pt, 10pt) -- (704.972pt, 10pt)
 -- (704.972pt, 10pt)
 -- (704.972pt, -529.972pt)
 -- (704.972pt, -529.972pt)
 -- (-15pt, -529.972pt)
 -- (-15pt, -529.972pt)
 -- (-15pt, 10pt) -- cycle
;
\end{tikzpicture}
\begin{picture}(-5,0)(2.5,0)
\put(-15,-529.972){\includegraphics[width=720pt,height=540pt]{latexImage_70f318a8f2de3f58eebb26eed4f2123e.png}}
\put(48.411,-60.69601){\fontsize{21.997}{1}\usefont{T1}{cmr}{b}{n}\selectfont\color{color_29791}S}
\put(63.083,-60.69601){\fontsize{21.997}{1}\usefont{T1}{cmr}{b}{n}\selectfont\color{color_29791}L}
\put(76.47917,-60.69601){\fontsize{21.997}{1}\usefont{T1}{cmr}{b}{n}\selectfont\color{color_29791}U}
\put(92.27302,-60.69601){\fontsize{21.997}{1}\usefont{T1}{cmr}{b}{n}\selectfont\color{color_29791}RM}
\put(48.213,-116.51){\fontsize{15.987}{1}\usefont{T1}{cmr}{m}{n}\selectfont\color{color_29791}J}
\put(56.28644,-116.51){\fontsize{15.987}{1}\usefont{T1}{cmr}{m}{n}\selectfont\color{color_29791}o}
\put(65.07928,-116.51){\fontsize{15.987}{1}\usefont{T1}{cmr}{m}{n}\selectfont\color{color_29791}b }
\put(78.15665,-116.51){\fontsize{15.987}{1}\usefont{T1}{cmr}{m}{n}\selectfont\color{color_29791}T}
\put(86.94949,-116.51){\fontsize{15.987}{1}\usefont{T1}{cmr}{m}{n}\selectfont\color{color_29791}ypes}
\put(84.213,-145.31){\fontsize{11.991}{1}\usefont{T1}{cmr}{m}{n}\selectfont\color{color_29791}•}
\put(98.187,-145.991){\fontsize{11.991}{1}\usefont{T1}{cmr}{m}{n}\selectfont\color{color_29791}B}
\put(106.2689,-145.991){\fontsize{11.991}{1}\usefont{T1}{cmr}{m}{n}\selectfont\color{color_29791}a}
\put(112.9479,-145.991){\fontsize{11.991}{1}\usefont{T1}{cmr}{m}{n}\selectfont\color{color_29791}t}
\put(116.2454,-145.991){\fontsize{11.991}{1}\usefont{T1}{cmr}{m}{n}\selectfont\color{color_29791}ch}
\put(129.0278,-145.991){\fontsize{11.991}{1}\usefont{T1}{cmr}{m}{n}\selectfont\color{color_29791} }
\put(132.3254,-145.991){\fontsize{11.991}{1}\usefont{T1}{cmr}{m}{n}\selectfont\color{color_29791}–}
\put(139.0044,-145.991){\fontsize{11.991}{1}\usefont{T1}{cmr}{m}{n}\selectfont\color{color_29791} }
\put(142.3858,-145.991){\fontsize{11.991}{1}\usefont{T1}{cmr}{m}{n}\selectfont\color{color_29791}St}
\put(153.7533,-145.991){\fontsize{11.991}{1}\usefont{T1}{cmr}{m}{n}\selectfont\color{color_29791}a}
\put(160.4323,-145.991){\fontsize{11.991}{1}\usefont{T1}{cmr}{m}{n}\selectfont\color{color_29791}n}
\put(167.1113,-145.991){\fontsize{11.991}{1}\usefont{T1}{cmr}{m}{n}\selectfont\color{color_29791}d}
\put(173.7903,-145.991){\fontsize{11.991}{1}\usefont{T1}{cmr}{m}{n}\selectfont\color{color_29791}a}
\put(180.5772,-145.991){\fontsize{11.991}{1}\usefont{T1}{cmr}{m}{n}\selectfont\color{color_29791}r}
\put(184.4742,-145.991){\fontsize{11.991}{1}\usefont{T1}{cmr}{m}{n}\selectfont\color{color_29791}d}
\put(191.1532,-145.991){\fontsize{11.991}{1}\usefont{T1}{cmr}{m}{n}\selectfont\color{color_29791} }
\put(194.5347,-145.991){\fontsize{11.991}{1}\usefont{T1}{cmr}{m}{n}\selectfont\color{color_29791}j}
\put(197.2207,-145.991){\fontsize{11.991}{1}\usefont{T1}{cmr}{m}{n}\selectfont\color{color_29791}o}
\put(203.8997,-145.991){\fontsize{11.991}{1}\usefont{T1}{cmr}{m}{n}\selectfont\color{color_29791}b}
\put(84.213,-174.11){\fontsize{11.991}{1}\usefont{T1}{cmr}{m}{n}\selectfont\color{color_29791}•}
\put(98.187,-174.791){\fontsize{11.991}{1}\usefont{T1}{cmr}{m}{n}\selectfont\color{color_29791}MPI}
\put(111.709,-198.29){\fontsize{7.994}{1}\usefont{T1}{cmr}{m}{n}\selectfont\color{color_29791}•}
\put(134.187,-198.687){\fontsize{7.994}{1}\usefont{T1}{cmr}{m}{n}\selectfont\color{color_29791}W}
\put(141.7733,-198.687){\fontsize{7.994}{1}\usefont{T1}{cmr}{m}{n}\selectfont\color{color_29791}h}
\put(146.17,-198.687){\fontsize{7.994}{1}\usefont{T1}{cmr}{m}{n}\selectfont\color{color_29791}e}
\put(150.6626,-198.687){\fontsize{7.994}{1}\usefont{T1}{cmr}{m}{n}\selectfont\color{color_29791}n}
\put(155.0593,-198.687){\fontsize{7.994}{1}\usefont{T1}{cmr}{m}{n}\selectfont\color{color_29791} }
\put(157.3536,-198.687){\fontsize{7.994}{1}\usefont{T1}{cmr}{m}{n}\selectfont\color{color_29791}y}
\put(161.4465,-198.687){\fontsize{7.994}{1}\usefont{T1}{cmr}{m}{n}\selectfont\color{color_29791}o}
\put(165.8432,-198.687){\fontsize{7.994}{1}\usefont{T1}{cmr}{m}{n}\selectfont\color{color_29791}u}
\put(170.2399,-198.687){\fontsize{7.994}{1}\usefont{T1}{cmr}{m}{n}\selectfont\color{color_29791} }
\put(172.5342,-198.687){\fontsize{7.994}{1}\usefont{T1}{cmr}{m}{n}\selectfont\color{color_29791}r}
\put(175.2282,-198.687){\fontsize{7.994}{1}\usefont{T1}{cmr}{m}{n}\selectfont\color{color_29791}e}
\put(179.6249,-198.687){\fontsize{7.994}{1}\usefont{T1}{cmr}{m}{n}\selectfont\color{color_29791}q}
\put(184.0216,-198.687){\fontsize{7.994}{1}\usefont{T1}{cmr}{m}{n}\selectfont\color{color_29791}u}
\put(188.5142,-198.687){\fontsize{7.994}{1}\usefont{T1}{cmr}{m}{n}\selectfont\color{color_29791}i}
\put(190.3049,-198.687){\fontsize{7.994}{1}\usefont{T1}{cmr}{m}{n}\selectfont\color{color_29791}r}
\put(192.9029,-198.687){\fontsize{7.994}{1}\usefont{T1}{cmr}{m}{n}\selectfont\color{color_29791}e}
\put(197.3956,-198.687){\fontsize{7.994}{1}\usefont{T1}{cmr}{m}{n}\selectfont\color{color_29791} }
\put(199.5939,-198.687){\fontsize{7.994}{1}\usefont{T1}{cmr}{m}{n}\selectfont\color{color_29791}l}
\put(201.3846,-198.687){\fontsize{7.994}{1}\usefont{T1}{cmr}{m}{n}\selectfont\color{color_29791}a}
\put(205.8772,-198.687){\fontsize{7.994}{1}\usefont{T1}{cmr}{m}{n}\selectfont\color{color_29791}r}
\put(208.4752,-198.687){\fontsize{7.994}{1}\usefont{T1}{cmr}{m}{n}\selectfont\color{color_29791}g}
\put(212.8719,-198.687){\fontsize{7.994}{1}\usefont{T1}{cmr}{m}{n}\selectfont\color{color_29791}e}
\put(217.3646,-198.687){\fontsize{7.994}{1}\usefont{T1}{cmr}{m}{n}\selectfont\color{color_29791} }
\put(219.6588,-198.687){\fontsize{7.994}{1}\usefont{T1}{cmr}{m}{n}\selectfont\color{color_29791}a}
\put(224.0555,-198.687){\fontsize{7.994}{1}\usefont{T1}{cmr}{m}{n}\selectfont\color{color_29791}m}
\put(230.7465,-198.687){\fontsize{7.994}{1}\usefont{T1}{cmr}{m}{n}\selectfont\color{color_29791}o}
\put(235.2392,-198.687){\fontsize{7.994}{1}\usefont{T1}{cmr}{m}{n}\selectfont\color{color_29791}u}
\put(239.6358,-198.687){\fontsize{7.994}{1}\usefont{T1}{cmr}{m}{n}\selectfont\color{color_29791}n}
\put(244.1285,-198.687){\fontsize{7.994}{1}\usefont{T1}{cmr}{m}{n}\selectfont\color{color_29791}t}
\put(246.3268,-198.687){\fontsize{7.994}{1}\usefont{T1}{cmr}{m}{n}\selectfont\color{color_29791}s}
\put(250.4198,-198.687){\fontsize{7.994}{1}\usefont{T1}{cmr}{m}{n}\selectfont\color{color_29791} }
\put(252.7141,-198.687){\fontsize{7.994}{1}\usefont{T1}{cmr}{m}{n}\selectfont\color{color_29791}o}
\put(257.1107,-198.687){\fontsize{7.994}{1}\usefont{T1}{cmr}{m}{n}\selectfont\color{color_29791}f}
\put(259.405,-198.687){\fontsize{7.994}{1}\usefont{T1}{cmr}{m}{n}\selectfont\color{color_29791} }
\put(261.6993,-198.687){\fontsize{7.994}{1}\usefont{T1}{cmr}{m}{n}\selectfont\color{color_29791}m}
\put(268.3903,-198.687){\fontsize{7.994}{1}\usefont{T1}{cmr}{m}{n}\selectfont\color{color_29791}e}
\put(272.8829,-198.687){\fontsize{7.994}{1}\usefont{T1}{cmr}{m}{n}\selectfont\color{color_29791}m}
\put(279.5739,-198.687){\fontsize{7.994}{1}\usefont{T1}{cmr}{m}{n}\selectfont\color{color_29791}o}
\put(283.9706,-198.687){\fontsize{7.994}{1}\usefont{T1}{cmr}{m}{n}\selectfont\color{color_29791}r}
\put(286.6646,-198.687){\fontsize{7.994}{1}\usefont{T1}{cmr}{m}{n}\selectfont\color{color_29791}y }
\put(292.9558,-198.687){\fontsize{7.994}{1}\usefont{T1}{cmr}{m}{n}\selectfont\color{color_29791}o}
\put(297.4485,-198.687){\fontsize{7.994}{1}\usefont{T1}{cmr}{m}{n}\selectfont\color{color_29791}r}
\put(300.0465,-198.687){\fontsize{7.994}{1}\usefont{T1}{cmr}{m}{n}\selectfont\color{color_29791} }
\put(302.3408,-198.687){\fontsize{7.994}{1}\usefont{T1}{cmr}{m}{n}\selectfont\color{color_29791}cp}
\put(310.8304,-198.687){\fontsize{7.994}{1}\usefont{T1}{cmr}{m}{n}\selectfont\color{color_29791}u}
\put(315.2271,-198.687){\fontsize{7.994}{1}\usefont{T1}{cmr}{m}{n}\selectfont\color{color_29791} }
\put(317.5214,-198.687){\fontsize{7.994}{1}\usefont{T1}{cmr}{m}{n}\selectfont\color{color_29791}c}
\put(321.6143,-198.687){\fontsize{7.994}{1}\usefont{T1}{cmr}{m}{n}\selectfont\color{color_29791}o}
\put(326.011,-198.687){\fontsize{7.994}{1}\usefont{T1}{cmr}{m}{n}\selectfont\color{color_29791}r}
\put(328.705,-198.687){\fontsize{7.994}{1}\usefont{T1}{cmr}{m}{n}\selectfont\color{color_29791}e}
\put(333.1017,-198.687){\fontsize{7.994}{1}\usefont{T1}{cmr}{m}{n}\selectfont\color{color_29791}s}
\put(337.1946,-198.687){\fontsize{7.994}{1}\usefont{T1}{cmr}{m}{n}\selectfont\color{color_29791}.}
\put(111.709,-221.506){\fontsize{7.994}{1}\usefont{T1}{cmr}{m}{n}\selectfont\color{color_29791}•}
\put(134.187,-221.902){\fontsize{7.994}{1}\usefont{T1}{cmr}{m}{n}\selectfont\color{color_29791}M}
\put(140.878,-221.902){\fontsize{7.994}{1}\usefont{T1}{cmr}{m}{n}\selectfont\color{color_29791}P}
\put(146.2659,-221.902){\fontsize{7.994}{1}\usefont{T1}{cmr}{m}{n}\selectfont\color{color_29791}I}
\put(148.5602,-221.902){\fontsize{7.994}{1}\usefont{T1}{cmr}{m}{n}\selectfont\color{color_29791} }
\put(150.8545,-221.902){\fontsize{7.994}{1}\usefont{T1}{cmr}{m}{n}\selectfont\color{color_29791}r}
\put(153.4525,-221.902){\fontsize{7.994}{1}\usefont{T1}{cmr}{m}{n}\selectfont\color{color_29791}e}
\put(157.9452,-221.902){\fontsize{7.994}{1}\usefont{T1}{cmr}{m}{n}\selectfont\color{color_29791}q}
\put(162.3419,-221.902){\fontsize{7.994}{1}\usefont{T1}{cmr}{m}{n}\selectfont\color{color_29791}u}
\put(166.7386,-221.902){\fontsize{7.994}{1}\usefont{T1}{cmr}{m}{n}\selectfont\color{color_29791}i}
\put(168.6331,-221.902){\fontsize{7.994}{1}\usefont{T1}{cmr}{m}{n}\selectfont\color{color_29791}r}
\put(171.2312,-221.902){\fontsize{7.994}{1}\usefont{T1}{cmr}{m}{n}\selectfont\color{color_29791}e}
\put(175.6279,-221.902){\fontsize{7.994}{1}\usefont{T1}{cmr}{m}{n}\selectfont\color{color_29791} }
\put(177.9222,-221.902){\fontsize{7.994}{1}\usefont{T1}{cmr}{m}{n}\selectfont\color{color_29791}i}
\put(179.7128,-221.902){\fontsize{7.994}{1}\usefont{T1}{cmr}{m}{n}\selectfont\color{color_29791}n}
\put(184.1095,-221.902){\fontsize{7.994}{1}\usefont{T1}{cmr}{m}{n}\selectfont\color{color_29791}f}
\put(186.4038,-221.902){\fontsize{7.994}{1}\usefont{T1}{cmr}{m}{n}\selectfont\color{color_29791}i}
\put(188.1945,-221.902){\fontsize{7.994}{1}\usefont{T1}{cmr}{m}{n}\selectfont\color{color_29791}n}
\put(192.6871,-221.902){\fontsize{7.994}{1}\usefont{T1}{cmr}{m}{n}\selectfont\color{color_29791}i}
\put(194.4777,-221.902){\fontsize{7.994}{1}\usefont{T1}{cmr}{m}{n}\selectfont\color{color_29791}b}
\put(198.8744,-221.902){\fontsize{7.994}{1}\usefont{T1}{cmr}{m}{n}\selectfont\color{color_29791}a}
\put(203.3671,-221.902){\fontsize{7.994}{1}\usefont{T1}{cmr}{m}{n}\selectfont\color{color_29791}n}
\put(207.7638,-221.902){\fontsize{7.994}{1}\usefont{T1}{cmr}{m}{n}\selectfont\color{color_29791}d}
\put(212.1605,-221.902){\fontsize{7.994}{1}\usefont{T1}{cmr}{m}{n}\selectfont\color{color_29791} }
\put(214.4547,-221.902){\fontsize{7.994}{1}\usefont{T1}{cmr}{m}{n}\selectfont\color{color_29791}c}
\put(218.5477,-221.902){\fontsize{7.994}{1}\usefont{T1}{cmr}{m}{n}\selectfont\color{color_29791}o}
\put(222.9444,-221.902){\fontsize{7.994}{1}\usefont{T1}{cmr}{m}{n}\selectfont\color{color_29791}n}
\put(227.437,-221.902){\fontsize{7.994}{1}\usefont{T1}{cmr}{m}{n}\selectfont\color{color_29791}n}
\put(231.8337,-221.902){\fontsize{7.994}{1}\usefont{T1}{cmr}{m}{n}\selectfont\color{color_29791}e}
\put(236.3263,-221.902){\fontsize{7.994}{1}\usefont{T1}{cmr}{m}{n}\selectfont\color{color_29791}ct}
\put(242.6176,-221.902){\fontsize{7.994}{1}\usefont{T1}{cmr}{m}{n}\selectfont\color{color_29791}i}
\put(244.4083,-221.902){\fontsize{7.994}{1}\usefont{T1}{cmr}{m}{n}\selectfont\color{color_29791}v}
\put(248.5012,-221.902){\fontsize{7.994}{1}\usefont{T1}{cmr}{m}{n}\selectfont\color{color_29791}i}
\put(250.2918,-221.902){\fontsize{7.994}{1}\usefont{T1}{cmr}{m}{n}\selectfont\color{color_29791}t}
\put(252.5861,-221.902){\fontsize{7.994}{1}\usefont{T1}{cmr}{m}{n}\selectfont\color{color_29791}y }
\put(258.8774,-221.902){\fontsize{7.994}{1}\usefont{T1}{cmr}{m}{n}\selectfont\color{color_29791}f}
\put(261.1717,-221.902){\fontsize{7.994}{1}\usefont{T1}{cmr}{m}{n}\selectfont\color{color_29791}o}
\put(265.6643,-221.902){\fontsize{7.994}{1}\usefont{T1}{cmr}{m}{n}\selectfont\color{color_29791}r}
\put(268.2624,-221.902){\fontsize{7.994}{1}\usefont{T1}{cmr}{m}{n}\selectfont\color{color_29791} }
\put(270.5566,-221.902){\fontsize{7.994}{1}\usefont{T1}{cmr}{m}{n}\selectfont\color{color_29791}M}
\put(277.2476,-221.902){\fontsize{7.994}{1}\usefont{T1}{cmr}{m}{n}\selectfont\color{color_29791}e}
\put(281.6443,-221.902){\fontsize{7.994}{1}\usefont{T1}{cmr}{m}{n}\selectfont\color{color_29791}s}
\put(285.7372,-221.902){\fontsize{7.994}{1}\usefont{T1}{cmr}{m}{n}\selectfont\color{color_29791}s}
\put(289.8302,-221.902){\fontsize{7.994}{1}\usefont{T1}{cmr}{m}{n}\selectfont\color{color_29791}a}
\put(294.2269,-221.902){\fontsize{7.994}{1}\usefont{T1}{cmr}{m}{n}\selectfont\color{color_29791}g}
\put(298.7195,-221.902){\fontsize{7.994}{1}\usefont{T1}{cmr}{m}{n}\selectfont\color{color_29791}i}
\put(300.5101,-221.902){\fontsize{7.994}{1}\usefont{T1}{cmr}{m}{n}\selectfont\color{color_29791}n}
\put(304.9068,-221.902){\fontsize{7.994}{1}\usefont{T1}{cmr}{m}{n}\selectfont\color{color_29791}g}
\put(111.709,-244.608){\fontsize{7.994}{1}\usefont{T1}{cmr}{m}{n}\selectfont\color{color_29791}•}
\put(134.187,-245.09){\fontsize{7.994}{1}\usefont{T1}{cmr}{m}{n}\selectfont\color{color_29791}A}
\put(139.575,-245.09){\fontsize{7.994}{1}\usefont{T1}{cmr}{m}{n}\selectfont\color{color_29791}l}
\put(141.3656,-245.09){\fontsize{7.994}{1}\usefont{T1}{cmr}{m}{n}\selectfont\color{color_29791}l}
\put(143.1563,-245.09){\fontsize{7.994}{1}\usefont{T1}{cmr}{m}{n}\selectfont\color{color_29791} }
\put(145.4505,-245.09){\fontsize{7.994}{1}\usefont{T1}{cmr}{m}{n}\selectfont\color{color_29791}n}
\put(149.8472,-245.09){\fontsize{7.994}{1}\usefont{T1}{cmr}{m}{n}\selectfont\color{color_29791}o}
\put(154.3399,-245.09){\fontsize{7.994}{1}\usefont{T1}{cmr}{m}{n}\selectfont\color{color_29791}d}
\put(158.7366,-245.09){\fontsize{7.994}{1}\usefont{T1}{cmr}{m}{n}\selectfont\color{color_29791}e}
\put(163.1333,-245.09){\fontsize{7.994}{1}\usefont{T1}{cmr}{m}{n}\selectfont\color{color_29791}s}
\put(167.2262,-245.09){\fontsize{7.994}{1}\usefont{T1}{cmr}{m}{n}\selectfont\color{color_29791} }
\put(169.5205,-245.09){\fontsize{7.994}{1}\usefont{T1}{cmr}{m}{n}\selectfont\color{color_29791}n}
\put(173.9172,-245.09){\fontsize{7.994}{1}\usefont{T1}{cmr}{m}{n}\selectfont\color{color_29791}e}
\put(178.4098,-245.09){\fontsize{7.994}{1}\usefont{T1}{cmr}{m}{n}\selectfont\color{color_29791}e}
\put(182.8065,-245.09){\fontsize{7.994}{1}\usefont{T1}{cmr}{m}{n}\selectfont\color{color_29791}d}
\put(187.2991,-245.09){\fontsize{7.994}{1}\usefont{T1}{cmr}{m}{n}\selectfont\color{color_29791} }
\put(189.5934,-245.09){\fontsize{7.994}{1}\usefont{T1}{cmr}{m}{n}\selectfont\color{color_29791}t}
\put(191.7917,-245.09){\fontsize{7.994}{1}\usefont{T1}{cmr}{m}{n}\selectfont\color{color_29791}o}
\put(196.2844,-245.09){\fontsize{7.994}{1}\usefont{T1}{cmr}{m}{n}\selectfont\color{color_29791} }
\put(198.4827,-245.09){\fontsize{7.994}{1}\usefont{T1}{cmr}{m}{n}\selectfont\color{color_29791}b}
\put(202.9754,-245.09){\fontsize{7.994}{1}\usefont{T1}{cmr}{m}{n}\selectfont\color{color_29791}e}
\put(207.3721,-245.09){\fontsize{7.994}{1}\usefont{T1}{cmr}{m}{n}\selectfont\color{color_29791} }
\put(209.6663,-245.09){\fontsize{7.994}{1}\usefont{T1}{cmr}{m}{n}\selectfont\color{color_29791}i}
\put(211.457,-245.09){\fontsize{7.994}{1}\usefont{T1}{cmr}{m}{n}\selectfont\color{color_29791}n}
\put(215.9496,-245.09){\fontsize{7.994}{1}\usefont{T1}{cmr}{m}{n}\selectfont\color{color_29791} }
\put(218.148,-245.09){\fontsize{7.994}{1}\usefont{T1}{cmr}{m}{n}\selectfont\color{color_29791}t}
\put(220.4422,-245.09){\fontsize{7.994}{1}\usefont{T1}{cmr}{m}{n}\selectfont\color{color_29791}h}
\put(224.9349,-245.09){\fontsize{7.994}{1}\usefont{T1}{cmr}{m}{n}\selectfont\color{color_29791}e}
\put(229.3316,-245.09){\fontsize{7.994}{1}\usefont{T1}{cmr}{m}{n}\selectfont\color{color_29791} }
\put(231.6259,-245.09){\fontsize{7.994}{1}\usefont{T1}{cmr}{m}{n}\selectfont\color{color_29791}s}
\put(235.7188,-245.09){\fontsize{7.994}{1}\usefont{T1}{cmr}{m}{n}\selectfont\color{color_29791}a}
\put(240.1155,-245.09){\fontsize{7.994}{1}\usefont{T1}{cmr}{m}{n}\selectfont\color{color_29791}m}
\put(246.8065,-245.09){\fontsize{7.994}{1}\usefont{T1}{cmr}{m}{n}\selectfont\color{color_29791}e}
\put(251.2991,-245.09){\fontsize{7.994}{1}\usefont{T1}{cmr}{m}{n}\selectfont\color{color_29791} }
\put(253.4974,-245.09){\fontsize{7.994}{1}\usefont{T1}{cmr}{m}{n}\selectfont\color{color_29791}q}
\put(257.9901,-245.09){\fontsize{7.994}{1}\usefont{T1}{cmr}{m}{n}\selectfont\color{color_29791}u}
\put(262.3867,-245.09){\fontsize{7.994}{1}\usefont{T1}{cmr}{m}{n}\selectfont\color{color_29791}e}
\put(266.7834,-245.09){\fontsize{7.994}{1}\usefont{T1}{cmr}{m}{n}\selectfont\color{color_29791}u}
\put(271.2761,-245.09){\fontsize{7.994}{1}\usefont{T1}{cmr}{m}{n}\selectfont\color{color_29791}e}
\put(275.6728,-245.09){\fontsize{7.994}{1}\usefont{T1}{cmr}{m}{n}\selectfont\color{color_29791}.}
\put(84.213,-272.387){\fontsize{11.991}{1}\usefont{T1}{cmr}{m}{n}\selectfont\color{color_29791}•}
\put(98.187,-273.011){\fontsize{11.991}{1}\usefont{T1}{cmr}{m}{n}\selectfont\color{color_29791}A}
\put(106.2689,-273.011){\fontsize{11.991}{1}\usefont{T1}{cmr}{m}{n}\selectfont\color{color_29791}r}
\put(110.166,-273.011){\fontsize{11.991}{1}\usefont{T1}{cmr}{m}{n}\selectfont\color{color_29791}ra}
\put(120.838,-273.011){\fontsize{11.991}{1}\usefont{T1}{cmr}{m}{n}\selectfont\color{color_29791}y }
\put(130.131,-273.011){\fontsize{11.991}{1}\usefont{T1}{cmr}{m}{n}\selectfont\color{color_29791}–}
\put(136.9179,-273.011){\fontsize{11.991}{1}\usefont{T1}{cmr}{m}{n}\selectfont\color{color_29791} }
\put(140.2155,-273.011){\fontsize{11.991}{1}\usefont{T1}{cmr}{m}{n}\selectfont\color{color_29791}W}
\put(151.511,-273.011){\fontsize{11.991}{1}\usefont{T1}{cmr}{m}{n}\selectfont\color{color_29791}h}
\put(158.19,-273.011){\fontsize{11.991}{1}\usefont{T1}{cmr}{m}{n}\selectfont\color{color_29791}e}
\put(164.9769,-273.011){\fontsize{11.991}{1}\usefont{T1}{cmr}{m}{n}\selectfont\color{color_29791}n}
\put(171.6559,-273.011){\fontsize{11.991}{1}\usefont{T1}{cmr}{m}{n}\selectfont\color{color_29791} }
\put(174.9534,-273.011){\fontsize{11.991}{1}\usefont{T1}{cmr}{m}{n}\selectfont\color{color_29791}yo}
\put(187.6279,-273.011){\fontsize{11.991}{1}\usefont{T1}{cmr}{m}{n}\selectfont\color{color_29791}u}
\put(194.4148,-273.011){\fontsize{11.991}{1}\usefont{T1}{cmr}{m}{n}\selectfont\color{color_29791} }
\put(197.7123,-273.011){\fontsize{11.991}{1}\usefont{T1}{cmr}{m}{n}\selectfont\color{color_29791}w}
\put(206.3938,-273.011){\fontsize{11.991}{1}\usefont{T1}{cmr}{m}{n}\selectfont\color{color_29791}a}
\put(213.0728,-273.011){\fontsize{11.991}{1}\usefont{T1}{cmr}{m}{n}\selectfont\color{color_29791}n}
\put(219.7518,-273.011){\fontsize{11.991}{1}\usefont{T1}{cmr}{m}{n}\selectfont\color{color_29791}t}
\put(223.1332,-273.011){\fontsize{11.991}{1}\usefont{T1}{cmr}{m}{n}\selectfont\color{color_29791} }
\put(226.4308,-273.011){\fontsize{11.991}{1}\usefont{T1}{cmr}{m}{n}\selectfont\color{color_29791}t}
\put(229.8122,-273.011){\fontsize{11.991}{1}\usefont{T1}{cmr}{m}{n}\selectfont\color{color_29791}o}
\put(236.4912,-273.011){\fontsize{11.991}{1}\usefont{T1}{cmr}{m}{n}\selectfont\color{color_29791} }
\put(239.8727,-273.011){\fontsize{11.991}{1}\usefont{T1}{cmr}{m}{n}\selectfont\color{color_29791}r}
\put(243.7697,-273.011){\fontsize{11.991}{1}\usefont{T1}{cmr}{m}{n}\selectfont\color{color_29791}u}
\put(250.5566,-273.011){\fontsize{11.991}{1}\usefont{T1}{cmr}{m}{n}\selectfont\color{color_29791}n}
\put(257.2356,-273.011){\fontsize{11.991}{1}\usefont{T1}{cmr}{m}{n}\selectfont\color{color_29791} }
\put(260.5331,-273.011){\fontsize{11.991}{1}\usefont{T1}{cmr}{m}{n}\selectfont\color{color_29791}u}
\put(267.2121,-273.011){\fontsize{11.991}{1}\usefont{T1}{cmr}{m}{n}\selectfont\color{color_29791}p}
\put(273.999,-273.011){\fontsize{11.991}{1}\usefont{T1}{cmr}{m}{n}\selectfont\color{color_29791}t}
\put(277.2965,-273.011){\fontsize{11.991}{1}\usefont{T1}{cmr}{m}{n}\selectfont\color{color_29791}o}
\put(284.0834,-273.011){\fontsize{11.991}{1}\usefont{T1}{cmr}{m}{n}\selectfont\color{color_29791} }
\put(287.381,-273.011){\fontsize{11.991}{1}\usefont{T1}{cmr}{m}{n}\selectfont\color{color_29791}a}
\put(294.0599,-273.011){\fontsize{11.991}{1}\usefont{T1}{cmr}{m}{n}\selectfont\color{color_29791} }
\put(297.4414,-273.011){\fontsize{11.991}{1}\usefont{T1}{cmr}{m}{n}\selectfont\color{color_29791}1}
\put(304.1204,-273.011){\fontsize{11.991}{1}\usefont{T1}{cmr}{m}{n}\selectfont\color{color_29791}0}
\put(310.7994,-273.011){\fontsize{11.991}{1}\usefont{T1}{cmr}{m}{n}\selectfont\color{color_29791}0}
\put(317.4784,-273.011){\fontsize{11.991}{1}\usefont{T1}{cmr}{m}{n}\selectfont\color{color_29791}0}
\put(324.2653,-273.011){\fontsize{11.991}{1}\usefont{T1}{cmr}{m}{n}\selectfont\color{color_29791} }
\put(327.5628,-273.011){\fontsize{11.991}{1}\usefont{T1}{cmr}{m}{n}\selectfont\color{color_29791}sma}
\put(350.2258,-273.011){\fontsize{11.991}{1}\usefont{T1}{cmr}{m}{n}\selectfont\color{color_29791}l}
\put(352.8278,-273.011){\fontsize{11.991}{1}\usefont{T1}{cmr}{m}{n}\selectfont\color{color_29791}l}
\put(355.5138,-273.011){\fontsize{11.991}{1}\usefont{T1}{cmr}{m}{n}\selectfont\color{color_29791} }
\put(358.8953,-273.011){\fontsize{11.991}{1}\usefont{T1}{cmr}{m}{n}\selectfont\color{color_29791}j}
\put(361.4973,-273.011){\fontsize{11.991}{1}\usefont{T1}{cmr}{m}{n}\selectfont\color{color_29791}o}
\put(368.1763,-273.011){\fontsize{11.991}{1}\usefont{T1}{cmr}{m}{n}\selectfont\color{color_29791}b}
\put(374.8553,-273.011){\fontsize{11.991}{1}\usefont{T1}{cmr}{m}{n}\selectfont\color{color_29791}s}
\put(84.213,-301.102){\fontsize{11.991}{1}\usefont{T1}{cmr}{m}{n}\selectfont\color{color_29791}•}
\put(98.187,-301.811){\fontsize{11.991}{1}\usefont{T1}{cmr}{m}{n}\selectfont\color{color_29791}G}
\put(107.564,-301.811){\fontsize{11.991}{1}\usefont{T1}{cmr}{m}{n}\selectfont\color{color_29791}PU}
\put(124.2314,-301.811){\fontsize{11.991}{1}\usefont{T1}{cmr}{m}{n}\selectfont\color{color_29791} }
\put(127.529,-301.811){\fontsize{11.991}{1}\usefont{T1}{cmr}{m}{n}\selectfont\color{color_29791}-}
\put(131.608,-301.811){\fontsize{11.991}{1}\usefont{T1}{cmr}{m}{n}\selectfont\color{color_29791} }
\put(134.896,-301.811){\fontsize{11.991}{1}\usefont{T1}{cmr}{m}{n}\selectfont\color{color_29791} }
\put(138.0976,-301.811){\fontsize{11.991}{1}\usefont{T1}{cmr}{m}{n}\selectfont\color{color_29791}T}
\put(144.0931,-301.811){\fontsize{11.991}{1}\usefont{T1}{cmr}{m}{n}\selectfont\color{color_29791}o}
\put(150.7721,-301.811){\fontsize{11.991}{1}\usefont{T1}{cmr}{m}{n}\selectfont\color{color_29791} }
\put(154.0696,-301.811){\fontsize{11.991}{1}\usefont{T1}{cmr}{m}{n}\selectfont\color{color_29791}u}
\put(160.8565,-301.811){\fontsize{11.991}{1}\usefont{T1}{cmr}{m}{n}\selectfont\color{color_29791}s}
\put(166.7561,-301.811){\fontsize{11.991}{1}\usefont{T1}{cmr}{m}{n}\selectfont\color{color_29791}e}
\put(173.543,-301.811){\fontsize{11.991}{1}\usefont{T1}{cmr}{m}{n}\selectfont\color{color_29791} }
\put(176.8405,-301.811){\fontsize{11.991}{1}\usefont{T1}{cmr}{m}{n}\selectfont\color{color_29791}G}
\put(186.2175,-301.811){\fontsize{11.991}{1}\usefont{T1}{cmr}{m}{n}\selectfont\color{color_29791}PU}
\put(202.885,-301.811){\fontsize{11.991}{1}\usefont{T1}{cmr}{m}{n}\selectfont\color{color_29791}'}
\put(205.1753,-301.811){\fontsize{11.991}{1}\usefont{T1}{cmr}{m}{n}\selectfont\color{color_29791}s }
\put(214.4683,-301.811){\fontsize{11.991}{1}\usefont{T1}{cmr}{m}{n}\selectfont\color{color_29791}yo}
\put(227.1428,-301.811){\fontsize{11.991}{1}\usefont{T1}{cmr}{m}{n}\selectfont\color{color_29791}u}
\put(233.9297,-301.811){\fontsize{11.991}{1}\usefont{T1}{cmr}{m}{n}\selectfont\color{color_29791} }
\put(237.2272,-301.811){\fontsize{11.991}{1}\usefont{T1}{cmr}{m}{n}\selectfont\color{color_29791}mu}
\put(253.8947,-301.811){\fontsize{11.991}{1}\usefont{T1}{cmr}{m}{n}\selectfont\color{color_29791}st}
\put(263.1877,-301.811){\fontsize{11.991}{1}\usefont{T1}{cmr}{m}{n}\selectfont\color{color_29791} }
\put(266.5692,-301.811){\fontsize{11.991}{1}\usefont{T1}{cmr}{m}{n}\selectfont\color{color_29791}i}
\put(269.2552,-301.811){\fontsize{11.991}{1}\usefont{T1}{cmr}{m}{n}\selectfont\color{color_29791}n}
\put(275.9342,-301.811){\fontsize{11.991}{1}\usefont{T1}{cmr}{m}{n}\selectfont\color{color_29791}cl}
\put(284.5317,-301.811){\fontsize{11.991}{1}\usefont{T1}{cmr}{m}{n}\selectfont\color{color_29791}u}
\put(291.2107,-301.811){\fontsize{11.991}{1}\usefont{T1}{cmr}{m}{n}\selectfont\color{color_29791}d}
\put(297.9976,-301.811){\fontsize{11.991}{1}\usefont{T1}{cmr}{m}{n}\selectfont\color{color_29791}e}
\put(304.6766,-301.811){\fontsize{11.991}{1}\usefont{T1}{cmr}{m}{n}\selectfont\color{color_29791} }
\put(307.9741,-301.811){\fontsize{11.991}{1}\usefont{T1}{cmr}{m}{n}\selectfont\color{color_29791}t}
\put(311.3556,-301.811){\fontsize{11.991}{1}\usefont{T1}{cmr}{m}{n}\selectfont\color{color_29791}h}
\put(318.0345,-301.811){\fontsize{11.991}{1}\usefont{T1}{cmr}{m}{n}\selectfont\color{color_29791}e}
\put(324.7135,-301.811){\fontsize{11.991}{1}\usefont{T1}{cmr}{m}{n}\selectfont\color{color_29791} }
\put(328.095,-301.811){\fontsize{11.991}{1}\usefont{T1}{cmr}{m}{n}\selectfont\color{color_29791}–}
\put(334.774,-301.811){\fontsize{11.991}{1}\usefont{T1}{cmr}{m}{n}\selectfont\color{color_29791}g}
\put(341.453,-301.811){\fontsize{11.991}{1}\usefont{T1}{cmr}{m}{n}\selectfont\color{color_29791}re}
\put(352.125,-301.811){\fontsize{11.991}{1}\usefont{T1}{cmr}{m}{n}\selectfont\color{color_29791}s }
\put(361.5019,-301.811){\fontsize{11.991}{1}\usefont{T1}{cmr}{m}{n}\selectfont\color{color_29791}o}
\put(368.1809,-301.811){\fontsize{11.991}{1}\usefont{T1}{cmr}{m}{n}\selectfont\color{color_29791}p}
\put(374.8599,-301.811){\fontsize{11.991}{1}\usefont{T1}{cmr}{m}{n}\selectfont\color{color_29791}t}
\put(378.2414,-301.811){\fontsize{11.991}{1}\usefont{T1}{cmr}{m}{n}\selectfont\color{color_29791}i}
\put(380.8434,-301.811){\fontsize{11.991}{1}\usefont{T1}{cmr}{m}{n}\selectfont\color{color_29791}o}
\put(387.5224,-301.811){\fontsize{11.991}{1}\usefont{T1}{cmr}{m}{n}\selectfont\color{color_29791}n}
\put(394.3093,-301.811){\fontsize{11.991}{1}\usefont{T1}{cmr}{m}{n}\selectfont\color{color_29791} }
\put(397.6068,-301.811){\fontsize{11.991}{1}\usefont{T1}{cmr}{m}{n}\selectfont\color{color_29791}w}
\put(406.2883,-301.811){\fontsize{11.991}{1}\usefont{T1}{cmr}{m}{n}\selectfont\color{color_29791}i}
\put(408.8904,-301.811){\fontsize{11.991}{1}\usefont{T1}{cmr}{m}{n}\selectfont\color{color_29791}t}
\put(412.2718,-301.811){\fontsize{11.991}{1}\usefont{T1}{cmr}{m}{n}\selectfont\color{color_29791}h}
\put(418.9508,-301.811){\fontsize{11.991}{1}\usefont{T1}{cmr}{m}{n}\selectfont\color{color_29791} }
\put(422.2483,-301.811){\fontsize{11.991}{1}\usefont{T1}{cmr}{m}{n}\selectfont\color{color_29791}t}
\put(425.6298,-301.811){\fontsize{11.991}{1}\usefont{T1}{cmr}{m}{n}\selectfont\color{color_29791}h}
\put(432.3088,-301.811){\fontsize{11.991}{1}\usefont{T1}{cmr}{m}{n}\selectfont\color{color_29791}e}
\put(438.9878,-301.811){\fontsize{11.991}{1}\usefont{T1}{cmr}{m}{n}\selectfont\color{color_29791} }
\put(442.3692,-301.811){\fontsize{11.991}{1}\usefont{T1}{cmr}{m}{n}\selectfont\color{color_29791}n}
\put(449.0482,-301.811){\fontsize{11.991}{1}\usefont{T1}{cmr}{m}{n}\selectfont\color{color_29791}u}
\put(455.7272,-301.811){\fontsize{11.991}{1}\usefont{T1}{cmr}{m}{n}\selectfont\color{color_29791}mb}
\put(472.3947,-301.811){\fontsize{11.991}{1}\usefont{T1}{cmr}{m}{n}\selectfont\color{color_29791}e}
\put(479.0737,-301.811){\fontsize{11.991}{1}\usefont{T1}{cmr}{m}{n}\selectfont\color{color_29791}r }
\put(486.3642,-301.811){\fontsize{11.991}{1}\usefont{T1}{cmr}{m}{n}\selectfont\color{color_29791}o}
\put(493.1511,-301.811){\fontsize{11.991}{1}\usefont{T1}{cmr}{m}{n}\selectfont\color{color_29791}f}
\put(496.4486,-301.811){\fontsize{11.991}{1}\usefont{T1}{cmr}{m}{n}\selectfont\color{color_29791} }
\put(499.8301,-301.811){\fontsize{11.991}{1}\usefont{T1}{cmr}{m}{n}\selectfont\color{color_29791}G}
\put(509.207,-301.811){\fontsize{11.991}{1}\usefont{T1}{cmr}{m}{n}\selectfont\color{color_29791}PU}
\put(525.7906,-301.811){\fontsize{11.991}{1}\usefont{T1}{cmr}{m}{n}\selectfont\color{color_29791}'}
\put(528.0809,-301.811){\fontsize{11.991}{1}\usefont{T1}{cmr}{m}{n}\selectfont\color{color_29791}s }
\put(537.4578,-301.811){\fontsize{11.991}{1}\usefont{T1}{cmr}{m}{n}\selectfont\color{color_29791}yo}
\put(550.1323,-301.811){\fontsize{11.991}{1}\usefont{T1}{cmr}{m}{n}\selectfont\color{color_29791}u}
\put(556.8113,-301.811){\fontsize{11.991}{1}\usefont{T1}{cmr}{m}{n}\selectfont\color{color_29791} }
\put(560.1928,-301.811){\fontsize{11.991}{1}\usefont{T1}{cmr}{m}{n}\selectfont\color{color_29791}r}
\put(564.0899,-301.811){\fontsize{11.991}{1}\usefont{T1}{cmr}{m}{n}\selectfont\color{color_29791}e}
\put(570.8768,-301.811){\fontsize{11.991}{1}\usefont{T1}{cmr}{m}{n}\selectfont\color{color_29791}q}
\put(577.5558,-301.811){\fontsize{11.991}{1}\usefont{T1}{cmr}{m}{n}\selectfont\color{color_29791}u}
\put(584.2349,-301.811){\fontsize{11.991}{1}\usefont{T1}{cmr}{m}{n}\selectfont\color{color_29791}i}
\put(586.8369,-301.811){\fontsize{11.991}{1}\usefont{T1}{cmr}{m}{n}\selectfont\color{color_29791}re}
\put(597.5089,-301.811){\fontsize{11.991}{1}\usefont{T1}{cmr}{m}{n}\selectfont\color{color_29791}:}
\put(156.213,-330.611){\fontsize{11.991}{1}\usefont{T1}{cmr}{m}{n}\selectfont\color{color_29791} }
\put(159.614,-330.611){\fontsize{11.991}{1}\usefont{T1}{cmr}{m}{n}\selectfont\color{color_29791}\#SBATCH --gres=gpu:2}
\end{picture}
\newpage
\begin{tikzpicture}[overlay]
\path(0pt,0pt);
\filldraw[color_283006][even odd rule]
(-15pt, 10pt) -- (704.972pt, 10pt)
 -- (704.972pt, 10pt)
 -- (704.972pt, -529.972pt)
 -- (704.972pt, -529.972pt)
 -- (-15pt, -529.972pt)
 -- (-15pt, -529.972pt)
 -- (-15pt, 10pt) -- cycle
;
\end{tikzpicture}
\begin{picture}(-5,0)(2.5,0)
\put(-15,-529.972){\includegraphics[width=720pt,height=540pt]{latexImage_15b9e6acda8be87c69673355bc94d374.png}}
\put(48.411,-60.69601){\fontsize{21.997}{1}\usefont{T1}{cmr}{b}{n}\selectfont\color{color_29791}S}
\put(63.083,-60.69601){\fontsize{21.997}{1}\usefont{T1}{cmr}{b}{n}\selectfont\color{color_29791}L}
\put(76.47917,-60.69601){\fontsize{21.997}{1}\usefont{T1}{cmr}{b}{n}\selectfont\color{color_29791}U}
\put(92.27302,-60.69601){\fontsize{21.997}{1}\usefont{T1}{cmr}{b}{n}\selectfont\color{color_29791}RM}
\put(126.3464,-60.69601){\fontsize{21.997}{1}\usefont{T1}{cmr}{b}{n}\selectfont\color{color_29791} }
\put(132.5275,-60.69601){\fontsize{21.997}{1}\usefont{T1}{cmr}{b}{n}\selectfont\color{color_29791}(}
\put(139.8305,-60.69601){\fontsize{21.997}{1}\usefont{T1}{cmr}{b}{n}\selectfont\color{color_29791}c}
\put(152.0389,-60.69601){\fontsize{21.997}{1}\usefont{T1}{cmr}{b}{n}\selectfont\color{color_29791}o}
\put(165.523,-60.69601){\fontsize{21.997}{1}\usefont{T1}{cmr}{b}{n}\selectfont\color{color_29791}n}
\put(178.9192,-60.69601){\fontsize{21.997}{1}\usefont{T1}{cmr}{b}{n}\selectfont\color{color_29791}t}
\put(186.2222,-60.69601){\fontsize{21.997}{1}\usefont{T1}{cmr}{b}{n}\selectfont\color{color_29791}in}
\put(205.7995,-60.69601){\fontsize{21.997}{1}\usefont{T1}{cmr}{b}{n}\selectfont\color{color_29791}u}
\put(219.1957,-60.69601){\fontsize{21.997}{1}\usefont{T1}{cmr}{b}{n}\selectfont\color{color_29791}e}
\put(231.47,-60.69601){\fontsize{21.997}{1}\usefont{T1}{cmr}{b}{n}\selectfont\color{color_29791}d}
\put(244.8662,-60.69601){\fontsize{21.997}{1}\usefont{T1}{cmr}{b}{n}\selectfont\color{color_29791})}
\put(48.213,-115.49){\fontsize{15.987}{1}\usefont{T1}{cmr}{m}{n}\selectfont\color{color_29791}•}
\put(62.187,-116.51){\fontsize{15.987}{1}\usefont{T1}{cmr}{m}{n}\selectfont\color{color_29791}J}
\put(70.26043,-116.51){\fontsize{15.987}{1}\usefont{T1}{cmr}{m}{n}\selectfont\color{color_29791}o}
\put(79.05328,-116.51){\fontsize{15.987}{1}\usefont{T1}{cmr}{m}{n}\selectfont\color{color_29791}b }
\put(92.41841,-116.51){\fontsize{15.987}{1}\usefont{T1}{cmr}{m}{n}\selectfont\color{color_29791}subm}
\put(131.4747,-116.51){\fontsize{15.987}{1}\usefont{T1}{cmr}{m}{n}\selectfont\color{color_29791}i}
\put(135.0557,-116.51){\fontsize{15.987}{1}\usefont{T1}{cmr}{m}{n}\selectfont\color{color_29791}ss}
\put(151.1227,-116.51){\fontsize{15.987}{1}\usefont{T1}{cmr}{m}{n}\selectfont\color{color_29791}i}
\put(154.7038,-116.51){\fontsize{15.987}{1}\usefont{T1}{cmr}{m}{n}\selectfont\color{color_29791}o}
\put(163.4966,-116.51){\fontsize{15.987}{1}\usefont{T1}{cmr}{m}{n}\selectfont\color{color_29791}n }
\put(176.7818,-116.51){\fontsize{15.987}{1}\usefont{T1}{cmr}{m}{n}\selectfont\color{color_29791}s}
\put(184.8553,-116.51){\fontsize{15.987}{1}\usefont{T1}{cmr}{m}{n}\selectfont\color{color_29791}cr}
\put(198.1405,-116.51){\fontsize{15.987}{1}\usefont{T1}{cmr}{m}{n}\selectfont\color{color_29791}i}
\put(201.7215,-116.51){\fontsize{15.987}{1}\usefont{T1}{cmr}{m}{n}\selectfont\color{color_29791}pt}
\put(215.0067,-116.51){\fontsize{15.987}{1}\usefont{T1}{cmr}{m}{n}\selectfont\color{color_29791}s}
\put(223.0802,-116.51){\fontsize{15.987}{1}\usefont{T1}{cmr}{m}{n}\selectfont\color{color_29791} }
\put(227.4766,-116.51){\fontsize{15.987}{1}\usefont{T1}{cmr}{m}{n}\selectfont\color{color_29791}– }
\put(240.7618,-116.51){\fontsize{15.987}{1}\usefont{T1}{cmr}{m}{n}\selectfont\color{color_29791}S}
\put(251.4411,-116.51){\fontsize{15.987}{1}\usefont{T1}{cmr}{m}{n}\selectfont\color{color_29791}i}
\put(255.0222,-116.51){\fontsize{15.987}{1}\usefont{T1}{cmr}{m}{n}\selectfont\color{color_29791}m}
\put(268.3074,-116.51){\fontsize{15.987}{1}\usefont{T1}{cmr}{m}{n}\selectfont\color{color_29791}pl}
\put(280.7773,-116.51){\fontsize{15.987}{1}\usefont{T1}{cmr}{m}{n}\selectfont\color{color_29791}e:}
\put(48.213,-287.298){\fontsize{15.987}{1}\usefont{T1}{cmr}{m}{n}\selectfont\color{color_29791}•}
\put(62.187,-288.205){\fontsize{15.987}{1}\usefont{T1}{cmr}{m}{n}\selectfont\color{color_29791}U}
\put(73.68165,-288.205){\fontsize{15.987}{1}\usefont{T1}{cmr}{m}{n}\selectfont\color{color_29791}s}
\put(81.75508,-288.205){\fontsize{15.987}{1}\usefont{T1}{cmr}{m}{n}\selectfont\color{color_29791}e}
\put(90.54793,-288.205){\fontsize{15.987}{1}\usefont{T1}{cmr}{m}{n}\selectfont\color{color_29791}f}
\put(95.02428,-288.205){\fontsize{15.987}{1}\usefont{T1}{cmr}{m}{n}\selectfont\color{color_29791}ul}
\put(107.4941,-288.205){\fontsize{15.987}{1}\usefont{T1}{cmr}{m}{n}\selectfont\color{color_29791} }
\put(111.8906,-288.205){\fontsize{15.987}{1}\usefont{T1}{cmr}{m}{n}\selectfont\color{color_29791}O}
\put(124.2805,-288.205){\fontsize{15.987}{1}\usefont{T1}{cmr}{m}{n}\selectfont\color{color_29791}pt}
\put(137.5657,-288.205){\fontsize{15.987}{1}\usefont{T1}{cmr}{m}{n}\selectfont\color{color_29791}i}
\put(141.1468,-288.205){\fontsize{15.987}{1}\usefont{T1}{cmr}{m}{n}\selectfont\color{color_29791}on}
\put(158.8284,-288.205){\fontsize{15.987}{1}\usefont{T1}{cmr}{m}{n}\selectfont\color{color_29791}s}
\put(84.213,-317.09){\fontsize{11.991}{1}\usefont{T1}{cmr}{m}{n}\selectfont\color{color_29791}•}
\put(98.187,-317.798){\fontsize{11.991}{1}\usefont{T1}{cmr}{m}{n}\selectfont\color{color_29791}E}
\put(106.2689,-317.798){\fontsize{11.991}{1}\usefont{T1}{cmr}{m}{n}\selectfont\color{color_29791}x}
\put(112.1685,-317.798){\fontsize{11.991}{1}\usefont{T1}{cmr}{m}{n}\selectfont\color{color_29791}c}
\put(118.2479,-317.798){\fontsize{11.991}{1}\usefont{T1}{cmr}{m}{n}\selectfont\color{color_29791}l}
\put(120.85,-317.798){\fontsize{11.991}{1}\usefont{T1}{cmr}{m}{n}\selectfont\color{color_29791}u}
\put(127.529,-317.798){\fontsize{11.991}{1}\usefont{T1}{cmr}{m}{n}\selectfont\color{color_29791}si}
\put(136.2105,-317.798){\fontsize{11.991}{1}\usefont{T1}{cmr}{m}{n}\selectfont\color{color_29791}ve}
\put(148.8849,-317.798){\fontsize{11.991}{1}\usefont{T1}{cmr}{m}{n}\selectfont\color{color_29791} }
\put(152.1825,-317.798){\fontsize{11.991}{1}\usefont{T1}{cmr}{m}{n}\selectfont\color{color_29791}n}
\put(158.8615,-317.798){\fontsize{11.991}{1}\usefont{T1}{cmr}{m}{n}\selectfont\color{color_29791}o}
\put(165.6484,-317.798){\fontsize{11.991}{1}\usefont{T1}{cmr}{m}{n}\selectfont\color{color_29791}d}
\put(172.3274,-317.798){\fontsize{11.991}{1}\usefont{T1}{cmr}{m}{n}\selectfont\color{color_29791}e}
\put(179.0063,-317.798){\fontsize{11.991}{1}\usefont{T1}{cmr}{m}{n}\selectfont\color{color_29791}:}
\put(182.603,-317.798){\fontsize{11.991}{1}\usefont{T1}{cmr}{m}{n}\selectfont\color{color_29791} }
\put(185.891,-317.798){\fontsize{11.991}{1}\usefont{T1}{cmr}{m}{n}\selectfont\color{color_29791} }
\put(189.293,-317.798){\fontsize{11.991}{1}\usefont{T1}{cmr}{m}{n}\selectfont\color{color_29791} }
\put(192.694,-317.798){\fontsize{11.991}{1}\usefont{T1}{cmr}{m}{n}\selectfont\color{color_29791} }
\put(196.011,-317.798){\fontsize{11.991}{1}\usefont{T1}{cmr}{m}{n}\selectfont\color{color_29791} }
\put(199.413,-317.798){\fontsize{11.991}{1}\usefont{T1}{cmr}{m}{n}\selectfont\color{color_29791} }
\put(202.701,-317.798){\fontsize{11.991}{1}\usefont{T1}{cmr}{m}{n}\selectfont\color{color_29791}\#SBATCH --exclusive}
\put(84.213,-345.89){\fontsize{11.991}{1}\usefont{T1}{cmr}{m}{n}\selectfont\color{color_29791}•}
\put(98.187,-346.598){\fontsize{11.991}{1}\usefont{T1}{cmr}{m}{n}\selectfont\color{color_29791}S}
\put(106.2689,-346.598){\fontsize{11.991}{1}\usefont{T1}{cmr}{m}{n}\selectfont\color{color_29791}p}
\put(112.9479,-346.598){\fontsize{11.991}{1}\usefont{T1}{cmr}{m}{n}\selectfont\color{color_29791}e}
\put(119.6269,-346.598){\fontsize{11.991}{1}\usefont{T1}{cmr}{m}{n}\selectfont\color{color_29791}ci}
\put(128.2245,-346.598){\fontsize{11.991}{1}\usefont{T1}{cmr}{m}{n}\selectfont\color{color_29791}f}
\put(131.6059,-346.598){\fontsize{11.991}{1}\usefont{T1}{cmr}{m}{n}\selectfont\color{color_29791}i}
\put(134.208,-346.598){\fontsize{11.991}{1}\usefont{T1}{cmr}{m}{n}\selectfont\color{color_29791}c }
\put(143.5849,-346.598){\fontsize{11.991}{1}\usefont{T1}{cmr}{m}{n}\selectfont\color{color_29791}n}
\put(150.2639,-346.598){\fontsize{11.991}{1}\usefont{T1}{cmr}{m}{n}\selectfont\color{color_29791}o}
\put(156.9429,-346.598){\fontsize{11.991}{1}\usefont{T1}{cmr}{m}{n}\selectfont\color{color_29791}d}
\put(163.7298,-346.598){\fontsize{11.991}{1}\usefont{T1}{cmr}{m}{n}\selectfont\color{color_29791}e}
\put(170.4088,-346.598){\fontsize{11.991}{1}\usefont{T1}{cmr}{m}{n}\selectfont\color{color_29791}:}
\put(173.901,-346.598){\fontsize{11.991}{1}\usefont{T1}{cmr}{m}{n}\selectfont\color{color_29791} }
\put(177.302,-346.598){\fontsize{11.991}{1}\usefont{T1}{cmr}{m}{n}\selectfont\color{color_29791} }
\put(180.591,-346.598){\fontsize{11.991}{1}\usefont{T1}{cmr}{m}{n}\selectfont\color{color_29791} }
\put(183.992,-346.598){\fontsize{11.991}{1}\usefont{T1}{cmr}{m}{n}\selectfont\color{color_29791} }
\put(187.394,-346.598){\fontsize{11.991}{1}\usefont{T1}{cmr}{m}{n}\selectfont\color{color_29791} }
\put(190.71,-346.598){\fontsize{11.991}{1}\usefont{T1}{cmr}{m}{n}\selectfont\color{color_29791} }
\put(194.112,-346.598){\fontsize{11.991}{1}\usefont{T1}{cmr}{m}{n}\selectfont\color{color_29791} }
\put(197.513,-346.598){\fontsize{11.991}{1}\usefont{T1}{cmr}{m}{n}\selectfont\color{color_29791} }
\put(200.802,-346.598){\fontsize{11.991}{1}\usefont{T1}{cmr}{m}{n}\selectfont\color{color_29791} }
\put(204.203,-346.598){\fontsize{11.991}{1}\usefont{T1}{cmr}{m}{n}\selectfont\color{color_29791}\#SBATCH --nodelist=node022}
\end{picture}
\newpage
\begin{tikzpicture}[overlay]
\path(0pt,0pt);
\filldraw[color_283006][even odd rule]
(-15pt, 10pt) -- (704.972pt, 10pt)
 -- (704.972pt, 10pt)
 -- (704.972pt, -529.972pt)
 -- (704.972pt, -529.972pt)
 -- (-15pt, -529.972pt)
 -- (-15pt, -529.972pt)
 -- (-15pt, 10pt) -- cycle
;
\end{tikzpicture}
\begin{picture}(-5,0)(2.5,0)
\put(-15,-529.972){\includegraphics[width=720pt,height=540pt]{latexImage_43fb9acc39826b86b6493904abaeacaa.png}}
\put(48.411,-60.69601){\fontsize{21.997}{1}\usefont{T1}{cmr}{b}{n}\selectfont\color{color_29791}S}
\put(63.083,-60.69601){\fontsize{21.997}{1}\usefont{T1}{cmr}{b}{n}\selectfont\color{color_29791}L}
\put(76.47917,-60.69601){\fontsize{21.997}{1}\usefont{T1}{cmr}{b}{n}\selectfont\color{color_29791}U}
\put(92.27302,-60.69601){\fontsize{21.997}{1}\usefont{T1}{cmr}{b}{n}\selectfont\color{color_29791}RM}
\put(126.3464,-60.69601){\fontsize{21.997}{1}\usefont{T1}{cmr}{b}{n}\selectfont\color{color_29791} }
\put(132.5275,-60.69601){\fontsize{21.997}{1}\usefont{T1}{cmr}{b}{n}\selectfont\color{color_29791}(}
\put(139.8305,-60.69601){\fontsize{21.997}{1}\usefont{T1}{cmr}{b}{n}\selectfont\color{color_29791}c}
\put(152.0389,-60.69601){\fontsize{21.997}{1}\usefont{T1}{cmr}{b}{n}\selectfont\color{color_29791}o}
\put(165.523,-60.69601){\fontsize{21.997}{1}\usefont{T1}{cmr}{b}{n}\selectfont\color{color_29791}n}
\put(178.9192,-60.69601){\fontsize{21.997}{1}\usefont{T1}{cmr}{b}{n}\selectfont\color{color_29791}t}
\put(186.2222,-60.69601){\fontsize{21.997}{1}\usefont{T1}{cmr}{b}{n}\selectfont\color{color_29791}in}
\put(205.7995,-60.69601){\fontsize{21.997}{1}\usefont{T1}{cmr}{b}{n}\selectfont\color{color_29791}u}
\put(219.1957,-60.69601){\fontsize{21.997}{1}\usefont{T1}{cmr}{b}{n}\selectfont\color{color_29791}e}
\put(231.47,-60.69601){\fontsize{21.997}{1}\usefont{T1}{cmr}{b}{n}\selectfont\color{color_29791}d}
\put(244.8662,-60.69601){\fontsize{21.997}{1}\usefont{T1}{cmr}{b}{n}\selectfont\color{color_29791})}
\put(48.213,-115.49){\fontsize{15.987}{1}\usefont{T1}{cmr}{m}{n}\selectfont\color{color_29791}•}
\put(62.187,-116.51){\fontsize{15.987}{1}\usefont{T1}{cmr}{m}{n}\selectfont\color{color_29791}J}
\put(70.26043,-116.51){\fontsize{15.987}{1}\usefont{T1}{cmr}{m}{n}\selectfont\color{color_29791}o}
\put(79.05328,-116.51){\fontsize{15.987}{1}\usefont{T1}{cmr}{m}{n}\selectfont\color{color_29791}b }
\put(92.41841,-116.51){\fontsize{15.987}{1}\usefont{T1}{cmr}{m}{n}\selectfont\color{color_29791}m}
\put(105.7036,-116.51){\fontsize{15.987}{1}\usefont{T1}{cmr}{m}{n}\selectfont\color{color_29791}on}
\put(123.3852,-116.51){\fontsize{15.987}{1}\usefont{T1}{cmr}{m}{n}\selectfont\color{color_29791}i}
\put(126.9663,-116.51){\fontsize{15.987}{1}\usefont{T1}{cmr}{m}{n}\selectfont\color{color_29791}t}
\put(131.4427,-116.51){\fontsize{15.987}{1}\usefont{T1}{cmr}{m}{n}\selectfont\color{color_29791}or}
\put(145.5272,-116.51){\fontsize{15.987}{1}\usefont{T1}{cmr}{m}{n}\selectfont\color{color_29791}i}
\put(149.1083,-116.51){\fontsize{15.987}{1}\usefont{T1}{cmr}{m}{n}\selectfont\color{color_29791}ng}
\put(84.213,-145.31){\fontsize{11.991}{1}\usefont{T1}{cmr}{m}{n}\selectfont\color{color_29791}•}
\put(98.187,-145.991){\fontsize{11.991}{1}\usefont{T1}{cmr}{m}{n}\selectfont\color{color_29791}sq}
\put(110.8615,-145.991){\fontsize{11.991}{1}\usefont{T1}{cmr}{m}{n}\selectfont\color{color_29791}u}
\put(117.6484,-145.991){\fontsize{11.991}{1}\usefont{T1}{cmr}{m}{n}\selectfont\color{color_29791}e}
\put(124.3274,-145.991){\fontsize{11.991}{1}\usefont{T1}{cmr}{m}{n}\selectfont\color{color_29791}u}
\put(131.0064,-145.991){\fontsize{11.991}{1}\usefont{T1}{cmr}{m}{n}\selectfont\color{color_29791}e}
\put(137.6853,-145.991){\fontsize{11.991}{1}\usefont{T1}{cmr}{m}{n}\selectfont\color{color_29791} }
\put(141.0668,-145.991){\fontsize{11.991}{1}\usefont{T1}{cmr}{m}{n}\selectfont\color{color_29791}-}
\put(144.9639,-145.991){\fontsize{11.991}{1}\usefont{T1}{cmr}{m}{n}\selectfont\color{color_29791}u}
\put(151.6429,-145.991){\fontsize{11.991}{1}\usefont{T1}{cmr}{m}{n}\selectfont\color{color_29791} }
\put(155.0243,-145.991){\fontsize{11.991}{1}\usefont{T1}{cmr}{m}{n}\selectfont\color{color_29791}<u}
\put(168.6941,-145.991){\fontsize{11.991}{1}\usefont{T1}{cmr}{m}{n}\selectfont\color{color_29791}se}
\put(181.3685,-145.991){\fontsize{11.991}{1}\usefont{T1}{cmr}{m}{n}\selectfont\color{color_29791}rn}
\put(192.0405,-145.991){\fontsize{11.991}{1}\usefont{T1}{cmr}{m}{n}\selectfont\color{color_29791}a}
\put(198.7195,-145.991){\fontsize{11.991}{1}\usefont{T1}{cmr}{m}{n}\selectfont\color{color_29791}me}
\put(215.387,-145.991){\fontsize{11.991}{1}\usefont{T1}{cmr}{m}{n}\selectfont\color{color_29791}>}
\put(84.213,-174.11){\fontsize{11.991}{1}\usefont{T1}{cmr}{m}{n}\selectfont\color{color_29791}•}
\put(98.187,-174.791){\fontsize{11.991}{1}\usefont{T1}{cmr}{m}{n}\selectfont\color{color_29791}sa}
\put(110.8615,-174.791){\fontsize{11.991}{1}\usefont{T1}{cmr}{m}{n}\selectfont\color{color_29791}cct}
\put(126.2339,-174.791){\fontsize{11.991}{1}\usefont{T1}{cmr}{m}{n}\selectfont\color{color_29791} }
\put(129.6154,-174.791){\fontsize{11.991}{1}\usefont{T1}{cmr}{m}{n}\selectfont\color{color_29791}-}
\put(133.5125,-174.791){\fontsize{11.991}{1}\usefont{T1}{cmr}{m}{n}\selectfont\color{color_29791}j}
\put(136.1985,-174.791){\fontsize{11.991}{1}\usefont{T1}{cmr}{m}{n}\selectfont\color{color_29791} }
\put(139.496,-174.791){\fontsize{11.991}{1}\usefont{T1}{cmr}{m}{n}\selectfont\color{color_29791}<j}
\put(149.0888,-174.791){\fontsize{11.991}{1}\usefont{T1}{cmr}{m}{n}\selectfont\color{color_29791}o}
\put(155.7678,-174.791){\fontsize{11.991}{1}\usefont{T1}{cmr}{m}{n}\selectfont\color{color_29791}b}
\put(162.4468,-174.791){\fontsize{11.991}{1}\usefont{T1}{cmr}{m}{n}\selectfont\color{color_29791}i}
\put(165.1328,-174.791){\fontsize{11.991}{1}\usefont{T1}{cmr}{m}{n}\selectfont\color{color_29791}d}
\put(171.8118,-174.791){\fontsize{11.991}{1}\usefont{T1}{cmr}{m}{n}\selectfont\color{color_29791}>}
\put(84.213,-202.91){\fontsize{11.991}{1}\usefont{T1}{cmr}{m}{n}\selectfont\color{color_29791}•}
\put(98.187,-203.591){\fontsize{11.991}{1}\usefont{T1}{cmr}{m}{n}\selectfont\color{color_29791}T}
\put(104.1825,-203.591){\fontsize{11.991}{1}\usefont{T1}{cmr}{m}{n}\selectfont\color{color_29791}o}
\put(110.8615,-203.591){\fontsize{11.991}{1}\usefont{T1}{cmr}{m}{n}\selectfont\color{color_29791} }
\put(114.243,-203.591){\fontsize{11.991}{1}\usefont{T1}{cmr}{m}{n}\selectfont\color{color_29791}se}
\put(126.9174,-203.591){\fontsize{11.991}{1}\usefont{T1}{cmr}{m}{n}\selectfont\color{color_29791}e}
\put(133.5964,-203.591){\fontsize{11.991}{1}\usefont{T1}{cmr}{m}{n}\selectfont\color{color_29791} }
\put(136.9779,-203.591){\fontsize{11.991}{1}\usefont{T1}{cmr}{m}{n}\selectfont\color{color_29791}d}
\put(143.6569,-203.591){\fontsize{11.991}{1}\usefont{T1}{cmr}{m}{n}\selectfont\color{color_29791}e}
\put(150.3358,-203.591){\fontsize{11.991}{1}\usefont{T1}{cmr}{m}{n}\selectfont\color{color_29791}t}
\put(153.7173,-203.591){\fontsize{11.991}{1}\usefont{T1}{cmr}{m}{n}\selectfont\color{color_29791}a}
\put(160.3963,-203.591){\fontsize{11.991}{1}\usefont{T1}{cmr}{m}{n}\selectfont\color{color_29791}i}
\put(162.9984,-203.591){\fontsize{11.991}{1}\usefont{T1}{cmr}{m}{n}\selectfont\color{color_29791}l}
\put(165.6843,-203.591){\fontsize{11.991}{1}\usefont{T1}{cmr}{m}{n}\selectfont\color{color_29791}s }
\put(175.0613,-203.591){\fontsize{11.991}{1}\usefont{T1}{cmr}{m}{n}\selectfont\color{color_29791}o}
\put(181.7403,-203.591){\fontsize{11.991}{1}\usefont{T1}{cmr}{m}{n}\selectfont\color{color_29791}n}
\put(188.4193,-203.591){\fontsize{11.991}{1}\usefont{T1}{cmr}{m}{n}\selectfont\color{color_29791} }
\put(191.7168,-203.591){\fontsize{11.991}{1}\usefont{T1}{cmr}{m}{n}\selectfont\color{color_29791}re}
\put(202.3888,-203.591){\fontsize{11.991}{1}\usefont{T1}{cmr}{m}{n}\selectfont\color{color_29791}so}
\put(215.1712,-203.591){\fontsize{11.991}{1}\usefont{T1}{cmr}{m}{n}\selectfont\color{color_29791}u}
\put(221.8502,-203.591){\fontsize{11.991}{1}\usefont{T1}{cmr}{m}{n}\selectfont\color{color_29791}r}
\put(225.7473,-203.591){\fontsize{11.991}{1}\usefont{T1}{cmr}{m}{n}\selectfont\color{color_29791}ce}
\put(238.4218,-203.591){\fontsize{11.991}{1}\usefont{T1}{cmr}{m}{n}\selectfont\color{color_29791}s}
\put(244.5012,-203.591){\fontsize{11.991}{1}\usefont{T1}{cmr}{m}{n}\selectfont\color{color_29791} }
\put(247.7987,-203.591){\fontsize{11.991}{1}\usefont{T1}{cmr}{m}{n}\selectfont\color{color_29791}u}
\put(254.4777,-203.591){\fontsize{11.991}{1}\usefont{T1}{cmr}{m}{n}\selectfont\color{color_29791}se}
\put(267.1522,-203.591){\fontsize{11.991}{1}\usefont{T1}{cmr}{m}{n}\selectfont\color{color_29791}d}
\put(273.9391,-203.591){\fontsize{11.991}{1}\usefont{T1}{cmr}{m}{n}\selectfont\color{color_29791} }
\put(277.2366,-203.591){\fontsize{11.991}{1}\usefont{T1}{cmr}{m}{n}\selectfont\color{color_29791}b}
\put(284.0235,-203.591){\fontsize{11.991}{1}\usefont{T1}{cmr}{m}{n}\selectfont\color{color_29791}y }
\put(293.3165,-203.591){\fontsize{11.991}{1}\usefont{T1}{cmr}{m}{n}\selectfont\color{color_29791}a}
\put(299.9955,-203.591){\fontsize{11.991}{1}\usefont{T1}{cmr}{m}{n}\selectfont\color{color_29791}l}
\put(302.6815,-203.591){\fontsize{11.991}{1}\usefont{T1}{cmr}{m}{n}\selectfont\color{color_29791}l}
\put(305.2835,-203.591){\fontsize{11.991}{1}\usefont{T1}{cmr}{m}{n}\selectfont\color{color_29791} }
\put(308.665,-203.591){\fontsize{11.991}{1}\usefont{T1}{cmr}{m}{n}\selectfont\color{color_29791}r}
\put(312.5621,-203.591){\fontsize{11.991}{1}\usefont{T1}{cmr}{m}{n}\selectfont\color{color_29791}u}
\put(319.349,-203.591){\fontsize{11.991}{1}\usefont{T1}{cmr}{m}{n}\selectfont\color{color_29791}n}
\put(326.028,-203.591){\fontsize{11.991}{1}\usefont{T1}{cmr}{m}{n}\selectfont\color{color_29791}n}
\put(332.7069,-203.591){\fontsize{11.991}{1}\usefont{T1}{cmr}{m}{n}\selectfont\color{color_29791}i}
\put(335.309,-203.591){\fontsize{11.991}{1}\usefont{T1}{cmr}{m}{n}\selectfont\color{color_29791}n}
\put(341.988,-203.591){\fontsize{11.991}{1}\usefont{T1}{cmr}{m}{n}\selectfont\color{color_29791}g}
\put(348.7749,-203.591){\fontsize{11.991}{1}\usefont{T1}{cmr}{m}{n}\selectfont\color{color_29791} }
\put(352.0724,-203.591){\fontsize{11.991}{1}\usefont{T1}{cmr}{m}{n}\selectfont\color{color_29791}j}
\put(354.7584,-203.591){\fontsize{11.991}{1}\usefont{T1}{cmr}{m}{n}\selectfont\color{color_29791}o}
\put(361.4373,-203.591){\fontsize{11.991}{1}\usefont{T1}{cmr}{m}{n}\selectfont\color{color_29791}b}
\put(368.1163,-203.591){\fontsize{11.991}{1}\usefont{T1}{cmr}{m}{n}\selectfont\color{color_29791}s:}
\put(377.4933,-203.591){\fontsize{11.991}{1}\usefont{T1}{cmr}{m}{n}\selectfont\color{color_29791} }
\put(381.51,-203.591){\fontsize{11.991}{1}\usefont{T1}{cmr}{m}{n}\selectfont\color{color_29791}scontrol show jobid –dd <jobid>}
\put(604.71,-203.591){\fontsize{11.991}{1}\usefont{T1}{cmr}{m}{n}\selectfont\color{color_29791} }
\put(611.91,-203.591){\fontsize{11.991}{1}\usefont{T1}{cmr}{m}{n}\selectfont\color{color_29791} }
\put(84.213,-231.71){\fontsize{11.991}{1}\usefont{T1}{cmr}{m}{n}\selectfont\color{color_29791}•}
\put(98.187,-232.306){\fontsize{11.991}{1}\usefont{T1}{cmr}{m}{n}\selectfont\color{color_29791}T}
\put(105.501,-232.306){\fontsize{11.991}{1}\usefont{T1}{cmr}{m}{n}\selectfont\color{color_29791}o}
\put(112.18,-232.306){\fontsize{11.991}{1}\usefont{T1}{cmr}{m}{n}\selectfont\color{color_29791} }
\put(115.5614,-232.306){\fontsize{11.991}{1}\usefont{T1}{cmr}{m}{n}\selectfont\color{color_29791}se}
\put(128.2359,-232.306){\fontsize{11.991}{1}\usefont{T1}{cmr}{m}{n}\selectfont\color{color_29791}e}
\put(134.9149,-232.306){\fontsize{11.991}{1}\usefont{T1}{cmr}{m}{n}\selectfont\color{color_29791} }
\put(138.2964,-232.306){\fontsize{11.991}{1}\usefont{T1}{cmr}{m}{n}\selectfont\color{color_29791}a}
\put(144.9754,-232.306){\fontsize{11.991}{1}\usefont{T1}{cmr}{m}{n}\selectfont\color{color_29791}l}
\put(147.6614,-232.306){\fontsize{11.991}{1}\usefont{T1}{cmr}{m}{n}\selectfont\color{color_29791}l}
\put(150.2634,-232.306){\fontsize{11.991}{1}\usefont{T1}{cmr}{m}{n}\selectfont\color{color_29791} }
\put(153.6449,-232.306){\fontsize{11.991}{1}\usefont{T1}{cmr}{m}{n}\selectfont\color{color_29791}yo}
\put(166.3194,-232.306){\fontsize{11.991}{1}\usefont{T1}{cmr}{m}{n}\selectfont\color{color_29791}u}
\put(172.9983,-232.306){\fontsize{11.991}{1}\usefont{T1}{cmr}{m}{n}\selectfont\color{color_29791}r }
\put(180.2889,-232.306){\fontsize{11.991}{1}\usefont{T1}{cmr}{m}{n}\selectfont\color{color_29791}re}
\put(190.9609,-232.306){\fontsize{11.991}{1}\usefont{T1}{cmr}{m}{n}\selectfont\color{color_29791}ce}
\put(203.6353,-232.306){\fontsize{11.991}{1}\usefont{T1}{cmr}{m}{n}\selectfont\color{color_29791}n}
\put(210.3143,-232.306){\fontsize{11.991}{1}\usefont{T1}{cmr}{m}{n}\selectfont\color{color_29791}t}
\put(213.6958,-232.306){\fontsize{11.991}{1}\usefont{T1}{cmr}{m}{n}\selectfont\color{color_29791} }
\put(217.0773,-232.306){\fontsize{11.991}{1}\usefont{T1}{cmr}{m}{n}\selectfont\color{color_29791}j}
\put(219.6793,-232.306){\fontsize{11.991}{1}\usefont{T1}{cmr}{m}{n}\selectfont\color{color_29791}o}
\put(226.3583,-232.306){\fontsize{11.991}{1}\usefont{T1}{cmr}{m}{n}\selectfont\color{color_29791}b}
\put(233.1452,-232.306){\fontsize{11.991}{1}\usefont{T1}{cmr}{m}{n}\selectfont\color{color_29791}s}
\put(239.0448,-232.306){\fontsize{11.991}{1}\usefont{T1}{cmr}{m}{n}\selectfont\color{color_29791}:}
\put(242.811,-232.306){\fontsize{11.991}{1}\usefont{T1}{cmr}{m}{n}\selectfont\color{color_29791} }
\put(246.213,-232.306){\fontsize{11.991}{1}\usefont{T1}{cmr}{m}{n}\selectfont\color{color_29791}sacct -u <username>}
\put(84.213,-260.397){\fontsize{11.991}{1}\usefont{T1}{cmr}{m}{n}\selectfont\color{color_29791}•}
\put(98.187,-261.106){\fontsize{11.991}{1}\usefont{T1}{cmr}{m}{n}\selectfont\color{color_29791}T}
\put(104.1825,-261.106){\fontsize{11.991}{1}\usefont{T1}{cmr}{m}{n}\selectfont\color{color_29791}o}
\put(110.8615,-261.106){\fontsize{11.991}{1}\usefont{T1}{cmr}{m}{n}\selectfont\color{color_29791} }
\put(114.243,-261.106){\fontsize{11.991}{1}\usefont{T1}{cmr}{m}{n}\selectfont\color{color_29791}ch}
\put(126.9174,-261.106){\fontsize{11.991}{1}\usefont{T1}{cmr}{m}{n}\selectfont\color{color_29791}e}
\put(133.5964,-261.106){\fontsize{11.991}{1}\usefont{T1}{cmr}{m}{n}\selectfont\color{color_29791}ck }
\put(148.9689,-261.106){\fontsize{11.991}{1}\usefont{T1}{cmr}{m}{n}\selectfont\color{color_29791}m}
\put(158.8615,-261.106){\fontsize{11.991}{1}\usefont{T1}{cmr}{m}{n}\selectfont\color{color_29791}e}
\put(165.6484,-261.106){\fontsize{11.991}{1}\usefont{T1}{cmr}{m}{n}\selectfont\color{color_29791}m}
\put(175.541,-261.106){\fontsize{11.991}{1}\usefont{T1}{cmr}{m}{n}\selectfont\color{color_29791}o}
\put(182.3279,-261.106){\fontsize{11.991}{1}\usefont{T1}{cmr}{m}{n}\selectfont\color{color_29791}r}
\put(186.2249,-261.106){\fontsize{11.991}{1}\usefont{T1}{cmr}{m}{n}\selectfont\color{color_29791}y }
\put(195.6019,-261.106){\fontsize{11.991}{1}\usefont{T1}{cmr}{m}{n}\selectfont\color{color_29791}a}
\put(202.2809,-261.106){\fontsize{11.991}{1}\usefont{T1}{cmr}{m}{n}\selectfont\color{color_29791}n}
\put(208.9599,-261.106){\fontsize{11.991}{1}\usefont{T1}{cmr}{m}{n}\selectfont\color{color_29791}d}
\put(215.6389,-261.106){\fontsize{11.991}{1}\usefont{T1}{cmr}{m}{n}\selectfont\color{color_29791} }
\put(219.0203,-261.106){\fontsize{11.991}{1}\usefont{T1}{cmr}{m}{n}\selectfont\color{color_29791}cp}
\put(231.6948,-261.106){\fontsize{11.991}{1}\usefont{T1}{cmr}{m}{n}\selectfont\color{color_29791}u}
\put(238.3738,-261.106){\fontsize{11.991}{1}\usefont{T1}{cmr}{m}{n}\selectfont\color{color_29791} }
\put(241.7552,-261.106){\fontsize{11.991}{1}\usefont{T1}{cmr}{m}{n}\selectfont\color{color_29791}u}
\put(248.4342,-261.106){\fontsize{11.991}{1}\usefont{T1}{cmr}{m}{n}\selectfont\color{color_29791}sa}
\put(261.1087,-261.106){\fontsize{11.991}{1}\usefont{T1}{cmr}{m}{n}\selectfont\color{color_29791}g}
\put(267.8956,-261.106){\fontsize{11.991}{1}\usefont{T1}{cmr}{m}{n}\selectfont\color{color_29791}e}
\put(274.5746,-261.106){\fontsize{11.991}{1}\usefont{T1}{cmr}{m}{n}\selectfont\color{color_29791} }
\put(277.8721,-261.106){\fontsize{11.991}{1}\usefont{T1}{cmr}{m}{n}\selectfont\color{color_29791}o}
\put(284.5511,-261.106){\fontsize{11.991}{1}\usefont{T1}{cmr}{m}{n}\selectfont\color{color_29791}n}
\put(291.338,-261.106){\fontsize{11.991}{1}\usefont{T1}{cmr}{m}{n}\selectfont\color{color_29791} }
\put(294.6355,-261.106){\fontsize{11.991}{1}\usefont{T1}{cmr}{m}{n}\selectfont\color{color_29791}a}
\put(301.4224,-261.106){\fontsize{11.991}{1}\usefont{T1}{cmr}{m}{n}\selectfont\color{color_29791} }
\put(304.7199,-261.106){\fontsize{11.991}{1}\usefont{T1}{cmr}{m}{n}\selectfont\color{color_29791}n}
\put(311.3989,-261.106){\fontsize{11.991}{1}\usefont{T1}{cmr}{m}{n}\selectfont\color{color_29791}o}
\put(318.1858,-261.106){\fontsize{11.991}{1}\usefont{T1}{cmr}{m}{n}\selectfont\color{color_29791}d}
\put(324.8648,-261.106){\fontsize{11.991}{1}\usefont{T1}{cmr}{m}{n}\selectfont\color{color_29791}e}
\put(331.5438,-261.106){\fontsize{11.991}{1}\usefont{T1}{cmr}{m}{n}\selectfont\color{color_29791}:}
\put(335.391,-261.106){\fontsize{11.991}{1}\usefont{T1}{cmr}{m}{n}\selectfont\color{color_29791} }
\put(338.792,-261.106){\fontsize{11.991}{1}\usefont{T1}{cmr}{m}{n}\selectfont\color{color_29791}scontrol show node <node>}
\end{picture}
\newpage
\begin{tikzpicture}[overlay]
\path(0pt,0pt);
\filldraw[color_283006][even odd rule]
(-15pt, 10pt) -- (704.972pt, 10pt)
 -- (704.972pt, 10pt)
 -- (704.972pt, -529.972pt)
 -- (704.972pt, -529.972pt)
 -- (-15pt, -529.972pt)
 -- (-15pt, -529.972pt)
 -- (-15pt, 10pt) -- cycle
;
\end{tikzpicture}
\begin{picture}(-5,0)(2.5,0)
\put(-15,-529.972){\includegraphics[width=720pt,height=540pt]{latexImage_1aef5d9a345e7259c360833b5fc6008f.png}}
\put(48.411,-60.69601){\fontsize{21.997}{1}\usefont{T1}{cmr}{b}{n}\selectfont\color{color_29791}S}
\put(63.083,-60.69601){\fontsize{21.997}{1}\usefont{T1}{cmr}{b}{n}\selectfont\color{color_29791}L}
\put(76.47917,-60.69601){\fontsize{21.997}{1}\usefont{T1}{cmr}{b}{n}\selectfont\color{color_29791}U}
\put(92.27302,-60.69601){\fontsize{21.997}{1}\usefont{T1}{cmr}{b}{n}\selectfont\color{color_29791}RM}
\put(126.3464,-60.69601){\fontsize{21.997}{1}\usefont{T1}{cmr}{b}{n}\selectfont\color{color_29791} }
\put(132.5275,-60.69601){\fontsize{21.997}{1}\usefont{T1}{cmr}{b}{n}\selectfont\color{color_29791}(}
\put(139.8305,-60.69601){\fontsize{21.997}{1}\usefont{T1}{cmr}{b}{n}\selectfont\color{color_29791}c}
\put(152.0389,-60.69601){\fontsize{21.997}{1}\usefont{T1}{cmr}{b}{n}\selectfont\color{color_29791}o}
\put(165.523,-60.69601){\fontsize{21.997}{1}\usefont{T1}{cmr}{b}{n}\selectfont\color{color_29791}n}
\put(178.9192,-60.69601){\fontsize{21.997}{1}\usefont{T1}{cmr}{b}{n}\selectfont\color{color_29791}t}
\put(186.2222,-60.69601){\fontsize{21.997}{1}\usefont{T1}{cmr}{b}{n}\selectfont\color{color_29791}in}
\put(205.7995,-60.69601){\fontsize{21.997}{1}\usefont{T1}{cmr}{b}{n}\selectfont\color{color_29791}u}
\put(219.1957,-60.69601){\fontsize{21.997}{1}\usefont{T1}{cmr}{b}{n}\selectfont\color{color_29791}e}
\put(231.47,-60.69601){\fontsize{21.997}{1}\usefont{T1}{cmr}{b}{n}\selectfont\color{color_29791}d}
\put(244.8662,-60.69601){\fontsize{21.997}{1}\usefont{T1}{cmr}{b}{n}\selectfont\color{color_29791})}
\put(48.213,-115.49){\fontsize{15.987}{1}\usefont{T1}{cmr}{m}{n}\selectfont\color{color_29791}•}
\put(62.187,-116.397){\fontsize{15.987}{1}\usefont{T1}{cmr}{m}{n}\selectfont\color{color_29791}T}
\put(71.37952,-116.397){\fontsize{15.987}{1}\usefont{T1}{cmr}{m}{n}\selectfont\color{color_29791}r}
\put(76.67122,-116.397){\fontsize{15.987}{1}\usefont{T1}{cmr}{m}{n}\selectfont\color{color_29791}oub}
\put(103.2416,-116.397){\fontsize{15.987}{1}\usefont{T1}{cmr}{m}{n}\selectfont\color{color_29791}l}
\put(106.8227,-116.397){\fontsize{15.987}{1}\usefont{T1}{cmr}{m}{n}\selectfont\color{color_29791}es}
\put(123.7849,-116.397){\fontsize{15.987}{1}\usefont{T1}{cmr}{m}{n}\selectfont\color{color_29791}h}
\put(132.5778,-116.397){\fontsize{15.987}{1}\usefont{T1}{cmr}{m}{n}\selectfont\color{color_29791}oot}
\put(154.7517,-116.397){\fontsize{15.987}{1}\usefont{T1}{cmr}{m}{n}\selectfont\color{color_29791}i}
\put(158.3328,-116.397){\fontsize{15.987}{1}\usefont{T1}{cmr}{m}{n}\selectfont\color{color_29791}ng }
\put(180.5867,-116.397){\fontsize{15.987}{1}\usefont{T1}{cmr}{m}{n}\selectfont\color{color_29791}f}
\put(184.9831,-116.397){\fontsize{15.987}{1}\usefont{T1}{cmr}{m}{n}\selectfont\color{color_29791}ai}
\put(197.453,-116.397){\fontsize{15.987}{1}\usefont{T1}{cmr}{m}{n}\selectfont\color{color_29791}l}
\put(201.0341,-116.397){\fontsize{15.987}{1}\usefont{T1}{cmr}{m}{n}\selectfont\color{color_29791}ed}
\put(218.7157,-116.397){\fontsize{15.987}{1}\usefont{T1}{cmr}{m}{n}\selectfont\color{color_29791} }
\put(223.1921,-116.397){\fontsize{15.987}{1}\usefont{T1}{cmr}{m}{n}\selectfont\color{color_29791}j}
\put(226.7732,-116.397){\fontsize{15.987}{1}\usefont{T1}{cmr}{m}{n}\selectfont\color{color_29791}ob}
\put(244.4548,-116.397){\fontsize{15.987}{1}\usefont{T1}{cmr}{m}{n}\selectfont\color{color_29791}s}
\put(84.213,-145.31){\fontsize{11.991}{1}\usefont{T1}{cmr}{m}{n}\selectfont\color{color_29791}•}
\put(98.187,-145.991){\fontsize{11.991}{1}\usefont{T1}{cmr}{m}{n}\selectfont\color{color_29791}Se}
\put(112.9599,-145.991){\fontsize{11.991}{1}\usefont{T1}{cmr}{m}{n}\selectfont\color{color_29791}t}
\put(116.2574,-145.991){\fontsize{11.991}{1}\usefont{T1}{cmr}{m}{n}\selectfont\color{color_29791} }
\put(119.6389,-145.991){\fontsize{11.991}{1}\usefont{T1}{cmr}{m}{n}\selectfont\color{color_29791}a}
\put(126.3179,-145.991){\fontsize{11.991}{1}\usefont{T1}{cmr}{m}{n}\selectfont\color{color_29791}n}
\put(132.9969,-145.991){\fontsize{11.991}{1}\usefont{T1}{cmr}{m}{n}\selectfont\color{color_29791} }
\put(136.3783,-145.991){\fontsize{11.991}{1}\usefont{T1}{cmr}{m}{n}\selectfont\color{color_29791}o}
\put(143.0573,-145.991){\fontsize{11.991}{1}\usefont{T1}{cmr}{m}{n}\selectfont\color{color_29791}u}
\put(149.7363,-145.991){\fontsize{11.991}{1}\usefont{T1}{cmr}{m}{n}\selectfont\color{color_29791}t}
\put(153.1178,-145.991){\fontsize{11.991}{1}\usefont{T1}{cmr}{m}{n}\selectfont\color{color_29791}p}
\put(159.7968,-145.991){\fontsize{11.991}{1}\usefont{T1}{cmr}{m}{n}\selectfont\color{color_29791}u}
\put(166.4757,-145.991){\fontsize{11.991}{1}\usefont{T1}{cmr}{m}{n}\selectfont\color{color_29791}t}
\put(169.8572,-145.991){\fontsize{11.991}{1}\usefont{T1}{cmr}{m}{n}\selectfont\color{color_29791} }
\put(173.1547,-145.991){\fontsize{11.991}{1}\usefont{T1}{cmr}{m}{n}\selectfont\color{color_29791}f}
\put(176.5362,-145.991){\fontsize{11.991}{1}\usefont{T1}{cmr}{m}{n}\selectfont\color{color_29791}i}
\put(179.1382,-145.991){\fontsize{11.991}{1}\usefont{T1}{cmr}{m}{n}\selectfont\color{color_29791}l}
\put(181.8242,-145.991){\fontsize{11.991}{1}\usefont{T1}{cmr}{m}{n}\selectfont\color{color_29791}e}
\put(188.5032,-145.991){\fontsize{11.991}{1}\usefont{T1}{cmr}{m}{n}\selectfont\color{color_29791},}
\put(191.8847,-145.991){\fontsize{11.991}{1}\usefont{T1}{cmr}{m}{n}\selectfont\color{color_29791} }
\put(195.1822,-145.991){\fontsize{11.991}{1}\usefont{T1}{cmr}{m}{n}\selectfont\color{color_29791}o}
\put(201.8612,-145.991){\fontsize{11.991}{1}\usefont{T1}{cmr}{m}{n}\selectfont\color{color_29791}f}
\put(205.2426,-145.991){\fontsize{11.991}{1}\usefont{T1}{cmr}{m}{n}\selectfont\color{color_29791}t}
\put(208.5402,-145.991){\fontsize{11.991}{1}\usefont{T1}{cmr}{m}{n}\selectfont\color{color_29791}e}
\put(215.3271,-145.991){\fontsize{11.991}{1}\usefont{T1}{cmr}{m}{n}\selectfont\color{color_29791}n}
\put(222.0061,-145.991){\fontsize{11.991}{1}\usefont{T1}{cmr}{m}{n}\selectfont\color{color_29791} }
\put(225.3036,-145.991){\fontsize{11.991}{1}\usefont{T1}{cmr}{m}{n}\selectfont\color{color_29791}h}
\put(232.0905,-145.991){\fontsize{11.991}{1}\usefont{T1}{cmr}{m}{n}\selectfont\color{color_29791}a}
\put(238.7695,-145.991){\fontsize{11.991}{1}\usefont{T1}{cmr}{m}{n}\selectfont\color{color_29791}s }
\put(248.0625,-145.991){\fontsize{11.991}{1}\usefont{T1}{cmr}{m}{n}\selectfont\color{color_29791}u}
\put(254.8494,-145.991){\fontsize{11.991}{1}\usefont{T1}{cmr}{m}{n}\selectfont\color{color_29791}se}
\put(267.5239,-145.991){\fontsize{11.991}{1}\usefont{T1}{cmr}{m}{n}\selectfont\color{color_29791}f}
\put(270.9054,-145.991){\fontsize{11.991}{1}\usefont{T1}{cmr}{m}{n}\selectfont\color{color_29791}u}
\put(277.5844,-145.991){\fontsize{11.991}{1}\usefont{T1}{cmr}{m}{n}\selectfont\color{color_29791}l}
\put(280.1864,-145.991){\fontsize{11.991}{1}\usefont{T1}{cmr}{m}{n}\selectfont\color{color_29791} }
\put(283.5679,-145.991){\fontsize{11.991}{1}\usefont{T1}{cmr}{m}{n}\selectfont\color{color_29791}i}
\put(286.17,-145.991){\fontsize{11.991}{1}\usefont{T1}{cmr}{m}{n}\selectfont\color{color_29791}n}
\put(292.8489,-145.991){\fontsize{11.991}{1}\usefont{T1}{cmr}{m}{n}\selectfont\color{color_29791}f}
\put(296.2304,-145.991){\fontsize{11.991}{1}\usefont{T1}{cmr}{m}{n}\selectfont\color{color_29791}o}
\put(302.9094,-145.991){\fontsize{11.991}{1}\usefont{T1}{cmr}{m}{n}\selectfont\color{color_29791}rm}
\put(316.795,-145.991){\fontsize{11.991}{1}\usefont{T1}{cmr}{m}{n}\selectfont\color{color_29791}a}
\put(323.474,-145.991){\fontsize{11.991}{1}\usefont{T1}{cmr}{m}{n}\selectfont\color{color_29791}t}
\put(326.8554,-145.991){\fontsize{11.991}{1}\usefont{T1}{cmr}{m}{n}\selectfont\color{color_29791}i}
\put(329.4575,-145.991){\fontsize{11.991}{1}\usefont{T1}{cmr}{m}{n}\selectfont\color{color_29791}o}
\put(336.2444,-145.991){\fontsize{11.991}{1}\usefont{T1}{cmr}{m}{n}\selectfont\color{color_29791}n}
\put(342.9234,-145.991){\fontsize{11.991}{1}\usefont{T1}{cmr}{m}{n}\selectfont\color{color_29791} }
\put(346.3048,-145.991){\fontsize{11.991}{1}\usefont{T1}{cmr}{m}{n}\selectfont\color{color_29791}w}
\put(354.9024,-145.991){\fontsize{11.991}{1}\usefont{T1}{cmr}{m}{n}\selectfont\color{color_29791}h}
\put(361.5814,-145.991){\fontsize{11.991}{1}\usefont{T1}{cmr}{m}{n}\selectfont\color{color_29791}e}
\put(368.2603,-145.991){\fontsize{11.991}{1}\usefont{T1}{cmr}{m}{n}\selectfont\color{color_29791}n}
\put(375.0472,-145.991){\fontsize{11.991}{1}\usefont{T1}{cmr}{m}{n}\selectfont\color{color_29791} }
\put(378.3448,-145.991){\fontsize{11.991}{1}\usefont{T1}{cmr}{m}{n}\selectfont\color{color_29791}j}
\put(380.9468,-145.991){\fontsize{11.991}{1}\usefont{T1}{cmr}{m}{n}\selectfont\color{color_29791}o}
\put(387.7337,-145.991){\fontsize{11.991}{1}\usefont{T1}{cmr}{m}{n}\selectfont\color{color_29791}b}
\put(394.4127,-145.991){\fontsize{11.991}{1}\usefont{T1}{cmr}{m}{n}\selectfont\color{color_29791}s }
\put(403.7896,-145.991){\fontsize{11.991}{1}\usefont{T1}{cmr}{m}{n}\selectfont\color{color_29791}f}
\put(407.0872,-145.991){\fontsize{11.991}{1}\usefont{T1}{cmr}{m}{n}\selectfont\color{color_29791}a}
\put(413.7661,-145.991){\fontsize{11.991}{1}\usefont{T1}{cmr}{m}{n}\selectfont\color{color_29791}i}
\put(416.4521,-145.991){\fontsize{11.991}{1}\usefont{T1}{cmr}{m}{n}\selectfont\color{color_29791}l}
\put(48.213,-178.589){\fontsize{15.987}{1}\usefont{T1}{cmr}{m}{n}\selectfont\color{color_29791}•}
\put(62.187,-179.609){\fontsize{15.987}{1}\usefont{T1}{cmr}{m}{n}\selectfont\color{color_29791}C}
\put(73.68165,-179.609){\fontsize{15.987}{1}\usefont{T1}{cmr}{m}{n}\selectfont\color{color_29791}om}
\put(95.85562,-179.609){\fontsize{15.987}{1}\usefont{T1}{cmr}{m}{n}\selectfont\color{color_29791}m}
\put(109.1408,-179.609){\fontsize{15.987}{1}\usefont{T1}{cmr}{m}{n}\selectfont\color{color_29791}on }
\put(131.3148,-179.609){\fontsize{15.987}{1}\usefont{T1}{cmr}{m}{n}\selectfont\color{color_29791}m}
\put(144.6959,-179.609){\fontsize{15.987}{1}\usefont{T1}{cmr}{m}{n}\selectfont\color{color_29791}i}
\put(148.277,-179.609){\fontsize{15.987}{1}\usefont{T1}{cmr}{m}{n}\selectfont\color{color_29791}st}
\put(160.6669,-179.609){\fontsize{15.987}{1}\usefont{T1}{cmr}{m}{n}\selectfont\color{color_29791}ak}
\put(177.6291,-179.609){\fontsize{15.987}{1}\usefont{T1}{cmr}{m}{n}\selectfont\color{color_29791}e}
\put(186.422,-179.609){\fontsize{15.987}{1}\usefont{T1}{cmr}{m}{n}\selectfont\color{color_29791}s}
\put(84.213,-208.409){\fontsize{11.991}{1}\usefont{T1}{cmr}{m}{n}\selectfont\color{color_29791}•}
\put(98.187,-209.09){\fontsize{11.991}{1}\usefont{T1}{cmr}{m}{n}\selectfont\color{color_29791}F}
\put(105.4775,-209.09){\fontsize{11.991}{1}\usefont{T1}{cmr}{m}{n}\selectfont\color{color_29791}o}
\put(112.1565,-209.09){\fontsize{11.991}{1}\usefont{T1}{cmr}{m}{n}\selectfont\color{color_29791}rg}
\put(122.8285,-209.09){\fontsize{11.991}{1}\usefont{T1}{cmr}{m}{n}\selectfont\color{color_29791}e}
\put(129.5075,-209.09){\fontsize{11.991}{1}\usefont{T1}{cmr}{m}{n}\selectfont\color{color_29791}t}
\put(132.8889,-209.09){\fontsize{11.991}{1}\usefont{T1}{cmr}{m}{n}\selectfont\color{color_29791}t}
\put(136.1865,-209.09){\fontsize{11.991}{1}\usefont{T1}{cmr}{m}{n}\selectfont\color{color_29791}i}
\put(138.8725,-209.09){\fontsize{11.991}{1}\usefont{T1}{cmr}{m}{n}\selectfont\color{color_29791}n}
\put(145.5515,-209.09){\fontsize{11.991}{1}\usefont{T1}{cmr}{m}{n}\selectfont\color{color_29791}g}
\put(152.2304,-209.09){\fontsize{11.991}{1}\usefont{T1}{cmr}{m}{n}\selectfont\color{color_29791} }
\put(155.6119,-209.09){\fontsize{11.991}{1}\usefont{T1}{cmr}{m}{n}\selectfont\color{color_29791}t}
\put(158.9934,-209.09){\fontsize{11.991}{1}\usefont{T1}{cmr}{m}{n}\selectfont\color{color_29791}o}
\put(165.6724,-209.09){\fontsize{11.991}{1}\usefont{T1}{cmr}{m}{n}\selectfont\color{color_29791} }
\put(168.9699,-209.09){\fontsize{11.991}{1}\usefont{T1}{cmr}{m}{n}\selectfont\color{color_29791}l}
\put(171.6559,-209.09){\fontsize{11.991}{1}\usefont{T1}{cmr}{m}{n}\selectfont\color{color_29791}o}
\put(178.3349,-209.09){\fontsize{11.991}{1}\usefont{T1}{cmr}{m}{n}\selectfont\color{color_29791}a}
\put(185.0139,-209.09){\fontsize{11.991}{1}\usefont{T1}{cmr}{m}{n}\selectfont\color{color_29791}d}
\put(191.6928,-209.09){\fontsize{11.991}{1}\usefont{T1}{cmr}{m}{n}\selectfont\color{color_29791} }
\put(195.0743,-209.09){\fontsize{11.991}{1}\usefont{T1}{cmr}{m}{n}\selectfont\color{color_29791}m}
\put(204.9669,-209.09){\fontsize{11.991}{1}\usefont{T1}{cmr}{m}{n}\selectfont\color{color_29791}o}
\put(211.7538,-209.09){\fontsize{11.991}{1}\usefont{T1}{cmr}{m}{n}\selectfont\color{color_29791}d}
\put(218.4328,-209.09){\fontsize{11.991}{1}\usefont{T1}{cmr}{m}{n}\selectfont\color{color_29791}u}
\put(225.1118,-209.09){\fontsize{11.991}{1}\usefont{T1}{cmr}{m}{n}\selectfont\color{color_29791}l}
\put(227.7978,-209.09){\fontsize{11.991}{1}\usefont{T1}{cmr}{m}{n}\selectfont\color{color_29791}e}
\put(234.4767,-209.09){\fontsize{11.991}{1}\usefont{T1}{cmr}{m}{n}\selectfont\color{color_29791}s}
\put(84.213,-237.209){\fontsize{11.991}{1}\usefont{T1}{cmr}{m}{n}\selectfont\color{color_29791}•}
\put(98.187,-237.89){\fontsize{11.991}{1}\usefont{T1}{cmr}{m}{n}\selectfont\color{color_29791}U}
\put(106.7845,-237.89){\fontsize{11.991}{1}\usefont{T1}{cmr}{m}{n}\selectfont\color{color_29791}si}
\put(115.466,-237.89){\fontsize{11.991}{1}\usefont{T1}{cmr}{m}{n}\selectfont\color{color_29791}n}
\put(122.145,-237.89){\fontsize{11.991}{1}\usefont{T1}{cmr}{m}{n}\selectfont\color{color_29791}g}
\put(128.824,-237.89){\fontsize{11.991}{1}\usefont{T1}{cmr}{m}{n}\selectfont\color{color_29791} }
\put(132.2055,-237.89){\fontsize{11.991}{1}\usefont{T1}{cmr}{m}{n}\selectfont\color{color_29791}st}
\put(141.4985,-237.89){\fontsize{11.991}{1}\usefont{T1}{cmr}{m}{n}\selectfont\color{color_29791}o}
\put(148.2854,-237.89){\fontsize{11.991}{1}\usefont{T1}{cmr}{m}{n}\selectfont\color{color_29791}r}
\put(152.1825,-237.89){\fontsize{11.991}{1}\usefont{T1}{cmr}{m}{n}\selectfont\color{color_29791}a}
\put(158.8615,-237.89){\fontsize{11.991}{1}\usefont{T1}{cmr}{m}{n}\selectfont\color{color_29791}g}
\put(165.5405,-237.89){\fontsize{11.991}{1}\usefont{T1}{cmr}{m}{n}\selectfont\color{color_29791}e}
\put(172.3274,-237.89){\fontsize{11.991}{1}\usefont{T1}{cmr}{m}{n}\selectfont\color{color_29791} }
\put(175.6249,-237.89){\fontsize{11.991}{1}\usefont{T1}{cmr}{m}{n}\selectfont\color{color_29791}w}
\put(184.3064,-237.89){\fontsize{11.991}{1}\usefont{T1}{cmr}{m}{n}\selectfont\color{color_29791}h}
\put(190.9854,-237.89){\fontsize{11.991}{1}\usefont{T1}{cmr}{m}{n}\selectfont\color{color_29791}i}
\put(193.5874,-237.89){\fontsize{11.991}{1}\usefont{T1}{cmr}{m}{n}\selectfont\color{color_29791}ch}
\put(206.3698,-237.89){\fontsize{11.991}{1}\usefont{T1}{cmr}{m}{n}\selectfont\color{color_29791} }
\put(209.6674,-237.89){\fontsize{11.991}{1}\usefont{T1}{cmr}{m}{n}\selectfont\color{color_29791}i}
\put(212.2694,-237.89){\fontsize{11.991}{1}\usefont{T1}{cmr}{m}{n}\selectfont\color{color_29791}s }
\put(221.6464,-237.89){\fontsize{11.991}{1}\usefont{T1}{cmr}{m}{n}\selectfont\color{color_29791}n}
\put(228.3254,-237.89){\fontsize{11.991}{1}\usefont{T1}{cmr}{m}{n}\selectfont\color{color_29791}o}
\put(235.1123,-237.89){\fontsize{11.991}{1}\usefont{T1}{cmr}{m}{n}\selectfont\color{color_29791}t}
\put(238.4098,-237.89){\fontsize{11.991}{1}\usefont{T1}{cmr}{m}{n}\selectfont\color{color_29791} }
\put(241.7913,-237.89){\fontsize{11.991}{1}\usefont{T1}{cmr}{m}{n}\selectfont\color{color_29791}vi}
\put(250.3888,-237.89){\fontsize{11.991}{1}\usefont{T1}{cmr}{m}{n}\selectfont\color{color_29791}si}
\put(258.9864,-237.89){\fontsize{11.991}{1}\usefont{T1}{cmr}{m}{n}\selectfont\color{color_29791}b}
\put(265.7733,-237.89){\fontsize{11.991}{1}\usefont{T1}{cmr}{m}{n}\selectfont\color{color_29791}l}
\put(268.3753,-237.89){\fontsize{11.991}{1}\usefont{T1}{cmr}{m}{n}\selectfont\color{color_29791}e}
\put(275.0543,-237.89){\fontsize{11.991}{1}\usefont{T1}{cmr}{m}{n}\selectfont\color{color_29791} }
\put(278.4358,-237.89){\fontsize{11.991}{1}\usefont{T1}{cmr}{m}{n}\selectfont\color{color_29791}t}
\put(281.8172,-237.89){\fontsize{11.991}{1}\usefont{T1}{cmr}{m}{n}\selectfont\color{color_29791}o}
\put(288.4962,-237.89){\fontsize{11.991}{1}\usefont{T1}{cmr}{m}{n}\selectfont\color{color_29791} }
\put(291.7937,-237.89){\fontsize{11.991}{1}\usefont{T1}{cmr}{m}{n}\selectfont\color{color_29791}t}
\put(295.1752,-237.89){\fontsize{11.991}{1}\usefont{T1}{cmr}{m}{n}\selectfont\color{color_29791}h}
\put(301.8542,-237.89){\fontsize{11.991}{1}\usefont{T1}{cmr}{m}{n}\selectfont\color{color_29791}e}
\put(308.5332,-237.89){\fontsize{11.991}{1}\usefont{T1}{cmr}{m}{n}\selectfont\color{color_29791} }
\put(311.9146,-237.89){\fontsize{11.991}{1}\usefont{T1}{cmr}{m}{n}\selectfont\color{color_29791}H}
\put(320.5122,-237.89){\fontsize{11.991}{1}\usefont{T1}{cmr}{m}{n}\selectfont\color{color_29791}P}
\put(328.5941,-237.89){\fontsize{11.991}{1}\usefont{T1}{cmr}{m}{n}\selectfont\color{color_29791}C}
\put(337.1917,-237.89){\fontsize{11.991}{1}\usefont{T1}{cmr}{m}{n}\selectfont\color{color_29791} }
\put(340.5731,-237.89){\fontsize{11.991}{1}\usefont{T1}{cmr}{m}{n}\selectfont\color{color_29791}n}
\put(347.2521,-237.89){\fontsize{11.991}{1}\usefont{T1}{cmr}{m}{n}\selectfont\color{color_29791}o}
\put(353.9311,-237.89){\fontsize{11.991}{1}\usefont{T1}{cmr}{m}{n}\selectfont\color{color_29791}d}
\put(360.6101,-237.89){\fontsize{11.991}{1}\usefont{T1}{cmr}{m}{n}\selectfont\color{color_29791}e}
\put(367.397,-237.89){\fontsize{11.991}{1}\usefont{T1}{cmr}{m}{n}\selectfont\color{color_29791}s}
\put(373.2965,-237.89){\fontsize{11.991}{1}\usefont{T1}{cmr}{m}{n}\selectfont\color{color_29791} }
\put(376.678,-237.89){\fontsize{11.991}{1}\usefont{T1}{cmr}{m}{n}\selectfont\color{color_29791}(u}
\put(387.35,-237.89){\fontsize{11.991}{1}\usefont{T1}{cmr}{m}{n}\selectfont\color{color_29791}se}
\put(400.0245,-237.89){\fontsize{11.991}{1}\usefont{T1}{cmr}{m}{n}\selectfont\color{color_29791} }
\put(403.406,-237.89){\fontsize{11.991}{1}\usefont{T1}{cmr}{m}{n}\selectfont\color{color_29791}e}
\put(410.085,-237.89){\fontsize{11.991}{1}\usefont{T1}{cmr}{m}{n}\selectfont\color{color_29791}i}
\put(412.687,-237.89){\fontsize{11.991}{1}\usefont{T1}{cmr}{m}{n}\selectfont\color{color_29791}t}
\put(416.0685,-237.89){\fontsize{11.991}{1}\usefont{T1}{cmr}{m}{n}\selectfont\color{color_29791}h}
\put(422.7475,-237.89){\fontsize{11.991}{1}\usefont{T1}{cmr}{m}{n}\selectfont\color{color_29791}e}
\put(429.4265,-237.89){\fontsize{11.991}{1}\usefont{T1}{cmr}{m}{n}\selectfont\color{color_29791}r }
\put(436.8009,-237.89){\fontsize{11.991}{1}\usefont{T1}{cmr}{m}{n}\selectfont\color{color_29791}/}
\put(440.0984,-237.89){\fontsize{11.991}{1}\usefont{T1}{cmr}{m}{n}\selectfont\color{color_29791}d}
\put(446.7774,-237.89){\fontsize{11.991}{1}\usefont{T1}{cmr}{m}{n}\selectfont\color{color_29791}a}
\put(453.5643,-237.89){\fontsize{11.991}{1}\usefont{T1}{cmr}{m}{n}\selectfont\color{color_29791}t}
\put(456.8618,-237.89){\fontsize{11.991}{1}\usefont{T1}{cmr}{m}{n}\selectfont\color{color_29791}a}
\put(463.5408,-237.89){\fontsize{11.991}{1}\usefont{T1}{cmr}{m}{n}\selectfont\color{color_29791}/}
\put(466.9223,-237.89){\fontsize{11.991}{1}\usefont{T1}{cmr}{m}{n}\selectfont\color{color_29791}h}
\put(473.6013,-237.89){\fontsize{11.991}{1}\usefont{T1}{cmr}{m}{n}\selectfont\color{color_29791}p}
\put(480.2802,-237.89){\fontsize{11.991}{1}\usefont{T1}{cmr}{m}{n}\selectfont\color{color_29791}cd}
\put(493.0627,-237.89){\fontsize{11.991}{1}\usefont{T1}{cmr}{m}{n}\selectfont\color{color_29791}a}
\put(499.7416,-237.89){\fontsize{11.991}{1}\usefont{T1}{cmr}{m}{n}\selectfont\color{color_29791}t}
\put(503.0392,-237.89){\fontsize{11.991}{1}\usefont{T1}{cmr}{m}{n}\selectfont\color{color_29791}a}
\put(509.7182,-237.89){\fontsize{11.991}{1}\usefont{T1}{cmr}{m}{n}\selectfont\color{color_29791} }
\put(513.0997,-237.89){\fontsize{11.991}{1}\usefont{T1}{cmr}{m}{n}\selectfont\color{color_29791}o}
\put(519.7787,-237.89){\fontsize{11.991}{1}\usefont{T1}{cmr}{m}{n}\selectfont\color{color_29791}r }
\put(527.0692,-237.89){\fontsize{11.991}{1}\usefont{T1}{cmr}{m}{n}\selectfont\color{color_29791}/}
\put(530.4507,-237.89){\fontsize{11.991}{1}\usefont{T1}{cmr}{m}{n}\selectfont\color{color_29791}d}
\put(537.1297,-237.89){\fontsize{11.991}{1}\usefont{T1}{cmr}{m}{n}\selectfont\color{color_29791}a}
\put(543.9166,-237.89){\fontsize{11.991}{1}\usefont{T1}{cmr}{m}{n}\selectfont\color{color_29791}t}
\put(547.2142,-237.89){\fontsize{11.991}{1}\usefont{T1}{cmr}{m}{n}\selectfont\color{color_29791}a}
\put(553.8932,-237.89){\fontsize{11.991}{1}\usefont{T1}{cmr}{m}{n}\selectfont\color{color_29791}/}
\put(557.2747,-237.89){\fontsize{11.991}{1}\usefont{T1}{cmr}{m}{n}\selectfont\color{color_29791}h}
\put(563.9537,-237.89){\fontsize{11.991}{1}\usefont{T1}{cmr}{m}{n}\selectfont\color{color_29791}p}
\put(570.6327,-237.89){\fontsize{11.991}{1}\usefont{T1}{cmr}{m}{n}\selectfont\color{color_29791}cf}
\put(580.0096,-237.89){\fontsize{11.991}{1}\usefont{T1}{cmr}{m}{n}\selectfont\color{color_29791}l}
\put(582.6117,-237.89){\fontsize{11.991}{1}\usefont{T1}{cmr}{m}{n}\selectfont\color{color_29791}a}
\put(589.2907,-237.89){\fontsize{11.991}{1}\usefont{T1}{cmr}{m}{n}\selectfont\color{color_29791}sh}
\put(602.0731,-237.89){\fontsize{11.991}{1}\usefont{T1}{cmr}{m}{n}\selectfont\color{color_29791})}
\put(84.213,-266.009){\fontsize{11.991}{1}\usefont{T1}{cmr}{m}{n}\selectfont\color{color_29791}•}
\put(98.187,-266.69){\fontsize{11.991}{1}\usefont{T1}{cmr}{m}{n}\selectfont\color{color_29791}A}
\put(105.9811,-266.69){\fontsize{11.991}{1}\usefont{T1}{cmr}{m}{n}\selectfont\color{color_29791}vo}
\put(118.6556,-266.69){\fontsize{11.991}{1}\usefont{T1}{cmr}{m}{n}\selectfont\color{color_29791}i}
\put(121.2577,-266.69){\fontsize{11.991}{1}\usefont{T1}{cmr}{m}{n}\selectfont\color{color_29791}d}
\put(128.0446,-266.69){\fontsize{11.991}{1}\usefont{T1}{cmr}{m}{n}\selectfont\color{color_29791} }
\put(131.3421,-266.69){\fontsize{11.991}{1}\usefont{T1}{cmr}{m}{n}\selectfont\color{color_29791}u}
\put(138.129,-266.69){\fontsize{11.991}{1}\usefont{T1}{cmr}{m}{n}\selectfont\color{color_29791}s}
\put(144.0286,-266.69){\fontsize{11.991}{1}\usefont{T1}{cmr}{m}{n}\selectfont\color{color_29791}i}
\put(146.7146,-266.69){\fontsize{11.991}{1}\usefont{T1}{cmr}{m}{n}\selectfont\color{color_29791}n}
\put(153.3936,-266.69){\fontsize{11.991}{1}\usefont{T1}{cmr}{m}{n}\selectfont\color{color_29791}g}
\put(160.0726,-266.69){\fontsize{11.991}{1}\usefont{T1}{cmr}{m}{n}\selectfont\color{color_29791} }
\put(163.454,-266.69){\fontsize{11.991}{1}\usefont{T1}{cmr}{m}{n}\selectfont\color{color_29791}syml}
\put(188.0356,-266.69){\fontsize{11.991}{1}\usefont{T1}{cmr}{m}{n}\selectfont\color{color_29791}i}
\put(190.6376,-266.69){\fontsize{11.991}{1}\usefont{T1}{cmr}{m}{n}\selectfont\color{color_29791}n}
\put(197.3166,-266.69){\fontsize{11.991}{1}\usefont{T1}{cmr}{m}{n}\selectfont\color{color_29791}k}
\put(203.3961,-266.69){\fontsize{11.991}{1}\usefont{T1}{cmr}{m}{n}\selectfont\color{color_29791}s}
\put(84.213,-294.809){\fontsize{11.991}{1}\usefont{T1}{cmr}{m}{n}\selectfont\color{color_29791}•}
\put(98.187,-295.405){\fontsize{11.991}{1}\usefont{T1}{cmr}{m}{n}\selectfont\color{color_29791}D}
\put(106.7845,-295.405){\fontsize{11.991}{1}\usefont{T1}{cmr}{m}{n}\selectfont\color{color_29791}i}
\put(109.4705,-295.405){\fontsize{11.991}{1}\usefont{T1}{cmr}{m}{n}\selectfont\color{color_29791}d}
\put(116.1495,-295.405){\fontsize{11.991}{1}\usefont{T1}{cmr}{m}{n}\selectfont\color{color_29791} }
\put(119.531,-295.405){\fontsize{11.991}{1}\usefont{T1}{cmr}{m}{n}\selectfont\color{color_29791}y}
\put(125.4305,-295.405){\fontsize{11.991}{1}\usefont{T1}{cmr}{m}{n}\selectfont\color{color_29791}o}
\put(132.2175,-295.405){\fontsize{11.991}{1}\usefont{T1}{cmr}{m}{n}\selectfont\color{color_29791}u}
\put(138.8964,-295.405){\fontsize{11.991}{1}\usefont{T1}{cmr}{m}{n}\selectfont\color{color_29791} }
\put(142.2779,-295.405){\fontsize{11.991}{1}\usefont{T1}{cmr}{m}{n}\selectfont\color{color_29791}l}
\put(144.8799,-295.405){\fontsize{11.991}{1}\usefont{T1}{cmr}{m}{n}\selectfont\color{color_29791}o}
\put(151.5589,-295.405){\fontsize{11.991}{1}\usefont{T1}{cmr}{m}{n}\selectfont\color{color_29791}a}
\put(158.2379,-295.405){\fontsize{11.991}{1}\usefont{T1}{cmr}{m}{n}\selectfont\color{color_29791}d}
\put(164.9169,-295.405){\fontsize{11.991}{1}\usefont{T1}{cmr}{m}{n}\selectfont\color{color_29791} }
\put(168.2984,-295.405){\fontsize{11.991}{1}\usefont{T1}{cmr}{m}{n}\selectfont\color{color_29791}a}
\put(174.9774,-295.405){\fontsize{11.991}{1}\usefont{T1}{cmr}{m}{n}\selectfont\color{color_29791}n}
\put(181.7643,-295.405){\fontsize{11.991}{1}\usefont{T1}{cmr}{m}{n}\selectfont\color{color_29791}y }
\put(191.0573,-295.405){\fontsize{11.991}{1}\usefont{T1}{cmr}{m}{n}\selectfont\color{color_29791}re}
\put(201.7293,-295.405){\fontsize{11.991}{1}\usefont{T1}{cmr}{m}{n}\selectfont\color{color_29791}q}
\put(208.4083,-295.405){\fontsize{11.991}{1}\usefont{T1}{cmr}{m}{n}\selectfont\color{color_29791}u}
\put(215.0873,-295.405){\fontsize{11.991}{1}\usefont{T1}{cmr}{m}{n}\selectfont\color{color_29791}i}
\put(217.7733,-295.405){\fontsize{11.991}{1}\usefont{T1}{cmr}{m}{n}\selectfont\color{color_29791}r}
\put(221.6703,-295.405){\fontsize{11.991}{1}\usefont{T1}{cmr}{m}{n}\selectfont\color{color_29791}e}
\put(228.4572,-295.405){\fontsize{11.991}{1}\usefont{T1}{cmr}{m}{n}\selectfont\color{color_29791} }
\put(231.7548,-295.405){\fontsize{11.991}{1}\usefont{T1}{cmr}{m}{n}\selectfont\color{color_29791}mo}
\put(248.4223,-295.405){\fontsize{11.991}{1}\usefont{T1}{cmr}{m}{n}\selectfont\color{color_29791}d}
\put(255.1013,-295.405){\fontsize{11.991}{1}\usefont{T1}{cmr}{m}{n}\selectfont\color{color_29791}u}
\put(261.7802,-295.405){\fontsize{11.991}{1}\usefont{T1}{cmr}{m}{n}\selectfont\color{color_29791}l}
\put(264.4662,-295.405){\fontsize{11.991}{1}\usefont{T1}{cmr}{m}{n}\selectfont\color{color_29791}e}
\put(271.1452,-295.405){\fontsize{11.991}{1}\usefont{T1}{cmr}{m}{n}\selectfont\color{color_29791}s}
\put(84.213,-323.496){\fontsize{11.991}{1}\usefont{T1}{cmr}{m}{n}\selectfont\color{color_29791}•}
\put(98.187,-324.205){\fontsize{11.991}{1}\usefont{T1}{cmr}{m}{n}\selectfont\color{color_29791}C}
\put(106.7845,-324.205){\fontsize{11.991}{1}\usefont{T1}{cmr}{m}{n}\selectfont\color{color_29791}h}
\put(113.5715,-324.205){\fontsize{11.991}{1}\usefont{T1}{cmr}{m}{n}\selectfont\color{color_29791}e}
\put(120.2504,-324.205){\fontsize{11.991}{1}\usefont{T1}{cmr}{m}{n}\selectfont\color{color_29791}ck }
\put(135.539,-324.205){\fontsize{11.991}{1}\usefont{T1}{cmr}{m}{n}\selectfont\color{color_29791}yo}
\put(148.3214,-324.205){\fontsize{11.991}{1}\usefont{T1}{cmr}{m}{n}\selectfont\color{color_29791}u}
\put(155.0004,-324.205){\fontsize{11.991}{1}\usefont{T1}{cmr}{m}{n}\selectfont\color{color_29791}r}
\put(158.8974,-324.205){\fontsize{11.991}{1}\usefont{T1}{cmr}{m}{n}\selectfont\color{color_29791} }
\put(162.2789,-324.205){\fontsize{11.991}{1}\usefont{T1}{cmr}{m}{n}\selectfont\color{color_29791}q}
\put(168.9579,-324.205){\fontsize{11.991}{1}\usefont{T1}{cmr}{m}{n}\selectfont\color{color_29791}u}
\put(175.6369,-324.205){\fontsize{11.991}{1}\usefont{T1}{cmr}{m}{n}\selectfont\color{color_29791}o}
\put(182.3159,-324.205){\fontsize{11.991}{1}\usefont{T1}{cmr}{m}{n}\selectfont\color{color_29791}t}
\put(185.6973,-324.205){\fontsize{11.991}{1}\usefont{T1}{cmr}{m}{n}\selectfont\color{color_29791}a}
\put(48.213,-356.803){\fontsize{15.987}{1}\usefont{T1}{cmr}{m}{n}\selectfont\color{color_29791}•}
\put(62.187,-357.795){\fontsize{15.987}{1}\usefont{T1}{cmr}{m}{n}\selectfont\color{color_29791}M}
\put(75.47219,-357.795){\fontsize{15.987}{1}\usefont{T1}{cmr}{m}{n}\selectfont\color{color_29791}or}
\put(89.65266,-357.795){\fontsize{15.987}{1}\usefont{T1}{cmr}{m}{n}\selectfont\color{color_29791}e }
\put(102.9379,-357.795){\fontsize{15.987}{1}\usefont{T1}{cmr}{m}{n}\selectfont\color{color_29791}i}
\put(106.5189,-357.795){\fontsize{15.987}{1}\usefont{T1}{cmr}{m}{n}\selectfont\color{color_29791}nf}
\put(119.8041,-357.795){\fontsize{15.987}{1}\usefont{T1}{cmr}{m}{n}\selectfont\color{color_29791}or}
\put(133.9846,-357.795){\fontsize{15.987}{1}\usefont{T1}{cmr}{m}{n}\selectfont\color{color_29791}m}
\put(147.2698,-357.795){\fontsize{15.987}{1}\usefont{T1}{cmr}{m}{n}\selectfont\color{color_29791}at}
\put(160.555,-357.795){\fontsize{15.987}{1}\usefont{T1}{cmr}{m}{n}\selectfont\color{color_29791}i}
\put(164.1361,-357.795){\fontsize{15.987}{1}\usefont{T1}{cmr}{m}{n}\selectfont\color{color_29791}on:}
\put(186.3101,-357.795){\fontsize{15.987}{1}\usefont{T1}{cmr}{m}{n}\selectfont\color{color_29791} }
\put(190.909,-357.795){\fontsize{15.987}{1}\usefont{T1}{cmr}{m}{n}\selectfont\color{color_167934}h}
\put(199.2063,-357.795){\fontsize{15.987}{1}\usefont{T1}{cmr}{m}{n}\selectfont\color{color_167934}tt}
\put(209.2941,-357.795){\fontsize{15.987}{1}\usefont{T1}{cmr}{m}{n}\selectfont\color{color_167934}p:}
\put(221.9717,-357.795){\fontsize{15.987}{1}\usefont{T1}{cmr}{m}{n}\selectfont\color{color_167934}/}
\put(228.0628,-357.795){\fontsize{15.987}{1}\usefont{T1}{cmr}{m}{n}\selectfont\color{color_167934}/}
\put(234.2498,-357.795){\fontsize{15.987}{1}\usefont{T1}{cmr}{m}{n}\selectfont\color{color_167934}i}
\put(237.9268,-357.795){\fontsize{15.987}{1}\usefont{T1}{cmr}{m}{n}\selectfont\color{color_167934}c}
\put(244.6253,-357.795){\fontsize{15.987}{1}\usefont{T1}{cmr}{m}{n}\selectfont\color{color_167934}t}
\put(249.8211,-357.795){\fontsize{15.987}{1}\usefont{T1}{cmr}{m}{n}\selectfont\color{color_167934}do}
\put(266.6075,-357.795){\fontsize{15.987}{1}\usefont{T1}{cmr}{m}{n}\selectfont\color{color_167934}c}
\put(273.306,-357.795){\fontsize{15.987}{1}\usefont{T1}{cmr}{m}{n}\selectfont\color{color_167934}s}
\put(279.493,-357.795){\fontsize{15.987}{1}\usefont{T1}{cmr}{m}{n}\selectfont\color{color_167934}/}
\put(285.584,-357.795){\fontsize{15.987}{1}\usefont{T1}{cmr}{m}{n}\selectfont\color{color_167934}w}
\put(297.0627,-357.795){\fontsize{15.987}{1}\usefont{T1}{cmr}{m}{n}\selectfont\color{color_167934}i}
\put(300.7397,-357.795){\fontsize{15.987}{1}\usefont{T1}{cmr}{m}{n}\selectfont\color{color_167934}k}
\put(307.9338,-357.795){\fontsize{15.987}{1}\usefont{T1}{cmr}{m}{n}\selectfont\color{color_167934}i}
\put(311.6108,-357.795){\fontsize{15.987}{1}\usefont{T1}{cmr}{m}{n}\selectfont\color{color_167934}/}
\put(317.7019,-357.795){\fontsize{15.987}{1}\usefont{T1}{cmr}{m}{n}\selectfont\color{color_167934}i}
\put(321.3789,-357.795){\fontsize{15.987}{1}\usefont{T1}{cmr}{m}{n}\selectfont\color{color_167934}nde}
\put(345.855,-357.795){\fontsize{15.987}{1}\usefont{T1}{cmr}{m}{n}\selectfont\color{color_167934}x}
\put(352.8413,-357.795){\fontsize{15.987}{1}\usefont{T1}{cmr}{m}{n}\selectfont\color{color_167934}.}
\put(356.8381,-357.795){\fontsize{15.987}{1}\usefont{T1}{cmr}{m}{n}\selectfont\color{color_167934}p}
\put(365.3112,-357.795){\fontsize{15.987}{1}\usefont{T1}{cmr}{m}{n}\selectfont\color{color_167934}h}
\put(373.6084,-357.795){\fontsize{15.987}{1}\usefont{T1}{cmr}{m}{n}\selectfont\color{color_167934}p?}
\put(389.4835,-357.795){\fontsize{15.987}{1}\usefont{T1}{cmr}{m}{n}\selectfont\color{color_167934}ti}
\put(398.3723,-357.795){\fontsize{15.987}{1}\usefont{T1}{cmr}{m}{n}\selectfont\color{color_167934}t}
\put(403.664,-357.795){\fontsize{15.987}{1}\usefont{T1}{cmr}{m}{n}\selectfont\color{color_167934}l}
\put(407.4529,-357.795){\fontsize{15.987}{1}\usefont{T1}{cmr}{m}{n}\selectfont\color{color_167934}e}
\put(415.3505,-357.795){\fontsize{15.987}{1}\usefont{T1}{cmr}{m}{n}\selectfont\color{color_167934}=}
\put(423.248,-357.795){\fontsize{15.987}{1}\usefont{T1}{cmr}{m}{n}\selectfont\color{color_167934}H}
\put(433.2239,-357.795){\fontsize{15.987}{1}\usefont{T1}{cmr}{m}{n}\selectfont\color{color_167934}P}
\put(441.5052,-357.795){\fontsize{15.987}{1}\usefont{T1}{cmr}{m}{n}\selectfont\color{color_167934}C}
\put(449.9943,-357.795){\fontsize{15.987}{1}\usefont{T1}{cmr}{m}{n}\selectfont\color{color_167934}:}
\put(454.2788,-357.795){\fontsize{15.987}{1}\usefont{T1}{cmr}{m}{n}\selectfont\color{color_167934}U}
\put(464.4626,-357.795){\fontsize{15.987}{1}\usefont{T1}{cmr}{m}{n}\selectfont\color{color_167934}s}
\put(470.6495,-357.795){\fontsize{15.987}{1}\usefont{T1}{cmr}{m}{n}\selectfont\color{color_167934}e}
\put(478.627,-357.795){\fontsize{15.987}{1}\usefont{T1}{cmr}{m}{n}\selectfont\color{color_167934}r}
\put(484.1266,-357.795){\fontsize{15.987}{1}\usefont{T1}{cmr}{m}{n}\selectfont\color{color_167934}\_}
\put(492.0242,-357.795){\fontsize{15.987}{1}\usefont{T1}{cmr}{m}{n}\selectfont\color{color_167934}Gui}
\put(514.1661,-357.795){\fontsize{15.987}{1}\usefont{T1}{cmr}{m}{n}\selectfont\color{color_167934}de}
\end{picture}
\begin{tikzpicture}[overlay]
\path(0pt,0pt);
\draw[color_167934,line width=0.992pt]
(190.909pt, -359.382pt) -- (531.01pt, -359.382pt)
;
\end{tikzpicture}
\begin{picture}(-5,0)(2.5,0)
\put(48.213,-391.187){\fontsize{15.987}{1}\usefont{T1}{cmr}{m}{n}\selectfont\color{color_29791}•}
\put(62.187,-392.094){\fontsize{15.987}{1}\usefont{T1}{cmr}{m}{n}\selectfont\color{color_29791}D}
\put(73.68165,-392.094){\fontsize{15.987}{1}\usefont{T1}{cmr}{m}{n}\selectfont\color{color_29791}em}
\put(95.85562,-392.094){\fontsize{15.987}{1}\usefont{T1}{cmr}{m}{n}\selectfont\color{color_29791}onst}
\put(126.0231,-392.094){\fontsize{15.987}{1}\usefont{T1}{cmr}{m}{n}\selectfont\color{color_29791}r}
\put(131.3148,-392.094){\fontsize{15.987}{1}\usefont{T1}{cmr}{m}{n}\selectfont\color{color_29791}at}
\put(144.6799,-392.094){\fontsize{15.987}{1}\usefont{T1}{cmr}{m}{n}\selectfont\color{color_29791}i}
\put(148.261,-392.094){\fontsize{15.987}{1}\usefont{T1}{cmr}{m}{n}\selectfont\color{color_29791}o}
\put(157.0538,-392.094){\fontsize{15.987}{1}\usefont{T1}{cmr}{m}{n}\selectfont\color{color_29791}n}
\end{picture}
\newpage
\begin{tikzpicture}[overlay]
\path(0pt,0pt);
\filldraw[color_283006][even odd rule]
(-15pt, 10pt) -- (704.972pt, 10pt)
 -- (704.972pt, 10pt)
 -- (704.972pt, -529.972pt)
 -- (704.972pt, -529.972pt)
 -- (-15pt, -529.972pt)
 -- (-15pt, -529.972pt)
 -- (-15pt, 10pt) -- cycle
;
\end{tikzpicture}
\begin{picture}(-5,0)(2.5,0)
\put(-15,-529.972){\includegraphics[width=720pt,height=540pt]{latexImage_70f318a8f2de3f58eebb26eed4f2123e.png}}
\put(48.411,-60.69601){\fontsize{21.997}{1}\usefont{T1}{cmr}{b}{n}\selectfont\color{color_29791}M}
\put(66.7125,-60.69601){\fontsize{21.997}{1}\usefont{T1}{cmr}{b}{n}\selectfont\color{color_29791}o}
\put(80.10867,-60.69601){\fontsize{21.997}{1}\usefont{T1}{cmr}{b}{n}\selectfont\color{color_29791}d}
\put(93.59283,-60.69601){\fontsize{21.997}{1}\usefont{T1}{cmr}{b}{n}\selectfont\color{color_29791}e}
\put(105.8012,-60.69601){\fontsize{21.997}{1}\usefont{T1}{cmr}{b}{n}\selectfont\color{color_29791}l E}
\put(132.6595,-60.69601){\fontsize{21.997}{1}\usefont{T1}{cmr}{b}{n}\selectfont\color{color_29791}n}
\put(146.0557,-60.69601){\fontsize{21.997}{1}\usefont{T1}{cmr}{b}{n}\selectfont\color{color_29791}s}
\put(158.33,-60.69601){\fontsize{21.997}{1}\usefont{T1}{cmr}{b}{n}\selectfont\color{color_29791}e}
\put(170.5383,-60.69601){\fontsize{21.997}{1}\usefont{T1}{cmr}{b}{n}\selectfont\color{color_29791}m}
\put(190.1157,-60.69601){\fontsize{21.997}{1}\usefont{T1}{cmr}{b}{n}\selectfont\color{color_29791}b}
\put(203.5119,-60.69601){\fontsize{21.997}{1}\usefont{T1}{cmr}{b}{n}\selectfont\color{color_29791}le}
\put(221.8134,-60.69601){\fontsize{21.997}{1}\usefont{T1}{cmr}{b}{n}\selectfont\color{color_29791}r}
\put(48.213,-115.49){\fontsize{15.987}{1}\usefont{T1}{cmr}{m}{n}\selectfont\color{color_29791}•}
\put(62.187,-116.51){\fontsize{15.987}{1}\usefont{T1}{cmr}{m}{n}\selectfont\color{color_29791}A}
\put(72.86631,-116.51){\fontsize{15.987}{1}\usefont{T1}{cmr}{m}{n}\selectfont\color{color_29791}i}
\put(76.4474,-116.51){\fontsize{15.987}{1}\usefont{T1}{cmr}{m}{n}\selectfont\color{color_29791}m}
\put(89.73259,-116.51){\fontsize{15.987}{1}\usefont{T1}{cmr}{m}{n}\selectfont\color{color_29791}ed }
\put(111.9066,-116.51){\fontsize{15.987}{1}\usefont{T1}{cmr}{m}{n}\selectfont\color{color_29791}at}
\put(125.1918,-116.51){\fontsize{15.987}{1}\usefont{T1}{cmr}{m}{n}\selectfont\color{color_29791} }
\put(129.6681,-116.51){\fontsize{15.987}{1}\usefont{T1}{cmr}{m}{n}\selectfont\color{color_29791}sol}
\put(150.1315,-116.51){\fontsize{15.987}{1}\usefont{T1}{cmr}{m}{n}\selectfont\color{color_29791}vi}
\put(161.7061,-116.51){\fontsize{15.987}{1}\usefont{T1}{cmr}{m}{n}\selectfont\color{color_29791}ng }
\put(183.88,-116.51){\fontsize{15.987}{1}\usefont{T1}{cmr}{m}{n}\selectfont\color{color_29791}c}
\put(191.9535,-116.51){\fontsize{15.987}{1}\usefont{T1}{cmr}{m}{n}\selectfont\color{color_29791}om}
\put(214.1274,-116.51){\fontsize{15.987}{1}\usefont{T1}{cmr}{m}{n}\selectfont\color{color_29791}pl}
\put(226.5973,-116.51){\fontsize{15.987}{1}\usefont{T1}{cmr}{m}{n}\selectfont\color{color_29791}e}
\put(235.3902,-116.51){\fontsize{15.987}{1}\usefont{T1}{cmr}{m}{n}\selectfont\color{color_29791}x}
\put(243.4636,-116.51){\fontsize{15.987}{1}\usefont{T1}{cmr}{m}{n}\selectfont\color{color_29791} }
\put(247.86,-116.51){\fontsize{15.987}{1}\usefont{T1}{cmr}{m}{n}\selectfont\color{color_29791}H}
\put(259.3547,-116.51){\fontsize{15.987}{1}\usefont{T1}{cmr}{m}{n}\selectfont\color{color_29791}P}
\put(270.034,-116.51){\fontsize{15.987}{1}\usefont{T1}{cmr}{m}{n}\selectfont\color{color_29791}C}
\put(281.5287,-116.51){\fontsize{15.987}{1}\usefont{T1}{cmr}{m}{n}\selectfont\color{color_29791} }
\put(286.005,-116.51){\fontsize{15.987}{1}\usefont{T1}{cmr}{m}{n}\selectfont\color{color_29791}w}
\put(297.4997,-116.51){\fontsize{15.987}{1}\usefont{T1}{cmr}{m}{n}\selectfont\color{color_29791}or}
\put(311.6801,-116.51){\fontsize{15.987}{1}\usefont{T1}{cmr}{m}{n}\selectfont\color{color_29791}kl}
\put(323.2547,-116.51){\fontsize{15.987}{1}\usefont{T1}{cmr}{m}{n}\selectfont\color{color_29791}oad}
\put(349.8251,-116.51){\fontsize{15.987}{1}\usefont{T1}{cmr}{m}{n}\selectfont\color{color_29791}s}
\put(357.8986,-116.51){\fontsize{15.987}{1}\usefont{T1}{cmr}{m}{n}\selectfont\color{color_29791}.}
\end{picture}
\newpage
\begin{tikzpicture}[overlay]
\path(0pt,0pt);
\filldraw[color_283006][even odd rule]
(-15pt, 10pt) -- (704.972pt, 10pt)
 -- (704.972pt, 10pt)
 -- (704.972pt, -529.972pt)
 -- (704.972pt, -529.972pt)
 -- (-15pt, -529.972pt)
 -- (-15pt, -529.972pt)
 -- (-15pt, 10pt) -- cycle
;
\end{tikzpicture}
\begin{picture}(-5,0)(2.5,0)
\put(-15,-529.972){\includegraphics[width=720pt,height=540pt]{latexImage_2828171a63492229494bbe96f89381b1.png}}
\put(48.411,-60.69601){\fontsize{21.997}{1}\usefont{T1}{cmr}{b}{n}\selectfont\color{color_29791}De}
\put(76.50117,-60.69601){\fontsize{21.997}{1}\usefont{T1}{cmr}{b}{n}\selectfont\color{color_29791}v}
\put(88.7095,-60.69601){\fontsize{21.997}{1}\usefont{T1}{cmr}{b}{n}\selectfont\color{color_29791}e}
\put(100.9838,-60.69601){\fontsize{21.997}{1}\usefont{T1}{cmr}{b}{n}\selectfont\color{color_29791}lo}
\put(120.5612,-60.69601){\fontsize{21.997}{1}\usefont{T1}{cmr}{b}{n}\selectfont\color{color_29791}p}
\put(133.9573,-60.69601){\fontsize{21.997}{1}\usefont{T1}{cmr}{b}{n}\selectfont\color{color_29791}e}
\put(146.1657,-60.69601){\fontsize{21.997}{1}\usefont{T1}{cmr}{b}{n}\selectfont\color{color_29791}r}
\put(154.7445,-60.69601){\fontsize{21.997}{1}\usefont{T1}{cmr}{b}{n}\selectfont\color{color_29791} E}
\put(175.4437,-60.69601){\fontsize{21.997}{1}\usefont{T1}{cmr}{b}{n}\selectfont\color{color_29791}n}
\put(188.9278,-60.69601){\fontsize{21.997}{1}\usefont{T1}{cmr}{b}{n}\selectfont\color{color_29791}v}
\put(201.1362,-60.69601){\fontsize{21.997}{1}\usefont{T1}{cmr}{b}{n}\selectfont\color{color_29791}ir}
\put(215.8082,-60.69601){\fontsize{21.997}{1}\usefont{T1}{cmr}{b}{n}\selectfont\color{color_29791}o}
\put(229.2043,-60.69601){\fontsize{21.997}{1}\usefont{T1}{cmr}{b}{n}\selectfont\color{color_29791}n}
\put(242.6005,-60.69601){\fontsize{21.997}{1}\usefont{T1}{cmr}{b}{n}\selectfont\color{color_29791}m}
\put(262.1778,-60.69601){\fontsize{21.997}{1}\usefont{T1}{cmr}{b}{n}\selectfont\color{color_29791}e}
\put(274.3862,-60.69601){\fontsize{21.997}{1}\usefont{T1}{cmr}{b}{n}\selectfont\color{color_29791}n}
\put(287.8703,-60.69601){\fontsize{21.997}{1}\usefont{T1}{cmr}{b}{n}\selectfont\color{color_29791}t}
\put(48.213,-115.49){\fontsize{15.987}{1}\usefont{T1}{cmr}{m}{n}\selectfont\color{color_29791}•}
\put(62.187,-116.51){\fontsize{15.987}{1}\usefont{T1}{cmr}{m}{n}\selectfont\color{color_29791}U}
\put(73.68165,-116.51){\fontsize{15.987}{1}\usefont{T1}{cmr}{m}{n}\selectfont\color{color_29791}s}
\put(81.75508,-116.51){\fontsize{15.987}{1}\usefont{T1}{cmr}{m}{n}\selectfont\color{color_29791}e}
\put(90.54793,-116.51){\fontsize{15.987}{1}\usefont{T1}{cmr}{m}{n}\selectfont\color{color_29791} }
\put(95.02428,-116.51){\fontsize{15.987}{1}\usefont{T1}{cmr}{m}{n}\selectfont\color{color_29791}gi}
\put(107.4941,-116.51){\fontsize{15.987}{1}\usefont{T1}{cmr}{m}{n}\selectfont\color{color_29791}t}
\put(111.8906,-116.51){\fontsize{15.987}{1}\usefont{T1}{cmr}{m}{n}\selectfont\color{color_29791}!}
\put(48.213,-149.902){\fontsize{15.987}{1}\usefont{T1}{cmr}{m}{n}\selectfont\color{color_29791}•}
\put(62.187,-150.809){\fontsize{15.987}{1}\usefont{T1}{cmr}{m}{n}\selectfont\color{color_29791}R}
\put(73.68165,-150.809){\fontsize{15.987}{1}\usefont{T1}{cmr}{m}{n}\selectfont\color{color_29791}epea}
\put(109.1408,-150.809){\fontsize{15.987}{1}\usefont{T1}{cmr}{m}{n}\selectfont\color{color_29791}t}
\put(113.6172,-150.809){\fontsize{15.987}{1}\usefont{T1}{cmr}{m}{n}\selectfont\color{color_29791}abl}
\put(134.9758,-150.809){\fontsize{15.987}{1}\usefont{T1}{cmr}{m}{n}\selectfont\color{color_29791}e}
\put(143.7686,-150.809){\fontsize{15.987}{1}\usefont{T1}{cmr}{m}{n}\selectfont\color{color_29791},}
\put(148.245,-150.809){\fontsize{15.987}{1}\usefont{T1}{cmr}{m}{n}\selectfont\color{color_29791} }
\put(152.6414,-150.809){\fontsize{15.987}{1}\usefont{T1}{cmr}{m}{n}\selectfont\color{color_29791}r}
\put(157.9331,-150.809){\fontsize{15.987}{1}\usefont{T1}{cmr}{m}{n}\selectfont\color{color_29791}epr}
\put(180.9064,-150.809){\fontsize{15.987}{1}\usefont{T1}{cmr}{m}{n}\selectfont\color{color_29791}oduci}
\put(219.1474,-150.809){\fontsize{15.987}{1}\usefont{T1}{cmr}{m}{n}\selectfont\color{color_29791}bl}
\put(231.6172,-150.809){\fontsize{15.987}{1}\usefont{T1}{cmr}{m}{n}\selectfont\color{color_29791}e }
\put(244.9024,-150.809){\fontsize{15.987}{1}\usefont{T1}{cmr}{m}{n}\selectfont\color{color_29791}\&}
\put(255.5817,-150.809){\fontsize{15.987}{1}\usefont{T1}{cmr}{m}{n}\selectfont\color{color_29791} }
\put(259.9781,-150.809){\fontsize{15.987}{1}\usefont{T1}{cmr}{m}{n}\selectfont\color{color_29791}s}
\put(268.0516,-150.809){\fontsize{15.987}{1}\usefont{T1}{cmr}{m}{n}\selectfont\color{color_29791}ha}
\put(285.7332,-150.809){\fontsize{15.987}{1}\usefont{T1}{cmr}{m}{n}\selectfont\color{color_29791}r}
\put(291.0249,-150.809){\fontsize{15.987}{1}\usefont{T1}{cmr}{m}{n}\selectfont\color{color_29791}eab}
\put(317.5953,-150.809){\fontsize{15.987}{1}\usefont{T1}{cmr}{m}{n}\selectfont\color{color_29791}l}
\put(321.1764,-150.809){\fontsize{15.987}{1}\usefont{T1}{cmr}{m}{n}\selectfont\color{color_29791}e }
\put(334.5415,-150.809){\fontsize{15.987}{1}\usefont{T1}{cmr}{m}{n}\selectfont\color{color_29791}-}
\put(339.8332,-150.809){\fontsize{15.987}{1}\usefont{T1}{cmr}{m}{n}\selectfont\color{color_29791} }
\put(344.2296,-150.809){\fontsize{15.987}{1}\usefont{T1}{cmr}{m}{n}\selectfont\color{color_29791}cont}
\put(374.3971,-150.809){\fontsize{15.987}{1}\usefont{T1}{cmr}{m}{n}\selectfont\color{color_29791}ai}
\put(386.8669,-150.809){\fontsize{15.987}{1}\usefont{T1}{cmr}{m}{n}\selectfont\color{color_29791}ner}
\put(409.9362,-150.809){\fontsize{15.987}{1}\usefont{T1}{cmr}{m}{n}\selectfont\color{color_29791}s }
\put(422.3261,-150.809){\fontsize{15.987}{1}\usefont{T1}{cmr}{m}{n}\selectfont\color{color_29791}c}
\put(430.3995,-150.809){\fontsize{15.987}{1}\usefont{T1}{cmr}{m}{n}\selectfont\color{color_29791}a}
\put(439.1924,-150.809){\fontsize{15.987}{1}\usefont{T1}{cmr}{m}{n}\selectfont\color{color_29791}n }
\put(452.5575,-150.809){\fontsize{15.987}{1}\usefont{T1}{cmr}{m}{n}\selectfont\color{color_29791}h}
\put(461.3504,-150.809){\fontsize{15.987}{1}\usefont{T1}{cmr}{m}{n}\selectfont\color{color_29791}el}
\put(473.8202,-150.809){\fontsize{15.987}{1}\usefont{T1}{cmr}{m}{n}\selectfont\color{color_29791}p}
\end{picture}
\newpage
\begin{tikzpicture}[overlay]
\path(0pt,0pt);
\filldraw[color_283006][even odd rule]
(-15pt, 10pt) -- (704.972pt, 10pt)
 -- (704.972pt, 10pt)
 -- (704.972pt, -529.972pt)
 -- (704.972pt, -529.972pt)
 -- (-15pt, -529.972pt)
 -- (-15pt, -529.972pt)
 -- (-15pt, 10pt) -- cycle
;
\end{tikzpicture}
\begin{picture}(-5,0)(2.5,0)
\put(-15,-529.972){\includegraphics[width=720pt,height=540pt]{latexImage_6f17c9ae575e6e3412bad9150a17b852.png}}
\put(48.411,-60.69601){\fontsize{21.997}{1}\usefont{T1}{cmr}{b}{n}\selectfont\color{color_29791}Be}
\put(76.50117,-60.69601){\fontsize{21.997}{1}\usefont{T1}{cmr}{b}{n}\selectfont\color{color_29791}s}
\put(88.7095,-60.69601){\fontsize{21.997}{1}\usefont{T1}{cmr}{b}{n}\selectfont\color{color_29791}t}
\put(96.0785,-60.69601){\fontsize{21.997}{1}\usefont{T1}{cmr}{b}{n}\selectfont\color{color_29791} P}
\put(116.7777,-60.69601){\fontsize{21.997}{1}\usefont{T1}{cmr}{b}{n}\selectfont\color{color_29791}r}
\put(125.3565,-60.69601){\fontsize{21.997}{1}\usefont{T1}{cmr}{b}{n}\selectfont\color{color_29791}a}
\put(137.5648,-60.69601){\fontsize{21.997}{1}\usefont{T1}{cmr}{b}{n}\selectfont\color{color_29791}c}
\put(149.7732,-60.69601){\fontsize{21.997}{1}\usefont{T1}{cmr}{b}{n}\selectfont\color{color_29791}t}
\put(157.0762,-60.69601){\fontsize{21.997}{1}\usefont{T1}{cmr}{b}{n}\selectfont\color{color_29791}i}
\put(163.2573,-60.69601){\fontsize{21.997}{1}\usefont{T1}{cmr}{b}{n}\selectfont\color{color_29791}c}
\put(175.4657,-60.69601){\fontsize{21.997}{1}\usefont{T1}{cmr}{b}{n}\selectfont\color{color_29791}e}
\put(48.213,-115.49){\fontsize{15.987}{1}\usefont{T1}{cmr}{m}{n}\selectfont\color{color_29791}•}
\put(62.187,-116.51){\fontsize{15.987}{1}\usefont{T1}{cmr}{m}{n}\selectfont\color{color_29791}U}
\put(73.68165,-116.51){\fontsize{15.987}{1}\usefont{T1}{cmr}{m}{n}\selectfont\color{color_29791}s}
\put(81.75508,-116.51){\fontsize{15.987}{1}\usefont{T1}{cmr}{m}{n}\selectfont\color{color_29791}e}
\put(90.54793,-116.51){\fontsize{15.987}{1}\usefont{T1}{cmr}{m}{n}\selectfont\color{color_29791}r}
\put(95.83962,-116.51){\fontsize{15.987}{1}\usefont{T1}{cmr}{m}{n}\selectfont\color{color_29791} }
\put(100.316,-116.51){\fontsize{15.987}{1}\usefont{T1}{cmr}{m}{n}\selectfont\color{color_29791}P}
\put(110.9953,-116.51){\fontsize{15.987}{1}\usefont{T1}{cmr}{m}{n}\selectfont\color{color_29791}o}
\put(119.7881,-116.51){\fontsize{15.987}{1}\usefont{T1}{cmr}{m}{n}\selectfont\color{color_29791}l}
\put(123.3692,-116.51){\fontsize{15.987}{1}\usefont{T1}{cmr}{m}{n}\selectfont\color{color_29791}i}
\put(126.9503,-116.51){\fontsize{15.987}{1}\usefont{T1}{cmr}{m}{n}\selectfont\color{color_29791}c}
\put(135.0238,-116.51){\fontsize{15.987}{1}\usefont{T1}{cmr}{m}{n}\selectfont\color{color_29791}y}
\put(84.213,-145.31){\fontsize{11.991}{1}\usefont{T1}{cmr}{m}{n}\selectfont\color{color_29791}•}
\put(98.187,-145.991){\fontsize{11.991}{1}\usefont{T1}{cmr}{m}{n}\selectfont\color{color_29791}L}
\put(104.866,-145.991){\fontsize{11.991}{1}\usefont{T1}{cmr}{m}{n}\selectfont\color{color_29791}i}
\put(107.552,-145.991){\fontsize{11.991}{1}\usefont{T1}{cmr}{m}{n}\selectfont\color{color_29791}n}
\put(114.231,-145.991){\fontsize{11.991}{1}\usefont{T1}{cmr}{m}{n}\selectfont\color{color_29791}k:}
\put(123.524,-145.991){\fontsize{11.991}{1}\usefont{T1}{cmr}{m}{n}\selectfont\color{color_29791} }
\put(48.213,-178.589){\fontsize{15.987}{1}\usefont{T1}{cmr}{m}{n}\selectfont\color{color_29791}•}
\put(62.187,-179.609){\fontsize{15.987}{1}\usefont{T1}{cmr}{m}{n}\selectfont\color{color_29791}D}
\put(73.68165,-179.609){\fontsize{15.987}{1}\usefont{T1}{cmr}{m}{n}\selectfont\color{color_29791}o!}
\put(84.213,-208.494){\fontsize{11.991}{1}\usefont{T1}{cmr}{m}{n}\selectfont\color{color_29791}•}
\put(98.187,-209.09){\fontsize{11.991}{1}\usefont{T1}{cmr}{m}{n}\selectfont\color{color_29791}A}
\put(106.2689,-209.09){\fontsize{11.991}{1}\usefont{T1}{cmr}{m}{n}\selectfont\color{color_29791}s}
\put(112.1685,-209.09){\fontsize{11.991}{1}\usefont{T1}{cmr}{m}{n}\selectfont\color{color_29791}k}
\put(118.2479,-209.09){\fontsize{11.991}{1}\usefont{T1}{cmr}{m}{n}\selectfont\color{color_29791} }
\put(121.5455,-209.09){\fontsize{11.991}{1}\usefont{T1}{cmr}{m}{n}\selectfont\color{color_29791}f}
\put(124.9269,-209.09){\fontsize{11.991}{1}\usefont{T1}{cmr}{m}{n}\selectfont\color{color_29791}o}
\put(131.6059,-209.09){\fontsize{11.991}{1}\usefont{T1}{cmr}{m}{n}\selectfont\color{color_29791}r }
\put(138.8965,-209.09){\fontsize{11.991}{1}\usefont{T1}{cmr}{m}{n}\selectfont\color{color_29791}h}
\put(145.5754,-209.09){\fontsize{11.991}{1}\usefont{T1}{cmr}{m}{n}\selectfont\color{color_29791}e}
\put(152.2544,-209.09){\fontsize{11.991}{1}\usefont{T1}{cmr}{m}{n}\selectfont\color{color_29791}l}
\put(154.9404,-209.09){\fontsize{11.991}{1}\usefont{T1}{cmr}{m}{n}\selectfont\color{color_29791}p}
\put(161.6194,-209.09){\fontsize{11.991}{1}\usefont{T1}{cmr}{m}{n}\selectfont\color{color_29791},}
\put(165.0009,-209.09){\fontsize{11.991}{1}\usefont{T1}{cmr}{m}{n}\selectfont\color{color_29791} }
\put(48.213,-270.488){\fontsize{15.987}{1}\usefont{T1}{cmr}{m}{n}\selectfont\color{color_29791}•}
\put(62.187,-271.509){\fontsize{15.987}{1}\usefont{T1}{cmr}{m}{n}\selectfont\color{color_29791}D}
\put(73.68165,-271.509){\fontsize{15.987}{1}\usefont{T1}{cmr}{m}{n}\selectfont\color{color_29791}on'}
\put(94.44876,-271.509){\fontsize{15.987}{1}\usefont{T1}{cmr}{m}{n}\selectfont\color{color_29791}t}
\put(98.84518,-271.509){\fontsize{15.987}{1}\usefont{T1}{cmr}{m}{n}\selectfont\color{color_29791}!}
\put(84.213,-300.394){\fontsize{11.991}{1}\usefont{T1}{cmr}{m}{n}\selectfont\color{color_29791}•}
\put(98.187,-300.989){\fontsize{11.991}{1}\usefont{T1}{cmr}{m}{n}\selectfont\color{color_29791}S}
\put(106.2689,-300.989){\fontsize{11.991}{1}\usefont{T1}{cmr}{m}{n}\selectfont\color{color_29791}u}
\put(112.9479,-300.989){\fontsize{11.991}{1}\usefont{T1}{cmr}{m}{n}\selectfont\color{color_29791}b}
\put(119.6269,-300.989){\fontsize{11.991}{1}\usefont{T1}{cmr}{m}{n}\selectfont\color{color_29791}mi}
\put(132.2175,-300.989){\fontsize{11.991}{1}\usefont{T1}{cmr}{m}{n}\selectfont\color{color_29791}t}
\put(135.5989,-300.989){\fontsize{11.991}{1}\usefont{T1}{cmr}{m}{n}\selectfont\color{color_29791} }
\put(138.8964,-300.989){\fontsize{11.991}{1}\usefont{T1}{cmr}{m}{n}\selectfont\color{color_29791}j}
\put(141.4985,-300.989){\fontsize{11.991}{1}\usefont{T1}{cmr}{m}{n}\selectfont\color{color_29791}o}
\put(148.2854,-300.989){\fontsize{11.991}{1}\usefont{T1}{cmr}{m}{n}\selectfont\color{color_29791}b}
\put(154.9644,-300.989){\fontsize{11.991}{1}\usefont{T1}{cmr}{m}{n}\selectfont\color{color_29791}s }
\put(164.3414,-300.989){\fontsize{11.991}{1}\usefont{T1}{cmr}{m}{n}\selectfont\color{color_29791}w}
\put(172.9389,-300.989){\fontsize{11.991}{1}\usefont{T1}{cmr}{m}{n}\selectfont\color{color_29791}h}
\put(179.6179,-300.989){\fontsize{11.991}{1}\usefont{T1}{cmr}{m}{n}\selectfont\color{color_29791}i}
\put(182.3039,-300.989){\fontsize{11.991}{1}\usefont{T1}{cmr}{m}{n}\selectfont\color{color_29791}c}
\put(188.2034,-300.989){\fontsize{11.991}{1}\usefont{T1}{cmr}{m}{n}\selectfont\color{color_29791}h}
\put(194.9904,-300.989){\fontsize{11.991}{1}\usefont{T1}{cmr}{m}{n}\selectfont\color{color_29791} }
\put(198.2879,-300.989){\fontsize{11.991}{1}\usefont{T1}{cmr}{m}{n}\selectfont\color{color_29791}u}
\put(205.0748,-300.989){\fontsize{11.991}{1}\usefont{T1}{cmr}{m}{n}\selectfont\color{color_29791}se}
\put(217.7493,-300.989){\fontsize{11.991}{1}\usefont{T1}{cmr}{m}{n}\selectfont\color{color_29791} }
\put(221.0468,-300.989){\fontsize{11.991}{1}\usefont{T1}{cmr}{m}{n}\selectfont\color{color_29791}mo}
\put(237.7143,-300.989){\fontsize{11.991}{1}\usefont{T1}{cmr}{m}{n}\selectfont\color{color_29791}re}
\put(248.3863,-300.989){\fontsize{11.991}{1}\usefont{T1}{cmr}{m}{n}\selectfont\color{color_29791} }
\put(251.7678,-300.989){\fontsize{11.991}{1}\usefont{T1}{cmr}{m}{n}\selectfont\color{color_29791}t}
\put(255.0653,-300.989){\fontsize{11.991}{1}\usefont{T1}{cmr}{m}{n}\selectfont\color{color_29791}h}
\put(261.7443,-300.989){\fontsize{11.991}{1}\usefont{T1}{cmr}{m}{n}\selectfont\color{color_29791}a}
\put(268.5312,-300.989){\fontsize{11.991}{1}\usefont{T1}{cmr}{m}{n}\selectfont\color{color_29791}n}
\put(275.2101,-300.989){\fontsize{11.991}{1}\usefont{T1}{cmr}{m}{n}\selectfont\color{color_29791} }
\put(278.5077,-300.989){\fontsize{11.991}{1}\usefont{T1}{cmr}{m}{n}\selectfont\color{color_29791}4}
\put(285.1866,-300.989){\fontsize{11.991}{1}\usefont{T1}{cmr}{m}{n}\selectfont\color{color_29791} }
\put(288.5681,-300.989){\fontsize{11.991}{1}\usefont{T1}{cmr}{m}{n}\selectfont\color{color_29791}n}
\put(295.2471,-300.989){\fontsize{11.991}{1}\usefont{T1}{cmr}{m}{n}\selectfont\color{color_29791}o}
\put(302.034,-300.989){\fontsize{11.991}{1}\usefont{T1}{cmr}{m}{n}\selectfont\color{color_29791}d}
\put(308.713,-300.989){\fontsize{11.991}{1}\usefont{T1}{cmr}{m}{n}\selectfont\color{color_29791}e}
\put(315.392,-300.989){\fontsize{11.991}{1}\usefont{T1}{cmr}{m}{n}\selectfont\color{color_29791}s }
\put(324.685,-300.989){\fontsize{11.991}{1}\usefont{T1}{cmr}{m}{n}\selectfont\color{color_29791}a}
\put(331.4719,-300.989){\fontsize{11.991}{1}\usefont{T1}{cmr}{m}{n}\selectfont\color{color_29791}t}
\put(334.7694,-300.989){\fontsize{11.991}{1}\usefont{T1}{cmr}{m}{n}\selectfont\color{color_29791} }
\put(338.1509,-300.989){\fontsize{11.991}{1}\usefont{T1}{cmr}{m}{n}\selectfont\color{color_29791}a}
\put(344.8299,-300.989){\fontsize{11.991}{1}\usefont{T1}{cmr}{m}{n}\selectfont\color{color_29791} }
\put(348.2113,-300.989){\fontsize{11.991}{1}\usefont{T1}{cmr}{m}{n}\selectfont\color{color_29791}t}
\put(351.5089,-300.989){\fontsize{11.991}{1}\usefont{T1}{cmr}{m}{n}\selectfont\color{color_29791}i}
\put(354.1948,-300.989){\fontsize{11.991}{1}\usefont{T1}{cmr}{m}{n}\selectfont\color{color_29791}m}
\put(364.0874,-300.989){\fontsize{11.991}{1}\usefont{T1}{cmr}{m}{n}\selectfont\color{color_29791}e}
\put(370.7664,-300.989){\fontsize{11.991}{1}\usefont{T1}{cmr}{m}{n}\selectfont\color{color_29791}.}
\end{picture}
\newpage
\begin{tikzpicture}[overlay]
\path(0pt,0pt);
\filldraw[color_283006][even odd rule]
(-15pt, 10pt) -- (704.972pt, 10pt)
 -- (704.972pt, 10pt)
 -- (704.972pt, -529.972pt)
 -- (704.972pt, -529.972pt)
 -- (-15pt, -529.972pt)
 -- (-15pt, -529.972pt)
 -- (-15pt, 10pt) -- cycle
;
\end{tikzpicture}
\begin{picture}(-5,0)(2.5,0)
\put(-15,-529.972){\includegraphics[width=720pt,height=540pt]{latexImage_bdabda61317014bb20dc1cc743a20acb.png}}
\put(48.411,-60.69601){\fontsize{21.997}{1}\usefont{T1}{cmr}{b}{n}\selectfont\color{color_29791}HE}
\put(78.89884,-60.69601){\fontsize{21.997}{1}\usefont{T1}{cmr}{b}{n}\selectfont\color{color_29791}L}
\put(92.29501,-60.69601){\fontsize{21.997}{1}\usefont{T1}{cmr}{b}{n}\selectfont\color{color_29791}P}
\put(106.967,-60.69601){\fontsize{21.997}{1}\usefont{T1}{cmr}{b}{n}\selectfont\color{color_29791}!}
\put(48.213,-115.49){\fontsize{15.987}{1}\usefont{T1}{cmr}{m}{n}\selectfont\color{color_29791}•}
\put(62.187,-116.51){\fontsize{15.987}{1}\usefont{T1}{cmr}{m}{n}\selectfont\color{color_29791}S}
\put(72.86631,-116.51){\fontsize{15.987}{1}\usefont{T1}{cmr}{m}{n}\selectfont\color{color_29791}er}
\put(86.95086,-116.51){\fontsize{15.987}{1}\usefont{T1}{cmr}{m}{n}\selectfont\color{color_29791}v}
\put(95.02429,-116.51){\fontsize{15.987}{1}\usefont{T1}{cmr}{m}{n}\selectfont\color{color_29791}i}
\put(98.60538,-116.51){\fontsize{15.987}{1}\usefont{T1}{cmr}{m}{n}\selectfont\color{color_29791}ce }
\put(119.8841,-116.51){\fontsize{15.987}{1}\usefont{T1}{cmr}{m}{n}\selectfont\color{color_29791}des}
\put(145.7351,-116.51){\fontsize{15.987}{1}\usefont{T1}{cmr}{m}{n}\selectfont\color{color_29791}k:}
\put(158.125,-116.51){\fontsize{15.987}{1}\usefont{T1}{cmr}{m}{n}\selectfont\color{color_29791} }
\put(162.789,-116.51){\fontsize{15.987}{1}\usefont{T1}{cmr}{m}{n}\selectfont\color{color_29791}ser}
\put(184.963,-116.51){\fontsize{15.987}{1}\usefont{T1}{cmr}{m}{n}\selectfont\color{color_29791}vi}
\put(196.5376,-116.51){\fontsize{15.987}{1}\usefont{T1}{cmr}{m}{n}\selectfont\color{color_29791}c}
\put(204.611,-116.51){\fontsize{15.987}{1}\usefont{T1}{cmr}{m}{n}\selectfont\color{color_29791}ed}
\put(222.2926,-116.51){\fontsize{15.987}{1}\usefont{T1}{cmr}{m}{n}\selectfont\color{color_29791}es}
\put(239.2548,-116.51){\fontsize{15.987}{1}\usefont{T1}{cmr}{m}{n}\selectfont\color{color_29791}k@}
\put(263.4272,-116.51){\fontsize{15.987}{1}\usefont{T1}{cmr}{m}{n}\selectfont\color{color_29791}bas.}
\put(293.5946,-116.51){\fontsize{15.987}{1}\usefont{T1}{cmr}{m}{n}\selectfont\color{color_29791}ac}
\put(310.5569,-116.51){\fontsize{15.987}{1}\usefont{T1}{cmr}{m}{n}\selectfont\color{color_29791}.}
\put(314.9533,-116.51){\fontsize{15.987}{1}\usefont{T1}{cmr}{m}{n}\selectfont\color{color_29791}uk}
\end{picture}
\begin{tikzpicture}[overlay]
\path(0pt,0pt);
\draw[color_29791,line width=0.907pt]
(162.789pt, -117.899pt) -- (332.102pt, -117.899pt)
;
\end{tikzpicture}
\begin{picture}(-5,0)(2.5,0)
\put(48.213,-149.902){\fontsize{15.987}{1}\usefont{T1}{cmr}{m}{n}\selectfont\color{color_29791}•}
\put(62.187,-150.809){\fontsize{15.987}{1}\usefont{T1}{cmr}{m}{n}\selectfont\color{color_29791}H}
\put(73.68165,-150.809){\fontsize{15.987}{1}\usefont{T1}{cmr}{m}{n}\selectfont\color{color_29791}P}
\put(84.36096,-150.809){\fontsize{15.987}{1}\usefont{T1}{cmr}{m}{n}\selectfont\color{color_29791}C}
\put(95.85561,-150.809){\fontsize{15.987}{1}\usefont{T1}{cmr}{m}{n}\selectfont\color{color_29791} }
\put(100.332,-150.809){\fontsize{15.987}{1}\usefont{T1}{cmr}{m}{n}\selectfont\color{color_29791}U}
\put(111.8266,-150.809){\fontsize{15.987}{1}\usefont{T1}{cmr}{m}{n}\selectfont\color{color_29791}ser}
\put(134.0006,-150.809){\fontsize{15.987}{1}\usefont{T1}{cmr}{m}{n}\selectfont\color{color_29791} }
\put(138.397,-150.809){\fontsize{15.987}{1}\usefont{T1}{cmr}{m}{n}\selectfont\color{color_29791}G}
\put(150.7869,-150.809){\fontsize{15.987}{1}\usefont{T1}{cmr}{m}{n}\selectfont\color{color_29791}ui}
\put(163.2568,-150.809){\fontsize{15.987}{1}\usefont{T1}{cmr}{m}{n}\selectfont\color{color_29791}de}
\put(180.9384,-150.809){\fontsize{15.987}{1}\usefont{T1}{cmr}{m}{n}\selectfont\color{color_29791}:}
\put(185.608,-150.809){\fontsize{15.987}{1}\usefont{T1}{cmr}{m}{n}\selectfont\color{color_29791} }
\put(190.002,-150.809){\fontsize{15.987}{1}\usefont{T1}{cmr}{m}{n}\selectfont\color{color_167934}h}
\put(198.2993,-150.809){\fontsize{15.987}{1}\usefont{T1}{cmr}{m}{n}\selectfont\color{color_167934}tt}
\put(208.3871,-150.809){\fontsize{15.987}{1}\usefont{T1}{cmr}{m}{n}\selectfont\color{color_167934}p:}
\put(221.0647,-150.809){\fontsize{15.987}{1}\usefont{T1}{cmr}{m}{n}\selectfont\color{color_167934}/}
\put(227.1558,-150.809){\fontsize{15.987}{1}\usefont{T1}{cmr}{m}{n}\selectfont\color{color_167934}/}
\put(233.2468,-150.809){\fontsize{15.987}{1}\usefont{T1}{cmr}{m}{n}\selectfont\color{color_167934}i}
\put(237.0358,-150.809){\fontsize{15.987}{1}\usefont{T1}{cmr}{m}{n}\selectfont\color{color_167934}c}
\put(243.7343,-150.809){\fontsize{15.987}{1}\usefont{T1}{cmr}{m}{n}\selectfont\color{color_167934}t}
\put(248.9301,-150.809){\fontsize{15.987}{1}\usefont{T1}{cmr}{m}{n}\selectfont\color{color_167934}do}
\put(265.6205,-150.809){\fontsize{15.987}{1}\usefont{T1}{cmr}{m}{n}\selectfont\color{color_167934}c}
\put(272.399,-150.809){\fontsize{15.987}{1}\usefont{T1}{cmr}{m}{n}\selectfont\color{color_167934}s}
\put(278.586,-150.809){\fontsize{15.987}{1}\usefont{T1}{cmr}{m}{n}\selectfont\color{color_167934}/}
\put(284.677,-150.809){\fontsize{15.987}{1}\usefont{T1}{cmr}{m}{n}\selectfont\color{color_167934}w}
\put(296.0598,-150.809){\fontsize{15.987}{1}\usefont{T1}{cmr}{m}{n}\selectfont\color{color_167934}i}
\put(299.8487,-150.809){\fontsize{15.987}{1}\usefont{T1}{cmr}{m}{n}\selectfont\color{color_167934}k}
\put(307.0428,-150.809){\fontsize{15.987}{1}\usefont{T1}{cmr}{m}{n}\selectfont\color{color_167934}i}
\put(310.7198,-150.809){\fontsize{15.987}{1}\usefont{T1}{cmr}{m}{n}\selectfont\color{color_167934}/}
\put(316.8109,-150.809){\fontsize{15.987}{1}\usefont{T1}{cmr}{m}{n}\selectfont\color{color_167934}i}
\put(320.4879,-150.809){\fontsize{15.987}{1}\usefont{T1}{cmr}{m}{n}\selectfont\color{color_167934}nde}
\put(344.964,-150.809){\fontsize{15.987}{1}\usefont{T1}{cmr}{m}{n}\selectfont\color{color_167934}x}
\put(351.8544,-150.809){\fontsize{15.987}{1}\usefont{T1}{cmr}{m}{n}\selectfont\color{color_167934}.}
\put(355.9311,-150.809){\fontsize{15.987}{1}\usefont{T1}{cmr}{m}{n}\selectfont\color{color_167934}php?}
\put(388.5925,-150.809){\fontsize{15.987}{1}\usefont{T1}{cmr}{m}{n}\selectfont\color{color_167934}ti}
\put(397.4813,-150.809){\fontsize{15.987}{1}\usefont{T1}{cmr}{m}{n}\selectfont\color{color_167934}t}
\put(402.7729,-150.809){\fontsize{15.987}{1}\usefont{T1}{cmr}{m}{n}\selectfont\color{color_167934}l}
\put(406.5619,-150.809){\fontsize{15.987}{1}\usefont{T1}{cmr}{m}{n}\selectfont\color{color_167934}e}
\put(414.4594,-150.809){\fontsize{15.987}{1}\usefont{T1}{cmr}{m}{n}\selectfont\color{color_167934}=}
\put(422.357,-150.809){\fontsize{15.987}{1}\usefont{T1}{cmr}{m}{n}\selectfont\color{color_167934}H}
\put(432.3329,-150.809){\fontsize{15.987}{1}\usefont{T1}{cmr}{m}{n}\selectfont\color{color_167934}P}
\put(440.6142,-150.809){\fontsize{15.987}{1}\usefont{T1}{cmr}{m}{n}\selectfont\color{color_167934}C}
\put(449.1033,-150.809){\fontsize{15.987}{1}\usefont{T1}{cmr}{m}{n}\selectfont\color{color_167934}:}
\put(453.3878,-150.809){\fontsize{15.987}{1}\usefont{T1}{cmr}{m}{n}\selectfont\color{color_167934}U}
\put(463.5715,-150.809){\fontsize{15.987}{1}\usefont{T1}{cmr}{m}{n}\selectfont\color{color_167934}s}
\put(469.7585,-150.809){\fontsize{15.987}{1}\usefont{T1}{cmr}{m}{n}\selectfont\color{color_167934}e}
\put(477.6561,-150.809){\fontsize{15.987}{1}\usefont{T1}{cmr}{m}{n}\selectfont\color{color_167934}r}
\put(483.2356,-150.809){\fontsize{15.987}{1}\usefont{T1}{cmr}{m}{n}\selectfont\color{color_167934}\_}
\put(491.1331,-150.809){\fontsize{15.987}{1}\usefont{T1}{cmr}{m}{n}\selectfont\color{color_167934}Gui}
\put(513.2751,-150.809){\fontsize{15.987}{1}\usefont{T1}{cmr}{m}{n}\selectfont\color{color_167934}de}
\end{picture}
\begin{tikzpicture}[overlay]
\path(0pt,0pt);
\draw[color_167934,line width=0.992pt]
(190.002pt, -152.396pt) -- (530.103pt, -152.396pt)
;
\end{tikzpicture}
\begin{picture}(-5,0)(2.5,0)
\put(48.213,-218.614){\fontsize{15.987}{1}\usefont{T1}{cmr}{m}{n}\selectfont\color{color_29791}•}
\put(62.187,-219.493){\fontsize{15.987}{1}\usefont{T1}{cmr}{m}{n}\selectfont\color{color_29791}S}
\put(72.86631,-219.493){\fontsize{15.987}{1}\usefont{T1}{cmr}{m}{n}\selectfont\color{color_29791}er}
\put(86.95086,-219.493){\fontsize{15.987}{1}\usefont{T1}{cmr}{m}{n}\selectfont\color{color_29791}v}
\put(95.02429,-219.493){\fontsize{15.987}{1}\usefont{T1}{cmr}{m}{n}\selectfont\color{color_29791}i}
\put(98.60538,-219.493){\fontsize{15.987}{1}\usefont{T1}{cmr}{m}{n}\selectfont\color{color_29791}ce }
\put(119.8841,-219.493){\fontsize{15.987}{1}\usefont{T1}{cmr}{m}{n}\selectfont\color{color_29791}D}
\put(131.4587,-219.493){\fontsize{15.987}{1}\usefont{T1}{cmr}{m}{n}\selectfont\color{color_29791}esk }
\put(160.8108,-219.493){\fontsize{15.987}{1}\usefont{T1}{cmr}{m}{n}\selectfont\color{color_29791}S}
\put(171.4901,-219.493){\fontsize{15.987}{1}\usefont{T1}{cmr}{m}{n}\selectfont\color{color_29791}o}
\put(180.283,-219.493){\fontsize{15.987}{1}\usefont{T1}{cmr}{m}{n}\selectfont\color{color_29791}l}
\put(183.864,-219.493){\fontsize{15.987}{1}\usefont{T1}{cmr}{m}{n}\selectfont\color{color_29791}ut}
\put(197.2292,-219.493){\fontsize{15.987}{1}\usefont{T1}{cmr}{m}{n}\selectfont\color{color_29791}i}
\put(200.8103,-219.493){\fontsize{15.987}{1}\usefont{T1}{cmr}{m}{n}\selectfont\color{color_29791}o}
\put(209.6031,-219.493){\fontsize{15.987}{1}\usefont{T1}{cmr}{m}{n}\selectfont\color{color_29791}ns}
\put(48.213,-252.913){\fontsize{15.987}{1}\usefont{T1}{cmr}{m}{n}\selectfont\color{color_29791}•}
\put(62.187,-253.906){\fontsize{15.987}{1}\usefont{T1}{cmr}{m}{n}\selectfont\color{color_29791}Y}
\put(71.68327,-253.906){\fontsize{15.987}{1}\usefont{T1}{cmr}{m}{n}\selectfont\color{color_29791}am}
\put(93.85724,-253.906){\fontsize{15.987}{1}\usefont{T1}{cmr}{m}{n}\selectfont\color{color_29791}m}
\put(107.1424,-253.906){\fontsize{15.987}{1}\usefont{T1}{cmr}{m}{n}\selectfont\color{color_29791}er}
\put(48.213,-287.298){\fontsize{15.987}{1}\usefont{T1}{cmr}{m}{n}\selectfont\color{color_29791}•}
\put(62.187,-288.205){\fontsize{15.987}{1}\usefont{T1}{cmr}{m}{n}\selectfont\color{color_29791}E}
\put(72.86631,-288.205){\fontsize{15.987}{1}\usefont{T1}{cmr}{m}{n}\selectfont\color{color_29791}m}
\put(86.1515,-288.205){\fontsize{15.987}{1}\usefont{T1}{cmr}{m}{n}\selectfont\color{color_29791}ai}
\put(98.62137,-288.205){\fontsize{15.987}{1}\usefont{T1}{cmr}{m}{n}\selectfont\color{color_29791}l}
\put(102.2025,-288.205){\fontsize{15.987}{1}\usefont{T1}{cmr}{m}{n}\selectfont\color{color_29791} }
\put(106.5989,-288.205){\fontsize{15.987}{1}\usefont{T1}{cmr}{m}{n}\selectfont\color{color_29791}Li}
\put(119.0687,-288.205){\fontsize{15.987}{1}\usefont{T1}{cmr}{m}{n}\selectfont\color{color_29791}st}
\end{picture}
\newpage
\begin{tikzpicture}[overlay]
\path(0pt,0pt);
\filldraw[color_283006][even odd rule]
(-15pt, 10pt) -- (704.972pt, 10pt)
 -- (704.972pt, 10pt)
 -- (704.972pt, -529.972pt)
 -- (704.972pt, -529.972pt)
 -- (-15pt, -529.972pt)
 -- (-15pt, -529.972pt)
 -- (-15pt, 10pt) -- cycle
;
\end{tikzpicture}
\begin{picture}(-5,0)(2.5,0)
\put(-15,-529.972){\includegraphics[width=720pt,height=540pt]{latexImage_340e7acc1534b9ba351a92c1a1081d49.png}}
\put(239.296,-207.814){\fontsize{28.006}{1}\usefont{T1}{cmr}{m}{n}\selectfont\color{color_29791}A}
\put(257.892,-207.814){\fontsize{28.006}{1}\usefont{T1}{cmr}{m}{n}\selectfont\color{color_29791}ny}
\put(287.5504,-207.814){\fontsize{28.006}{1}\usefont{T1}{cmr}{m}{n}\selectfont\color{color_29791} }
\put(295.336,-207.814){\fontsize{28.006}{1}\usefont{T1}{cmr}{m}{n}\selectfont\color{color_29791}Q}
\put(317.0407,-207.814){\fontsize{28.006}{1}\usefont{T1}{cmr}{m}{n}\selectfont\color{color_29791}uest}
\put(369.972,-207.814){\fontsize{28.006}{1}\usefont{T1}{cmr}{m}{n}\selectfont\color{color_29791}i}
\put(376.2454,-207.814){\fontsize{28.006}{1}\usefont{T1}{cmr}{m}{n}\selectfont\color{color_29791}ons? }
\end{picture}
\newpage
\begin{tikzpicture}[overlay]
\path(0pt,0pt);
\filldraw[color_283006][even odd rule]
(-15pt, 10pt) -- (704.972pt, 10pt)
 -- (704.972pt, 10pt)
 -- (704.972pt, -529.972pt)
 -- (704.972pt, -529.972pt)
 -- (-15pt, -529.972pt)
 -- (-15pt, -529.972pt)
 -- (-15pt, 10pt) -- cycle
;
\filldraw[color_283006][even odd rule]
(-15pt, 10pt) -- (705pt, 10pt)
 -- (705pt, 10pt)
 -- (705pt, -530pt)
 -- (705pt, -530pt)
 -- (-15pt, -530pt)
 -- (-15pt, -530pt)
 -- (-15pt, 10pt) -- cycle
;
\end{tikzpicture}
\begin{picture}(-5,0)(2.5,0)
\put(41.098,-74.189){\fontsize{43.994}{1}\usefont{T1}{cmr}{m}{n}\selectfont\color{color_29791}T}
\put(62.56707,-74.189){\fontsize{43.994}{1}\usefont{T1}{cmr}{m}{n}\selectfont\color{color_29791}he}
\put(107.6169,-74.189){\fontsize{43.994}{1}\usefont{T1}{cmr}{m}{n}\selectfont\color{color_29791} F}
\put(137.1369,-74.189){\fontsize{43.994}{1}\usefont{T1}{cmr}{m}{n}\selectfont\color{color_29791}o}
\put(160.4097,-74.189){\fontsize{43.994}{1}\usefont{T1}{cmr}{m}{n}\selectfont\color{color_29791}llo}
\put(203.5678,-74.189){\fontsize{43.994}{1}\usefont{T1}{cmr}{m}{n}\selectfont\color{color_29791}w}
\put(235.0235,-74.189){\fontsize{43.994}{1}\usefont{T1}{cmr}{m}{n}\selectfont\color{color_29791}i}
\put(245.0102,-74.189){\fontsize{43.994}{1}\usefont{T1}{cmr}{m}{n}\selectfont\color{color_29791}n}
\put(268.151,-74.189){\fontsize{43.994}{1}\usefont{T1}{cmr}{m}{n}\selectfont\color{color_29791}g}
\put(288.9162,-74.189){\fontsize{43.994}{1}\usefont{T1}{cmr}{m}{n}\selectfont\color{color_29791} }
\put(298.8148,-74.189){\fontsize{43.994}{1}\usefont{T1}{cmr}{m}{n}\selectfont\color{color_29791}A}
\put(324.1994,-74.189){\fontsize{43.994}{1}\usefont{T1}{cmr}{m}{n}\selectfont\color{color_29791}r}
\put(338.9814,-74.189){\fontsize{43.994}{1}\usefont{T1}{cmr}{m}{n}\selectfont\color{color_29791}e}
\put(360.9344,-74.189){\fontsize{43.994}{1}\usefont{T1}{cmr}{m}{n}\selectfont\color{color_29791} }
\put(370.833,-74.189){\fontsize{43.994}{1}\usefont{T1}{cmr}{m}{n}\selectfont\color{color_29791}S}
\put(391.0703,-74.189){\fontsize{43.994}{1}\usefont{T1}{cmr}{m}{n}\selectfont\color{color_29791}lide}
\put(456.2694,-74.189){\fontsize{43.994}{1}\usefont{T1}{cmr}{m}{n}\selectfont\color{color_29791} T}
\put(483.7216,-74.189){\fontsize{43.994}{1}\usefont{T1}{cmr}{m}{n}\selectfont\color{color_29791}em}
\put(540.7379,-74.189){\fontsize{43.994}{1}\usefont{T1}{cmr}{m}{n}\selectfont\color{color_29791}p}
\put(563.8787,-74.189){\fontsize{43.994}{1}\usefont{T1}{cmr}{m}{n}\selectfont\color{color_29791}l}
\put(573.8654,-74.189){\fontsize{43.994}{1}\usefont{T1}{cmr}{m}{n}\selectfont\color{color_29791}a}
\put(594.5425,-74.189){\fontsize{43.994}{1}\usefont{T1}{cmr}{m}{n}\selectfont\color{color_29791}t}
\put(608.8406,-74.189){\fontsize{43.994}{1}\usefont{T1}{cmr}{m}{n}\selectfont\color{color_29791}es}
\end{picture}
\newpage
\begin{tikzpicture}[overlay]
\path(0pt,0pt);
\filldraw[color_283006][even odd rule]
(-15pt, 10pt) -- (704.972pt, 10pt)
 -- (704.972pt, 10pt)
 -- (704.972pt, -529.972pt)
 -- (704.972pt, -529.972pt)
 -- (-15pt, -529.972pt)
 -- (-15pt, -529.972pt)
 -- (-15pt, 10pt) -- cycle
;
\end{tikzpicture}
\begin{picture}(-5,0)(2.5,0)
\put(-15,-529.972){\includegraphics[width=720pt,height=540pt]{latexImage_340e7acc1534b9ba351a92c1a1081d49.png}}
\put(48.411,-207.814){\fontsize{28.006}{1}\usefont{T1}{cmr}{m}{n}\selectfont\color{color_29791}T}
\put(64.51445,-207.814){\fontsize{28.006}{1}\usefont{T1}{cmr}{m}{n}\selectfont\color{color_29791}i}
\put(70.61976,-207.814){\fontsize{28.006}{1}\usefont{T1}{cmr}{m}{n}\selectfont\color{color_29791}t}
\put(78.40543,-207.814){\fontsize{28.006}{1}\usefont{T1}{cmr}{m}{n}\selectfont\color{color_29791}l}
\put(84.67877,-207.814){\fontsize{28.006}{1}\usefont{T1}{cmr}{m}{n}\selectfont\color{color_29791}e}
\end{picture}
\newpage
\begin{tikzpicture}[overlay]
\path(0pt,0pt);
\filldraw[color_283006][even odd rule]
(-15pt, 10pt) -- (704.972pt, 10pt)
 -- (704.972pt, 10pt)
 -- (704.972pt, -529.972pt)
 -- (704.972pt, -529.972pt)
 -- (-15pt, -529.972pt)
 -- (-15pt, -529.972pt)
 -- (-15pt, 10pt) -- cycle
;
\end{tikzpicture}
\begin{picture}(-5,0)(2.5,0)
\put(-15,-529.972){\includegraphics[width=720pt,height=540pt]{latexImage_70f318a8f2de3f58eebb26eed4f2123e.png}}
\put(48.411,-60.69601){\fontsize{21.997}{1}\usefont{T1}{cmr}{b}{n}\selectfont\color{color_29791}T}
\put(61.41122,-60.69601){\fontsize{21.997}{1}\usefont{T1}{cmr}{b}{n}\selectfont\color{color_29791}it}
\put(74.8074,-60.69601){\fontsize{21.997}{1}\usefont{T1}{cmr}{b}{n}\selectfont\color{color_29791}l}
\put(80.98856,-60.69601){\fontsize{21.997}{1}\usefont{T1}{cmr}{b}{n}\selectfont\color{color_29791}e}
\end{picture}
\newpage
\begin{tikzpicture}[overlay]
\path(0pt,0pt);
\filldraw[color_283006][even odd rule]
(-15pt, 10pt) -- (704.972pt, 10pt)
 -- (704.972pt, 10pt)
 -- (704.972pt, -529.972pt)
 -- (704.972pt, -529.972pt)
 -- (-15pt, -529.972pt)
 -- (-15pt, -529.972pt)
 -- (-15pt, 10pt) -- cycle
;
\end{tikzpicture}
\begin{picture}(-5,0)(2.5,0)
\put(-15,-529.972){\includegraphics[width=720pt,height=540pt]{latexImage_bdabda61317014bb20dc1cc743a20acb.png}}
\put(48.411,-60.69601){\fontsize{21.997}{1}\usefont{T1}{cmr}{b}{n}\selectfont\color{color_29791}T}
\put(61.41122,-60.69601){\fontsize{21.997}{1}\usefont{T1}{cmr}{b}{n}\selectfont\color{color_29791}it}
\put(74.8074,-60.69601){\fontsize{21.997}{1}\usefont{T1}{cmr}{b}{n}\selectfont\color{color_29791}l}
\put(80.98856,-60.69601){\fontsize{21.997}{1}\usefont{T1}{cmr}{b}{n}\selectfont\color{color_29791}e}
\end{picture}
\newpage
\begin{tikzpicture}[overlay]
\path(0pt,0pt);
\filldraw[color_283006][even odd rule]
(-15pt, 10pt) -- (704.972pt, 10pt)
 -- (704.972pt, 10pt)
 -- (704.972pt, -529.972pt)
 -- (704.972pt, -529.972pt)
 -- (-15pt, -529.972pt)
 -- (-15pt, -529.972pt)
 -- (-15pt, 10pt) -- cycle
;
\end{tikzpicture}
\begin{picture}(-5,0)(2.5,0)
\put(-15,-529.972){\includegraphics[width=720pt,height=540pt]{latexImage_1aef5d9a345e7259c360833b5fc6008f.png}}
\put(48.411,-60.69601){\fontsize{21.997}{1}\usefont{T1}{cmr}{b}{n}\selectfont\color{color_29791}T}
\put(61.41122,-60.69601){\fontsize{21.997}{1}\usefont{T1}{cmr}{b}{n}\selectfont\color{color_29791}it}
\put(74.8074,-60.69601){\fontsize{21.997}{1}\usefont{T1}{cmr}{b}{n}\selectfont\color{color_29791}l}
\put(80.98856,-60.69601){\fontsize{21.997}{1}\usefont{T1}{cmr}{b}{n}\selectfont\color{color_29791}e}
\end{picture}
\newpage
\begin{tikzpicture}[overlay]
\path(0pt,0pt);
\filldraw[color_283006][even odd rule]
(-15pt, 10pt) -- (704.972pt, 10pt)
 -- (704.972pt, 10pt)
 -- (704.972pt, -529.972pt)
 -- (704.972pt, -529.972pt)
 -- (-15pt, -529.972pt)
 -- (-15pt, -529.972pt)
 -- (-15pt, 10pt) -- cycle
;
\end{tikzpicture}
\begin{picture}(-5,0)(2.5,0)
\put(-15,-529.972){\includegraphics[width=720pt,height=540pt]{latexImage_15b9e6acda8be87c69673355bc94d374.png}}
\put(48.411,-60.69601){\fontsize{21.997}{1}\usefont{T1}{cmr}{b}{n}\selectfont\color{color_29791}T}
\put(61.41122,-60.69601){\fontsize{21.997}{1}\usefont{T1}{cmr}{b}{n}\selectfont\color{color_29791}it}
\put(74.8074,-60.69601){\fontsize{21.997}{1}\usefont{T1}{cmr}{b}{n}\selectfont\color{color_29791}l}
\put(80.98856,-60.69601){\fontsize{21.997}{1}\usefont{T1}{cmr}{b}{n}\selectfont\color{color_29791}e}
\end{picture}
\newpage
\begin{tikzpicture}[overlay]
\path(0pt,0pt);
\filldraw[color_283006][even odd rule]
(-15pt, 10pt) -- (704.972pt, 10pt)
 -- (704.972pt, 10pt)
 -- (704.972pt, -529.972pt)
 -- (704.972pt, -529.972pt)
 -- (-15pt, -529.972pt)
 -- (-15pt, -529.972pt)
 -- (-15pt, 10pt) -- cycle
;
\end{tikzpicture}
\begin{picture}(-5,0)(2.5,0)
\put(-15,-529.972){\includegraphics[width=720pt,height=540pt]{latexImage_d5964f0fe7de530ea43707adcb11ad51.png}}
\put(48.411,-60.69601){\fontsize{21.997}{1}\usefont{T1}{cmr}{b}{n}\selectfont\color{color_29791}T}
\put(61.41122,-60.69601){\fontsize{21.997}{1}\usefont{T1}{cmr}{b}{n}\selectfont\color{color_29791}it}
\put(74.8074,-60.69601){\fontsize{21.997}{1}\usefont{T1}{cmr}{b}{n}\selectfont\color{color_29791}l}
\put(80.98856,-60.69601){\fontsize{21.997}{1}\usefont{T1}{cmr}{b}{n}\selectfont\color{color_29791}e}
\end{picture}
\newpage
\begin{tikzpicture}[overlay]
\path(0pt,0pt);
\filldraw[color_283006][even odd rule]
(-15pt, 10pt) -- (704.972pt, 10pt)
 -- (704.972pt, 10pt)
 -- (704.972pt, -529.972pt)
 -- (704.972pt, -529.972pt)
 -- (-15pt, -529.972pt)
 -- (-15pt, -529.972pt)
 -- (-15pt, 10pt) -- cycle
;
\end{tikzpicture}
\begin{picture}(-5,0)(2.5,0)
\put(-15,-529.972){\includegraphics[width=720pt,height=540pt]{latexImage_2828171a63492229494bbe96f89381b1.png}}
\put(48.411,-60.69601){\fontsize{21.997}{1}\usefont{T1}{cmr}{b}{n}\selectfont\color{color_29791}T}
\put(61.41122,-60.69601){\fontsize{21.997}{1}\usefont{T1}{cmr}{b}{n}\selectfont\color{color_29791}it}
\put(74.8074,-60.69601){\fontsize{21.997}{1}\usefont{T1}{cmr}{b}{n}\selectfont\color{color_29791}l}
\put(80.98856,-60.69601){\fontsize{21.997}{1}\usefont{T1}{cmr}{b}{n}\selectfont\color{color_29791}e}
\end{picture}
\newpage
\begin{tikzpicture}[overlay]
\path(0pt,0pt);
\filldraw[color_283006][even odd rule]
(-15pt, 10pt) -- (704.972pt, 10pt)
 -- (704.972pt, 10pt)
 -- (704.972pt, -529.972pt)
 -- (704.972pt, -529.972pt)
 -- (-15pt, -529.972pt)
 -- (-15pt, -529.972pt)
 -- (-15pt, 10pt) -- cycle
;
\end{tikzpicture}
\begin{picture}(-5,0)(2.5,0)
\put(-15,-529.972){\includegraphics[width=720pt,height=540pt]{latexImage_43fb9acc39826b86b6493904abaeacaa.png}}
\put(48.411,-60.69601){\fontsize{21.997}{1}\usefont{T1}{cmr}{b}{n}\selectfont\color{color_29791}T}
\put(61.41122,-60.69601){\fontsize{21.997}{1}\usefont{T1}{cmr}{b}{n}\selectfont\color{color_29791}it}
\put(74.8074,-60.69601){\fontsize{21.997}{1}\usefont{T1}{cmr}{b}{n}\selectfont\color{color_29791}l}
\put(80.98856,-60.69601){\fontsize{21.997}{1}\usefont{T1}{cmr}{b}{n}\selectfont\color{color_29791}e}
\end{picture}
\newpage
\begin{tikzpicture}[overlay]
\path(0pt,0pt);
\filldraw[color_283006][even odd rule]
(-15pt, 10pt) -- (704.972pt, 10pt)
 -- (704.972pt, 10pt)
 -- (704.972pt, -529.972pt)
 -- (704.972pt, -529.972pt)
 -- (-15pt, -529.972pt)
 -- (-15pt, -529.972pt)
 -- (-15pt, 10pt) -- cycle
;
\end{tikzpicture}
\begin{picture}(-5,0)(2.5,0)
\put(-15,-529.972){\includegraphics[width=720pt,height=540pt]{latexImage_3f3a920d1f055ee96cf3917d47957291.png}}
\put(48.411,-60.69601){\fontsize{21.997}{1}\usefont{T1}{cmr}{b}{n}\selectfont\color{color_29791}T}
\put(61.41122,-60.69601){\fontsize{21.997}{1}\usefont{T1}{cmr}{b}{n}\selectfont\color{color_29791}it}
\put(74.8074,-60.69601){\fontsize{21.997}{1}\usefont{T1}{cmr}{b}{n}\selectfont\color{color_29791}l}
\put(80.98856,-60.69601){\fontsize{21.997}{1}\usefont{T1}{cmr}{b}{n}\selectfont\color{color_29791}e}
\end{picture}
\newpage
\begin{tikzpicture}[overlay]
\path(0pt,0pt);
\filldraw[color_283006][even odd rule]
(-15pt, 10pt) -- (704.972pt, 10pt)
 -- (704.972pt, 10pt)
 -- (704.972pt, -529.972pt)
 -- (704.972pt, -529.972pt)
 -- (-15pt, -529.972pt)
 -- (-15pt, -529.972pt)
 -- (-15pt, 10pt) -- cycle
;
\end{tikzpicture}
\begin{picture}(-5,0)(2.5,0)
\put(-15,-529.972){\includegraphics[width=720pt,height=540pt]{latexImage_6f17c9ae575e6e3412bad9150a17b852.png}}
\put(48.411,-60.69601){\fontsize{21.997}{1}\usefont{T1}{cmr}{b}{n}\selectfont\color{color_29791}T}
\put(61.41122,-60.69601){\fontsize{21.997}{1}\usefont{T1}{cmr}{b}{n}\selectfont\color{color_29791}it}
\put(74.8074,-60.69601){\fontsize{21.997}{1}\usefont{T1}{cmr}{b}{n}\selectfont\color{color_29791}l}
\put(80.98856,-60.69601){\fontsize{21.997}{1}\usefont{T1}{cmr}{b}{n}\selectfont\color{color_29791}e}
\end{picture}
\end{document}

\end{document}
